% !TeX program = pdflatex
% Tom IV: Statystyka matematyczna, rozdzialy 1--32 (po audycie 29--30).
% Samodzielny skrypt. Czesc klasyczna i wnioskowanie bayesowskie.
\documentclass[11pt,a4paper,oneside]{report}
\usepackage[T1]{fontenc}
\usepackage[utf8]{inputenc}
\usepackage[polish]{babel}
\usepackage{lmodern}
\usepackage{amsmath,amssymb,amsthm,mathtools}
\usepackage{microtype}
\usepackage{geometry}
\usepackage{enumitem}
\usepackage{booktabs,array,longtable}
\usepackage{xcolor}
\usepackage{needspace}
\usepackage{listings}
\definecolor{codebg}{RGB}{248,249,251}
\definecolor{codeborder}{RGB}{180,191,203}
\definecolor{codecomment}{RGB}{61,113,79}
\definecolor{codekeyword}{RGB}{41,78,132}
\lstdefinelanguage{Julia}{
  morekeywords={using,import,begin,end,if,else,elseif,for,while,return,function,
    struct,mutable,where,let,do,try,catch,finally,const,local,global,println,print},
  sensitive=true,morecomment=[l]{\#},morestring=[b]",
}
\lstdefinestyle{skrypt}{
 backgroundcolor=\color{codebg},frame=single,rulecolor=\color{codeborder},
 framesep=5pt,frameround=tttt,
 basicstyle=\ttfamily\footnotesize,keywordstyle=\color{codekeyword}\bfseries,
 commentstyle=\color{codecomment}\itshape,stringstyle=\color{purple!65!black},
 numbers=left,numberstyle=\tiny\color{gray!75!black},numbersep=8pt,
 stepnumber=1,breaklines=true,breakatwhitespace=true,
 columns=fullflexible,keepspaces=true,showstringspaces=false,
 linewidth=.90\textwidth,
 captionpos=b,aboveskip=11pt,belowskip=12pt,
 xleftmargin=12pt,xrightmargin=2pt,
}
\lstset{style=skrypt}
\usepackage{tikz}
\usetikzlibrary{arrows.meta,positioning,calc,shapes.geometric,fit,backgrounds}
\usepackage[unicode]{hyperref}
% Odsylacze miedzytomowe; zbudowac wszystkie .aux i przechowywac PDF-y obok siebie.
\usepackage{xr-hyper}
\externaldocument[I-]{rachunek_prawdopodobienstwa_tom_I}[rachunek_prawdopodobienstwa_tom_I.pdf]
\externaldocument[II-]{rachunek_prawdopodobienstwa_tom_II}[rachunek_prawdopodobienstwa_tom_II.pdf]
\externaldocument[III-]{rachunek_prawdopodobienstwa_tom_III}[rachunek_prawdopodobienstwa_tom_III.pdf]

\geometry{margin=2.55cm}
\hypersetup{colorlinks=true,linkcolor=blue!45!black,urlcolor=blue!45!black}
\setlength{\parindent}{0pt}
\setlength{\parskip}{4pt}
\setlength{\emergencystretch}{2em}
\setlist[enumerate]{topsep=4pt,itemsep=2pt}
\setlist[itemize]{topsep=4pt,itemsep=2pt}
\numberwithin{equation}{chapter}
\theoremstyle{plain}
\newtheorem{twierdzenie}{Twierdzenie}[chapter]
\newtheorem{lemat}[twierdzenie]{Lemat}
\newtheorem{wniosek}[twierdzenie]{Wniosek}
\newtheorem{stwierdzenie}[twierdzenie]{Stwierdzenie}
\theoremstyle{definition}
\newtheorem{definicja}{Definicja}[chapter]
\newtheorem{przyklad}{Przykład}[chapter]
\newtheorem{zadanie}{Zadanie rozwiązane}[chapter]
\theoremstyle{remark}
\newtheorem{uwaga}{Uwaga}[chapter]
\renewcommand{\proofname}{Dowód}
\newcommand{\R}{\mathbb R}
\newcommand{\N}{\mathbb N}
\newcommand{\Prob}{\mathbb P}
\newcommand{\E}{\mathbb E}
\newcommand{\Var}{\operatorname{Var}}
\newcommand{\Cov}{\operatorname{Cov}}
\newcommand{\Law}{\operatorname{Law}}
\newcommand{\calF}{\mathcal F}
\newcommand{\calA}{\mathcal A}
\newcommand{\calB}{\mathcal B}
\newcommand{\calP}{\mathcal P}
\newcommand{\ind}[1]{\mathbf 1_{#1}}
\newcommand{\dd}{\mathrm d}
\newcommand{\norm}[1]{\lVert#1\rVert}
\newcommand{\abs}[1]{\lvert#1\rvert}
\newcommand{\Normal}{\mathcal N}
\newcommand{\iid}{\mathrel{\overset{\mathrm{iid}}\sim}}
\DeclareMathOperator{\Bin}{Bin}
\DeclareMathOperator{\Pois}{Pois}
\DeclareMathOperator{\Med}{Med}
\DeclareMathOperator{\MSE}{MSE}
\DeclareMathOperator{\supp}{supp}
\DeclareMathOperator{\diag}{diag}
\renewcommand{\lstlistingname}{Kod}
\renewcommand{\lstlistlistingname}{Spis kodów}
\begin{document}
\hypersetup{pageanchor=false}
\begin{titlepage}
\centering\vspace*{3.2cm}
{\Large Podstawy teorii prawdopodobieństwa i statystyki\par}
\vspace{1.1cm}
{\Huge\bfseries Tom IV\\[0.25cm]Statystyka matematyczna\par}
\vspace{0.8cm}
{\LARGE Modele, estymacja, testowanie i decyzje\par}
\vspace{0.7cm}
{\Large Rozdziały 1--32\par}
\vfill
{\large Teoria, dowody, intuicje, dobór metod i przykłady praktyczne\par}
\vspace{1cm}
\end{titlepage}
\hypersetup{pageanchor=true}
\tableofcontents
\chapter*{Wprowadzenie do tomu IV}
\addcontentsline{toc}{chapter}{Wprowadzenie do tomu IV}
W rachunku prawdopodobieństwa ustalaliśmy model i badaliśmy jego konsekwencje. W statystyce kierunek rozumowania jest w pewnym sensie odwrotny: obserwujemy dane i pytamy, jakie właściwości nieznanego mechanizmu ich powstawania możemy na ich podstawie odtworzyć. Nie jest to jednak dosłowne odwracanie funkcji: różne mechanizmy mogą prowadzić do podobnych lub nawet identycznych danych, a przypadkowość oznacza, że na ogół nie otrzymamy pewnego wniosku.

Przykładowo: zaobserwowaliśmy siedem wadliwych produktów w próbie stu sztuk. Liczba $7/100$ jest łatwa do obliczenia, lecz to dopiero początek. Czy produkty wybrano losowo? Czy obserwacje są niezależne? Czy urządzenie pracowało w tym samym trybie? Czy pytamy o całą produkcję, czy o partię z określonej zmiany? Jak zmienia się niepewność wniosku, gdy obserwujemy siedem wad na dziesięć sztuk? W kolejnych rozdziałach będziemy do tego zadania wracać.

\paragraph{Zasada pracy z metodami statystycznymi.}
Dla każdej nowej metody ustalimy: \emph{cel}, \emph{model i założenia}, \emph{wyprowadzenie}, \emph{sposób obliczenia}, \emph{interpretację}, \emph{alternatywy} oraz \emph{sytuacje, w których metoda zawodzi}. Porównywanie metod na tych samych danych nie jest dodatkiem do teorii, tylko częścią jej rozumienia. W tym tomie będziemy rozwijać zarówno wnioskowanie klasyczne, jak i bayesowskie; osobną, systematyczną teorię bayesowską przedstawimy w dalszych częściach.

\paragraph{Co uznajemy za znane.}
Zakładamy pojęcia przestrzeni probabilistycznej (tom I, rozdział~\ref{I-ch:probability}), niezależności i rozkładu (tom I, rozdziały~\ref{I-ch:conditioning} i~\ref{I-ch:random-variables}), wartości oczekiwanej (tom I, rozdział~\ref{I-ch:expectation}), warunkowej wartości oczekiwanej (tom I, rozdział~\ref{I-ch:condexp}) oraz praw wielkich liczb i centralnego twierdzenia granicznego (tom I, rozdziały~\ref{I-ch:lln} i~\ref{I-ch:clt}). Gdy skorzystamy z bardziej zaawansowanego twierdzenia tego tomu (np. Glivenki--Cantellego), podamy dokładnie jego hipotezy. Statystyka nie wymaga przyjmowania nieznanych parametrów za zmienne losowe: w podejściu klasycznym parametr jest ustalony, a losowość przypisujemy obserwacjom.


\section*{Mapa odsyłaczy do pozostałych tomów}
\begin{itemize}
\item \textbf{Tom I — Rachunek prawdopodobieństwa:} rozdział~\ref{I-ch:conditioning} (warunkowanie), rozdział~\ref{I-ch:condexp} (warunkowa wartość oczekiwana), rozdziały~\ref{I-ch:lln}--\ref{I-ch:clt} (twierdzenia graniczne) i rozdział~\ref{I-ch:markov} (łańcuchy Markowa).
\item \textbf{Tom II — Procesy stochastyczne:} rozdział~\ref{II-ch:procesy} (modele trajektorii), rozdział~\ref{II-ch:poisson-process} (proces Poissona), rozdział~\ref{II-ch:wiener} (ruch Browna) oraz rozdział~\ref{II-ch:quadratic-cont} (wariacja kwadratowa).
\item \textbf{Tom III — Analiza stochastyczna:} rozdział~\ref{III-ch:sde-existence} (SDE), rozdział~\ref{III-ch:sde-properties} (aproksymacja numeryczna), rozdział~\ref{III-ch:levy} (procesy Lévy'ego) i rozdział~\ref{III-ch:path-integrals} (funkcjonały trajektorii).
\end{itemize}
Odsyłacze w tej sekcji oraz w akapitach \emph{Powiązania i odsyłacze} odwołują się do pozostałych tomów. W oddzielnych plikach PDF są zapisane jako odsyłacze \texttt{GoToR}, których obsługa zależy od czytnika i jego ustawień bezpieczeństwa. Aby mieć pewność działania wszystkich odsyłaczy także w przeglądarce internetowej, należy używać wydania zbiorczego \texttt{Cztery\_tomy\_polaczone\_aktywne\_linki.pdf}, gdzie są one zwykłymi odsyłaczami wewnętrznymi. Wydanie zbiorcze przygotowuje dostarczony skrypt \texttt{polacz\_tomy.py}. Odsyłacze lokalne w pojedynczym tomie nie wymagają innych plików.

\chapter{Eksperyment statystyczny i model danych}
\label{ch:eksperyment}
\paragraph{Powiązania i odsyłacze.} Przypomnienie konstrukcji przestrzeni probabilistycznej i zmiennych losowych znajduje się w tomie I, rozdziałach~\ref{I-ch:probability} i~\ref{I-ch:random-variables}. Zależność czasową obserwacji formalizujemy w tomie II, rozdziale~\ref{II-ch:procesy}.

\section{Od problemu praktycznego do modelu probabilistycznego}
Rozważmy urządzenie, które przy każdym użyciu może ulec awarii. Jeżeli interesuje nas \emph{tylko} liczba awarii w $n$ niezależnych próbach o jednakowym prawdopodobieństwie awarii $p$, naturalnym modelem jest rozkład dwumianowy. W praktyce nie wiemy, czy wartość $p$ wynosi $0{,}01$, $0{,}05$ czy $0{,}2$. Rozważamy więc całą rodzinę prawdopodobieństw. Każda z nich opisuje inny, możliwy mechanizm powstawania danych.

\begin{definicja}[Eksperyment statystyczny]\label{def:eksperyment}
Niech $(\mathsf X,\calA)$ będzie przestrzenią mierzalną, $\Theta$ niepustym zbiorem oraz $\{P_\theta:\theta\in\Theta\}$ rodziną miar probabilistycznych na $(\mathsf X,\calA)$. Trójkę
\[
\mathcal E=(\mathsf X,\calA,\{P_\theta:\theta\in\Theta\})
\]
nazywamy \emph{eksperymentem statystycznym}. Zbiór $\Theta$ jest \emph{przestrzenią parametrów}, a $(\mathsf X,\calA)$ \emph{przestrzenią obserwacji}. W modelu klasycznym nieznany $\theta\in\Theta$ jest ustalony: losowe są obserwacje o prawie $P_\theta$.
\end{definicja}

Przestrzenią obserwacji może być zbiór $\{0,1\}^{100}$, przestrzeń $\R^n$, przestrzeń wszystkich tabel kontyngencji albo przestrzeń trajektorii procesu. Nie wolno utożsamiać parametru z konkretną obserwacją: parametr indeksuje \emph{rozkład}, z którego może powstać każda z wielu prób.

\begin{przyklad}[Produkt sprawny lub wadliwy]\label{ex:bernoulli-exp}
Dla jednej sztuki niech $\mathsf X=\{0,1\}$, gdzie $1$ oznacza wadę. Model Bernoulliego przyjmuje $\Theta=[0,1]$ oraz $P_p(\{1\})=p$, $P_p(\{0\})=1-p$. Dla $n$ niezależnych sztuk model ma postać
\[
\mathsf X_n=\{0,1\}^n,\qquad
P_p^{\otimes n}(\{(x_1,\ldots,x_n)\})=p^{\sum_i x_i}(1-p)^{n-\sum_i x_i}.
\]
Dla $p\in\{0,1\}$ stosujemy naturalną interpretację prawdopodobieństw zdarzeń niemożliwych i pewnych. Statystyka $S=\sum_iX_i$ ma rozkład $\Bin(n,p)$, ale pełna próba to uporządkowany wektor $(X_1,\ldots,X_n)$: przejście do $S$ jest osobną decyzją o redukcji danych, którą uzasadnimy później teorią dostateczności.
\end{przyklad}

\begin{przyklad}[Pomiar ciągły]
Zakładamy, że wynik $X$ pomiaru temperatury ma rozkład $\Normal(\mu,\sigma^2)$, gdzie $\mu\in\R$ i $\sigma>0$ są nieznane. Parametr jest dwuwymiarowy $\theta=(\mu,\sigma)$, a jego przestrzeń wynosi $\Theta=\R\times(0,\infty)$. Jeśli znamy dokładność przyrządu i chcemy estymować wyłącznie $\mu$, możemy uznać $\sigma$ za znaną. Jest to jednak \emph{inny eksperyment statystyczny} i inne założenie, a nie tylko wygodniejszy sposób zapisania tego samego problemu.
\end{przyklad}

\begin{przyklad}[Zliczanie z ekspozycją]\label{ex:exposure}
Niech $Y_j$ oznacza liczbę awarii maszyny w czasie pracy $t_j>0$. W modelu Poissona $Y_j\sim\Pois(\lambda t_j)$, gdzie $\lambda>0$ jest nieznaną intensywnością na jednostkę czasu. Liczb $Y_j$ nie należy bezrefleksyjnie porównywać, gdy $t_j$ są różne. Gdyby dwie maszyny miały po $5$ awarii, lecz pracowały odpowiednio $10$ i $100$ godzin, surowe liczby byłyby identyczne, a obserwowane intensywności dziesięciokrotnie różne. Ten przykład motywuje rolę \emph{ekspozycji} i późniejsze modele z offsetem.
\end{przyklad}

\section{Obserwacja, próba, schemat losowania}
\begin{definicja}[Próba losowa]\label{def:sample}
Niech $(\mathsf S,\mathcal S)$ będzie przestrzenią mierzalną, a $Q_\theta$ miarą probabilistyczną na niej dla każdego $\theta\in\Theta$. Mówimy, że $X_1,\ldots,X_n$ są \emph{próbą prostą} z rozkładu $Q_\theta$, jeśli są niezależne i każda ma rozkład $Q_\theta$. Wówczas prawo wektora $X=(X_1,\ldots,X_n)$ jest iloczynem $Q_\theta^{\otimes n}$.
\end{definicja}

Próba prosta to założenie probabilistyczne, nie automatyczna cecha tabeli danych. Powtarzane pomiary tej samej osoby, kolejne dni działania jednej maszyny czy transakcje pochodzące z jednego silnego zdarzenia rynkowego zwykle nie są niezależnymi obserwacjami. Liczba wierszy tabeli nie musi zatem być liczbą niezależnych jednostek informacji.

\begin{przyklad}[Dwie konstrukcje próby z urny]\label{ex:replacement}
W urnie znajduje się $N$ kul, z czego $K$ jest oznaczonych. Losujemy $n\le N$ kul i zapisujemy liczbę oznaczonych jako $S$. Jeśli po każdym losowaniu zwracamy kulę do urny, to kolejne wskaźniki oznaczenia są niezależne o prawdopodobieństwie $K/N$ i $S\sim\Bin(n,K/N)$. Jeśli losujemy \emph{bez zwracania}, wskaźniki nie są niezależne i
\[
\Prob(S=k)=\frac{\binom Kk\binom{N-K}{n-k}}{\binom Nn},\qquad
\max(0,n-N+K)\le k\le\min(n,K).
\]
\emph{Wyprowadzenie.} Wśród $\binom Nn$ jednakowo prawdopodobnych nieuporządkowanych prób dokładnie $\binom Kk\binom{N-K}{n-k}$ zawiera $k$ kul oznaczonych. Warunków brzegowych nie można pominąć: zapis z symbolami Newtona jest równoważny jedynie przy przyjęciu, że $\binom ab=0$ poza zakresem $0\le b\le a$.
\end{przyklad}

\begin{stwierdzenie}[Zależność przy losowaniu bez zwracania]\label{prop:without-replacement-cov}
Niech $N\ge2$ i $K\in\{0,\ldots,N\}$. Jeśli $X_1,X_2$ są wskaźnikami oznaczenia w pierwszym i drugim losowaniu bez zwracania z opisanej wyżej urny, to
\[
\E X_i=\frac KN=:p,\qquad
\Cov(X_1,X_2)=-\frac{p(1-p)}{N-1}.
\]
\end{stwierdzenie}
\begin{proof}
Warunkowo względem $X_1=1$ w urnie pozostaje $K-1$ oznaczonych kul spośród $N-1$, zatem $\E(X_1X_2)=\Prob(X_1=1,X_2=1)=K(K-1)/(N(N-1))$. Odejmując $p^2=K^2/N^2$ i sprowadzając do wspólnego mianownika, otrzymujemy
$\Cov(X_1,X_2)=-K(N-K)/(N^2(N-1))$, co jest równoważne tezie.
\end{proof}

\section{Jakie dane mamy? Jednostka obserwacji i typ zmiennej}
Przed wyborem wzoru trzeba ustalić, czym jest obserwowana jednostka. \emph{Zmienne jakościowe nominalne} (np. producent, klasa urządzenia) mają kategorie bez naturalnego porządku; \emph{porządkowe} (np. poziom satysfakcji) mają porządek, lecz odległość między kategoriami nie musi mieć znaczenia liczbowego. Zmienne \emph{liczbowe dyskretne} zliczają zdarzenia; \emph{ciągłe} mierzą wielkości. Czas do awarii jest szczególną wielkością: może być \emph{cenzorowany}, gdy dla niektórych urządzeń wiadomo jedynie, że przeżyły do końca obserwacji.

\begin{center}
\begin{tikzpicture}[>=Latex,node distance=6mm and 8mm,every node/.style={font=\small},box/.style={draw,rounded corners,align=center,inner sep=5pt}]
\node[box] (goal) {Pytanie o populację\\i parametr};
\node[box,right=of goal] (design) {Schemat wyboru\\jednostek i pomiaru};
\node[box,right=of design] (data) {Dane $X_1,\ldots,X_n$\\i ich zależność};
\node[box,below=10mm of data] (method) {Statystyka lub metoda\\dopasowana do celu};
\node[box,left=of method] (check) {Diagnostyka założeń\\i niepewność};
\node[box,left=of check] (decision) {Wniosek dla populacji\\z ograniczeniami};
\draw[->,thick] (goal)--(design);
\draw[->,thick] (design)--(data);
\draw[->,thick] (data)--(method);
\draw[->,thick] (method)--(check);
\draw[->,thick] (check)--(decision);
\end{tikzpicture}
\end{center}
\noindent\emph{Rysunek 1.1.} Łańcuch wnioskowania statystycznego. Źle dobranego schematu obserwacji nie naprawi późniejszy wybór bardziej wyrafinowanego testu.

\paragraph{Trzy układy danych.}
Dla \emph{niezależnych grup} każda osoba lub maszyna należy do jednej grupy. Dla \emph{prób sparowanych} ten sam obiekt obserwujemy przed i po interwencji, a właściwą jednostką analizy jest para (lub różnica w parze). Dla \emph{szeregu czasowego} obserwacje są uporządkowane i zależność może dotyczyć wielu opóźnień. W tomie IV wrócimy do tego rozróżnienia przy teście Welcha, teście t dla prób zależnych, ANOVA z powtarzanymi pomiarami i modelach mieszanych.


\begin{twierdzenie}[Ile niezależnej informacji zawierają pomiary w grupach?]\label{thm:cluster-design}
Niech $g,m\ge1$ i niech $X_{ji}$, $j=1,\ldots,g$, $i=1,\ldots,m$,
będą rzeczywistymi kwadratowo całkowalnymi zmiennymi losowymi.
Załóżmy, że wszystkie mają średnią $\mu$ i wariancję $\sigma^2>0$,
że różne grupy są niezależne, a kowariancja dwóch różnych pomiarów
w tej samej grupie jest równa $\rho\sigma^2$ (dla wartości $\rho$
zgodnych z nieujemną określonością macierzy kowariancji).
Dla średniej $\overline X=(gm)^{-1}\sum_{j,i}X_{ji}$ zachodzi
\[
\Var(\overline X)=\frac{\sigma^2}{gm}\bigl(1+(m-1)\rho\bigr).
\]
\end{twierdzenie}
\begin{proof}
Oznaczmy sumę wszystkich obserwacji przez $S$. Rozwijając wariancję,
uwzględniamy $gm$ składników wariancji $\sigma^2$ i
$gm(m-1)$ uporządkowanych par różnych obserwacji z tej samej grupy,
każdą o kowariancji $\rho\sigma^2$. Pary z różnych grup wnoszą zero.
Stąd $\Var(S)=gm\sigma^2+gm(m-1)\rho\sigma^2$.
Dzieląc przez $(gm)^2$, otrzymujemy tezę.
\end{proof}
\begin{przyklad}[Sto pomiarów nie musi oznaczać stu niezależnych obserwacji]
Zmierzono $20$ urządzeń po $5$ razy, razem $100$ wyników.
Jeżeli w modelu poprzedniego twierdzenia $\rho=0{,}6$, wariancja
średniej jest $1+4\cdot0{,}6=3{,}4$ razy większa niż przy
niezależności wszystkich stu wyników. Liczba około $100/3{,}4\approx29{,}4$
opisuje wyłącznie \emph{równoważną liczebność ze względu na wariancję
średniej}; nie oznacza istnienia $29{,}4$ niezależnych obserwacji.
W praktyce należy uwzględnić strukturę grupową lub sprowadzić
analizę do poziomu niezależnych jednostek, jeżeli odpowiada to celowi.
\end{przyklad}

\section{Identyfikowalność: czy parametr da się w ogóle odzyskać?}
\begin{definicja}[Identyfikowalność]\label{def:ident}
Eksperyment statystyczny $\mathcal E=(\mathsf X,\calA,\{P_\theta:\theta\in\Theta\})$ nazywamy \emph{identyfikowalnym}, jeżeli dla dowolnych $\theta_1,\theta_2\in\Theta$ zachodzi implikacja
\[
P_{\theta_1}=P_{\theta_2}\quad\Longrightarrow\quad \theta_1=\theta_2.
\]
Równość po lewej jest równością miar na $\calA$, nie tylko zbieżnością kilku wybranych momentów.
\end{definicja}

\begin{przyklad}[Nieobserwowalny znak parametru]\label{ex:not-identifiable}
Niech $\Theta=\R$ i $X\sim\Normal(\theta^2,1)$. Wówczas dla dowolnego $a>0$ rozkłady $P_a$ i $P_{-a}$ są identyczne. Z dowolnie wielu niezależnych obserwacji można próbować odzyskać $\theta^2$, ale nie znak $\theta$: dane nie rozróżniają dwóch wartości parametru. Nie jest to problem małej liczby obserwacji ani słabego algorytmu. Jest to ograniczenie samego eksperymentu.
\end{przyklad}

\begin{twierdzenie}[Brak zgodności przy braku identyfikowalności]\label{thm:no-consistency}
Niech $\Theta$ będzie przestrzenią metryczną z metryką $d$, a dla każdego $n\ge1$ dane $(X_1,\ldots,X_n)$ mają prawo $P_\theta^{\otimes n}$ na przestrzeni $\mathsf X^n$. Załóżmy, że istnieją różne $\theta_1,\theta_2\in\Theta$ takie, że $P_{\theta_1}=P_{\theta_2}$. Nie istnieje ciąg estymatorów $T_n:\mathsf X^n\to\Theta$, który byłby zgodny według prawdopodobieństwa jednocześnie dla $\theta_1$ i $\theta_2$, tzn. dla obu wartości miałby $d(T_n,\theta_j)\to0$ według $P_{\theta_j}^{\otimes n}$.
\end{twierdzenie}
\begin{proof}
Niech $\varepsilon=d(\theta_1,\theta_2)/3>0$ i zdefiniujmy zdarzenia $A_{n,j}=\{d(T_n,\theta_j)<\varepsilon\}$ dla $j=1,2$. Nierówność trójkąta daje $A_{n,1}\cap A_{n,2}=\varnothing$, ponieważ na ich przecięciu mielibyśmy $d(\theta_1,\theta_2)<2\varepsilon<d(\theta_1,\theta_2)$. Z drugiej strony rozkłady danych dla obu parametrów są takie same dla każdego $n$, gdyż równość miar zachowuje się przy iloczynie. Oznaczmy to wspólne prawo przez $Q_n$. Zgodność w $\theta_1$ oznacza $Q_n(A_{n,1})\to1$, a zgodność w $\theta_2$ oznacza $Q_n(A_{n,2})\to1$. To przeczy $Q_n(A_{n,1})+Q_n(A_{n,2})\le1$.
\end{proof}

\begin{uwaga}
Identyfikowalność nie wystarcza sama w sobie do istnienia estymatora o dobrych własnościach. Jest \emph{warunkiem koniecznym}, lecz do zgodności potrzebujemy również odpowiednio rozróżnialnych rozkładów oraz kontroli zachowania statystyk przy rosnącej próbie.
\end{uwaga}

\section{Model zdominowany, gęstość i dlaczego nie zaczynamy od wiarygodności}
\begin{definicja}[Model zdominowany]
Eksperyment $\mathcal E$ nazywamy \emph{zdominowanym}, jeżeli istnieje miara $\sigma$-skończona $\mu$ na $(\mathsf X,\calA)$ taka, że $P_\theta\ll\mu$ dla każdego $\theta\in\Theta$. Wówczas istnieją pochodne Radona--Nikodyma $p_\theta=\dd P_\theta/\dd\mu$. Dla każdej mierzalnej $A\subseteq\mathsf X$ zachodzi $P_\theta(A)=\int_Ap_\theta(x)\,\dd\mu(x)$.
\end{definicja}

W rodzinie Bernoulliego wybieramy miarę liczącą, a w rodzinie normalnej miarę Lebesgue'a. Dominacja pozwala później zdefiniować funkcję wiarygodności $L(\theta;x)$ poprzez $p_\theta(x)$ dla zaobserwowanego $x$. Trzeba jednak pamiętać, że funkcja wiarygodności \emph{nie jest rozkładem prawdopodobieństwa parametru} w podejściu klasycznym. Rozkłady a priori i a posteriori pojawią się w części bayesowskiej; do tej pory zmienna losowa jest obserwacją, a parametr pozostaje ustalony.

\begin{przyklad}[Parametr wpływający na nośnik]\label{ex:uniform-support}
Dla $\theta>0$ niech $X\sim\operatorname{Unif}(0,\theta)$, czyli $p_\theta(x)=\theta^{-1}\ind{(0,\theta)}(x)$. Obserwacja $x>0$ wyklucza wszystkie $\theta<x$. Nie wolno więc w analizie tego modelu pomijać zależności nośnika od parametru. Później zobaczymy, dlaczego standardowe twierdzenia asymptotyczne o MLE wymagają osobnego sprawdzenia warunków regularności.
\end{przyklad}

\section{Próba reprezentatywna, selekcja i obciążenie}
\paragraph{Dwie losowości.}
W modelu \emph{superpopulacyjnym} wartości badanych jednostek są zmiennymi losowymi o określonym rozkładzie. W modelu \emph{losowania z populacji skończonej} wartości $y_1,\ldots,y_N$ mogą być ustalone, a losowy jest wybór jednostek. Rozróżnienie ma znaczenie przy interpretacji wartości oczekiwanej i przy obliczaniu błędu estymatora.

\begin{przyklad}[Nierówna reprezentacja grup]\label{ex:selection}
Populacja zawiera 90 pracowników z grupy $A$, których średni wynik wynosi $10$, oraz 10 pracowników z grupy $B$, których średni wynik wynosi $50$. Średnia populacyjna to $0{,}9\cdot10+0{,}1\cdot50=14$. Jeśli z każdej grupy weźmiemy po 10 osób, otrzymamy w tych dwóch próbach średnie równe odpowiednio $10$ i $50$ oraz policzymy nieważoną średnią tych średnich grupowych, dostaniemy $30$, a nie $14$. To \emph{nie} oznacza, że średnia arytmetyczna jest złym wzorem; oznacza, że w próbie inna jest struktura grup niż w populacji. Poprawna średnia ważona udziałami populacyjnymi odzyskuje $14$, o ile średnie w grupach są reprezentatywne dla tych grup.
\end{przyklad}

\begin{twierdzenie}[Nieobciążoność estymatora Horvitza--Thompsona]\label{thm:HT}
Niech populacja składa się z $N\ge1$ jednostek o ustalonych wartościach $y_1,\ldots,y_N\in\R$. Niech $I_i\in\{0,1\}$ będzie wskaźnikiem włączenia jednostki $i$ do próby według dowolnego probabilistycznego schematu, dla którego $\pi_i=\E I_i=\Prob(I_i=1)>0$. Estymator sumy populacyjnej
\[
\widehat T_{\mathrm{HT}}=\sum_{i=1}^N\frac{I_i y_i}{\pi_i}
\]
spełnia $\E\widehat T_{\mathrm{HT}}=\sum_{i=1}^Ny_i$. Nie wymaga to niezależności wskaźników $I_i$.
\end{twierdzenie}
\begin{proof}
Każdy składnik jest całkowalny, ponieważ liczba jednostek jest skończona, a $\pi_i>0$. Z liniowości wartości oczekiwanej,
$\E\widehat T_{\mathrm{HT}}=\sum_i(y_i/\pi_i)\E I_i=\sum_i y_i$. W dowodzie nie obliczamy żadnego iloczynu $I_iI_j$, dlatego założenie niezależności nie jest potrzebne. Przy liczeniu \emph{wariancji} trzeba byłoby uwzględnić wspólne prawdopodobieństwa włączenia jednostek.
\end{proof}

\begin{uwaga}[Kiedy nawet ważenie nie rozwiązuje problemu?]
Jeśli nie znamy prawdopodobieństw włączenia lub ich oszacowania są błędne, waga $1/\pi_i$ może nie korygować selekcji. Jeśli dla pewnej jednostki $\pi_i=0$, nie można w ten sposób odzyskać jej wkładu do sumy z dostępnego schematu. Ponadto ekstremalnie duże wagi zwiększają wariancję. Statystyczna poprawność i precyzja to dwa różne zagadnienia.
\end{uwaga}

\section{Obserwacje brakujące, cenzorowanie i mechanizm zbierania danych}
Brak danej nie zawsze oznacza pustą komórkę, którą można bez szkody usunąć. Niech $X$ oznacza wartość, którą chcielibyśmy zaobserwować, $Z$ zmienne pomocnicze dostępne dla każdej jednostki, a $R\in\{0,1\}$ wskaźnik obserwacji $X$. Rozróżniamy trzy klasyczne sytuacje opisowe:
\begin{itemize}
\item \emph{MCAR} (brak całkowicie losowy): $R$ jest niezależny od $(X,Z)$; usuwanie rekordów nie powoduje selekcji względem tych zmiennych, lecz zmniejsza liczebność próby.
\item \emph{MAR} (brak losowy warunkowo): $R$ i $X$ są warunkowo niezależne przy zadanym $Z$; po uwzględnieniu obserwowanych zmiennych $Z$ brak nie zależy już od niezaobserwowanego $X$.
\item \emph{MNAR} (brak zależny od niezaobserwowanych wartości): zależność braków od $X$ pozostaje nawet po warunkowaniu względem $Z$; z danych obserwowanych zwykle nie da się jej jednoznacznie rozpoznać bez dodatkowych założeń.
\end{itemize}
Powyższe sformułowania dotyczą uproszczonego przypadku, w którym $Z$ jest zawsze obserwowana, a brak może dotyczyć tylko $X$. Nie wynika z nich, że zwykłe usuwanie brakujących rekordów przy MAR daje nieobciążone oszacowania: potrzebne jest poprawne uwzględnienie $Z$ i mechanizmu selekcji. Terminologia jest zależna od tego, które zmienne uznajemy za obserwowane. W rzeczywistych problemach z wieloma brakującymi zmiennymi formalne definicje odnoszą się do całej maski braków i części obserwowanej danych.

\begin{przyklad}[Cenzorowanie a brak danych]
W badaniu awaryjności obserwujemy maszynę przez $T=100$ godzin. Jeśli do chwili $T$ nie uległa awarii, nie wiemy, kiedy awaria wystąpi, ale wiemy, że czas życia przekracza $100$. Jest to \emph{prawostronne cenzorowanie}, a nie obserwacja, którą można utożsamiać z czasem awarii równym $100$. Późniejsze metody analizy przeżycia muszą włączyć tę informację do modelu.
\end{przyklad}

\section{Czego nie zapewniają same dane: zależność, mylenie korelacji z przyczynowością i błędna specyfikacja}
Przyjmijmy, że mamy dane o liczbie reklamacji i czasie pracy urządzenia. Obserwowana korelacja może pochodzić z tego, że dłużej pracujące urządzenia częściej ulegają awarii, bez różnicy w ich niezawodności. Wnioskowanie o przyczynach wymaga osobnego modelu mechanizmu interwencji lub kontroli zmiennych zakłócających. Zależność statystyczna nie jest równoznaczna ze związkiem przyczynowym.

\begin{przyklad}[Przybliżony model dwumianowy może zawodzić]
Jeśli wadliwość zależy od partii, to warunkowo na losowym parametrze partii $P$ wyniki mogą być niezależnymi zmiennymi Bernoulliego, ale bez warunkowania dodatnio skorelowane. Przy $X_i\mid P\sim\operatorname{Bernoulli}(P)$ i niezależności warunkowej,
\[
\Cov(X_i,X_j)=\Var(P),\qquad i\neq j.
\]
\emph{Uzasadnienie.} Wzór na kowariancję całkowitą daje $\Cov(X_i,X_j)=\E\Cov(X_i,X_j\mid P)+\Cov(\E[X_i\mid P],\E[X_j\mid P])=0+\Var(P)$. Jeśli $P$ nie jest stałe prawie na pewno, liczba obserwacji nie zapewnia informacji takiej jak w niezależnym modelu $\Bin(n,p)$.
\end{przyklad}

\section{Jak w praktyce wybrać model?}
Poniższa tabela jest \emph{procedurą rozpoznania problemu}, a nie automatycznym doborem testu. W rozdziale 2 zobaczymy odpowiednie statystyki opisowe; w dalszych rozdziałach zbudujemy estymatory, testy i przedziały.

\begin{center}\small
\begin{tabular}{>{\raggedright\arraybackslash}p{0.27\textwidth}>{\raggedright\arraybackslash}p{0.29\textwidth}>{\raggedright\arraybackslash}p{0.32\textwidth}}
\toprule
\textbf{Co obserwujemy?} & \textbf{Model do rozważenia} & \textbf{Co sprawdzić przed użyciem?}\\
\midrule
Wada: tak/nie & Bernoulli, dwumianowy & Losowanie, zależność w partiach, jednakowe $p$\\
Liczba zgłoszeń w czasie & Poisson, ujemny dwumianowy & Ekspozycja, nadmierna dyspersja, sezonowość\\
Wynik pomiaru & Normalny, lognormalny, model odporny & Symetria, ogony, granice skali, wartości odstające\\
Czas do awarii & Wykładniczy, Weibulla, analiza przeżycia & Cenzorowanie, zmienność intensywności awarii\\
Pomiar przed/po & Model sparowany & Powiązanie par i rozkład różnic\\
Wielokrotne pomiary & Model zależności, model mieszany & Identyfikator jednostki, korelacja w czasie\\
\bottomrule
\end{tabular}
\end{center}

\begin{zadanie}[Pierwsza analiza jakości produkcji]\label{zad:quality-case}
W próbie $100$ produktów stwierdzono $7$ wadliwych. Wiadomo, że $80$ produktów pochodzi z jednej zmiany, a pozostałe $20$ z innej. Jakie informacje mamy, co możemy obliczyć i czego jeszcze nie wolno wnioskować?
\par\emph{Rozwiązanie.} Wskaźnik wadliwości próby wynosi $\widehat p=7/100=0{,}07$. Nie jest on automatycznie prawdopodobieństwem wady w całej produkcji. Do tego potrzebujemy określić populację, sposób losowania i zależność między wynikami. Jeśli produkty pobrano losowo i niezależnie z procesu o wspólnym parametrze $p$, liczba wad ma prawo $\Bin(100,p)$; w innym przypadku może być potrzebny model z efektem zmiany lub model hierarchiczny. Bez podania liczby wad na każdej zmianie nie możemy również porównać zmian. Testowanie hipotezy $p=0{,}05$ albo przedział ufności dla $p$ należy do kolejnych rozdziałów, ponieważ wymaga jeszcze zdefiniowania poziomu błędu i metody konstrukcji wniosku.
\end{zadanie}

\begin{zadanie}[Kiedy próba prosta nie jest właściwym modelem?]
Dla każdego z $40$ pacjentów wykonano po $3$ pomiary ciśnienia w kolejnych dniach. Czy dysponujemy próbą prostą o liczności $120$?
\par\emph{Rozwiązanie.} Nie wynika to z opisu. Pomiary tej samej osoby mogą dzielić wspólną nieobserwowaną cechę. Nie wolno więc traktować liczby $120$ jako liczby niezależnych pacjentów. Potrzebujemy identyfikatora osoby i modelu korelacji wewnątrzosobniczej (np. modelu efektów losowych) lub statystyki sprowadzającej dane do poziomu pacjenta, jeśli to odpowiada celowi analizy. Dalsze porównania powtórzonych pomiarów przedstawimy przy modelach mieszanych.
\end{zadanie}

\paragraph{Podsumowanie praktyczne rozdziału 1.}
Statystyka zaczyna się \emph{przed} obliczeniem pierwszej średniej. Trzeba ustalić: (i) populację docelową, (ii) mechanizm pomiaru i selekcji, (iii) przestrzeń możliwych rozkładów, (iv) parametr lub wielkość, którą chcemy poznać, (v) rodzaj zależności między obserwacjami i (vi) możliwe naruszenia założeń. Gdy te elementy nie są określone, nawet poprawny algebraicznie test może odpowiadać na niewłaściwe pytanie.

\chapter{Statystyki i rozkłady z próby}
\label{ch:statystyki}
\paragraph{Powiązania i odsyłacze.} Rozkłady klasyczne i momenty omawia tom I, rozdziały~\ref{I-ch:distributions} i~\ref{I-ch:expectation}; zbieżności i granice próbek — rozdziały~\ref{I-ch:convergence}--\ref{I-ch:clt}.

\section{Po co zastępować próbę kilkoma liczbami?}
Próba $x_1,\ldots,x_n$ może być bardzo długa. Nie zawsze potrzebujemy przechowywać cały wektor, aby udzielić odpowiedzi na ustalone pytanie. Jeśli interesuje nas typowy poziom pomiaru, naturalne są średnia lub mediana; jeśli zmienność, wariancja lub rozstęp międzykwartylowy; jeśli ryzyko przekroczenia progu, odsetek wartości powyżej progu. Ta redukcja jest jednak \emph{zależna od celu}. Średnia nie opisuje wystarczająco rozkładu z dwoma wyraźnymi skupieniami, a mediana nie jest na ogół najlepszym narzędziem do szacowania całkowitej sumy.

\begin{definicja}[Statystyka i jej rozkład]\label{def:statistic}
Niech $\mathcal E=(\mathsf X,\calA,\{P_\theta:\theta\in\Theta\})$ będzie eksperymentem statystycznym, a $(\mathsf Y,\mathcal C)$ przestrzenią mierzalną. Każde mierzalne odwzorowanie $T:(\mathsf X,\calA)\to(\mathsf Y,\mathcal C)$ nazywamy \emph{statystyką}. Jej \emph{rozkład z próby} przy parametrze $\theta$ to miara $P_\theta\circ T^{-1}$ na $\mathcal C$.
\end{definicja}

W praktyce słowo „statystyka” bywa używane zarówno dla funkcji $T$ obliczanej z \emph{losowej} próby, jak i dla zaobserwowanej wartości $T(x)$. Rozróżnienie jest ważne. Przy tworzeniu przedziałów ufności wykorzystujemy rozkład losowej statystyki, natomiast po zebraniu danych podstawiamy konkretny wynik.

\begin{przyklad}[Średnia kontra liczba przekroczeń]
Dla $X=(X_1,\ldots,X_n)\in\R^n$ funkcje $\overline X_n=n^{-1}\sum_iX_i$ oraz $N_a=\sum_i\ind{\{X_i>a\}}$ są statystykami. Jeśli chcemy poznać średni poziom $\E X_1$, pierwsza jest naturalnym kandydatem. Jeśli pytamy o prawdopodobieństwo przekroczenia dopuszczalnej temperatury $a$, druga statystyka jest dopasowana do celu, a średnia może ukryć częste przekroczenia progu.
\end{przyklad}

\section{Średnia, wariancja i błąd standardowy}
\begin{definicja}[Podstawowe statystyki liczbowe]\label{def:meanvar}
Dla $n\ge2$ i próby rzeczywistej $X_1,\ldots,X_n$ definiujemy
\[
\overline X_n=\frac1n\sum_{i=1}^nX_i,\qquad
S_n^2=\frac1{n-1}\sum_{i=1}^n(X_i-\overline X_n)^2,\qquad S_n=\sqrt{S_n^2}.
\]
Dla $n=1$ średnia jest określona, lecz wzór na $S_n^2$ nie jest zdefiniowany. Wzór z mianownikiem $n$ oznaczymy przez $V_n=n^{-1}\sum_i(X_i-\overline X_n)^2$ i nie będziemy utożsamiać tych dwóch wielkości.
\end{definicja}

\begin{lemat}[Tożsamość sumy kwadratów]\label{lem:sum-of-squares}
Dla dowolnych liczb rzeczywistych $x_1,\ldots,x_n$ i dowolnego $a\in\R$, przy $\bar x=n^{-1}\sum_i x_i$, zachodzi
\[
\sum_{i=1}^n(x_i-a)^2
=\sum_{i=1}^n(x_i-\bar x)^2+n(\bar x-a)^2.
\]
\end{lemat}
\begin{proof}
Podstawiamy $x_i-a=(x_i-\bar x)+(\bar x-a)$ i rozwijamy kwadraty. Suma wyrazów mieszanych to $2(\bar x-a)\sum_i(x_i-\bar x)=0$ z definicji średniej. Pozostają dwa pozostałe składniki.
\end{proof}

\emph{Intuicja.} Średnia $\bar x$ minimalizuje sumę kwadratów odległości od stałej: w poprzedniej tożsamości składnik $n(\bar x-a)^2$ jest nieujemny i znika tylko przy $a=\bar x$. Stąd związek średniej z metodą najmniejszych kwadratów, którą rozwiniemy w rozdziale o regresji.

\begin{twierdzenie}[Nieobciążoność średniej i wariancji]\label{thm:unbiased-variance}
Niech $n\ge2$, a $X_1,\ldots,X_n$ będą niezależnymi zmiennymi o jednakowym rozkładzie rzeczywistym, z $\mu=\E X_1\in\R$ oraz $\sigma^2=\Var(X_1)<\infty$. Wówczas
\[
\E\overline X_n=\mu,\qquad
\Var(\overline X_n)=\frac{\sigma^2}{n},\qquad
\E S_n^2=\sigma^2,
\]
podczas gdy $\E V_n=(n-1)\sigma^2/n$.
\end{twierdzenie}
\begin{proof}
Z liniowości wartości oczekiwanej $\E\overline X_n=n^{-1}\sum_i\E X_i=\mu$. Z niezależności $\Var(\overline X_n)=n^{-2}\sum_i\Var(X_i)=\sigma^2/n$. Stosując lemat~\ref{lem:sum-of-squares} do losowych wartości $x_i=X_i$ i $a=\mu$, dostajemy tożsamość prawie na pewno
\[
\sum_{i=1}^n(X_i-\overline X_n)^2
=\sum_{i=1}^n(X_i-\mu)^2-n(\overline X_n-\mu)^2.
\]
Wszystkie składniki są całkowalne, gdyż istnieje drugi moment. Po wzięciu wartości oczekiwanych prawa strona ma wartość $n\sigma^2-n\sigma^2/n=(n-1)\sigma^2$. Podzielenie przez $n-1$, a następnie przez $n$, daje pozostałe równości.
\end{proof}

\begin{uwaga}[Dlaczego $n-1$?]
Nie chodzi o umowną poprawkę, która zawsze zwiększa dokładność każdej możliwej procedury. Podczas liczenia odchyleń $X_i-\overline X_n$ oszacowaliśmy średnią z tych samych danych i otrzymaliśmy liniowy warunek $\sum_i(X_i-\overline X_n)=0$. Dla modelu i celu opisanego w twierdzeniu mianownik $n-1$ koryguje oczekiwane zmniejszenie sumy kwadratów. Nie wynika stąd, że $S_n^2$ zawsze ma mniejszy błąd średniokwadratowy niż $V_n$; są to różne kryteria jakości.
\end{uwaga}

\begin{wniosek}[Ograniczenie precyzji średniej]\label{cor:se}
Przy założeniach twierdzenia~\ref{thm:unbiased-variance} odchylenie standardowe losowej średniej wynosi $\sigma/\sqrt n$. Wielkość $S_n/\sqrt n$ jest jego naturalnym estymatorem i bywa nazywana \emph{szacowanym błędem standardowym średniej}. Nie należy mylić jej z $S_n$, które mierzy rozproszenie pojedynczych obserwacji.
\end{wniosek}

\begin{stwierdzenie}[Błąd estymatora: wariancja i obciążenie]\label{prop:biasvar}
Niech $T$ będzie kwadratowo całkowalną rzeczywistą statystyką w modelu z parametrem rzeczywistym $\theta$. Definiując $\operatorname{Bias}_\theta(T)=\E_\theta T-\theta$ oraz $\MSE_\theta(T)=\E_\theta(T-\theta)^2$, mamy
\[
\MSE_\theta(T)=\Var_\theta(T)+\bigl(\operatorname{Bias}_\theta(T)\bigr)^2.
\]
\end{stwierdzenie}
\begin{proof}
Piszemy $T-\theta=(T-\E_\theta T)+(\E_\theta T-\theta)$ i rozwijamy kwadrat. Wyraz mieszany znika po wzięciu wartości oczekiwanej, ponieważ $\E_\theta(T-\E_\theta T)=0$.
\end{proof}

\paragraph{Wykorzystanie w praktyce.}
Gdy estymujemy $\sigma^2$ w opisanym modelu i wymagamy nieobciążoności, naturalnym wyborem jest $S_n^2$. Nie jest to jednak automatyczny wybór dla wszystkich innych parametrów i funkcji straty. Gdy celem jest predykcja albo minimalizowanie błędu kwadratowego, porównujemy MSE i nie wykluczamy estymatorów obciążonych. Dopiero dalsza teoria estymacji pozwoli formalnie oceniać tę wymianę wariancji na obciążenie.

\section{Prawa wielkich liczb i centralne twierdzenie graniczne w praktyce}
\begin{twierdzenie}[Przypomnienie SPWL i CTG dla średniej]\label{thm:clt-review}
Niech $X_1,X_2,\ldots$ będą niezależne i o jednakowym rozkładzie rzeczywistym, z $\mu=\E X_1$ oraz $0<\sigma^2=\Var(X_1)<\infty$. Wówczas $\overline X_n\to\mu$ według prawdopodobieństwa oraz
\[
\frac{\sqrt n(\overline X_n-\mu)}{\sigma}\;\xrightarrow{\mathcal D}\;\Normal(0,1).
\]
Jeśli $S_n^2$ jest wariancją z próby, to również
$\sqrt n(\overline X_n-\mu)/S_n\xrightarrow{\mathcal D}\Normal(0,1)$, gdzie iloraz definiujemy arbitralnie na zdarzeniu $S_n=0$.
\end{twierdzenie}
\begin{proof}
Pierwszy wniosek wynika z prawa wielkich liczb, a drugi z klasycznego CTG udowodnionego w tomie I. Aby uzasadnić zastąpienie $\sigma$ przez $S_n$, korzystamy z tożsamości
$S_n^2=\tfrac n{n-1}\bigl(n^{-1}\sum_iX_i^2-\overline X_n^2\bigr)$. Z prawa wielkich liczb zastosowanego do $X_i$ i $X_i^2$ (oba mają skończone pierwsze momenty) prawa strona zbiega według prawdopodobieństwa do $\E X_1^2-\mu^2=\sigma^2$. Zatem $S_n/\sigma\to1$ według prawdopodobieństwa. Twierdzenie Slutsky'ego daje tezę. Ponieważ $\sigma>0$, prawdopodobieństwo $S_n=0$ dąży do zera, więc ustalona tam konwencja nie wpływa na granicę.
\end{proof}

\begin{uwaga}[CTG nie stwierdza, że dane są normalne]
Centralne twierdzenie graniczne mówi o \emph{odpowiednio standaryzowanej średniej}, a nie o rozkładzie każdej pojedynczej obserwacji. Dla bardzo ciężkich ogonów wariancja może nie istnieć, a klasyczna wersja CTG nie ma zastosowania. Dla danych zależnych mogą być potrzebne inne skalowania i inna wariancja graniczna. W szczególności nie ma ogólnej reguły, że dla $n=30$ każda próba staje się „normalna”.
\end{uwaga}

\section{Wskaźniki i częstości: próba Bernoulliego}
\begin{twierdzenie}[Rozkład i momenty odsetka sukcesów]\label{thm:prop-bin}
Jeżeli $X_1,\ldots,X_n$ są niezależne, o jednakowym rozkładzie Bernoulliego z parametrem $p\in[0,1]$, to $S_n^*=\sum_iX_i\sim\Bin(n,p)$ oraz $\widehat p_n=\overline X_n$ spełnia
\[
\E\widehat p_n=p,\qquad \Var(\widehat p_n)=\frac{p(1-p)}n.
\]
Jeśli $0<p<1$, to $\sqrt n(\widehat p_n-p)/\sqrt{p(1-p)}\xrightarrow{\mathcal D}\Normal(0,1)$.
\end{twierdzenie}
\begin{proof}
Dla każdego $k\in\{0,\ldots,n\}$ jest $\binom nk$ wektorów z dokładnie $k$ jedynkami. Każdy ma prawdopodobieństwo $p^k(1-p)^{n-k}$, co daje rozkład dwumianowy. Wyniki o średniej i wariancji są szczególnym przypadkiem twierdzenia~\ref{thm:unbiased-variance}, ponieważ $\E X_i=p$ i $\Var(X_i)=p(1-p)$. Dla $0<p<1$ założenia CTG są spełnione i otrzymujemy podaną zbieżność.
\end{proof}

\begin{przyklad}[Prawdopodobieństwo a obserwowany odsetek]\label{ex:failure-rate}
Zaobserwowaliśmy $7$ wad na $100$ niezależnie losowanych produktach i przyjmujemy model Bernoulliego. Odsetek w próbie to $\widehat p=0{,}07$. Gdyby prawdziwe $p$ rzeczywiście wynosiło $0{,}07$, odchylenie standardowe losowego odsetka wynosiłoby $\sqrt{0{,}07\cdot0{,}93/100}\approx0{,}0255$. Jest to \emph{rachunek zależny od przyjętego $p$}, a nie wyznaczenie jego nieznanej wartości. W praktyce zastępuje się $p$ przez $\widehat p$, uzyskując szacowany błąd standardowy. Dla niewielkiej liczby sukcesów albo wartości bliskich $0$ i $1$ lepiej jednak używać odpowiednich przedziałów dla dwumianu, np. Wilsona lub metod dokładnych, zamiast bezkrytycznie korzystać z symetrycznego przybliżenia normalnego.
\end{przyklad}

\section{Dokładny rozkład średniej i wariancji w próbie normalnej}
\begin{lemat}[Niezmienniczość rozkładu normalnego względem przekształceń ortogonalnych]\label{lem:gauss-rotation}
Niech $Z=(Z_1,\ldots,Z_n)\sim\Normal_n(0,I_n)$, czyli $Z_i$ są niezależne i $\Normal(0,1)$. Dla dowolnej rzeczywistej macierzy ortogonalnej $U\in\R^{n\times n}$ wektor $UZ$ ma także rozkład $\Normal_n(0,I_n)$, a więc jego współrzędne są niezależnymi standardowymi zmiennymi normalnymi.
\end{lemat}
\begin{proof}
Gęstość wektora $Z$ względem $n$-wymiarowej miary Lebesgue'a ma postać
\[
 p_Z(z)=(2\pi)^{-n/2}\exp(-\norm{z}^2/2).
\]
Zmiana zmiennych $y=Uz$ ma wyznacznik Jacobiego o module $1$, zaś $\norm{U^{-1}y}=\norm y$. Gęstość wektora $UZ$ jest więc identyczna z gęstością $Z$. Ponieważ gęstość ta rozkłada się na iloczyn gęstości standardowych normalnych, współrzędne $UZ$ są niezależne.
\end{proof}

\begin{twierdzenie}[Rozkłady statystyk normalnych i ich niezależność]\label{thm:normal-sampling}
Niech $n\ge2$ oraz $X_1,\ldots,X_n$ będą niezależne o rozkładzie $\Normal(\mu,\sigma^2)$, gdzie $\mu\in\R$ i $\sigma>0$. Wówczas
\begin{enumerate}
\item $\overline X_n\sim\Normal(\mu,\sigma^2/n)$;
\item $\dfrac{(n-1)S_n^2}{\sigma^2}\sim\chi^2_{n-1}$;
\item $\overline X_n$ i $S_n^2$ są niezależne;
\item $\dfrac{\sqrt n(\overline X_n-\mu)}{S_n}\sim t_{n-1}$.
\end{enumerate}
\end{twierdzenie}
\begin{proof}
Niech $Z_i=(X_i-\mu)/\sigma$ i $Z=(Z_1,\ldots,Z_n)$. Wybierzmy wektor $e_1=n^{-1/2}(1,\ldots,1)$ i uzupełnijmy go do bazy ortonormalnej $e_1,\ldots,e_n$ w $\R^n$. Istnieje więc macierz ortogonalna $U$, której wiersze są tymi wektorami. Niech $Y=UZ$. Na mocy lematu~\ref{lem:gauss-rotation} wszystkie $Y_1,\ldots,Y_n$ są niezależne i $\Normal(0,1)$. Pierwsza współrzędna to $Y_1=\sqrt n(\overline X_n-\mu)/\sigma$, zatem pierwsza teza wynika bezpośrednio z przeskalowania rozkładu normalnego.

Rzut ortogonalny $Z$ na prostą rozpiętą przez $e_1$ ma postać $\langle Z,e_1\rangle e_1$. Kwadrat długości składowej prostopadłej wynosi
\[
\sum_{i=1}^n\left(Z_i-\overline Z\right)^2
=\norm Z^2-\langle Z,e_1\rangle^2
=\sum_{j=2}^nY_j^2
=\frac{(n-1)S_n^2}{\sigma^2}.
\]
Z definicji rozkładu $\chi^2$ suma kwadratów $n-1$ niezależnych standardowych normalnych ma rozkład $\chi^2_{n-1}$, co dowodzi punktu 2. Statystyka $\overline X_n$ zależy wyłącznie od $Y_1$, a $S_n^2$ wyłącznie od $(Y_2,\ldots,Y_n)$; niezależność współrzędnych daje punkt 3. Wreszcie iloraz $Y_1/\sqrt{(\sum_{j=2}^nY_j^2)/(n-1)}$ ma, z definicji, rozkład Studenta o $n-1$ stopniach swobody. Ten iloraz jest równy wyrażeniu z punktu 4. Ponieważ suma kwadratów $n-1$ standardowych normalnych jest dodatnia prawie na pewno, mianownik nie znika prawie na pewno.
\end{proof}


\begin{stwierdzenie}[W modelu normalnym nieobciążoność może kosztować MSE]
\label{prop:mse-samplevariance}
Niech $n\ge2$ i $X_1,\ldots,X_n$ będą niezależne o rozkładzie
$\Normal(\mu,\sigma^2)$ z $\sigma^2>0$. Niech
$S_n^2=(n-1)^{-1}\sum_i(X_i-\overline X_n)^2$ i
$V_n=n^{-1}\sum_i(X_i-\overline X_n)^2$.
Wówczas dla estymacji $\sigma^2$ przy stracie kwadratowej
\[
\MSE(S_n^2)=\frac{2\sigma^4}{n-1},\qquad
\MSE(V_n)=\frac{(2n-1)\sigma^4}{n^2},
\]
więc $\MSE(V_n)<\MSE(S_n^2)$ dla każdego $n\ge2$.
\end{stwierdzenie}
\begin{proof}
Na mocy twierdzenia~\ref{thm:normal-sampling} zmienna
$Y=(n-1)S_n^2/\sigma^2$ ma rozkład $\chi^2_{n-1}$.
Z jego momentów $\E Y=n-1$ i $\Var(Y)=2(n-1)$ otrzymujemy
$\Var(S_n^2)=2\sigma^4/(n-1)$, a estymator ten jest nieobciążony.
Ponieważ $V_n=(n-1)S_n^2/n$, jego wariancja wynosi
$2(n-1)\sigma^4/n^2$, a obciążenie $-\sigma^2/n$.
Dodanie kwadratu obciążenia daje drugi wzór.
Różnica \mbox{$2/(n-1)-(2n-1)/n^2=(3n-1)/(n^2(n-1))>0$},
co kończy dowód.
\end{proof}
\emph{Wniosek praktyczny.} Jeśli celem jest oszacowanie wariancji z
zerowym obciążeniem, $S_n^2$ spełnia ten warunek. Jeśli przyjmiemy
błąd średniokwadratowy jako kryterium, nawet prosty estymator
$V_n$ ma mniejsze ryzyko w modelu normalnym. Nie oznacza to, że
jest najlepszy spośród \emph{wszystkich} możliwych estymatorów.

\begin{uwaga}[Co dokładnie jest założone?]
Dokładny rozkład $t_{n-1}$ wynika z \emph{normalności całej próby}, a nie wyłącznie z istnienia wariancji. Gdy próba jest nienormalna, odpowiedni iloraz może być asymptotycznie normalny, ale nie musi mieć dokładnego rozkładu Studenta przy skończonym $n$. Test Shapiro--Wilka, diagnostyka wykresów Q--Q i alternatywy dla testu t pojawią się w rozdziałach o diagnostyce i testowaniu hipotez. Nie należy testować normalności surowych danych z różnych grup jako zamiennika sprawdzania założeń modelu błędów.
\end{uwaga}

\begin{przyklad}[Obliczenie statystyki Studenta, bez przedwczesnego testowania]\label{ex:t-demo}
Załóżmy, że mamy próbę $n=9$ z modelu $\Normal(\mu,\sigma^2)$, a wyniki dają $\bar x=102$ i $s=6$. Dla hipotetycznej wartości $\mu_0=100$ otrzymujemy $t=\sqrt9(102-100)/6=1$. Gdy prawdziwe $\mu=\mu_0$, statystyka ma dokładny rozkład $t_8$. Sam wynik $t=1$ nie jest jeszcze decyzją: trzeba z góry ustalić hipotezę alternatywną, poziom istotności i regułę testowania. W ten sposób zapobiegamy myleniu opisu statystyki z gotowym wnioskiem badawczym.
\end{przyklad}

\section{Dystrybuanta empiryczna: dane opisują cały rozkład}\label{sec:ecdf}
\begin{definicja}[Miara i dystrybuanta empiryczna]
Niech $X_1,\ldots,X_n$ będą rzeczywistymi obserwacjami, gdzie $n\ge1$. \emph{Miarą empiryczną} na $\calB(\R)$ nazywamy
$\widehat P_n=n^{-1}\sum_{i=1}^n\delta_{X_i}$, a \emph{dystrybuantą empiryczną} funkcję
\[
\widehat F_n(x)=\widehat P_n(( -\infty,x])=\frac1n\sum_{i=1}^n\ind{\{X_i\le x\}},\qquad x\in\R.
\]
\end{definicja}

Funkcja $\widehat F_n$ jest niemalejąca, prawostronnie ciągła, ma wartości w $\{0,1/n,\ldots,1\}$ i przyjmuje granice $0$ oraz $1$ na końcach prostej. To nie jest założenie o kształcie nieznanego rozkładu: dystrybuantę empiryczną możemy obliczyć z dowolnych danych rzeczywistych.

\begin{stwierdzenie}[Rozkład dystrybuanty empirycznej w ustalonym punkcie]\label{prop:ecdf-pointwise}
Niech $X_1,\ldots,X_n$ będą niezależne i o wspólnej dystrybuancie $F$. Dla dowolnego ustalonego $x\in\R$ zachodzi
\[
n\widehat F_n(x)\sim\Bin(n,F(x)),\quad
\E\widehat F_n(x)=F(x),\quad
\Var(\widehat F_n(x))=\frac{F(x)(1-F(x))}{n}.
\]
\end{stwierdzenie}
\begin{proof}
Zmienne $I_i(x)=\ind{\{X_i\le x\}}$ są niezależnymi zmiennymi Bernoulliego o parametrze $F(x)$; ich suma jest równa $n\widehat F_n(x)$. Stosujemy twierdzenie~\ref{thm:prop-bin}.
\end{proof}

\begin{twierdzenie}[Jednostajna zgodność dystrybuanty empirycznej]\label{thm:gc}
Jeżeli $X_1,X_2,\ldots$ są niezależne i mają wspólną dystrybuantę $F$, to
\[
\sup_{x\in\R}\abs{\widehat F_n(x)-F(x)}\longrightarrow0
\qquad\text{prawie na pewno}.
\]
\end{twierdzenie}
\begin{proof}
Jest to twierdzenie Glivenki--Cantellego, którego dowód przedstawiono w tomie I. Zwracamy uwagę na jego treść: zbieżność jest \emph{jednostajna po wszystkich progach $x$} i zachodzi bez założenia ciągłości dystrybuanty $F$. W szczególności wynik obejmuje rozkłady dyskretne, ciągłe i mieszane.
\end{proof}

\begin{uwaga}[Co to znaczy dla praktyki?]
Dystrybuanta empiryczna umożliwia oszacowanie \emph{dowolnego} prawdopodobieństwa przekroczenia progu: $\Prob(X>a)=1-F(a)$ estymujemy przez $1-\widehat F_n(a)$. W diagnostyce modelu porównujemy $\widehat F_n$ z dystrybuantą dopasowanego rozkładu. Jednak samo twierdzenie Glivenki--Cantellego nie daje przedziału o określonym poziomie ufności przy skończonym $n$; w tym celu potrzebne są dodatkowe nierówności lub kwantyle rozkładu odpowiedniej statystyki.
\end{uwaga}

\begin{center}
\begin{tikzpicture}[x=0.75cm,y=4.1cm,>=Latex]
\draw[->] (-.3,0)--(10.6,0) node[right] {$x$};
\draw[->] (0,-.05)--(0,1.13) node[above] {$\widehat F_5(x)$};
\foreach \y/\lab in {.2/{1/5},.4/{2/5},.6/{3/5},.8/{4/5},1/{1}} {
  \draw (-.08,\y)--(.08,\y) node[left=3pt] {$\lab$};
}
\foreach \x in {1,2,4,7,9} {\draw (\x,-.015)--(\x,.015) node[below=4pt] {\small $\x$};}
\draw[very thick,blue!65!black] (0,0)--(1,0) (1,.2)--(2,.2) (2,.4)--(4,.4) (4,.6)--(7,.6) (7,.8)--(9,.8) (9,1)--(10.2,1);
\foreach \x/\yl/\yh in {1/0/.2,2/.2/.4,4/.4/.6,7/.6/.8,9/.8/1} {
\draw[blue!65!black,fill=white] (\x,\yl) circle (1.6pt);
\fill[blue!65!black] (\x,\yh) circle (1.6pt);
}
\end{tikzpicture}
\end{center}
\noindent\emph{Rysunek 2.1.} Dystrybuanta empiryczna pięciu obserwacji $1,2,4,7,9$. Skok następuje w zaobserwowanej wartości, a zaznaczenie punktu pełnego po prawej obrazuje prawostronną ciągłość.

\section{Statystyki pozycyjne: minimum, maksimum, kwantyle}\label{sec:orderstats}
\begin{definicja}[Statystyki porządkowe]\label{def:order}
Dla rzeczywistej próby $(X_1,\ldots,X_n)$ niech $X_{(1)}\le X_{(2)}\le\cdots\le X_{(n)}$ oznaczają wartości tej próby uporządkowane niemalejąco, z uwzględnieniem powtórzeń. Są one \emph{statystykami porządkowymi}. Dla $p\in(0,1)$ \emph{empirycznym kwantylem lewostronnym} nazywamy $\widehat q_p=\inf\{x:\widehat F_n(x)\ge p\}=X_{(\lceil np\rceil)}$. Przy parzystym $n$ przez \emph{medianę obliczeniową} często rozumie się $\bigl(X_{(n/2)}+X_{(n/2+1)}\bigr)/2$; ta konwencja może różnić się od lewostronnego kwantyla $\widehat q_{1/2}$.
\end{definicja}

\begin{twierdzenie}[Rozkład statystyk porządkowych]\label{thm:order-cdf}
Niech $X_1,\ldots,X_n$ będą niezależne o jednakowej dystrybuancie $F$. Dla dowolnych $k\in\{1,\ldots,n\}$ oraz $x\in\R$ zachodzi
\[
\Prob(X_{(k)}\le x)=\sum_{j=k}^n\binom nj F(x)^j\bigl(1-F(x)\bigr)^{n-j}.
\]
W szczególności $\Prob(X_{(1)}>x)=(1-F(x))^n$ oraz $\Prob(X_{(n)}\le x)=F(x)^n$. Jeśli $F$ ma gęstość $f$ względem miary Lebesgue'a, to $X_{(k)}$ ma gęstość
\[
f_{(k)}(x)=\frac{n!}{(k-1)!(n-k)!}F(x)^{k-1}(1-F(x))^{n-k}f(x)
\]
dla prawie wszystkich $x\in\R$.
\end{twierdzenie}
\begin{proof}
Zdarzenie $\{X_{(k)}\le x\}$ zachodzi wtedy i tylko wtedy, gdy co najmniej $k$ spośród $X_i$ jest mniejsze lub równe $x$. Liczba takich obserwacji ma rozkład $\Bin(n,F(x))$, co daje pierwszy wzór. W szczególności minimum jest większe od $x$ dokładnie wtedy, gdy wszystkie $X_i>x$, a maksimum jest mniejsze lub równe $x$ dokładnie wtedy, gdy wszystkie $X_i\le x$.

Przy założeniu istnienia gęstości dystrybuanta $F$ jest absolutnie ciągła. Funkcja $G(u)=\sum_{j=k}^n\binom nj u^j(1-u)^{n-j}$ jest wielomianem, więc $G\circ F$ jest absolutnie ciągła na każdym przedziale zwartym i ma prawie wszędzie pochodną $G'(F(x))f(x)$. Różniczkujemy skończoną sumę. Przy kolejnym wyrazie wykorzystujemy $j\binom nj=n\binom{n-1}{j-1}$ oraz $(n-j)\binom nj=n\binom{n-1}j$; wszystkie składniki wewnętrzne skracają się parami, pozostawiając
$G'(u)=n\binom{n-1}{k-1}u^{k-1}(1-u)^{n-k}$. Otrzymujemy podaną gęstość.
\end{proof}

\begin{wniosek}[Jednostajny rozkład statystyk porządkowych]\label{cor:uniform-beta}
Jeśli $U_1,\ldots,U_n$ są niezależne i jednostajne na $(0,1)$, to $U_{(k)}\sim\operatorname{Beta}(k,n+1-k)$ oraz $\E U_{(k)}=k/(n+1)$.
\end{wniosek}
\begin{proof}
W twierdzeniu~\ref{thm:order-cdf} przy $0<x<1$ podstawiamy $F(x)=x$ oraz $f(x)=1$. Otrzymujemy gęstość $\frac{n!}{(k-1)!(n-k)!}x^{k-1}(1-x)^{n-k}$, która jest gęstością beta z parametrami $k,n+1-k$. Wzór na wartość oczekiwaną rozkładu beta, udowodniony w tomie I, daje $k/(n+1)$.
\end{proof}

\begin{przyklad}[Kiedy maksimum jest informatywne?]\label{ex:max-uniform}
Niech $X_1,\ldots,X_n$ będą niezależne i jednostajne na $(0,\theta)$, gdzie $\theta>0$. Z poprzedniego twierdzenia $\Prob(X_{(n)}\le x)=(x/\theta)^n$ dla $0<x<\theta$, a z przeskalowania jednostajnego rozkładu wynika $\E X_{(n)}=n\theta/(n+1)$. Zatem statystyka $\frac{n+1}{n}X_{(n)}$ jest nieobciążonym estymatorem $\theta$. W tym modelu maksimum jest bardzo użyteczne, ponieważ parametr jest końcem nośnika. W modelu normalnym maksimum zachowuje się całkiem inaczej i z reguły nie jest pierwszym wyborem do estymowania średniej.
\end{przyklad}

\begin{twierdzenie}[Zgodność kwantyla empirycznego]\label{thm:quant-consistency}
Niech $X_1,X_2,\ldots$ będą niezależne i o jednakowej dystrybuancie $F$. Ustalmy $p\in(0,1)$ i $q\in\R$ takie, że dla każdego $\varepsilon>0$ zachodzi
\[
F(q-\varepsilon)<p<F(q+\varepsilon).
\]
Wówczas empiryczny kwantyl lewostronny $\widehat q_p=\inf\{x:\widehat F_n(x)\ge p\}$ zbiega do $q$ prawie na pewno.
\end{twierdzenie}
\begin{proof}
Na zdarzeniu prawdopodobieństwa $1$, na którym zbieżność jednostajna z twierdzenia~\ref{thm:gc} zachodzi, ustalmy $\varepsilon>0$. Liczby $p-F(q-\varepsilon)$ i $F(q+\varepsilon)-p$ są dodatnie. Dla dostatecznie dużych $n$ mamy więc $\widehat F_n(q-\varepsilon)<p<\widehat F_n(q+\varepsilon)$. Z definicji lewostronnego kwantyla wynika wówczas $q-\varepsilon<\widehat q_p\le q+\varepsilon$. Ponieważ $\varepsilon$ było dowolne, zbieżność wynika na tym samym zdarzeniu. Wystarczy formalnie wybrać najpierw przeliczalną rodzinę $\varepsilon=1/m$, $m\in\N$, aby uzyskać jedno zdarzenie pełnej miary spełniające wszystkie potrzebne oszacowania.
\end{proof}

\begin{uwaga}[Dlaczego założenie ma znaczenie?]
Gdy dystrybuanta ma płaski fragment na poziomie $p$, „kwantyl” nie musi być stabilny względem niewielkich zaburzeń dystrybuanty. Nie wystarcza zatem powiedzieć, że dystrybuanta empiryczna zbiega jednostajnie; trzeba jeszcze zadbać o jednoznaczne, stabilne przecięcie poziomu $p$.
\end{uwaga}

\section{Mediana, średnia obcięta i obserwacje odstające}\label{sec:robust}
\begin{stwierdzenie}[Mediana minimalizuje sumę odchyleń bezwzględnych]\label{prop:median-min}
Niech $x_1,\ldots,x_n\in\R$, a $x_{(1)}\le\cdots\le x_{(n)}$ będą wartościami uporządkowanymi. Funkcja $a\mapsto\sum_{i=1}^n\abs{x_i-a}$ przyjmuje minimum dokładnie dla $a\in[x_{(\lceil n/2\rceil)},x_{(\lfloor n/2+1\rfloor)}]$. Dla nieparzystego $n$ zbiór ten jest jednopunktowy; dla parzystego jest przedziałem między dwiema środkowymi obserwacjami.
\end{stwierdzenie}
\begin{proof}
Niech $g(a)=\sum_i\abs{x_i-a}$. Jeśli $a$ nie jest równy żadnej obserwacji, pochodna wynosi $g'(a)=\#\{i:x_i<a\}-\#\{i:x_i>a\}$: składnik $\abs{x_i-a}$ ma pochodną $1$ przy $a>x_i$ i $-1$ przy $a<x_i$. Gdy ponad połowa obserwacji leży po prawej, pochodna jest ujemna i przesunięcie $a$ w prawo zmniejsza $g$; gdy ponad połowa leży po lewej, pochodna jest dodatnia. W punktach $a=x_i$ używamy jednostronnych pochodnych funkcji wypukłej. Minimum występuje tam, gdzie lewa pochodna jest niedodatnia, a prawa nieujemna, czyli dokładnie między środkowymi obserwacjami. Ten argument obejmuje także wartości powtarzające się.
\end{proof}

W ten sposób otrzymujemy podstawowe geometryczne porównanie: \emph{średnia minimalizuje sumę kwadratów}, a \emph{mediana minimalizuje sumę wartości bezwzględnych}. Dlatego pojedyncza ekstremalna obserwacja ma silniejszy wpływ na średnią niż na medianę.

\begin{definicja}[Wskaźniki odpornego opisu]
Dla rzeczywistych danych $x_1,\ldots,x_n$ definiujemy \emph{rozstęp międzykwartylowy} $\operatorname{IQR}=\widehat q_{3/4}-\widehat q_{1/4}$ w zadanej wcześniej konwencji kwantyli. \emph{Medianowe odchylenie bezwzględne} to $\operatorname{MAD}=\Med_i\abs{x_i-m}$, gdzie $m$ jest wybraną medianą danych. \emph{Średnią obciętą} nazywamy średnią pozostającą po usunięciu ustalonej liczby obserwacji z obu krańców uporządkowanej próby.
\end{definicja}

\begin{zadanie}[Średnia, mediana i odporność na skrajną obserwację]\label{zad:outlier}
Wyniki kolejnych dwunastu pomiarów wynoszą
\[
8,10,10,11,11,12,12,13,13,14,15,31.
\]
Oblicz średnią, wariancję z próby, medianę, rozstęp międzykwartylowy i średnią po usunięciu jednej najmniejszej oraz jednej największej wartości. Zinterpretuj wyniki.
\par\emph{Rozwiązanie.} Suma obserwacji to $160$, a suma ich kwadratów $2514$. Stąd
\[
\bar x=\frac{160}{12}=13\tfrac13,\qquad
s^2=\frac1{11}\left(2514-\frac{160^2}{12}\right)=\frac{1142}{33}\approx34{,}606,
\qquad s\approx5{,}883.
\]
Środkowe obserwacje są obie równe $12$, więc mediana wynosi $12$. W konwencji kwantyli lewostronnych $\widehat q_{1/4}=x_{(3)}=10$, $\widehat q_{3/4}=x_{(9)}=13$, zatem $\operatorname{IQR}=3$. Średnia obcięta po usunięciu $8$ i $31$ wynosi $(160-8-31)/10=12{,}1$.

Wartość $31$ podnosi średnią i odchylenie standardowe, podczas gdy mediana oraz środkowy zakres danych są znacznie mniej wrażliwe. Jeśli $31$ jest skutkiem błędu zapisu, należy najpierw zweryfikować pomiar. Jeśli jest rzeczywistą obserwacją, nie wolno jej arbitralnie usunąć: może mieć istotne znaczenie np. dla oceny ryzyka przekroczenia progu. Wybór między średnią i medianą zależy od pytania, nie od automatycznej reguły usuwania wartości odstających.
\end{zadanie}

\begin{uwaga}[Różne konwencje kwantyli]
Programy statystyczne stosują różne sposoby interpolacji kwantyli przy skończonym $n$. W powyższym zadaniu użyliśmy definicji przez odwrotność dystrybuanty empirycznej. Jeśli oprogramowanie zwraca inne kwartyle, należy sprawdzić jego ustawienia, a nie uznawać jednego z wyników za błąd matematyczny.
\end{uwaga}

\section{Wielowymiarowa próba: kowariancja i korelacja}\label{sec:sample-covariance}
\begin{definicja}[Kowariancja z próby]\label{def:samplecov}
Dla $n\ge2$ i obserwacji wektorowych $X_1,\ldots,X_n\in\R^d$ definiujemy $\overline X=n^{-1}\sum_iX_i$ oraz macierz
\[
\widehat\Sigma=\frac1{n-1}\sum_{i=1}^n(X_i-\overline X)(X_i-\overline X)^{\mathsf T}.
\]
Jeżeli $d=2$, pozadiagonalny współczynnik jest kowariancją z próby; współczynnik korelacji Pearsona uzyskujemy przez podzielenie jej przez iloczyn dodatnich odchyleń standardowych dwóch współrzędnych. Korelacja nie jest zdefiniowana, gdy odchylenie standardowe którejś współrzędnej wynosi zero.
\end{definicja}

\begin{twierdzenie}[Nieobciążoność macierzy kowariancji]\label{thm:samplecov-unbiased}
Niech $X_1,\ldots,X_n$ będą niezależnymi, jednakowo rozłożonymi wektorami w $\R^d$ z $\E\norm{X_1}^2<\infty$, średnią $m=\E X_1$ i macierzą kowariancji $\Sigma=\E[(X_1-m)(X_1-m)^{\mathsf T}]$. Dla $n\ge2$ zachodzi $\E\widehat\Sigma=\Sigma$. Ponadto $\widehat\Sigma$ jest symetryczna i nieujemnie określona dla każdej realizacji próby.
\end{twierdzenie}
\begin{proof}
Rozwijamy macierzową tożsamość
\[
\sum_i(X_i-\overline X)(X_i-\overline X)^{\mathsf T}
=\sum_i(X_i-m)(X_i-m)^{\mathsf T}
-n(\overline X-m)(\overline X-m)^{\mathsf T}.
\]
Wyrazy mieszane znikają, gdyż $\sum_i(X_i-\overline X)=0$. Wartość oczekiwana pierwszej sumy wynosi $n\Sigma$, a drugiego składnika $n\cdot\Sigma/n=\Sigma$ dzięki niezależności wektorów i faktowi $\E\overline X=m$. Po podzieleniu przez $n-1$ dostajemy nieobciążoność. Dla dowolnego $v\in\R^d$ mamy $v^{\mathsf T}\widehat\Sigma v=(n-1)^{-1}\sum_i\langle v,X_i-\overline X\rangle^2\ge0$, więc macierz jest nieujemnie określona.
\end{proof}

\begin{przyklad}[Korelacja zerowa, lecz zależność istotna]
Niech $X\sim\operatorname{Unif}(-1,1)$ i $Y=X^2$. Wówczas $\Cov(X,Y)=\E X^3-(\E X)(\E X^2)=0$. Mimo to $Y$ jest całkowicie wyznaczona przez $X$ i zmienne nie są niezależne. W praktyce wykres rozrzutu może ujawnić wyraźny związek kwadratowy mimo korelacji Pearsona bliskiej zeru. Z tego powodu nie wolno używać pojedynczego współczynnika korelacji jako ogólnego miernika dowolnej zależności.
\end{przyklad}

\section{Co zmienia próba bez zwracania? Poprawka na skończoną populację}\label{sec:fpc}
\begin{twierdzenie}[Wariancja średniej przy losowaniu bez zwracania]\label{thm:fpc}
Niech $N\ge2$ i $y_1,\ldots,y_N\in\R$ oznaczają ustalone wartości w populacji skończonej. Wybieramy $n\in\{1,\ldots,N\}$ jednostek jednostajnie spośród wszystkich uporządkowanych prób bez zwracania i zapisujemy ich wartości jako $X_1,\ldots,X_n$. Niech $\mu_N=N^{-1}\sum_i y_i$ oraz $v_N=N^{-1}\sum_i(y_i-\mu_N)^2$. Wówczas
\[
\E\overline X_n=\mu_N,\qquad
\Var(\overline X_n)=\frac{v_N}{n}\frac{N-n}{N-1}.
\]
\end{twierdzenie}
\begin{proof}
Każda pozycja losowania ma rozkład jednostajny na populacji, więc $\E X_i=\mu_N$ oraz $\Var(X_i)=v_N$. Dla $i\ne j$ każda uporządkowana para dwóch różnych jednostek występuje jednakowo często, stąd
\[
\E(X_iX_j)=\frac1{N(N-1)}\sum_{a\ne b}y_ay_b
=\frac{(\sum_a y_a)^2-\sum_a y_a^2}{N(N-1)}.
\]
Po odjęciu $\mu_N^2$ i uporządkowaniu otrzymujemy $\Cov(X_i,X_j)=-v_N/(N-1)$. Zatem
\[
\Var(\overline X_n)=\frac1{n^2}\left(nv_N+n(n-1)\frac{-v_N}{N-1}\right)
=\frac{v_N}{n}\frac{N-n}{N-1}.
\]
Wartość oczekiwana średniej jest równa $\mu_N$ z liniowości wartości oczekiwanej. Dla $n=N$ wariancja rzeczywiście znika: poznaliśmy całą populację.
\end{proof}

\begin{uwaga}[Porównanie z próbą prostą z populacji nieskończonej]
Czynnik $(N-n)/(N-1)$ zmniejsza wariancję średniej względem losowania ze zwracaniem. W podręcznikach spotyka się równoważny zapis $\Var(\overline X_n)=(1-n/N)S_N^2/n$, gdzie $S_N^2=(N-1)^{-1}\sum_i(y_i-\mu_N)^2$. Poprawka jest istotna przy dużym udziale próby w populacji; nie należy jej stosować bezrefleksyjnie do danych zależnych z innych powodów, np. do szeregu czasowego.
\end{uwaga}

\section{Oszacowania koncentracji: ile obserwacji potrzeba?}
\begin{twierdzenie}[Przypomnienie nierówności Hoeffdinga]\label{thm:hoeffding-demo}
Niech $X_1,\ldots,X_n$ będą niezależnymi zmiennymi losowymi przyjmującymi wartości w $[0,1]$, niekoniecznie o jednakowych rozkładach. Niech $m_n=n^{-1}\sum_i\E X_i$. Dla każdego $\varepsilon>0$,
\[
\Prob\!\left(\abs{\overline X_n-m_n}\ge\varepsilon\right)
\le 2e^{-2n\varepsilon^2}.
\]
\end{twierdzenie}
\begin{proof}
Jest to szczególny przypadek nierówności Hoeffdinga, udowodnionej w tomie I. Jej zaletą jest brak konieczności znajomości dokładnego rozkładu zmiennych poza niezależnością i ograniczeniem przedziałem $[0,1]$.
\end{proof}

\begin{zadanie}[Orientacyjna liczebność próby]\label{zad:sample-size}
Chcemy, aby przy niezależnych pomiarach Bernoulliego prawdopodobieństwo odchylenia $\widehat p_n$ od $p$ o co najmniej $0{,}1$ było co najwyżej $0{,}05$. Jaką liczebność $n$ gwarantuje nierówność Hoeffdinga?
\par\emph{Rozwiązanie.} Wystarczy zapewnić $2\exp(-2n\cdot0{,}1^2)\le0{,}05$, czyli
\[
n\ge\frac{\log(2/0{,}05)}{2\cdot0{,}1^2}=\frac{\log40}{0{,}02}\approx184{,}44.
\]
Zatem $n=185$ jest wystarczające. Jest to konserwatywna gwarancja dla każdego $p\in[0,1]$ przy niezależności, a nie optymalna liczba obserwacji w każdym modelu. W praktyce należy dodatkowo uwzględnić schemat losowania, koszty pomiaru i to, jak precyzyjny wniosek jest potrzebny.
\end{zadanie}

\section{Różne cele, różne statystyki: porównanie metod}\label{sec:choose-stats}
W praktyce nie zaczynamy od pytania „który test jest najlepszy?”, lecz od pytania, \emph{co chcemy oszacować lub porównać}. Na przykład dla danych z zadania~\ref{zad:outlier} odpowiedź zależy od celu: średnia opisuje sumę podzieloną przez liczbę pomiarów, mediana opisuje środek uporządkowanej próby, a odsetek przekroczeń progu $20$ mierzy ryzyko dużego wyniku.

\begin{center}\small
\begin{tabular}{>{\raggedright\arraybackslash}p{0.25\textwidth}>{\raggedright\arraybackslash}p{0.27\textwidth}>{\raggedright\arraybackslash}p{0.36\textwidth}}
\toprule
\textbf{Pytanie} & \textbf{Co obliczyć najpierw?} & \textbf{Kiedy uważać?}\\
\midrule
Średni wynik lub łączna produkcja & Średnia, odchylenie standardowe, błąd standardowy & Zależne obserwacje, brak reprezentatywności, ciężkie ogony\\
Typowa obserwacja & Mediana, kwartyle, IQR & Przy bardzo małej próbie kwantyle są niestabilne; definicje różnią się\\
Przekroczenia progu & Odsetek, dystrybuanta empiryczna & Zbyt mało przekroczeń dla przybliżeń normalnych\\
Rozproszenie pomiarów & $S_n$, IQR, MAD & Jednostka, asymetria, wartości skrajne\\
Zależność dwóch wielkości & Wykres rozrzutu, kowariancja, korelacja & Nieliniowość, obserwacje odstające, wspólna przyczyna\\
Porównanie grup & Statystyki w każdej grupie osobno & Różne liczebności, heteroskedastyczność, sparowanie\\
\bottomrule
\end{tabular}
\end{center}

\paragraph{Diagnostyka jest częścią analizy.}
Przed późniejszym testem Shapiro--Wilka lub ANOVA warto obejrzeć histogram, wykres Q--Q, wykres reszt i wykresy poszczególnych grup. W modelach ANOVA normalność dotyczy rozkładu błędów przy danych wartościach czynników (w klasycznym modelu normalnym), a nie koniecznie połączonego histogramu wszystkich obserwacji. Wynik testu normalności z wartością $p>0{,}05$ nie dowodzi normalności; wynik $p<0{,}05$ nie oznacza automatycznie, że każda metoda oparta na średnich jest bezużyteczna. O właściwym doborze metody decydują także odporność na naruszenia założeń, wielkość próby i cel badania.

\section{Pełne studium przypadku: produkcja, progi i wybór statystyki}\label{sec:case-study}
\begin{zadanie}[Dwa pytania o te same dane]\label{zad:full-case}
Zmierzyliśmy w sekundach czas wykonania zadania przez dwanaście identycznych urządzeń, uzyskując dane z zadania~\ref{zad:outlier}. Dział kontroli jakości chce odpowiedzi na dwa pytania: (i) jaki jest typowy czas wykonania? (ii) jak często czas przekracza $20$ sekund? Jakie statystyki i założenia należy zaproponować?
\par\emph{Rozwiązanie.} Dla pytania (i) najpierw przedstawiamy zarówno średnią $13{,}33$ s, jak i medianę $12$ s, IQR równy $3$ s oraz wartość odstającą $31$ s. Jeśli priorytetem jest przewidywanie \emph{całkowitego czasu wszystkich operacji}, średnia może być właściwą wielkością docelową mimo wartości skrajnej. Jeśli priorytetem jest \emph{typowa operacja}, mediana jest czytelniejsza. Jednocześnie sprawdzamy, czy $31$ s jest pomyłką, czy rzeczywistym opóźnieniem.

Dla pytania (ii) zliczamy $\sum_i\ind{\{x_i>20\}}=1$, stąd odsetek w próbie wynosi $1/12\approx0{,}0833$. Nie należy uznawać tego za dokładne prawdopodobieństwo przekroczenia progu w populacji; próba jest mała, a jeden przypadek daje ograniczoną informację. Jeśli urządzenia pracowały w różnych warunkach lub nie zostały losowo dobrane, potrzebujemy modelu uwzględniającego te okoliczności. Gdy niezależność i jednakowe warunki są wiarygodne, wskaźnik przekroczenia można modelować Bernoullim i skorzystać później z odpowiednich przedziałów dwumianowych.

\emph{Wniosek praktyczny:} jeden zbiór danych uzasadnia \emph{różne} statystyki zależnie od pytania. Średnia, mediana i odsetek przekroczeń nie są konkurentami o jeden uniwersalny tytuł „najlepszej statystyki”; mierzą różne aspekty rozkładu.
\end{zadanie}

\section{Przykład numeryczny w Pythonie i Julii}
Poniższe skrypty obliczają opis próbki z zadania~\ref{zad:outlier}. Nie przeprowadzają żadnego testu hipotez: najpierw trzeba poprawnie zdefiniować cel i model. Do kwartyli używamy tutaj konwencji odwrotności dystrybuanty empirycznej, nie domyślnej interpolacji oprogramowania.

\Needspace{16\baselineskip}
\noindent\textbf{Python (NumPy).}\par
\par
\Needspace{13\baselineskip}
\smallskip
\begin{lstlisting}[language=Python,caption={Opis próby: średnia, mediana, wariancja i kwantyle (Python).}]
import numpy as np
x = np.array([8, 10, 10, 11, 11, 12, 12, 13, 13, 14, 15, 31])
mean = x.mean()
median = np.median(x)
variance_unbiased = x.var(ddof=1)
q25 = np.quantile(x, 0.25, method="inverted_cdf")
q75 = np.quantile(x, 0.75, method="inverted_cdf")
share_above_20 = np.mean(x > 20)
print(mean, median, variance_unbiased, q25, q75, share_above_20)
\end{lstlisting}

\Needspace{16\baselineskip}
\noindent\textbf{Julia (Statistics).}\par
\par
\Needspace{13\baselineskip}
\smallskip
\begin{lstlisting}[language=Julia,caption={Opis próby z użyciem pakietu Statistics (Julia).}]
using Statistics
x = [8, 10, 10, 11, 11, 12, 12, 13, 13, 14, 15, 31]
mu = mean(x)
med = median(x)
s2 = var(x; corrected=true)
y = sort(x)
q25 = y[ceil(Int, length(y) * 0.25)]
q75 = y[ceil(Int, length(y) * 0.75)]
share_above_20 = count(v -> v > 20, x) / length(x)
println((mu, med, s2, q25, q75, share_above_20))
\end{lstlisting}

\paragraph{Jak interpretować wyniki programu?}
Obliczenia nie sprawdzają założeń modelu ani nie rozstrzygają, czy wartość $31$ jest błędem pomiarowym. Kod nie stanowi również uzasadnienia, że obserwacje są niezależne. Jeśli próba pochodzi z wielu partii lub z jednego szeregu czasowego, należy zmienić model niepewności, zamiast jedynie wybrać inną funkcję z biblioteki statystycznej.

\section{Wskazówki do dalszej lektury: od opisu do wnioskowania}
Z rozdziału 1 wiemy, jaki model obserwacji jest rozważany, a z rozdziału 2 -- jakie statystyki można obliczyć i jakie mają rozkłady. Kolejnym krokiem będzie \emph{funkcja wiarygodności}, następnie \emph{ocena decyzji przez funkcję straty}, a w dalszych rozdziałach teoria estymatorów, testów i przedziałów. W szczególności porównamy:
\begin{itemize}
\item średnią jako estymator $\mu$ i medianę jako estymator typowego położenia, w tym ich odporność i własności asymptotyczne;
\item estymatory parametryczne i nieparametryczne, MLE, metodę momentów i metody odporne;
\item dokładne wnioskowanie z rozkładu t-Studenta i przybliżenia asymptotyczne;
\item ANOVA, ANCOVA, testy Welcha i testy nieparametryczne, podając jawnie warunki ich stosowania;
\item klasyczne przedziały ufności i bayesowskie przedziały wiarygodności, interpretowane w odpowiednich modelach.
\end{itemize}

\paragraph{Lista kontrolna przed obliczeniami.}
Niezależnie od dziedziny zastosowania należy umieć odpowiedzieć na cztery pytania: \emph{kto jest jednostką obserwacji?}, \emph{co jest parametrem docelowym?}, \emph{jakie losowanie i jaka zależność są założone?}, \emph{czy wybrana statystyka rzeczywiście mierzy tę wielkość?}. Dopiero wtedy konkretny test lub algorytm nabiera właściwego znaczenia.

% ================================================================
% ROZDZIAL 3 -- FUNKCJA WIARYGODNOSCI
% ================================================================
\chapter{Funkcja wiarygodności i informacja zawarta w danych}
\label{ch:likelihood}
\paragraph{Powiązania i odsyłacze.} Dalsze zastosowania wiarygodności w MLE i informacji Fishera podajemy w rozdziałach~\ref{ch:estimators} i~\ref{ch:fisher-information}; warunkowanie jest rozwinięte w tomie I, rozdziale~\ref{I-ch:condexp}.


\section{Dlaczego sam rozkład próby nie wystarcza?}
W tomie I obliczaliśmy prawdopodobieństwa \emph{przy ustalonym modelu}. Teraz
obserwujemy wynik $x$ i porównujemy wartości parametru $\theta$, pod którymi
wynik ten jest bardziej albo mniej zgodny z modelem. Nie zamieniamy przy tym
automatycznie parametru w zmienną losową. W podejściu klasycznym $\theta$
pozostaje nieznaną, ale ustaloną wartością. W szczególności wyrażenie
$L_x(\theta)$ nie oznacza $\Prob(\theta\mid X=x)$.

\begin{przyklad}[Dwie hipotezy o wadliwości]
Jeśli w $100$ niezależnych kontrolach stwierdzono $7$ wad, model
$S\sim\Bin(100,p)$ przypisuje prawdopodobieństwo
$\binom{100}{7}p^7(1-p)^{93}$ obserwacji $S=7$. Dla $p=0{,}05$ i $p=0{,}07$
możemy porównać te prawdopodobieństwa. Iloraz jest miarą \emph{względnej
zgodności zaobserwowanego wyniku z dwiema ustalonymi wartościami parametru},
a nie prawdopodobieństwem, że drugi parametr jest prawdziwy.
\end{przyklad}

\begin{definicja}[Funkcja wiarygodności]\label{def:likelihood}
Niech $\mathcal E=(\mathsf X,\calA,\{P_\theta:\theta\in\Theta\})$ będzie
modelem zdominowanym przez ustaloną miarę $\sigma$-skończoną $\mu$.
Wybierzmy dla wszystkich $\theta\in\Theta$ określone wersje gęstości
$p_\theta=\dd P_\theta/\dd\mu$. Dla zaobserwowanej wartości $x\in\mathsf X$
\emph{funkcją wiarygodności} nazywamy funkcję
$L_x:\Theta\to[0,\infty]$, $L_x(\theta)=p_\theta(x)$.
Jeśli $p_\theta(x)>0$, logarytmiczną funkcją wiarygodności jest
$\ell_x(\theta)=\log p_\theta(x)$; dla $p_\theta(x)=0$ przyjmujemy
$\ell_x(\theta)=-\infty$.
\end{definicja}

\begin{uwaga}[Wersje gęstości, nośnik i interpretacja]
Gęstość jest określona wyłącznie $\mu$-prawie wszędzie dla każdego
\emph{ustalonego} $\theta$. Zmieniając jej wersję na zbiorze zerowym
zależnym od $\theta$, można w nieprzeliczalnej rodzinie sztucznie zmienić
wartości punktowe wiarygodności. W praktyce wybiera się wspólną,
jednoznacznie opisaną postać $p_\theta(x)$, a w twierdzeniach o MLE
wprowadza odpowiednie warunki regularności i mierzalności.
Jeśli na całej przestrzeni parametrów $L_x$ jest zerowa, model przypisuje
obserwacji zerową gęstość w każdej rozważanej wartości parametru; nie
wyciągamy wtedy wniosku o maksimum wiarygodności bez osobnej analizy modelu.
\end{uwaga}

\begin{stwierdzenie}[Zmiana miary dominującej i względna wiarygodność]\label{prop:change-reference}
Załóżmy, że $\mu$ i $\nu$ są równoważnymi miarami $\sigma$-skończonymi na
$(\mathsf X,\calA)$ i że wybrano zgodne wersje gęstości $p_\theta^\mu$
i $p_\theta^\nu$, dla których w ustalonym punkcie $x$ zachodzi
$p_\theta^\nu(x)=c(x)p_\theta^\mu(x)$ dla wszystkich $\theta\in\Theta$,
gdzie $0<c(x)<\infty$ nie zależy od $\theta$. Wtedy porządek wartości
wiarygodności, ilorazy tych wartości i zbiory maksimów pozostają niezmienione.
\end{stwierdzenie}
\begin{proof}
Dla dowolnych $\theta_1,\theta_2$ mamy
$L_x^\nu(\theta_1)/L_x^\nu(\theta_2)
=L_x^\mu(\theta_1)/L_x^\mu(\theta_2)$, o ile mianowniki są dodatnie.
Ponadto mnożenie wszystkich wartości przez tę samą dodatnią liczbę nie
zmienia relacji $\le$ ani argumentów maksimum. Związek wersji gęstości
wynika z reguły łańcuchowej Radona--Nikodyma, ale jej równość prawie
wszędzie sama przez się nie daje równości punktowej we wszystkich parametrach.
\end{proof}


\begin{uwaga}[Maksimum punktowe a estymator losowy]
W definicji MLE nie wystarczy dla każdej obserwacji $x$ wskazać
liczbę maksymalizującą $L_x$: odwzorowanie $x\mapsto\widehat\theta(x)$
musi być także mierzalne. Przy zwartej przestrzeni parametrów,
ciągłości odpowiedniej funkcji względem parametru i wspólnych
warunkach mierzalności można stosować twierdzenia o mierzalnym
wyborze; nie wynikają one jednak z samej reguły „wyzeruj score”.
W wszystkich jawnych modelach rozważanych poniżej podajemy
bezpośrednie, mierzalne formuły.
\end{uwaga}

\section{Próba niezależna i logarytm wiarygodności}
\begin{stwierdzenie}[Rozkład iloczynowy]\label{prop:product-likelihood}
Niech $X_1,\ldots,X_n$ będą niezależne i mają wspólny rozkład z gęstością
$f_\theta$ względem miary $\sigma$-skończonej $\lambda$ na $(\mathsf S,\mathcal S)$.
Wybierzmy wspólną wersję gęstości iloczynowej
$p_\theta(x_1,\ldots,x_n)=\prod_{i=1}^nf_\theta(x_i)$.
Wówczas $L_x(\theta)=\prod_i f_\theta(x_i)$, a dla dodatnich czynników
$\ell_x(\theta)=\sum_i\log f_\theta(x_i)$.
\end{stwierdzenie}
\begin{proof}
Dla mierzalnych prostokątów $A_1\times\cdots\times A_n$ niezależność
daje $P_\theta(X_i\in A_i,\ i=1,\ldots,n)=\prod_i\int_{A_i}f_\theta\,\dd\lambda$.
Z twierdzenia Tonellego prawa strona jest całką $\prod_i f_\theta(x_i)$
na prostokącie względem $\lambda^{\otimes n}$. Prostokąty tworzą klasę
$\pi$ generującą \mbox{$\mathcal S^{\otimes n}$}, zatem jednoznaczność
miar probabilistycznych rozszerza tę równość na całe $\sigma$-ciało.
Ostatnia równość wynika z $\log\prod_i a_i=\sum_i\log a_i$ dla $a_i>0$.
\end{proof}

\emph{Motywacja obliczeniowa.} Iloczyn wielu małych wartości może prowadzić
do niedomiaru numerycznego. Suma logarytmów jest stabilniejsza, a maksimum
pozostaje niezmienione, ponieważ logarytm jest ściśle rosnący. Wyrazy
stałe względem parametru można pominąć \emph{przy maksymalizacji}, ale nie
zawsze przy obliczaniu bezwzględnych wartości funkcji wiarygodności.

\section{Estymator największej wiarygodności}
\begin{definicja}[MLE]\label{def:mle}
W eksperymencie z definicji~\ref{def:likelihood} \emph{estymatorem
największej wiarygodności} (MLE) nazywamy mierzalny wybór
$\widehat\theta:\mathsf X\to\Theta$ taki, że dla obserwacji $x$ należących
do określonego zbioru o pełnej mierze dla każdego $P_\theta$ zachodzi
$L_x(\widehat\theta(x))=\sup_{\vartheta\in\Theta}L_x(\vartheta)>0$.
Jeśli maksimum nie jest osiągane lub nie istnieje mierzalny wybór,
MLE w tej postaci może nie istnieć. Gdy maksimów jest kilka, estymator
nie musi być jednoznaczny.
\end{definicja}

\begin{twierdzenie}[MLE w modelu Bernoulliego]\label{thm:mle-bernoulli}
Niech $X_1,\ldots,X_n$ będą niezależne o rozkładzie $\operatorname{Bernoulli}(p)$,
$p\in[0,1]$, i $S=\sum_iX_i$. Wybierając wersję funkcji wiarygodności
$L_x(p)=p^{S(x)}(1-p)^{n-S(x)}$ z interpretacją wartości brzegowych przez
ciągłość, otrzymujemy jednoznaczny MLE $\widehat p=S/n$.
\end{twierdzenie}
\begin{proof}
Jeżeli $0<S<n$, to dla $p\in(0,1)$ mamy
$\ell_x(p)=S\log p+(n-S)\log(1-p)$ i
$\ell_x'(p)=S/p-(n-S)/(1-p)$, zatem jedyny punkt stacjonarny to
$p=S/n$. Ponadto $\ell_x''(p)=-S/p^2-(n-S)/(1-p)^2<0$, więc jest to
maksimum ścisłe. Przy $p\downarrow0$ lub $p\uparrow1$ funkcja dąży do
$-\infty$. Jeżeli $S=0$, wiarygodność równa $(1-p)^n$ jest ściśle malejąca
i ma maksimum w $p=0$; gdy $S=n$, funkcja $p^n$ jest ściśle rosnąca i ma
maksimum w $p=1$. Rozpatrzyliśmy wszystkie możliwe wyniki.
\end{proof}

\begin{zadanie}[Sto produktów: co naprawdę mówi wiarygodność?]\label{zad:lik-quality}
W $100$ niezależnych kontrolach znaleziono $7$ wad. Oblicz MLE $p$ i iloraz
wiarygodności hipotez $p=0{,}07$ i $p=0{,}05$. Czy iloraz jest prawdopodobieństwem
prawdziwości hipotezy?
\par\emph{Rozwiązanie.} Z twierdzenia~\ref{thm:mle-bernoulli} mamy
$\widehat p=7/100=0{,}07$. Dla dwóch ustalonych parametrów
\[
\frac{L_x(0{,}07)}{L_x(0{,}05)}
=\left(\frac{0{,}07}{0{,}05}\right)^7
 \left(\frac{0{,}93}{0{,}95}\right)^{93}
\approx1{,}457.
\]
Dane są około $1{,}46$ raza bardziej wiarygodne pod parametrem $0{,}07$
niż pod parametrem $0{,}05$. To porównanie \emph{dwóch wartości parametru},
a nie rozkład prawdopodobieństwa na wartościach $p$. Sama liczba nie
stanowi jeszcze testu o kontrolowanym prawdopodobieństwie błędu ani
odpowiedzi na pytanie, z jakim prawdopodobieństwem $p=0{,}07$.
\end{zadanie}

\begin{center}
\begin{tikzpicture}[x=28cm,y=3.0cm,>=Latex]
  \draw[->] (0,0)--(.185,0) node[right] {$p$};
  \draw[->] (0,0)--(0,1.16) node[above] {$L(p)/L(0{,}07)$};
  \foreach \x/\lab in {.03/{0{,}03},.05/{0{,}05},.07/{0{,}07},.10/{0{,}10},.15/{0{,}15}} {
    \draw (\x,-.019)--(\x,.019) node[below=3pt] {\scriptsize $\lab$};
  }
  \foreach \y in {.5,1} {\draw (-.0015,\y)--(.0015,\y) node[left=3pt] {\scriptsize $\y$};}
  \draw[thick,blue!65!black,domain=.004:.175,samples=120,smooth]
    plot (\x,{exp(7*ln(\x/.07)+93*ln((1-\x)/.93))});
  \draw[dashed] (.07,0)--(.07,1);
  \fill[blue!65!black] (.07,1) circle (1.7pt);
\end{tikzpicture}
\end{center}
\noindent\emph{Rysunek 3.1.} Względna wiarygodność w modelu
Bernoulliego dla $k=7$ wad i $n=100$ prób. Maksimum wskazuje
$\widehat p=0{,}07$, lecz wykres nie jest rozkładem prawdopodobieństwa
na osi parametru.


\begin{twierdzenie}[MLE intensywności z uwzględnieniem ekspozycji]\label{thm:mle-pois}
Niech $Y_1,\ldots,Y_m$ będą niezależne, $Y_i\sim\Pois(\lambda t_i)$,
gdzie $t_i>0$ są znane, a $\lambda\in[0,\infty)$. Niech
$S=\sum_iY_i$ i $T=\sum_i t_i$. Wtedy
$\widehat\lambda=S/T$ jest jednoznacznym MLE.
\end{twierdzenie}
\begin{proof}
Po pominięciu wyrazów niezależnych od $\lambda$ log-wiarygodność ma dla
$\lambda>0$ postać $\ell(\lambda)=S\log\lambda-\lambda T+C$.
Jeżeli $S>0$, jej pochodna znika wyłącznie w $\lambda=S/T$, a druga
pochodna $-S/\lambda^2$ jest ujemna. Gdy $S=0$, wiarygodność jest
proporcjonalna do $e^{-\lambda T}$ i przyjmuje maksimum przy $\lambda=0$.
Pozostałe czynniki $\prod_i t_i^{y_i}/y_i!$ są dodatnie i niezależne
od $\lambda$, więc nie zmieniają maksymalizatora.
\end{proof}

\begin{zadanie}[Zgłoszenia przy różnych czasach ekspozycji]
Przez $10$, $20$ i $40$ godzin obserwowano trzy maszyny i zliczono
odpowiednio $2$, $5$ i $7$ awarii. Oszacuj intensywność w modelu
$Y_i\sim\Pois(\lambda t_i)$ niezależnie.
\par\emph{Rozwiązanie.} Suma zdarzeń to $14$, a czas ekspozycji $70$
godzin, więc $\widehat\lambda=14/70=0{,}2$ awarii na godzinę.
Nie uśredniamy bez wag trzech intensywności $2/10$, $5/20$, $7/40$:
otrzymalibyśmy inną liczbę, która nie jest MLE tego modelu.
Jeżeli maszyny pracowały w różnych warunkach, model ze wspólnym
$\lambda$ może być niewiarygodny, niezależnie od poprawności rachunku.
\end{zadanie}

\begin{twierdzenie}[MLE w modelu normalnym]\label{thm:mle-normal}
Niech $n\ge2$, a $X_1,\ldots,X_n$ będą niezależne o rozkładzie
$\Normal(\mu,\sigma^2)$, gdzie $\mu\in\R$, $\sigma^2>0$. Dla każdej
próby $x$ o niejednakowych współrzędnych funkcja wiarygodności ma
jednoznaczne maksimum w
$\widehat\mu=\bar x$ oraz
$\widehat{\sigma^2}=n^{-1}\sum_i(x_i-\bar x)^2$.
Jeśli wszystkie obserwacje są identyczne, supremum wiarygodności jest
nieskończone i nie jest osiągane w przestrzeni parametrów.
\end{twierdzenie}
\begin{proof}
Oznaczmy $v=\sigma^2$. Z dokładnością do stałej log-wiarygodność jest
$\ell(\mu,v)=-\frac n2\log v-(2v)^{-1}\sum_i(x_i-\mu)^2$.
Dla każdego ustalonego $v>0$ maksymalizacja względem $\mu$ sprowadza się
do minimalizacji sumy kwadratów. Z lematu~\ref{lem:sum-of-squares}
jedynym minimizatorem jest $\mu=\bar x$. Po podstawieniu
$Q=\sum_i(x_i-\bar x)^2$ otrzymujemy
$\ell(\bar x,v)=-\frac n2\log v-\frac Q{2v}+C$.
Jeśli $Q>0$, pochodna wynosi $(Q-nv)/(2v^2)$: jest dodatnia dla
$v<Q/n$ i ujemna dla $v>Q/n$. Gdy $v\downarrow0$ lub $v\uparrow\infty$,
wartość funkcji dąży do $-\infty$. Jedynym maksimum jest więc $v=Q/n$.
Gdy $Q=0$, mamy $\ell(\bar x,v)=-\frac n2\log v+C\to\infty$
przy $v\downarrow0$, co dowodzi ostatniego twierdzenia.
\end{proof}

\begin{uwaga}[MLE i nieobciążoność to różne kryteria]
Estymator MLE wariancji w modelu normalnym ma mianownik $n$;
nieobciążony estymator z rozdziału 2 ma mianownik $n-1$.
Oba są matematycznie poprawnymi odpowiedziami na \emph{różne kryteria}:
pierwszy maksymalizuje wiarygodność, drugi spełnia warunek
$\E_\theta S_n^2=\sigma^2$. Mniejsza wartość mianownika nie jest
argumentem przeciw MLE, podobnie jak obciążenie MLE nie oznacza jego
bezużyteczności.
\end{uwaga}

\section{Nieregularny model: jednostajny rozkład i koniec nośnika}
\begin{twierdzenie}[MLE końca przedziału]\label{thm:mle-uniform}
Niech $X_1,\ldots,X_n$ będą niezależne o rozkładzie jednostajnym na
$(0,\theta)$, gdzie $\theta>0$. Dla obserwacji $x_i>0$ z
$M=\max_i x_i>0$ oraz wybranej wersji gęstości
$f_\theta(x)=\theta^{-1}\ind{\{0<x\le\theta\}}$ jedynym MLE jest
$\widehat\theta=M$. Co więcej, $\E_\theta\widehat\theta=n\theta/(n+1)$.
\end{twierdzenie}
\begin{proof}
Dla obserwacji dodatnich wiarygodność ma postać
$L_x(\theta)=\theta^{-n}\ind{\{\theta\ge M\}}$.
Dla $\theta<M$ jest zerowa; dla $\theta\ge M$ ściśle maleje, więc
osiąga maksimum w $M$. Wartość oczekiwaną maksimum policzono
w przykładzie~\ref{ex:max-uniform}. Zwróćmy uwagę, że inna wersja
gęstości z ostrą nierównością $x<\theta$ prowadziłaby do supremum
w $\theta\downarrow M$, którego nie osiągano by w tym konkretnym punkcie.
Właśnie dlatego w niektórych modelach z nośnikiem zależnym od parametru
trzeba jawnie ustalić wersję gęstości.
\end{proof}

\begin{uwaga}[Kiedy nie używać zwykłych wzorów score?]
W modelu jednostajnym nośnik zależy od $\theta$. Nie wolno więc
formalnie różniczkować $\int_0^\theta\theta^{-1}\,\dd x=1$
tak, jakby granice całkowania były niezależne od parametru.
W szczególności brak średniego score równego zero przy naiwnym
różniczkowaniu $\log(1/\theta)$. Ta rodzina nie spełnia typowych
warunków regularności klasycznej teorii informacji Fishera.
Dla $0<X<\theta$ formalny score pochodzący z samej gęstości
wewnątrz nośnika wynosi $-1/\theta$, więc
$\E_\theta s_\theta(X)=-1/\theta\ne0$.
Nie jest to sprzeczne z twierdzeniem~\ref{thm:fisher-identities}:
pochodna całki po \emph{ruchomym} nośniku zawiera dodatkowy wyraz
brzegowy, którego tamte hipotezy wykluczały.
\end{uwaga}

\section{Score i informacja Fishera: po co mierzyć krzywiznę wiarygodności?}
\begin{definicja}[Score i informacja Fishera]\label{def:score-fisher}
Niech $\Theta\subset\R$ będzie otwartym przedziałem, a $p_\theta(x)>0$
wersją gęstości dla $\mu$-prawie wszystkich $x$ na wspólnym, niezależnym od
$\theta$ nośniku. W punktach, gdzie istnieje pochodna, definiujemy
\emph{score} $s_\theta(x)=\partial_\theta\log p_\theta(x)$.
Jeśli $s_\theta(X)\in L^2(P_\theta)$, \emph{informacją Fishera} nazywamy
$I(\theta)=\E_\theta[s_\theta(X)^2]$.
Jeśli użyjemy całej próby, informacja jej wspólnego prawa może być różna
od informacji jednej obserwacji; będziemy oznaczać ją przez $I_n(\theta)$.
\end{definicja}

\begin{twierdzenie}[Tożsamości informacyjne przy jawnych założeniach]\label{thm:fisher-identities}
Niech $\Theta\subset\R$ będzie otwartym przedziałem, a $(\mathsf X,\calA,\mu)$
przestrzenią z miarą $\sigma$-skończoną. Załóżmy, że dla każdego
$\theta\in\Theta$ funkcja $x\mapsto p_\theta(x)$ jest dodatnią gęstością
probabilistyczną, że istnieje wspólny zbiór $E\in\calA$ o $\mu(\mathsf X\setminus E)=0$,
na którym $\theta\mapsto p_\theta(x)$ jest dwukrotnie
różniczkowalna dla każdego $x\in E$, i że w otoczeniu każdej
$\theta_0$ istnieją całkowalne względem $\mu$ funkcje dominujące
$|\partial_\theta p_\theta(x)|$ oraz $|\partial^2_\theta p_\theta(x)|$.
Dodatkowo załóżmy skończoność $\E_\theta[s_\theta(X)^2]$ i
$\E_\theta|\partial_\theta^2\log p_\theta(X)|$.
Wówczas dla każdego $\theta\in\Theta$ zachodzi
\[
\E_\theta s_\theta(X)=0,\qquad
I(\theta)=-\E_\theta[\partial_\theta^2\log p_\theta(X)].
\]
Jeśli $X_1,\ldots,X_n$ są niezależne z tym samym prawem $P_\theta$,
to $I_n(\theta)=nI(\theta)$.
\end{twierdzenie}
\begin{proof}
Z dominacji i twierdzenia o różniczkowaniu pod całką
$\int_{\mathsf X}\partial_\theta p_\theta\,\dd\mu
=\partial_\theta\int_{\mathsf X}p_\theta\,\dd\mu=0$ oraz
$\int_{\mathsf X}\partial_\theta^2p_\theta\,\dd\mu=0$.
Ponieważ $s_\theta=(\partial_\theta p_\theta)/p_\theta$, mamy
$\E_\theta s_\theta(X)=\int_{\mathsf X}\partial_\theta p_\theta\,\dd\mu=0$.
Kolejne różniczkowanie logarytmu daje
$\partial_\theta^2\log p_\theta
=(\partial_\theta^2p_\theta)/p_\theta
-((\partial_\theta p_\theta)/p_\theta)^2$.
Mnożymy przez $p_\theta$, całkujemy i stosujemy poprzednią równość.
Dla niezależnej próby score jest sumą $\sum_i s_\theta(X_i)$.
Każdy składnik ma średnią zero i skończony drugi moment, a różne składniki
są niezależne, toteż $\E(\sum_i s_\theta(X_i))^2=nI(\theta)$.
\end{proof}

\begin{przyklad}[Informacja dla Bernoulliego i modelu normalnego]
Dla pojedynczej obserwacji Bernoulliego z $p\in(0,1)$
$s_p(X)=X/p-(1-X)/(1-p)=(X-p)/(p(1-p))$, więc
$I(p)=1/(p(1-p))$. Dla $n$ obserwacji informacja wynosi
$n/(p(1-p))$. Gdy obserwujemy $X\sim\Normal(\mu,\sigma^2)$ przy znanym
$\sigma>0$, score dla $\mu$ to $(X-\mu)/\sigma^2$, toteż
$I(\mu)=1/\sigma^2$ i $I_n(\mu)=n/\sigma^2$.
Wyniki informacyjne odnoszą się do określonej parametryzacji i modelu:
nie wolno porównywać samych liczbowych wartości informacji dla parametrów
o różnych jednostkach bez uwzględnienia transformacji.
\end{przyklad}

\section{Rodziny wykładnicze i skąd biorą się proste statystyki}
\begin{definicja}[Jednoparametrowa naturalna rodzina wykładnicza]\label{def:exp-family}
Niech $J\subset\R$ będzie otwartym przedziałem, $h:\mathsf X\to[0,\infty)$
i $T:\mathsf X\to\R$ funkcjami mierzalnymi, a $\mu$ ustaloną miarą.
Załóżmy, że
$0<\int_{\mathsf X}h(x)e^{\eta T(x)}\,\dd\mu(x)<\infty$ dla wszystkich $\eta\in J$.
Definiujemy
$A(\eta)=\log\int_{\mathsf X}h(x)e^{\eta T(x)}\,\dd\mu(x)$ oraz
$p_\eta(x)=h(x)\exp\{\eta T(x)-A(\eta)\}$.
Rodzinę $\{p_\eta:\eta\in J\}$ nazywamy naturalną rodziną wykładniczą.
\end{definicja}

\begin{twierdzenie}[Pochodne funkcji normalizującej]\label{thm:exp-family-moments}
W definicji~\ref{def:exp-family} załóżmy, że $J$ jest otwarty.
Wtedy $A$ jest dwukrotnie różniczkowalna na $J$ i
$A'(\eta)=\E_\eta T(X)$, $A''(\eta)=\Var_\eta(T(X))\ge0$.
Dla niezależnej próby z tej rodziny log-wiarygodność ma postać
$\ell(\eta)=\eta\sum_iT(x_i)-nA(\eta)+\sum_i\log h(x_i)$
na zbiorze obserwacji o dodatnim $h(x_i)$.
\end{twierdzenie}
\begin{proof}
Niech $Z(\eta)=\int_{\mathsf X}h e^{\eta T}\,\dd\mu$. Dla ustalonego
$\eta\in J$ wybierzmy $\delta>0$ tak, że
$[\eta-2\delta,\eta+2\delta]\subset J$. Dla $k=1,2$ wyrażenie
$|T|^k e^{\eta T}$ jest z góry ograniczone przez stałą razy
$e^{(\eta-\delta)T}+e^{(\eta+\delta)T}$.
Po pomnożeniu przez $h$ otrzymujemy całkowalną funkcję dominującą,
więc możemy dwukrotnie różniczkować pod całką:
$Z'(\eta)=\int_{\mathsf X}T h e^{\eta T}\,\dd\mu$,
$Z''(\eta)=\int_{\mathsf X}T^2 h e^{\eta T}\,\dd\mu$.
Ponieważ $A=\log Z$, mamy $A'=Z'/Z=\E_\eta T(X)$ i
$A''=Z''/Z-(Z'/Z)^2=\Var_\eta T(X)$.
Ostatnią równość otrzymujemy przez zsumowanie logarytmów gęstości.
\end{proof}

\begin{przyklad}[Postać kanoniczna Bernoulliego]
Jeżeli $p\in(0,1)$ i $\eta=\log(p/(1-p))$, to
$\Prob_\eta(X=x)=\exp\{\eta x-\log(1+e^\eta)\}$ dla $x\in\{0,1\}$.
Zatem $T(x)=x$, $A(\eta)=\log(1+e^\eta)$, a równanie score dla
$n$ niezależnych obserwacji ma postać
$\sum_i x_i-nA'(\eta)=0$. Ponieważ
$A'(\eta)=e^\eta/(1+e^\eta)=p$, wracamy do
$\widehat p=\bar x$ dla $0<\bar x<1$.
Dla $\bar x\in\{0,1\}$ optimum występuje na granicy
$\eta\to\mp\infty$ i nie ma skończonego MLE w parametryzacji $\eta\in\R$.
\end{przyklad}

\section{Dlaczego MLE ma sens asymptotyczny? Dywergencja Kullbacka--Leiblera}
\begin{definicja}[Dywergencja Kullbacka--Leiblera]\label{def:kl}
Niech $P$ i $Q$ będą miarami probabilistycznymi z gęstościami $p$ i $q$
względem wspólnej miary $\mu$. Jeśli $P\ll Q$, definiujemy
$D_{\mathrm{KL}}(P\|Q)=\int_{\mathsf X}p(x)\log\frac{p(x)}{q(x)}\,\dd\mu(x)$,
z konwencją, że całkujemy po zbiorze $\{p>0\}$, a iloraz
$p/q$ ma znaczenie tylko tam, gdzie $q>0$; $0\log0$ interpretujemy jako $0$.
Jeśli $P$ nie jest absolutnie ciągła względem $Q$, definiujemy
$D_{\mathrm{KL}}(P\|Q)=+\infty$.
\end{definicja}
\begin{twierdzenie}[Nieujemność i znaczenie dywergencji KL]\label{thm:kl-basic}
Dla dowolnych miar probabilistycznych $P,Q$ zachodzi
$D_{\mathrm{KL}}(P\|Q)\in[0,\infty]$. Równość z zerem zachodzi wtedy
i tylko wtedy, gdy $P=Q$. Jeżeli $P_{\theta_0}$ i $P_\theta$ mają
odpowiednie gęstości, a $\log p_{\theta_0}(X)$ oraz
$\log p_\theta(X)$ są całkowalne przy $P_{\theta_0}$, to
\[
\E_{\theta_0}\log p_\theta(X)
=\E_{\theta_0}\log p_{\theta_0}(X)
-D_{\mathrm{KL}}(P_{\theta_0}\|P_\theta).
\]
\end{twierdzenie}
\begin{proof}
W przypadku $P\ll Q$ zdefiniujmy pochodną $r=\dd P/\dd Q$.
Funkcja $\phi(t)=t\log t-t+1$ dla $t\ge0$ jest nieujemna;
osiąga zero tylko w $t=1$ (pochodna dla $t>0$ to $\log t$, a druga
pochodna $1/t>0$). Ponieważ $\int_{\mathsf X}r\,\dd Q=1$, mamy
$\int_{\mathsf X}\phi(r)\,\dd Q=\int_{\mathsf X}r\log r\,\dd Q=D_{\mathrm{KL}}(P\|Q)\ge0$.
Jeżeli całka jest zerowa, $\phi(r)=0$ $Q$-prawie wszędzie,
czyli $r=1$ $Q$-prawie wszędzie i $P=Q$.
Dla $P\not\ll Q$ dywergencja jest nieskończona z definicji.
Ostatnią równość otrzymujemy przez odjęcie obu całkowalnych logarytmów
i zastosowanie definicji $D_{\mathrm{KL}}$.
\end{proof}

\emph{Intuicja.} Przy prawdziwym parametrze $\theta_0$ oczekiwany
logarytm wiarygodności jest największy (przy identyfikowalności jednoznacznie)
dla $\theta=\theta_0$. Z prawa wielkich liczb dla \emph{ustalonego}
$\theta$ mamy zbieżność uśrednionej log-wiarygodności. To jednak nie
wystarcza do zgodności MLE: maksimum jest wyznaczane po \emph{wszystkich}
parametrach naraz. Potrzebne jest odpowiednie \emph{jednostajne}
kontrolowanie log-wiarygodności.

\begin{twierdzenie}[Elementarne twierdzenie o zgodności maksymalizatora]\label{thm:argmax-consistency}
Niech $(\Theta,d)$ będzie zwartą przestrzenią metryczną,
$M:\Theta\to\R$ funkcją ciągłą o jedynym maksimum w $\theta_0$.
Niech $M_n:\Theta\to\R$ będą losowymi funkcjami ciągłymi, dla których
$\sup_{\theta\in\Theta}|M_n(\theta)-M(\theta)|\to0$ prawie na pewno.
Jeśli $\widehat\theta_n$ są mierzalnymi maksymalizatorami $M_n$, to
$\widehat\theta_n\to\theta_0$ prawie na pewno.
\end{twierdzenie}
\begin{proof}
Ustalmy $\varepsilon>0$. Zbiór
$K_\varepsilon=\{\theta\in\Theta:d(\theta,\theta_0)\ge\varepsilon\}$
jest zwarty. Jeżeli jest pusty, twierdzenie dla tego $\varepsilon$
jest trywialne. W przeciwnym razie, przez jednoznaczność maksimum
i ciągłość $M$, liczba
$\gamma=M(\theta_0)-\max_{\theta\in K_\varepsilon}M(\theta)$
jest dodatnia. Na zdarzeniu zbieżności jednostajnej dla dostatecznie
dużego $n$ zachodzi $\sup|M_n-M|<\gamma/3$.
Wówczas $M_n(\theta_0)>M(\theta_0)-\gamma/3$, natomiast dla
$\theta\in K_\varepsilon$ mamy
$M_n(\theta)<M(\theta_0)-2\gamma/3$.
Maksymalizator $M_n$ nie może więc należeć do $K_\varepsilon$.
Ponieważ $\varepsilon>0$ było dowolne, otrzymujemy zbieżność.
\end{proof}

\begin{przyklad}[Zgodność MLE dla Bernoulliego wewnątrz parametrycznego przedziału]
Niech $p_0\in[\delta,1-\delta]$, gdzie $0<\delta<1/2$, oraz
$\Theta=[\delta,1-\delta]$. Dla próby Bernoulliego
$M_n(p)=n^{-1}\ell(p)=\bar X_n\log p+(1-\bar X_n)\log(1-p)$.
Prawo wielkich liczb daje $\bar X_n\to p_0$ prawie na pewno, a
\[
\sup_{p\in\Theta}|M_n(p)-M(p)|
\le |\bar X_n-p_0|
\sup_{p\in\Theta}\left|\log\frac p{1-p}\right|\longrightarrow0,
\]
gdzie $M(p)=p_0\log p+(1-p_0)\log(1-p)$.
Ścisła wklęsłość $M$ i obliczenie pochodnej dają jedyne maksimum
w $p_0$. Poprzednie twierdzenie dowodzi zatem mocnej zgodności
ograniczonego MLE. Nie użyliśmy tu nieudowodnionego ogólnego twierdzenia
o zbieżności jednostajnej: oszacowanie jest jawne.
\end{przyklad}

\section{Parametry zakłócające, profilowanie i niezmienniczość}
\begin{definicja}[Wiarygodność profilowa]
Niech $\theta=(\psi,\lambda)\in\Psi\times\Lambda$, gdzie $\psi$ jest
parametrem zainteresowania, a $\lambda$ parametrem zakłócającym.
Definiujemy $L_{\mathrm{prof}}(\psi;x)=\sup_{\lambda\in\Lambda}
L_x(\psi,\lambda)$, jeśli supremum jest dobrze określone.
Nie traktujemy $L_{\mathrm{prof}}$ jako rozkładu prawdopodobieństwa na $\Psi$.
\end{definicja}

\begin{przyklad}[Profilowanie wariancji w modelu normalnym]
Dla obserwacji $x$ takich, że $Q(\mu)=\sum_i(x_i-\mu)^2>0$
dla każdej $\mu$, maksimum wiarygodności względem wariancji przy ustalonym
$\mu$ osiąga się w $\widehat v_\mu=Q(\mu)/n$. Po podstawieniu
otrzymujemy $L_{\mathrm{prof}}(\mu;x)\propto Q(\mu)^{-n/2}$.
Ponieważ $Q(\mu)$ jest najmniejsze w $\bar x$, maksimum profilowe
wypada w $\mu=\bar x$. Profilowanie zachowuje w tym przypadku MLE $\mu$,
lecz rozkład statystyki ilorazu wiarygodności wymaga odrębnej teorii.
\end{przyklad}

\begin{stwierdzenie}[Niezmienniczość MLE względem zmiany parametryzacji]
Niech $g:\Theta\to\Phi$ będzie bijekcją, a $\widehat\theta$ maksymalizatorem
$L_x$ w $\Theta$. W parametryzacji $\varphi=g(\theta)$ funkcja wiarygodności
ma postać $\widetilde L_x(\varphi)=L_x(g^{-1}(\varphi))$.
Wtedy $g(\widehat\theta)$ jest maksymalizatorem $\widetilde L_x$.
\end{stwierdzenie}
\begin{proof}
Dla każdego $\varphi\in\Phi$ istnieje $\theta=g^{-1}(\varphi)$.
Z definicji maksimum
$\widetilde L_x(\varphi)=L_x(\theta)\le L_x(\widehat\theta)
=\widetilde L_x(g(\widehat\theta))$.
\end{proof}

\begin{uwaga}[Praktyczna pułapka: transformacja parametru]
Przy zastąpieniu parametru $\sigma$ parametrem $\sigma^2$ \emph{nie}
mnożymy wiarygodności przez Jacobian zmiany zmiennej: jest to zmiana
\emph{parametryzacji modelu}, a nie transformacja gęstości rozkładu
losowego parametru. Jacobian pojawiłby się natomiast przy przekształcaniu
gęstości \emph{prioru} w teorii bayesowskiej.
\end{uwaga}

\section{Jak dobierać narzędzie w praktyce? Porównania i ograniczenia}
\begin{center}\small
\begin{tabular}{>{\raggedright\arraybackslash}p{0.26\textwidth}>{\raggedright\arraybackslash}p{0.32\textwidth}>{\raggedright\arraybackslash}p{0.30\textwidth}}
\toprule
\textbf{Problem} & \textbf{Metoda / pierwsze obliczenie} & \textbf{Co koniecznie sprawdzić?}\\
\midrule
Odsetek wad & $\widehat p=k/n$ z modelu Bernoulliego & Zależność w partiach; reprezentatywność\\
Liczba awarii w różnych okresach & $\widehat\lambda=\sum_i y_i/\sum_i t_i$ & Ekspozycję; nadmierną dyspersję\\
Pomiar około stałego poziomu & MLE średniej i wariancji w modelu normalnym & Ogony; odstające pomiary; dopasowanie rozkładu\\
Nieznany górny kres rozkładu & Maksimum próby w modelu $U(0,\theta)$ & Czy rzeczywiście występuje ostry kres nośnika?\\
Parametry zakłócające & Profilowanie lub wspólna optymalizacja & Wielomodalność; różnice identyfikowalności\\
Optimum na brzegu & Analiza ograniczeń, nie tylko równania score & Czy maksimum jest osiągane?\\
\bottomrule
\end{tabular}
\end{center}

\paragraph{Praktyczna sekwencja obliczeń.}
Należy wskazać jednostkę obserwacji i parametry, zapisać \emph{pełną}
wiarygodność z ograniczeniami na $\theta$, a następnie (i) sprawdzić
punkty stacjonarne, (ii) brzegi i zależność nośnika od parametru,
(iii) istnienie maksimum, (iv) wpływ specyfikacji modelu na wynik.
Jeśli dane są zależne, nie wolno bez uzasadnienia zastąpić ich prawa
wspólnego iloczynem rozkładów marginalnych.

\Needspace{16\baselineskip}
\noindent\textbf{Przykład numeryczny w Pythonie.}\par
\par
\Needspace{13\baselineskip}
\smallskip
\begin{lstlisting}[language=Python,caption={Wiarygodność modelu Bernoulliego i iloraz wiarygodności (Python).}]
import numpy as np
k, n = 7, 100
p_hat = k / n
logL = lambda p: k * np.log(p) + (n-k) * np.log1p(-p)
log_ratio = logL(0.07) - logL(0.05)
print(p_hat, np.exp(log_ratio))
# 0.07, ok. 1.457...
\end{lstlisting}
\Needspace{16\baselineskip}
\noindent\textbf{To samo w Julii (bez dodatkowych pakietów).}\par
\par
\Needspace{13\baselineskip}
\smallskip
\begin{lstlisting}[language=Julia,caption={Wiarygodność Bernoulliego i intensywność Poissona (Julia).}]
k, n = 7, 100
p_hat = k / n
logL(p) = k * log(p) + (n-k) * log1p(-p)
println((p_hat, exp(logL(0.07) - logL(0.05))))
y, t = [2, 5, 7], [10.0, 20.0, 40.0]
println(sum(y) / sum(t))  # intensywnosc: 0.2 na godzine
\end{lstlisting}

Funkcja używa $\log(1-p)$ w wariancie numerycznie stabilniejszym
$\operatorname{log1p}(-p)$. Dla obserwacji brzegowych potrzebujemy
osobnej obsługi $p=0$ i $p=1$, aby nie obliczać wyrażeń typu $0\cdot\log0$.

% ================================================================
% ROZDZIAL 4 -- STATYSTYCZNA TEORIA DECYZJI
% ================================================================
\chapter{Statystyczna teoria decyzji: od danych do działania}
\label{ch:decision}
\paragraph{Powiązania i odsyłacze.} Różnice między regułą klasyczną a bayesowską stają się szczególnie istotne w rozdziałach~\ref{ch:bayes-foundations} i~\ref{ch:bayes-estimation-prediction}.


\section{Dlaczego estymator nie jest jeszcze decyzją?}
Poznanie nieznanego parametru jest zwykle środkiem do celu: musimy
zdecydować, czy zatrzymać linię produkcyjną, wymienić urządzenie,
wdrożyć określony wariant leczenia lub zakwalifikować transakcję do
weryfikacji. Dwa sposoby szacowania parametru mogą prowadzić do
różnych decyzji, a ten sam błąd liczbowy może mieć bardzo różną cenę
w zależności od kierunku pomyłki.

\begin{przyklad}[Awaria i koszt decyzji]
Rozważmy dwie pomyłki: niepotrzebne zatrzymanie sprawnego urządzenia
kosztuje $100$ jednostek, zaś zignorowanie groźnej awarii kosztuje
$5000$ jednostek. Reguła minimalizująca \emph{liczbę błędnych decyzji}
nie musi minimalizować \emph{oczekiwany koszt}. Właśnie dlatego
w statystyce trzeba jawnie sformułować funkcję straty.
\end{przyklad}

\begin{definicja}[Problem decyzji statystycznej]\label{def:decision-problem}
Niech $\mathcal E=(\mathsf X,\calA,\{P_\theta:\theta\in\Theta\})$ będzie
eksperymentem statystycznym, $(\mathsf A,\mathcal C)$ przestrzenią
mierzalną \emph{działań}, a $L:\Theta\times\mathsf A\to[0,\infty]$
\emph{funkcją straty}; zakładamy mierzalność $a\mapsto L(\theta,a)$
dla każdego $\theta\in\Theta$. \emph{Nielosową regułą decyzyjną}
nazywamy odwzorowanie mierzalne $\delta:\mathsf X\to\mathsf A$.
Jej \emph{funkcją ryzyka} jest
\[
R(\theta,\delta)
=\E_\theta[L(\theta,\delta(X))]
=\int_{\mathsf X}L(\theta,\delta(x))\,\dd P_\theta(x),
\qquad\theta\in\Theta.
\]
Dopuszczamy wartość $+\infty$. W szczególności reguła estymacji
$\delta(X)\in\Theta$ jest przypadkiem decyzji, ale przestrzeń działań
nie musi być równa przestrzeni parametrów.
\end{definicja}

\begin{uwaga}[Na czym polega częstotliwościowy charakter ryzyka?]
Dla każdego ustalonego $\theta$ liczymy wartość oczekiwaną względem
losowości \emph{danych} $X\sim P_\theta$. Wynikiem jest funkcja
$\theta\mapsto R(\theta,\delta)$, a nie jedna liczba niezależna od
nieznanego parametru. Stąd różne zasady wyboru reguły: można
porównywać całe funkcje ryzyka, minimalizować ryzyko najgorszego
przypadku albo uśrednić ryzyko względem jawnie przyjętego rozkładu
parametru. Każda odpowiada na inne pytanie.
\end{uwaga}

\begin{definicja}[Ważne funkcje straty]
Dla rzeczywistego parametru $\theta$ i decyzji liczbowej $a\in\R$
rozważamy stratę kwadratową $L_2(\theta,a)=(a-\theta)^2$ i absolutną
$L_1(\theta,a)=|a-\theta|$. Dla decyzji binarnej
$\theta,a\in\{0,1\}$ strata błędu klasyfikacji ma postać
$L_{01}(\theta,a)=\ind{\{\theta\ne a\}}$.
Stratę asymetryczną definiujemy, ustalając koszty
$c_{\mathrm{FP}},c_{\mathrm{FN}}>0$: wartość to
$c_{\mathrm{FP}}$ przy $(\theta,a)=(0,1)$,
$c_{\mathrm{FN}}$ przy $(1,0)$ i zero przy poprawnej decyzji.
\end{definicja}

\section{Strata kwadratowa, ryzyko i kompromis obciążenia z wariancją}
\begin{twierdzenie}[Dekompozycja ryzyka kwadratowego]\label{thm:decision-mse}
Niech $\theta\in\R$, $T:\mathsf X\to\R$ będzie statystyką taką, że
$\E_\theta T^2<\infty$, a $\delta_T(x)=T(x)$ regułą estymacyjną.
Dla straty $L_2$ zachodzi
\[
R(\theta,\delta_T)
=\Var_\theta(T)+(\E_\theta T-\theta)^2.
\]
W szczególności mniejsze obciążenie nie musi oznaczać mniejszego ryzyka.
\end{twierdzenie}
\begin{proof}
Rozkładamy błąd jako
$T-\theta=(T-\E_\theta T)+(\E_\theta T-\theta)$.
Po podniesieniu do kwadratu i wzięciu wartości oczekiwanej
wyraz mieszany znika. Pozostałe składniki to odpowiednio wariancja
i kwadrat obciążenia. Warunek drugiego momentu zapewnia skończoność
każdego użytego składnika. Jest to tożsamość już poznana w rozdziale 2,
teraz interpretowana jako funkcja ryzyka reguły decyzyjnej.
\end{proof}

\begin{przyklad}[Dwie reguły estymacji prawdopodobieństwa sukcesu]
Mamy tylko jedną obserwację $X\sim\operatorname{Bernoulli}(p)$,
$p\in[0,1]$. Reguła $\delta_1(X)=X$ jest nieobciążona i ma
$R(p,\delta_1)=p(1-p)$. Reguła stała $\delta_0(X)=1/2$ ma
$R(p,\delta_0)=(p-1/2)^2$.
Gdy $p=1/2$, druga reguła ma ryzyko zero, podczas gdy pierwsza ma
ryzyko $1/4$. Gdy $p=0$ lub $p=1$, pierwsza reguła ma ryzyko zero,
a druga $1/4$. \emph{Nie ma więc dominacji jednej reguły nad drugą}.
Bez wskazania kryterium nie istnieje jedna odpowiedź na pytanie,
która z nich jest „najlepsza”.
\end{przyklad}

\begin{definicja}[Dominacja i dopuszczalność]\label{def:admissible}
Niech $\delta_1,\delta_2$ będą regułami o ryzykach określonych
na $\Theta$. Mówimy, że $\delta_1$ \emph{dominuje} $\delta_2$,
jeśli $R(\theta,\delta_1)\le R(\theta,\delta_2)$ dla każdego
$\theta\in\Theta$ oraz nierówność jest ostra dla co najmniej jednego
$\theta$. Reguła $\delta$ jest \emph{dopuszczalna}, jeśli w rozważanej
klasie reguł nie istnieje reguła ją dominująca.
\end{definicja}

\begin{uwaga}[Interpretacja geometryczna]
Możemy wyobrazić sobie wykresy $\theta\mapsto R(\theta,\delta)$.
Dominacja oznacza, że jeden wykres leży wszędzie nie wyżej od
drugiego, a gdzieś ściśle niżej. Krzyżowanie się wykresów oznacza,
że preferencja zależy od wartości nieznanego parametru.
\end{uwaga}

\section{Reguła minimaksowa: zabezpieczenie przed najgorszym przypadkiem}
\begin{definicja}[Ryzyko maksymalne i minimaksowość]\label{def:minimax}
Dla reguły decyzyjnej $\delta$ definiujemy
$R^{\max}(\delta)=\sup_{\theta\in\Theta}R(\theta,\delta)$.
Jeśli reguła $\delta^*$ spełnia
$R^{\max}(\delta^*)=\inf_\delta R^{\max}(\delta)$, to nazywamy ją
\emph{minimaksową} w danej klasie dopuszczalnych reguł.
Dla ogólnej przestrzeni $\Theta$ ani istnienie, ani jednoznaczność
takiej reguły nie są automatyczne.
\end{definicja}

\begin{twierdzenie}[Dokładna reguła minimaksowa z jednej próby Bernoulliego]
\label{thm:minimax-bernoulli-one}
Niech $X\sim\operatorname{Bernoulli}(p)$, $p\in[0,1]$.
Szacujemy $p$ za pomocą dowolnej mierzalnej reguły
$\delta:\{0,1\}\to\R$ przy stracie kwadratowej.
Wtedy reguła
\[
\delta^*(0)=\frac14,\qquad \delta^*(1)=\frac34
\]
ma stałe ryzyko $R(p,\delta^*)=1/16$ dla każdego $p\in[0,1]$
i jest minimaksowa w klasie wszystkich takich reguł.
Ten sam dolny kres obowiązuje także dla reguł losowych, jeżeli
losowanie dodatkowe jest niezależne od parametru przy ustalonych danych.
\end{twierdzenie}
\begin{proof}
Najpierw obliczamy bezpośrednio
\[
R(p,\delta^*)=(1-p)(\tfrac14-p)^2+p(\tfrac34-p)^2
=\frac1{16},
\]
gdyż po rozwinięciu składniki zawierające $p$, $p^2$ i $p^3$
wzajemnie się redukują. To daje górne ograniczenie wartości
minimaksowej przez $1/16$.

Potrzebujemy \emph{dolnego} ograniczenia dla dowolnej konkurencyjnej
reguły. Wykorzystamy do tego pomocniczy rozkład na przestrzeni
parametrów: niech $\Pi$ ma rozkład
$\operatorname{Beta}(1/2,1/2)$, a warunkowo na $\Pi=p$ niech
$X\sim\operatorname{Bernoulli}(p)$. Jest to sposób ważenia ryzyk,
nie zmiana definicji częstotliwościowego ryzyka $R(p,\delta)$.
Z rachunku rozkładu beta dostajemy $\Prob(X=1)=\E\Pi=1/2$.
Dla $x=0,1$ warunkowy rozkład $\Pi$ po ustaleniu $X=x$ ma gęstość
proporcjonalną do
$p^{x-1/2}(1-p)^{1-x-1/2}$, czyli jest równy
$\operatorname{Beta}(x+1/2,3/2-x)$.
Jego średnia wynosi $(x+1/2)/2$, a wariancja
\[
\Var(\Pi\mid X=x)
=\frac{(x+1/2)(3/2-x)}{2^2(2+1)}
=\frac1{16}.
\]
Dla każdej reguły $\delta$ zachodzi warunkowa tożsamość Pitagorasa
\[
\E[(\delta(X)-\Pi)^2\mid X]
=(\delta(X)-\E[\Pi\mid X])^2+\Var(\Pi\mid X)
\ge\frac1{16}.
\]
Po wzięciu wartości oczekiwanej i użyciu Tonellego otrzymujemy
$\int_0^1R(p,\delta)\,\dd\Pi(p)\ge1/16$.
Zawsze
$\sup_pR(p,\delta)\ge\int_{p\in[0,1]}R(p,\delta)\,\dd\Pi(p)$,
zatem $R^{\max}(\delta)\ge1/16$.
Dodatkowe losowanie nie może obniżyć warunkowego błędu kwadratowego:
dla losowej akcji $A$ i danej $X$ nierówność Jensena daje
$\E[(A-\Pi)^2\mid X,\Pi]\ge(\E[A\mid X]-\Pi)^2$,
jeśli losowanie akcji warunkowo na $X$ jest niezależne od $\Pi$.
To kończy dowód także dla reguł losowych.
\end{proof}

\emph{Co jest zaskakujące?} MLE $\widehat p=X$ przy jednej obserwacji
jest nieobciążony, ale jego największe ryzyko wynosi $1/4$.
Reguła $\delta^*$ jest obciążona, jednak zmniejsza największy błąd
średniokwadratowy \emph{czterokrotnie}. Nie oznacza to, że zawsze jest
lepsza: przy $p=0$ lub $p=1$ MLE ma ryzyko zero. Reguła minimaksowa
chroni przed najgorszym przypadkiem, nie jest zaś punktowo najlepsza.

\begin{center}
\begin{tikzpicture}[x=10cm,y=9cm,>=Latex]
  \draw[->] (0,0)--(1.07,0) node[right] {$p$};
  \draw[->] (0,0)--(0,.30) node[above] {$R(p,\delta)$};
  \foreach \x/\lab in {0/{0},.25/{1/4},.5/{1/2},.75/{3/4},1/{1}} {
    \draw (\x,-.004)--(\x,.004) node[below=3pt] {\scriptsize $\lab$};
  }
  \foreach \y/\lab in {.0625/{1/16},.25/{1/4}} {
    \draw (-.005,\y)--(.005,\y) node[left=3pt] {\scriptsize $\lab$};
  }
  \draw[very thick,blue!70!black,domain=0:1,samples=100]
    plot (\x,{\x*(1-\x)});
  \draw[very thick,red!65!black,domain=0:1]
    plot (\x,{.0625});
  \node[blue!70!black,anchor=west] at (.63,.22) {\scriptsize $\delta(X)=X$};
  \node[red!65!black,anchor=west] at (.58,.077) {\scriptsize $\delta^*(X)$};
\end{tikzpicture}
\end{center}
\noindent\emph{Rysunek 4.1.} Ryzyko kwadratowe MLE i reguły
minimaksowej dla jednej próby Bernoulliego. Reguła minimaksowa ma
mniejsze ryzyko maksymalne, ale nie jest najlepsza przy każdym $p$.


\section{Ryzyko uśrednione i zasada wyboru przy danej wiedzy a priori}
\begin{definicja}[Ryzyko Bayesa w teorii decyzji]\label{def:bayes-risk}
Niech $(\Theta,\mathcal T)$ będzie przestrzenią mierzalną,
$\pi$ miarą probabilistyczną na niej. Zakładamy, że
$\theta\mapsto P_\theta(B)$ jest $\mathcal T$-mierzalna dla
każdego $B\in\calA$ (a więc $(P_\theta)$ stanowi jądro
probabilistyczne) oraz że $L:\Theta\times\mathsf A\to[0,\infty]$
jest mierzalna względem $\mathcal T\otimes\mathcal C$.
Wtedy $\theta\mapsto R(\theta,\delta)$ jest mierzalna;
wynika to z twierdzenia o klasie monotonicznej, najpierw dla
wskaźników prostokątów, a następnie z aproksymacji dodatnimi
funkcjami prostymi i zbieżności monotonicznej.
\emph{Ryzykiem Bayesa} reguły $\delta$ nazywamy
$\rho_\pi(\delta)=\int_\Theta R(\theta,\delta)\,\dd\pi(\theta)$.
Reguła minimalizująca to ryzyko jest \emph{regułą Bayesa względem $\pi$}.
W dalszej części tomu IV wprowadzimy pełną teorię rozkładów
a priori i a posteriori. Tutaj rozpatrujemy tylko niezbędną
konstrukcję decyzyjną.
\end{definicja}

\begin{stwierdzenie}[Punktowa minimalizacja warunkowej straty]
\label{prop:conditional-decision}
Niech $\Theta$ i $\mathsf X$ będą skończone, $\pi$ rozkładem na
$\Theta$, a $P_\theta$ prawami obserwacji na $\mathsf X$.
Niech $\mathsf A$ będzie skończonym zbiorem akcji i
$L:\Theta\times\mathsf A\to[0,\infty)$ funkcją straty.
Dla $x$ o dodatnim prawdopodobieństwie brzegowym
$w_x=\sum_{\theta}\pi(\theta)P_\theta(\{x\})$ zdefiniujmy
$q_x(\theta)=\pi(\theta)P_\theta(\{x\})/w_x$.
Reguła $\delta$ minimalizuje ryzyko Bayesa wtedy i tylko wtedy,
gdy dla każdego $x$ z $w_x>0$ wybiera akcję minimalizującą
$\sum_{\theta}L(\theta,a)q_x(\theta)$.
\end{stwierdzenie}
\begin{proof}
Rozwijamy skończoną sumę i grupujemy składniki według obserwacji:
\[
\rho_\pi(\delta)
=\sum_{\theta,x}\pi(\theta)P_\theta(\{x\})L(\theta,\delta(x))
=\sum_{x:w_x>0}w_x\sum_{\theta}q_x(\theta)L(\theta,\delta(x)).
\]
Każdy współczynnik $w_x$ jest dodatni, a wybór $\delta(x)$ dla
jednego $x$ nie ogranicza wyboru $\delta(y)$ dla $y\ne x$.
Zatem suma jest minimalna dokładnie wtedy, gdy minimalny jest każdy
jej niezerowy składnik. Na obserwacjach o $w_x=0$ regułę
można określić dowolnie.
\end{proof}

\begin{uwaga}[Kiedy której reguły używać?]
Ryzyko Bayesa zależy od wagi $\pi$, a minimaksowe od najgorszej
wartości parametru. Priory zbyt arbitralne lub niewiarygodne mogą
prowadzić do decyzji nieodpowiednich dla celu. Jeśli nie umiemy
uzasadnić żadnego rozkładu a priori, nie wolno przedstawiać
reguły Bayesa jako decyzji wynikającej wyłącznie z danych.
Jednocześnie zasadzie minimaksowej może towarzyszyć nadmierna
ostrożność w typowych sytuacjach.
\end{uwaga}

\section{Od testu diagnostycznego do decyzji: koszty błędów i częstość bazowa}
\begin{twierdzenie}[Optymalny próg dla dwóch stanów]\label{thm:decision-lr}
Niech nieznany stan $\Theta\in\{0,1\}$ ma w modelu decyzyjnym
prawdopodobieństwo a priori $\Prob(\Theta=1)=q\in(0,1)$.
Niech $X$ ma warunkową funkcję masy $f_j(x)$ przy $\Theta=j$
na wspólnym przeliczalnym zbiorze obserwacji $\mathsf X$.
Możliwe akcje to $a=0$ i $a=1$; koszt pomyłki przy
$(\Theta,a)=(0,1)$ wynosi $c_{\mathrm{FP}}>0$, a przy $(1,0)$
wynosi $c_{\mathrm{FN}}>0$; poprawne decyzje kosztują zero.
Dla $x$ z dodatnią mieszaną masą
$(1-q)f_0(x)+qf_1(x)$ reguła minimalizująca ryzyko Bayesa
wybiera $a=1$ wtedy i tylko wtedy, gdy
\[
 c_{\mathrm{FN}}\,q f_1(x)
>c_{\mathrm{FP}}\,(1-q) f_0(x),
\]
a przy równości obie akcje są optymalne. Dla $f_0(x)>0$ jest to
równoważne nierówności ilorazu wiarygodności
$f_1(x)/f_0(x)>c_{\mathrm{FP}}(1-q)/(c_{\mathrm{FN}}q)$.
\end{twierdzenie}
\begin{proof}
Po obserwacji $X=x$ warunkowy oczekiwany koszt akcji $a=1$ to
$c_{\mathrm{FP}}\Prob(\Theta=0\mid X=x)$; dla $a=0$ jest to
$c_{\mathrm{FN}}\Prob(\Theta=1\mid X=x)$.
Dzieląc przez dodatni mianownik rozkładu warunkowego, widzimy, że
pierwszy koszt jest mniejszy od drugiego dokładnie przy podanej
nierówności. Z poprzedniego stwierdzenia minimalizacja warunkowa dla
każdej wartości $x$ minimalizuje całkowite ryzyko Bayesa.
\end{proof}

\begin{zadanie}[Rzadkie zdarzenie i dodatni wynik testu]\label{zad:rare-classification}
W populacji częstość groźnej usterki to $q=0{,}01$.
Test wykrywa usterkę z czułością $0{,}90$ i ma swoistość $0{,}95$.
Dodatni wynik pojawił się na sprawdzanym urządzeniu. Porównaj decyzję
„interweniować” przy jednakowych kosztach pomyłek z sytuacją,
gdy przeoczenie usterki kosztuje dziesięć razy więcej niż fałszywy alarm.
\par\emph{Rozwiązanie.}
Dla wyniku dodatniego $f_1(+)=0{,}90$ i $f_0(+)=0{,}05$.
Iloraz wiarygodności wynosi $18$, lecz stosując rachunek warunkowy,
\[
\Prob(\Theta=1\mid +)
=\frac{0{,}01\cdot0{,}90}
 {0{,}01\cdot0{,}90+0{,}99\cdot0{,}05}
=\frac{0{,}009}{0{,}0585}\approx0{,}154.
\]
Przy jednakowych kosztach błędów mniejszy oczekiwany koszt ma
brak interwencji, ponieważ prawdopodobieństwo usterki jest mniejsze niż
$1/2$. Gdy $c_{\mathrm{FN}}=10$ i $c_{\mathrm{FP}}=1$,
próg ilorazu wiarygodności wynosi
$0{,}99/(10\cdot0{,}01)=9{,}9$, więc $18>9{,}9$ i należy
interweniować. Równoważnie koszt zignorowania usterki po wyniku
pozytywnym wynosi około $10\cdot0{,}154=1{,}54$, podczas gdy koszt
fałszywego alarmu to około $1\cdot0{,}846=0{,}846$.
\emph{Wniosek:} jeden wynik badania może uzasadniać różne działania przy
różnych kosztach. Test nie jest decyzją; jest źródłem informacji dla decyzji.
\end{zadanie}

\section{Wybór straty: średnia, mediana i kwantyle}
\begin{stwierdzenie}[Najlepsza stała przy różnych stratach]
\label{prop:optimal-constant-loss}
Niech $Y$ będzie całkowalną rzeczywistą zmienną losową.
Jeśli $Y\in L^2$, jedyną stałą $a\in\R$ minimalizującą
$\E(Y-a)^2$ jest $a=\E Y$.
Zbiór stałych minimalizujących $\E|Y-a|$ jest zbiorem median
rozkładu $Y$, tj. zbiorem tych $a$, dla których
$\Prob(Y\le a)\ge1/2$ i $\Prob(Y\ge a)\ge1/2$.
\end{stwierdzenie}
\begin{proof}
Dla $Y\in L^2$ tożsamość
$\E(Y-a)^2=\Var(Y)+(a-\E Y)^2$ dowodzi pierwszej tezy.
Dla straty absolutnej funkcja $g(a)=\E|Y-a|$ jest skończona
i wypukła. Różnicę ilorazową
$\bigl(|Y-a-h|-|Y-a|\bigr)/h$ dla $h>0$
można ograniczyć co do modułu przez $1$.
Twierdzenie o zbieżności zdominowanej daje pochodną prawostronną
$g'_+(a)=\Prob(Y\le a)-\Prob(Y>a)=2F_Y(a)-1$.
Analogicznie pochodna lewostronna wynosi
$g'_-(a)=\Prob(Y<a)-\Prob(Y\ge a)=2F_Y(a-)-1$.
Dla funkcji wypukłej minimum występuje dokładnie wtedy, gdy
$g'_-(a)\le0\le g'_+(a)$, co jest równoważne warunkom mediany.
\end{proof}


\begin{stwierdzenie}[Asymetryczna strata absolutna prowadzi do kwantyla]
\label{prop:pinball-quantile}
Niech $Y\in L^1$ będzie rzeczywistą zmienną losową, a $\tau\in(0,1)$.
Dla decyzji $a\in\R$ określmy
$L_\tau(Y,a)=\tau(Y-a)_++(1-\tau)(a-Y)_+$.
Wtedy zbiór minimizatorów $a\mapsto\E L_\tau(Y,a)$
składa się dokładnie z tych $a$, które spełniają
$P(Y<a)\le\tau\le P(Y\le a)$.
\end{stwierdzenie}
\begin{proof}
Funkcja $g(a)=\E L_\tau(Y,a)$ jest skończona, ponieważ
$Y\in L^1$, i wypukła jako średnia funkcji wypukłych.
Dla każdego $h\ne0$ bezwzględna wartość ilorazu
$[L_\tau(Y,a+h)-L_\tau(Y,a)]/h$ jest ograniczona przez $1$.
Dla $h\downarrow0$ pochodna punktowa względem $a$ równa się
$1-\tau$ na $\{Y\le a\}$ i $-\tau$ na $\{Y>a\}$.
Zbieżność zdominowana daje zatem $g'_+(a)=F_Y(a)-\tau$.
Analogicznie $g'_-(a)=F_Y(a-)-\tau$.
Wypukła funkcja rzeczywista przyjmuje minimum w $a$ wtedy i tylko
wtedy, gdy $g'_-(a)\le0\le g'_+(a)$. Po podstawieniu obu
pochodnych otrzymujemy wskazane nierówności dla dystrybuanty.
\end{proof}

\begin{przyklad}[Zapas magazynowy: koszt nadmiaru i niedoboru]
Jeśli popyt $Y$ jest nieujemny, koszt każdej niezaspokojonej
jednostki wynosi $c_u>0$, a nadmiarowej $c_o>0$, oczekiwany koszt
przy zamówieniu $a$ jest równy
$c_u\E(Y-a)_++c_o\E(a-Y)_+$. Jest to dodatnia wielokrotność
straty z poprzedniego stwierdzenia, z poziomem
$\tau=c_u/(c_u+c_o)$. Przy $c_u=9c_o$ optymalny zapas wyznacza
\emph{kwantyl rzędu $0{,}9$}, a nie średnia ani mediana popytu.
Wynik wymaga przyjętego modelu kosztów; przy ograniczeniu
magazynowym rozwiązanie trzeba dodatkowo ograniczyć do
zbioru działań dopuszczalnych.
\end{przyklad}

\begin{uwaga}[Nie ma jednej „najlepszej statystyki położenia”]
Jeżeli model strat nagradza dokładność średniokwadratową, naturalnym
celem jest średnia warunkowa. Jeżeli koszt rośnie liniowo z
bezwzględnym błędem, naturalna jest mediana warunkowa.
Przy asymetrycznych kosztach niedoszacowania i przeszacowania
pojawiają się kwantyle. Przykład zarządzania zapasami: jeśli brak
towaru jest wielokrotnie droższy niż jego nadmiar, poziom zapasu
powinien leżeć powyżej mediany prognozowanego popytu.
\end{uwaga}

\section{Kiedy stosować poszczególne kryteria?}
\begin{center}\small
\begin{tabular}{>{\raggedright\arraybackslash}p{0.25\textwidth}>{\raggedright\arraybackslash}p{0.32\textwidth}>{\raggedright\arraybackslash}p{0.31\textwidth}}
\toprule
\textbf{Cel} & \textbf{Kryterium / reguła} & \textbf{Ograniczenie i alternatywa}\\
\midrule
Dobre oszacowanie w błędzie kwadratowym & MSE; porównanie obciążenia i wariancji & Estymator obciążony może być lepszy; zależy od $\theta$\\
Ochrona najgorszego przypadku & Minimax: $\inf_\delta\sup_\theta R(\theta,\delta)$ & Bywa konserwatywny w typowych sytuacjach\\
Ważenie realistycznych scenariuszy & Średnie ryzyko przy jawnej $\pi$ & Wynik zależy od przyjętego rozkładu scenariuszy\\
Wykrywanie groźnego zdarzenia & Strata asymetryczna; koszty FP i FN & Częstość bazowa i błędna kalibracja zmieniają decyzję\\
Prognoza typowej wartości & Strata absolutna; mediana & Inny cel niż prognoza sumy i średniej\\
Prognoza całkowitego kosztu & Strata kwadratowa lub odpowiednio dobrana ekonomiczna & Ekstremalne obserwacje mogą być istotne dla kosztu\\
\bottomrule
\end{tabular}
\end{center}

\paragraph{Schemat praktycznej analizy decyzji.}
Najpierw określamy \emph{działania}, które rzeczywiście możemy podjąć.
Następnie ustalamy koszty poszczególnych pomyłek, model obserwacji,
możliwe parametry i sposób porównywania funkcji ryzyka. Dopiero
na końcu obliczamy regułę. Często dobrą praktyką jest \emph{analiza
wrażliwości}: powtarzamy rachunki przy kilku wiarygodnych funkcjach
straty, kosztach i parametrach modelu. Jeżeli decyzja się zmienia,
trzeba to jawnie zasygnalizować.

\Needspace{16\baselineskip}
\noindent\textbf{Przykład obliczeniowy: analiza kosztów w Pythonie.}\par
\par
\Needspace{13\baselineskip}
\smallskip
\begin{lstlisting}[language=Python,caption={Wpływ kosztów pomyłek na decyzję (Python).}]
prevalence = 0.01
sensitivity, specificity = 0.90, 0.95
p_pos_if_fault = sensitivity
p_pos_if_ok = 1 - specificity
posterior = prevalence * p_pos_if_fault / (
    prevalence * p_pos_if_fault + (1-prevalence) * p_pos_if_ok
)
false_positive_cost, false_negative_cost = 1, 10
cost_intervene = false_positive_cost * (1-posterior)
cost_ignore = false_negative_cost * posterior
print(posterior, cost_intervene, cost_ignore)
print('interweniowac:', cost_intervene < cost_ignore)
\end{lstlisting}
Ten kod pozwala przeprowadzić \emph{analizę wrażliwości}: należy
zmieniać częstość bazową, czułość, swoistość i koszty, a następnie
obserwować, kiedy optymalna decyzja ulega zmianie.

\section{Zadanie łączące wiarygodność i teorię decyzji}
\begin{zadanie}[Ta sama próba, różne kryteria]\label{zad:compare-likelihood-decision}
Niech $X\sim\operatorname{Bernoulli}(p)$, $p\in[0,1]$.
Wyznacz MLE $p$ i porównaj jego ryzyko kwadratowe z regułą minimaksową
z twierdzenia~\ref{thm:minimax-bernoulli-one} dla
$p\in\{0,1/4,1/2,3/4,1\}$.
\par\emph{Rozwiązanie.} MLE wynosi $\widehat p=X$.
Jego ryzyko to $p(1-p)$; wartości przy wymienionych parametrach to
odpowiednio $0$, $3/16$, $1/4$, $3/16$, $0$.
Dla reguły $\delta^*(X)=1/4+X/2$ ryzyko jest stale równe $1/16$.
To, że minimaks wygrywa przy $p=1/2$, nie unieważnia MLE przy
$p=0$ albo $p=1$. MLE odpowiada na pytanie, który parametr
maksymalizuje wiarygodność zaobserwowanego wyniku, a reguła
minimaksowa na pytanie o najniższy możliwy największy koszt
szacowania. Są to inne cele matematyczne.
\end{zadanie}

\paragraph{Zakończenie części wprowadzającej.}
Rozdziały 1--4 rozdzielają cztery poziomy wnioskowania:
\emph{eksperyment} określa rodzinę możliwych praw danych,
\emph{statystyka} redukuje obserwacje,
\emph{wiarygodność} porównuje parametry względem obserwowanych danych,
a \emph{reguła decyzyjna} przypisuje danym działanie oceniane przez
ryzyko. Dopiero po takim rozróżnieniu można sensownie mówić o
optymalnych estymatorach, testach i metodach bayesowskich.
Rozdziały 5--6 rozwiną formalne porównywanie estymatorów,
dostateczność i zupełność; później przejdziemy do granic efektywności.


% ================================================================
% ROZDZIAL 5 -- KONSTRUKCJA ESTYMATOROW
% ================================================================
\chapter{Konstrukcja i porównywanie estymatorów}
\label{ch:estimators}
\paragraph{Powiązania i odsyłacze.} Wcześniejsze definicje funkcji wiarygodności i ryzyka pochodzą z rozdziałów~\ref{ch:likelihood} i~\ref{ch:decision}; o UMVU mówimy w rozdziale~\ref{ch:sufficiency}.


\section{Co dokładnie chcemy oszacować?}
W rozdziale o wiarygodności szukaliśmy parametru, który najlepiej
wyjaśnia zaobserwowane dane \emph{w obrębie założonego modelu}.
Jednak zastosowanie nie zawsze wymaga odtworzenia całego parametru.
W kontroli jakości może nas interesować $p$, prawdopodobieństwo
wystąpienia wady, albo $g(p)=1-(1-p)^m$, prawdopodobieństwo
co najmniej jednej wady w następnych $m$ sztukach. W prognozowaniu
awarii interesuje nas nie tylko intensywność $\lambda$, lecz także
prawdopodobieństwo $e^{-\lambda t}$ bezawaryjnej pracy przez czas $t$.
Odróżnimy zatem \emph{parametr modelu} od \emph{estymowanej wielkości}.

\begin{definicja}[Estymand i estymator]\label{def:estimand}
Niech $\mathcal E=(\mathsf X,\calA,\{P_\theta:\theta\in\Theta\})$
będzie eksperymentem statystycznym, $g:\Theta\to\R^k$ ustaloną
funkcją, a $T:(\mathsf X,\calA)\to(\R^k,\calB(\R^k))$
mierzalnym odwzorowaniem. Wielkość $g(\theta)$ nazywamy
\emph{estymandem}, a $T(X)$ \emph{estymatorem} tej wielkości.
Jeśli $\E_\theta\|T\|<\infty$, obciążeniem jest
$b_T(\theta)=\E_\theta T-g(\theta)$. Estymator jest
\emph{nieobciążony}, jeżeli $b_T(\theta)=0$ dla każdego
$\theta\in\Theta$. Dla $k=1$ i $T\in L^2(P_\theta)$
$\operatorname{MSE}_\theta(T)=\E_\theta(T-g(\theta))^2$.
\end{definicja}

Definicja nie uzależnia estymatora od wybranej metody: średnia,
mediana, maksymalizator wiarygodności i rozwiązanie równania
momentowego są różnymi funkcjami \emph{tych samych danych}.
Praktycznie należy najpierw sprecyzować estymand, dopiero potem
wybrać algorytm. Dla przykładu liczba awarii przypadająca na godzinę
i prawdopodobieństwo przynajmniej jednej awarii w ciągu godziny
nie są tym samym parametrem, nawet w modelu Poissona.

\begin{definicja}[Zgodność i szybkość zbieżności]\label{def:est-consistency}
Dla eksperymentów $\mathcal E_n$ z parametrem $\theta\in\Theta$
i estymatorów $T_n$ wielkości $g(\theta)$ mówimy, że
$T_n$ jest \emph{zgodny}, jeżeli dla każdego $\theta$ i $\varepsilon>0$
zachodzi $P_\theta(\|T_n-g(\theta)\|>\varepsilon)\to0$.
Zgodność \emph{mocna} oznacza zbieżność prawie na pewno, gdy
estymatory są określone na wspólnej przestrzeni probabilistycznej.
Jeżeli $\sqrt n(T_n-g(\theta))$ zbiega według rozkładu do pewnej
zmiennej $Z_\theta$, mówimy, że ciąg $\sqrt n(T_n-g(\theta))$ ma granicę rozkładową; nie wyklucza to szybszej zbieżności w przypadku granicy zdegenerowanej.
\end{definicja}

Nieobciążoność jest własnością wartości oczekiwanej dla danego $n$;
zgodność opisuje zachowanie przy rosnącej liczbie danych; MSE
mierzy błąd w określonej funkcji straty. Żadna z tych własności
nie implikuje automatycznie pozostałych. Na przykład dla
$X_1,\ldots,X_n\iid\Normal(\mu,1)$ estymator $T_n=X_1$
jest nieobciążony dla każdego $n$, ale nie jest zgodny, ponieważ
jego rozkład błędu nie zależy od $n$.

\section{Metoda momentów: dopasowujemy proste funkcjonały rozkładu}
\label{sec:moments-method}
Przypuśćmy, że nie umiemy łatwo zmaksymalizować wiarygodności,
ale potrafimy obliczyć $\E_\theta X$, $\E_\theta X^2$ lub inne
momenty. Intuicja metody momentów jest prosta: zastępujemy
wartości teoretyczne odpowiednimi średnimi z próby i rozwiązujemy
otrzymane równania. Istotą metody są jednak warunki istnienia momentów,
identyfikowalność oraz istnienie rozwiązania w przestrzeni parametrów.

\begin{definicja}[Estymator metodą momentów]\label{def:mom}
Niech $X_1,\ldots,X_n$ będą próbą prostą z rozkładu $P_\theta$,
$\theta\in\Theta\subset\R^d$. Niech
$h:\mathsf S\to\R^d$ będzie mierzalna i całkowalna przy każdym
$P_\theta$, a $m(\theta)=\E_\theta h(X_1)$.
Mierzalne rozwiązanie $\widehat\theta_{\rm MM}$ równania
\[
m(\widehat\theta_{\rm MM})=\frac1n\sum_{i=1}^n h(X_i)
\]
nazywamy \emph{estymatorem metodą momentów}.
Rozwiązanie może nie istnieć, nie być jednoznaczne albo nie należeć
do $\Theta$; wtedy trzeba jawnie określić dalszą procedurę.
Wybór $h(x)=(x,x^2,\ldots,x^d)$ daje klasyczną metodę momentów.
\end{definicja}

\begin{twierdzenie}[Prosta zasada zgodności metody momentów]
\label{thm:mom-consistency}
Niech $X_i$ będą niezależne o jednakowym rozkładzie $P_{\theta_0}$,
$\theta_0\in\Theta\subset\R^d$, i niech $\E_{\theta_0}\|h(X_1)\|<\infty$.
Załóżmy, że istnieje zbiór otwarty $U\subset\R^d$ zawierający
$m(\theta_0)$ oraz funkcja ciągła $q:U\to\Theta$ spełniająca
$q(m(\theta_0))=\theta_0$ i $m(q(y))=y$ dla wszystkich $y\in U$.
Definiujemy $\widehat\theta_n=q(n^{-1}\sum_i h(X_i))$ na zdarzeniu,
na którym średnia należy do $U$, a poza nim przypisujemy ustaloną
wartość z $\Theta$. Wtedy $\widehat\theta_n\to\theta_0$ prawie na pewno.
\end{twierdzenie}
\begin{proof}
Stosujemy mocne prawo wielkich liczb do każdej z $d$ współrzędnych
$h(X_i)$. Na wspólnym zdarzeniu prawdopodobieństwa jeden
$\bar h_n:=n^{-1}\sum_i h(X_i)\to m(\theta_0)$.
Ponieważ $U$ jest otwarte, od pewnego losowego indeksu $N(\omega)$
mamy $\bar h_n\in U$. Dla tych $n$ estymator jest równy
$q(\bar h_n)$. Z ciągłości $q$ w $m(\theta_0)$ dostajemy
$q(\bar h_n)\to q(m(\theta_0))=\theta_0$. Dowód uwzględnia
także możliwość, że dla małych prób układ momentowy nie ma rozwiązania.
\end{proof}


\begin{twierdzenie}[Asymptotyczny rozkład estymatora momentowego]
\label{thm:mom-clt}
Niech $X_1,X_2,\ldots$ będą niezależne o jednakowym rozkładzie
$P_{\theta_0}$, niech $h(X_1)$ będzie rzeczywista,
$\E_{\theta_0}h(X_1)=m_0$, $0<\Var_{\theta_0}(h(X_1))=v<\infty$.
Załóżmy, że funkcja $q$ jest określona i różniczkowalna w otoczeniu
$m_0$, przy czym $q(m_0)=\theta_0$. Niech
$\widehat\theta_n=q(n^{-1}\sum_i h(X_i))$ wszędzie, gdzie argument
należy do tego otoczenia, a poza nim przypiszmy stałą wartość.
Wtedy
\[
\sqrt n(\widehat\theta_n-\theta_0)
\xrightarrow{d}\Normal(0,v[q'(m_0)]^2).
\]
Dopuszczamy rozkład graniczny skupiony w zerze, gdy $q'(m_0)=0$.
\end{twierdzenie}
\begin{proof}
Z centralnego twierdzenia granicznego zastosowanego do zmiennej
$h(X_i)$ wynika
$\sqrt n(\bar h_n-m_0)\xrightarrow{d}\Normal(0,v)$,
a więc w szczególności $\bar h_n\to m_0$ według
prawdopodobieństwa. Z różniczkowalności $q$ w $m_0$ mamy
$q(y)=q(m_0)+q'(m_0)(y-m_0)+r(y)(y-m_0)$,
gdzie $r(y)\to0$ przy $y\to m_0$.
Stąd
\[
\sqrt n(\widehat\theta_n-\theta_0)
=q'(m_0)\sqrt n(\bar h_n-m_0)
+r(\bar h_n)\sqrt n(\bar h_n-m_0)
\]
na zdarzeniu, którego prawdopodobieństwo dąży do jedności.
Czynnik $r(\bar h_n)$ dąży do zera według prawdopodobieństwa,
a ciąg $\sqrt n(\bar h_n-m_0)$ jest ciasny, bo ma granicę
według rozkładu. Ich iloczyn dąży więc do zera według
prawdopodobieństwa. Z twierdzenia Slutsky'ego otrzymujemy tezę.
\end{proof}
\begin{przyklad}[Błąd standardowy estymatora intensywności]
Dla czasów wykładniczych $X_i\sim\operatorname{Exp}(\lambda)$
wybieramy $h(x)=x$ i $q(m)=1/m$. Wówczas
$m_0=1/\lambda$, $v=1/\lambda^2$ oraz
$q'(m_0)=-\lambda^2$. Poprzednie twierdzenie daje
$\sqrt n(\widehat\lambda-\lambda)\xrightarrow{d}
\Normal(0,\lambda^2)$. Zatem orientacyjny błąd standardowy MLE
wynosi $\widehat\lambda/\sqrt n$ przy dużym $n$;
nie jest to wzór na jego dokładny rozkład dla małej próby.
\end{przyklad}

\begin{przyklad}[Bernoulli i Poisson]
Dla $X_i\iid\operatorname{Bernoulli}(p)$ mamy $\E X_i=p$, stąd
$\widehat p_{\rm MM}=\bar X$; jest to także MLE.
Dla $Y_i\iid\Pois(\lambda)$ średnia wynosi $\lambda$, więc
$\widehat\lambda_{\rm MM}=\bar Y$, ponownie tak samo jak MLE.
Zbieżność obu wynika z MPWL. Równość tych dwóch procedur jest
własnością modeli, a nie ogólnym twierdzeniem o metodach estymacji.
\end{przyklad}

\begin{zadanie}[Beta z dwóch momentów]
\label{zad:beta-mom}
Niech $X_1,\ldots,X_n\iid\operatorname{Beta}(\alpha,\beta)$,
$\alpha,\beta>0$, o średniej $\mu$ i wariancji $v$.
Wyznacz estymatory momentowe $\alpha,\beta$ i sprawdź,
kiedy istnieją w dodatniej przestrzeni parametrów.
\par\emph{Rozwiązanie.} Dla $k=\alpha+\beta$ zachodzi
$\mu=\alpha/k$ oraz $v=\mu(1-\mu)/(k+1)$. Zatem
$k=\mu(1-\mu)/v-1$ i $\alpha=\mu k$, $\beta=(1-\mu)k$.
Wstawiamy $\bar X$ i $v_n=n^{-1}\sum_i(X_i-\bar X)^2$:
\[
\widehat k=\frac{\bar X(1-\bar X)}{v_n}-1,\qquad
\widehat\alpha=\bar X\widehat k,\quad
\widehat\beta=(1-\bar X)\widehat k.
\]
Wymagamy $0<\bar X<1$ i $0<v_n<\bar X(1-\bar X)$.
Dla rzeczywistych danych z wnętrza $(0,1)$ i niezerowej wariancji
nierówność $v_n<\bar X(1-\bar X)$ wynika z $x_i^2<x_i$:
$v_n=n^{-1}\sum_i x_i^2-\bar X^2<\bar X-\bar X^2$.
Jeżeli wszystkie obserwacje są jednakowe, $v_n=0$ i metoda nie
daje skończonego rozwiązania. W praktyce dla danych przy brzegu
albo nadmiernie skoncentrowanych warto porównać wynik z numerycznym MLE.
\end{zadanie}

\section{Różne metody, ten sam model: rozkład jednostajny}
\begin{twierdzenie}[Trzy estymatory końca przedziału i ich ryzyka]
\label{thm:uniform-comparison}
Niech $n\ge2$ i $X_1,\ldots,X_n\iid U(0,\theta)$, $\theta>0$.
Oznaczmy $M=X_{(n)}$. Dla estymacji $\theta$ przy stracie kwadratowej:
\begin{enumerate}
\item metoda momentów daje $T_{\rm MM}=2\bar X$ oraz
$\E_\theta(T_{\rm MM}-\theta)^2=\theta^2/(3n)$;
\item estymator MLE (dla wersji gęstości z $x\le\theta$) to
$T_{\rm ML}=M$, o obciążeniu $-\theta/(n+1)$ i ryzyku
$2\theta^2/((n+1)(n+2))$;
\item estymator $T_U=(n+1)M/n$ jest nieobciążony i ma ryzyko
$\theta^2/(n(n+2))$.
\end{enumerate}
Dla $n>1$ estymator $T_U$ ma mniejsze MSE od obu pozostałych.
\end{twierdzenie}
\begin{proof}
Z $\E X_i=\theta/2$ i $\Var(X_i)=\theta^2/12$
otrzymujemy $\E2\bar X=\theta$ i
$\Var(2\bar X)=4(\theta^2/12)/n=\theta^2/(3n)$.
Maksimum ma dystrybuantę $P_\theta(M\le t)=(t/\theta)^n$
dla $0\le t\le\theta$, a więc gęstość
$n t^{n-1}/\theta^n$. Stąd przez całkowanie
$\E M=n\theta/(n+1)$ oraz
$\E M^2=n\theta^2/(n+2)$. Zatem
\[
\E(M-\theta)^2=\frac{n\theta^2}{n+2}
-\frac{2n\theta^2}{n+1}+\theta^2
=\frac{2\theta^2}{(n+1)(n+2)}.
\]
Po przeskalowaniu $M$ czynnik $(n+1)/n$ usuwa obciążenie;
wykorzystując powyższe momenty, otrzymujemy
$\Var(M)=n\theta^2/((n+1)^2(n+2))$ i
$\Var(T_U)=\theta^2/(n(n+2))$.
Dla $n>1$ zachodzą nierówności
$1/(n(n+2))<2/((n+1)(n+2))$ i
$1/(n(n+2))<1/(3n)$, co kończy porównanie.
\end{proof}


\begin{center}
\begin{tikzpicture}[x=.47cm,y=15cm,>=Latex]
\draw[->] (0,0)--(21,0) node[right] {$n$};
\draw[->] (0,0)--(0,.188) node[above] {$\mathrm{MSE}/\theta^2$};
\foreach \x in {2,5,10,15,20} {\draw (\x,-.003)--(\x,.003) node[below=2pt] {\scriptsize $\x$};}
\foreach \y/\lab in {.05/{0{,}05},.10/{0{,}10},.15/{0{,}15}} {
  \draw (-.25,\y)--(.25,\y) node[left=2pt] {\scriptsize $\lab$};
}
\draw[very thick,blue!70!black,domain=2:20,samples=120]
  plot (\x,{1/(3*\x)});
\draw[very thick,red!70!black,domain=2:20,samples=120]
  plot (\x,{2/((\x+1)*(\x+2))});
\draw[very thick,green!45!black,domain=2:20,samples=120]
  plot (\x,{1/(\x*(\x+2))});
\node[align=left,anchor=west,font=\scriptsize] at (8,.148) {
  \textcolor{blue!70!black}{--- metoda momentów}\\
  \textcolor{red!70!black}{--- MLE}\\
  \textcolor{green!45!black}{--- skorygowane maksimum}
};
\end{tikzpicture}
\end{center}
\noindent\emph{Rysunek 5.1.} Znormalizowany błąd średniokwadratowy
trzech estymatorów końca przedziału jednostajnego. Wykresy
przedstawiają wzory ścisłe z poprzedniego twierdzenia,
nie wyniki symulacji.

\begin{uwaga}[Dlaczego maksimum jest lepsze od średniej?]
Model zakłada twardą granicę nośnika. Największa obserwacja
informuje o tej granicy \emph{bezpośrednio}, natomiast średnia
wykorzystuje typowy poziom próby i rozprasza informację o końcu.
Dla rozkładu jednostajnego obcięcie jednego skrajnego pomiaru
może silnie zmienić estymator; nie jest to wada rachunku, tylko
konsekwencja bardzo mocnego założenia o ostrym końcu przedziału.
Jeżeli pomiary mogą zawierać błąd przekraczający $\theta$,
model jednostajny przestaje być odpowiedni.
\end{uwaga}

\section{Estymacja parametru z rozkładu wykładniczego}
\begin{twierdzenie}[Estymator intensywności i jego obciążenie]
\label{thm:exp-rate}
Niech $X_1,\ldots,X_n\iid\operatorname{Exp}(\lambda)$,
$\lambda>0$, w parametryzacji z gęstością
$f_\lambda(x)=\lambda e^{-\lambda x}\ind{\{x>0\}}$.
Dla $S=\sum_iX_i$ estymator metodą momentów i MLE wynosi
$\widehat\lambda_{\rm ML}=n/S$ prawie na pewno.
Jeśli $n>1$, to $\E_\lambda\widehat\lambda_{\rm ML}
=n\lambda/(n-1)$, a $\widetilde\lambda=(n-1)/S$ jest
nieobciążony. Dla $n>2$ zachodzi
$\Var_\lambda(\widetilde\lambda)=\lambda^2/(n-2)$.
\end{twierdzenie}
\begin{proof}
Mamy $\E X_i=1/\lambda$, a więc równanie momentów
$\bar X=1/\lambda$ prowadzi do $\lambda=n/S$.
Log-wiarygodność $\ell(\lambda)=n\log\lambda-\lambda S$
ma pochodną $n/\lambda-S$, stąd to samo maksimum; wartości
na brzegach $0$ i $+\infty$ są mniejsze dla $S>0$.
Suma niezależnych czasów wykładniczych ma rozkład gamma o
parametrach kształtu $n$ i intensywności $\lambda$:
$f_S(s)=\lambda^n s^{n-1}e^{-\lambda s}/\Gamma(n)$ dla $s>0$.
Dla $n>k$ obliczenie całki gamma daje
\[
\E_\lambda S^{-k}
=\frac{\lambda^n}{\Gamma(n)}
 \int_0^\infty s^{n-k-1}e^{-\lambda s}\,\dd s
=\lambda^k\frac{\Gamma(n-k)}{\Gamma(n)}.
\]
Podstawiając $k=1,2$, dostajemy
$\E S^{-1}=\lambda/(n-1)$ oraz
$\E S^{-2}=\lambda^2/((n-1)(n-2))$.
Pomnożenie przez odpowiednie czynniki kończy dowód.
\end{proof}

\begin{zadanie}[Pięć czasów do awarii: czy poprawiać obciążenie?]
Zaobserwowano czasy $0{,}4$, $0{,}6$, $1{,}2$, $1{,}8$, $2{,}0$
godziny. W modelu wykładniczym oszacuj intensywność awarii
oraz średni czas do awarii.
\par\emph{Rozwiązanie.} $S=6$ godzin, więc
$\widehat\lambda_{\rm ML}=5/6\approx0{,}833$ na godzinę,
a estymator nieobciążony intensywności to
$\widetilde\lambda=4/6\approx0{,}667$ na godzinę.
Średni czas do awarii $1/\lambda$ estymujemy naturalnie przez
$\bar X=1{,}2$ godziny, który w tym modelu jest nieobciążony.
\emph{Uwaga:} $1/\widetilde\lambda=1{,}5$ godziny nie jest tym
samym estymatorem; nieliniowe przekształcenie estymatora
nie zachowuje na ogół nieobciążoności. Przy tylko pięciu
obserwacjach ocena sensowności modelu wykładniczego jest
ważniejsza od trzeciej cyfry oszacowania.
\end{zadanie}

\section{Estymacja odporna: kiedy momenty i MLE zawodzą?}
Załóżmy, że obserwujemy temperatury w stabilnym procesie, ale
jeden czujnik okresowo zwraca ogromne wartości. Średnia jest
niezwykle czuła na ten problem: wystarczy jeden pomiar $M\to\infty$,
aby średnia rosła bez ograniczeń. Mediana pozostaje ograniczona,
jeśli mniej niż połowa wyników została zastąpiona skrajnymi wartościami.
Nie oznacza to jednak, że mediana jest zawsze lepsza od średniej.
W idealnym modelu normalnym średnia wykorzystuje informację
efektywniej w sensie wariancji asymptotycznej.

\begin{definicja}[Mediana, średnia obcięta, estymator Hubera]
Niech $x_{(1)}\le\cdots\le x_{(n)}$ będą uporządkowanymi obserwacjami.
Mediana empiryczna jest dowolnym mediano-minimalizatorem
$\sum_i|x_i-a|$ (w przykładach przy $n$ parzystym wybieramy
środek dwóch środkowych obserwacji). Przy całkowitym $k<n/2$
\emph{dwustronną średnią obciętą} nazywamy
$\overline x_{\rm trim}=(n-2k)^{-1}\sum_{i=k+1}^{n-k}x_{(i)}$.
Dla $c>0$ wprowadzamy funkcję
$\psi_c(u)=\max(-c,\min(u,c))$. \emph{Estymator M Hubera położenia}
to mierzalne rozwiązanie równania
$\sum_{i=1}^n\psi_c((x_i-a)/s)=0$ przy ustalonej skali $s>0$.
W przypadku wielu rozwiązań potrzebna jest reguła wyboru.
\end{definicja}

\begin{przyklad}[Cauchy: dlaczego metoda momentów nie ma sensu?]
Dla $X\sim\operatorname{Cauchy}(x_0,\gamma)$, $\gamma>0$,
nie istnieje wartość oczekiwana ani wariancja. Równanie
$\bar X=\E X$ nie jest więc zdefiniowane. Zamiast tego można
użyć kwantyli: $q_{1/2}=x_0$,
$q_{1/4}=x_0-\gamma$ i $q_{3/4}=x_0+\gamma$. Stąd
$\widehat x_0=\widehat q_{1/2}$ oraz
$\widehat\gamma=(\widehat q_{3/4}-\widehat q_{1/4})/2$.
Wybór empirycznej konwencji kwantyla należy ustalić przed
obliczeniem. Istnieje również MLE oparty na pełnej gęstości,
ale zwykle wymaga optymalizacji numerycznej. Porównanie metod
w praktyce powinno uwzględniać ciężkość ogonów i odporność
na odległe obserwacje.
\end{przyklad}

\begin{zadanie}[Dane z wartością odstającą: trzy odpowiedzi]
Dla próby $4,5,5,6,6,7,7,8,100$ wyznacz średnią, medianę
i średnią obciętą po jednej obserwacji z każdej strony.
\par\emph{Rozwiązanie.} Suma wynosi $148$, więc
$\bar x=148/9\approx16{,}44$. Mediana to piąty wynik, równy $6$.
Po usunięciu minimum $4$ i maksimum $100$ zostaje suma $44$,
a zatem średnia obcięta to $44/7\approx6{,}29$.
Jeżeli celem jest oszacowanie typowego poziomu przy możliwej
awarii czujnika, mediana lub estymator odporny są rozsądne.
Jeżeli $100$ oznacza rzeczywisty koszt, którego nie wolno ignorować,
średnia może być odpowiednia. Dobór statystyki wynika z definicji
problemu, a nie wyłącznie z testu normalności.
\end{zadanie}

\section{Jak porównywać estymatory? Teoria i praktyka}
\begin{stwierdzenie}[Przekształcenie zgodnych estymatorów]
\label{prop:continuous-estimator}
Niech $T_n\to t_0$ według prawdopodobieństwa przy ustalonym
$P_{\theta_0}$ i niech $g$ będzie ciągła w $t_0$. Wtedy
$g(T_n)\to g(t_0)$ według prawdopodobieństwa.
\end{stwierdzenie}
\begin{proof}
Dla dowolnego $\varepsilon>0$ z ciągłości $g$ w $t_0$
wybieramy $\delta>0$ takie, że $\|x-t_0\|<\delta$
implikuje $\|g(x)-g(t_0)\|<\varepsilon$.
Wówczas
$P(\|g(T_n)-g(t_0)\|\ge\varepsilon)
\le P(\|T_n-t_0\|\ge\delta)\to0$.
\end{proof}

Zgodność przekształconego MLE nie oznacza zachowania
nieobciążoności. Podobnie \emph{mniejsza wariancja} nie wystarcza
do uznania estymatora obciążonego za lepszy: trzeba porównywać MSE
w odniesieniu do \emph{tego samego estymandu}.

\begin{center}\small
\begin{tabular}{>{\raggedright\arraybackslash}p{0.24\textwidth}>{\raggedright\arraybackslash}p{0.32\textwidth}>{\raggedright\arraybackslash}p{0.31\textwidth}}
\toprule
\textbf{Podejście} & \textbf{Kiedy warto rozważyć?} & \textbf{Co sprawdzić przed użyciem?}\\
\midrule
Metoda momentów & Momenty są znane i łatwe do odwrócenia; potrzebujemy punktu startowego optymalizacji & Istnienie momentów, identyfikowalność, rozwiązanie w $\Theta$\\
MLE & Wiarygodny pełny model parametryczny, dobrze określone maksimum & Brzegi i zależność nośnika, multimodalność, dopasowanie modelu\\
Estymator nieobciążony & Wymóg braku systematycznego błędu przy wszystkich $\theta$ & Wariancja i MSE; nieobciążoność nie gwarantuje stabilności\\
Mediana i metody odporne & Obserwacje odstające, ciężkie ogony, cel „typowy poziom” & Koszt utraty efektywności, model skażenia, dobór skali\\
Reguła minimaksowa & Liczy się zabezpieczenie najgorszego ryzyka & Czy najgorszy przypadek ma znaczenie praktyczne?\\
\bottomrule
\end{tabular}
\end{center}

\paragraph{Praktyczny schemat pracy.}
Najpierw wskazujemy estymand i jednostkę obserwacji. Następnie
sprawdzamy, czy wszystkie wykorzystywane momenty istnieją;
wyprowadzamy kandydatów oraz ich ograniczenia. Porównujemy obciążenie,
MSE i zgodność \emph{pod zadanym modelem}; badamy również wrażliwość
na odchylenia od modelu. Na koniec raportujemy jednostkę parametru
i niepewność oszacowania. Same estymaty punktowe nie zastępują
przedziałów ufności, które wprowadzimy w kolejnych rozdziałach.

\Needspace{16\baselineskip}
\noindent\textbf{Minimalny kod w Pythonie (bez zewnętrznych pakietów).}\par
\par
\Needspace{13\baselineskip}
\smallskip
\begin{lstlisting}[language=Python,caption={Porównanie statystyk odpornych i estymatorów intensywności (Python).}]
from statistics import mean, median
x = [4, 5, 5, 6, 6, 7, 7, 8, 100]
print('srednia:', mean(x), 'mediana:', median(x))
print('srednia obcieta:', mean(sorted(x)[1:-1]))
waiting_times = [0.4, 0.6, 1.2, 1.8, 2.0]
S, n = sum(waiting_times), len(waiting_times)
print('MLE lambda:', n/S, 'nieobciazony:', (n-1)/S)
\end{lstlisting}
\Needspace{16\baselineskip}
\noindent\textbf{To samo w Julii.}\par
\par
\Needspace{13\baselineskip}
\smallskip
\begin{lstlisting}[language=Julia,caption={Średnia obcięta i estymatory intensywności (Julia).}]
using Statistics
x = [4, 5, 5, 6, 6, 7, 7, 8, 100]
println((mean(x), median(x), mean(sort(x)[2:end-1])))
w = [0.4, 0.6, 1.2, 1.8, 2.0]
println((length(w)/sum(w), (length(w)-1)/sum(w)))
\end{lstlisting}


% ================================================================
% ROZDZIAL 6 -- DOSTATECZNOSC, ZUPELNOSC I UMVU
% ================================================================
\chapter{Dostateczność, zupełność i optymalna estymacja nieobciążona}
\label{ch:sufficiency}
\paragraph{Powiązania i odsyłacze.} Punktem wyjścia są statystyki i estymatory z rozdziałów~\ref{ch:statystyki} i~\ref{ch:estimators}; ograniczenia efektywności opisuje rozdział~\ref{ch:fisher-information}.


\section{Czy wolno zastąpić całą próbę jedną liczbą?}
Dla $n$ niezależnych prób Bernoulliego zapisujemy wektor
$X=(X_1,\ldots,X_n)$, ale estymujemy parametr jedynie przez
liczbę sukcesów $S=\sum_i X_i$. Dlaczego wolno wyrzucić
informację o kolejności wyników? W tym konkretnym modelu,
\emph{warunkowo na $S=k$}, każda sekwencja zawierająca $k$
sukcesów ma to samo prawdopodobieństwo $\binom nk^{-1}$,
które \emph{nie zależy od $p$}. Znajomość kolejności nie wnosi
więc nowej informacji o parametrze. Taką redukcję formalizuje
dostateczność. Nie każda statystyka ma tę własność: w rozkładzie
jednostajnym ze zmiennym kresem informację skupia maksimum,
a w ogólnym modelu zależnym od czasu kolejność może być istotna.

\begin{definicja}[Statystyka dostateczna]\label{def:sufficiency}
Niech $X$ ma prawo $P_\theta$ na przestrzeni standardowej borelowskiej
$(\mathsf X,\calA)$ dla każdego $\theta\in\Theta$,
a $T:\mathsf X\to\mathsf T$ będzie mierzalna, przy czym
$\mathsf T$ także jest przestrzenią standardową borelowską.
Mówimy, że $T$ jest \emph{dostateczna} dla rodziny
$\{P_\theta\}$, jeśli istnieje jądro probabilistyczne
$K(t,A)$, mierzalne w $t$ dla każdego $A\in\calA$,
takie, że dla wszystkich $\theta$ i $A\in\calA$
\[
P_\theta(X\in A\mid\sigma(T))=K(T,A)
\quad P_\theta\text{-prawie na pewno}.
\]
\end{definicja}

Warunkowy rozkład pełnych danych przy ustalonej statystyce musi mieć
\emph{jedną wspólną wersję} niezależną od parametru. Nie wystarcza,
że dla każdego parametru oddzielnie istnieje pewien rozkład
warunkowy: to zawsze jest prawdą w przestrzeniach standardowych
borelowskich. Dostateczność jest własnością \emph{rodziny rozkładów},
a nie jednej obserwacji.

\begin{zadanie}[Dlaczego liczba sukcesów wystarcza?]\label{zad:suff-bernoulli}
Niech $X_1,\ldots,X_n\iid\operatorname{Bernoulli}(p)$,
$p\in[0,1]$, $S=\sum_iX_i$. Oblicz prawo pełnej sekwencji
warunkowo na $S=k$.
\par\emph{Rozwiązanie.} Dla $0<p<1$ i sekwencji
$x\in\{0,1\}^n$ o sumie $k$ wzór Bayesa daje
\[
P_p(X=x\mid S=k)=\frac{p^k(1-p)^{n-k}}
 {\binom nk p^k(1-p)^{n-k}}=\binom nk^{-1}.
\]
Dla sekwencji o innej sumie prawdopodobieństwo jest zerowe.
Przy $p=0$ obserwujemy $S=0$ prawie na pewno, a przy $p=1$
$S=n$ prawie na pewno; ta sama formuła zdefiniowana dla wszystkich
$k$ nadal stanowi wspólną wersję jądra. Jest więc ono niezależne
od $p$, zatem $S$ jest dostateczna. Wniosek praktyczny:
w modelu \emph{iid Bernoulliego} do wnioskowania o $p$ można
zachować tylko $S$ i $n$; jeśli prawdopodobieństwo sukcesu zmienia
się z czasem, redukcja może zgubić informację o trendzie.
\end{zadanie}


\begin{center}
\begin{tikzpicture}[>=Latex,node distance=10mm,
 box/.style={draw,rounded corners,align=center,inner sep=5pt,
 font=\small,minimum height=13mm}]
\node[box,text width=30mm] (full) {Pełna próba $X$\\kolejność i wyniki};
\node[box,text width=38mm,right=of full] (stat) {Statystyka $T(X)$\\informacja o parametrze};
\node[box,text width=41mm,right=of stat] (cond) {Warunkowe prawo $X\mid T$\\nie zależy od $\theta$};
\draw[->,thick] (full)--(stat);
\draw[->,thick] (stat)--(cond);
\end{tikzpicture}
\end{center}
\noindent\emph{Rysunek 6.1.} Dostateczność nie oznacza możliwości
odtworzenia pełnej próby z $T$. Oznacza, że pozostała losowość
próby po ustaleniu $T$ nie zależy od nieznanego parametru.

\section{Dlaczego istnieje wspólna miara odniesienia?}
Aby uzasadnić kryterium faktoryzacyjne bez pomijania warunków
mierzalności, przypomnimy jeden lemat. W całym podrozdziale
zakładamy standardowe przestrzenie borelowskie i zdominowaną
rodzinę $P_\theta\ll\mu$, gdzie $\mu$ jest $\sigma$-skończona.
Jądra warunkowe względem $T$ istnieją w tych przestrzeniach,
co przyjmujemy jako standardowy wynik o regularnych rozkładach
warunkowych.

\begin{lemat}[Przeliczalna mieszanina dominująca]\label{lem:halmos-savage}
Dla dowolnej rodziny miar probabilistycznych $\{P_\theta\}_{\theta\in\Theta}$
zdominowanej przez jedną miarę $\sigma$-skończoną $\mu$
istnieją parametry $\theta_1,\theta_2,\ldots$ oraz dodatnie wagi
$a_j$ o sumie $1$ takie, że prawdopodobieństwo
$P_*:=\sum_{j\ge1}a_jP_{\theta_j}$ dominuje każde $P_\theta$.
\end{lemat}
\begin{proof}
Wybierzmy miarę skończoną $\nu$ równoważną $\mu$: jeśli
$\mathsf X=\bigcup_{j\ge1}A_j$ i $0<\mu(A_j)<\infty$
(po usunięciu zbiorów zerowych i rozłączeniu), możemy położyć
$\nu(B)=\sum_j2^{-j}\mu(B\cap A_j)/(1+\mu(A_j))$,
a następnie znormalizować. Dla każdego $\theta$ niech
$B_\theta=\{\dd P_\theta/\dd\mu>0\}$.
Niech $c$ będzie supremum $\nu$-miar skończonych sum
$\bigcup_{i=1}^m B_{\theta_i}$. Dla każdego $k$ wybieramy
skończoną rodzinę, której suma $D_k$ spełnia
$\nu(D_k)>c-1/k$. Po zebraniu wszystkich wybranych parametrów
w jeden ciąg otrzymujemy $D=\bigcup_kD_k$ z $\nu(D)=c$
(z jednej strony $D$ jest rosnącą granicą skończonych sum,
każda o mierze nie większej niż $c$). Dla każdego $\theta$,
gdyby $\nu(B_\theta\setminus D)>0$, odpowiednio duża skończona
część $D$ powiększona o $B_\theta$ miałaby miarę większą niż $c$,
sprzeczność. Stąd $\nu(B_\theta\setminus D)=0$.
Niech $P_*=\sum_j a_jP_{\theta_j}$, np. $a_j=2^{-j}$.
Jego gęstość jest dodatnia na $D$ z dokładnością do zbioru
$\mu$-zerowego. Jeżeli $P_*(A)=0$, to
$\mu(A\cap D)=0$, a zatem dla każdego $\theta$
$P_\theta(A)=P_\theta(A\cap D)=0$. To oznacza $P_\theta\ll P_*$.
\end{proof}

\begin{twierdzenie}[Kryterium faktoryzacyjne Fishera--Neymana]
\label{thm:factorization}
Niech $\{P_\theta:\theta\in\Theta\}$ będzie rodziną miar na
standardowej przestrzeni borelowskiej $(\mathsf X,\calA)$,
zdominowaną przez miarę $\sigma$-skończoną $\mu$.
Niech $p_\theta=\dd P_\theta/\dd\mu$ będą wybranymi wersjami
gęstości. Statystyka $T:\mathsf X\to\mathsf T$ o wartościach
w przestrzeni standardowej borelowskiej jest dostateczna wtedy
i tylko wtedy, gdy istnieją mierzalne funkcje $h\ge0$ na
$\mathsf X$ i $g_\theta\ge0$ na $\mathsf T$ takie, że
\[
p_\theta(x)=g_\theta(T(x))h(x)
\quad\mu\text{-prawie wszędzie dla każdego }\theta.
\]
Zbiór, na którym równość zawodzi, może zależeć od $\theta$.
\end{twierdzenie}
\begin{proof}
Niech $P_*$ będzie mieszaniną z lematu~\ref{lem:halmos-savage},
a $p_*$ jej gęstością względem $\mu$. Dla każdego $\theta$
mamy $P_\theta\ll P_*$ i pochodną
$r_\theta=\dd P_\theta/\dd P_*$. Załóżmy najpierw,
że $r_\theta(x)=g_\theta(T(x))$ $P_*$-prawie wszędzie.
Niech $K(t,A)$ będzie wspólną wersją regularnego prawa warunkowego
$X$ względem $T$ pod $P_*$. Dla każdego $A\in\calA$
i mierzalnego $B\subset\mathsf T$, korzystając z zależności
$r_\theta=g_\theta(T)$, otrzymujemy
\[
P_\theta(A\cap\{T\in B\})
=\E_*[\ind{\{X\in A\}}g_\theta(T)\ind{\{T\in B\}}]
=\E_*[K(T,A)g_\theta(T)\ind{\{T\in B\}}]
=\E_\theta[K(T,A)\ind{\{T\in B\}}].
\]
Wobec definicji warunkowania $K$ jest wspólnym jądrem warunkowym,
więc $T$ jest dostateczna.

Odwrotnie, jeśli $T$ jest dostateczna z jądrem $K$, dla każdej
ograniczonej funkcji mierzalnej $f$ mamy
$\E_\theta f(X)=\E_\theta\int_{\mathsf X}f(x)K(T,\dd x)$.
Ta sama równość zachodzi pod mieszaniną $P_*$ (sumujemy ją
z wagami $a_j$); oznacza to, że
$\int_{\mathsf X}f(x)K(T,\dd x)=\E_*[f(X)\mid T]$ $P_*$-p.n.
Ponieważ $P_\theta\ll P_*$, dla każdego ograniczonego $f$
\[
\E_*[f(X)r_\theta(X)]
=\E_\theta f(X)
=\E_\theta\E_*[f(X)\mid T]
=\E_*[r_\theta(X)\E_*[f(X)\mid T]].
\]
Z własności warunkowej wartości oczekiwanej ostatnie wyrażenie
równa się $\E_*[f(X)\E_*(r_\theta(X)\mid T)]$.
Dowolność ograniczonego $f$ implikuje
$r_\theta=\E_*(r_\theta\mid T)$ $P_*$-p.n., więc z lematu
Dooba--Dynkina $r_\theta=g_\theta(T)$ $P_*$-p.n.
Równoważność z podanym rozkładem gęstości wynika następująco.
Z zależności $p_\theta=r_\theta p_*$ dostajemy faktoryzację
z $h=p_*$. Jeżeli zaś
$p_\theta=g_\theta(T)h$, to mieszanka ma gęstość
$p_*=(\sum_j a_jg_{\theta_j}(T))h$.
Na zbiorze $p_*>0$ iloraz $p_\theta/p_*$ zależy tylko od $T$,
więc jest $\sigma(T)$-mierzalny. Na zbiorze zerowym $P_*$
wersję ilorazu definiujemy dowolnie. Oba kierunki są dowiedzione.
\end{proof}

\begin{przyklad}[Trzy modele i trzy różne statystyki dostateczne]
\begin{enumerate}
\item W próbie Bernoulliego statystyką dostateczną jest
$S=\sum_iX_i$, ponieważ dla obserwacji $x$ mamy
\[
p_p(x)=p^{S(x)}(1-p)^{n-S(x)}.
\]
\item Poisson: dla $X_i\iid\Pois(\lambda)$,
$p_\lambda(x)=\exp(-n\lambda)\lambda^{\sum_i x_i}
\prod_i(x_i!)^{-1}$; wystarczy $S=\sum_iX_i$.
\item Normalny: dla próby $\Normal(\mu,\sigma^2)$ z \emph{dwoma}
nieznanymi parametrami gęstość zależy od danych przez
$S_1=\sum_iX_i$ i $S_2=\sum_iX_i^2$.
Kiedy $\sigma^2$ jest znane, wystarcza samo $\bar X$.
\end{enumerate}
\end{przyklad}

\section{Minimalna dostateczność: kiedy nie zachowaliśmy za dużo?}
Dostateczną statystykę można zawsze sztucznie rozszerzyć:
wektor $(S,X_1)$ jest dostateczny, jeżeli $S$ jest dostateczny.
Chcielibyśmy zachować \emph{jak najmniej} informacji niezbędnej
do opisu parametru. W modelach z dodatnimi gęstościami na wspólnym
nośniku prowadzi do tego proste kryterium ilorazów wiarygodności.

\begin{definicja}[Minimalna dostateczność]
Dostateczna statystyka $T$ jest \emph{minimalnie dostateczna},
jeśli dla każdej innej dostatecznej statystyki $S$ istnieje mierzalna
funkcja $\varphi$ taka, że $T=\varphi(S)$ $P_\theta$-prawie na pewno
dla każdego $\theta$.
\end{definicja}

\begin{stwierdzenie}[Kryterium ilorazowe w modelu dyskretnym]
\label{prop:minimal-discrete}
Niech $\mathsf X$ będzie przeliczalna i niech wszystkie funkcje
masy $p_\theta(x)$ będą \emph{ściśle dodatnie} dla $x\in\mathsf X$.
Załóżmy, że istnieje statystyka $T$ spełniająca
\[
T(x)=T(y)\quad\Longleftrightarrow\quad
\frac{p_\theta(x)}{p_\theta(y)}\text{ nie zależy od }\theta.
\]
Wówczas $T$ jest minimalnie dostateczna.
\end{stwierdzenie}
\begin{proof}
Wybierzmy ustalone $\theta_0$. Jeśli $T(x)=T(y)$, to
$p_\theta(x)/p_\theta(y)=p_{\theta_0}(x)/p_{\theta_0}(y)$.
W każdej klasie wartości $t$ wybierzmy reprezentanta $x_t$
(co jest możliwe dla przeliczalnej przestrzeni). Definiując
$g_\theta(t)=p_\theta(x_t)/p_{\theta_0}(x_t)$ oraz
$h(x)=p_{\theta_0}(x)$ dostajemy
$p_\theta(x)=g_\theta(T(x))h(x)$, zatem $T$ jest dostateczna
z twierdzenia~\ref{thm:factorization}.
Niech teraz $S$ będzie dowolną inną dostateczną statystyką.
Faktoryzacja daje
$p_\theta(x)=u_\theta(S(x))v(x)$, gdzie $v(x)>0$
na całej przeliczalnej przestrzeni (po ewentualnej zmianie
wersji na zbiorach zerowych: takich punktów nie ma, bo wszystkie
$p_\theta$ są dodatnie). Zatem $S(x)=S(y)$ implikuje
$p_\theta(x)/p_\theta(y)=v(x)/v(y)$, niezależnie od $\theta$,
a więc $T(x)=T(y)$. To znaczy, że $T$ jest stała na włóknach
$S$; dla przeliczalnej dziedziny definiujemy mierzalną funkcję
$\varphi$ na zbiorze wartości $S$ przez $\varphi(S(x))=T(x)$,
a poza nim dowolnie. Otrzymujemy $T=\varphi(S)$.
\end{proof}

\begin{przyklad}[W modelu Bernoulliego liczba sukcesów jest minimalnie dostateczna]
Niech $p\in(0,1)$, $x,y\in\{0,1\}^n$. Wówczas
\[
\frac{p_p(x)}{p_p(y)}
=\left(\frac p{1-p}\right)^{\sum_i x_i-\sum_i y_i}.
\]
Iloraz jest niezależny od $p$ dokładnie wtedy, gdy
$\sum_i x_i=\sum_i y_i$. Zatem liczba sukcesów $S$ jest
minimalnie dostateczna. Oznacza to minimalność w sensie
relacji mierzalności, \emph{nie} najmniejszej liczby bitów w pliku
ani automatycznie minimalnego błędu estymacji.
\end{przyklad}

\section{Rao--Blackwell: wykorzystanie informacji bez dodawania danych}
Wyobraźmy sobie, że estymujemy $p$ na podstawie pierwszego wyniku
$X_1$, choć mamy do dyspozycji całą próbę. Statystyka dostateczna
$S$ pozwala zastąpić ten niestabilny estymator jego \emph{średnią
warunkową przy znanym $S$}. Prowadzi to do mniejszej wariancji:
operacja warunkowania usuwa przypadkowość, która przy znanym $S$
nie niesie informacji o parametrze.

\begin{twierdzenie}[Twierdzenie Rao--Blackwella]\label{thm:rao-blackwell}
Niech $T$ będzie dostateczną statystyką dla rodziny
$\{P_\theta:\theta\in\Theta\}$ na standardowych przestrzeniach
borelowskich, a $U(X)$ rzeczywistym estymatorem takim, że
$\E_\theta|U|<\infty$ dla każdego $\theta$. Istnieje statystyka
$U^*=u(T)$, \emph{niezależna od wyboru parametru}, taka że
$U^*=\E_\theta[U\mid T]$ $P_\theta$-p.n. dla każdego $\theta$.
Jeśli $\E_\theta U=g(\theta)$, to także
$\E_\theta U^*=g(\theta)$.
Dla każdej wypukłej funkcji straty $a\mapsto L(\theta,a)$,
przy której obie wartości oczekiwane są określone,
\[
\E_\theta L(\theta,U^*)\le\E_\theta L(\theta,U).
\]
Dla straty kwadratowej, jeśli $U\in L^2(P_\theta)$, zachodzi
ściślejsza tożsamość
\[
\E_\theta(U-g(\theta))^2
=\E_\theta(U^*-g(\theta))^2
 +\E_\theta\bigl(U-U^*\bigr)^2.
\]
\end{twierdzenie}
\begin{proof}
Z dostateczności istnieje wspólne jądro warunkowe $K(t,\dd x)$.
Definiujemy $u(t)=\int_{\mathsf X}U(x)K(t,\dd x)$ w punktach, gdzie całka
jest bezwzględnie zbieżna, a poza nimi przyjmujemy $u(t)=0$.
Funkcja $u$ jest mierzalna (najpierw dla funkcji prostych,
potem przez aproksymację i zbieżność monotoniczną).
Ponieważ $U\in L^1(P_\theta)$ dla każdego $\theta$,
$u(T)=\E_\theta(U\mid T)$ prawie na pewno pod każdym $P_\theta$.
Z własności wieży $\E_\theta U^*=\E_\theta U$.
Warunkowa nierówność Jensena daje
$L(\theta,\E_\theta(U\mid T))\le
\E_\theta(L(\theta,U)\mid T)$;
po scałkowaniu otrzymujemy nierówność ryzyk.
Dla straty kwadratowej piszemy
$U-g(\theta)=(U-U^*)+(U^*-g(\theta))$.
Człon mieszany ma wartość oczekiwaną zero, bo
$\E_\theta[U-U^*\mid T]=0$ i $U^*-g(\theta)$ jest
$\sigma(T)$-mierzalny oraz całkowalny w $L^2$.
Rozwinięcie kwadratu daje tezę.
\end{proof}

\begin{zadanie}[Od pierwszej obserwacji do udziału sukcesów]
Dla $X_i\iid\operatorname{Bernoulli}(p)$ zacznij od
$U=X_1$, który estymuje $p$ nieobciążenie.
Wyznacz $\E(U\mid S)$, gdzie $S=\sum_iX_i$, i porównaj wariancje.
\par\emph{Rozwiązanie.} Warunkowo na $S=k$ każda pozycja
ma prawdopodobieństwo sukcesu $k/n$ z symetrii sekwencji,
więc $\E(X_1\mid S)=S/n=\bar X$.
Zatem $\Var_p(U)=p(1-p)$,
a $\Var_p(U^*)=p(1-p)/n$.
Zysk jest szczególnie duży, gdy $n$ rośnie: nie dodaliśmy
żadnej nowej obserwacji, wykorzystaliśmy tylko \emph{informację
już obecną w całej próbie}.
\end{zadanie}

\section{Zupełność: dlaczego nie ma dwóch różnych UMVU?}
Dostateczność mówi, że $T$ zachowuje informację o parametrze.
Zupełność zapewnia coś innego: nie istnieje niezerowa funkcja
$T$ o wartości oczekiwanej identycznie równej zero dla wszystkich
parametrów. To właśnie gwarantuje jednoznaczność estymatora
nieobciążonego opartego na statystyce dostatecznej.

\begin{definicja}[Zupełność rodziny statystyki]\label{def:completeness}
Rodzina rozkładów $\{\Law_\theta(T):\theta\in\Theta\}$
jest \emph{zupełna}, jeżeli dla każdej mierzalnej $a$ takiej, że
$\E_\theta|a(T)|<\infty$ dla wszystkich $\theta$, z warunku
$\E_\theta a(T)=0$ dla wszystkich $\theta$ wynika
$a(T)=0$ $P_\theta$-prawie na pewno dla wszystkich $\theta$.
\end{definicja}

\begin{twierdzenie}[Zupełność liczby sukcesów Bernoulliego]
\label{thm:complete-binomial}
Niech $S=\sum_{i=1}^nX_i$, gdzie
$X_i\iid\operatorname{Bernoulli}(p)$, $p\in(0,1)$.
Wtedy rodzina rozkładów $S$ jest zupełna.
\end{twierdzenie}
\begin{proof}
Niech $a:\{0,\ldots,n\}\to\R$ spełnia
$\E_p a(S)=0$ dla każdego $p\in(0,1)$. Mamy
\[
0=\sum_{k=0}^n a(k)\binom nk p^k(1-p)^{n-k}.
\]
Dzieląc przez $(1-p)^n>0$ i podstawiając
$z=p/(1-p)\in(0,\infty)$, otrzymujemy
$\sum_{k=0}^n a(k)\binom nk z^k=0$ dla każdego $z>0$.
Wielomian mający nieskończenie wiele zer jest zerowy,
więc wszystkie jego współczynniki znikają:
$a(k)\binom nk=0$. Ponieważ $\binom nk>0$,
$a(k)=0$ dla wszystkich $k$. Dowiedliśmy nawet silniejszego
wniosku niż wymagana równość prawie na pewno.
\end{proof}

\begin{twierdzenie}[Zupełność sumy zmiennych Poissona]
\label{thm:complete-poisson}
Niech $X_i\iid\Pois(\lambda)$, $\lambda\in(0,\infty)$,
a $S=\sum_iX_i\sim\Pois(n\lambda)$. Wtedy rodzina praw
$S$ jest zupełna.
\end{twierdzenie}
\begin{proof}
Jeśli $\E_\lambda|a(S)|<\infty$ i
$\E_\lambda a(S)=0$ dla każdego $\lambda>0$, to
\[
\sum_{k=0}^\infty a(k)\frac{(n\lambda)^k}{k!}=0,
\qquad\lambda>0.
\]
Szereg potęgowy po lewej stronie jest absolutnie zbieżny
dla każdego dodatniego argumentu, ponieważ odpowiednie
wartości bezwzględne są równe
$e^{n\lambda}\E_\lambda|a(S)|<\infty$.
Ma więc nieskończony promień zbieżności i reprezentuje
funkcję analityczną. Skoro znika na $(0,\infty)$,
wszystkie jej współczynniki Taylora w zerze muszą być zerowe.
To daje $a(k)=0$ dla każdego $k$.
\end{proof}


\begin{twierdzenie}[Zupełność pełnej jednoparametrowej rodziny wykładniczej]
\label{thm:complete-exponential-family}
Niech $Q_\eta$, $\eta\in J$, będzie rodziną rozkładów rzeczywistej
statystyki $T$ o postaci
$\dd Q_\eta(t)=e^{\eta t-A(\eta)}\,\dd\rho(t)$,
gdzie $\rho$ jest nieujemną miarą $\sigma$-skończoną,
$J\subset\R$ niepustym \emph{otwartym przedziałem}, a
$0<\int_{\R}e^{\eta t}\,\dd\rho(t)<\infty$ dla każdego $\eta\in J$.
Wtedy rodzina $\{Q_\eta:\eta\in J\}$ jest zupełna.
\end{twierdzenie}
\begin{proof}
Niech $a$ będzie mierzalna, $\E_\eta|a(T)|<\infty$
i $\E_\eta a(T)=0$ dla wszystkich $\eta\in J$.
Ustalmy $\eta_0\in J$ i wybierzmy $d>0$ z
$[\eta_0-2d,\eta_0+2d]\subset J$.
Zdefiniujmy skończoną miarę ze znakiem
$\nu(\dd t)=a(t)Q_{\eta_0}(\dd t)$. Dla $|u|\le d$
całka $\int_{\R}e^{u t}|\nu|(\dd t)$ jest skończona:
na $t\ge0$ ograniczamy ją przez stałą razy całkę
$|a(t)|e^{dt}$ względem $Q_{\eta_0}$,
a na $t<0$ analogicznie przez $e^{-dt}$;
obie są skończone dzięki całkowalności $a$ przy
$\eta_0+d$ i $\eta_0-d$. Funkcja
$F(z)=\int_{\R}e^{zt}\,\nu(\dd t)$ jest więc holomorficzna
w pasku $|\operatorname{Re}z|<d$; różniczkowanie pod całką
uzasadniają analogiczne oszacowania na mniejszych paskach.
Dla rzeczywistego $u\in(-d,d)$ mamy
\[
F(u)=e^{A(\eta_0+u)-A(\eta_0)}\E_{\eta_0+u}a(T)=0.
\]
Z twierdzenia o identyczności funkcji holomorficznych
$F(z)=0$ na całym pasku. W szczególności
$F(iv)=\int_{\R}e^{ivt}\,\nu(\dd t)=0$ dla każdego $v\in\R$.
Z jednoznaczności transformaty Fouriera skończonej miary
ze znakiem wynika $\nu=0$, czyli $a=0$ $Q_{\eta_0}$-p.n.
Gęstości $Q_\eta$ względem $\rho$ są dodatnie na tym samym
nośniku, więc wszystkie $Q_\eta$ są równoważne. Zatem
$a=0$ $Q_\eta$-p.n. dla każdego $\eta\in J$.
\end{proof}
\begin{przyklad}[Zupełność średniej w modelu normalnym]
Gdy $X_i\iid\Normal(\mu,\sigma^2)$ przy znanym $\sigma^2>0$
i $\mu\in\R$, rozkład $T=\sum_iX_i$ jest rodziną wykładniczą
z parametrem naturalnym $\eta=\mu/\sigma^2$ przebiegającym
całą prostą. Z poprzedniego twierdzenia wynika zupełność
statystyki $T$, a z faktoryzacji jej dostateczność.
Wobec tego $\bar X$ jest UMVU dla $\mu$. Nie należy
mieszać tego przypadku z sytuacją, w której również
$\sigma^2$ jest parametrem nieznanym.
\end{przyklad}

\begin{twierdzenie}[Lehmann--Scheffé: jednoznaczność UMVU]
\label{thm:lehmann-scheffe}
Niech $T$ będzie statystyką dostateczną i zupełną dla
$\{P_\theta:\theta\in\Theta\}$, a $g:\Theta\to\R$ estymandem.
Załóżmy, że istnieje estymator $U(X)$ o skończonym drugim
momencie pod każdym $P_\theta$ i taki, że
$\E_\theta U=g(\theta)$ dla każdego $\theta$.
Wtedy $U^*=\E_\theta(U\mid T)$, wybrany jako wspólna
wersja z twierdzenia Rao--Blackwella, jest \emph{jednostajnie
najlepszym nieobciążonym estymatorem} (UMVU): dla każdego
nieobciążonego estymatora $V\in L^2(P_\theta)$
$\Var_\theta(U^*)\le\Var_\theta(V)$ dla wszystkich $\theta$.
Ponadto jest jedyny z dokładnością do równości
$P_\theta$-prawie na pewno dla każdego $\theta$.
\end{twierdzenie}
\begin{proof}
Rao--Blackwell daje $U^*=u(T)$ nieobciążony i o skończonym
momencie kwadratowym. Dla dowolnego innego nieobciążonego
$V$ jego uśrednienie $V^*=v(T)$ jest nieobciążone i
$\Var_\theta V^*\le\Var_\theta V$. Ponieważ
$\E_\theta[u(T)-v(T)]=0$ dla każdego $\theta$ i różnica
jest całkowalna, zupełność implikuje
$u(T)=v(T)$ $P_\theta$-p.n. dla każdego $\theta$.
Stąd $\Var_\theta U^*=\Var_\theta V^*\le\Var_\theta V$.
Jeśli $V$ też jest UMVU, to w szczególności
$\Var_\theta(V)=\Var_\theta(U^*)$ i równość w tożsamości
Pitagorasa Rao--Blackwella wymusza
$V=V^*=U^*$ $P_\theta$-p.n.
\end{proof}

\section{Przykłady UMVU: praktyczne obliczenia}
\begin{zadanie}[Prawdopodobieństwo dwóch sukcesów]
\label{zad:umvu-p2}
Niech $X_1,\ldots,X_n\iid\operatorname{Bernoulli}(p)$,
$n\ge2$. Chcemy oszacować $p^2$ nieobciążenie.
\par\emph{Rozwiązanie.} Estymator $U=X_1X_2$ jest nieobciążony,
gdyż $\E U=\E X_1\E X_2=p^2$.
Niech $S=\sum_iX_i$. Warunkowo na $S=k$ prawdopodobieństwo,
że pierwsze dwie pozycje są sukcesami, wynosi
$k(k-1)/(n(n-1))$: wybieramy dwa z $k$ sukcesów i
porównujemy z wyborem dwóch spośród $n$ pozycji. Wobec tego
\[
\widehat{p^2}_{\rm UMVU}=\E(X_1X_2\mid S)
=\frac{S(S-1)}{n(n-1)}.
\]
Z twierdzeń~\ref{thm:complete-binomial} i~\ref{thm:lehmann-scheffe}
ma najmniejszą wariancję ze wszystkich estymatorów
nieobciążonych $p^2$. Jeżeli $n=4$ i zaobserwowano
$x=(1,0,1,0)$, to $S=2$ i wynik wynosi $1/6$, podczas gdy
wstawienie MLE $\widehat p=1/2$ do $p^2$ daje $1/4$.
To nie sprzeczność: $\widehat p^{\,2}$ jest obciążony jako
estymator $p^2$, ponieważ
$\E\bar X^2=p^2+p(1-p)/n$.
\end{zadanie}

\begin{zadanie}[Prognoza braku zgłoszeń w przyszłym okresie]
\label{zad:umvu-pois-zero}
W $n$ niezależnych okresach obserwujemy liczby zgłoszeń
$X_i\sim\Pois(\lambda)$, $\lambda>0$. Oszacuj nieobciążenie
prawdopodobieństwo braku zgłoszeń w \emph{następnym} takim
okresie, równe $e^{-\lambda}$.
\par\emph{Rozwiązanie.} Statystyka $S=\sum_iX_i\sim\Pois(n\lambda)$
jest dostateczna i zupełna. Funkcja tworząca rozkładu Poissona
daje $\E_\lambda a^S=\exp(n\lambda(a-1))$ dla $a\ge0$.
Dla $a=1-1/n$ (przy $n\ge1$) otrzymujemy
\[
\widehat{e^{-\lambda}}=\left(1-\frac1n\right)^S,
\quad \E_\lambda\widehat{e^{-\lambda}}=e^{-\lambda}.
\]
Dla $n=1$ przyjmujemy $0^0=1$, więc estymator jest wskaźnikiem
$\{S=0\}$. Z twierdzenia Lehmann--Scheffégo jest UMVU.
Przykładowo przy $n=5$ i sumie zgłoszeń $S=3$ wynik to
$(4/5)^3=0{,}512$, natomiast estymator przez wstawienie MLE
$\lambda=S/n$ daje $e^{-3/5}\approx0{,}549$.
Oba są sensowne dla innych kryteriów; UMVU nie gwarantuje
najmniejszego błędu średniokwadratowego wśród \emph{wszystkich}
estymatorów, tylko wśród nieobciążonych.
\end{zadanie}

\section{Statystyki poboczne i twierdzenie Basu}
\begin{definicja}[Statystyka poboczna]
Statystykę $A(X)$ nazywamy \emph{poboczną} (ang. \emph{ancillary}),
jeśli jej rozkład nie zależy od nieznanego parametru $\theta$.
\end{definicja}

\begin{twierdzenie}[Twierdzenie Basu]\label{thm:basu}
Niech $T$ będzie dostateczna i zupełna dla rodziny
$\{P_\theta:\theta\in\Theta\}$ na przestrzeni standardowej
borelowskiej, a $A$ statystyką poboczną o wartościach w przestrzeni
standardowej borelowskiej. Wtedy $T$ i $A$ są niezależne
pod każdym $P_\theta$.
\end{twierdzenie}
\begin{proof}
Ustalmy zdarzenie $B$ z $\sigma$-ciała wartości statystyki $A$.
Z dostateczności istnieje funkcja mierzalna $q_B$ \emph{niezależna
od $\theta$} taka, że
$\E_\theta[\ind{\{A\in B\}}\mid T]=q_B(T)$.
Z poboczności $P_\theta(A\in B)=c_B$ nie zależy od $\theta$.
Zatem dla każdego $\theta$
$\E_\theta[q_B(T)-c_B]=0$.
Funkcja $q_B-c_B$ jest ograniczona, więc zupełność wymusza
$q_B(T)=c_B$ $P_\theta$-p.n. dla wszystkich $\theta$.
Oznacza to, że
$P_\theta(A\in B\mid T)=P_\theta(A\in B)$;
ponieważ zachodzi dla dowolnego $B$, z definicji jest to
niezależność $T$ i $A$.
\end{proof}

\begin{przyklad}[Niezależność średniej i wariancji normalnej]
Dla $X_i\iid\Normal(\mu,\sigma^2)$ z \emph{znanym} $\sigma^2>0$
i nieznaną $\mu\in\R$, statystyka $T=\bar X$ jest
w tym pełnym modelu normalnym dostateczna i zupełna.
Statystyka $S^2=(n-1)^{-1}\sum_i(X_i-\bar X)^2$
ma rozkład niezależny od $\mu$; jest zatem poboczna dla
\emph{tego konkretnego modelu}. Twierdzenie Basu daje
niezależność $\bar X$ i $S^2$; w rozdziale 2 udowodniliśmy
ją już bezpośrednio przez obrót ortogonalny wektora normalnego.
Gdy $\sigma^2$ jest również nieznana, $S^2$ nie jest poboczna,
ponieważ jej rozkład zależy od $\sigma^2$. Nie wolno więc
mechanicznie przenosić argumentu między różnymi modelami.
\end{przyklad}


\begin{twierdzenie}[Maksimum próby jednostajnej jest zupełne]
\label{thm:uniform-complete}
Niech $n\ge2$ i $X_1,\ldots,X_n\iid U(0,\theta)$, $\theta>0$.
Wówczas $M=X_{(n)}$ jest statystyką dostateczną i zupełną.
Co więcej, statystyka $A=X_{(1)}/X_{(n)}$ jest poboczna;
na mocy twierdzenia Basu $A$ i $M$ są niezależne.
\end{twierdzenie}
\begin{proof}
Dla dodatnich obserwacji gęstość próby jest
$\theta^{-n}\ind{\{M\le\theta\}}$, więc faktoryzacja dowodzi
dostateczności $M$. Jego gęstość wynosi
$f_{M,\theta}(t)=n t^{n-1}\theta^{-n}\ind{\{0<t<\theta\}}$.
Niech $a$ spełnia $\E_\theta|a(M)|<\infty$ i
$\E_\theta a(M)=0$ dla każdego $\theta>0$.
Po pomnożeniu przez $\theta^n/n$ otrzymujemy
$G(\theta)=\int_0^\theta a(t)t^{n-1}\,\dd t=0$.
Funkcja $G$ jest lokalnie absolutnie ciągła na $(0,\infty)$,
gdyż $a(t)t^{n-1}$ jest lokalnie całkowalna; z twierdzenia
o różniczkowaniu całki mamy
$G'(\theta)=a(\theta)\theta^{n-1}=0$ dla niemal każdego
$\theta>0$. Zatem $a(M)=0$ $P_\theta$-p.n. dla każdego
$\theta$ i rodzina jest zupełna.
Ponadto przy zapisie $X_i=\theta U_i$, gdzie $U_i\iid U(0,1)$,
$A=U_{(1)}/U_{(n)}$ ma rozkład niezależny od $\theta$;
jest poboczna. Basu daje niezależność $A$ i $M$.
\end{proof}
\begin{uwaga}[Zaskakująca niezależność dwóch statystyk]
Maksimum i iloraz minimum do maksimum powstają z tych samych
obserwacji, a mimo to są niezależne w tym szczególnym modelu.
Twierdzenie nie mówi, że minimum i maksimum są niezależne
(na ogół nie są). Przy innych rozkładach lub niedokładnym
założeniu jednostajności omawiana niezależność może zniknąć.
\end{uwaga}

\section{Jaką metodę wybrać w konkretnym zadaniu?}
\begin{center}\small
\begin{tabular}{>{\raggedright\arraybackslash}p{0.25\textwidth}>{\raggedright\arraybackslash}p{0.31\textwidth}>{\raggedright\arraybackslash}p{0.31\textwidth}}
\toprule
\textbf{Problem} & \textbf{Najpierw sprawdź} & \textbf{Narzędzie i ograniczenie}\\
\midrule
Zredukować przechowywane dane & Czy istnieje statystyka dostateczna? & Faktoryzacja Fishera--Neymana; tylko w ustalonym modelu\\
Poprawić estymator nieobciążony & Czy znamy statystykę dostateczną? & Rao--Blackwell; poprawa MSE pod stratą kwadratową\\
Uzyskać UMVU & Czy statystyka jest \emph{zupełna i dostateczna}? & Lehmann--Scheffé; nie optymalizuje MSE wśród wszystkich obciążonych\\
Sprawdzić zbędność kolejności obserwacji & Czy próba rzeczywiście jest iid? & W modelach Bernoulliego i Poissona wystarcza suma; przy trendzie nie\\
Sprawdzić niezależność dwóch statystyk & Czy jedna jest zupełna dostateczna, druga poboczna? & Twierdzenie Basu; brak założeń wyklucza wniosek\\
Przewidzieć przyszłe zdarzenie & Jaki funkcjonał parametru chcemy oszacować? & UMVU lub estymator przez podstawienie MLE; zależnie od straty\\
\bottomrule
\end{tabular}
\end{center}

\paragraph{Granice stosowania.}
Dostateczność i zupełność są ściśle zależne od rodziny rozkładów.
Zachowanie sumy sukcesów jest bezstratne w modelu iid Bernoulliego,
lecz niekoniecznie przy zależności pomiędzy próbami lub niejednakowych
prawdopodobieństwach sukcesu. Ponadto \emph{najlepszy nieobciążony}
nie znaczy \emph{najlepszy przy dowolnej funkcji straty}.
Gdy wymagana jest kontrola ryzyka najgorszego przypadku,
wybór może paść na regułę minimaksową z poprzedniego rozdziału;
gdy najważniejsza jest odporność na odstające obserwacje,
na estymator odporny. W teorii statystycznej nie istnieje
optymalność niezależna od modelu i kryterium.

\Needspace{16\baselineskip}
\noindent\textbf{Przykład kodu: estymator UMVU dla $p^2$.}\par
\par
\Needspace{13\baselineskip}
\smallskip
\begin{lstlisting}[language=Python,caption={UMVU dla kwadratu prawdopodobieństwa sukcesu (Python).}]
x = [1, 0, 1, 0]
n, S = len(x), sum(x)
p2_umvu = S * (S - 1) / (n * (n - 1))
p2_plugin = (S/n)**2
print(p2_umvu, p2_plugin)  # 0.166666..., 0.25
\end{lstlisting}
\par
\Needspace{13\baselineskip}
\smallskip
\begin{lstlisting}[language=Julia,caption={UMVU dla kwadratu prawdopodobieństwa sukcesu (Julia).}]
x = [1, 0, 1, 0]
n, S = length(x), sum(x)
println((S*(S-1)/(n*(n-1)), (S/n)^2))
\end{lstlisting}

\paragraph{Podsumowanie.}
Metoda momentów i MLE \emph{konstruują} estymatory; teoria
dostateczności mówi, które aspekty danych zachowują informację
o parametrze; Rao--Blackwell poprawia nieobciążone estymatory,
a zupełność wraz z twierdzeniem Lehmanna--Scheffégo pozwala
uzasadnić ich optymalność w odpowiedniej klasie. Rozdział następny
poświęcimy informacji Fishera i dolnym ograniczeniom wariancji,
a następnie statystycznym twierdzeniom asymptotycznym.


% ================================================================
% ROZDZIAL 7 -- INFORMACJA FISHERA I DOLNE GRANICE
% ================================================================
\chapter{Informacja Fishera, efektywność i granice dokładności estymacji}
\label{ch:fisher-information}
\paragraph{Powiązania i odsyłacze.} Korzystamy z funkcji score z rozdziału~\ref{ch:likelihood} oraz z teorii estymatorów w rozdziale~\ref{ch:estimators}; przybliżenia asymptotyczne omawia rozdział~\ref{ch:asymptotic-estimation}.


\section{Motywacja: jak rozróżnialne są pobliskie parametry?}
Rozważmy dwa modele Bernoulliego: $p=0{,}50$ oraz $p=0{,}51$.
Jedna obserwacja nie pozwala ich wiarygodnie rozdzielić, nawet jeżeli
znamy dokładnie formułę rozkładu. Podobnie dwie wartości średniej
rozkładu normalnego są trudniejsze do rozróżnienia, gdy wariancja
pomiaru jest duża. \emph{Informacja Fishera} formalizuje lokalną
wrażliwość rozkładu obserwacji na zmianę parametru; nie jest miarą
informacji zawartej w pojedynczym, już zaobserwowanym wyniku.

\begin{definicja}[Score i informacja Fishera]\label{def:fisher-score}
Niech $\Theta\subset\R$ będzie przedziałem otwartym, a
$\{P_\theta:\theta\in\Theta\}$ rodziną rozkładów zdominowaną
przez miarę $\sigma$-skończoną $\mu$. Załóżmy, że istnieje wspólny zbiór $S\in\mathcal A$
taki, że $p_\theta(x)>0$ dla $x\in S$, $p_\theta(x)=0$
poza $S$ dla wszystkich $\theta$, oraz dla $\mu$-prawie
każdego $x\in S$ funkcja $\theta\mapsto p_\theta(x)$
jest różniczkowalna. Załóżmy ponadto, że
$s_\theta(x):=\partial_\theta\log p_\theta(x)$ należy do
$L^2(P_\theta)$. Funkcję $s_\theta$ nazywamy \emph{score}, a
$I(\theta):=\E_\theta[s_\theta(X)^2]$ \emph{informacją Fishera
w jednej obserwacji}. Informację w próbie zapisujemy $I_n(\theta)$.
\end{definicja}

Nie zakładamy automatycznie, że $\E_\theta s_\theta=0$.
Tę własność trzeba uzasadnić regularnością modelu. Nie należy też
mylić losowej \emph{informacji obserwowanej}
$j_\theta(X)=-\partial_\theta^2\log p_\theta(X)$ z jej
wartością oczekiwaną: dla ustalonej obserwacji wartości mogą być
bardzo różne.

\begin{twierdzenie}[Tożsamość score i addytywność informacji]
\label{thm:score-information}
Niech $p_\theta(x)$ będzie rodziną gęstości względem wspólnej
miary $\mu$ na przedziale parametrów $\Theta$. Załóżmy, że
$\partial_\theta p_\theta$ istnieje $\mu$-p.n. oraz że
$\partial_\theta\int_{\mathsf X} p_\theta(x)\,\dd\mu(x)
=\int_{\mathsf X}\partial_\theta p_\theta(x)\,\dd\mu(x)$; ponadto
$\int_{\mathsf X}(\partial_\theta p_\theta(x))^2/p_\theta(x)\,\dd\mu(x)<\infty$.
Wtedy $\E_\theta s_\theta=0$. Jeśli $X_1,\ldots,X_n$ są
niezależne o tym samym prawie $P_\theta$, to
$I_n(\theta)=nI(\theta)$.
Jeśli dodatkowo istnieje $\partial_\theta^2p_\theta$,
można dwukrotnie różniczkować całkę normalizującą i
$\E_\theta|\partial_\theta^2\log p_\theta(X)|<\infty$,
to $I(\theta)=-\E_\theta\partial_\theta^2\log p_\theta(X)$.
\end{twierdzenie}
\begin{proof}
Z definicji $s_\theta p_\theta=\partial_\theta p_\theta$,
więc $\E_\theta s_\theta=\int_{\mathsf X}\partial_\theta p_\theta(x)
\,\dd\mu(x)=\partial_\theta1=0$. Dla próby score jest sumą
$S_n=\sum_i s_\theta(X_i)$: wynika to z logarytmowania
iloczynu gęstości. Po podniesieniu do kwadratu
$\E S_n^2=\sum_i\E s_\theta(X_i)^2+
2\sum_{i<j}\E s_\theta(X_i)s_\theta(X_j)$.
Niezależność i zerowanie średniej eliminują wszystkie wyrazy
mieszane, zatem $I_n=nI$.
Z równości $\partial_\theta^2\log p=(\partial_\theta^2p)/p-
(\partial_\theta p)^2/p^2$ dostajemy
$\E_\theta\partial_\theta^2\log p=
\int_{\mathsf X}\partial_\theta^2p_\theta(x)\,\dd\mu(x)-I(\theta)=-I(\theta)$.
Ostatnia równość wykorzystuje dopuszczoną zamianę drugiej
pochodnej z całką i nie jest poprawna bez niej.
\end{proof}

\begin{zadanie}[Bernoulli: dlaczego skrajne prawdopodobieństwa są szczególne?]
Niech $X\sim\operatorname{Bernoulli}(p)$, $p\in(0,1)$.
Oblicz score i informację, a następnie podaj przybliżony
błąd standardowy $\widehat p$ z $n$ niezależnych prób.
\par\emph{Rozwiązanie.} Mamy
$\ell(p;x)=x\log p+(1-x)\log(1-p)$,
$s_p(x)=x/p-(1-x)/(1-p)=(x-p)/(p(1-p))$.
Ponieważ $\Var_p X=p(1-p)$, informacja wynosi
$I(p)=1/[p(1-p)]$ i $I_n(p)=n/[p(1-p)]$.
Dla $\widehat p=\bar X$ zachodzi dokładnie
$\Var_p\widehat p=p(1-p)/n$. Stąd po podstawieniu
$\widehat p$ orientacyjny błąd standardowy to
$\sqrt{\widehat p(1-\widehat p)/n}$.
Blisko $p=0$ lub $1$ wzór wymaga ostrożności: punkty brzegowe
nie należą do otwartej przestrzeni parametrów, a dla $\widehat p=0$
nie można wnioskować, że niepewność parametru faktycznie znika.
\end{zadanie}

\begin{przyklad}[Poisson z ekspozycją i normalny pomiar]
Jeśli niezależne $Y_i\sim\operatorname{Pois}(\lambda t_i)$,
$t_i>0$, to score jest $S(\lambda)=\sum_iY_i/\lambda-
\sum_it_i$, a informacja w całej próbie wynosi
$I(\lambda)=\sum_it_i/\lambda$. To \emph{łączny czas ekspozycji},
a nie liczba okresów, wyznacza podstawową skalę dokładności
estymacji intensywności. Jeśli $X_i\sim\Normal(\mu,\sigma^2)$
przy znanym $\sigma$, to $s_\mu(X_i)=(X_i-\mu)/\sigma^2$,
$I_n(\mu)=n/\sigma^2$: większy szum oznacza mniejszą informację.
\end{przyklad}

\section{Nierówność Craméra--Rao i warunki jej stosowania}
\begin{twierdzenie}[Nierówność Craméra--Rao dla funkcji parametru]
\label{thm:cramer-rao}
Niech $\Theta\subset\R$ będzie przedziałem otwartym,
$T$ rzeczywistą statystyką z $\E_\theta T^2<\infty$,
a $g(\theta)=\E_\theta T$. Załóżmy, że score
$S_\theta=\partial_\theta\log p_\theta(X)$ całej obserwacji
jest kwadratowo całkowalny, $\E_\theta S_\theta=0$,
$0<I_n(\theta)<\infty$ oraz że wolno różniczkować
$\int_{\mathsf X} T(x)p_\theta(x)\,\dd\mu(x)$ pod całką. Wtedy
\[
\Var_\theta(T)\ge\frac{[g'(\theta)]^2}{I_n(\theta)}.
\]
Równość dla ustalonego $\theta$ zachodzi (gdy $g'(\theta)\ne0$)
dokładnie wtedy, gdy $T-g(\theta)$ jest $P_\theta$-p.n.
proporcjonalna do $S_\theta$.
\end{twierdzenie}
\begin{proof}
Z dozwolonego różniczkowania i tożsamości
$\partial_\theta p_\theta=S_\theta p_\theta$ otrzymujemy
$g'(\theta)=\E_\theta[T S_\theta]$.
Ponieważ $\E_\theta S_\theta=0$, ta wartość jest kowariancją
$\Cov_\theta(T,S_\theta)$. Nierówność Cauchy'ego--Schwarza
daje $|g'(\theta)|^2\le\Var_\theta(T)\Var_\theta(S_\theta)$.
Ostatnia wariancja równa się $I_n(\theta)$, co po podzieleniu
przez dodatnią liczbę daje tezę. Warunek równości w
Cauchym--Schwarzu jest równoważny liniowej zależności
zmiennych scentrowanych w $L^2(P_\theta)$.
\end{proof}

\begin{wniosek}[Znak równości w modelu normalnym]
Jeśli $X_1,\ldots,X_n\iid\Normal(\mu,\sigma^2)$ przy
znanej $\sigma^2>0$, to dla nieobciążonego estymatora
$\bar X$ średniej $\mu$ mamy
$\Var_\mu(\bar X)=\sigma^2/n=1/I_n(\mu)$.
Zatem $\bar X$ osiąga granicę Craméra--Rao przy każdym $\mu$.
\end{wniosek}

\begin{uwaga}[Czemu rozkład jednostajny łamie naiwny rachunek?]
Dla próby $X_i\iid U(0,\theta)$ nośnik $(0,\theta)^n$
zależy od $\theta$. Na wnętrzu nośnika log-wiarygodność ma
pochodną $-n/\theta$, której wartość oczekiwana jest
$-n/\theta\ne0$. Zatem krok $\E_\theta S_\theta=0$
nie zachodzi i poprzednie twierdzenie nie ma zastosowania.
Co więcej, skorygowany estymator maksimum ma wariancję rzędu
$n^{-2}$, a nie $n^{-1}$. Nie jest to paradoks, lecz efekt
informacji pochodzącej z \emph{brzegu nośnika}.
\end{uwaga}

\section{Wiele parametrów i parametry zakłócające}
\begin{twierdzenie}[Macierzowa nierówność Craméra--Rao]
\label{thm:cr-matrix}
Niech $\Theta\subset\R^d$ będzie otwarta, a score całej próby
$S_\theta=\nabla_\theta\log p_\theta(X)$ spełnia
$\E_\theta S_\theta=0$, $I_n(\theta)=\E_\theta[S_\theta
S_\theta^{\mathsf T}]$ jest dodatnio określona oraz
$\nabla_\theta\E_\theta T=\E_\theta[T S_\theta]$ dla
$T\in L^2(P_\theta)$. Oznaczmy $g(\theta)=\E_\theta T$.
Wtedy
\[
\Var_\theta(T)\ge \nabla g(\theta)^{\mathsf T}
I_n(\theta)^{-1}\nabla g(\theta).
\]
\end{twierdzenie}
\begin{proof}
Niech $a=I_n^{-1}\nabla g$ i $Z=T-g-a^{\mathsf T}S_\theta$.
Ponieważ $\E Z=0$ oraz $\E[(T-g)S_\theta]=\nabla g$,
rozwinięcie kwadratu i $\E[S_\theta S_\theta^{\mathsf T}]=I_n$
dają $\E Z^2=\Var(T)-2a^{\mathsf T}\nabla g+a^{\mathsf T}I_na
=\Var(T)-\nabla g^{\mathsf T}I_n^{-1}\nabla g$.
Nieujemność $\E Z^2$ kończy dowód.
\end{proof}

Jeśli $\theta=(\psi,\eta)$, gdzie $\psi$ jest interesującym
nas parametrem, a $\eta$ parametrem zakłócającym, należy użyć
pełnej macierzy informacji. Dla odpowiedniego rozkładu blokowego
$I=\begin{psmallmatrix}I_{\psi\psi}&I_{\psi\eta}\\
I_{\eta\psi}&I_{\eta\eta}\end{psmallmatrix}$, gdy dolny blok
jest odwracalny, pierwszy element $I^{-1}$ to
$(I_{\psi\psi}-I_{\psi\eta}I_{\eta\eta}^{-1}I_{\eta\psi})^{-1}$.
Wyrażenie w nawiasie nazywamy informacją efektywną o $\psi$.
Nie wolno więc automatycznie używać $1/I_{\psi\psi}$,
jak gdyby $\eta$ było znane.

\begin{zadanie}[Normalny model z nieznaną wariancją]
Dla próby $X_i\iid\Normal(\mu,v)$, $\mu\in\R$, $v>0$,
oblicz macierz informacji parametrów $(\mu,v)$.
\par\emph{Rozwiązanie.} Pojedynczy score ma składowe
$s_\mu=(X-\mu)/v$ i
$s_v=-1/(2v)+(X-\mu)^2/(2v^2)$.
Z symetrii rozkładu normalnego składnik mieszany wynosi zero.
Z momentów centralnych $\E(X-\mu)^2=v$ i
$\E(X-\mu)^4=3v^2$ otrzymujemy
$I_1=\diag(1/v,1/(2v^2))$.
Dla $n$ niezależnych pomiarów $I_n=nI_1$.
Zatem parametr $v$ jest tu informacyjnie ortogonalny do $\mu$:
nieznana wariancja nie zmienia elementu $(I_n^{-1})_{\mu\mu}$,
choć istotnie zmienia sposób szacowania błędu standardowego z danych.
\end{zadanie}


\section{Lokalna geometria modeli: informacja a dywergencja KL}
Tożsamość informacyjna ma dodatkową interpretację, która jest
przydatna przy projektowaniu doświadczenia. Dwa bardzo bliskie
parametry nie dają dwóch zupełnie różnych rozkładów; w modelach
regularnych ich odległość w sensie dywergencji KL jest
\emph{drugiego rzędu} względem różnicy parametrów. Informacja
Fishera jest współczynnikiem tego rozwinięcia.

\begin{stwierdzenie}[Lokalne rozwinięcie dywergencji KL]
\label{prop:kl-fisher-expansion}
Niech $\Theta\subset\R$ będzie otwartym przedziałem, a
$\{p_\eta:\eta\in\Theta\}$ modelem zdominowanym o dodatnich
gęstościach na wspólnym nośniku. Ustalmy $\theta\in\Theta$.
Załóżmy, że funkcja
$K(h)=\E_\theta[\log p_\theta(X)-\log p_{\theta+h}(X)]$
jest dwukrotnie różniczkowalna w otoczeniu $h=0$, przy czym
wolno różniczkować pod wartością oczekiwaną, a tożsamości
informacyjne obowiązują dla $\eta$ w tym otoczeniu.
Wtedy
$K(h)=\tfrac12 I(\theta)h^2+o(h^2)$, gdy $h\to0$.
\end{stwierdzenie}
\begin{proof}
Ponieważ $K(0)=0$ i $K'(0)=-\E_\theta s_\theta(X)=0$,
wystarczy policzyć drugą pochodną w zerze.
Mamy $K''(0)=-\E_\theta\partial_\theta^2\log p_\theta(X)
=I(\theta)$. Rozwinięcie Taylora z resztą Peana daje tezę.
Przytoczone założenia są istotne: przy nośniku zależnym od
parametru nawet wartość $K(h)$ może być nieskończona w jednej
stronie zera.
\end{proof}

\begin{przyklad}[Dlaczego powtarzanie pomiaru pomaga?]
Dla $P_\mu=\Normal(\mu,\sigma^2)$ przy znanym $\sigma>0$
rachunek gaussowski daje dokładnie
$D_{\rm KL}(P_\mu\|P_{\mu+h})=h^2/(2\sigma^2)$.
Dla $n$ niezależnych pomiarów dywergencja praw wspólnych jest
$n$ razy większa. Dlatego przy ustalonym stosunku sygnału
do szumu wzrost liczby niezależnych pomiarów zwiększa
rozróżnialność parametrów. Jeśli jednak wielokrotnie mierzymy
ten sam obiekt przy wspólnym błędzie systematycznym, model
iloczynowy nie jest poprawny i wniosek o addytywności nie musi
pozostać prawdziwy.
\end{przyklad}

\section{Efektywność i ryzyko: co rzeczywiście mierzymy?}
W statystyce mówi się czasem, że estymator jest \emph{efektywny},
gdy osiąga granicę Craméra--Rao w danym regularnym modelu.
To jest własność \emph{względem konkretnej funkcji straty i klasy
estymatorów}, a nie stwierdzenie, że jest to najlepsza procedura
do wszystkich zastosowań. Granica Craméra--Rao dotyczy w podanej
wersji estymatorów o ustalonej wartości oczekiwanej $g(\theta)$.
Estymator obciążony może mieć mniejsze MSE; metody odporne mogą
być lepsze, gdy część pomiarów jest zanieczyszczona.

\begin{zadanie}[Dwa sposoby estymacji wariancji normalnej]
Dla $X_i\iid\Normal(\mu,\sigma^2)$ o nieznanych $\mu$ i
$\sigma^2$ porównaj estymator nieobciążony
$S^2=(n-1)^{-1}\sum_i(X_i-\bar X)^2$ i MLE
$V_n=n^{-1}\sum_i(X_i-\bar X)^2$ przy stracie kwadratowej.
\par\emph{Rozwiązanie.} Z rozkładu
$(n-1)S^2/\sigma^2\sim\chi^2_{n-1}$ wynika
$\E S^2=\sigma^2$ i $\Var S^2=2\sigma^4/(n-1)$.
Dla $V_n=(n-1)S^2/n$ mamy
$\E V_n=(n-1)\sigma^2/n$ i
$\Var V_n=2(n-1)\sigma^4/n^2$. Dlatego
\[
\MSE_\theta(V_n)=\frac{2(n-1)+1}{n^2}\sigma^4
=\frac{2n-1}{n^2}\sigma^4,
\qquad
\MSE_\theta(S^2)=\frac{2}{n-1}\sigma^4.
\]
Dla $n\ge2$ pierwszy z tych błędów jest mniejszy.
Ten przykład przypomina, że minimalna wariancja w klasie
estymatorów nieobciążonych nie rozstrzyga optymalności przy
uwzględnieniu estymatorów obciążonych.
\end{zadanie}

\section{Kiedy granica informacyjna daje praktyczną odpowiedź?}
\begin{center}\small
\begin{tabular}{>{\raggedright\arraybackslash}p{0.25\textwidth}>{\raggedright\arraybackslash}p{0.32\textwidth}>{\raggedright\arraybackslash}p{0.30\textwidth}}
\toprule
\textbf{Cel} & \textbf{Co można obliczyć?} & \textbf{Czego nie wolno wnioskować?}\\
\midrule
Porównanie dwóch nieobciążonych estymatorów & Wariancje, MSE i dolną granicę Craméra--Rao & Osiąganie granicy nie jest gwarantowane\\
Projektowanie pomiaru & Informację na obserwację lub ekspozycję & Więcej rekordów nie zawsze znaczy więcej informacji przy zależności\\
Parametr z nuisance & Macierz informacji i dopełnienie Schura & Nie wolno pominąć współparametrów\\
Mała próba i punkt brzegowy & Dokładny rozkład estymatora, przedział dokładny & Asymptotyczna normalność może być złym przybliżeniem\\
Model o zmiennym nośniku & Bezpośredni rachunek rozkładu statystyki & Nie stosować regularnego Craméra--Rao bez założeń\\
\bottomrule
\end{tabular}
\end{center}

\begin{zadanie}[Jednostki informacji przy obserwacji zdarzeń]
Dwie maszyny pracowały odpowiednio $20$ i $80$ godzin;
zaobserwowano łącznie $7$ awarii. W modelu niezależnych zliczeń
Poissona o wspólnej intensywności $\lambda$ oblicz MLE i
przybliżony błąd standardowy.
\par\emph{Rozwiązanie.} Łączna ekspozycja to $T=100$ godzin,
więc $\widehat\lambda=7/100=0{,}07$ awarii na godzinę.
Obserwowana krzywizna log-wiarygodności przy maksimum wynosi
$j(\widehat\lambda)=7/\widehat\lambda^2$;
informacja Fishera to $T/\lambda$.
Oba sposoby po podstawieniu MLE prowadzą do
$\widehat{\mathrm{se}}=\sqrt{\widehat\lambda/T}
\approx0{,}0265$ na godzinę. Jest to oszacowanie asymptotyczne:
przy zaledwie siedmiu zdarzeniach warto porównać je z dokładną
procedurą opartą na rozkładzie Poissona.
\end{zadanie}

% ================================================================
% ROZDZIAL 8 -- ASYMPTOTYCZNA TEORIA ESTYMACJI
% ================================================================
\chapter{Asymptotyczna teoria estymacji i niepewność w praktyce}
\label{ch:asymptotic-estimation}
\paragraph{Powiązania i odsyłacze.} Wymagane prawa wielkich liczb i CTG przypomina tom I, rozdziały~\ref{I-ch:lln} i~\ref{I-ch:clt}; informację Fishera wprowadziliśmy w rozdziale~\ref{ch:fisher-information}.


\section{Dlaczego rozkład dokładny bywa nieosiągalny?}
Dla średniej z próby normalnej dysponujemy dokładnym rozkładem.
Dla MLE skomplikowanego modelu regresyjnego znamy natomiast zwykle
jedynie równanie optymalizacji. Zamiast próbować obliczać pełną
gęstość estymatora, pytamy, co dzieje się przy $n\to\infty$.
Otrzymujemy wtedy narzędzia do konstrukcji błędów standardowych,
przedziałów i testów. Nie należy jednak mylić \emph{asymptotycznej
normalności} ze zdaniem, że w dowolnej małej próbie rozkład błędu
jest w przybliżeniu normalny.

\section{Jednostajna zbieżność kryterium i zgodność estymatora}
\begin{lemat}[Jednostajne prawo wielkich liczb przy warunku Lipschitza]
\label{lem:uniform-lln}
Niech $\Theta\subset\R^d$ będzie niepustym zbiorem zwartym,
$X_1,X_2,\ldots$ niezależnymi zmiennymi o tym samym prawie $P$,
a $m:\mathsf X\times\Theta\to\R$ funkcją mierzalną względem $x$
i ciągłą względem $\theta$. Załóżmy, że dla ustalonego
$\theta_*\in\Theta$ zachodzi $\E|m(X_1,\theta_*)|<\infty$
oraz że istnieje nieujemna zmienna $L(X_1)$ z $\E L(X_1)<\infty$
taka, że dla wszystkich $\theta,\eta\in\Theta$
$|m(X_1,\theta)-m(X_1,\eta)|\le L(X_1)\|\theta-\eta\|$
prawie na pewno na jednym wspólnym zdarzeniu. Wtedy
\[
\sup_{\theta\in\Theta}\left|\frac1n\sum_{i=1}^n
m(X_i,\theta)-\E m(X_1,\theta)\right|
\longrightarrow0\quad\text{prawie na pewno}.
\]
\end{lemat}
\begin{proof}
Dzięki zwartości $\Theta$ dla każdego $\delta>0$ istnieje
skończona $\delta$-sieć $\{\theta_1,\ldots,\theta_N\}$.
Dla $\theta\in\Theta$ wybieramy węzeł odległy o mniej niż
$\delta$. Nierówność trójkąta i wspólne oszacowanie Lipschitza
dają
\[
\sup_{\theta\in\Theta}|P_nm_\theta-Pm_\theta|
\le\max_{1\le j\le N}|P_nm_{\theta_j}-Pm_{\theta_j}|
+\delta(P_nL+PL),
\]
gdzie $P_nf=n^{-1}\sum_i f(X_i)$ i $Pf=\E f(X_1)$.
Mocne prawo wielkich liczb stosowane do skończenie wielu $m_{\theta_j}$
i do $L$ zapewnia, że granica górna lewej strony jest nie większa
niż $2\delta\E L$ prawie na pewno. Bierzemy przeliczalny
ciąg $\delta=1/k$ i przecinamy odpowiadające mu zdarzenia
prawdopodobieństwa jeden; przejście $k\to\infty$ kończy dowód.
\end{proof}

\begin{twierdzenie}[Zgodność maksymalizatora]\label{thm:m-est-consistency}
Przy założeniach lematu~\ref{lem:uniform-lln} oznaczmy
$M(\theta)=\E m(X_1,\theta)$ i
$M_n(\theta)=n^{-1}\sum_i m(X_i,\theta)$.
Jeżeli $M$ ma dokładnie jeden maksymalizator $\theta_0\in\Theta$
i $\widehat\theta_n$ jest mierzalnym maksymalizatorem $M_n$
na $\Theta$ (albo mierzalnym przybliżonym maksymalizatorem
z błędem $o_P(1)$), to $\widehat\theta_n\to\theta_0$
według prawdopodobieństwa. Dla maksymalizatorów dokładnych
zachodzi nawet zbieżność prawie na pewno.
\end{twierdzenie}
\begin{proof}
Ponieważ $M$ jest ciągła na zwartej $\Theta$, dla każdego
$\varepsilon>0$ zbiór
$F_\varepsilon=\{\theta\in\Theta:\|\theta-\theta_0\|
\ge\varepsilon\}$ jest zwarty. Jeśli jest niepusty,
jedyność maksimum daje dodatnią lukę
$\Delta_\varepsilon=M(\theta_0)-\max_{F_\varepsilon}M>0$.
Na zdarzeniu $\sup_\theta|M_n-M|<\Delta_\varepsilon/3$
dla $\theta\in F_\varepsilon$ mamy
$M_n(\theta)\le M(\theta_0)-2\Delta_\varepsilon/3$,
a $M_n(\theta_0)\ge M(\theta_0)-\Delta_\varepsilon/3$.
Żaden dokładny maksymalizator nie należy więc do
$F_\varepsilon$. Lemat zapewnia tę nierówność od pewnego
indeksu prawie na pewno. Dla przybliżonych maksymalizatorów
potrzebujemy dodatkowo, by ich błąd optymalizacji był mniejszy niż
$\Delta_\varepsilon/3$; zachodzi to z prawdopodobieństwem dążącym
do jedności. Dowód jest zakończony.
\end{proof}

\begin{przyklad}[Zgodność MLE przez dywergencję KL]
W identyfikowalnym, poprawnie określonym modelu zdominowanym
$P_\theta$ średnie log-wiarygodności spełniają, przy warunkach
całkowalności, $\E_{\theta_0}\log p_\theta(X)
=\E_{\theta_0}\log p_{\theta_0}(X)
-D_{\mathrm{KL}}(P_{\theta_0}\|P_\theta)$.
Zatem funkcja graniczna ma jednoznaczne maksimum w $\theta_0$.
Aby zastosować poprzednie twierdzenie, trzeba \emph{dodatkowo}
wykazać jednostajną zbieżność po $\theta$; sama identyfikowalność
ani punktowe MPWL nie wystarczają.
\end{przyklad}

\section{Asymptotyczna normalność M-estymatorów}
\begin{twierdzenie}[Rozwinięcie równania estymacyjnego]
\label{thm:m-est-normality}
Niech $\theta_0\in\R^d$ i $\widehat\theta_n\to\theta_0$
według prawdopodobieństwa. Niech $\Psi_n:\Theta\to\R^d$
będzie losową funkcją klasy $C^1$ w ustalonym wypukłym otoczeniu
$U$ punktu $\theta_0\in U\subset\Theta$. Załóżmy, że
$\Prob(\widehat\theta_n\in U)\to1$,
$\Prob(\Psi_n(\widehat\theta_n)=0)\to1$,
$\sqrt n\Psi_n(\theta_0)\xrightarrow{d}\Normal_d(0,B)$,
gdzie $B$ jest nieujemnie określona, oraz że istnieje
ciągła funkcja $H:U\to\R^{d\times d}$ taka, że
$A:=H(\theta_0)$ jest nieosobliwa i
\[
\sup_{\theta\in U}\|D\Psi_n(\theta)-H(\theta)\|
\xrightarrow{P}0.
\]
Wtedy
\[
\sqrt n(\widehat\theta_n-\theta_0)
\xrightarrow{d}\Normal_d(0,A^{-1}B A^{-\mathsf T}).
\]
\end{twierdzenie}
\begin{proof}
Na zdarzeniu $\widehat\theta_n\in U$ i
$\Psi_n(\widehat\theta_n)=0$ wzór Newtona--Leibniza dla
funkcji wektorowej na odcinku daje
$0=\Psi_n(\theta_0)+A_n(\widehat\theta_n-\theta_0)$,
gdzie $A_n=\int_0^1D\Psi_n(\theta_0+t(\widehat\theta_n-
\theta_0))\,\dd t$.
Jednostajna zbieżność $D\Psi_n$ do $H$, ciągłość $H$ w
$\theta_0$ oraz zgodność $\widehat\theta_n$ zapewniają
$A_n\to H(\theta_0)=A$ według prawdopodobieństwa.
Z nieosobliwości $A$ wynika odwracalność
$A_n$ z prawdopodobieństwem dążącym do jedności.
Wobec tego $\sqrt n(\widehat\theta_n-\theta_0)
=-A_n^{-1}\sqrt n\Psi_n(\theta_0)$ na takim zdarzeniu.
Twierdzenie Slutsky'ego i liniowa transformacja rozkładu normalnego
dają tezę. Dowód nie wymaga założenia normalności danych.
\end{proof}

W praktyce $\Psi_n(\theta)=n^{-1}\sum_i\psi(X_i,\theta)$,
$A=\E D_\theta\psi(X_i,\theta_0)$ i
$B=\E[\psi(X_i,\theta_0)\psi(X_i,\theta_0)^{\mathsf T}]$,
przy założeniu zerowej średniej i skończonego drugiego momentu
$\psi$. Wówczas \emph{macierz kanapkowa}
$A^{-1}BA^{-\mathsf T}$ jest podstawą odpornego na błędną
specyfikację oszacowania niepewności, o ile estymand jest
właściwie określony przez równanie $\E\psi=0$.

\begin{wniosek}[Regularna asymptotyka MLE]
\label{cor:mle-asymptotic}
Niech $\Theta\subset\R^d$ będzie otwarta, a
$X_1,X_2,\ldots$ będą iid o rozkładzie $P_{\theta_0}$,
$\theta_0\in\Theta$. Niech $s_\theta(x)$ będzie score,
różniczkowalnym względem $\theta$ w wypukłym otoczeniu
$U$ punktu $\theta_0$, i załóżmy:
\begin{enumerate}
\item $\E_{\theta_0}s_{\theta_0}(X_1)=0$,
$\E_{\theta_0}\|s_{\theta_0}(X_1)\|^2<\infty$ oraz
$I(\theta_0)=\E_{\theta_0}[s_{\theta_0}s_{\theta_0}^{\mathsf T}]$
jest dodatnio określona;
\item zachodzi druga tożsamość informacyjna
$\E_{\theta_0}D_\theta s_{\theta_0}(X_1)=-I(\theta_0)$;
\item dla $\Psi_n(\theta)=n^{-1}\sum_{i=1}^n s_\theta(X_i)$
istnieje ciągła funkcja $H:U\to\R^{d\times d}$ z
$H(\theta_0)=-I(\theta_0)$, taka że
$\sup_{\theta\in U}\|D\Psi_n(\theta)-H(\theta)\|\to0$
według prawdopodobieństwa;
\item istnieje mierzalny estymator $\widehat\theta_n$ taki,
że $\widehat\theta_n\to\theta_0$ według prawdopodobieństwa,
$P(\widehat\theta_n\in U,\;\Psi_n(\widehat\theta_n)=0)\to1$.
\end{enumerate}
Wtedy
\[
\sqrt n(\widehat\theta_n-\theta_0)
\xrightarrow{d}\Normal_d(0,I(\theta_0)^{-1}).
\]
\end{wniosek}
\begin{proof}
W poprzednim twierdzeniu wybieramy $\Psi_n(\theta)=
n^{-1}\sum_i s_\theta(X_i)$. Z CLT wynika $B=I(\theta_0)$,
a z tożsamości informacyjnej
$A=\E\partial_\theta s_{\theta_0}=-I(\theta_0)$.
Po wstawieniu do macierzy kanapkowej otrzymujemy
$(-I)^{-1}I(-I)^{-\mathsf T}=I^{-1}$.
\end{proof}

\section{Metoda delty i przedziały asymptotyczne}
\begin{twierdzenie}[Wielowymiarowa metoda delty]
Niech $\sqrt n(T_n-\theta_0)\xrightarrow{d}Z$ w $\R^d$,
a $g:\R^d\to\R^k$ będzie różniczkowalna w $\theta_0$.
Wtedy
$\sqrt n(g(T_n)-g(\theta_0))\xrightarrow{d}Dg(\theta_0)Z$.
Jeśli $Z\sim\Normal_d(0,V)$, granica jest normalna o macierzy
kowariancji $Dg(\theta_0)V Dg(\theta_0)^{\mathsf T}$.
\end{twierdzenie}
\begin{proof}
Z różniczkowalności
$g(\theta_0+h)=g(\theta_0)+Dg(\theta_0)h+o(\|h\|)$.
Ciąg $\sqrt n(T_n-\theta_0)$ jest ciasny, a więc
$T_n\to\theta_0$ według prawdopodobieństwa.
Mnożąc resztę Taylora przez $\sqrt n$ dostajemy iloczyn
$\sqrt n\|T_n-\theta_0\|=O_P(1)$ i
$o_P(1)$, czyli składnik dążący do zera według
prawdopodobieństwa. Teza wynika ze Slutsky'ego.
\end{proof}

\begin{stwierdzenie}[Asymptotyczny przedział Waldowski]
\label{prop:wald-ci}
Niech $\widehat\theta_n$ estymuje rzeczywisty parametr
$\theta_0$ i załóżmy
$\sqrt n(\widehat\theta_n-\theta_0)\xrightarrow{d}
\Normal(0,V)$ dla $0<V<\infty$.
Jeśli $\widehat V_n\to V$ według prawdopodobieństwa
oraz $\widehat V_n>0$ z prawdopodobieństwem dążącym do
jedności, to przedział
$[\widehat\theta_n-z_{1-\alpha/2}\sqrt{\widehat V_n/n},
\widehat\theta_n+z_{1-\alpha/2}\sqrt{\widehat V_n/n}]$
ma asymptotyczne prawdopodobieństwo pokrycia $1-\alpha$,
gdzie $z_q$ jest kwantylem rzędu $q$ rozkładu $\Normal(0,1)$.
\end{stwierdzenie}
\begin{proof}
Z oszacowania $\widehat V_n\to V>0$ i twierdzenia Slutsky'ego
mamy $(\widehat\theta_n-\theta_0)/\sqrt{\widehat V_n/n}
\xrightarrow{d}\Normal(0,1)$. Ponieważ rozkład graniczny
jest ciągły, prawdopodobieństwo znalezienia się pomiędzy
$\pm z_{1-\alpha/2}$ dąży do $1-\alpha$; jest to dokładnie
zdarzenie pokrycia.
\end{proof}

\begin{uwaga}[Wariancja estymowana i zdarzenia wyjątkowe]
Warunek dodatniości $\widehat V_n$ z prawdopodobieństwem
dążącym do jedności nie wymaga, by oszacowanie było dodatnie
dla każdej próby. Na zdarzeniu, na którym nie jest dodatnie,
należy jawnie określić konwencję postępowania, np. nie
raportować przedziału Waldowskiego. Przy małej liczbie sukcesów
w próbie Bernoulliego metoda Wilsona zwykle zachowuje się
lepiej niż symetryczny przedział Walda.
\end{uwaga}

\begin{zadanie}[7 wad w próbie 100 sztuk: Wald czy Wilson?]
\label{zad:wilson-100}
Niech $X_i\iid\operatorname{Bernoulli}(p)$, $n=100$,
$S=7$ i $\widehat p=0{,}07$. Dla $z=1{,}96$ (poziom około
$95\%$) porównaj asymptotyczny przedział Waldowski i
przedział Wilsona uzyskany przez odwrócenie nierówności
$|\widehat p-p|\le z\sqrt{p(1-p)/n}$.
\par\emph{Rozwiązanie.} Przedział Waldowski to
$\widehat p\pm z\sqrt{\widehat p(1-\widehat p)/n}$.
Dla Wilsona przenosimy wszystkie składniki nierówności na
jedną stronę i otrzymujemy nierówność kwadratową względem $p$.
Jej pierwiastki są
\[
\frac{\widehat p+z^2/(2n)\pm
z\sqrt{\widehat p(1-\widehat p)/n+z^2/(4n^2)}}{1+z^2/n}.
\]
W odróżnieniu od Walda Wilson uwzględnia zależność wariancji
od testowanej wartości $p$ i pozostaje w $[0,1]$.
Nie jest jednak dokładnym przedziałem o nominalnym pokryciu
dla dowolnego $n$. Gdy $S=0$, Wald daje przedział zdegenerowany
w zerze, podczas gdy Wilson zachowuje dodatni górny kraniec.
\end{zadanie}


\section{Metoda delty drugiego rzędu: gdy pierwsza pochodna znika}
\begin{twierdzenie}[Jednowymiarowa metoda delty drugiego rzędu]
\label{thm:delta-second}
Niech $\sqrt n(T_n-\theta_0)\xrightarrow{d}Z$ oraz niech
$g:\R\to\R$ będzie dwukrotnie różniczkowalna w otoczeniu
$\theta_0$ z ciągłą drugą pochodną w $\theta_0$ i
$g'(\theta_0)=0$. Wtedy
\[
n\bigl(g(T_n)-g(\theta_0)\bigr)
\xrightarrow{d}\frac12g''(\theta_0)Z^2.
\]
\end{twierdzenie}
\begin{proof}
Wzór Taylora daje
$g(\theta_0+h)-g(\theta_0)=\frac12g''(\theta_0)h^2+
o(h^2)$. Ponieważ $\sqrt n(T_n-\theta_0)$ jest ciasny,
$T_n\to\theta_0$ według prawdopodobieństwa oraz
$n(T_n-\theta_0)^2=O_P(1)$. Przemnożenie rozwinięcia
przez $n$ daje resztę $o_P(1)$ i wskazaną granicę na mocy
ciągłego odwzorowania i Slutsky'ego.
\end{proof}
\begin{zadanie}[Estymacja kwadratu średniej w zerze]
Niech $X_i\iid\Normal(0,1)$ i $T_n=\bar X$. Jak zachowuje
się estymator $T_n^2$ wartości $\mu^2$ przy $\mu=0$?
\par\emph{Rozwiązanie.} Z dokładnego rozkładu
$\sqrt n\bar X\sim\Normal(0,1)$ otrzymujemy
$n\bar X^2\sim\chi_1^2$ dla każdego $n$.
Zwykła metoda delty dla $g(\mu)=\mu^2$ daje zdegenerowaną
granicę, gdyż $g'(0)=0$. Metoda drugiego rzędu odzyskuje
niegaussowski rozkład graniczny. W tym modelu warto również
rozważyć nieobciążony estymator $\bar X^2-S^2/n$,
który może przyjmować wartości ujemne: sama nieobciążoność
nie gwarantuje zachowania ograniczeń estymandu.
\end{zadanie}

\section{Model błędnie określony i wariancja kanapkowa}
\begin{zadanie}[Zliczenia z nadmierną dyspersją]
Zaobserwowano liczby zgłoszeń w dziesięciu porównywalnych
okresach: $0,0,0,0,1,1,2,2,2,10$. Oblicz estymator średniej
intensywności i porównaj błędy standardowe pod modelem Poissona
oraz bez zakładania równości średniej i wariancji.
\par\emph{Rozwiązanie.} Średnia wynosi $\bar x=1{,}8$,
a wariancja z mianownikiem $n-1$ wynosi
$s^2=9{,}0667$. Gdy przyjmiemy model
$X_i\iid\Pois(\lambda)$, MLE to $\widehat\lambda=\bar x$,
a modelowy błąd standardowy wynosi
$\sqrt{\widehat\lambda/n}\approx0{,}424$.
Jeżeli obserwacje są nadal iid, ale ich wspólny rozkład nie
musi być poissonowski, z CLT dla średniej i spójności
$s^2$ wynika oszacowanie błędu standardowego
$s/\sqrt n\approx0{,}952$.
Rozbieżność jest bardzo duża i wskazuje na nadmierną dyspersję.
Nie dowodzi sama w sobie, jaki jest jej mechanizm. Przyczyną
mogą być zmienne warunki pracy, heterogeniczność, nadmiar zer
albo obserwacja odstająca. W analizie należy porównać model
ujemnego dwumianowego, modele hierarchiczne i diagnostykę danych.
\end{zadanie}

\begin{stwierdzenie}[Kanapkowa wariancja estymatora średniej]
Niech $X_1,X_2,\ldots$ będą iid z $\E X_i=\mu$ oraz
$0<\Var(X_i)=v<\infty$. Dla równania estymacyjnego
$\Psi_n(\theta)=n^{-1}\sum_i(X_i-\theta)$ otrzymujemy
$\widehat\theta_n=\bar X$, $A=-1$, $B=v$ i
asymptotyczną wariancję $v/n$. Gdy dodatkowo błędnie
założymy model Poissona, modelowa informacja daje zamiast
tego $\mu/n$, co jest poprawne tylko przy $v=\mu$.
\end{stwierdzenie}
\begin{proof}
Score zastępczy $\psi(x,\theta)=x-\theta$ ma
$\E\psi(X,\mu)=0$, pochodną $\partial_\theta\psi=-1$,
a wariancję $\Var(\psi(X,\mu))=v$. Twierdzenie
\ref{thm:m-est-normality} daje
$\sqrt n(\bar X-\mu)\xrightarrow{d}\Normal(0,v)$.
Pod modelem Poissona $v=\mu$; poza nim równanie średniej
pozostaje sensowne, ale przy błędnej wariancji nie otrzymujemy
poprawnej kalibracji niepewności.
\end{proof}

\section{Trzy testy asymptotyczne z tej samej wiarygodności}
\begin{twierdzenie}[Jednoparametrowe twierdzenie Wilksa przy regularności]
\label{thm:wilks-one}
Niech $\theta_0$ leży we wnętrzu otwartego przedziału
$\Theta$. Dla próby iid rozważamy dwukrotnie różniczkowalną
log-wiarygodność $\ell_n$. Załóżmy istnienie zgodnego
maksymalizatora wewnętrznego $\widehat\theta_n$, takiego że
$\ell_n'(\widehat\theta_n)=0$, i niech
$\ell_n'(\theta_0)/\sqrt n\xrightarrow{d}\Normal(0,I_0)$,
gdzie $I_0>0$, oraz że istnieje wypukłe otoczenie
$U\subset\Theta$ parametru $\theta_0$ i ciągła funkcja
$J:U\to\R$ z $J(\theta_0)=I_0$, takie że
$\sup_{\theta\in U}|-\ell_n''(\theta)/n-J(\theta)|
\xrightarrow{P}0$. Wtedy
\[
2[\ell_n(\widehat\theta_n)-\ell_n(\theta_0)]
\xrightarrow{d}\chi_1^2.
\]
\end{twierdzenie}
\begin{proof}
Z twierdzenia Taylora dla pochodnej
$0=\ell_n'(\widehat\theta_n)
=\ell_n'(\theta_0)+\ell_n''(\xi_n)(\widehat\theta_n-
\theta_0)$, gdzie $\xi_n$ leży pomiędzy parametrami.
Zatem $\sqrt n(\widehat\theta_n-\theta_0)
=[\ell_n'(\theta_0)/\sqrt n]/[-\ell_n''(\xi_n)/n]$
i ma rozkład graniczny $\Normal(0,I_0^{-1})$.
Rozwijając z kolei $\ell_n(\widehat\theta_n)$ wokół
$\widehat\theta_n$, gdzie pierwsza pochodna znika, otrzymujemy
$2(\ell_n(\widehat\theta_n)-\ell_n(\theta_0))
= -\ell_n''(\zeta_n)(\widehat\theta_n-\theta_0)^2$,
przy $\zeta_n$ między $\theta_0$ a $\widehat\theta_n$.
Ponieważ $\widehat\theta_n\to\theta_0$, również punkty
$\xi_n$ i $\zeta_n$ dążą do $\theta_0$ według
prawdopodobieństwa. Ciągłość $J$ i jednostajna zbieżność
krzywizny dają zatem
$-\ell_n''(\xi_n)/n\to I_0$ i
$-\ell_n''(\zeta_n)/n\to I_0$ według
prawdopodobieństwa. Ostatnia postać ilorazu wiarygodności
zbiega więc do kwadratu zmiennej $\Normal(0,1)$,
czyli do $\chi_1^2$.
\end{proof}

\begin{przyklad}[Wald, score i LR: kiedy wyniki się rozchodzą?]
W regularnym modelu jednoparametrowym test Waldowski porównuje
$|\widehat\theta-\theta_0|$ z oszacowanym błędem standardowym,
test score używa $\ell_n'(\theta_0)$ i informacji przy hipotezie
zerowej, a test ilorazu wiarygodności używa
$2[\ell_n(\widehat\theta)-\ell_n(\theta_0)]$.
Asymptotycznie wszystkie trzy mają typową kalibrację normalną lub
$\chi^2_1$ pod $H_0$, lecz w małych próbach mogą dawać różne
wartości. Wald jest szczególnie wrażliwy na parametryzację i
przybliżenia w pobliżu brzegu. Testy dokładne i permutacyjne
omówimy w następnych częściach tomu.
\end{przyklad}

\section{Kiedy asymptotyka zawodzi? Zmienny nośnik i rzadkie zdarzenia}
\begin{twierdzenie}[Niegaussowska granica maksimum w modelu jednostajnym]
Niech $X_1,\ldots,X_n\iid U(0,\theta)$, $\theta>0$, i
$M_n=\max_iX_i$. Wtedy
$\frac n\theta(\theta-M_n)\xrightarrow{d}\operatorname{Exp}(1)$.
\end{twierdzenie}
\begin{proof}
Dla $x\ge0$ i $n>x$ mamy
\[
\Prob_\theta\left(\frac n\theta(\theta-M_n)>x\right)
=\Prob_\theta(M_n<\theta(1-x/n))=(1-x/n)^n.
\]
Granica jest $e^{-x}$, funkcją przeżycia rozkładu
wykładniczego z intensywnością jeden; w punktach $x<0$
prawdopodobieństwo przeżycia jest równe $1$.
Scharakteryzowanie funkcji rozkładu przez granice
w punktach jej ciągłości kończy dowód.
\end{proof}

Wyraz $n^{-1}$, a nie $n^{-1/2}$, jest naturalną skalą błędu
maksimum. Dlatego mechaniczne zastosowanie informacji Fishera
lub przedziału Waldowskiego byłoby błędem. Innym wyzwaniem jest
mała liczba sukcesów w modelu Bernoulliego: mimo liczebności
$n=100$ obserwacja $S=0$ lub $S=1$ może powodować bardzo słabe
działanie symetrycznego przybliżenia normalnego.

\section{Przykład obliczeniowy i tabela doboru metod}
\begin{center}\small
\begin{tabular}{>{\raggedright\arraybackslash}p{0.25\textwidth}>{\raggedright\arraybackslash}p{0.31\textwidth}>{\raggedright\arraybackslash}p{0.31\textwidth}}
\toprule
\textbf{Sytuacja} & \textbf{Pierwsza metoda} & \textbf{Weryfikacja lub alternatywa}\\
\midrule
Duża próba iid, parametr wewnętrzny & MLE i informacja Fishera & Krzywizna i diagnostyka modelu, przedział profilowy\\
Brzeg przestrzeni parametrów & Dokładna analiza rozkładu lub jednostronny przedział & Nie polegać na symetrycznym Waldzie\\
Liczba zdarzeń i ekspozycja & MLE intensywności Poissona & Dla małych liczebności przedział dokładny\\
Model może być błędnie określony & M-estymator i macierz kanapkowa & Interpretacja parametru pseudo-prawdziwego\\
Zależne obserwacje & Model zależności, korekta struktury kowariancji & Nie używać bezrefleksyjnie wariancji iid\\
Parametr zmienia nośnik & Statystyki porządkowe, rachunek dokładny & Cramér--Rao i Wilks wymagają regularności\\
\bottomrule
\end{tabular}
\end{center}

\Needspace{16\baselineskip}
\noindent\textbf{Przykład numeryczny w Pythonie (tylko biblioteka standardowa).}\par
\par
\Needspace{13\baselineskip}
\smallskip
\begin{lstlisting}[language=Python,caption={Porównanie przedziałów Walda i Wilsona dla odsetka wad.}]
from math import sqrt
n, k, z = 100, 7, 1.96
p = k / n
wald = (p - z * sqrt(p * (1-p)/n),
        p + z * sqrt(p * (1-p)/n))
den = 1 + z*z/n
center = (p + z*z/(2*n)) / den
half = z * sqrt(p*(1-p)/n + z*z/(4*n*n)) / den
wilson = (center-half, center+half)
print("Wald:", wald)
print("Wilson:", wilson)
\end{lstlisting}

\Needspace{16\baselineskip}
\noindent\textbf{Ten sam rachunek w Julii.}\par
\par
\Needspace{13\baselineskip}
\smallskip
\begin{lstlisting}[language=Julia,caption={Obliczanie przedziałów Walda i Wilsona w Julii.}]
n, k, z = 100, 7, 1.96
p = k / n
wald = (p-z*sqrt(p*(1-p)/n), p+z*sqrt(p*(1-p)/n))
den = 1 + z^2/n
center = (p + z^2/(2n)) / den
half = z*sqrt(p*(1-p)/n + z^2/(4n^2)) / den
println((wald=wald, wilson=(center-half, center+half)))
\end{lstlisting}

\paragraph{Jak używać tych wyników?}
Po pierwsze, podajemy założenie modelu niezależnych prób o wspólnym
$p$. Po drugie, raportujemy liczbę obserwacji i licznik sukcesów.
Po trzecie, sprawdzamy zachowanie przedziału na brzegu i zgodność
jego interpretacji z celem eksperymentu. Po czwarte, odróżniamy
\emph{niepewność parametru przy ustalonym modelu} od błędu
selekcji i błędów pomiaru, których żaden przedział wyznaczony z
błędnego modelu automatycznie nie naprawi.

% ================================================================
% ROZDZIAL 9 -- TESTOWANIE HIPOTEZ
% ================================================================
\chapter{Testowanie hipotez: konstrukcja, moc i interpretacja}
\label{ch:testing}
\paragraph{Powiązania i odsyłacze.} Teoria decyzji pojawia się w rozdziale~\ref{ch:decision}; asymptotyczne narzędzia testowania — w rozdziale~\ref{ch:asymptotic-estimation}.


\section{Od pytania praktycznego do testu}
Zakład wykazuje w ostatnim kwartale $7$ wad na $100$ losowo wybranych
produktów. Standard jakości wymaga, by prawdopodobieństwo wady nie
przekraczało $0{,}1$. Czy zaobserwowany odsetek $0{,}07$ świadczy o
poprawie? Trzeba odróżnić trzy pytania: (i) czy prawdopodobieństwo
wady jest mniejsze od $0{,}1$; (ii) czy jest ono \emph{znacząco}
mniejsze w sensie statystycznym; (iii) czy różnica jest wystarczająco
duża, by uzasadnić zmianę procesu produkcyjnego. Trzecie pytanie
wymaga także kryterium decyzji i kosztów, omówionych w rozdziale~\ref{ch:decision}.

Test jest regułą zależną od próby, a nie prawdopodobieństwem
prawdziwości hipotezy. Aby test miał sens, musimy zadeklarować
model, kierunek alternatywy, poziom istotności i sposób zebrania danych
\emph{przed} oglądaniem wyniku, jeśli chcemy zachować deklarowaną
kontrolę błędów.

\begin{definicja}[Hipotezy i test randomizowany]\label{def:statistical-test}
Niech $\mathcal E=(\mathsf X,\mathcal A,\{P_\theta:\theta\in\Theta\})$
będzie eksperymentem statystycznym, a $\Theta_0,\Theta_1$
niepustymi, rozłącznymi podzbiorami $\Theta$.
Hipotezę $H_0:\theta\in\Theta_0$ nazywamy zerową, a
$H_1:\theta\in\Theta_1$ alternatywną. \emph{Testem
randomizowanym} jest mierzalna funkcja
$\varphi:\mathsf X\to[0,1]$: po obserwacji $x$ hipotezę $H_0$
odrzucamy z prawdopodobieństwem $\varphi(x)$, korzystając w razie
potrzeby z niezależnego losowania pomocniczego. Test jest
\emph{nierandomizowany}, jeśli przyjmuje tylko wartości $0$ i $1$.
Funkcja $\beta_\varphi(\theta):=\E_\theta\varphi(X)$ jest
\emph{funkcją mocy}, natomiast
$\sup_{\theta\in\Theta_0}\beta_\varphi(\theta)$ jest
\emph{rozmiarem testu}. Test ma \emph{poziom istotności} co najwyżej
$\alpha\in[0,1]$, jeżeli $\sup_{\theta\in\Theta_0}\beta_\varphi(\theta)
\le\alpha$.
\end{definicja}

Jeśli $\theta\in\Theta_0$, odrzucenie jest błędem I rodzaju; jego
prawdopodobieństwo jest kontrolowane przez poziom testu. Jeśli
$\theta\in\Theta_1$, nieodrzucenie jest błędem II rodzaju;
prawdopodobieństwo tego błędu wynosi $1-\beta_\varphi(\theta)$.
Sama wartość $\alpha$ nie mówi nic o mocy ani o rozmiarze efektu.
Nieodrzucenie $H_0$ nie oznacza jego zaakceptowania jako prawdziwego.

\begin{przyklad}[Trzy pytania, trzy hipotezy]
Dla $X_i\iid\operatorname{Bernoulli}(p)$ można rozważać
$H_0:p\ge 0{,}1$ przeciw $H_1:p<0{,}1$ (poprawa),
$H_0:p\le 0{,}1$ przeciw $H_1:p>0{,}1$ (pogorszenie),
lub $H_0:p=0{,}1$ przeciw $H_1:p\ne0{,}1$ (dowolna zmiana).
To są różne problemy; wynik dwustronnego testu nie zastępuje
jednostronnego testu zaplanowanego ze względu na ryzyko wad.
\end{przyklad}

\section{Lemat Neymana--Pearsona: skąd bierze się najlepsza reguła?}
\begin{twierdzenie}[Lemat Neymana--Pearsona, z randomizacją]
\label{thm:neyman-pearson}
Niech $P_0,P_1$ będą dwoma miarami probabilistycznymi na
$(\mathsf X,\mathcal A)$ i niech $P_1\ll P_0$. Oznaczmy
$L=\dd P_1/\dd P_0$ ich iloraz wiarygodności. Niech
$0<\alpha<1$. Istnieją $c\ge0$ i $\gamma\in[0,1]$ takie, że
\[
 \varphi_*(x)=\ind{\{L(x)>c\}}+
 \gamma\ind{\{L(x)=c\}},\qquad \E_0\varphi_*=\alpha.
\]
Wtedy dla każdego mierzalnego testu $\varphi:\mathsf X\to[0,1]$
spełniającego $\E_0\varphi\le\alpha$ zachodzi
$\E_1\varphi\le\E_1\varphi_*$. A zatem $\varphi_*$ ma
największą moc wśród testów poziomu co najwyżej $\alpha$ dla
hipotez prostych $H_0:P=P_0$ i $H_1:P=P_1$.
\end{twierdzenie}
\begin{proof}
Ponieważ $L\ge0$ i $\E_0L=1$, istnieje próg $c\ge0$ spełniający
$P_0(L>c)\le\alpha\le P_0(L\ge c)$. Można go dobrać
jako odpowiedni kwantyl rozkładu $L$ pod $P_0$; gdy $c=0$,
warunek ma sens, bo $P_0(L\ge0)=1$. Jeśli
$P_0(L=c)>0$, ustalamy
$\gamma=[\alpha-P_0(L>c)]/P_0(L=c)$, a w przeciwnym razie
$\gamma=0$. Nierówności zapewniają $\gamma\in[0,1]$ i
$\E_0\varphi_*=\alpha$.

W punktach, w których $L>c$, mamy $\varphi_*=1$, zatem
$(L-c)(\varphi_*-\varphi)\ge0$. Przy $L<c$ mamy
$\varphi_*=0$, co znów daje nieujemny iloczyn; na zbiorze
$L=c$ iloczyn znika. W rezultacie
\[
 \E_1(\varphi_*-\varphi)
 =\E_0[L(\varphi_*-\varphi)]
 \ge c\E_0(\varphi_*-\varphi)\ge0.
\]
Wykorzystaliśmy absolutną ciągłość, by zapisać wartość
oczekiwaną pod $P_1$ jako całkę względem $P_0$, oraz kontrolę
poziomu testu, a nie szczególną postać obserwacji.
\end{proof}

\begin{uwaga}[Dlaczego randomizacja?]
Dla dyskretnych rozkładów nie zawsze można znaleźć zbiór
odrzucenia o prawdopodobieństwie dokładnie $\alpha$ pod $H_0$.
Test bez randomizacji zwykle ma \emph{rzeczywisty} rozmiar
mniejszy niż deklarowany poziom; jest to poprawny test
konserwatywny. Randomizacja na granicy pozwala wykorzystać
cały dopuszczony budżet błędu I rodzaju, lecz w typowych analizach
praktycznych częściej wybiera się prostszą procedurę
nierandomizowaną.
\end{uwaga}

\section{Testy jednostajnie najmocniejsze i statystyki monotoniczne}
\begin{definicja}[Test jednostajnie najmocniejszy]
Test $\varphi_*$ o poziomie co najwyżej $\alpha$ jest
\emph{jednostajnie najmocniejszy} (UMP) dla
$H_0:\theta\in\Theta_0$ przeciw $H_1:\theta\in\Theta_1$,
jeśli dla każdego $\theta\in\Theta_1$ i każdego testu
$\varphi$ tego samego poziomu mamy
$\E_\theta\varphi\le\E_\theta\varphi_*$.
\end{definicja}

\begin{twierdzenie}[Jednostronny test dwumianowy jest UMP]
\label{thm:binomial-ump}
Niech $X_1,\ldots,X_n\iid\operatorname{Bernoulli}(p)$,
$0<p<1$, $S=\sum_iX_i$, a $p_0\in(0,1)$.
Dla $0<\alpha<1$ wybierzmy $k\in\{0,\ldots,n\}$ oraz
$\gamma\in[0,1]$ tak, żeby funkcja
$\varphi_*(X)=\ind{\{S>k\}}+\gamma\ind{\{S=k\}}$
spełniała $\E_{p_0}\varphi_*=\alpha$.
Wtedy $\varphi_*$ jest testem UMP poziomu $\alpha$ dla
$H_0:p\le p_0$ przeciw $H_1:p>p_0$.
\end{twierdzenie}
\begin{proof}
Dla dowolnego $p_1>p_0$ iloraz wiarygodności próby wynosi
\[
 L(x)=\left(\frac{p_1}{p_0}\right)^{S(x)}
 \left(\frac{1-p_1}{1-p_0}\right)^{n-S(x)}.
\]
Stosunek $L(S+1)/L(S)=p_1(1-p_0)/[p_0(1-p_1)]>1$,
więc test Neymana--Pearsona przeciw $p_1$ odrzuca dla
dużych wartości $S$ i ma graniczną randomizację.
Ponieważ rozkład $S$ przy $p_0$ nie zależy od $p_1$,
można wybrać ten sam próg $k$ i $\gamma$ dla wszystkich
$p_1>p_0$.

Trzeba jeszcze sprawdzić poziom względem całego $H_0$.
Niech $U_1,\ldots,U_n$ będą niezależne z rozkładu
$U(0,1)$ i zdefiniujmy $S_p=\sum_i\ind{\{U_i\le p\}}$.
Dla $p\le p_0$ zachodzi $S_p\le S_{p_0}$ prawie na pewno.
Ponieważ funkcja $s\mapsto\ind{\{s>k\}}+
\gamma\ind{\{s=k\}}$ jest niemalejąca, moc przy $p$
nie przewyższa mocy przy $p_0$, czyli $\alpha$.
Lemat Neymana--Pearsona zapewnia następnie największą moc
przy każdym $p_1>p_0$ wśród testów poziomu $\alpha$.
\end{proof}

\begin{uwaga}[Test dwustronny]
W modelu normalnym o znanej wariancji test obustronny
$H_0:\mu=\mu_0$ przeciw $H_1:\mu\ne\mu_0$
nie ma na ogół testu UMP: najbardziej mocna reguła przeciw
$\mu_1>\mu_0$ odrzuca dla dużej średniej, natomiast przeciw
$\mu_1<\mu_0$ dla małej średniej. Wspólny test o maksymalnej
mocy dla obu kierunków nie może jednocześnie być obydwoma
jednostronnymi testami Neymana--Pearsona. Potrzebne są dodatkowe
kryteria, np. nieobciążoność testu lub niezmienniczość.
\end{uwaga}

\section{Wartość p: definicja i częste błędy interpretacji}
\begin{definicja}[Wartość p dla statystyki testowej]\label{def:pvalue}
Niech $T:\mathsf X\to\R$ będzie mierzalną statystyką,
a większe jej wartości świadczą przeciw $H_0$.
Dla obserwacji $x$ definiujemy
\[
 p(x):=\sup_{\theta\in\Theta_0}
 P_\theta\bigl(T(X)\ge T(x)\bigr).
\]
Zakładamy, że funkcja $x\mapsto p(x)$ jest mierzalna.
Dla hipotezy prostej supremum jest niepotrzebne.
\end{definicja}

\begin{twierdzenie}[Kontrola błędu przez wartość p]
\label{thm:pvalue-valid}
Jeśli $p$ jest określona jak w definicji~\ref{def:pvalue},
to dla każdego $u\in[0,1]$ oraz każdego $\theta\in\Theta_0$
zachodzi $P_\theta(p(X)\le u)\le u$.
\end{twierdzenie}
\begin{proof}
Ustalmy $\theta\in\Theta_0$ i oznaczmy
$G(t)=P_\theta(T(X)\ge t)$.
Z definicji $p(X)\ge G(T(X))$ prawie na pewno, a więc
$\{p(X)\le u\}\subseteq\{G(T(X))\le u\}$.
Dla dowolnej rzeczywistej zmiennej losowej $Y$ funkcja
$G_Y(t)=P(Y\ge t)$ jest nierosnąca. Zbiór
$A_u=\{t\in\R:G_Y(t)\le u\}$ jest zatem pusty,
cały lub jest półprostą postaci $[a,\infty)$ bądź
$(a,\infty)$. W pierwszym przypadku
$P(Y\in A_u)=0$. W drugim, o ile $A_u=[a,\infty)$,
mamy $P(Y\in A_u)=G_Y(a)\le u$.
Jeśli $A_u=(a,\infty)$, to z ciągłości prawdopodobieństwa
przy granicach monotonicznych
$P(Y>a)=\lim_{m\to\infty}P(Y\ge a+1/m)\le u$.
Dla $A_u=\R$ zachodzi $u=1$.
Zatem zawsze $P(G_Y(Y)\le u)\le u$. Stosując to do
$Y=T(X)$, otrzymujemy
$P_\theta(G(T(X))\le u)\le u$, a więc tezę.
\end{proof}

Dla ciągłej statystyki o ciągłym rozkładzie pod prostym
$H_0$ wartość p jest jednostajna na $[0,1]$; dla statystyk
dyskretnych najczęściej jest \emph{nadjednostajna} w sensie
powyższej nierówności, zatem dokładna nominalna równość nie
musi zachodzić. W żadnym przypadku $p=0{,}03$ nie oznacza
$3\%$ prawdopodobieństwa, że hipoteza zerowa jest prawdziwa.

\begin{uwaga}[Jedno- i dwustronne wartości p nie są tym samym]
Definicja~\ref{def:pvalue} jest jednostronnym ogonem statystyki,
której duże wartości przemawiają przeciw hipotezie zerowej.
W symetrycznym ciągłym modelu normalnym dla
$H_0:\mu=\mu_0$ i statystyki $Z$ o rozkładzie
$\Normal(0,1)$ dwustronna wartość p ma postać
$P(|Z'|\ge|z_{\rm obs}|)=2[1-\Phi(|z_{\rm obs}|)]$,
gdzie $Z'\sim\Normal(0,1)$. W modelu dyskretnym,
np. dwumianowym, istnieje kilka konwencji dotyczących
„tak samo ekstremalnych” wyników dwustronnych, które
nie muszą dawać identycznych wartości p. Należy więc
zawsze nazwać zastosowaną regułę odrzucenia, a nie tylko
podać słowo „dwustronny”.
\end{uwaga}

\begin{zadanie}[7 wad na 100 produktów: test dokładny]\label{zad:binom-test}
W modelu $X_i\iid\operatorname{Bernoulli}(p)$
zaobserwowano $S=7$ przy $n=100$. Testujemy
$H_0:p\ge0{,}1$ przeciw $H_1:p<0{,}1$.
Uzasadnij postać wartości p i oblicz ją.
\par\emph{Rozwiązanie.} Mniejsza liczba wad przemawia za
alternatywą, zatem z monotoniczności rozkładu dwumianowego
w parametrze $p$ mamy
\[
 p\text{-value}=\sup_{p\ge0{,}1}P_p(S\le7)
 =P_{0{,}1}(S\le7)
 =\sum_{j=0}^7\binom{100}{j}(0{,}1)^j(0{,}9)^{100-j}
 \approx0{,}2061.
\]
Na poziomie $5\%$ nie odrzucamy $H_0$. Nie oznacza to,
że udowodniliśmy $p\ge0{,}1$: badanie może mieć
niewystarczającą moc dla niewielkiej poprawy.
\end{zadanie}

\section{Moc, wielkość efektu i dobór liczebności próby}
\begin{twierdzenie}[Dokładna moc jednostronnego testu normalnego]
\label{thm:normal-power}
Niech $X_1,\ldots,X_n\iid\Normal(\mu,\sigma^2)$,
$\sigma>0$ będzie znana, a $\mu_0\in\R$.
Test odrzucający $H_0:\mu\le\mu_0$ na rzecz
$H_1:\mu>\mu_0$ wtedy i tylko wtedy, gdy
$\bar X>\mu_0+z_{1-\alpha}\sigma/\sqrt n$, ma
poziom $\alpha$ i funkcję mocy
\[
 \beta(\mu)=1-\Phi\left(z_{1-\alpha}
 -\frac{\sqrt n(\mu-\mu_0)}{\sigma}\right),
\]
gdzie $\Phi$ jest dystrybuantą rozkładu $\Normal(0,1)$.
\end{twierdzenie}
\begin{proof}
Z dokładnego rozkładu $\bar X\sim\Normal(\mu,\sigma^2/n)$
po standaryzacji otrzymujemy podaną funkcję mocy.
Jest ona rosnąca w $\mu$, a więc jej największa wartość na
$\{\mu\le\mu_0\}$ występuje w $\mu_0$ i wynosi $\alpha$.
\end{proof}

\begin{wniosek}[Dobór liczebności dla zadanego efektu]
\label{cor:normal-sample-size}
Przy założeniach twierdzenia~\ref{thm:normal-power}, dla
$0<\alpha<1/2$, $0<\beta_0<1/2$ i efektu $\delta>0$
warunek mocy co najmniej $1-\beta_0$ przy $\mu=\mu_0+\delta$
jest równoważny
\[
 n\ge\left(\frac{\sigma(z_{1-\alpha}+z_{1-\beta_0})}
 {\delta}\right)^2.
\]
Najmniejsze dopuszczalne $n$ uzyskujemy przez
zaokrąglenie prawej strony w górę (z zastrzeżeniem $n\ge1$).
\end{wniosek}
\begin{proof}
Nierówność $1-\Phi(z_{1-\alpha}-\sqrt n\delta/\sigma)
\ge1-\beta_0$ oznacza
$z_{1-\alpha}-\sqrt n\delta/\sigma\le z_{\beta_0}
=-z_{1-\beta_0}$. Po przestawieniu otrzymujemy tezę.
\end{proof}

\begin{zadanie}[Ile pomiarów potrzeba?]
Załóżmy znane odchylenie standardowe $\sigma=10$ jednostek.
Chcemy z poziomem istotności $\alpha=0{,}05$ wykryć wzrost
średniej o $\delta=3$ z mocą co najmniej $0{,}80$.
\par\emph{Rozwiązanie.} Dla $z_{0{,}95}\approx1{,}645$
oraz $z_{0{,}80}\approx0{,}842$ wystarcza
$n\ge[10(1{,}645+0{,}842)/3]^2\approx68{,}73$,
czyli $n=69$. To rachunek dla \emph{znanej} $\sigma$,
niezależnych obserwacji i testu jednostronnego. Przy danych
zależnych, większym rzeczywistym rozproszeniu lub kilku
porównaniach potrzebny jest inny projekt.
\end{zadanie}

Nie należy określać $\delta$ jako „istotnego efektu” wyłącznie
dlatego, że dla dostatecznie dużego $n$ daje małą wartość p.
Trzeba wcześniej ustalić \emph{minimalny efekt praktycznie ważny},
czyli najmniejszą zmianę, która miałaby znaczenie dla decyzji.
Raportowanie przedziału ufności i miary wielkości efektu jest
zwykle znacznie bardziej informacyjne niż samo odrzucenie hipotezy.

\section{Pułapki praktyczne: selekcja, wielokrotność i zatrzymywanie}
\begin{przyklad}[Zatrzymywanie, gdy pojawi się mała wartość p]
Załóżmy, że pięć niezależnych badań prowadzonych pod prawdziwą
hipotezą zerową daje niezależne wartości p o rozkładzie
jednostajnym na $[0,1]$. Jeśli ogłosimy sukces, gdy
\emph{którakolwiek} wartość będzie mniejsza od $0{,}05$,
prawdopodobieństwo przynajmniej jednego fałszywego alarmu wynosi
$1-(1-0{,}05)^5\approx0{,}2262$, a nie $0{,}05$.
Ten sam problem pojawia się przy wielokrotnym testowaniu albo
ciągłym podglądaniu danych i zakończeniu eksperymentu w
pierwszej korzystnej chwili. Zależności między testami mogą
zmienić tę liczbę, ale nie usprawiedliwiają ignorowania problemu.
\end{przyklad}

\begin{center}\small
\begin{tabular}{>{\raggedright\arraybackslash}p{0.30\textwidth}>{\raggedright\arraybackslash}p{0.57\textwidth}}
\toprule
\textbf{Pytanie przed testem} & \textbf{Dlaczego jest ważne?}\\
\midrule
Czy hipotezę i kierunek wybrano wcześniej? & Wybranie kierunku po zobaczeniu danych może naruszać kontrolę poziomu.\\
Czy obserwacje są niezależne? & Korelacja wewnątrz grup zwykle zmienia rozkład statystyki i moc.\\
Czy testujemy różnicę, przewagę czy równoważność? & Brak istotnej różnicy nie jest dowodem równoważności.\\
Czy rozkład statystyki jest dokładny czy asymptotyczny? & Małe próby i rzadkie zdarzenia mogą wymagać testu dokładnego.\\
Ile hipotez sprawdzamy? & Wielokrotność wymaga jawnej strategii kontroli błędów.\\
Czy efekt jest użyteczny? & Wartość p nie mierzy wielkości ani praktycznej ważności efektu.\\
\bottomrule
\end{tabular}
\end{center}

\begin{lstlisting}[language=Python,caption={Dokładny test dwumianowy i obliczenie mocy testu normalnego.}]
from math import ceil, sqrt
from scipy.stats import binomtest, norm

k, n, p0 = 7, 100, 0.10
result = binomtest(k, n, p0, alternative="less")
print("p-value:", result.pvalue)

alpha, power, sigma, delta = 0.05, 0.80, 10.0, 3.0
z_alpha = norm.ppf(1 - alpha)
z_power = norm.ppf(power)
n_required = ceil((sigma * (z_alpha + z_power) / delta)**2)
print("minimalne n:", n_required)
\end{lstlisting}

\paragraph{Co sprawdzamy przed uruchomieniem programu?}
Pakiet \texttt{scipy} należy zainstalować w środowisku Pythona;
kod zakłada niezależne próby Bernoulliego oraz znaną wariancję
w planowaniu testu normalnego. Żadne wywołanie funkcji testowej
nie sprawdza automatycznie reprezentatywności próby ani
prawidłowości projektu badania.


% ================================================================
% ROZDZIAL 10 -- TESTY PARAMETRYCZNE I WPROWADZENIE DO ANOVA
% ================================================================
\chapter{Testy parametryczne: średnie, wariancje, częstości i ANOVA}
\label{ch:parametric-tests}
\paragraph{Powiązania i odsyłacze.} Testy opierają się na rozkładach statystyk z rozdziału~\ref{ch:statystyki} i konstrukcji testów z rozdziału~\ref{ch:testing}.


\section{Jedna średnia: kiedy test normalny, a kiedy t-Studenta?}
Rozpocznijmy od pomiarów czasu pracy urządzenia. Interesuje nas,
czy średni czas przekracza wartość referencyjną $\mu_0$.
Jeśli znamy odchylenie standardowe modelu normalnego, statystyka
$Z=\sqrt n(\bar X-\mu_0)/\sigma$ ma dokładnie rozkład
$\Normal(0,1)$ pod hipotezą $H_0:\mu=\mu_0$.
Najczęściej $\sigma$ nie jest jednak znane i trzeba je oszacować
z tych samych danych. Wtedy pojawia się rozkład t-Studenta.

\begin{twierdzenie}[Dokładny test t dla jednej średniej]
\label{thm:one-sample-t}
Niech $n\ge2$ i $X_1,\ldots,X_n\iid\Normal(\mu,\sigma^2)$,
gdzie $\mu\in\R$, $\sigma>0$. Niech
$\bar X=n^{-1}\sum_iX_i$ i
$S^2=(n-1)^{-1}\sum_i(X_i-\bar X)^2$.
Pod $H_0:\mu=\mu_0$ statystyka
\[
 T=\frac{\bar X-\mu_0}{S/\sqrt n}
\]
ma rozkład t-Studenta z $n-1$ stopniami swobody, z wyjątkiem
zdarzenia $S=0$, które ma prawdopodobieństwo zero.
Tutaj $t_{\nu,q}$ oznacza kwantyl rzędu $q\in(0,1)$
rozkładu t-Studenta z $\nu>0$ stopniami swobody.
Dla $0<\alpha<1/2$ test dwustronny odrzucający dla
$|T|>t_{n-1,1-\alpha/2}$ ma poziom dokładnie $\alpha$,
a jednostronny dla $T>t_{n-1,1-\alpha}$ również ma
poziom dokładnie $\alpha$ przy $H_0:\mu\le\mu_0$.
\end{twierdzenie}
\begin{proof}
Z rozkładu próby normalnej wyprowadzonego w
rozdziale~\ref{ch:statystyki} mamy
$Z=\sqrt n(\bar X-\mu)/\sigma\sim\Normal(0,1)$ oraz
$V=(n-1)S^2/\sigma^2\sim\chi^2_{n-1}$, przy czym
$Z$ i $V$ są niezależne. Pod $H_0$ statystyka $T$ jest
równa $Z/\sqrt{V/(n-1)}$, co definiuje rozkład t-Studenta.
Dla alternatywy jednostronnej przy $\mu\le\mu_0$
rozkład $T$ jest przesunięty przez niedodatni składnik
$\sqrt n(\mu-\mu_0)/\sigma$ w liczniku, przy zachowanej
niezależności mianownika. Stąd prawdopodobieństwo przekroczenia
ustalonego dodatniego progu jest największe przy $\mu=\mu_0$.
Symetria i ciągłość rozkładu t zapewniają dokładny poziom
dla obu reguł.
\end{proof}

\begin{uwaga}[Normalność czego?]
Test t dla jednej średniej ma dokładny rozkład pod
założeniem, że \emph{obserwacje} są iid normalne. W teście t dla
różnic sparowanych normalności wymagamy od \emph{różnic}, nie
koniecznie od każdej serii oddzielnie. W regresji i ANOVA
normalność dotyczy rozkładu błędów warunkowych, a nie połączonej
próby ze wszystkich grup. Przy dużych próbach przybliżenia
na podstawie CLT mogą działać bez normalności danych, ale
zależą od ogonów rozkładu, skończoności wariancji i schematu
próby.
\end{uwaga}

\begin{stwierdzenie}[Przedział dokładny dla średniej normalnej]
Pod założeniami twierdzenia~\ref{thm:one-sample-t}
przedział losowy
\[
 \left[\bar X-t_{n-1,1-\alpha/2}\frac S{\sqrt n},\;
 \bar X+t_{n-1,1-\alpha/2}\frac S{\sqrt n}\right]
\]
obejmuje $\mu$ z prawdopodobieństwem dokładnie $1-\alpha$.
\end{stwierdzenie}
\begin{proof}
To zdarzenie jest równoważne
$|\sqrt n(\bar X-\mu)/S|\le t_{n-1,1-\alpha/2}$.
Twierdzenie~\ref{thm:one-sample-t} dla prawdziwej średniej
zapewnia prawdopodobieństwo $1-\alpha$.
\end{proof}

\section{Pomiary powtarzane: różnice par zamiast dwóch prób niezależnych}
\begin{twierdzenie}[Test t dla par]\label{thm:paired-t}
Niech pary $(X_i,Y_i)$, $i=1,\ldots,n$, będą niezależne
między jednostkami, a różnice $D_i=X_i-Y_i$ mają wspólny
rozkład $\Normal(\delta,\tau^2)$ z $\tau>0$ i $n\ge2$.
Wtedy przy $H_0:\delta=0$ statystyka
$T_D=\sqrt n\bar D/S_D$ ma dokładnie rozkład $t_{n-1}$.
Nie wymagamy niezależności $X_i$ i $Y_i$ wewnątrz pary.
\end{twierdzenie}
\begin{proof}
Zmienne $D_i$ tworzą próbę prostą normalną, więc jest to
bezpośrednie zastosowanie twierdzenia~\ref{thm:one-sample-t}
do obserwacji $D_1,\ldots,D_n$.
\end{proof}

\begin{zadanie}[Przed i po zmianie procedury]
Dla ośmiu osób wyniki przed zmianą wynoszą
$54,58,57,60,56,55,61,59$, a po zmianie
$51,56,53,57,55,53,56,57$. Czy średnia różnica jest dodatnia?
\par\emph{Rozwiązanie.} Różnice to
$3,2,4,3,1,2,5,2$, stąd
$\bar d=2{,}75$, $s_D^2=23/14\approx1{,}6429$,
$T_D=\sqrt8\,\bar d/s_D\approx6{,}068$,
przy siedmiu stopniach swobody. W teście dwustronnym
wartość p jest około $0{,}00051$; w zaplanowanym teście
jednostronnym na spadek wyniku po interwencji jest około
$0{,}00025$. To jednak nie dowodzi przyczynowości, jeżeli
dane nie pochodzą z odpowiednio zaprojektowanego eksperymentu.
\end{zadanie}

\section{Dwie niezależne grupy: Student czy Welch?}
\begin{twierdzenie}[Dwie średnie przy wspólnej wariancji]
\label{thm:two-sample-t}
Niech $X_1,\ldots,X_n\iid\Normal(\mu_X,\sigma^2)$ i
$Y_1,\ldots,Y_m\iid\Normal(\mu_Y,\sigma^2)$ będą niezależnymi
próbami, gdzie $n,m\ge2$, a $\sigma>0$ jest nieznane.
Niech
$S_p^2=[(n-1)S_X^2+(m-1)S_Y^2]/(n+m-2)$.
Pod $H_0:\mu_X-\mu_Y=\delta_0$ zachodzi
\[
 T_p=\frac{(\bar X-\bar Y)-\delta_0}
 {S_p\sqrt{1/n+1/m}}\sim t_{n+m-2}.
\]
\end{twierdzenie}
\begin{proof}
Ponieważ próby są niezależne,
$Z=[(\bar X-\bar Y)-(\mu_X-\mu_Y)]/
[\sigma\sqrt{1/n+1/m}]$ ma rozkład $\Normal(0,1)$.
Dla każdej próby z osobna resztowa suma kwadratów podzielona
przez $\sigma^2$ jest zmienną $\chi^2$ o odpowiednio
$n-1$ i $m-1$ stopniach swobody; obie są niezależne
między sobą i od obu średnich. Ich suma
$(n+m-2)S_p^2/\sigma^2$ ma rozkład
$\chi^2_{n+m-2}$. Iloraz $Z$ i pierwiastka ze
średniej tej sumy daje dokładny rozkład t.
\end{proof}

Gdy wariancje są różne, podstawienie wariancji wspólnej
może istotnie zniekształcić test, szczególnie jeśli
liczebności grup też się różnią. Wtedy naturalną alternatywą
jest \emph{test Welcha}:
\[
 T_W=\frac{(\bar X-\bar Y)-\delta_0}
 {\sqrt{S_X^2/n+S_Y^2/m}},\qquad
 \nu=\frac{(S_X^2/n+S_Y^2/m)^2}
 {(S_X^2/n)^2/(n-1)+(S_Y^2/m)^2/(m-1)}.
\]
Rozkład $T_W$ jest zwykle \emph{przybliżany} rozkładem
$t_\nu$, gdzie $\nu$ to losowe, niekoniecznie całkowite
przybliżenie liczby stopni swobody Satterthwaite'a.
Nie jest to twierdzenie o dokładnym rozkładzie t przy
nierównych wariancjach, nawet jeśli obie próby są normalne.

\begin{zadanie}[Dwie technologie produkcji]\label{zad:welch-case}
Czasy obróbki w technologii A to
$18,19,20,19,21,20,18,22$; w technologii B:
$22,23,21,24,22,25,23,24$. Porównaj średnie.
\par\emph{Rozwiązanie.} Mamy $n=m=8$,
$\bar x=19{,}625$, $\bar y=23$, $s_X^2\approx1{,}9821$,
$s_Y^2\approx1{,}7143$. Statystyka Welcha wynosi
$T_W\approx-4{,}9651$, $\nu\approx13{,}927$,
a dwustronna wartość p około $0{,}00021$.
W tym przykładzie test ze wspólną wariancją daje bardzo
zbliżony wynik, ale to \emph{nie dowodzi}, że w dowolnym
niezrównoważonym badaniu obie procedury są równoważne.
Do raportu należy dołączyć różnicę średnich
$\bar x-\bar y=-3{,}375$ i jej przedział ufności.
\end{zadanie}

\begin{uwaga}[Nie wybieramy testu wyłącznie przez pretest wariancji]
Procedura „najpierw test równości wariancji, potem Student lub
Welch zależnie od p-value” zmienia rozkład końcowej reguły i
nie jest automatycznie testem o nominalnym poziomie.
Gdy nie mamy silnego powodu, by zakładać równość wariancji,
Welch jest zwykle rozsądnym punktem wyjścia dla dwóch
niezależnych grup. Z kolei silna asymetria, ciężkie ogony,
zależność obserwacji czy mała próba wymagają dodatkowej oceny:
sam Welch nie usuwa wszystkich tych problemów.
\end{uwaga}

\section{Testy wariancji i kiedy nie należy ich używać}
\begin{twierdzenie}[Test chi-kwadrat jednej wariancji]
Niech $X_i\iid\Normal(\mu,\sigma^2)$, $n\ge2$,
$\mu\in\R$ nieznana i $\sigma>0$. Niech $v_0>0$. Pod $H_0:\sigma^2=v_0$
statystyka $Q=(n-1)S^2/v_0$ ma rozkład
$\chi^2_{n-1}$. Dwustronny test odrzuca dla
$Q<\chi^2_{n-1,\alpha/2}$ lub
$Q>\chi^2_{n-1,1-\alpha/2}$, gdzie dolne i górne
kwantyle odpowiadają rozkładowi $\chi^2$.
\end{twierdzenie}
\begin{proof}
Wynik o rozkładzie sumy kwadratów scentrowanej próby
normalnej został udowodniony w rozdziale~\ref{ch:statystyki}.
Pod hipotezą zerową podstawiamy $\sigma^2=v_0$.
\end{proof}

\begin{twierdzenie}[Test F równości dwóch wariancji]
Niech niezależne próby normalne mają odpowiednio
liczebności $n,m\ge2$ i wariancje $\sigma_X^2,\sigma_Y^2>0$.
Pod $H_0:\sigma_X^2=\sigma_Y^2$ iloraz
$F=S_X^2/S_Y^2$ ma rozkład Fishera--Snedecora
$F_{n-1,m-1}$.
\end{twierdzenie}
\begin{proof}
Z niezależności prób i rozkładu wariancji normalnej
$U=(n-1)S_X^2/\sigma_X^2\sim\chi^2_{n-1}$ i
$V=(m-1)S_Y^2/\sigma_Y^2\sim\chi^2_{m-1}$ są
niezależne. Pod $H_0$ zachodzi
$F=[U/(n-1)]/[V/(m-1)]$, co z definicji daje
rozkład Fishera--Snedecora.
\end{proof}

Rozkład dokładny jest tutaj konsekwencją normalności.
W praktyce test F jest bardzo wrażliwy na jej naruszenie.
Jeśli zadaniem jest sprawdzenie założenia jednorodności
rozproszenia w wielu grupach, często używa się procedur
Levene'a, Browna--Forsythe'a lub Flignera--Killeena.
Ich konstrukcje, ograniczenia i porównania z testami
normalności (w tym Shapiro--Wilka) omówimy osobno w części
o diagnostyce modeli.

\section{Jednoczynnikowa ANOVA: skąd pochodzi statystyka F?}
\label{sec:anova-introduction}
\begin{definicja}[Model jednoczynnikowej ANOVA]
Niech $k\ge2$ i $n_j\ge2$ dla $j=1,\ldots,k$.
Obserwacje $Y_{ji}$ są niezależne, a
\[
 Y_{ji}=\mu_j+\varepsilon_{ji},\qquad
 \varepsilon_{ji}\iid\Normal(0,\sigma^2),\qquad\sigma>0.
\]
Hipoteza zerowa ma postać
$H_0:\mu_1=\cdots=\mu_k$.
Niech $N=\sum_jn_j$,
$\bar Y_j=n_j^{-1}\sum_{i=1}^{n_j}Y_{ji}$ i
$\bar Y=N^{-1}\sum_jn_j\bar Y_j$.
Definiujemy sumy kwadratów
$SS_B=\sum_jn_j(\bar Y_j-\bar Y)^2$ (między grupami)
oraz $SS_W=\sum_j\sum_i(Y_{ji}-\bar Y_j)^2$
(wewnątrz grup).
\end{definicja}

\begin{lemat}[Ortogonalny rozkład sum kwadratów]
\label{lem:anova-decomp}
Dla dowolnych liczb rzeczywistych $y_{ji}$ przy
zdefiniowanych średnich zachodzi
\[
 \sum_{j=1}^k\sum_{i=1}^{n_j}(y_{ji}-\bar y)^2
 =\sum_j\sum_i(y_{ji}-\bar y_j)^2
 +\sum_jn_j(\bar y_j-\bar y)^2.
\]
\end{lemat}
\begin{proof}
Dla obserwacji z grupy $j$ zapisujemy
$y_{ji}-\bar y=(y_{ji}-\bar y_j)+(\bar y_j-\bar y)$.
Po podniesieniu do kwadratu i zsumowaniu wyraz mieszany
znika, gdyż $\sum_i(y_{ji}-\bar y_j)=0$.
Pozostałe dwa składniki są dokładnie $SS_W$ i $SS_B$.
\end{proof}

\begin{twierdzenie}[Dokładny rozkład statystyki ANOVA]
\label{thm:anova-f}
W modelu jednoczynnikowej ANOVA pod $H_0$ zmienne
$SS_B/\sigma^2$ i $SS_W/\sigma^2$ są niezależne,
a ich rozkłady wynoszą odpowiednio $\chi^2_{k-1}$
oraz $\chi^2_{N-k}$. W konsekwencji
\[
 F=\frac{SS_B/(k-1)}{SS_W/(N-k)}\sim F_{k-1,N-k}
\]
i test odrzucający przy dużych wartościach $F$ ma
poziom dokładnie $\alpha$, jeśli użyjemy kwantyla
$F_{k-1,N-k,1-\alpha}$.
\end{twierdzenie}
\begin{proof}
Zapiszmy wszystkie $N$ obserwacji jako wektor
$Y\in\R^N$ o postaci $Y=\mu\mathbf1+\varepsilon$
pod $H_0$, gdzie $\varepsilon\sim\Normal_N(0,\sigma^2I_N)$.
Niech $G\subset\R^N$ oznacza podprzestrzeń wektorów
stałych w obrębie każdej grupy. Ma ona wymiar $k$ i zawiera
wektor $\mathbf1$. Oznaczmy przez $P_G$ rzut
ortogonalny na $G$, a przez $P_0$ rzut na
$\operatorname{span}\{\mathbf1\}$. Wówczas
$P_G-P_0$ i $I_N-P_G$ są rzutami ortogonalnymi na
wzajemnie prostopadłe podprzestrzenie o wymiarach
odpowiednio $k-1$ i $N-k$.

Wektor $P_GY$ ma na pozycjach grupy $j$ wartość
$\bar Y_j$, podczas gdy $P_0Y=\bar Y\mathbf1$.
Zatem $SS_B=\|(P_G-P_0)Y\|^2$ i
$SS_W=\|(I_N-P_G)Y\|^2$.
Oba rzuty zerują $\mathbf1$, stąd pod $H_0$ możemy w tych
wyrażeniach zastąpić $Y$ przez $\varepsilon$.
Po wyborze ortonormalnych baz obu podprzestrzeni
współrzędne rzutów $\varepsilon/\sigma$ są niezależnymi
standardowymi zmiennymi normalnymi. Wynika to z
ortogonalnej niezmienniczości $\Normal_N(0,I_N)$:
współrzędne obu rzutów są funkcjami rozłącznych zbiorów
niezależnych współrzędnych. Sumy ich kwadratów mają
zatem wskazane rozkłady chi-kwadrat i są niezależne.
Iloraz niezależnych średnich takich sum ma z definicji
rozkład $F_{k-1,N-k}$.
\end{proof}

\begin{zadanie}[Trzy metody produkcji]\label{zad:anova-three}
Dla trzech niezależnych grup po osiem obserwacji mamy:
$A=(18,19,20,19,21,20,18,22)$,
$B=(22,23,21,24,22,25,23,24)$ i
$C=(25,27,26,24,28,26,27,25)$.
Przetestuj równość średnich i zinterpretuj wynik.
\par\emph{Rozwiązanie.} Średnie grup wynoszą odpowiednio
$19{,}625$, $23$ i $26$. Dla $N=24$ mamy $SS_B=162{,}75$, $SS_W=37{,}875$,
zatem $F\approx45{,}119$ o stopniach swobody $(2,21)$ i
wartości p około $2{,}50\cdot10^{-8}$.
Przykładowa miara wielkości efektu to
$\eta^2=SS_B/(SS_B+SS_W)\approx0{,}811$; opisuje część
zmienności empirycznej wyjaśnianą przez podział na grupy,
nie dowodzi zaś przyczynowego wpływu metody.
Hipoteza wszystkich równych średnich zostaje odrzucona przy
konwencjonalnych poziomach istotności. Samo ANOVA nie
wskazuje jednak, \emph{które} pary średnich się różnią.
Na podstawie testu zbiorczego nie wolno też bez dodatkowej
procedury deklarować istotności wszystkich trzech porównań.
\end{zadanie}

\begin{uwaga}[Wielkość efektu a estymator poprawiony]
W jednoczynnikowej ANOVA miara $\eta^2=SS_B/(SS_B+SS_W)$
jest prostym opisem danych, lecz ma dodatnią skłonność do
zawyżania udziału wariancji w małych próbach, także gdy
wszystkie średnie populacyjne są równe. Często raportuje się
również miarę $\omega^2$, zdefiniowaną (przed ewentualnym
obcięciem ujemnego wyniku do zera) przez
\[
 \widehat{\omega}^2=
 \frac{SS_B-(k-1)MS_W}{SS_B+SS_W+MS_W},\qquad
 MS_W:=\frac{SS_W}{N-k}.
\]
Sama wielkość efektu nie koryguje błędów doboru próby ani
nie dowodzi przyczynowości. W raportach należy rozróżnić
miarę empiryczną, estymator parametru populacyjnego oraz
jego niepewność.
\end{uwaga}

\begin{uwaga}[ANOVA, Welch--ANOVA, ANCOVA i MANOVA]
Jednoczynnikowa ANOVA zakłada wspólną wariancję błędów i
niezależność obserwacji. Przy różnicach wariancji można rozważyć
ANOVA Welcha. Jeśli wynik zależy również od ciągłego
kowariantu (np. wieku albo wartości wyjściowej), w określonym
modelu można stosować ANCOVA. Jeśli każda jednostka ma kilka
współzależnych wyników, pojawia się MANOVA. Dla pomiarów
powtarzanych konieczne jest uwzględnienie zależności
wewnątrz jednostki. Każda z tych konstrukcji będzie miała
osobny, szczegółowy rozdział; nie są one zamiennymi etykietami.
\end{uwaga}

\section{Test dokładny dla częstości i testy chi-kwadrat}
\begin{stwierdzenie}[Test dwumianowy jako test dokładny]
Niech $S\sim\Bin(n,p)$, $n\in\N$, $n\ge1$, oraz
$p_0\in(0,1)$. Przy hipotezie $H_0:p\le p_0$ i
alternatywie $H_1:p>p_0$ wartość p dla obserwacji $S=s$
wynosi $P_{p_0}(S\ge s)$. Przy alternatywie
$H_1:p<p_0$ wobec $H_0:p\ge p_0$ wynosi
$P_{p_0}(S\le s)$. Odrzucanie odpowiednio dla małych
wartości p daje poziom co najwyżej $\alpha$.
\end{stwierdzenie}
\begin{proof}
Wynika to z monotoniczności stochastycznej
$\Bin(n,p)$ względem $p$, udowodnionej przez
sprzężenie w twierdzeniu~\ref{thm:binomial-ump}, oraz
z twierdzenia~\ref{thm:pvalue-valid}.
\end{proof}

\begin{twierdzenie}[Statystyka Pearsona dla wielomianowego modelu prostego]
\label{thm:pearson-chi-square}
Niech $m\ge2$ i $p_1,\ldots,p_m>0$, przy
$\sum_jp_j=1$. Niech $(N_1,\ldots,N_m)$ ma rozkład
wielomianowy z parametrami $n$ i
$(p_1,\ldots,p_m)$. Wtedy przy $n\to\infty$
\[
 Q_n=\sum_{j=1}^m\frac{(N_j-np_j)^2}{np_j}
 \xrightarrow{d}\chi^2_{m-1}.
\]
\end{twierdzenie}
\begin{proof}
Każdą z $n$ niezależnych obserwacji reprezentujemy wektorem
jednostkowym $Z_i\in\R^m$ o rozkładzie
$P(Z_i=e_j)=p_j$. Mamy
$\E Z_i=p=(p_1,\ldots,p_m)^{\mathsf T}$ oraz
$\Cov(Z_i)=\diag(p)-pp^{\mathsf T}$.
Z wielowymiarowego CLT
$\sqrt n(N/n-p)\xrightarrow{d}\Normal_m(0,
\diag(p)-pp^{\mathsf T})$.
Niech $D=\diag(p)$ i $u=(\sqrt{p_1},\ldots,\sqrt{p_m})$.
Po przekształceniu liniowym $D^{-1/2}$ macierz
kowariancji granicznej ma postać
$I_m-uu^{\mathsf T}$, czyli rzut ortogonalny na
$(m-1)$-wymiarową podprzestrzeń $u^\perp$, bo
$\|u\|^2=1$. Zgodnie z twierdzeniem o ciągłym
odwzorowaniu $Q_n$ zbiega do kwadratu normy wektora
normalnego z tą macierzą kowariancji. W ortonormalnej
bazie $u^\perp$ jest to suma kwadratów $m-1$
niezależnych zmiennych $\Normal(0,1)$, czyli
$\chi^2_{m-1}$.
\end{proof}

\begin{uwaga}[Estymowane prawdopodobieństwa i małe liczebności]
Powyższe twierdzenie dotyczy \emph{ustalonych}, znanych
wartości $p_j$ pod $H_0$. Jeśli estymujemy parametry z tej
samej próby, liczba stopni swobody może się zmienić; wymaga
to osobnej analizy regularności i rzędu modelu.
Gdy oczekiwane liczebności $np_j$ są małe, asymptotyczne
przybliżenie może być niedokładne. Reguła typu
„wszystkie oczekiwane liczebności co najmniej 5” jest
praktyczną heurystyką, a nie twierdzeniem zapewniającym
określony błąd poziomu testu. Testy dokładne i procedury
symulacyjne bywają wówczas lepsze.
\end{uwaga}

\begin{uwaga}[Poziom 5\% i dwie strony]
Test jednostronny na poziomie $\alpha$ ma dualną
\emph{jednostronną} granicę ufności poziomu $1-\alpha$.
Dwustronny przedział równych ogonów poziomu $1-\alpha$
wykorzystuje zwykle kwantyle odpowiadające $\alpha/2$.
Nie wolno utożsamiać automatycznie obu procedur. W
rozdziale~\ref{ch:confidence-intervals} przeprowadzimy
formalną konstrukcję przez odwracanie testów.
\end{uwaga}

\section{Praktyczne drzewo wyboru testu}
\begin{center}\small
\begin{tabular}{>{\raggedright\arraybackslash}p{0.31\textwidth}>{\raggedright\arraybackslash}p{0.27\textwidth}>{\raggedright\arraybackslash}p{0.31\textwidth}}
\toprule
\textbf{Dane i pytanie} & \textbf{Punkt wyjścia} & \textbf{Co trzeba dodatkowo sprawdzić?}\\
\midrule
Jedna średnia, znana $\sigma$, iid normalne & Test z & Czy $\sigma$ jest naprawdę znana z modelu?\\
Jedna średnia, nieznana $\sigma$ & Test t & Rozkład, ciężkie ogony, niezależność\\
Dwa pomiary tych samych osób & Test t dla różnic & Zależność pomiarów i rozkład \emph{różnic}\\
Dwie niezależne grupy & Test Welcha & Dobór próby, asymetria, ogony, liczebności\\
Więcej niż dwie niezależne grupy & ANOVA lub ANOVA Welcha & Wariancje, niezależność, efekty i post hoc\\
Jedna wariancja w modelu normalnym & Test $\chi^2$ & Wrażliwość na brak normalności\\
Równość dwóch wariancji w modelu normalnym & Test F & Bardzo duża wrażliwość na ogony\\
Odsetek i mała liczba sukcesów & Dokładny test dwumianowy & Schemat próby i kierunek alternatywy\\
Kilka kategorii i znane proporcje & Test Pearsona & Oczekiwane liczebności; parametry estymowane\\
Wiele pomiarów tej samej osoby & Model mieszany lub metody dla powtarzanych pomiarów & Struktura korelacji, braki danych\\
\bottomrule
\end{tabular}
\end{center}

\paragraph{Raportowanie wyników.}
Poprawny raport powinien zawierać rodzaj testu, hipotezę
zerową i alternatywną, jednostkę analizy, liczebności,
oszacowanie efektu i jego przedział, statystykę wraz ze
stopniami swobody, wartość p, sprawdzone założenia i
ograniczenia. Nie należy pisać „test potwierdził normalność”
ani „brak istotności dowodzi równości”.

\begin{lstlisting}[language=Python,caption={Kompletny przykład: Welch, ANOVA i diagnostyka w Pythonie (SciPy).}]
from scipy.stats import f_oneway, levene, shapiro, ttest_ind

A = [18, 19, 20, 19, 21, 20, 18, 22]
B = [22, 23, 21, 24, 22, 25, 23, 24]
C = [25, 27, 26, 24, 28, 26, 27, 25]

welch = ttest_ind(A, B, equal_var=False)
anova = f_oneway(A, B, C)
normality_A = shapiro(A)
variances = levene(A, B, C, center="median")

print("Welch A vs B:", welch)
print("ANOVA A, B, C:", anova)
print("Shapiro-Wilk A:", normality_A)
print("Brown-Forsythe (Levene median):", variances)
\end{lstlisting}

Ten program ilustruje sposób wywołania procedur, ale jego
wyniki wymagają \emph{analizy projektu}, a nie tylko odczytania
kolumny p-value. Shapiro--Wilk dla surowej grupy A nie jest
uniwersalnym pretestem, od którego zależy wybór ANOVA;
jednorodność wariancji też nie wynika z dużej wartości p.

\Needspace{20\baselineskip}
\begin{lstlisting}[language=Julia,caption={Porównanie grup metodą Welcha w Julii: obliczenie statystyki i p-value.}]
using Statistics, Distributions

A = [18, 19, 20, 19, 21, 20, 18, 22]
B = [22, 23, 21, 24, 22, 25, 23, 24]
n, m = length(A), length(B)
vA, vB = var(A), var(B)  # wariancja z mianownikiem n-1
se2 = vA / n + vB / m
tstat = (mean(A) - mean(B)) / sqrt(se2)
df = se2^2 / ((vA/n)^2/(n-1) + (vB/m)^2/(m-1))
pvalue = 2 * cdf(TDist(df), -abs(tstat))
println((t=tstat, df=df, p=pvalue))
\end{lstlisting}

Pakiet \texttt{Distributions.jl} jest zależnością zewnętrzną;
przed uruchomieniem należy go zainstalować w używanym projekcie.
Wynik tego programu jest aproksymacją Satterthwaite'a,
a nie dokładnym rozkładem t przy nierównych wariancjach.

\section{Studium decyzyjne: od wyniku testu do zalecenia}
Rozważmy ponownie trzy technologie z zadania~\ref{zad:anova-three}.
Wynik ANOVA daje bardzo mocne świadectwo przeciw hipotezie
jednakowych średnich, ale sam nie określa najlepszej decyzji.
Przed wyborem technologii należy odpowiedzieć na następujące
pytania.

\begin{enumerate}
\item \emph{Co jest parametrem celu?} Jeśli mierzony czas
pracy chcemy minimalizować, sama statystyczna istotność
różnicy nie przesądza, czy oszczędność czasu rekompensuje
koszty urządzeń. Warto podać różnice średnich w jednostkach
fizycznych, a nie jedynie statystykę $F$.
\item \emph{Jak powstały grupy?} Jeśli te same urządzenia
badano wielokrotnie, nie można ich traktować jako trzech
prób niezależnych. Jeśli technologię przydzielano zależnie
od typu maszyny, nie można przypisać stwierdzonej różnicy
wyłącznie wpływowi technologii.
\item \emph{Jak sprawdzono model?} Wykresy reszt, rozkład
błędów warunkowych i możliwa heterogeniczność wariancji
są ważniejsze niż mechaniczne uruchomienie testu
Shapiro--Wilka osobno dla każdej grupy. Przy nierównych
wariancjach trzeba rozważyć ANOVA Welcha.
\item \emph{Które porównania są celem?} Jeśli wcześniej
zaplanowano porównanie A z B, można analizować ten kontrast.
Jeżeli po obejrzeniu danych wybieramy wszystkie trzy pary,
trzeba kontrolować wielokrotność (np. procedurą Tukeya
w odpowiednim modelu lub Gamesa--Howella przy nierównych
wariancjach).
\item \emph{Jak raportować?} Należy podać liczebności,
średnie i rozproszenie w grupach, sumy kwadratów, stopnie
swobody, wartość $F$ i p-value, miarę efektu oraz
przedziały ufności dla istotnych praktycznie kontrastów.
Dodatkowo trzeba wskazać ograniczenia modelu i możliwe
czynniki zakłócające.
\end{enumerate}

\begin{uwaga}[Brak istotności a równoważność]
Nawet jeśli test różnicy średnich nie odrzuca hipotezy
zerowej, nie wolno mówić, że dwie metody są równoważne.
Aby wykazać równoważność, należy z góry określić
akceptowalny margines $\Delta>0$ oraz przetestować
odpowiednio hipotezę braku równoważności. Omówimy
formalną procedurę dwóch jednostronnych testów (TOST)
w części poświęconej rozszerzeniom testowania hipotez.
\end{uwaga}



% ================================================================
% ROZDZIAL 11 -- PRZEDZIALY UFNOSCI
% ================================================================
\chapter{Przedziały ufności: konstrukcja, pokrycie i interpretacja}
\label{ch:confidence-intervals}
\paragraph{Powiązania i odsyłacze.} Związek przedziałów z testami wyjaśnia rozdział~\ref{ch:testing}; przykłady dokładnych rozkładów statystyk podaje rozdział~\ref{ch:parametric-tests}.


\section{Co dokładnie oznacza 95\% ufności?}
Zaobserwowany odsetek wadliwych produktów wynosi $7/100$. Samo
$\widehat p=0{,}07$ nie mówi, jaka jest precyzja wnioskowania.
Przedział ufności jest sposobem opisania tej niepewności, ale
nie jest rozkładem prawdopodobieństwa nieznanego parametru.

\begin{definicja}[Zbiór ufności i pokrycie]\label{def:confidence-set}
Niech $(\mathsf X,\mathcal A,\{P_\theta:\theta\in\Theta\})$
będzie eksperymentem statystycznym, a
$g:\Theta\to\mathsf T$ ustaloną funkcją parametru,
gdzie $(\mathsf T,\mathcal T)$ jest przestrzenią mierzalną.
\emph{Zbiorem ufności} na poziomie $1-\alpha$, gdzie
$0<\alpha<1$, nazywamy losowy zbiór $C(X)\subseteq\mathsf T$
taki, że dla każdego $\theta\in\Theta$ zdarzenie
$\{g(\theta)\in C(X)\}$ jest mierzalne i
\[
 P_\theta\bigl(g(\theta)\in C(X)\bigr)\ge1-\alpha.
\]
Jeśli $\mathsf T=\R$ i $C(X)=[L(X),U(X)]$ dla
mierzalnych $L\le U$, nazywamy go \emph{przedziałem ufności}.
Powyższe prawdopodobieństwo jest \emph{pokryciem} przy $\theta$.
Gdy dla każdego parametru zachodzi równość, mówimy o
\emph{dokładnym pokryciu nominalnym}; procedura z pokryciem
większym od nominalnego jest \emph{konserwatywna}.
\end{definicja}

Losowość należy do przedziału $C(X)$, zanim zaobserwujemy próbę.
Po podstawieniu konkretnych danych $x$ otrzymujemy ustalony zbiór
$C(x)$. W interpretacji klasycznej nie stwierdzamy wtedy, że
„parametr z prawdopodobieństwem $95\%$ leży w tym konkretnym
przedziale”. Stwierdzamy, że \emph{procedura} konstruowania
przedziałów ma zadane pokrycie w powtarzanych doświadczeniach.
Rozkład wiarygodności a posteriori prowadzący do innego rodzaju
przedziału zostanie wprowadzony w części bayesowskiej.

\begin{uwaga}[Pokrycie punktowe, jednostajne i przybliżone]
Zapis $P_\theta(g(\theta)\in C_n)\to1-\alpha$ dla każdego
ustalonego $\theta$ oznacza pokrycie asymptotyczne \emph{punktowe}.
Silniejsze żądanie
$\liminf_{n\to\infty}\inf_{\theta\in\Theta}
P_\theta(g(\theta)\in C_n)\ge1-\alpha$
gwarantuje pokrycie asymptotyczne \emph{jednostajne}.
Pierwsze nie implikuje drugiego; przy granicach przestrzeni
parametrów klasyczne przybliżenie normalne może działać
bardzo źle. W zastosowaniach małych prób pytamy zatem nie tylko
o nominalny procent, lecz także o zakres parametrów, dla których
gwarancja jest prawdziwa.
\end{uwaga}

\section{Dualność testów i przedziałów ufności}
\begin{twierdzenie}[Odwracanie rodziny testów]\label{thm:test-inversion}
Niech $\{P_\theta:\theta\in\Theta\}$ będzie rodziną miar
na $(\mathsf X,\mathcal A)$ oraz $g:\Theta\to\mathsf T$.
Dla każdego $t\in g(\Theta)$ niech $\varphi_t:\mathsf X\to\{0,1\}$
będzie mierzalnym testem hipotezy
$H_{0,t}:g(\theta)=t$ o poziomie co najwyżej $\alpha$:
\[
 \sup_{\theta:g(\theta)=t}\E_\theta\varphi_t(X)\le\alpha.
\]
Zakładamy ponadto, że zbiór
$C(x)=\{t\in g(\Theta):\varphi_t(x)=0\}$ ma mierzalne
zdarzenia przynależności. Wtedy $C(X)$ jest zbiorem ufności
dla $g(\theta)$ o pokryciu co najmniej $1-\alpha$.
\end{twierdzenie}
\begin{proof}
Ustalmy $\theta\in\Theta$ i przyjmijmy $t=g(\theta)$.
Z definicji konstrukcji $g(\theta)\notin C(X)$ wtedy i tylko
wtedy, gdy $\varphi_t(X)=1$. Zatem
$P_\theta(g(\theta)\notin C(X))
=\E_\theta\varphi_t(X)\le\alpha$, co daje tezę.
Nie trzeba zakładać, że test dla każdego $t$ odrzuca
jednakowe wartości obserwowanej statystyki: istotna jest
\emph{jednoczesna definicja rodziny testów}, a nie pojedynczy
test dopasowany do danych.
\end{proof}

\begin{uwaga}[Co oznacza odwrotna implikacja?]
Dla danego zbioru ufności $C(X)$ można zdefiniować test
hipotezy $g(\theta)=t$ przez
$\varphi_t(X)=\ind{\{t\notin C(X)\}}$. Jeśli
pokrycie jest co najmniej $1-\alpha$, każdy z tych testów ma
poziom co najwyżej $\alpha$. Testy i przedziały są więc
różnymi reprezentacjami tej samej reguły wnioskowania,
\emph{o ile} zostały skonstruowane z tą samą rodziną hipotez
i tymi samymi założeniami.
\end{uwaga}

\section{Prawidłowy model normalny: przedział dla średniej i wariancji}
\begin{twierdzenie}[Średnia i wariancja przy nieznanych parametrach]
\label{thm:normal-mean-var-ci}
Niech $n\ge2$, $X_1,\ldots,X_n\iid\Normal(\mu,\sigma^2)$,
$\mu\in\R$, $\sigma>0$, oraz niech $\bar X$ i $S^2$ oznaczają
średnią i wariancję z mianownikiem $n-1$. Dla $0<\alpha<1$
przedziały
\[
 C_\mu(X)=\left[\bar X-t_{n-1,1-\alpha/2}\frac{S}{\sqrt n},
 \bar X+t_{n-1,1-\alpha/2}\frac{S}{\sqrt n}\right],
\]
\[
 C_{\sigma^2}(X)=\left[
 \frac{(n-1)S^2}{\chi^2_{n-1,1-\alpha/2}},
 \frac{(n-1)S^2}{\chi^2_{n-1,\alpha/2}}\right]
\]
mają \emph{każdy osobno} pokrycie dokładnie $1-\alpha$.
Symbol $\chi^2_{\nu,q}$ oznacza kwantyl rzędu $q$ rozkładu
chi-kwadrat o $\nu>0$ stopniach swobody; symbol
$t_{\nu,q}$ wprowadzono w rozdziale~\ref{ch:parametric-tests}.
\end{twierdzenie}
\begin{proof}
Rozkład próby normalnej zapewnia niezależne zmienne
$Z=\sqrt n(\bar X-\mu)/\sigma\sim\Normal(0,1)$ i
$V=(n-1)S^2/\sigma^2\sim\chi^2_{n-1}$. Stąd
$Z/\sqrt{V/(n-1)}\sim t_{n-1}$. Zdarzenie
$\{|Z/\sqrt{V/(n-1)}|\le t_{n-1,1-\alpha/2}\}$ jest
dokładnie zdarzeniem $\{\mu\in C_\mu(X)\}$ i ma
prawdopodobieństwo $1-\alpha$. Analogicznie
$\{\chi^2_{n-1,\alpha/2}\le V\le
\chi^2_{n-1,1-\alpha/2}\}$ ma prawdopodobieństwo
$1-\alpha$ i po rozwiązaniu nierówności względem
$\sigma^2$ daje drugi przedział. Użyliśmy ciągłości
rozkładów; dla obu przedziałów razem nie zadeklarowano
tu pokrycia jednoczesnego.
\end{proof}

\begin{stwierdzenie}[Przedział dla ilorazu dwóch wariancji]
\label{prop:ratio-var-ci}
Niech próby $X_1,\ldots,X_n\iid\Normal(\mu_X,\sigma_X^2)$
i $Y_1,\ldots,Y_m\iid\Normal(\mu_Y,\sigma_Y^2)$ będą
niezależne, $n,m\ge2$, $\sigma_X,\sigma_Y>0$.
Jeśli $R=S_X^2/S_Y^2$ oraz $q=\sigma_X^2/\sigma_Y^2$,
to przedział
\[
 \left[\frac{R}{F_{n-1,m-1,1-\alpha/2}},
       \frac{R}{F_{n-1,m-1,\alpha/2}}\right]
\]
obejmuje $q$ z prawdopodobieństwem $1-\alpha$,
gdzie $F_{a,b,u}$ jest kwantylem rzędu $u$
rozkładu Fishera--Snedecora z $a,b>0$ stopniami swobody.
\end{stwierdzenie}
\begin{proof}
Zmienne
$(n-1)S_X^2/\sigma_X^2$ i
$(m-1)S_Y^2/\sigma_Y^2$ są niezależne i mają
rozkłady chi-kwadrat. Z definicji rozkładu F
$R/q\sim F_{n-1,m-1}$, więc odwrócenie nierówności
między kwantylami daje wskazany przedział.
\end{proof}

\begin{przyklad}[Dwie niezależne technologie a pomiary sparowane]
W zadaniu~\ref{zad:welch-case} dane pochodzą z dwóch
\emph{niezależnych} grup. Przy nierównych wariancjach
przedział dla $\mu_X-\mu_Y$ uzyskany metodą Welcha ma
postać
$[\bar X-\bar Y\pm t_{\nu,1-\alpha/2}
\sqrt{S_X^2/n+S_Y^2/m}]$, gdzie losowe $\nu$
zdefiniowano w rozdziale~\ref{ch:parametric-tests}.
Jest to przedział \emph{przybliżony}, nie dokładny.
Jeżeli te same jednostki zmierzono dwukrotnie i różnice
$D_i$ są niezależne między parami oraz normalne,
przedział oparty na $\bar D\pm t_{n-1,1-\alpha/2}S_D/\sqrt n$
jest \emph{dokładny}. Nie wolno podmieniać tych procedur
bez zmiany założeń o zależności.
\end{przyklad}

\section{Odsetek wad: trzy przedziały dla parametru Bernoulliego}
Niech $X_1,\ldots,X_n\iid\operatorname{Bernoulli}(p)$,
$n\ge1$, $p\in[0,1]$, $S=\sum_iX_i$ i
$\widehat p=S/n$. Trzy popularne konstrukcje odpowiadają
na to samo pytanie, ale opierają się na różnych argumentach.

\paragraph{Wald: prosta formuła bez gwarancji przy małych częstościach.}
Dla $z=z_{1-\alpha/2}$ przedział Walda ma postać
\[
 C_{\rm Wald}(X)=\left[\widehat p-z\sqrt{\frac{\widehat p(1-\widehat p)}n},
\widehat p+z\sqrt{\frac{\widehat p(1-\widehat p)}n}\right]
\cap[0,1].
\]
Z twierdzenia centralnego granicznego i twierdzenia Slutsky'ego
wynika pokrycie zbieżne do $1-\alpha$ \emph{dla każdego
ustalonego} $p\in(0,1)$. Nie jest to jednak gwarancja
jednostajna, szczególnie blisko $p=0$ i $p=1$.
Przy $S=0$ przedział redukuje się do $\{0\}$, więc dla
każdego prawdziwego $p>0$ ten wynik próby nie pokrywa
parametru. W praktyce nie zalecamy Walda dla rzadkich
zdarzeń ani małych prób.

\begin{twierdzenie}[Przedział Wilsona przez odwrócenie testu score]
\label{thm:wilson-ci}
Przy powyższych założeniach i $0<\alpha<1$ niech
$z=z_{1-\alpha/2}$. Zbiór parametrów $p\in[0,1]$
spełniających
\[
 \frac{n(\widehat p-p)^2}{p(1-p)}\le z^2,
 \qquad 0<p<1,
\]
z domknięciem rozwiązań przy $p=0,1$, jest przedziałem
$[c-h,c+h]$, gdzie
\[
 c=\frac{\widehat p+z^2/(2n)}{1+z^2/n},\qquad
 h=\frac{z}{1+z^2/n}\sqrt{
 \frac{\widehat p(1-\widehat p)}n+\frac{z^2}{4n^2}}.
\]
To \emph{przedział Wilsona}. Jego pokrycie wynosi
$1-\alpha+o(1)$ przy $n\to\infty$ dla każdego
ustalonego $p\in(0,1)$, lecz nie musi wynosić dokładnie
$1-\alpha$ dla skończonego $n$.
\end{twierdzenie}
\begin{proof}
W modelu Bernoulliego dla ustalonego $p\in(0,1)$
statystyka score $Z_p=\sqrt n(\widehat p-p)/\sqrt{p(1-p)}$
ma na mocy CLT rozkład asymptotycznie standardowy normalny.
Nierówność $|Z_p|\le z$ jest równoważna
$n(\widehat p-p)^2\le z^2p(1-p)$, czyli
\[
 (n+z^2)p^2-(2n\widehat p+z^2)p+n\widehat p^2\le0.
\]
Trójmian kwadratowy o dodatnim współczynniku przy $p^2$
przyjmuje wartości niedodatnie między pierwiastkami.
Wzory na środek i połowę ich różnicy upraszczają się
odpowiednio do $c$ i $h$. Dla $\widehat p=0$ albo $1$
odpowiednie domknięcie w punkcie brzegowym daje tę samą
formułę; otrzymane końce zawsze należą do $[0,1]$.
Przy ustalonym $p\in(0,1)$ na mocy CLT prawdopodobieństwo
nierówności dąży do $1-\alpha$. Nie wyciągamy wniosku
o jednostajnym pokryciu bez dodatkowego dowodu.
\end{proof}

\begin{twierdzenie}[Dokładny konserwatywny przedział Cloppera--Pearsona]
\label{thm:clopper-pearson}
Dla $0<\alpha<1$, $s\in\{0,\ldots,n\}$ definiujemy
$L(s)$ i $U(s)$ następująco:
\begin{itemize}
\item $L(0)=0$, a dla $s\ge1$ liczba $L(s)\in(0,1)$
jest jedynym rozwiązaniem $P_{L(s)}(S\ge s)=\alpha/2$;
\item $U(n)=1$, a dla $s\le n-1$ liczba $U(s)\in(0,1)$
jest jedynym rozwiązaniem $P_{U(s)}(S\le s)=\alpha/2$,
\end{itemize}
gdzie w obu równaniach $S\sim\Bin(n,p)$ z parametrem
$p$ wskazanym w indeksie. Wtedy dla każdego $p\in[0,1]$
\[
 P_p\bigl(L(S)\le p\le U(S)\bigr)\ge1-\alpha.
\]
\end{twierdzenie}
\begin{proof}
Dla $1\le s\le n$ funkcja
$p\mapsto P_p(S\ge s)$ jest ciągła i ściśle rosnąca
na $(0,1)$, z granicami $0$ i $1$; wynika to na przykład
z różniczkowania sumy dwumianowej albo ze sprzężenia
wskaźnikami $\ind{\{U_i\le p\}}$ z rozdziału~\ref{ch:testing}.
Analogicznie $p\mapsto P_p(S\le s)$ dla
$0\le s\le n-1$ jest ściśle malejąca od $1$ do $0$.
Zatem końce istnieją i są jednoznaczne.

Ustalmy prawdziwe $p\in(0,1)$ i oznaczmy przez
$S'$ pomocniczą zmienną o rozkładzie $\Bin(n,p)$.
Jeśli $p<L(S)$, to z monotoniczności ogona w parametrze
otrzymujemy $P_p(S'\ge S)<\alpha/2$.
Dla każdego dyskretnego rozkładu funkcja
$u(s)=P_p(S'\ge s)$ jest nierosnąca, a zbiór
$\{s:u(s)\le t\}$ jest końcowym przedziałem liczb
całkowitych (albo jest pusty). Jeżeli jego najmniejszy
element to $s_t$, mamy
$P_p(u(S)\le t)=P_p(S\ge s_t)=u(s_t)\le t$.
Stąd $P_p(p<L(S))\le\alpha/2$.
Analogicznie używając lewego ogona $P_p(S'\le S)$
otrzymujemy $P_p(p>U(S))\le\alpha/2$.
Nierówność sumy prawdopodobieństw daje
$P_p(p\notin[L(S),U(S)])\le\alpha$.
Dla $p=0$ lub $p=1$ właściwy koniec przedziału jest
równy parametrowi z prawdopodobieństwem jeden.
\end{proof}

\begin{uwaga}[Dokładny test, lecz niekoniecznie pokrycie równe nominalnemu]
Słowo „dokładny” odnosi się tu do gwarancji skończonopróbkowej
bez przybliżenia normalnego. Wskutek dyskretności rozkładu
pokrycie Cloppera--Pearsona jest często \emph{większe} niż
$1-\alpha$, a przedział może być szerszy niż Wilsona.
Nie należy mówić, że „Clopper--Pearson zawsze pokrywa dokładnie
$95\%$” w znaczeniu równości dla każdego $p$.
\end{uwaga}

\begin{zadanie}[7 wad na 100 sztuk: trzy różne odpowiedzi]
\label{zad:ci-seven}
Dla $n=100$, $s=7$ i $\alpha=0{,}05$ oblicz trzy
przedziały dla $p$ i zdecyduj, czy istnieją mocne podstawy,
by przyjąć $p<0{,}1$.
\par\emph{Rozwiązanie.} Punktowy estymator wynosi
$\widehat p=0{,}07$. Dla $z_{0{,}975}\approx1{,}96$
po podstawieniu do wzorów otrzymujemy
\[
 \begin{aligned}
 C_{\mathrm{Wald}}&\approx[0{,}01999;\ 0{,}12001],\\
 C_{\mathrm{Wilson}}&\approx[0{,}03432;\ 0{,}13750],\\
 C_{\mathrm{CP}}&\approx[0{,}02861;\ 0{,}13892].
 \end{aligned}
\]
Przedziały różnią się istotnie dolnym końcem, choć
wszystkie obejmują wartość $0{,}1$. Dla hipotezy
$H_0:p\ge0{,}1$ przeciw $p<0{,}1$ z rozdziału~\ref{ch:testing}
\emph{dokładna jednostronna górna granica} poziomu $95\%$
to $U_1(s)$ rozwiązujące
$P_{U_1(s)}(S\le s)=0{,}05$; przy $s=7$
wynosi $U_1(7)\approx0{,}12746$. Wobec tego
nie możemy na poziomie $5\%$ stwierdzić, że
$p<0{,}1$. Równoważne jest to temu, że jednostronny
test dokładny daje $p\text{-value}\approx0{,}20605$.
\emph{Dwustronny} przedział ufności poziomu $95\%$
nie jest automatycznie dualny do \emph{jednostronnego}
testu na poziomie $5\%$; ich zgodność tutaj nie
usprawiedliwia utożsamiania obu procedur.
\end{zadanie}

\begin{zadanie}[Przedział dla średniej różnicy w pomiarach sparowanych]
W zadaniu z rozdziału~\ref{ch:parametric-tests} osiem
niezależnych par miało różnice
$D=(3,2,4,3,1,2,5,2)$. Pod założeniem normalności
różnic przedział $95\%$ dla średniej różnicy $\delta$
ma postać $\bar D\pm t_{7,0{,}975}S_D/\sqrt8$.
Ponieważ $\bar D=2{,}75$, $S_D^2=23/14$ i
$t_{7,0{,}975}\approx2{,}365$, otrzymujemy
przedział w przybliżeniu $[1{,}68;\ 3{,}82]$.
Wnioskujemy o średniej \emph{różnicy wewnątrz jednostki},
a nie o różnicy średnich z dwóch niezależnych prób.
Interpretacja przyczynowa nadal wymaga odpowiedniego
projektu badania.
\end{zadanie}

\section{Proces Poissona: przedział dla intensywności zdarzeń}
Jeśli w czasie ekspozycji $T>0$ zaobserwowano $S=s$
zdarzeń, a model zakłada $S\sim\Pois(\lambda T)$,
to estymator $\widehat\lambda=s/T$ nie opisuje niepewności
przy małych $s$. W szczególności $s=0$ nie dowodzi,
że $\lambda=0$.

\begin{twierdzenie}[Dokładny przedział dla intensywności Poissona]
\label{thm:poisson-ci}
Niech $T>0$, $S\sim\Pois(\lambda T)$, $\lambda\ge0$,
oraz $0<\alpha<1$. Zdefiniujmy $a(0)=0$,
$a(s)=\chi^2_{2s,\alpha/2}/2$ dla $s\ge1$, a
$b(s)=\chi^2_{2(s+1),1-\alpha/2}/2$ dla $s\ge0$.
Wówczas
\[
 C_\lambda(S)=\left[\frac{a(S)}T,\frac{b(S)}T\right]
\]
ma pokrycie co najmniej $1-\alpha$ dla wszystkich
$\lambda\ge0$.
\end{twierdzenie}
\begin{proof}
Dla $s\ge1$ i $u>0$ związek pomiędzy rozkładami
Poissona i chi-kwadrat daje
\[
 P_{u}(S\ge s)=P(\chi^2_{2s}\le2u),\qquad
 P_{u}(S\le s)=P(\chi^2_{2(s+1)}>2u),
\]
gdzie w ogonach po lewej $S\sim\Pois(u)$.
Tożsamości wynikają z całkowania gęstości gamma albo
z pochodnej odpowiednich skończonych sum
$e^{-u}\sum_{j=0}^{s-1}u^j/j!$ i
$e^{-u}\sum_{j=0}^{s}u^j/j!$ z warunkiem brzegowym
przy $u=0$. Stąd $a(s)$ rozwiązuje
$P_{a(s)}(S\ge s)=\alpha/2$, a $b(s)$ rozwiązuje
$P_{b(s)}(S\le s)=\alpha/2$.
Tak jak w dowodzie twierdzenia~\ref{thm:clopper-pearson},
przekroczenie dolnego lub górnego końca oznacza wystąpienie
odpowiedniej jednostronnej wartości p mniejszej niż
$\alpha/2$. Nierówność sumy prawdopodobieństw kończy
dowód. Dla $\lambda=0$ mamy $S=0$ prawie na pewno
i przedział zawiera zero.
\end{proof}

\begin{przyklad}[Zero awarii nie oznacza zerowej intensywności]
Przy $s=0$ dolny koniec przedziału wynosi zero,
a górny wynosi $-\log(\alpha/2)/T$. Dla
$\alpha=0{,}05$ i $T=100$ godzin daje to około
$0{,}03689$ awarii na godzinę. Jeśli z góry
interesuje nas wyłącznie \emph{górna} granica, można
uzyskać krótszy, jednostronny przedział z końcem
$-\log(\alpha)/T\approx0{,}02996$; trzeba jednak
podać w raporcie, że chodzi o granicę jednostronną.
\end{przyklad}

\section{Gdy nie ufamy rozkładowi: przedział z nierówności Hoeffdinga}
\begin{twierdzenie}[Przedział rozkładowo niezależny dla średniej]
\label{thm:hoeffding-ci}
Niech $X_1,\ldots,X_n$ będą niezależnymi zmiennymi
losowymi o wspólnej wartości oczekiwanej $\mu$,
a przy tym $a\le X_i\le b$ prawie na pewno dla
pewnych znanych liczb $a<b$. Dla $0<\alpha<1$
przedział
\[
 \left[\bar X-(b-a)\sqrt{\frac{\log(2/\alpha)}{2n}},
       \bar X+(b-a)\sqrt{\frac{\log(2/\alpha)}{2n}}\right]
\]
ma pokrycie co najmniej $1-\alpha$.
\end{twierdzenie}
\begin{proof}
Nierówność Hoeffdinga dla niezależnych zmiennych
ograniczonych z tym samym przedziałem wartości zapewnia
$P(|\bar X-\mu|\ge\varepsilon)
\le2\exp[-2n\varepsilon^2/(b-a)^2]$.
Po podstawieniu
$\varepsilon=(b-a)\sqrt{\log(2/\alpha)/(2n)}$
prawa strona jest równa $\alpha$. Ponieważ zdarzenie
niepokrycia przedziałem implikuje
$|\bar X-\mu|\ge\varepsilon$, otrzymujemy tezę.
Nie wymagana jest normalność, ale \emph{znane ograniczenie}
$a\le X_i\le b$ jest istotnym założeniem.
\end{proof}

\begin{uwaga}[Praktyczne ograniczenia]
Przedział Hoeffdinga bywa szerszy od przedziału t lub Wilsona,
ponieważ kupujemy gwarancję przy słabszych założeniach
o rozkładzie. Nie działa automatycznie dla nieograniczonego
czasu awarii, wielokrotnych pomiarów tej samej osoby ani
danych z nieznanym mechanizmem selekcji.
Przedział $1-\alpha$ dla jednego ustalonego $n$ nie jest
automatycznie \emph{sekwencją ufności} zachowującą pokrycie
przy dowolnym, wybranym adaptacyjnie czasie zatrzymania.
\end{uwaga}

\section{Jak wybrać przedział i jak go raportować?}
\begin{center}\small
\begin{tabular}{>{\raggedright\arraybackslash}p{0.31\textwidth}>{\raggedright\arraybackslash}p{0.29\textwidth}>{\raggedright\arraybackslash}p{0.30\textwidth}}
\toprule
\textbf{Cel i dane} & \textbf{Konstrukcja} & \textbf{Co trzeba sprawdzić?}\\
\midrule
Średnia przy iid normalnych, $\sigma$ nieznana & Przedział t & Niezależność, normalność modelu; przy małych $n$ szczególnie istotne.\\
Średnia w dużej próbie o skończonej wariancji & Asymptotyczny przedział ze średniej i SE & Ciężkie ogony, zależność, wiarygodność oszacowania SE.\\
Średnia przy znanych granicach obserwacji & Hoeffding & Czy każde $X_i$ naprawdę mieści się w $[a,b]$?\\
Prawdopodobieństwo sukcesu w próbie Bernoulliego & Wilson & Przybliżone pokrycie; ostrożność przy skrajnych $p$.\\
Prawdopodobieństwo, mała próba lub skrajne częstości & Clopper--Pearson & Konserwatywność, możliwie duża szerokość.\\
Wariancja/iloraz wariancji próby normalnej & Chi-kwadrat/F & Wrażliwość na brak normalności.\\
Różnica średnich dwóch niezależnych grup & Welch (przybliżony) & Niezależność grup, ogony, próby niezrównoważone.\\
Intensywność zliczeń i znana ekspozycja & Dokładny Poisson & Model stałej intensywności; nadmierna dyspersja.\\
\bottomrule
\end{tabular}
\end{center}

\paragraph{Co koniecznie podajemy?}
Podajemy estymowany parametr i jego jednostkę, schemat
losowania, zastosowaną konstrukcję, poziom ufności,
liczebność lub ekspozycję, punktowe oszacowanie oraz
oba końce przedziału. W metodach asymptotycznych warto
podać przybliżenie i przeprowadzić analizę wrażliwości.
Jeżeli porównujemy wiele efektów naraz, przedziały
\emph{indywidualne} nie zapewniają automatycznie
jednoczesnego pokrycia całej rodziny.

\begin{lstlisting}[language=Python,caption={Porównanie przedziałów Walda, Wilsona i Cloppera--Pearsona oraz granica dla Poissona (SciPy).}]
from math import sqrt
from scipy.stats import beta, chi2, norm

n, s, alpha = 100, 7, 0.05
phat = s / n
z = norm.ppf(1 - alpha / 2)
wald = (max(0.0, phat - z * sqrt(phat * (1-phat)/n)),
        min(1.0, phat + z * sqrt(phat * (1-phat)/n)))
center = (phat + z*z/(2*n)) / (1 + z*z/n)
half = z/(1 + z*z/n) * sqrt(phat*(1-phat)/n + z*z/(4*n*n))
wilson = (center - half, center + half)
clopper = (0.0 if s == 0 else beta.ppf(alpha/2, s, n-s+1),
           1.0 if s == n else beta.ppf(1-alpha/2, s+1, n-s))
print('Wald:', wald)
print('Wilson:', wilson)
print('Clopper-Pearson:', clopper)

s_events, exposure = 0, 100.0
rate_upper_2s = chi2.ppf(1-alpha/2, 2*s_events+2) / (2*exposure)
print('Poisson upper (two-sided):', rate_upper_2s)
\end{lstlisting}

\begin{lstlisting}[language=Julia,caption={Przedział Wilsona i dokładny przedział dwumianowy w Julii (Distributions.jl).}]
using Distributions

n, s, alpha = 100, 7, 0.05
phat = s / n
z = quantile(Normal(), 1 - alpha/2)
center = (phat + z^2/(2n))/(1+z^2/n)
half = z/(1+z^2/n)*sqrt(phat*(1-phat)/n + z^2/(4n^2))
wilson = (center-half, center+half)
lower = s == 0 ? 0.0 : quantile(Beta(s, n-s+1), alpha/2)
upper = s == n ? 1.0 : quantile(Beta(s+1, n-s), 1-alpha/2)
println((wilson=wilson, clopper=(lower, upper)))
\end{lstlisting}

\paragraph{Interpretacja kodu.}
Funkcje \texttt{beta.ppf} oraz \texttt{quantile(Beta(...))}
służą wyłącznie jako wygodny sposób numerycznego
rozwiązania równań definiujących Cloppera--Pearsona;
nie oznacza to, że z góry przyjęliśmy bayesowski rozkład
a priori. Kompletne uzasadnienie pokrycia zostało podane
w twierdzeniu~\ref{thm:clopper-pearson}.


% ================================================================
% ROZDZIAL 12 -- WIELOKROTNE TESTOWANIE
% ================================================================
\chapter{Wielokrotne testowanie: FWER, FDR i decyzje po selekcji}
\label{ch:multiple-testing}
\paragraph{Powiązania i odsyłacze.} Przed użyciem korekt warto przeczytać rozdziały~\ref{ch:testing} i~\ref{ch:confidence-intervals}; do projektowania badań przejdziemy w rozdziale~\ref{ch:causal-ab}.


\section{Dlaczego test wykonany poprawnie może prowadzić do błędnej procedury?}
Testowanie jednej hipotezy z poziomem $\alpha=0{,}05$ jest
różne od testowania stu hipotez, z których następnie
wybieramy tylko te o małych wartościach p. Każdy test
może być poprawny oddzielnie, a cała procedura może mieć
duże prawdopodobieństwo przynajmniej jednego
fałszywego odkrycia. Trzeba więc ustalić, \emph{jaki rodzaj
błędu kontrolujemy w zbiorze decyzji}.

\begin{definicja}[Rodzina hipotez i miary błędów]\label{def:fwer-fdr}
Niech $m\ge1$, a $H_1,\ldots,H_m$ będą ustaloną rodziną
hipotez. Niech $I_0\subseteq\{1,\ldots,m\}$ oznacza
zbiór indeksów prawdziwych hipotez zerowych i
$m_0=|I_0|$ (nieznane badaczowi). Dla testów
$\varphi_i\in\{0,1\}$ definiujemy liczbę wszystkich
odrzuceń $R=\sum_{i=1}^m\varphi_i$ oraz liczbę
błędnych odrzuceń $V=\sum_{i\in I_0}\varphi_i$.
\emph{Rodzinne prawdopodobieństwo błędu I rodzaju} to
$\mathrm{FWER}=P(V\ge1)$, a \emph{oczekiwany udział
fałszywych odkryć} to
\[
 \mathrm{FDR}=\E\left[\frac{V}{R\vee1}\right],
 \qquad R\vee1:=\max\{R,1\}.
\]
Konwencja $V/(R\vee1)=0$ dla $R=0$ jest częścią definicji.
Dodatkowo $\mathrm{FDP}=V/(R\vee1)$ oznacza
\emph{zrealizowany} udział fałszywych odkryć.
\end{definicja}

FWER dotyczy zdarzenia, że \emph{jakikolwiek} wniosek
jest fałszywy. FDR dopuszcza pewną liczbę błędnych
odkryć, ale kontroluje \emph{średni} ich udział.
Nie oznacza to, że w danym konkretnym badaniu udział
błędnych odkryć jest nie większy niż $5\%$.

\begin{definicja}[Silna i słaba kontrola FWER]
Procedura ma \emph{silną kontrolę FWER} na poziomie
$\alpha$, jeśli $P(V\ge1)\le\alpha$ dla każdego
rozmieszczenia prawdziwych i fałszywych hipotez w rozważanej
rodzinie modeli. \emph{Słaba kontrola FWER} wymaga tej
nierówności tylko w konfiguracji, w której wszystkie
hipotezy zerowe są prawdziwe.
\end{definicja}

\begin{przyklad}[Dwadzieścia niezależnych testów]
Jeśli pod prawdziwymi hipotezami zerowymi
$p_1,\ldots,p_{20}$ są niezależne i jednostajne na
$[0,1]$, a każde $H_i$ odrzucamy przy $p_i\le0{,}05$,
to $\mathrm{FWER}=1-(1-0{,}05)^{20}\approx0{,}6415$.
Nawet jeśli każdy pojedynczy test ma poprawny poziom,
cała procedura nie kontroluje rodzinnego błędu na $5\%$.
Ten konkretny wzór wykorzystuje \emph{niezależność}; bez
niej nie można go automatycznie stosować.
\end{przyklad}

\section{Bonferroni: prosta gwarancja bez niezależności}
\begin{twierdzenie}[Procedura Bonferroniego]\label{thm:bonferroni}
Niech $p_1,\ldots,p_m$ będą mierzalnymi wartościami p
dla ustalonej rodziny hipotez i niech dla każdego
$i\in I_0$ oraz każdego $u\in[0,1]$ zachodzi
$P(p_i\le u)\le u$. Bez żadnych założeń o wzajemnej
zależności wartości p reguła
\[
 \text{odrzuć }H_i\quad\Longleftrightarrow\quad
 p_i\le\frac\alpha m
\]
ma silną kontrolę $\mathrm{FWER}\le\alpha$.
\end{twierdzenie}
\begin{proof}
Zdarzenie $\{V\ge1\}$ jest zawarte w unii zdarzeń
$\{p_i\le\alpha/m\}$ po $i\in I_0$.
Z nierówności sumy prawdopodobieństw i ważności
poszczególnych wartości p otrzymujemy
\[
 P(V\ge1)\le\sum_{i\in I_0}P(p_i\le\alpha/m)
 \le\frac{m_0\alpha}m\le\alpha.
\]
Dowód nie wykorzystuje niezależności ani liczby fałszywych
hipotez w rodzinie.
\end{proof}

\begin{wniosek}[Jednoczesne przedziały ufności]
\label{cor:simultaneous-ci}
Niech $\theta_1,\ldots,\theta_m$ oznaczają ustalone
parametry (lub funkcjonały parametrów), a losowe
przedziały $C_i(X)$ spełniają
$P(\theta_i\notin C_i(X))\le\alpha/m$ dla każdego
$i=1,\ldots,m$. Wówczas
\[
 P\bigl(\theta_i\in C_i(X)\ \text{dla wszystkich }
 i=1,\ldots,m\bigr)\ge1-\alpha.
\]
\end{wniosek}
\begin{proof}
Zdarzenie braku jednoczesnego pokrycia jest unią zdarzeń
$\{\theta_i\notin C_i(X)\}$. Nierówność sumy
prawdopodobieństw ogranicza jego prawdopodobieństwo przez
$\sum_i\alpha/m=\alpha$. Żadna niezależność przedziałów
nie jest potrzebna.
\end{proof}

\begin{uwaga}[Cena ochrony rodziny]
Bonferroni zachowuje gwarancję przy dowolnej zależności,
ale może być konserwatywny, zwłaszcza przy wielu silnie
skorelowanych testach. Jeżeli mamy z góry wybrany
jeden \emph{główny} wynik badania i wiele wyników
pomocniczych, rozsądny plan może oddzielać rodziny
hipotez o różnych celach; nie wolno jednak tworzyć
takich rodzin dopiero po obejrzeniu wartości p, tylko
po to, by uniknąć korekty.
\end{uwaga}

\section{Procedura Holma: lepsza od Bonferroniego bez dodatkowych założeń}
\begin{definicja}[Procedura krokowa Holma]\label{def:holm}
Niech $p_{(1)}\le\cdots\le p_{(m)}$ będą uporządkowanymi
wartościami p, a $H_{(i)}$ odpowiadającymi im hipotezami;
przy remisach wybieramy z góry ustaloną regułę porządku.
Dla $0<\alpha<1$ sprawdzamy kolejno
$p_{(i)}\le\alpha/(m-i+1)$. Odrzucamy hipotezy
$H_{(1)},\ldots,H_{(r)}$ aż do pierwszej niespełnionej
nierówności, po czym zatrzymujemy procedurę.
Jeśli pierwsza nierówność nie zachodzi, nie odrzucamy
żadnej hipotezy.
\end{definicja}

\begin{twierdzenie}[Silna kontrola FWER przez Holma]
\label{thm:holm}
Załóżmy, że każda prawdziwa hipoteza zerowa w rodzinie
$H_1,\ldots,H_m$ ma ważną wartość p, to znaczy
$P(p_i\le u)\le u$ dla każdego $u\in[0,1]$ i
$i\in I_0$. Procedura Holma z definicji~\ref{def:holm}
ma silną kontrolę $\mathrm{FWER}\le\alpha$ bez względu
na zależności między wartościami p. Ponadto zawsze
odrzuca wszystkie hipotezy, które odrzuca Bonferroni
przy tym samym poziomie $\alpha$.
\end{twierdzenie}
\begin{proof}
Jeśli $m_0=0$, to $V=0$ i pierwsza teza jest oczywista.
Niech $m_0\ge1$. Jeżeli Holm odrzuci przynajmniej jedną
prawdziwą hipotezę, niech $j$ będzie jej pierwszą pozycją
w uporządkowanym ciągu. Przed pozycją $j$ znajdują się
wyłącznie hipotezy fałszywe, więc $j\le m-m_0+1$ i
$m-j+1\ge m_0$. Aby Holm odrzucił $H_{(j)}$, musiała
zajść nierówność
\[
 p_{(j)}\le\frac\alpha{m-j+1}\le\frac\alpha{m_0}.
\]
Zatem zdarzenie przynajmniej jednego fałszywego
odrzucenia implikuje, że dla pewnego $i\in I_0$ mamy
$p_i\le\alpha/m_0$. Stosując nierówność sumy do
$m_0$ prawdziwych hipotez, dostajemy $P(V\ge1)\le\alpha$.

Jeśli Bonferroni odrzuca $H_{(j)}$, to
$p_{(j)}\le\alpha/m$. Dla każdego $i\le j$
mamy $p_{(i)}\le p_{(j)}\le\alpha/m\le
\alpha/(m-i+1)$, a zatem Holm nie zatrzyma się
przed pozycją $j$ i również odrzuci $H_{(j)}$.
\end{proof}

\begin{przyklad}[Cztery porównania zaplanowane z góry]
Przypuśćmy, że $m=4$, $\alpha=0{,}05$ i obserwowane
wartości p wynoszą
$0{,}003$, $0{,}018$, $0{,}041$, $0{,}20$.
Bonferroni wymaga $p\le0{,}0125$, więc odrzuca
jedną hipotezę. Holm porównuje kolejno wartości p z
$0{,}0125$, $0{,}016\overline6$, $0{,}025$ i
$0{,}05$; zatrzymuje się na drugiej pozycji i również
odrzuca jedną hipotezę. W innej konfiguracji Holm może
odrzucić więcej hipotez, ale nigdy mniej niż Bonferroni.
\end{przyklad}

\begin{uwaga}[Korekta Šidáka]
Jeżeli \emph{wszystkie} wartości p z prawdziwych hipotez
zerowych są niezależne, to odrzucanie przy
$p_i\le1-(1-\alpha)^{1/m}$ daje silną kontrolę FWER.
Rzeczywiście, dla $m_0\le m$ niezależnych i ważnych
wartości p prawdopodobieństwo braku błędnych odrzuceń
jest co najmniej $(1-t)^{m_0}\ge(1-t)^m=1-\alpha$.
Bez odpowiednich założeń o zależności tej korekty nie
należy stosować jako bezwarunkowej gwarancji.
\end{uwaga}

\section{Benjamini--Hochberg: kontrolowanie odsetka fałszywych odkryć}
\begin{definicja}[Procedura Benjamini--Hochberg (BH)]
\label{def:bh}
Niech $p_{(1)}\le\cdots\le p_{(m)}$ i ustalmy
$q\in(0,1)$. Wyznaczamy
\[
 R=\max\left(\{0\}\cup
 \left\{i\in\{1,\ldots,m\}:
 p_{(i)}\le\frac{iq}{m}\right\}\right).
\]
Jeżeli $R=0$, nie odrzucamy żadnej hipotezy.
Jeżeli $R\ge1$, odrzucamy wszystkie hipotezy z
$p_i\le p_{(R)}$. Liczba tych odrzuceń wynosi dokładnie
$R$: jeśli $R<m$ i $p_{(R+1)}=p_{(R)}$, to również
$p_{(R+1)}\le qR/m<q(R+1)/m$, wbrew maksymalności $R$.
Dla $R=m$ odrzucamy wszystkie $m$ hipotez.
\end{definicja}

\begin{uwaga}[Dlaczego procedura jest krokowa, a nie zbiorem
niezależnych testów?]
Wartości progowe $iq/m$ rosną wraz z indeksem.
Może się zdarzyć, że nierówność nie zachodzi dla pewnego
$i<R$, a jednak zachodzi dla $R$. Procedura BH nie
przerywa badania przy pierwszym niepowodzeniu; wybiera
\emph{największy} indeks spełniający nierówność.
W przeciwieństwie do Holma nie jest procedurą
zatrzymywaną na pierwszym przekroczeniu.
\end{uwaga}

\begin{twierdzenie}[Kontrola FDR przez BH przy niezależności]
\label{thm:bh}
Niech $p_1,\ldots,p_m$ będą niezależnymi zmiennymi
losowymi o wartościach w $[0,1]$. Zakładamy, że
każde $p_i$ odpowiadające prawdziwej hipotezie zerowej
jest ważną wartością p:
$P(p_i\le u)\le u$ dla każdego $u\in[0,1]$.
Wtedy procedura BH na poziomie $q\in(0,1)$ spełnia
\[
 \mathrm{FDR}\le\frac{m_0}{m}q\le q.
\]
\end{twierdzenie}
\begin{proof}
Jeśli $m_0=0$, to FDR wynosi zero. Ustalmy
$i\in I_0$, a następnie pozostawmy wszystkie pozostałe
wartości p bez zmian i zastąpmy $p_i$ przez $0$.
Niech $R_i^0\ge1$ oznacza liczbę odrzuceń w tej
zmodyfikowanej próbie. Z własności krokowej konstrukcji
BH wynika, że na zdarzeniu, w którym oryginalna
procedura odrzuca $H_i$, zachodzi $R=R_i^0$.
Rzeczywiście, ustawienie już odrzuconej wartości p na
zero nie może zmniejszyć maksymalnego indeksu $R$;
nie może też go zwiększyć: pozostałe wartości nie
zmieniły się, a żaden indeks większy niż $R$ nie
spełniał wcześniej warunku progowego.

Warunek odrzucenia $H_i$ pociąga
$p_i\le qR_i^0/m$. Zatem, warunkując względem
wszystkich $p_j$ z $j\ne i$ i używając ich
niezależności od $p_i$, otrzymujemy
\[
 \E\left[\frac{\ind{\{H_i\ \text{odrzucona}\}}}{R\vee1}
 \ \middle|\ (p_j)_{j\ne i}\right]
 \le\frac1{R_i^0}P\left(p_i\le\frac{qR_i^0}{m}
 \ \middle|\ (p_j)_{j\ne i}\right)
 \le\frac q m.
\]
Ostatnia nierówność wynika z ważności $p_i$ i
$1\le R_i^0\le m$. Sumując po $i\in I_0$
i biorąc wartości oczekiwane, dostajemy
$\mathrm{FDR}=\sum_{i\in I_0}
\E[\ind{\{H_i\ \text{odrzucona}\}}/(R\vee1)]
\le m_0q/m$, co dowodzi twierdzenia.
\end{proof}

\begin{przyklad}[Te same cztery testy, inny cel]
Dla wartości p $0{,}003$, $0{,}018$, $0{,}041$, $0{,}20$
i $q=0{,}05$ progi BH to
$0{,}0125$, $0{,}025$, $0{,}0375$, $0{,}05$.
Największym poprawnym indeksem jest $R=2$, więc BH
odrzuca dwie hipotezy, podczas gdy Bonferroni i Holm
odrzucają jedną. Nie oznacza to, że dokładnie
$5\%$ tych dwóch odkryć jest fałszywych; kontrola
dotyczy wartości oczekiwanej FDP w powtarzanym
schemacie badania i opiera się na założeniach
twierdzenia~\ref{thm:bh}.
\end{przyklad}

\section{Zależne testy: co pozostaje prawdziwe?}
Wiele testów może korzystać z tych samych obserwacji,
na przykład z tej samej grupy kontrolnej albo z silnie
skorelowanych pomiarów. Nie wolno wtedy bez komentarza
stosować dowodu BH, który wykorzystywał niezależność
$p_i$ od pozostałych wartości p.

\begin{twierdzenie}[Korekta Benjamini--Yekutieli]\label{thm:by}
Niech $p_1,\ldots,p_m$ będą dowolnie zależnymi
wartościami p, przy czym dla każdego $i\in I_0$
i $u\in[0,1]$ zachodzi $P(p_i\le u)\le u$.
Niech $H_m=\sum_{r=1}^m1/r$. Jeśli w procedurze BH
z definicji~\ref{def:bh} zastąpimy próg $iq/m$ progiem
$iq/(mH_m)$, to otrzymana procedura spełnia
$\mathrm{FDR}\le(m_0/m)q\le q$.
\end{twierdzenie}
\begin{proof}
Niech $c=q/(mH_m)$ i dla ustalonego $i\in I_0$
rozważmy zdarzenia $A_r=\{(r-1)c<p_i\le rc\}$
dla $r=1,\ldots,m$. Jeśli hipoteza $H_i$ jest
odrzucona i $p_i\in A_r$, to z warunku krokowego
$p_i\le cR$ wynika $R\ge r$. Dla $p_i=0$
też zachodzi $R\ge1$ na zdarzeniu odrzucenia;
ponieważ rozkład ważnej wartości p spełnia
$P(p_i=0)=0$, możemy ten przypadek pominąć.
Wobec tego
\[
 \E\frac{\ind{\{H_i\ \text{odrzucona}\}}}{R\vee1}
 \le\sum_{r=1}^m\frac1rP(A_r).
\]
Niech $F_i(t)=P(p_i\le t)$, a $F_i(0)=0$.
Ważność oznacza $F_i(rc)\le rc$. Z przekształcenia
sumy przez części (współczynniki $1/r$ maleją) wynika
\[
 \sum_{r=1}^m\frac{F_i(rc)-F_i((r-1)c)}r
 \le\sum_{r=1}^m\frac{rc-(r-1)c}{r}
 =cH_m=\frac qm.
\]
Dla jasności, nierówność wynika z tożsamości
$\sum_{r=1}^m a_r(F_r-F_{r-1})=
 a_mF_m+\sum_{r=1}^{m-1}(a_r-a_{r+1})F_r$
i oszacowań $F_r\le rc$; suma podstawionych górnych
ograniczeń teleskopuje się do $c\sum_{r=1}^m1/r$.
Sumujemy teraz po $i\in I_0$, otrzymując
$\mathrm{FDR}\le m_0q/m$.
\end{proof}

Korekta BY jest bezpieczna przy dowolnej zależności,
ale może istotnie zmniejszać liczbę odkryć. Dla pewnych
rodzin zależnych testów BH również kontroluje FDR
(np. przy określonych warunkach dodatniej zależności),
lecz nie należy przypisywać tej gwarancji wszystkim
możliwym zależnościom.

\section{Po ANOVA: porównania post hoc i pytania pierwotne}
W rozdziale~\ref{ch:parametric-tests} test ANOVA badał
$H_0:\mu_1=\cdots=\mu_k$. Po jego odrzuceniu często
zadajemy bardziej szczegółowe pytania: czy A różni się od B,
czy A różni się od C, czy B różni się od C? To są
\emph{trzy oddzielne hipotezy}, a ich zbiór powinien zostać
uwzględniony w kontroli błędów.

\begin{center}\small
\begin{tabular}{>{\raggedright\arraybackslash}p{0.27\textwidth}>{\raggedright\arraybackslash}p{0.34\textwidth}>{\raggedright\arraybackslash}p{0.29\textwidth}}
\toprule
\textbf{Sytuacja} & \textbf{Metoda} & \textbf{Dlaczego / ograniczenia}\\
\midrule
Kilka zaplanowanych kontrastów & Bonferroni lub Holm & Silna kontrola FWER przy ważnych wartościach p; nie wymagają niezależności.\\
Wszystkie pary grup w normalnej ANOVA ze wspólną wariancją & Tukey HSD & Wykorzystuje wspólny rozkład maksimum porównań i zwykle jest mniej konserwatywny niż Bonferroni.\\
Wiele grup przeciw jednej kontroli & Dunnett & Korzysta z zależności porównań do wspólnej kontroli.\\
Wiele par przy różnych wariancjach & Games--Howell & Aproksymacja dla wariancji nierównych; należy ocenić liczebności i ogony.\\
Setki testów przesiewowych & BH, ewentualnie BY & Wybór FDR zamiast FWER jest decyzją o tolerowanej skali błędów.\\
\bottomrule
\end{tabular}
\end{center}

Wszystkie te metody wymagają sprawdzenia modelu
prawdopodobieństwa. Nie należy interpretować
„braku istotności pary A--B” jako dowodu równoważności;
podobnie jak w pojedynczym teście konieczne są
odrębne hipotezy równoważności i uzasadniony margines.
Pełną teorię procedur Tukeya, Dunnetta i Gamesa--Howella,
w tym ich statystyki i rozkłady, rozwiniemy w rozdziale
o rozbudowanej ANOVA i ANCOVA.

\section{Praktyczne ryzyka: eksploracja, selekcja i podglądanie danych}
\begin{przyklad}[Wybieranie najlepszego wyniku]
Analityk oblicza pięćdziesiąt korelacji i raportuje
jedynie trzy najmniejsze wartości p, nie ujawniając
pozostałych testów. Odbiorca mógłby błędnie odczytać je
jako trzy z góry zaplanowane porównania. W rzeczywistości
zbiór hipotez jest szerszy; aby ocenić wiarygodność
odkryć, potrzebujemy informacji o całej rodzinie
przeprowadzonych testów. Procedura korekty oparta tylko
na trzech wybranych wynikach nie kontroluje błędów
w pierwotnym schemacie wyboru.
\end{przyklad}

\begin{uwaga}[Sekwencyjne zaglądanie do danych]
Korekty Bonferroniego, Holma i BH rozważane tutaj
dotyczą ustalonej rodziny ważnych wartości p.
Jeśli dla \emph{tej samej} hipotezy wielokrotnie zaglądamy
do narastającej próby i kończymy badanie, gdy zwykłe
$p\le0{,}05$, to końcowa wartość p nie musi już być
ważna przy tej regule zatrzymania. Potrzebne są osobno
zaplanowane metody sekwencyjne: wydatkowanie poziomu
istotności, granice grupowo-sekwencyjne albo testy ważne
w czasie. Samo dodanie korekty na liczbę raportowanych
hipotez nie naprawia takiego problemu.
\end{uwaga}

\section{Dobór procedury w praktyce}
\begin{center}\small
\begin{tabular}{>{\raggedright\arraybackslash}p{0.31\textwidth}>{\raggedright\arraybackslash}p{0.28\textwidth}>{\raggedright\arraybackslash}p{0.29\textwidth}}
\toprule
\textbf{Pytanie} & \textbf{Zwykle rozważamy} & \textbf{Na co uważać?}\\
\midrule
Czy choć jeden fałszywy alarm byłby kosztowny? & FWER: Holm, Bonferroni & Silna kontrola całej rodziny.\\
Czy wykonujemy liczne badania przesiewowe? & FDR: BH & Niezależność lub odpowiednie warunki zależności.\\
Czy zależności p-value są nieznane? & Holm dla FWER, BY dla FDR & BY może mieć małą moc.\\
Czy porównujemy wszystkie pary po ANOVA? & Tukey lub Games--Howell & Wspólna wariancja lub jej brak, liczebności.\\
Czy wszystkie grupy porównujemy z kontrolą? & Dunnett & Wspólny model i współzależność kontrastów.\\
Czy obserwowaliśmy dane wielokrotnie w czasie? & Procedura sekwencyjna & Predefiniowany plan, ważność przy zatrzymaniu.\\
Czy raportujemy wiele oszacowań? & Jednoczesne przedziały & Własności pokrycia całej rodziny.\\
\bottomrule
\end{tabular}
\end{center}

\begin{lstlisting}[language=Python,caption={Kompletny program w Pythonie: Bonferroni, Holm, BH i BY bez zewnętrznych pakietów.}]
from math import fsum

pvals = [0.003, 0.018, 0.041, 0.20]
alpha = 0.05
m = len(pvals)
order = sorted(range(m), key=lambda i: pvals[i])

bonf = [p <= alpha/m for p in pvals]
holm = [False] * m
for rank, i in enumerate(order):
    if pvals[i] > alpha/(m-rank):
        break
    holm[i] = True

def step_up(pvals, q):
    m = len(pvals)
    order = sorted(range(m), key=lambda i: pvals[i])
    r = max((j for j in range(1, m+1)
             if pvals[order[j-1]] <= q*j/m), default=0)
    return [p <= pvals[order[r-1]] for p in pvals] if r else [False]*m

harmonic = fsum(1/j for j in range(1, m+1))
bh = step_up(pvals, alpha)
by = step_up(pvals, alpha/harmonic)
print('Bonferroni:', bonf)
print('Holm:', holm)
print('BH:', bh)
print('BY:', by)
\end{lstlisting}

\begin{lstlisting}[language=Julia,caption={Obliczenie krokowych korekt Holma i BH w Julii bez pakietów zewnętrznych.}]
pvals = [0.003, 0.018, 0.041, 0.20]
alpha = 0.05
m = length(pvals)
idx = sortperm(pvals)

holm = falses(m)
for j in 1:m
    i = idx[j]
    pvals[i] > alpha/(m-j+1) && break
    holm[i] = true
end

threshold = [alpha*j/m for j in 1:m]
valid = [j for j in 1:m if pvals[idx[j]] <= threshold[j]]
r = isempty(valid) ? 0 : maximum(valid)
bh = falses(m)
if r > 0
    bh .= pvals .<= pvals[idx[r]]
end
println((holm=holm, bh=bh))
\end{lstlisting}

\paragraph{Jak czytać wyniki programu?}
Obie procedury uruchomione na tych samych danych mogą
zwrócić różne decyzje, mimo że każda jest poprawna
\emph{względem swojego celu i założeń}. Nie należy
wybierać między FWER a FDR dopiero na podstawie
tego, która procedura odrzuciła więcej hipotez.

\section{Zadanie podsumowujące: rekomendacje dla analityka}
\begin{zadanie}[Cztery lokalizacje, trzy rodzaje pytań]
Firma porównuje cztery lokalizacje pod kątem wskaźników
awaryjności. Rozważa trzy cele: (a) potwierdzić, że
\emph{każda} wskazana lokalizacja spełnia normę;
(b) znaleźć możliwie wiele lokalizacji wymagających
interwencji; (c) porównać każdą lokalizację z jedną
wspólną lokalizacją kontrolną. Jakie metody rozważyć?
\par\emph{Rozwiązanie.} W (a) wybieramy procedurę
kontrolującą rodzinny błąd zgodnie z kierunkiem hipotez,
np. Holm, oraz odpowiednio jednostronne przedziały
jednoczesne. W (b), jeśli akceptujemy kontrolę
oczekiwanego udziału błędnych sygnałów i rozumiemy
koszty pomyłek, możemy rozważyć BH lub BY zależnie od
struktury zależności. W (c) porównania do wspólnej
kontroli sugerują procedurę Dunnetta \emph{w odpowiednim
modelu normalnym}; przy niezależnych prawidłowych
wartościach p albo bez szczególnego modelu korelacji
można użyć Holma. W każdym przypadku od wyboru
procedury ważniejsza jest poprawność jednostki analizy,
założeń modelu i interpretacji wyniku.
\end{zadanie}



% ================================================================
% ROZDZIAL 13 -- REGRESJA LINIOWA
% ================================================================
\chapter{Regresja liniowa: estymacja, diagnostyka i predykcja}
\label{ch:linear-regression}
\paragraph{Powiązania i odsyłacze.} Metody predykcji i doboru złożoności rozwija rozdział~\ref{ch:model-selection}, a teorię testów~\ref{ch:parametric-tests}.


\section{Kiedy statystyka średniej przestaje wystarczać?}
Porównując trzy technologie, w rozdziale~\ref{ch:parametric-tests}
analizowaliśmy różnice średnich między grupami. Co jednak zrobić, gdy
sprawność urządzenia zależy także od temperatury, czasu pracy i wieku
maszyny? Potrzebujemy modelu opisującego \emph{warunkową średnią}
wyniku przy ustalonych cechach. Regresja liniowa jest najprostszą
konstrukcją o tej własności. Jej nazwa odnosi się do liniowości
\emph{współczynników}, niekoniecznie do liniowości wszystkich funkcji
wykorzystanych jako zmienne objaśniające: można używać $x^2$ lub
interakcji $xz$, a nadal otrzymać model liniowy względem parametrów.

\begin{definicja}[Model liniowy przy ustalonym planie doświadczenia]
\label{def:linear-fixed}
Niech $n,p\in\N$, $n>p\ge1$, a
$X\in\R^{n\times p}$ będzie znaną macierzą o rzędzie $p$.
Niech $\beta\in\R^p$, zaś $\varepsilon\in\R^n$ będzie
losowym wektorem o $\E\varepsilon=0$ i
$\Cov(\varepsilon)=\sigma^2I_n$, gdzie $\sigma^2>0$.
W \emph{liniowym modelu homoskedastycznym} obserwujemy
$Y=X\beta+\varepsilon$. Jeśli dodatkowo
$\varepsilon\sim\mathcal N_n(0,\sigma^2I_n)$, mówimy o
\emph{normalnym modelu liniowym}. Macierz $X$ to
\emph{macierz planu}, a $\beta$ to wektor współczynników.
Wiersz $x_i^\top$ macierzy $X$ opisuje cechy $i$-tej jednostki.
\end{definicja}

Przy losowych cechach $X$ wszystkie poniższe twierdzenia można
stosować \emph{warunkowo względem $X$}, o ile odpowiednie założenia
spełnione są prawie na pewno. Sam fakt losowości cech nie unieważnia
regresji. Potrzebujemy jednak warunkowych założeń
$\E(\varepsilon\mid X)=0$ i
$\Cov(\varepsilon\mid X)=\sigma^2I_n$,
aby bez zmian stosować twierdzenia o nieobciążoności
i wariancji warunkowej. Same warunki bezwarunkowe
$\E\varepsilon=0$ i $\Cov(\varepsilon)=\sigma^2I_n$
nie wystarczają. Przykładowo dla skalarnych
$X\sim\Normal(0,1)$ i $\varepsilon=X$ mamy
$\E\varepsilon=0$, ale $\E(\varepsilon\mid X)=X$.
Problem ten występuje m.in. przy pominięciu
istotnej zmiennej zakłócającej.

\begin{przyklad}[Co oznacza współczynnik w modelu?]
Jeśli $\E(Y_i\mid X)=\beta_0+\beta_1t_i+\beta_2g_i$,
gdzie $t_i$ to temperatura, a $g_i\in\{0,1\}$ oznacza
technologię B, to $\beta_1$ jest zmianą średniej odpowiedzi
na wzrost temperatury o jedną jednostkę \emph{przy tej samej
technologii}, a $\beta_2$ jest różnicą średnich B--A
\emph{przy tej samej temperaturze}. Gdy znaczenie temperatury
może zależeć od technologii, sprawdzamy rozszerzenie
$\beta_3t_ig_i$. Wtedy różnica B--A to
$\beta_2+\beta_3t_i$ i nie wolno przypisywać
$\beta_2$ znaczenia stałej różnicy między grupami.
To zapowiada analizę kowariancji (ANCOVA).
\end{przyklad}

\section{Metoda najmniejszych kwadratów jako rzut ortogonalny}
\begin{twierdzenie}[Istnienie, jednoznaczność i postać MNK]
\label{thm:ols-formula}
Załóżmy, że $X\in\R^{n\times p}$ ma rząd $p\le n$.
Dla dowolnego $y\in\R^n$ funkcja
$Q_y(b)=\|y-Xb\|_2^2$ ma dokładnie jedno minimum
na $\R^p$, osiągane w punkcie
\[
 \widehat\beta=(X^\top X)^{-1}X^\top y.
\]
Przy $H=X(X^\top X)^{-1}X^\top$ wektor
$\widehat y=Hy$ jest rzutem ortogonalnym $y$ na
$\operatorname{im}X$, a reszta $e=(I_n-H)y$ jest
prostopadła do każdej kolumny $X$.
\end{twierdzenie}
\begin{proof}
Macierz $X^\top X$ jest dodatnio określona, ponieważ
dla $v\ne0$ zachodzi
$v^\top X^\top Xv=\|Xv\|_2^2>0$.
Zdefiniujmy $b_*=(X^\top X)^{-1}X^\top y$.
Wtedy $X^\top(y-Xb_*)=0$, więc $y-Xb_*$ jest
prostopadły do wszystkich wektorów postaci $Xv$.
Dla dowolnego $b$ z twierdzenia Pitagorasa otrzymujemy
$\|y-Xb\|_2^2=\|y-Xb_*\|_2^2+\|X(b-b_*)\|_2^2$.
Drugi składnik jest zerowy tylko dla $b=b_*$.
Ponadto $H^\top=H$ i $H^2=H$, zatem $H$ jest
rzutem ortogonalnym na $\operatorname{im}X$.
\end{proof}

\begin{uwaga}[Dlaczego nie obliczamy odwrotności w programie?]
Wzór teoretyczny zawiera $(X^\top X)^{-1}$, ale
jawne wyznaczanie tej odwrotności jest niepotrzebne i
bywa niestabilne numerycznie, ponieważ liczba uwarunkowania
$X^\top X$ jest kwadratem liczby uwarunkowania $X$.
W praktyce stosujemy rozkład QR lub SVD (rozkład według
wartości osobliwych). W przykładzie programistycznym
użyjemy \texttt{numpy.linalg.lstsq}. Analogicznie,
przy regularizacji ridge w rozdziale~\ref{ch:model-selection}
nie będziemy formować odwrotności macierzy wprost.
\end{uwaga}

\section{Dlaczego współczynniki MNK bywają optymalne?}
\begin{twierdzenie}[Nieobciążoność i macierz kowariancji]
\label{thm:ols-cov}
W modelu~\ref{def:linear-fixed} przy pełnym rzędzie $X$
estymator $\widehat\beta=(X^\top X)^{-1}X^\top Y$
spełnia
\[
 \E\widehat\beta=\beta,\qquad
 \Cov(\widehat\beta)=\sigma^2(X^\top X)^{-1}.
\]
\end{twierdzenie}
\begin{proof}
Z tożsamości $\widehat\beta-\beta=A\varepsilon$ dla
$A=(X^\top X)^{-1}X^\top$ wynika
$\E\widehat\beta-\beta=A\E\varepsilon=0$.
Ponadto $\Cov(A\varepsilon)=A(\sigma^2I_n)A^\top
=\sigma^2(X^\top X)^{-1}$, ponieważ
$AA^\top=(X^\top X)^{-1}X^\top X(X^\top X)^{-1}$.
\end{proof}

\begin{twierdzenie}[Gauss--Markow: najlepszy liniowy nieobciążony]
\label{thm:gauss-markov}
W założeniach twierdzenia~\ref{thm:ols-cov}, niech
$\widetilde\beta=CY$ dla ustalonej macierzy
$C\in\R^{p\times n}$ takiej, że $CX=I_p$.
Wtedy różnica
$\Cov(\widetilde\beta)-\Cov(\widehat\beta)$
jest nieujemnie określona. Zatem MNK ma najmniejszą
macierz kowariancji w porządku Loewnera \emph{w klasie
liniowych nieobciążonych estymatorów}.
Do twierdzenia nie jest potrzebna normalność błędów.
\end{twierdzenie}
\begin{proof}
Niech $A=(X^\top X)^{-1}X^\top$ i $D=C-A$.
Z równości $CX=AX=I_p$ wynika $DX=0$.
Ponieważ $AD^\top=(X^\top X)^{-1}X^\top D^\top
=(X^\top X)^{-1}(DX)^\top=0$, mamy
$CC^\top=AA^\top+DD^\top$.
W konsekwencji
$\Cov(CY)-\Cov(AY)=\sigma^2DD^\top$,
a dla każdego $v\in\R^p$ zachodzi
$v^\top DD^\top v=\|D^\top v\|_2^2\ge0$.
\end{proof}

\begin{uwaga}[Czego twierdzenie Gaussa--Markowa nie dowodzi?]
Nie oznacza ono, że MNK jest najlepszy spośród wszystkich
estymatorów, odporny na dowolne odstające obserwacje ani że
współczynniki mają rozkład normalny. To ostatnie wymaga
założenia normalności (lub odpowiedniej asymptotyki).
Jeśli $\Cov(\varepsilon)=\Omega\ne\sigma^2I_n$,
to klasyczna minimalność MNK nie wynika z twierdzenia;
po właściwym opisaniu $\Omega$ można rozważyć ważone
najmniejsze kwadraty.
\end{uwaga}

\section{Reszty, wariancja i dokładne rozkłady w modelu normalnym}
\begin{twierdzenie}[Rozkład współczynników i sumy kwadratów reszt]
\label{thm:ols-gaussian}
Niech $n>p\ge1$, $X$ ma pełny rząd kolumnowy i
$Y\sim\mathcal N_n(X\beta,\sigma^2I_n)$, $\sigma>0$.
Niech $e=(I_n-H)Y$, $\mathrm{RSS}=\|e\|_2^2$ oraz
$s^2=\mathrm{RSS}/(n-p)$. Wtedy:
\begin{enumerate}
\item $\widehat\beta\sim\mathcal N_p\!\bigl(\beta,
 \sigma^2(X^\top X)^{-1}\bigr)$;
\item $\mathrm{RSS}/\sigma^2\sim\chi^2_{n-p}$;
\item $\widehat\beta$ i $\mathrm{RSS}$ są niezależne,
 a $\E(s^2)=\sigma^2$.
\end{enumerate}
\end{twierdzenie}
\begin{proof}
Estymator jest liniowym przekształceniem wektora normalnego,
więc pierwszy punkt wynika z obliczonej już macierzy kowariancji.
Macierz $H$ jest rzutem ortogonalnym rzędu $p$, zatem
$I_n-H$ jest rzutem ortogonalnym rzędu $n-p$.
Niech $U$ będzie macierzą ortogonalną, która diagonalizuje
$I_n-H$ do $\operatorname{diag}(I_{n-p},0_p)$.
Wektor $U^\top\varepsilon/\sigma$ ma niezależne
współrzędne $\mathcal N(0,1)$. Ponieważ
$(I_n-H)X\beta=0$, otrzymujemy
$\mathrm{RSS}/\sigma^2=\sum_{j=1}^{n-p}Z_j^2$.
Niezależność wynika z ortogonalności rzutów:
$\Cov(H\varepsilon,(I_n-H)\varepsilon)=0$,
a wspólna normalność oznacza niezależność wektorów.
Estymator zależy jedynie od $HY$, zaś RSS jedynie od
$(I_n-H)Y$. Wreszcie średnia rozkładu
$\chi^2_{n-p}$ wynosi $n-p$.
\end{proof}

\begin{wniosek}[Dokładny przedział dla współczynnika]
\label{cor:ols-ci}
W założeniach twierdzenia~\ref{thm:ols-gaussian}
niech $c\in\R^p\setminus\{0\}$ będzie ustalonym
wektorem. Wtedy
\[
 T_c=\frac{c^\top\widehat\beta-c^\top\beta}
 {s\sqrt{c^\top(X^\top X)^{-1}c}}\sim t_{n-p}.
\]
W rezultacie przedział
$c^\top\widehat\beta\pm t_{n-p,1-\alpha/2}
 s\sqrt{c^\top(X^\top X)^{-1}c}$
ma pokrycie $1-\alpha$, dla $0<\alpha<1$.
\end{wniosek}
\begin{proof}
Licznik podzielony przez
$\sigma\sqrt{c^\top(X^\top X)^{-1}c}$ ma rozkład
$\mathcal N(0,1)$; zgodnie z twierdzeniem
\ref{thm:ols-gaussian} jest niezależny od
$(n-p)s^2/\sigma^2\sim\chi^2_{n-p}$.
Ich iloraz ma rozkład $t_{n-p}$.
\end{proof}

\begin{przyklad}[Dlaczego zmiana jednostek nie zmienia istotności?]
Jeżeli temperatura jest mierzona w stopniach Celsjusza,
a następnie zastąpimy ją różnicą skali, np. podzielimy
przez $10$, odpowiedni współczynnik zostanie pomnożony
przez $10$, podobnie jak jego błąd standardowy.
Statystyka t dla testu zerowego współczynnika pozostanie
niezmieniona. Interpretując efekt, zawsze podajemy
\emph{jednostkę} zmiennej objaśniającej.
\end{przyklad}

\section{Predykcja przyszłej obserwacji a szacowanie średniej}
\begin{twierdzenie}[Przedział ufności i przedział predykcji]
\label{thm:prediction-vs-mean}
Przy założeniach twierdzenia~\ref{thm:ols-gaussian}
ustalmy niezerowy wektor cech $x_0\in\R^p\setminus\{0\}$ i niech
$Y_0=x_0^\top\beta+\varepsilon_0$, gdzie
$\varepsilon_0\sim\Normal(0,\sigma^2)$ jest niezależna
od $Y$. Niech $h_0=x_0^\top(X^\top X)^{-1}x_0$.
Wówczas przedział ufności dla średniej warunkowej
$\mu_0=x_0^\top\beta$ wynosi
$x_0^\top\widehat\beta\pm t_{n-p,1-\alpha/2}s\sqrt{h_0}$,
a przedział predykcji dla nowego $Y_0$ ma postać
\[
 x_0^\top\widehat\beta\pm
 t_{n-p,1-\alpha/2}s\sqrt{1+h_0}.
\]
Każdy z przedziałów ma odpowiednio pokrycie
$1-\alpha$ dla własnego celu; drugi dotyczy \emph{losowej
nowej obserwacji}, nie stałej wartości parametru.
\end{twierdzenie}
\begin{proof}
Z poprzedniego twierdzenia
$x_0^\top(\widehat\beta-\beta)$ jest normalny o
wariancji $\sigma^2h_0$. Dla nowego wyniku błąd
$Y_0-x_0^\top\widehat\beta$ jest sumą niezależnych
składników $\varepsilon_0$ i
$-x_0^\top(\widehat\beta-\beta)$, ma więc wariancję
$\sigma^2(1+h_0)$ i rozkład normalny.
Oba błędy są niezależne od $s^2$, bo $\varepsilon_0$
jest niezależne od próby, a estymator i RSS są niezależne.
Standaryzacja przez $s$ prowadzi do rozkładu t.
\end{proof}

\begin{uwaga}[Ekstrapolacja]
Gdy $x_0$ jest daleko od zebranych obserwacji, wartość
$h_0$ może być duża; dodatkowo liniowa zależność może
w tym obszarze nie obowiązywać. Nawet węższy przedział
ufności dla średniej nie gwarantuje prawidłowości modelu
poza zakresem danych. \emph{Ekstrapolacja} wymaga
uzasadnienia dziedzinowego, nie tylko małej wartości p.
\end{uwaga}

\section{Ogólny test F i dlaczego ANOVA jest regresją}
\begin{twierdzenie}[Test zagnieżdżonych modeli normalnych]
\label{thm:partial-f-test}
Niech $Y\sim\mathcal N_n(\mu,\sigma^2I_n)$,
$\sigma>0$, a $V_0\subsetneq V_1$ będą ustalonymi
podprzestrzeniami $\R^n$ wymiarów odpowiednio
$p_0<p_1<n$. Oznaczmy ich rzuty ortogonalne przez
$H_0,H_1$ i resztowe sumy kwadratów przez
$\mathrm{RSS}_j=\|(I_n-H_j)Y\|_2^2$.
Przy hipotezie zerowej $\mu\in V_0$ statystyka
\[
 F=\frac{(\mathrm{RSS}_0-\mathrm{RSS}_1)/(p_1-p_0)}
 {\mathrm{RSS}_1/(n-p_1)}\sim F_{p_1-p_0,n-p_1}.
\]
\end{twierdzenie}
\begin{proof}
Ponieważ $V_0\subset V_1$, macierz $H_1-H_0$ jest
rzutem ortogonalnym na $V_1\cap V_0^\perp$ rzędu
$p_1-p_0$. Macierz $I_n-H_1$ jest rzutem na
$V_1^\perp$ rzędu $n-p_1$. Są to rzuty ortogonalne
o iloczynie zero; wektor średniej $\mu\in V_0$ leży
w jądrach obu. Z twierdzenia o ortogonalnych składowych
wektora normalnego kwadraty norm przeskalowane przez
$\sigma^2$ są niezależnymi zmiennymi o rozkładach
$\chi^2_{p_1-p_0}$ i $\chi^2_{n-p_1}$.
Ponieważ $\mathrm{RSS}_0-\mathrm{RSS}_1
=\|(H_1-H_0)Y\|_2^2$, ich iloraz ma podany rozkład.
\end{proof}

\begin{przyklad}[Jednoczynnikowa ANOVA jako test F regresji]
Niech $G_i\in\{1,\ldots,k\}$ oznacza grupę obserwacji
i załóżmy, że w każdej grupie jest przynajmniej jedna
jednostka. W modelu ograniczonym $V_0$ wszystkie grupy
mają wspólną średnią; w pełnym $V_1$ każda grupa ma
osobną średnią. Wymiary tych przestrzeni wynoszą
odpowiednio $1$ i $k$. Resztowe sumy kwadratów to
$\mathrm{RSS}_0=\sum_i(Y_i-\bar Y)^2$ i
$\mathrm{RSS}_1=\sum_g\sum_{i:G_i=g}(Y_i-\bar Y_g)^2$.
Różnica jest sumą kwadratów między grupami,
$\sum_gn_g(\bar Y_g-\bar Y)^2$.
Podstawienie do twierdzenia~\ref{thm:partial-f-test}
odtwarza statystykę ANOVA z rozdziału~\ref{ch:parametric-tests}.
\end{przyklad}

\begin{uwaga}[ANCOVA jest kolejnym krokiem]
Jeśli do macierzy planu dodamy kolumnę temperatury,
test ogólnej różnicy grup kontrolując temperaturę jest
porównaniem zagnieżdżonych przestrzeni kolumn. Warto
sprawdzić interakcję grupy z temperaturą przed przyjęciem
jednego wspólnego nachylenia. Nie należy utożsamiać
istotności współczynnika z przyczynowym efektem
zmiany technologii bez argumentu o randomizacji
lub identyfikacji przyczynowej.
\end{uwaga}

\section{Diagnostyka modelu: co może zepsuć poprawne wzory?}
\begin{definicja}[Dźwignia obserwacji i reszta studencka]
Dla macierzy planu $X$ o pełnym rzędzie i rzutowania
$H=X(X^\top X)^{-1}X^\top$ liczbę $h_{ii}=H_{ii}$
nazywamy \emph{dźwignią} $i$-tej obserwacji.
Przy $h_{ii}<1$ i oszacowaniu $s>0$ definiujemy
\emph{wewnętrznie standaryzowaną resztę}
$r_i=e_i/(s\sqrt{1-h_{ii}})$.
\end{definicja}

\begin{stwierdzenie}[Dlaczego obserwacje skrajne mogą mieć małe reszty?]
W modelu $Y=X\beta+\varepsilon$ z
$\Cov(\varepsilon)=\sigma^2I_n$ reszta $e=(I_n-H)Y$
spełnia $\E e=0$ i
$\Cov(e)=\sigma^2(I_n-H)$.
W szczególności $\Var(e_i)=\sigma^2(1-h_{ii})$,
a $0\le h_{ii}\le1$ oraz $\sum_ih_{ii}=p$.
\end{stwierdzenie}
\begin{proof}
Równość $(I_n-H)X=0$ zapewnia
$e=(I_n-H)\varepsilon$. Korzystając z
$(I_n-H)^2=I_n-H=(I_n-H)^\top$,
otrzymujemy macierz kowariancji. Z własności rzutu
$h_{ii}=e_i^\top He_i=\|He_i\|_2^2\in[0,1]$,
a jego ślad jest równy rzędowi $p$.
\end{proof}

\paragraph{Pięć pytań diagnostycznych.}
\begin{enumerate}
\item \emph{Czy średnia jest właściwie opisana?}
Rysujemy reszty względem wartości przewidywanych i cech,
sprawdzając krzywiznę i interakcje. To kontrola poprawności
postaci modelu, a nie tylko test pojedynczego współczynnika.
\item \emph{Czy wariancja jest stała?} Jeśli rozrzut rośnie
wraz z poziomem sygnału, klasyczne błędy standardowe mogą być
niewłaściwe. Rozważamy ważone MNK, odporne oszacowania
kowariancji lub inny model rozkładu.
\item \emph{Czy obserwacje są niezależne?} Przy powtarzanych
pomiarach osoby lub silnej autokorelacji potrzebna może być
kowariancja klastrowa albo jawny model zależności.
\item \emph{Czy istnieją punkty wpływowe?} Analizujemy
\emph{dźwignie}, reszty i zmianę współczynników po usunięciu
obserwacji; sam duży $|r_i|$ nie wyczerpuje diagnostyki.
\item \emph{Czy kolumny planu są rozróżnialne?}
Silna współliniowość może znacznie powiększyć
$\operatorname{diag}((X^\top X)^{-1})$, a przy dokładnej
zależności liniowej $\widehat\beta$ nie jest jednoznaczny.
\end{enumerate}

\begin{stwierdzenie}[Korekta odporna na heteroskedastyczność]
\label{prop:hetero-ols}
Jeśli $Y=X\beta+\varepsilon$, $\E\varepsilon=0$ oraz
$\Cov(\varepsilon)=\Omega$ dla znanej nieujemnie
określonej macierzy $\Omega$, to
\[
 \Cov(\widehat\beta)=(X^\top X)^{-1}X^\top
 \Omega X(X^\top X)^{-1}.
\]
Gdy błędy są nieskorelowane, lecz mają różne wariancje,
można \emph{oszacowywać} $\Omega$ diagonalnie poprzez
$e_i^2/(1-h_{ii})^2$ dla $h_{ii}<1$. Otrzymana
macierz kanapkowa HC3 jest narzędziem asymptotycznym;
nie daje automatycznie dokładnych małopróbkowych testów t.
\end{stwierdzenie}
\begin{proof}
Wciąż $\widehat\beta-\beta
=(X^\top X)^{-1}X^\top\varepsilon$.
Wzór na kowariancję liniowego przekształcenia daje tezę.
Zastąpienie nieznanych elementów $\Omega$ oszacowaniami
nie jest tożsamością skończonopróbkową i wymaga osobnej
analizy asymptotycznej.
\end{proof}

\begin{przyklad}[Pominięta zmienna i pozorna istotność]
Niech prawdziwa zależność warunkowej średniej ma postać
$\E(Y\mid x,z)=\beta_0+\beta_1x+\gamma z$.
Jeżeli regresujemy $Y$ wyłącznie na $1$ i $x$,
to średnia estymowanego nachylenia, przy ustalonych
wektorach $x,z$ i błędach o średniej zero, wynosi
\[
 \E\widehat b_1=\beta_1+
 \gamma\frac{\sum_i(x_i-\bar x)(z_i-\bar z)}
                   {\sum_i(x_i-\bar x)^2}.
\]
Współczynnik może być istotnie różny od zera nawet
gdy $\beta_1=0$, jeśli $z$ jest związane z $x$.
Sama większa próba nie usuwa tego rodzaju obciążenia.
\end{przyklad}

\section{Studium praktyczne: jak czytać wynik regresji?}
Rozważmy dane z pomiarów procesu technologicznego:
czas działania w godzinach $t_i=1,\ldots,8$ i wyniki
$y=(4{,}8,5{,}1,6{,}0,6{,}1,7{,}2,7{,}4,8{,}1,8{,}9)$.
Modelujemy $\E(Y_i\mid t_i)=\beta_0+\beta_1t_i$.
Przed interpretacją parametru $\beta_1$ pytamy, czy
każdy pomiar dotyczy innej, niezależnej partii, czy tej
samej maszyny śledzonej w czasie. W drugim przypadku
standardowa analiza niezależnych błędów może nie być
poprawna, choć dopasowana prosta będzie wyglądała tak samo.

\begin{lstlisting}[language=Python,caption={Kompletny program: MNK, przedział dla nachylenia, diagnostyka dźwigni i predykcja (NumPy, SciPy).}]
import numpy as np
from scipy.stats import t as student

t = np.arange(1.0, 9.0)
y = np.array([4.8, 5.1, 6.0, 6.1, 7.2, 7.4, 8.1, 8.9])
X = np.column_stack((np.ones_like(t), t))
n, p = X.shape

beta, _, rank, _ = np.linalg.lstsq(X, y, rcond=None)
assert rank == p
residual = y - X @ beta
rss = residual @ residual
s2 = rss / (n-p)
xtx_inv = np.linalg.solve(X.T @ X, np.eye(p))
se = np.sqrt(np.diag(s2 * xtx_inv))
quant = student.ppf(0.975, n-p)
ci_slope = (beta[1] - quant*se[1], beta[1] + quant*se[1])

H_diag = np.sum((X @ xtx_inv) * X, axis=1)
standardized_residual = residual / np.sqrt(s2 * (1-H_diag))
x0 = np.array([1.0, 9.0])
mean_prediction = x0 @ beta
h0 = x0 @ xtx_inv @ x0
pi = (mean_prediction - quant*np.sqrt(s2*(1+h0)),
      mean_prediction + quant*np.sqrt(s2*(1+h0)))
print('coefficients:', beta)
print('95% CI slope:', ci_slope)
print('leverages:', H_diag)
print('standardized residuals:', standardized_residual)
print('prediction interval at t=9:', pi)
\end{lstlisting}

\begin{lstlisting}[language=Julia,caption={MNK poprzez rozkład QR i błąd standardowy w Julii (standardowe LinearAlgebra, Statistics).}]
using LinearAlgebra, Statistics

t = collect(1.0:8.0)
y = [4.8, 5.1, 6.0, 6.1, 7.2, 7.4, 8.1, 8.9]
X = hcat(ones(length(t)), t)
beta = X \ y  # stabilniejsze niz inv(X'X) * X'y
r = y - X * beta
n, p = size(X)
s2 = sum(abs2, r) / (n - p)
xtx_inverse = (X' * X) \ Matrix{Float64}(I, p, p)
se = sqrt.(diag(s2 * xtx_inverse))
println((beta=beta, se=se, rss=sum(abs2, r)))
\end{lstlisting}

\paragraph{Interpretacja programu.}
Dla podanych danych program otrzymuje \(\widehat\beta_0\approx4{,}075\)
i \(\widehat\beta_1\approx0{,}5833\), a przedział ufności
$95\%$ dla nachylenia to w przybliżeniu
$[0{,}5035;\ 0{,}6632]$. W modelu zakładającym
niezależne normalne błędy przedział predykcji nowego
wyniku dla $t_0=9$ wynosi około $[8{,}669;\ 9{,}981]$.
Są to \emph{współczynniki modelu warunkowej średniej},
nie gotowe twierdzenia przyczynowe. Przedział predykcji zawiera składnik
wariancji nowej obserwacji i jest zwykle szerszy
niż przedział ufności dla średniej. Jeżeli reszty
wykazują autokorelację lub heteroskedastyczność,
nie należy raportować klasycznych wartości p bez
ponownego oszacowania niepewności.

\section{Wybór metody i kontrola jakości analizy}
\begin{center}\small
\begin{tabular}{>{\raggedright\arraybackslash}p{.29\textwidth}>{\raggedright\arraybackslash}p{.31\textwidth}>{\raggedright\arraybackslash}p{.29\textwidth}}
\toprule
\textbf{Sytuacja} & \textbf{Rozważana metoda} & \textbf{Ograniczenia}\\
\midrule
Wynik ciągły, średnia liniowa, błędy o stałej wariancji & Zwykłe MNK; przy normalności dokładne t i F & Zależność, nieliniowość średniej, punkty wpływowe.\\
Zmienna wynikowa ciągła, niejednakowe wariancje & MNK z odpornym SE lub WLS & WLS wymaga uzasadnienia wag; HC3 jest przybliżeniem.\\
Silna zależność w czasie albo grupowanie pomiarów & Model korelacji, GLS, modele mieszane lub odporny SE klastrowy & Liczba niezależnych jednostek i specyfikacja zależności.\\
Wpływ skrajnych wyników & Regresja odporna, kontrola wpływu, analiza wrażliwości & Nie usuwać punktów wyłącznie na podstawie p-value.\\
Dane binarne lub zliczenia & GLM z odpowiednią rodziną & Przedział wartości średniej, wariancja zależna od średniej.\\
\bottomrule
\end{tabular}
\end{center}

\begin{zadanie}[Czy dwa modele mówią to samo?]
W badaniu stwierdzono, że współczynnik temperatury jest
istotny w regresji bez zmiennej opisującej wiek maszyny,
a nieistotny po jej dodaniu. Czy to dowód, że temperatura
nie wpływa na wynik?
\par\emph{Rozwiązanie.} Nie. W pierwszym modelu współczynnik
może odzwierciedlać również skorelowany z temperaturą wiek.
W drugim rośnie liczba współczynników i może wzrosnąć
błąd standardowy. Właściwe jest porównanie założeń,
przedziałów ufności, rozkładu cech oraz celu analizy,
a przy twierdzeniach przyczynowych także projektu badania.
Sama zmiana przekroczenia progu $p=0{,}05$ nie dowodzi
braku efektu.
\end{zadanie}

% ================================================================
% ROZDZIAL 14 -- GLM
% ================================================================
\chapter{Uogólnione modele liniowe: regresja logistyczna i Poissona}
\label{ch:glm}
\paragraph{Powiązania i odsyłacze.} Bazujemy na wiarygodności z rozdziału~\ref{ch:likelihood}; model liniowy i jego diagnostykę wprowadziliśmy w rozdziale~\ref{ch:linear-regression}.


\section{Dlaczego liniowy model średniej nie wystarcza dla każdego wyniku?}
Prawdopodobieństwo wady należy do $[0,1]$, a liczba
awarii jest nieujemną liczbą całkowitą. Regresja
liniowa z dowolnymi współczynnikami może przewidywać
$p<0$ lub ujemną liczbę awarii. Ponadto wariancja
zmiennej Bernoulliego wynosi $p(1-p)$, a Poissona
$\lambda$: zmienia się z warunkową średnią.
\emph{Uogólnione modele liniowe} (GLM) rozdzielają
trzy elementy: rozkład odpowiedzi warunkowo względem
cech, model średniej oraz \emph{funkcję łączącą}
pozwalającą opisać tę średnią za pomocą kombinacji
liniowej cech. Pozornie podobne dane wymagają więc
różnych modeli w zależności od ich mechanizmu powstawania.

\begin{definicja}[Rodzina wykładnicza z parametrem dyspersji]
\label{def:glm-family}
Niech $\Theta_0\subset\R$ będzie otwartym przedziałem,
$\Phi\subset(0,\infty)$ oraz $\nu$ miarą na przestrzeni
obserwacji $\mathsf Y\subset\R$. Rozważmy rodzinę
rozkładów z gęstością lub funkcją masy względem $\nu$:
\[
 f(y;\theta,\phi)=\exp\!\left\{
 \frac{y\theta-b(\theta)}{a(\phi)}+c(y,\phi)\right\},
 \quad \theta\in\Theta_0,\ \phi\in\Phi,
\]
gdzie $a(\phi)>0$, a nośnik względem $\nu$ nie zależy
od $\theta$. Zakładamy, że gęstości całkują się do $1$
oraz że wolno dwukrotnie różniczkować pod całką względem
$\theta$. Parametr $\theta$ nazywamy \emph{kanonicznym},
zaś $a(\phi)$ jest skalą dyspersji.
\end{definicja}

\begin{twierdzenie}[Średnia i wariancja rodziny wykładniczej]
\label{thm:glm-moments}
W założeniach definicji~\ref{def:glm-family}, jeżeli
$Y$ ma rozkład $f(\cdot;\theta,\phi)$, to
\[
 \E_{\theta,\phi}Y=b'(\theta),\qquad
 \Var_{\theta,\phi}(Y)=a(\phi)b''(\theta).
\]
\end{twierdzenie}
\begin{proof}
Niech $\ell(\theta;y)=\log f(y;\theta,\phi)$.
Różniczkując równość
$\int_{\mathsf Y}f(y;\theta,\phi)\,\dd\nu(y)=1$
otrzymujemy
$0=\E\{\partial_\theta\ell(\theta;Y)\}$,
a więc $0=(\E Y-b'(\theta))/a(\phi)$.
Po ponownym różniczkowaniu normalizacji i
wykorzystaniu
$\partial_\theta^2\ell=-b''(\theta)/a(\phi)$
mamy
$0=\E[(\partial_\theta\ell)^2]+
\E[\partial_\theta^2\ell]$.
Po podstawieniu $\E Y=b'(\theta)$ i pomnożeniu przez
$a(\phi)^2$ otrzymujemy żądany wzór na wariancję.
\end{proof}

\begin{definicja}[Uogólniony model liniowy]
\label{def:glm}
Niech $x_i\in\R^p$, $i=1,\ldots,n$, będą ustalonymi
wektorami cech, a $Y_1,\ldots,Y_n$ niezależnymi
odpowiedziami, których rozkłady należą do rodziny
z definicji~\ref{def:glm-family} przy wspólnym $\phi$.
Dla każdego $i$ niech $\mu_i=\E Y_i$ i niech
$g$ będzie ustaloną różniczkowalną, odwracalną funkcją
na dopuszczalnym zbiorze średnich. Model
\[
 g(\mu_i)=\eta_i=x_i^\top\beta,
 \quad \beta\in B\subseteq\R^p
\]
nazywamy GLM, jeśli wszystkie $\mu_i=g^{-1}(x_i^\top\beta)$
należą do zbioru średnich rodziny. Funkcja $g$
to \emph{funkcja łącząca}; łącze
$g=b'^{-1}$ na odpowiedniej dziedzinie
nazywamy \emph{kanonicznym}.
\end{definicja}

\begin{center}\small
\begin{tabular}{>{\raggedright\arraybackslash}p{.20\textwidth}>{\raggedright\arraybackslash}p{.18\textwidth}>{\raggedright\arraybackslash}p{.23\textwidth}>{\raggedright\arraybackslash}p{.23\textwidth}}
\toprule
\textbf{Rozkład} & \textbf{Średnia} & \textbf{Łącze kanoniczne} & \textbf{Wariancja warunkowa}\\
\midrule
Normalny, stała wariancja & $\mu\in\R$ & $g(\mu)=\mu$ & $\sigma^2$\\
Bernoulliego & $p\in(0,1)$ & $\log(p/(1-p))$ & $p(1-p)$\\
Poissona & $\mu>0$ & $g(\mu)=\log\mu$ & $\mu$\\
\bottomrule
\end{tabular}
\end{center}

\section{Regresja logistyczna: prawdopodobieństwa, ilorazy szans i decyzje}
Niech $Y_i\in\{0,1\}$ będzie wskaźnikiem wady i
$Y_i\mid x_i\sim\operatorname{Bernoulli}(p_i)$.
Regresja logistyczna zakłada
\[
 p_i=\frac{\exp(x_i^\top\beta)}
          {1+\exp(x_i^\top\beta)},\qquad
 \log\frac{p_i}{1-p_i}=x_i^\top\beta.
\]
Dla modelu z wyrazem wolnym i zmienną $x$ zwiększenie
$x$ o jednostkę mnoży \emph{szanse} $p/(1-p)$ przez
$\exp(\beta_1)$; nie mnoży na ogół samego
prawdopodobieństwa przez tę liczbę.

\begin{przyklad}[Dlaczego szansa i prawdopodobieństwo są różne?]
Niech $\beta_1=\log2$. Szanse ulegają podwojeniu.
Jeśli prawdopodobieństwo początkowe wynosi $0{,}1$,
to szanse wzrastają z $1/9$ do $2/9$,
a prawdopodobieństwo do $2/11\approx0{,}1818$.
Jeżeli prawdopodobieństwo początkowe wynosi $0{,}6$,
to szanse wzrastają z $3/2$ do $3$,
a prawdopodobieństwo do $0{,}75$.
Taki sam iloraz szans odpowiada więc różnym zmianom
prawdopodobieństwa.
\end{przyklad}

\begin{twierdzenie}[Log-wiarygodność i wklęsłość w regresji logistycznej]
\label{thm:logit-score}
Niech $Y_1,\ldots,Y_n$ będą niezależnymi zmiennymi
Bernoulliego z parametrami
$p_i(\beta)=\exp(x_i^\top\beta)/
(1+\exp(x_i^\top\beta))$, $\beta\in\R^p$, gdzie
$X$ ma wiersze $x_i^\top$. Wówczas log-wiarygodność
(z pominięciem składników niezależnych od $\beta$),
jej gradient i hesjan wynoszą
\[
 \ell(\beta)=\sum_{i=1}^n\bigl[Y_i x_i^\top\beta
        -\log(1+e^{x_i^\top\beta})\bigr],\quad
 \nabla\ell=X^\top(Y-p),\quad
 \nabla^2\ell=-X^\top W X,
\]
gdzie $p=(p_1,\ldots,p_n)^\top$ oraz
$W=\operatorname{diag}(p_i(1-p_i))$.
Funkcja $\ell$ jest wklęsła, a jeśli $X$ ma rząd $p$,
to jest ściśle wklęsła dla każdego skończonego $\beta$.
\end{twierdzenie}
\begin{proof}
Dla obserwacji Bernoulliego
$\log P_\beta(Y_i=y_i)=y_i\log p_i+(1-y_i)\log(1-p_i)$.
Po podstawieniu logitu otrzymujemy
$y_ix_i^\top\beta-\log(1+e^{x_i^\top\beta})$.
Pochodna funkcji $\log(1+e^u)$ wynosi
$e^u/(1+e^u)=p_i$, a druga pochodna to $p_i(1-p_i)$.
Dlatego dla każdego $v\in\R^p$
$v^\top(-\nabla^2\ell)v=
\sum_i p_i(1-p_i)(x_i^\top v)^2\ge0$.
Przy pełnym rzędzie $X$ i skończonym $\beta$ wszystkie
wagi są dodatnie, co daje ścisłą nierówność dla $v\ne0$.
\end{proof}

\begin{uwaga}[Ścisła wklęsłość nie zapewnia istnienia MLE]
\label{rem:glm-separation}
Jeżeli obserwacje z $Y_i=0$ oraz $Y_i=1$ można
idealnie rozdzielić pewną liniową kombinacją cech,
log-wiarygodność może rosnąć wzdłuż kierunku
$\|\beta\|\to\infty$, nie osiągając maksimum
w $\R^p$. Nazywamy to \emph{separacją}.
Przykładowo dla dwóch obserwacji $(x,y)=(0,0)$
i $(1,1)$ w modelu z wyrazem wolnym wartości
$\beta_0=-t$, $\beta_1=2t$ przy $t\to\infty$
dają prawdopodobieństwa sukcesu dążące odpowiednio
do $0$ i $1$. Log-wiarygodność dąży do supremum
$0$, lecz nie osiąga go w skończonym punkcie.
Analogiczny problem może wystąpić przy separacji
quasi-zupełnej; brak doskonałej klasyfikacji wszystkich
obserwacji nie gwarantuje więc istnienia skończonego MLE.
W praktyce sprawdzamy separację; można zastosować
regularizację lub metody redukcji obciążenia, ale
wtedy jawnie zmienia się zadanie estymacji.
\end{uwaga}

\section{Regresja Poissona: intensywność, liczba zdarzeń i ekspozycja}
\begin{twierdzenie}[Model Poissona i jego log-wiarygodność]
\label{thm:poisson-glm-score}
Niech $t_i>0$ oznacza znaną ekspozycję, a
$Y_i$, $i=1,\ldots,n$, będą niezależne oraz
$Y_i\sim\Pois(\mu_i)$, gdzie
$\mu_i=t_i\exp(x_i^\top\beta)$, $\beta\in\R^p$.
Wówczas log-wiarygodność, jej gradient i hesjan są równe
\[
 \begin{aligned}
 \ell(\beta)&=\sum_{i=1}^n
 [Y_i(\log t_i+x_i^\top\beta)
      -t_i e^{x_i^\top\beta}-\log(Y_i!)],\\
 \nabla\ell(\beta)&=X^\top(Y-\mu),\\
 \nabla^2\ell(\beta)&=-X^\top\operatorname{diag}(\mu_i)X.
 \end{aligned}
\]
Jeśli $X$ ma pełny rząd kolumnowy, log-wiarygodność
jest ściśle wklęsła w każdym punkcie $\R^p$.
\end{twierdzenie}
\begin{proof}
W modelu Poissona
$\log P(Y_i=y_i)=y_i\log\mu_i-\mu_i-\log(y_i!)$.
Podstawiamy $\log\mu_i=\log t_i+x_i^\top\beta$.
Ponieważ $\partial_\beta\mu_i=\mu_ix_i$,
różniczkowanie po $\beta$ daje gradient.
Ponowne różniczkowanie daje
$-\sum_i\mu_ix_ix_i^\top$.
Każda $\mu_i>0$, a więc pełny rząd $X$ zapewnia
ujemną określoność hesjanu.
\end{proof}

\begin{przyklad}[Nie wolno pomijać ekspozycji]
Dwie maszyny mają odpowiednio $4$ awarie w ciągu
$20$ godzin oraz $8$ awarii w ciągu $100$ godzin.
Surowo druga ma dwa razy więcej awarii, ale
obserwowane częstości wynoszą $0{,}20$ i $0{,}08$
na godzinę. Model $\mu_i=t_i e^{x_i^\top\beta}$
traktuje $\log t_i$ jako \emph{offset}, czyli
znany składnik predyktora o współczynniku równym $1$.
Nie estymujemy współczynnika przy ekspozycji,
chyba że świadomie zmieniamy model procesu.
\end{przyklad}

\begin{uwaga}[Zera w zliczeniach i brzeg przestrzeni parametrów]
Dla niezależnych $Y_i\sim\Pois(t_i\lambda)$,
$t_i>0$, z ograniczeniem $\lambda>0$, jeśli
$\sum_iY_i=0$, log-wiarygodność jest ściśle
malejąca jako funkcja $\lambda$. Jej supremum jest
osiągane jedynie w granicy $\lambda\downarrow0$;
w otwartej przestrzeni parametrów nie istnieje MLE.
Jeśli model dopuszcza $\lambda=0$, MLE wynosi $0$,
ale jest to punkt brzegowy i zwykła normalna asymptotyka
Walda nie musi tam obowiązywać. Problem dotyczy
zwłaszcza rzadkich zdarzeń.
\end{uwaga}

\section{Score i informacja Fishera w ogólnym GLM}
\begin{twierdzenie}[Równania score i informacja dla ustalonej dyspersji]
\label{thm:glm-general-score}
Niech model spełnia definicję~\ref{def:glm}.
Załóżmy, że $b''(\theta_i)>0$, że
$\mu_i(\eta)=g^{-1}(\eta)$ jest różniczkowalna,
$\dd\mu_i/\dd\eta_i\ne0$ i że $a(\phi)>0$
jest znane. Oznaczmy $V_i=b''(\theta_i)$,
a przez $\dot\mu_i$ pochodną $\mu_i$ względem
$\eta_i=x_i^\top\beta$. Wówczas
\[
 \nabla\ell(\beta)=\sum_{i=1}^n
 \frac{Y_i-\mu_i}{a(\phi)V_i}\dot\mu_i x_i,
 \qquad
 \mathcal I(\beta)=\sum_{i=1}^n
 \frac{\dot\mu_i^2}{a(\phi)V_i}x_ix_i^\top.
\]
W przypadku łącza kanonicznego
$\dot\mu_i=V_i$, a równanie score przyjmuje postać
$\sum_i x_i(Y_i-\mu_i)=0$.
\end{twierdzenie}
\begin{proof}
Pochodna log-gęstości względem $\theta_i$ wynosi
$(Y_i-b'(\theta_i))/a(\phi)=(Y_i-\mu_i)/a(\phi)$.
Ponieważ $\dd\mu_i/\dd\theta_i=V_i>0$,
reguła łańcuchowa daje
$\dd\theta_i/\dd\eta_i=\dot\mu_i/V_i$.
Pomnożenie przez $x_i=\nabla_\beta\eta_i$ prowadzi
do pierwszego wzoru. Niezależność obserwacji i
$\Var(Y_i)=a(\phi)V_i$ oznaczają, że macierz
kowariancji score jest sumą
$\dot\mu_i^2x_ix_i^\top/[a(\phi)V_i]$.
Przy łączu kanonicznym $\eta_i=\theta_i$, skąd
$\dot\mu_i=b''(\theta_i)=V_i$.
\end{proof}

\section{Skąd bierze się IRLS? Iteracyjnie ważone najmniejsze kwadraty}
Z punktu widzenia obliczeń MLE rozwiązujemy na ogół
nieliniowe równania score. \emph{Fisher scoring} to
iteracja $\beta^{(k+1)}=\beta^{(k)}+
\mathcal I(\beta^{(k)})^{-1}\nabla\ell(\beta^{(k)})$,
o ile macierz informacji jest odwracalna.
Niech $\eta_i=x_i^\top\beta^{(k)}$, $\mu_i=g^{-1}(\eta_i)$,
$\dot\mu_i\ne0$ i zdefiniujmy
\[
 w_i=\frac{\dot\mu_i^2}{a(\phi)V_i},\qquad
 z_i=\eta_i+\frac{Y_i-\mu_i}{\dot\mu_i}.
\]
Równanie normalne dla ważonych najmniejszych kwadratów
$\min_{b\in\R^p}\sum_iw_i(z_i-x_i^\top b)^2$
przy pełnym rzędzie macierzy ważonej daje
$\beta^{(k+1)}=(X^\top WX)^{-1}X^\top Wz$.
Ponieważ
$X^\top W(z-X\beta^{(k)})=
\sum_i[(Y_i-\mu_i)\dot\mu_i/(a(\phi)V_i)]x_i$,
iteracja jest dokładnie krokiem Fisher scoring.
Jest to \emph{IRLS} (iteratively reweighted least squares).

\begin{uwaga}[Kiedy algorytm może nie zbiegać?]
Krok Fisher scoring nie musi automatycznie zwiększać
wiarygodności, a iteracja może być niestabilna przy
separacji, niemal osobliwej macierzy planu lub zbyt
dużym kroku. Produkcyjne implementacje stosują
kontrolę kroku, ograniczenia parametrów i diagnostykę
zbieżności. Żaden wynik programu nie jest wiarygodny,
jeśli algorytm jedynie zakończył iteracje z ostrzeżeniem.
\end{uwaga}

\section{Testowanie współczynników i diagnostyka modelu GLM}
\begin{stwierdzenie}[Informacja i przybliżone błędy standardowe]
\label{prop:glm-asymptotic-se}
Jeśli w ciągu modeli GLM spełnione są warunki
regularności z rozdziału~\ref{ch:asymptotic-estimation} zapewniające
istnienie zgodnego MLE $\widehat\beta$, asymptotyczną
normalność oraz zgodność oszacowania informacji,
to przybliżona macierz kowariancji estymatora wynosi
$\mathcal I(\widehat\beta)^{-1}$. Dla współczynnika
$\beta_j$ daje to orientacyjny błąd standardowy
$\mathrm{SE}_j=\sqrt{[\mathcal I(\widehat\beta)^{-1}]_{jj}}$
oraz przedział Walda $\widehat\beta_j\pm
z_{1-\alpha/2}\mathrm{SE}_j$.
Dla ilorazu szans $\exp(\beta_j)$ przekształcamy
\emph{oba końce} przedziału przez funkcję wykładniczą.
Nie są to dokładne przedziały dla małych prób.
\end{stwierdzenie}

\begin{definicja}[Dewiancja i reszty Pearsona]
Dla modelu o funkcji wiarygodności $L(\beta)$ oraz
modelu nasyconego, który dopasowuje osobną dopuszczalną
średnią każdej obserwacji, definiujemy \emph{dewiancję}
$D=2(\sup\ell_{\rm sat}-\ell(\widehat\beta))$
przy tej samej ustalonej dyspersji; w modelach
Poissona i Bernoulliego bez dodatkowej dyspersji
jest ona nieujemna. Dla danych $Y_i$ i wartości
$\widehat\mu_i$ resztę Pearsona określamy wzorem
$r_{P,i}=(Y_i-\widehat\mu_i)/\sqrt{\widehat{\Var}(Y_i)}$,
gdy oszacowana wariancja jest dodatnia.
\end{definicja}

\begin{stwierdzenie}[Dewiancja w regresji Poissona]
\label{prop:poisson-deviance}
Dla niezależnych zliczeń $Y_i\sim\Pois(\mu_i)$ i
oszacowań $\widehat\mu_i>0$ otrzymujemy
\[
 D=2\sum_{i=1}^n\left[
 Y_i\log\frac{Y_i}{\widehat\mu_i}
 -(Y_i-\widehat\mu_i)\right],
\]
przy konwencji $0\log(0/a)=0$ dla $a>0$.
Każdy składnik sumy jest nieujemny.
\end{stwierdzenie}
\begin{proof}
Dla ustalonego $y_i\ge0$ maksimum
$y_i\log\mu_i-\mu_i-\log(y_i!)$
po $\mu_i\ge0$ osiąga się w punkcie $\mu_i=y_i$;
przy $y_i=0$ rozumiemy maksimum na domknięciu
$[0,\infty)$, zaś rozkład Poissona z parametrem
$\mu=0$ jest zdegenerowany. Odejmując log-wiarygodność
modelu ograniczonego od supremum nasyconego, otrzymujemy wzór.
Nierówność $z\log z-z+1\ge0$ dla $z>0$,
z granicą w $z=0$, dowodzi nieujemności składników.
\end{proof}

\paragraph{Jak używać diagnostyki?}
Dla regresji logistycznej porównujemy przewidywane
prawdopodobieństwa z obserwowanymi częstościami
w odpowiednich grupach, oceniamy kalibrację,
zdolność rozróżniania klas i ryzyko separacji.
Dla zliczeń sprawdzamy nie tylko średnią, lecz również
zróżnicowanie: model Poissona zakłada
$\Var(Y_i\mid x_i)=\E(Y_i\mid x_i)$.
W obu przypadkach sprawdzamy punkty wpływowe,
strukturę reszt i poprawność niezależności obserwacji.
Brak widocznej struktury na wykresie reszt nie dowodzi
niezależności; dane uporządkowane chronologicznie należy
również analizować w funkcji czasu. Szczegóły diagnostyki
autokorelacji rozwiniemy w rozdziale~\ref{ch:time-series}.

\section{Nadmierna dyspersja i błędna specyfikacja}
\begin{przyklad}[Dlaczego standardowy Poisson może zaniżać niepewność?]
W identycznych warunkach średnia liczba awarii
wynosi $2$, natomiast empiryczna wariancja z
wielu niezależnych okresów wynosi około $9$.
Poisson wymaga równości średniej i wariancji,
więc sam model warunkowej średniej nie opisuje
wariancji danych. Klasyczne błędy standardowe
oparte na informacji Poissona mogą być zaniżone.
Trzeba sprawdzić, czy wynika to ze zmienności
niezaobserwowanych intensywności, zależności,
nadmiernej liczby zer lub innej specyfikacji.
\end{przyklad}

\begin{stwierdzenie}[Mieszanina Poisson--gamma prowadzi do nadmiernej dyspersji]
\label{prop:poisson-gamma}
Niech $\Lambda$ będzie dodatnią zmienną losową o
$\E\Lambda=\mu>0$ i $\Var(\Lambda)=\kappa\mu^2$,
$\kappa>0$. Niech warunkowo względem $\Lambda$
zmienna $Y$ ma rozkład $\Pois(\Lambda)$.
Wtedy $\E Y=\mu$ i
$\Var(Y)=\mu+\kappa\mu^2$.
Jeżeli dodatkowo $\Lambda$ ma rozkład gamma z
parametrem kształtu $1/\kappa$ i parametrem skali
$\kappa\mu$, to marginalny rozkład $Y$ jest
ujemnym dwumianowym.
\end{stwierdzenie}
\begin{proof}
Z prawa całkowitej wartości oczekiwanej mamy
$\E Y=\E(\E[Y\mid\Lambda])=\E\Lambda=\mu$.
Z prawa całkowitej wariancji
$\Var(Y)=\E(\Var(Y\mid\Lambda))+
\Var(\E[Y\mid\Lambda])=
\E\Lambda+\Var(\Lambda)=\mu+\kappa\mu^2$.
W wypadku gamma funkcja tworząca dla $Y$ jest równa
$\E[z^Y]=\E[e^{\Lambda(z-1)}]
=(1+\kappa\mu(1-z))^{-1/\kappa}$
dla $0\le z\le1$. Rozwinięcie dwumianowe tej
funkcji daje rozkład ujemny dwumianowy.
\end{proof}

\begin{uwaga}[Poisson, quasi-Poisson czy ujemny dwumianowy?]
Quasi-Poisson zachowuje funkcję średniej Poissona,
ale zakłada $\Var(Y_i\mid x_i)=\phi\mu_i$ dla
$\phi>0$ i nie wyznacza pełnego rozkładu obserwacji.
Nie używamy zatem klasycznego AIC opartego na
pełnej wiarygodności quasi-Poissona tak, jak dla
zwykłego modelu Poissona. Model ujemny dwumianowy
z wariancją $\mu_i+\kappa\mu_i^2$ wymaga określenia
pełnego rozkładu i estymacji $\kappa$.
Jeśli źródłem problemu są skorelowane obserwacje,
sama zmiana rodziny rozkładów może nie wystarczyć.
\end{uwaga}

\section{Studium przypadku: od modelu do użytecznej decyzji}
Rozważmy dwie technologie, A i B. W każdej przebadano
po $100$ niezależnych produktów: w A stwierdzono
$7$ wad, a w B $15$. Model z jedną zmienną binarną
$g_i=0$ dla A i $g_i=1$ dla B ma
$\operatorname{logit}(p_i)=\beta_0+\beta_1g_i$.
Ponieważ obie grupy mają dodatnią liczbę sukcesów
i porażek, log-wiarygodność osiąga maksimum
przy $\widehat p_A=0{,}07$ i $\widehat p_B=0{,}15$.
Stąd
\[
 \widehat\beta_1=
 \log\frac{0{,}15/(1-0{,}15)}{0{,}07/(1-0{,}07)}
 =\log\frac{15\cdot93}{85\cdot7}
 \approx0{,}852,\qquad
 \exp(\widehat\beta_1)\approx2{,}345.
\]
Interpretujemy: szanse wady w grupie B są oszacowane
na około $2{,}35$ razy większe niż w A. Nie jest to
tożsame z \emph{ryzykiem względnym}, równym
$0{,}15/0{,}07\approx2{,}14$, ani z różnicą
prawdopodobieństw, wynoszącą $0{,}08$.
Dla decyzji o wdrożeniu technologii zwykle potrzebujemy
również przedziałów niepewności, kosztów jednostkowych
wad i oceny porównywalności obu prób. Przybliżony
przedział $95\%$ dla ilorazu szans obliczony w programie
poniżej wynosi $[0{,}912;\ 6{,}027]$: obejmuje $1$.
Mimo znacznej różnicy oszacowań punktowych ten zbiór
danych nie uzasadnia zatem odrzucenia hipotezy równości
szans przy dwustronnej klasycznej procedurze Walda
na poziomie $5\%$. Jest to ilustracja różnicy między
rozmiarem oszacowanego efektu a jego niepewnością.

\begin{lstlisting}[language=Python,caption={Kompletny program: dwie grupy binarne, iloraz szans i przedział przybliżony (Python).}]
from math import exp, log, sqrt
from scipy.stats import norm

bad_a, total_a = 7, 100
bad_b, total_b = 15, 100
ok_a, ok_b = total_a - bad_a, total_b - bad_b

odds_ratio = bad_b * ok_a / (ok_b * bad_a)
log_or = log(odds_ratio)
se_log_or = sqrt(1/bad_b + 1/ok_b + 1/bad_a + 1/ok_a)
z = norm.ppf(0.975)
ci_or = (exp(log_or - z*se_log_or),
         exp(log_or + z*se_log_or))
risk_ratio = (bad_b/total_b) / (bad_a/total_a)
print('odds ratio:', odds_ratio)
print('risk ratio:', risk_ratio)
print('approximate 95% CI odds ratio:', ci_or)
# For zero cells consider exact/penalized approaches.
\end{lstlisting}

\begin{lstlisting}[language=Python,caption={Kompletny program: MLE intensywności i offset w modelu Poissona (SciPy).}]
import numpy as np
from scipy.optimize import minimize

exposure = np.array([10., 30., 15., 40., 20., 50.])
group_b = np.array([0., 0., 0., 1., 1., 1.])
counts = np.array([2., 4., 1., 4., 3., 5.])
X = np.column_stack((np.ones_like(exposure), group_b))

def negative_loglik(beta):
    eta = np.log(exposure) + X @ beta
    return float(np.exp(eta).sum() - counts @ eta)

def gradient(beta):
    eta = np.log(exposure) + X @ beta
    return X.T @ (np.exp(eta) - counts)

fit = minimize(negative_loglik, x0=np.array([-2., 0.]),
               method='BFGS', jac=gradient, tol=1e-8)
assert np.isfinite(fit.x).all()
rate_ratio_b_to_a = np.exp(fit.x[1])
rate_a = np.exp(fit.x[0])
rate_b = np.exp(fit.x[0] + fit.x[1])
print('rates per unit exposure:', rate_a, rate_b)
print('rate ratio B/A:', rate_ratio_b_to_a)
print('converged:', fit.success, fit.message)
if not fit.success:
    raise RuntimeError('Optimizer did not converge')
\end{lstlisting}

\paragraph{Jak interpretować obliczenia?}
Przy połączonych licznikach grup A i B powyższy model
Poissona ma dwa parametry i dzięki dostateczności sum
w grupach jego MLE można otrzymać również bez optymalizatora:
$\widehat\lambda_A=\sum_{i:g_i=0}Y_i/
\sum_{i:g_i=0}t_i$ oraz analogicznie dla B.
Program numeryczny ilustruje użycie offsetu: daje intensywności
$\widehat\lambda_A\approx0{,}1273$ i
$\widehat\lambda_B\approx0{,}1091$ na jednostkę
czasu, stąd iloraz intensywności około $0{,}8571$.
Nie jest to dowód różnicy między technologiami.
Wynik należy sprawdzić z podanymi wzorami, a ponadto
zbadać, czy wariancja danych odpowiada modelowi.

\section{Dobór modelu: praktyczna tabela porównawcza}
\begin{center}\small
\begin{tabular}{>{\raggedright\arraybackslash}p{.28\textwidth}>{\raggedright\arraybackslash}p{.30\textwidth}>{\raggedright\arraybackslash}p{.30\textwidth}}
\toprule
\textbf{Cel i rodzaj odpowiedzi} & \textbf{Model do rozważenia} & \textbf{Kontrola i ograniczenia}\\
\midrule
Wynik ciągły o liniowej średniej & Regresja liniowa normalna / MNK & Rozkład reszt dla dokładnej inferencji, wariancja, zależność.\\
Wynik binarny, np. wada / brak wady & Logistyczny GLM & Separacja, kalibracja, liczba zdarzeń, iloraz szans $\ne$ iloraz ryzyka.\\
Zliczenia przy znanej ekspozycji & Poisson z offsetem & Średnia $=$ wariancja warunkowa, niezależność, poprawna ekspozycja.\\
Zliczenia z nadmierną dyspersją & Ujemny dwumianowy lub quasi-Poisson & Źródło dyspersji, zależność i nadmiar zer; AIC nie dla quasi-likelihood.\\
Powtarzane pomiary binarne lub zliczenia & GLMM, GEE lub jawny model zależności & Właściwa jednostka niezależności; inna interpretacja parametrów.\\
Dane o czasie do awarii z cenzorowaniem & Modele przeżycia, np. Cox & Nie zastępować cenzorowanych czasów zliczeniami bez uzasadnienia.\\
\bottomrule
\end{tabular}
\end{center}

\begin{zadanie}[Czy zawsze logistyczna?]
W zakładzie zbadano $n=200$ produktów, ale nie tylko
liczbę wad, lecz również \emph{czas do uszkodzenia},
przy czym obserwacja wielu produktów zakończyła się
przed awarią. Analityk proponuje regresję logistyczną
wskaźnika „awaria wystąpiła do dnia zakończenia badania”.
Czy jest ona wystarczająca?
\par\emph{Rozwiązanie.} Taka regresja może być właściwa,
jeśli wszystkie produkty miały ten sam ustalony horyzont
obserwacji i interesuje nas wyłącznie zdarzenie do
tego horyzontu. Jeśli horyzonty są różne, a część
czasów cenzorowana, zwykła regresja wskaźnika zignoruje
czas ekspozycji oraz mechanizm cenzorowania. Wtedy
rozważamy model przeżycia albo model intensywności
zdarzeń z prawidłowym opisem czasu zagrożenia.
\end{zadanie}


% ================================================================
% ROZDZIAL 15 -- WYBOR I WALIDACJA MODELI
% ================================================================
\chapter{Wybór, walidacja i regularizacja modeli statystycznych}
\label{ch:model-selection}
\paragraph{Powiązania i odsyłacze.} Modele bazowe pochodzą z rozdziałów~\ref{ch:linear-regression} i~\ref{ch:glm}; o ryzyku predykcji i funkcjach straty zobacz rozdział~\ref{ch:decision}.


\section{Dopasowanie nie jest jeszcze uogólnianiem}
W rozdziałach~\ref{ch:linear-regression} i~\ref{ch:glm} nauczyliśmy się
estymować współczynniki modelu oraz wnioskować o nich \emph{pod założeniem
poprawności modelu}. Praktyczny problem jest trudniejszy: czy rozważany
model dobrze przewiduje nowe dane, czy nie dopasował jedynie przypadkowego
szumu? Jeśli dodamy do regresji liniowej kolejne kolumny, resztowa suma
kwadratów na danych uczących nie wzrośnie, ale jakość predykcji może się
pogorszyć. Odróżnimy trzy cele: \emph{opis zależności}, \emph{predykcję}
i \emph{wnioskowanie przyczynowe}. Dla każdego z nich dobór modelu może być
inny. W szczególności niska wartość błędu predykcji sama nie uzasadnia
interpretacji współczynnika jako efektu przyczynowego.

\begin{definicja}[Ryzyko predykcyjne i ryzyko empiryczne]
\label{def:predictive-risk}
Niech $(\mathsf X,\mathcal X)$ oraz $(\mathsf Y,\mathcal Y)$
będą przestrzeniami mierzalnymi, $(X,Y)$ parą losową
o prawie $P$ na ich iloczynie, a
$\ell:\mathsf Y\times\mathsf Y\to[0,\infty)$ mierzalną funkcją
straty. Dla mierzalnej funkcji predykcyjnej
$f:\mathsf X\to\mathsf Y$ definiujemy
$R_P(f)=\E_P\ell(Y,f(X))$, zakładając skończoność wartości oczekiwanej.
Dla obserwacji $(X_i,Y_i)_{i=1}^n$ ryzyko empiryczne wynosi
$\widehat R_n(f)=n^{-1}\sum_{i=1}^n\ell(Y_i,f(X_i))$.
Jeżeli $\widehat f_n$ jest wyznaczona z tej samej próby, to
$\widehat R_n(\widehat f_n)$ nazywamy \emph{błędem treningowym}.
Dla niezależnej od próby uczącej pary $(X_*,Y_*)\sim P$ błąd
nowej obserwacji ma warunkową wartość oczekiwaną
$R_P(\widehat f_n)=\E[\ell(Y_*,\widehat f_n(X_*))\mid\mathcal D_n]$,
gdzie $\mathcal D_n$ oznacza dane uczące.
\end{definicja}

\begin{uwaga}[Co jest losowe?]
$R_P(f)$ jest liczbą dla ustalonego $f$; $R_P(\widehat f_n)$ jest
zmienną losową, ponieważ procedura treningowa używa losowej próby.
Błąd na zbiorze testowym estymuje ryzyko \emph{już wytrenowanego}
modelu, pod warunkiem że testowe dane reprezentują docelowy rozkład
i nie były wykorzystywane przy wyborze tego modelu.
\end{uwaga}

\begin{twierdzenie}[Optymizm błędu treningowego w normalnym modelu liniowym]
\label{thm:training-optimism}
Niech $n>p\ge1$, $X\in\R^{n\times p}$ ma rząd $p$ oraz
$Y=X\beta+\varepsilon$, gdzie
$\varepsilon\sim\mathcal N_n(0,\sigma^2I_n)$ i $\sigma>0$.
Niech $H=X(X^\top X)^{-1}X^\top$ oraz $\widehat Y=HY$.
Niech niezależny wektor $Y^{\rm new}=X\beta+\varepsilon^{\rm new}$
ma ten sam rozkład co $Y$. Wtedy
\[
 \E\|Y-\widehat Y\|_2^2=(n-p)\sigma^2,
 \qquad
 \E\|Y^{\rm new}-\widehat Y\|_2^2=(n+p)\sigma^2.
\]
Zatem optymizm błędu treningowego wynosi dokładnie
$2p\sigma^2/n$ po podzieleniu sum kwadratów przez $n$.
\end{twierdzenie}
\begin{proof}
Ponieważ $(I_n-H)X\beta=0$, mamy
$Y-\widehat Y=(I_n-H)\varepsilon$. Stąd
$\E\|Y-\widehat Y\|^2=\sigma^2\operatorname{tr}(I_n-H)
=(n-p)\sigma^2$.
Z kolei $Y^{\rm new}-\widehat Y=\varepsilon^{\rm new}-H\varepsilon$.
Składniki są niezależne, wycentrowane, więc średni kwadrat normy to
$\sigma^2\operatorname{tr}(I_n)+\sigma^2\operatorname{tr}(H^\top H)
=(n+p)\sigma^2$, gdyż $H$ jest rzutem rzędu $p$.
\end{proof}

\begin{przyklad}[Więcej parametrów nie musi pomagać]
Dwa modele liniowe o wymiarach $p=2$ i $p=12$ dopasowano do tej samej
próby $n=50$ przy poprawności obu modeli i tej samej wariancji błędu.
Różnica oczekiwanego optymizmu, liczona na obserwację, wynosi
$2(12-2)\sigma^2/50=0{,}4\sigma^2$. To nie mówi, który model
jest bezwzględnie lepszy: model prostszy może pomijać ważną
zależność, a model bogatszy zwiększa wariancję predykcji.
\end{przyklad}

\section{AIC i BIC: dwie kary za złożoność}
\begin{definicja}[Kryteria informacyjne]
Niech $\{P_\theta:\theta\in\Theta\subset\R^k\}$ będzie
regularnym modelem parametrycznym o log-wiarygodności $\ell_n(\theta)$
dla \emph{tych samych obserwacji} $y_1,\ldots,y_n$.
Jeżeli istnieje maksimum $\widehat\theta$, definiujemy
\[
 \operatorname{AIC}=-2\ell_n(\widehat\theta)+2k,
 \qquad
 \operatorname{BIC}=-2\ell_n(\widehat\theta)+k\log n.
\]
AIC to \emph{kryterium informacyjne Akaikego} (Akaike information
criterion), a BIC to \emph{bayesowskie kryterium informacyjne}
(Bayesian information criterion). W obu przypadkach mniejsza
wartość oznacza preferencję kryterium wśród porównywanych modeli.
Liczba $k$ obejmuje wszystkie
estymowane parametry, także parametr wariancji, jeżeli jest nieznany.
\end{definicja}

\begin{twierdzenie}[AIC i BIC w regresji normalnej]
\label{thm:aic-gaussian}
Niech $Y\sim\mathcal N_n(X\beta,\sigma^2I_n)$, gdzie
$X\in\R^{n\times p}$ ma pełny rząd kolumnowy, $n>p$,
a $\beta\in\R^p$ i $\sigma^2>0$ są nieznane.
Jeśli $\mathrm{RSS}=\min_b\|Y-Xb\|^2>0$, to MLE wariancji
wynosi $\widehat\sigma^2=\mathrm{RSS}/n$ i
\[
 -2\ell_n(\widehat\beta,\widehat\sigma^2)
 =n\bigl(\log(2\pi)+1+\log(\mathrm{RSS}/n)\bigr).
\]
W szczególności, przy porównywaniu modeli dla tego samego $Y$,
$\operatorname{AIC}=n\log(\mathrm{RSS}/n)+2(p+1)+C_n$ oraz
$\operatorname{BIC}=n\log(\mathrm{RSS}/n)+(p+1)\log n+C_n$,
gdzie $C_n=n(\log(2\pi)+1)$ nie zależy od modelu.
\end{twierdzenie}
\begin{proof}
Log-wiarygodność wynosi
$\ell_n(b,v)=-\tfrac n2\log(2\pi v)-\|Y-Xb\|^2/(2v)$.
Dla każdego $v>0$ maksymalizacja po $b$ prowadzi do MNK.
Dla $v>0$ funkcja $-\tfrac n2\log v-\mathrm{RSS}/(2v)$
ma pochodną $(\mathrm{RSS}-nv)/(2v^2)$ i osiąga
jedyne maksimum w $v=\mathrm{RSS}/n$, pod warunkiem
$\mathrm{RSS}>0$. Podstawienie i policzenie parametrów daje tezę.
\end{proof}

\begin{uwaga}[Dlaczego kryteria różnią się celem?]
AIC ma uzasadnienie asymptotyczne w szacowaniu pozapróbkowego
błędu logarytmicznego (straty związanej z dywergencją
Kullbacka--Leiblera). BIC powstaje, przy dodatkowych
założeniach, z przybliżenia logarytmu \emph{zintegrowanej
wiarygodności} i silniej karze liczbę parametrów, gdy $n$
jest duże. Nie są to równoważne cele. Nie należy
interpretować $\Delta\mathrm{AIC}$ jako p-value ani niższego BIC
jako dowodu prawdziwości modelu.
\end{uwaga}

\begin{stwierdzenie}[Źródło kary BIC: przybliżenie Laplace'a]
\label{prop:bic-laplace}
Niech $\Theta\subset\R^k$ będzie otwarta, a $\pi$
gęstością a priori dodatnią i ciągłą w otoczeniu
$\theta_0$. Załóżmy, że log-wiarygodność $\ell_n$
ma wewnętrzny maksymalizator $\widehat\theta_n$ zbieżny do
$\theta_0$ według prawdopodobieństwa, że lokalnie
$-n^{-1}\nabla^2\ell_n(\theta)\to J$ jednostajnie według
prawdopodobieństwa, gdzie $J$ jest dodatnio określona,
i że pozostała część całki jest zaniedbywalna względem
wkładu otoczenia maksimum. Zakładamy również
odpowiednią dominację umożliwiającą asymptotykę Laplace'a.
Wtedy dla $m_n=\int_\Theta e^{\ell_n(\theta)}\pi(\theta)
\,\dd\lambda_k(\theta)$ zachodzi
$-2\log m_n=-2\ell_n(\widehat\theta_n)+k\log n
+O_{\mathbb P}(1)$.
\end{stwierdzenie}
\begin{proof}
W otoczeniu maksimum rozwinięcie Taylora daje
$\ell_n(\theta)=\ell_n(\widehat\theta_n)-\tfrac n2
(\theta-\widehat\theta_n)^\top J
(\theta-\widehat\theta_n)+o(n\|\theta-\widehat\theta_n\|^2)$
na skali $\|\theta-\widehat\theta_n\|=O(n^{-1/2})$,
w sensie prawdopodobieństwa.
Po podstawieniu $u=\sqrt n(\theta-\widehat\theta_n)$
całka lokalna jest asymptotycznie równa
$e^{\ell_n(\widehat\theta_n)}\pi(\theta_0)n^{-k/2}
(2\pi)^{k/2}(\det J)^{-1/2}$.
Założenie o zaniedbywalności reszty pozwala zastąpić
całą całkę wkładem lokalnym; logarytmowanie kończy dowód.
\end{proof}

\begin{uwaga}[Uczciwe granice twierdzenia]
Do powyższego przybliżenia nie wystarcza samo istnienie
hesjanu w maksimum: potrzeba kontroli całki poza nim,
regularności skali $n^{-1/2}$ i właściwego prioru.
Przy parametrach na brzegu, w modelach osobliwych
lub przy silnej nieidentyfikowalności kara $k\log n$
może nie być właściwa. W szczególności pełne
uzasadnienie AIC w modelu nieregularnym wymaga
osobnej analizy; prosty wzór z karą $2k$ nie jest
uniwersalną korektą optymizmu.
\end{uwaga}

\section{Rozdzielenie danych, walidacja krzyżowa i wyciek informacji}
\begin{definicja}[Walidacja $K$-krotna]
Niech $(Z_i)_{i=1}^n$, $Z_i=(X_i,Y_i)$, będzie próbą
niezależną o wspólnym rozkładzie. Dzielimy indeksy
$\{1,\ldots,n\}$ na niepuste, rozłączne zbiory
$V_1,\ldots,V_K$, $2\le K\le n$, o sumie równej
całemu zbiorowi. Dla algorytmu uczącego $\mathcal A$
niech $\widehat f^{(-j)}$ oznacza model wytrenowany
wyłącznie na danych o indeksach spoza $V_j$.
Wtedy \emph{błąd walidacji krzyżowej} jest równy
\[
 \widehat R_{\rm CV}=
 \frac1n\sum_{j=1}^K\sum_{i\in V_j}
 \ell(Y_i,\widehat f^{(-j)}(X_i)).
\]
\end{definicja}

\begin{stwierdzenie}[Bezstronność błędu testowego przy niezależnym teście]
\label{prop:test-unbiased}
Niech $(\mathsf X,\mathcal X)$ i $(\mathsf Y,\mathcal Y)$ będą przestrzeniami mierzalnymi, $P$ prawem na $(\mathsf X\times\mathsf Y,\mathcal X\otimes\mathcal Y)$, a $m\ge1$. Niech $\mathcal D_{\rm train}$ oznacza dane uczące, $\widehat f=\mathcal A(\mathcal D_{\rm train})$ będzie mierzalną funkcją predykcyjną $\mathsf X\to\mathsf Y$, a $(X_j^*,Y_j^*)_{j=1}^m$ niezależną próbą testową o wspólnym prawie $P$, niezależną od $\mathcal D_{\rm train}$. Niech strata $\ell:\mathsf Y\times\mathsf Y\to\R$ będzie mierzalna i $\E|\ell(Y_j^*,\widehat f(X_j^*))|<\infty$. Definiujemy losowe ryzyko $R_P(\widehat f)=\int_{\mathsf X\times\mathsf Y}\ell(y,\widehat f(x))\,\dd P(x,y)$. Wówczas estymator
$\widehat R_{\rm test}=m^{-1}\sum_{j=1}^m
\ell(Y_j^*,\widehat f(X_j^*))$
spełnia $\E(\widehat R_{\rm test}\mid\mathcal D_{\rm train})
=R_P(\widehat f)$.
\end{stwierdzenie}
\begin{proof}
Po warunkowaniu względem danych uczących $\widehat f$
jest ustalona, zaś każda para $(X_j^*,Y_j^*)$ ma rozkład $P$.
Liniowość warunkowej wartości oczekiwanej daje
$\E[\ell(Y_j^*,\widehat f(X_j^*))\mid\mathcal D_{\rm train}]
=\int_{\mathsf X\times\mathsf Y}\ell(y,\widehat f(x))\,\dd P(x,y)$.
Uśrednienie po $j$ kończy dowód.
\end{proof}

\begin{przyklad}[Wyciek informacji przed walidacją]
Jeśli standaryzujemy kolumny całej tabeli, a dopiero potem
dzielimy dane na trening i test, średnie i odchylenia użyte
do transformacji korzystają z obserwacji testowych. Podobny
problem pojawia się przy imputacji braków danych, wyborze cech
na podstawie korelacji z wynikiem albo selekcji hiperparametrów
na zbiorze testowym. Każda taka operacja musi zostać wykonana
\emph{wewnątrz} odpowiedniego fragmentu uczącego.
Dla próby niezależnej służy do tego \texttt{Pipeline}; przy
selekcji hiperparametrów i równoczesnym szacowaniu jakości
stosujemy zagnieżdżoną walidację krzyżową albo oddzielny test.
\end{przyklad}

\begin{center}
\begin{tikzpicture}[>=Latex, every node/.style={font=\small},
  s/.style={draw,rounded corners,align=center,minimum height=9mm,inner sep=3pt},
  node distance=11mm and 8mm]
\node[s] (d) {Dane dostępne do analizy};
\node[s,right=of d] (tr) {Trening i walidacja\\strojenie modelu};
\node[s,right=of tr] (te) {Odłożony test\\jednorazowa ocena};
\draw[->,thick] (d)--(tr);
\draw[->,thick] (tr)--(te);
\node[below=6mm of tr,align=center,text=black!70] {Wszystkie transformacje i dobór parametrów\\wyłącznie w obrębie danych treningowych};
\end{tikzpicture}
\end{center}

\begin{uwaga}[Czego walidacja $K$-krotna nie obiecuje?]
$\widehat R_{\rm CV}$ nie jest na ogół dokładnie
nieobciążonym oszacowaniem ryzyka modelu trenowanego
na całych $n$ obserwacjach: w kolejnych podziałach
uczymy na mniejszych próbach. Przy danych zależnych
losowe mieszanie obserwacji może dodatkowo zaniżać błąd.
Dla szeregów czasowych stosujemy podziały zachowujące
chronologię, opisane w rozdziale~\ref{ch:time-series}.
\end{uwaga}

\section{Ridge i lasso: kiedy wprowadzenie obciążenia pomaga predykcji?}
\begin{definicja}[Regresja ridge i lasso]
Dla ustalonej macierzy $X\in\R^{n\times p}$, wektora
$y\in\R^n$ i $\lambda>0$ definiujemy:
\begin{itemize}
\item \emph{ridge}: dowolny minimizator
$Q_{\rm R}(b)=\frac12\|y-Xb\|_2^2+\frac\lambda2\|b\|_2^2$;
\item \emph{lasso}: dowolny minimizator
$Q_{\rm L}(b)=\frac12\|y-Xb\|_2^2+\lambda\sum_{j=1}^p|b_j|$.
\end{itemize}
W zastosowaniach z wyrazem wolnym zwykle \emph{nie karzemy}
tego wyrazu i standaryzujemy cechy liczbowe, ponieważ
wielkość kary zależy od jednostek współczynników.
\end{definicja}

\begin{twierdzenie}[Rozwiązanie ridge i jego obciążenie]
\label{thm:ridge}
Dla dowolnego $X\in\R^{n\times p}$, $y\in\R^n$
i $\lambda>0$ problem ridge ma dokładnie jedno rozwiązanie:
$\widehat\beta_\lambda=(X^\top X+\lambda I_p)^{-1}X^\top y$.
Jeżeli $Y=X\beta+\varepsilon$ z $\E\varepsilon=0$
i $\Cov(\varepsilon)=\sigma^2I_n$, to
\[
 \E\widehat\beta_\lambda=(X^\top X+\lambda I_p)^{-1}
 X^\top X\beta,
 \qquad
 \Cov(\widehat\beta_\lambda)=\sigma^2A_\lambda X^\top X A_\lambda,
 \quad A_\lambda=(X^\top X+\lambda I_p)^{-1}.
\]
\end{twierdzenie}
\begin{proof}
Hesjan funkcji $Q_{\rm R}$ wynosi
$X^\top X+\lambda I_p$, jest dodatnio określony, więc
kryterium jest ściśle wypukłe i koercywne.
Jego pochodna $X^\top(Xb-y)+\lambda b$ zeruje się
wyłącznie w podanym punkcie. Ostatnie równości
wynikają z liniowości estymatora i tożsamości
$\widehat\beta_\lambda=A_\lambda X^\top(X\beta+\varepsilon)$.
\end{proof}

\begin{twierdzenie}[Lasso z ortogonalnymi kolumnami]
\label{thm:lasso-soft}
Jeżeli $X\in\R^{n\times p}$ spełnia $X^\top X=I_p$,
a $z=X^\top y$, to jedyny minimizator $Q_{\rm L}$ ma
współrzędne
$\widehat\beta_{\lambda,j}=\operatorname{sgn}(z_j)
\max(|z_j|-\lambda,0)$, przy $\operatorname{sgn}(0)=0$.
\end{twierdzenie}
\begin{proof}
Po odrzuceniu składnika niezależnego od $b$ kryterium
jest sumą funkcji $q_j(t)=\frac12(t-z_j)^2+\lambda|t|$.
Dla $t>0$ pochodna to $t-z_j+\lambda$, dla $t<0$
wynosi $t-z_j-\lambda$. W punkcie $0$ warunek minimum
dla funkcji wypukłej to $0\in-z_j+\lambda[-1,1]$.
Z tych trzech przypadków otrzymujemy formułę.
Ścisła wypukłość składnika kwadratowego zapewnia
jedyność rozwiązania.
\end{proof}

\begin{uwaga}[Wybór ridge czy lasso]
Ridge zmniejsza współczynniki w sposób ciągły, ale na ogół
nie zeruje ich dokładnie; jest szczególnie pomocny przy
współliniowości i wielu słabych sygnałach.
Lasso może dawać współczynniki równe zeru, jednak przy
silnej korelacji cech wybór konkretnych zmiennych bywa
niestabilny. Żadna z metod nie naprawia błędnej definicji
jednostki obserwacji lub zależności czasowej. Wybór
$\lambda$ powinien korzystać wyłącznie z danych uczących.
\end{uwaga}

\section{Laboratorium: dobór stopnia wielomianu bez przecieku danych}
W kolejnym programie używamy wielomianowych regresorów
o stopniach od $1$ do $8$. Każdy kandydat ma oddzielny
potok transformacji: rozwinięcie cech i dopasowanie
regresji zachodzą tylko na danej części treningowej.
Na końcu porównujemy wybór walidacyjny z wynikiem na
odłożonym teście. W pokazanej, odtwarzalnej symulacji
najlepiej wypada stopień $2$, z walidacyjnym MSE około
$0{,}319$ i testowym MSE około $0{,}189$. Nie oczekujemy,
że wartości te będą identyczne: oba oszacowania opierają
się na różnych fragmentach skończonej próby.
Poniższy przykład jest symulacją
odtwarzalną po ustaleniu ziarna generatora; rzeczywiste
dane powinny być dzielone zgodnie z mechanizmem ich
powstawania.

\begin{lstlisting}[language=Python,caption={Kompletny program: wybór stopnia modelu, walidacja i niezależny test (Python, scikit-learn).}]
import numpy as np
from sklearn.model_selection import KFold, cross_val_score
from sklearn.pipeline import make_pipeline
from sklearn.preprocessing import PolynomialFeatures, StandardScaler
from sklearn.linear_model import Ridge
from sklearn.metrics import mean_squared_error

rng = np.random.default_rng(2026)
x = rng.uniform(-2.0, 2.0, size=150)
y = 1.5 + 0.8*x - 0.7*x**2 + rng.normal(0, 0.5, 150)
indices = rng.permutation(len(x))
train, test = indices[:110], indices[110:]
Xtr, Xte = x[train, None], x[test, None]
ytr, yte = y[train], y[test]
cv = KFold(n_splits=5, shuffle=True, random_state=11)

results = []
for degree in range(1, 9):
    pipe = make_pipeline(PolynomialFeatures(degree, include_bias=False),
                         StandardScaler(), Ridge(alpha=1.0))
    # Negative score because scikit-learn maximizes its scoring metric.
    mse_cv = -cross_val_score(pipe, Xtr, ytr, cv=cv,
                             scoring="neg_mean_squared_error").mean()
    results.append((degree, mse_cv))

best_degree, best_cv = min(results, key=lambda item: item[1])
best = make_pipeline(PolynomialFeatures(best_degree, include_bias=False),
                     StandardScaler(), Ridge(alpha=1.0))
best.fit(Xtr, ytr)
test_mse = mean_squared_error(yte, best.predict(Xte))
print("CV results:", results)
print("degree chosen by CV:", best_degree)
print("CV MSE:", best_cv)
print("holdout MSE:", test_mse)
\end{lstlisting}

\paragraph{Interpretacja.}
Program nie wybiera stopnia na podstawie wyniku testowego.
Nie porównujemy wielu konfiguracji na teście, a następnie
nie nazywamy najmniejszego błędu \emph{niezależną oceną}.
Przy rzeczywistych danych czasowych zwykły \texttt{KFold}
z mieszaniem byłby niewłaściwy; użyjemy innego podziału
w rozdziale~\ref{ch:time-series}.

\section{Porównania metod i zadania}
\begin{center}\small
\begin{tabular}{>{\raggedright\arraybackslash}p{.24\textwidth}>{\raggedright\arraybackslash}p{.32\textwidth}>{\raggedright\arraybackslash}p{.31\textwidth}}
\toprule
\textbf{Cel} & \textbf{Metoda} & \textbf{Warunek i ryzyko}\\
\midrule
Dopasowanie porównywalnych modeli o pełnej wiarygodności & AIC & Te same dane i odpowiedź; regularna asymptotyka; ryzyko przetrenowania.\\
Preferowanie prostszego modelu parametrycznego & BIC & Regularność i interpretacja w kategoriach aproksymacji ewidencji modelu.\\
Wybór jakości predykcji dla danych iid & Walidacja krzyżowa & Transformacje w ramach foldów; dobór liczby foldów i metryki.\\
Niezależna ocena końcowego modelu & Odłożony zbiór testowy & Test nie może uczestniczyć w strojeniu.\\
Wiele skorelowanych zmiennych & Ridge & Obciążenie za mniejszą wariancję; skala cech.\\
Szukamy rzadkiego zbioru predyktorów & Lasso & Niestabilna selekcja przy współliniowości.\\
\bottomrule
\end{tabular}
\end{center}

\begin{zadanie}[Dlaczego nie można porównać AIC wprost?]
Analityk dopasował regresję logistyczną do $100$ produktów,
z których $5$ ma braki danych, oraz regresję Poissona do
$95$ obserwacji dotyczących innego rodzaju odpowiedzi.
Pierwszy model ma AIC równe $89$, drugi $76$. Czy drugi
model jest lepszy?
\par\emph{Rozwiązanie.} Nie. Wiarygodności dotyczą innych
zmiennych i potencjalnie różnych zbiorów danych;
różnice AIC nie odpowiadają wspólnej mierze ryzyka.
Najpierw trzeba ustalić ten sam cel predykcji, rozkład
odpowiedzi, obserwacje i stratę. Dla dwóch konkurencyjnych
modeli tego samego wyniku można zastosować wspólną
walidację na identycznych podziałach danych.
\end{zadanie}

\begin{zadanie}[Ridge a nieobciążoność]
Niech $X^\top X=I_p$, a $Y=X\beta+\varepsilon$,
$\E\varepsilon=0$, $\Cov(\varepsilon)=\sigma^2I_n$.
Wyznacz obciążenie ridge oraz MSE wektorowy
$\E\|\widehat\beta_\lambda-\beta\|_2^2$.
\par\emph{Rozwiązanie.} Z twierdzenia~\ref{thm:ridge}
$\E\widehat\beta_\lambda=\beta/(1+\lambda)$,
więc obciążenie wynosi $-\lambda\beta/(1+\lambda)$.
Kowariancja jest równa
$\sigma^2(1+\lambda)^{-2}I_p$, skąd
$\E\|\widehat\beta_\lambda-\beta\|_2^2
=(\lambda^2\|\beta\|_2^2+p\sigma^2)/(1+\lambda)^2$.
Zależnie od nieznanego $\beta$ pewne dodatnie
$\lambda$ mogą poprawiać MSE względem MNK, ale nie
wynika z tego uniwersalna optymalność jednej kary.
\end{zadanie}

% ================================================================
% ROZDZIAL 16 -- DANE ZALEZNE W CZASIE
% ================================================================
\chapter{Wnioskowanie z procesów zależnych: autokorelacja i modele szeregów czasowych}
\label{ch:time-series}
\paragraph{Powiązania i odsyłacze.} Podstawy procesów Markowa i stacjonarności przedstawia tom II (rozdziały~\ref{II-ch:ctmc} i~\ref{II-ch:gauss}); klasyczne modele regresji — rozdział~\ref{ch:linear-regression}.


\section{Dlaczego sto dni to nie zawsze sto niezależnych obserwacji?}
Pomiar temperatury w kolejnych godzinach, dzienne straty
finansowe lub odczyty jednego urządzenia tworzą \emph{szereg czasowy}:
uporządkowane obserwacje jednego procesu. Zamiana ich kolejności
zazwyczaj zmienia sens modelu. Jeżeli zależność w czasie jest
pomijana, nawet poprawna średnia z próby może mieć źle oszacowaną
wariancję, a losowa walidacja krzyżowa może wykorzystywać
informację z przyszłości. Podstawowe pojęcia procesów i
filtracji wprowadziliśmy w tomie II; obecnie naszym zadaniem
jest \emph{estymacja i predykcja} na podstawie jednej trajektorii.

\begin{definicja}[Słaba stacjonarność i autokowariancja]
\label{def:weak-stationarity}
Niech $(X_t)_{t\in\mathbb Z}$ będzie procesem rzeczywistym
o skończonych drugich momentach. Nazywamy go \emph{słabo
stacjonarnym}, jeśli istnieje $\mu\in\R$ takie, że
$\E X_t=\mu$ dla wszystkich $t$, a
$\Cov(X_t,X_s)=\gamma(t-s)$ dla pewnej funkcji
$\gamma:\mathbb Z\to\R$. Wtedy $\gamma(h)$ to
\emph{autokowariancja} opóźnienia $h$, a przy
$\gamma(0)>0$ wielkość $\rho(h)=\gamma(h)/\gamma(0)$
to \emph{autokorelacja}. Zachodzą $\gamma(-h)=\gamma(h)$
i $|\rho(h)|\le1$. Proces ściśle stacjonarny ma
niezmiennicze względem przesunięcia czasu rozkłady
skończeniewymiarowe; przy skończonych drugich momentach
jest też słabo stacjonarny.
\end{definicja}

\begin{uwaga}[Stacjonarność a trend]
Jeśli średnia rośnie w czasie, dane nie są słabo
stacjonarne, choć reszty po właściwym usunięciu trendu
mogą być stacjonarne. Również szeregi sezonowe wymagają
jawnego modelowania kalendarza. Automatyczne różnicowanie
każdej serii bez diagnozy może usunąć użyteczny sygnał.
\end{uwaga}

\begin{twierdzenie}[Dokładna wariancja średniej z szeregu]
\label{thm:mean-time-variance}
Niech $(X_t)_{t\in\mathbb Z}$ będzie słabo stacjonarny
o autokowariancji $\gamma$. Dla
$\overline X_n=n^{-1}\sum_{t=1}^nX_t$ zachodzi
\[
 \Var(\overline X_n)=\frac1{n^2}
 \left(n\gamma(0)+2\sum_{h=1}^{n-1}(n-h)\gamma(h)\right).
\]
Jeśli dodatkowo $\sum_{h\in\mathbb Z}|\gamma(h)|<\infty$,
to
$\lim_{n\to\infty}n\Var(\overline X_n)
=\sum_{h\in\mathbb Z}\gamma(h)=:\tau^2\ge0$.
\end{twierdzenie}
\begin{proof}
Z definicji wariancji sumy
$n^2\Var(\overline X_n)=\sum_{s=1}^n\sum_{t=1}^n
\Cov(X_s,X_t)$. Dla $s=t$ występuje $n$ składników
$\gamma(0)$; dla każdego $h\in\{1,\ldots,n-1\}$
istnieje $n-h$ par o $t-s=h$ i tyle samo o $s-t=h$.
Daje to pierwszy wzór. Mnożąc go przez $1/n$,
otrzymujemy
$n\Var(\overline X_n)=\sum_{|h|<n}
(1-|h|/n)\gamma(h)$. Bezwzględna sumowalność
pozwala zastosować zbieżność zdominowaną dla szeregów.
Nieujemność granicy wynika z nieujemności wariancji.
\end{proof}

\begin{przyklad}[Jak korelacja zmienia efektywną liczebność?]
Jeśli $\gamma(h)=\gamma(0)\phi^{|h|}$ dla $|\phi|<1$,
to $\tau^2=\gamma(0)(1+\phi)/(1-\phi)$.
Dla $\phi=0{,}8$ błąd standardowy średniej w dużej próbie
jest w przybliżeniu trzykrotnie większy niż przy błędnym
założeniu niezależności. Odpowiada to heurystycznej
liczebności efektywnej $n_{\mathrm{eff}}\approx n/9$.
Jest to \emph{porównanie wariancji}, nie dosłowne
stwierdzenie, że usunięto obserwacje. Przy ujemnej
autokorelacji liczebność efektywna może przekroczyć $n$.
\end{przyklad}

\section{Model AR(1): konstrukcja, estymacja i prognoza}
\begin{definicja}[Biały szum i model autoregresyjny]
Mówimy, że $(\varepsilon_t)_{t\in\mathbb Z}$ jest
\emph{białym szumem} o wariancji $\sigma^2>0$, jeśli
$\E\varepsilon_t=0$ oraz
$\Cov(\varepsilon_t,\varepsilon_s)=\sigma^2\ind{t=s}$.
W poniższym twierdzeniu założymy ponadto niezależność
oraz identyczny rozkład innowacji. \emph{AR(1)} opisuje
proces $(X_t)$ spełniający
$X_t-\mu=\phi(X_{t-1}-\mu)+\varepsilon_t$,
gdzie $\mu\in\R$ i $\phi\in\R$ są parametrami.
\end{definicja}

\begin{twierdzenie}[Stacjonarne rozwiązanie AR(1)]
\label{thm:ar1-stationary}
Niech $(\varepsilon_t)_{t\in\mathbb Z}$ będą
niezależnymi zmiennymi o wspólnym rozkładzie,
$\E\varepsilon_t=0$ i
$\Var(\varepsilon_t)=\sigma^2\in(0,\infty)$.
Niech $\mu\in\R$ i $|\phi|<1$. Wtedy szereg
$X_t=\mu+\sum_{j=0}^{\infty}\phi^j\varepsilon_{t-j}$
jest zbieżny w $L^2$ dla każdego $t$ i definiuje
słabo stacjonarne rozwiązanie AR(1) z
\[
 \E X_t=\mu,\qquad
 \gamma(h)=\frac{\sigma^2}{1-\phi^2}\phi^{|h|},
 \quad h\in\mathbb Z.
\]
Jest to jedyne słabo stacjonarne rozwiązanie o skończonym
drugim momencie na tej samej przestrzeni, napędzane tym samym
szumem i spełniające to równanie w $L^2$.
\end{twierdzenie}
\begin{proof}
Niezależność i centrowanie dają dla $m>k$
$\E|\sum_{j=k}^m\phi^j\varepsilon_{t-j}|^2
=\sigma^2\sum_{j=k}^m\phi^{2j}$.
Ponieważ $|\phi|<1$, sumy częściowe są ciągiem
Cauchy'ego w $L^2$. Przesunięcie indeksów pokazuje
$X_t-\mu=\phi(X_{t-1}-\mu)+\varepsilon_t$ w $L^2$.
Dla $h\ge0$ w iloczynie dwóch szeregów przeżywają
jedynie składniki ze wspólnym indeksem innowacji:
$\gamma(h)=\sigma^2\sum_{j\ge0}\phi^{2j+h}
=\sigma^2\phi^h/(1-\phi^2)$. Pozostałe opóźnienia
wynikają z symetrii kowariancji.
Jeżeli $X_t$ i $X'_t$ są dwoma słabo stacjonarnymi
rozwiązaniami z tymi samymi innowacjami, to
$D_t=X_t-X'_t=\phi^mD_{t-m}$ dla każdego $m$.
Słaba stacjonarność obu procesów zapewnia
$\sup_t\E D_t^2<\infty$ przez nierówność
$D_t^2\le2X_t^2+2(X'_t)^2$.
Stąd $\E D_t^2\le\phi^{2m}\sup_s\E D_s^2\to0$.
\end{proof}

\begin{uwaga}[Jaki jest sens parametru $\phi$?]
Dodatnie $\phi$ opisuje utrzymywanie się odchyleń
od poziomu $\mu$, a ujemne naprzemienność znaków.
Dla $\phi=0$ otrzymujemy niezależny szum wokół średniej.
Dla $\phi=1$ zwykły proces z sumowaniem innowacji
jest błądzeniem losowym, a nie stacjonarnym AR(1).
Przy $|\phi|\ge1$ nie istnieje opisane powyżej
przyczynowe stacjonarne rozwiązanie napędzane
jednostronną sumą przeszłych innowacji; nie wyklucza
to nietypowych rozwiązań z informacją o przyszłości.
\end{uwaga}

\begin{twierdzenie}[Prognoza AR(1) przy znanych parametrach]
\label{thm:ar1-forecast}
W założeniach twierdzenia~\ref{thm:ar1-stationary}
niech $\mathcal F_t=\sigma(\varepsilon_s:s\le t)$.
Dla każdego $h\in\N$, $h\ge1$, zachodzi
\[
 \E[X_{t+h}\mid\mathcal F_t]
 =\mu+\phi^h(X_t-\mu),\qquad
 \Var(X_{t+h}\mid\mathcal F_t)
 =\sigma^2\sum_{j=0}^{h-1}\phi^{2j}.
\]
Jeżeli innowacje są normalne, warunkowy rozkład jest
normalny i podane parametry wyznaczają dokładny
przedział predykcji przy \emph{znanych} $\mu,\phi,\sigma$.
\end{twierdzenie}
\begin{proof}
Iteracja rekurencji prowadzi do
$X_{t+h}-\mu=\phi^h(X_t-\mu)+
\sum_{j=0}^{h-1}\phi^j\varepsilon_{t+h-j}$.
Innowacje w ostatniej sumie są niezależne od
$\mathcal F_t$, zatem ich warunkowa średnia
wynosi zero, a wariancja to suma wariancji.
Przy normalności suma jest normalna.
\end{proof}

\begin{stwierdzenie}[Estymacja warunkowa AR(1)]
\label{prop:ar1-conditional-mle}
Załóżmy, że obserwujemy $x_1,\ldots,x_n$, $n\ge3$,
z modelu $X_t=\phi X_{t-1}+\varepsilon_t$ dla $t\ge2$,
gdzie $\varepsilon_t$ są niezależne
$\mathcal N(0,\sigma^2)$ i niezależne od $X_1$.
Przy znanej średniej zerowej i nieznanych
$\phi\in\R$, $\sigma^2>0$, dla każdej ustalonej
$\sigma^2$ log-wiarygodność warunkowa względem $X_1=x_1$
osiąga maksimum po $\phi$ w punkcie
\[
 \widehat\phi_{\mathrm{CLS}}=
 \frac{\sum_{t=2}^n x_tx_{t-1}}{\sum_{t=2}^n x_{t-1}^2},
\]
jeśli mianownik jest dodatni. Estymator ten jest
również estymatorem MNK dla regresji bez wyrazu wolnego.
Jeżeli dodatkowo resztowa suma kwadratów jest dodatnia,
wspólne maksimum po $\phi$ i $\sigma^2>0$ istnieje,
a warunkowy MLE wariancji to $\mathrm{RSS}/(n-1)$.
Jeżeli zaś $\mathrm{RSS}=0$, wiarygodność rośnie
bez ograniczenia przy $\sigma^2\downarrow0$ i wspólny MLE
w otwartej przestrzeni parametrów nie istnieje.
\end{stwierdzenie}
\begin{proof}
Warunkowa log-wiarygodność, z pominięciem stałej,
wynosi $-\frac{n-1}{2}\log v-
\frac{1}{2v}\sum_{t=2}^n(x_t-\phi x_{t-1})^2$.
Dla ustalonego $v$ jej maksimum po $\phi$ odpowiada
minimum wypukłej funkcji kwadratowej; równanie
normalne daje podany iloraz. Maksymalizacja po
$v>0$ odbywa się tak samo jak w twierdzeniu
\ref{thm:aic-gaussian}, z liczebnością $n-1$.
\end{proof}

\begin{uwaga}[Warianty białego szumu]
Biały szum w definicji powyżej wymaga jedynie nieskorelowania
innowacji. Z samego nieskorelowania nie wolno wyprowadzić
niezależności warunkowej, formuły wiarygodności Gaussowskiej
ani prognozy jako pełnego rozkładu warunkowego. W twierdzeniach
o stacjonarnym AR(1) i prognozie założyliśmy dlatego
jawnie niezależność innowacji; do warunkowej wiarygodności
potrzebna jest ponadto ich normalność.
\end{uwaga}

\begin{uwaga}[Warunkowe a dokładne MLE]
Gdy model wymaga $|\phi|<1$, powyższe rozwiązanie
bez ograniczenia może leżeć poza przestrzenią parametrów;
wtedy należy maksymalizować na dopuszczalnej dziedzinie
(w otwartym przedziale supremum może nie być osiągnięte).
Dokładna wiarygodność stacjonarnego normalnego AR(1)
uwzględnia ponadto rozkład $X_1$ i może dawać inny
estymator. Niepewność \emph{oszacowanych} parametrów
nie jest uwzględniona w przedziale z twierdzenia
\ref{thm:ar1-forecast}; w małych próbach ma to znaczenie.
\end{uwaga}

\section{MA, ARMA, ACF i identyfikacja struktury}
\begin{definicja}[Procesy MA($q$) i ARMA($p,q$)]
Niech $(\varepsilon_t)_{t\in\mathbb Z}$ będzie
białym szumem o wariancji $\sigma^2>0$.
Dla $q\in\N\cup\{0\}$ i współczynników
$\theta_1,\ldots,\theta_q\in\R$ proces
$X_t=\mu+\varepsilon_t+
\sum_{j=1}^q\theta_j\varepsilon_{t-j}$
to \emph{MA($q$)}. Wprowadzamy operator opóźnienia
$BX_t=X_{t-1}$ i wielomiany
$\Phi(z)=1-\sum_{i=1}^p\phi_i z^i$ oraz
$\Theta(z)=1+\sum_{j=1}^q\theta_j z^j$.
Model \emph{ARMA($p,q$)} (autoregressive moving average,
czyli połączenie autoregresji i średniej ruchomej) jest równaniem
$\Phi(B)(X_t-\mu)=\Theta(B)\varepsilon_t$.
\emph{Przyczynowość} oznacza możliwość wyrażenia
$X_t-\mu$ jako zbieżnego w $L^2$ szeregu
z przeszłych i bieżących innowacji; \emph{odwracalność}
oznacza analogiczną możliwość odzyskania innowacji
z przeszłych i bieżących wartości procesu.
\end{definicja}

\begin{twierdzenie}[Autokowariancja MA($q$)]
\label{thm:maq-cov}
Niech $X_t-\mu=\sum_{j=0}^q\theta_j\varepsilon_{t-j}$,
$\theta_0=1$, a innowacje będą nieskorelowane,
wycentrowane i mają wspólną wariancję $\sigma^2$.
Wtedy proces jest słabo stacjonarny i dla
$0\le h\le q$ zachodzi
$\gamma(h)=\sigma^2\sum_{j=0}^{q-h}\theta_j\theta_{j+h}$,
natomiast $\gamma(h)=0$ przy $h>q$.
\end{twierdzenie}
\begin{proof}
Rozwijamy
$\Cov(X_t,X_{t-h})=\sum_{j,k=0}^q
\theta_j\theta_k\Cov(\varepsilon_{t-j},\varepsilon_{t-h-k})$.
Dodatnie składniki otrzymujemy tylko, gdy
$j=h+k$, co daje $\sigma^2\sum_{k=0}^{q-h}
\theta_{h+k}\theta_k$. Gdy $h>q$, żadne indeksy
nie mogą być równe. Kowariancja nie zależy od $t$.
\end{proof}

\begin{uwaga}[Charakterystyki ACF i PACF]
\emph{ACF} (autocorrelation function) to funkcja autokorelacji
$\rho(h)$. \emph{PACF} (partial autocorrelation function),
czyli częściowa funkcja autokorelacji, na opóźnieniu $h$
to korelacja liniowych
reszt z regresji $X_t$ i $X_{t-h}$ na pośrednie
$X_{t-1},\ldots,X_{t-h+1}$, jeżeli kowariancja
wektora pośredniego jest odwracalna; dla $h=1$
PACF jest zwykłą autokorelacją opóźnienia $1$.
Dla MA($q$) teoretyczna ACF zeruje się po $q$,
zaś dla AR($p$) w standardowym modelu przyczynowym
teoretyczna PACF zeruje się po $p$. W skończonej
próbie \emph{oszacowane} wartości nie zerują się dokładnie,
więc nie jest to mechaniczna reguła identyfikacji rzędu.
\end{uwaga}

\begin{stwierdzenie}[Warunki przyczynowości i odwracalności]
Jeśli wielomian $\Phi$ modelu ARMA($p,q$) nie ma zer
w kole domkniętym $\{z\in\mathbb C:|z|\le1\}$,
to rozwinięcie $\Theta(z)/\Phi(z)=\sum_{j=0}^\infty\psi_jz^j$
ma współczynniki bezwzględnie sumowalne i definiuje
przyczynowe rozwiązanie
$X_t-\mu=\sum_{j\ge0}\psi_j\varepsilon_{t-j}$
w $L^2$, dla białego szumu o skończonej wariancji.
Jeśli dodatkowo $\Theta$ nie ma zer w tym kole,
a wielomiany $\Phi$ i $\Theta$ nie mają wspólnych czynników,
innowacje można odzyskać przez bezwzględnie sumowalny
filtr z obserwacji (odwracalność).
\end{stwierdzenie}
\begin{proof}
Jeśli wielomian nie zeruje się na kole o promieniu
większym od $1$, odwrotność $1/\Phi(z)$ jest
holomorficzna na tym kole. Szereg Taylora ma wtedy
współczynniki malejące co najmniej geometrycznie
(ze wzoru Cauchy'ego na pochodne), więc jest
bezwzględnie sumowalny. Iloczyn z wielomianem
$\Theta$ zachowuje tę własność. Dla białego szumu
odpowiadający szereg jest Cauchy'ego w $L^2$.
Mnożenie formalnych szeregów, uzasadnione
sumowalnością, odtwarza równanie ARMA.
Dla $\Theta$ stosujemy identyczny argument do
$\Phi/\Theta$, otrzymując filtr odwrotny.
\end{proof}

\section{Diagnostyka autokorelacji i niepewność estymacji}
Dla obserwacji $x_1,\ldots,x_n$ definiujemy
$\widehat\gamma(h)=n^{-1}\sum_{t=h+1}^n
(x_t-\overline x)(x_{t-h}-\overline x)$,
$0\le h<n$, a przy $\widehat\gamma(0)>0$
$\widehat\rho(h)=\widehat\gamma(h)/\widehat\gamma(0)$.
Zwyczajowo rysujemy ACF zarówno oryginalnych danych,
jak i reszt dopasowanego modelu. Użycie wspólnego
mianownika $n$ jest jedną z konwencji; oszacowania ACF
mogą zależeć od wyboru normalizacji i średniej.

\begin{definicja}[Statystyka Ljunga--Boxa]
Dla szeregu długości $n$ i opóźnień $1\le m<n$
przy $\widehat\gamma(0)>0$ statystykę portmanteau
definiujemy wzorem
\[
 Q_m=n(n+2)\sum_{h=1}^{m}\frac{\widehat\rho(h)^2}{n-h}.
\]
Stosujemy ją do badania \emph{łącznej} obecności
autokorelacji do opóźnienia $m$. W klasycznym
przybliżeniu dla reszt dopasowanego ARMA($p,q$)
porównuje się ją z rozkładem $\chi^2_{m-p-q}$,
o ile $m>p+q$ i spełnione są dodatkowe warunki
asymptotyczne modelu i procedury estymacji.
\end{definicja}

\begin{uwaga}[Jak interpretować test Ljunga--Boxa?]
Duża wartość p nie dowodzi braku autokorelacji;
przy małej próbie test ma ograniczoną moc.
Z kolei istotność może wynikać ze zmiennej
wariancji, niestacjonarności lub błędnej specyfikacji.
Wybór $m$ należy uzasadnić częstotliwością danych,
sezonowością i horyzontem prognozy, nie zaś dokonać
go po wielokrotnym przeszukiwaniu wartości p.
\end{uwaga}

\begin{stwierdzenie}[Korekta HAC: sens estymatora wariancji długookresowej]
\label{prop:hac}
Niech $(X_t)$ będzie słabo stacjonarny o bezwzględnie
sumowalnej autokowariancji. Wówczas wariancja
asymptotyczna $\sqrt n(\overline X_n-\mu)$ wynosi
$\tau^2=\sum_{h\in\mathbb Z}\gamma(h)$, w sensie
granicy wariancji z twierdzenia~\ref{thm:mean-time-variance}.
Typową estymatą tej wielkości jest jądrowy estymator
Bartletta--Neweya--Westa; określenie HAC oznacza
\emph{heteroskedasticity and autocorrelation consistent},
czyli odporność na heteroskedastyczność i autokorelację
przy odpowiednich warunkach asymptotycznych:
\[
 \widehat\tau^2_b=\widehat\gamma(0)+
 2\sum_{h=1}^{b}\left(1-\frac{h}{b+1}\right)
 \widehat\gamma(h),\qquad 0\le b<n.
\]
Jego \emph{zgodność} nie wynika z samej stacjonarności:
wymaga dodatkowych warunków zależności, momentów oraz
odpowiednio wolnego wzrostu $b$ z liczebnością próby.
\end{stwierdzenie}

\begin{uwaga}[Wariancja nie jest centralnym twierdzeniem granicznym]
Z równości $n\Var(\overline X_n)\to\tau^2$
nie wynika sama z siebie asymptotyczna normalność.
Potrzebne jest osobne centralne twierdzenie graniczne
dla klasy rozważanych procesów, np. poprawnie określonej
klasy procesów mieszających. Bez niego przedział
$\overline X_n\pm z_{1-\alpha/2}\widehat\tau_b/\sqrt n$
nie ma automatycznej gwarancji pokrycia. Dla
normalnego AR(1) normalność średniej jest dokładna,
ale szacowanie $\tau$ pozostaje dodatkowym problemem.
\end{uwaga}

\section{Praktyczny wybór modelu i uczciwa walidacja w czasie}
Nie dzielimy danych chronologicznych poprzez losowe
przemieszanie, jeśli celem jest prognoza przyszłości.
Jeśli $\Delta X_t=X_t-X_{t-1}$, to model
ARIMA($p,d,q$) zakłada, że $\Delta^d X_t$ spełnia
odpowiedni model ARMA($p,q$); $d$ jest liczbą
różnicowań, a w odmianie sezonowej uwzględnia się
dodatkowo opóźnienia sezonowe. Takiego różnicowania
nie należy wykonywać bez sprawdzenia potrzeby i
przyjętej interpretacji procesu.
Na odłożonym bloku przyszłych obserwacji często mierzymy
RMSE (root mean squared error), czyli pierwiastek
średniego kwadratu błędu prognozy:
$\operatorname{RMSE}=\sqrt{m^{-1}\sum_{j=1}^m(y_j-\widehat y_j)^2}$,
gdzie $m$ jest liczbą sprawdzanych prognoz.
W \emph{walidacji przesuwającego się początku} uczymy
na przeszłości i prognozujemy późniejszy odcinek.
Dwa typowe warianty to poszerzające się okno treningowe
oraz przesuwające się okno o stałej długości.
To drugie może lepiej adaptować się do zmian
warunków, ale korzysta z mniejszej ilości danych.

\begin{center}
\begin{tikzpicture}[x=0.85cm,y=0.53cm,every node/.style={font=\small}]
\foreach \r/\k in {0/4,1/6,2/8} {
  \draw[fill=blue!13,draw=blue!50!black] (0,-\r*1.55) rectangle (\k,-\r*1.55+0.85);
  \draw[fill=orange!22,draw=orange!60!black] (\k,-\r*1.55) rectangle (\k+2,-\r*1.55+0.85);
}
\node[anchor=west] at (8.3,0.48) {trening};
\node[anchor=west] at (8.3,0.02) {prognoza};
\draw[fill=blue!13,draw=blue!50!black] (7.65,0.57) rectangle (8.0,0.79);
\draw[fill=orange!22,draw=orange!60!black] (7.65,0.10) rectangle (8.0,0.32);
\node[anchor=west] at (0,-4.5) {Każde okno testowe następuje po swoim oknie treningowym.};
\end{tikzpicture}
\end{center}

\begin{lstlisting}[language=Python,caption={Kompletny program: symulacja AR(1), estymacja, test reszt i prognozy chronologiczne (Python).}]
import numpy as np
from scipy.stats import norm
from statsmodels.stats.diagnostic import acorr_ljungbox

rng = np.random.default_rng(2026)
n = 260
phi_true, sigma_true = 0.75, 1.0
noise = rng.normal(0.0, sigma_true, size=n + 300)
series = np.zeros(n + 300)
for t in range(1, len(series)):
    series[t] = phi_true * series[t-1] + noise[t]
x = series[300:]  # burn-in: reduces dependence on x[0] = 0

# Chronological holdout (last 40 points) and one-step forecast.
train = x[:220]
test = x[220:]
phi_hat = np.dot(train[1:], train[:-1]) / np.dot(train[:-1], train[:-1])
resid = train[1:] - phi_hat * train[:-1]
sigma_hat = np.sqrt(np.mean(resid**2))
last_seen = np.empty_like(test)
last_seen[0] = train[-1]
last_seen[1:] = test[:-1]
ar_pred = phi_hat * last_seen  # forecast 1 step ahead with updating
rmse_ar = np.sqrt(np.mean((test - ar_pred)**2))
rmse_persist = np.sqrt(np.mean((test - last_seen)**2))
rmse_zero = np.sqrt(np.mean(test**2))
print("estimated AR parameter:", phi_hat)
print("residual standard deviation:", sigma_hat)
print("RMSE AR, persistence, zero:", rmse_ar, rmse_persist, rmse_zero)
print(acorr_ljungbox(resid, lags=[5, 10], return_df=True))

# Forecast from the last available observation at time n-1,
# with parameters fixed. Parameter-estimation uncertainty omitted.
h = 3
mean_h = phi_hat**h * x[-1]
sd_h = sigma_hat * np.sqrt(sum(phi_hat**(2*j) for j in range(h)))
ci_h = (mean_h - norm.ppf(0.975)*sd_h,
        mean_h + norm.ppf(0.975)*sd_h)
print("3-step forecast and approximate interval:", mean_h, ci_h)
\end{lstlisting}

\paragraph{Praktyczna interpretacja.}
Ocenę predykcji wykonujemy wyłącznie na obserwacjach
późniejszych niż trening. Warto porównać model AR(1)
z prognozą stałej średniej, prognozą ostatniej wartości,
a dla danych sezonowych także z prognozą sezonowo naiwną.
Dla ziarna generatora użytego w programie otrzymujemy
$\widehat\phi\approx0{,}731$ i RMSE prognozy AR(1)
około $1{,}014$. Wynik nie gwarantuje przewagi
na innych zbiorach danych.
W programie użyto \emph{aktualizacji} przy prognozie
jednokrokowej: do przewidzenia $X_t$ wolno użyć już
zaobserwowanego $X_{t-1}$, ale nie wartości $X_t$
ani późniejszych. Jeśli wymagamy jednorazowej prognozy
całych $40$ przyszłych okresów, trzeba zastosować
prognozę wielokrokową z \emph{jednego} momentu startowego,
która jest innym zadaniem i zwykle ma większy błąd.

\section{Zależne zdarzenia: model Markowa jako alternatywa dla iid}
\begin{przyklad}[Awaria w kolejnych dniach i trwałe stany]
Niech $S_t\in\{0,1\}$ oznacza stan maszyny w dniu $t$:
$0$ to stan prawidłowy, $1$ to stan alarmowy. Załóżmy,
że $(S_t)$ jest jednorodnym łańcuchem Markowa o macierzy
przejścia
$P=\begin{pmatrix}1-a&a\\ b&1-b\end{pmatrix}$,
gdzie $0<a,b<1$. Prawdopodobieństwo alarmu w stanie
stacjonarnym wynosi $\pi_1=a/(a+b)$, lecz kolejne
alarmy nie są niezależne, chyba że $a+b=1$.
Warunkowa wiarygodność obserwowanej ścieżki
$s_1,\ldots,s_n$ jest proporcjonalna do
$(1-a)^{N_{00}}a^{N_{01}}b^{N_{10}}(1-b)^{N_{11}}$,
gdzie $N_{ij}=\#\{t\in\{1,\ldots,n-1\}:s_t=i,s_{t+1}=j\}$.
W konsekwencji, jeśli odpowiednie mianowniki są dodatnie,
warunkowe MLE wynoszą
$\widehat a=N_{01}/(N_{00}+N_{01})$ i
$\widehat b=N_{10}/(N_{10}+N_{11})$.
Wyprowadzenie wynika z rozdzielenia log-wiarygodności
na dwa składniki dwumianowe; gdy obserwujemy zerową
liczbę przejść danego typu, trzeba uwzględnić brzeg
przestrzeni parametrów. Próba zignorowania kolejności
i dopasowania modelu niezależnych wskaźników Bernoulliego
prowadziłaby do innego modelu niepewności.
\end{przyklad}

\begin{stwierdzenie}[Korelacja w dwustanowym łańcuchu]
\label{prop:markov-corr}
W powyższym łańcuchu przy rozkładzie początkowym
stacjonarnym zachodzi
$\rho(h)=(1-a-b)^h$ dla $h\ge0$.
\end{stwierdzenie}
\begin{proof}
Oznaczmy $\lambda=1-a-b$.
Dla $s\in\{0,1\}$ mamy
$\E[S_{t+1}\mid S_t=s]=a+(1-a-b)s
=\pi_1+\lambda(s-\pi_1)$.
Iteracja własności Markowa daje
$\E[S_{t+h}-\pi_1\mid S_t]=\lambda^h(S_t-\pi_1)$.
Stąd
$\Cov(S_t,S_{t+h})=
\E[(S_t-\pi_1)\E[S_{t+h}-\pi_1\mid S_t]]
=\lambda^h\Var(S_t)$, co kończy dowód.
\end{proof}

\section{Dobór narzędzia i błędy praktyczne}
\begin{center}\small
\begin{tabular}{>{\raggedright\arraybackslash}p{.25\textwidth}>{\raggedright\arraybackslash}p{.32\textwidth}>{\raggedright\arraybackslash}p{.30\textwidth}}
\toprule
\textbf{Sytuacja} & \textbf{Metoda do rozważenia} & \textbf{Co zweryfikować}\\
\midrule
Jedna trajektoria o stałej średniej, krótka pamięć & AR/MA/ARMA & Stacjonarność, rzędy i ACF reszt.\\
Trend i sezonowość & Regresja z czasem, sezonowe ARIMA & Brak przecieku i stabilność reżimu.\\
Średnia danych zależnych & Długookresowy SE, HAC, model procesu & Sumowalność korelacji i CLT.\\
Dwie wartości stanu, istotna kolejność & Łańcuch Markowa & Jednorodność przejść, liczba przejść.\\
Zliczenia zdarzeń w kolejnych okresach & Poisson tylko przy uzasadnionej niezależności; modele zależnych zliczeń w pozostałych sytuacjach & Sezonowość, nadmierna dyspersja, korelacja.\\
Prognoza w czasie & Rolling origin, odłożona końcówka szeregu & Wyłącznie historyczne cechy, właściwy horyzont.\\
\bottomrule
\end{tabular}
\end{center}

\begin{zadanie}[Dlaczego okna mogą być skorelowane mimo procesu Poissona?]
Niech $(N_t)_{t\ge0}$ będzie jednorodnym procesem
Poissona o intensywności $\lambda>0$. Definiujemy
$A=N_2-N_0$ i $B=N_3-N_1$. Czy $A$ i $B$ są niezależne?
\par\emph{Rozwiązanie.} Nie. Zapiszmy
$A=(N_1-N_0)+(N_2-N_1)$ oraz
$B=(N_2-N_1)+(N_3-N_2)$.
Rozłączne przyrosty są niezależne, a wspólny składnik
$N_2-N_1\sim\Pois(\lambda)$ ma wariancję $\lambda$.
Stąd $\Cov(A,B)=\lambda>0$.
Model niezależnych \emph{nachodzących na siebie} okien
zliczeń byłby błędny, chociaż bazowy proces ma
niezależne przyrosty na rozłącznych przedziałach.
\end{zadanie}

\begin{zadanie}[Dobór błędu standardowego]
Dla $n=200$ kolejnych obserwacji stabilnego procesu
oszacowano autokorelację jako w przybliżeniu
$\rho(h)=0{,}6^{|h|}$. Analityk obliczył błąd
standardowy średniej $s/\sqrt n$. Jaką poprawkę
motywuje model AR(1)?
\par\emph{Rozwiązanie.} W przybliżeniu
$\tau^2/\gamma(0)=(1+0{,}6)/(1-0{,}6)=4$.
Błąd standardowy średniej jest około dwa razy
większy od klasycznego, o ile model i przybliżenie
asymptotyczne są odpowiednie. W analizie danych należy
oszacować parametr, wariancję długookresową i ocenić
niepewność tych oszacowań; nie wolno stosować
tego mnożnika bez sprawdzenia struktury zależności.
\end{zadanie}


% ================================================================
% ROZDZIAL 17 -- WNIOSKOWANIE BAYESOWSKIE
% ================================================================
\chapter{Podstawy wnioskowania bayesowskiego: model, prior i posterior}
\label{ch:bayes-foundations}
\paragraph{Powiązania i odsyłacze.} Formalne warunkowanie omawia tom I, rozdział~\ref{I-ch:condexp}; funkcję wiarygodności i teorię decyzji — rozdziały~\ref{ch:likelihood} i~\ref{ch:decision}.


\section{Od wiarygodności do rozkładu warunkowego parametru}
W modelu klasycznym parametr $\theta$ jest ustalony, ale nieznany,
a rozkład danych $P_\theta$ służy do oceny procedury estymacji lub
testu przy każdym $\theta$. W podejściu bayesowskim określamy dodatkowo
\emph{rozkład a priori} $\pi$ na przestrzeni parametrów. Wtedy parametr
$\Theta$ jest zmienną losową w modelu wspólnym, a obserwacja $X=x$
aktualizuje naszą informację o jego wartości. Samo przypisanie prioru
nie dowodzi, że wybrany model opisuje rzeczywistość; w szczególności
niedopasowane prawdopodobieństwo obserwacji może prowadzić do mylącego
posterioru niezależnie od poprawności rachunków.

\begin{definicja}[Eksperyment bayesowski i rozkład a posteriori]
\label{def:bayes-model}
Niech $(\Theta_0,\mathcal T)$ będzie przestrzenią parametrów,
$(\mathsf X,\mathcal X)$ przestrzenią obserwacji, $\pi$ miarą
probabilistyczną na $(\Theta_0,\mathcal T)$, a
$K:\Theta_0\times\mathcal X\to[0,1]$ jądrem probabilistycznym:
$K(\theta,\cdot)$ jest rozkładem obserwacji przy parametrze $\theta$,
a $\theta\mapsto K(\theta,A)$ jest mierzalna dla każdego $A\in\mathcal X$.
\emph{Rozkład a priori} to $\pi$, a \emph{rozkład predykcyjny przed
obserwacją} jest miarą $M(A)=\int_{\Theta_0}K(\theta,A)\,\dd\pi(\theta)$.
\emph{Rozkładem a posteriori} nazywamy jądro probabilistyczne
$\Pi:\mathsf X\times\mathcal T\to[0,1]$ spełniające dla wszystkich
$B\in\mathcal T$, $A\in\mathcal X$ tożsamość
\[
 \int_A\Pi(x,B)\,\dd M(x)
 =\int_B K(\theta,A)\,\dd\pi(\theta).
\]
Jest to regularny rozkład warunkowy $\mathcal L(\Theta\mid X=x)$
w modelu wspólnym $\dd Q(\theta,x)=\dd\pi(\theta)K(\theta,\dd x)$.
Jeżeli obie przestrzenie są standardowymi przestrzeniami borelowskimi,
taki rozkład istnieje; jego wersja jest jednoznaczna tylko $M$-prawie
wszędzie. Twierdzenie o istnieniu regularnych rozkładów warunkowych
przyjmujemy z teorii miary, a w poniższych modelach podajemy ich
konstrukcję jawną.
\end{definicja}

\begin{twierdzenie}[Wzór Bayesa dla modelu zdominowanego]
\label{thm:bayes-dominated}
Niech $(\Theta_0,\mathcal T)$ i $(\mathsf X,\mathcal X)$ będą
standardowymi przestrzeniami borelowskimi. Niech $\pi$ będzie
rozkładem a priori, $\nu$ miarą $\sigma$-skończoną na $\mathsf X$,
a $p:\Theta_0\times\mathsf X\to[0,\infty)$ wspólnie mierzalną
funkcją taką, że
$K(\theta,A)=\int_A p(x\mid\theta)\,\dd\nu(x)$
jest miarą probabilistyczną dla każdego $\theta\in\Theta_0$.
Niech $m(x)=\int_{\Theta_0}p(x\mid\theta)\,\dd\pi(\theta)$.
Dla $x$ takich, że $0<m(x)<\infty$, wzór
\[
 \Pi(x,B)=\frac{\int_B p(x\mid\theta)\,\dd\pi(\theta)}{m(x)},
 \qquad B\in\mathcal T,
\]
definiuje rozkład a posteriori. Na pozostałym zbiorze $x$ możemy
wybrać dowolną ustaloną miarę probabilistyczną na $\Theta_0$;
wybór ten nie zmienia tożsamości warunkowania.
\end{twierdzenie}
\begin{proof}
Z twierdzenia Tonellego $m$ jest mierzalna i
$\int_{\mathsf X}m(x)\,\dd\nu(x)=
\int_{\Theta_0}\int_{\mathsf X}p(x\mid\theta)\,\dd\nu(x)\,\dd\pi(\theta)=1$.
W szczególności $m(x)$ jest dodatnia i skończona dla
$M$-prawie każdego $x$. Dla takich $x$ licznik jest miarą w
zmiennej $B$, a mianownik normalizuje ją do jedności.
Mierzalność $x\mapsto\Pi(x,B)$ wynika z Tonellego i mierzalności
$p$. Dla $A\in\mathcal X$ i $B\in\mathcal T$ otrzymujemy
\[
 \int_A\Pi(x,B)m(x)\,\dd\nu(x)
 =\int_A\int_B p(x\mid\theta)\,\dd\pi(\theta)\,\dd\nu(x)
 =\int_BK(\theta,A)\,\dd\pi(\theta),
\]
co jest definicją rozkładu warunkowego. Zmiany na zbiorze
$M$-miary zero nie mają wpływu na tę tożsamość.
\end{proof}

\begin{uwaga}[Gęstość posterioru i mianownik]
Jeśli $\pi$ ma gęstość $q(\theta)$ względem miary $\mu$,
to gęstość posterioru ma postać
$q(\theta\mid x)=p(x\mid\theta)q(\theta)/m(x)$,
gdzie $m(x)=\int_{\Theta_0}p(x\mid\theta)q(\theta)\,\dd\mu(\theta)$.
Oznaczenie $p(x\mid\theta)$ to \emph{gęstość obserwacji}, a nie
prawdopodobieństwo konkretnego punktu, jeśli $X$ jest ciągła.
W jednym modelu $m(x)$ jest też gęstością predykcyjną obserwacji;
w porównywaniu modeli służy jako \emph{wiarygodność zintegrowana}
(ang. \emph{marginal likelihood}, nazywana również \emph{evidence}).
Zapis $\Pi(\dd\theta\mid x)\propto p(x\mid\theta)\pi(\dd\theta)$
ma sens dopiero po sprawdzeniu, że całka normalizująca jest dodatnia
i skończona. Symbol $\propto$ oznacza równość z dokładnością do
\emph{dodatniego} czynnika niezależnego od $\theta$.
\end{uwaga}

\section{Rozkład a priori: informacja, wrażliwość i uczciwość modelowania}
Rozkład a priori może kodować rzeczywiste wcześniejsze informacje,
wiedzę o skali parametru lub ograniczenia fizyczne. Może też pełnić rolę
regularizacji. Należy jednak pamiętać, że „słaby” prior nie musi być
nieinformatywny po zmianie parametryzacji, a szeroki prior bywa
bardzo informatywny przy porównywaniu modeli.

\begin{przyklad}[Dwa priory, te same dane]
Niech $X_1,\ldots,X_n\mid p\iid\operatorname{Bernoulli}(p)$,
$p\in(0,1)$ i $S=\sum_{i=1}^nX_i$. Dla prioru $p\sim
\operatorname{Beta}(a,b)$ posterior to $\operatorname{Beta}(a+S,b+n-S)$;
ścisły dowód podamy w rozdziale~\ref{ch:conjugate-families}.
Przy $S=7$, $n=100$ prior jednostajny $\operatorname{Beta}(1,1)$
daje średnią posterioru $8/102\approx0{,}0784$, a prior
$\operatorname{Beta}(20,80)$ daje $27/200=0{,}135$.
Oba wyniki są poprawne w \emph{różnych} modelach wspólnych, ale
znacząco się różnią. W raporcie należy więc podawać prior i
analizę wrażliwości, a nie tylko jedną liczbę „bayesowską”.
\end{przyklad}

\begin{definicja}[Prior niewłaściwy i poprawność posterioru]
Niech $\mu$ będzie miarą $\sigma$-skończoną na przestrzeni
parametrów $\Theta_0$, lecz $\mu(\Theta_0)=\infty$.
Jeśli dla obserwacji $x$ zachodzi
$0<Z_\mu(x)=\int_{\Theta_0}p(x\mid\theta)\,\dd\mu(\theta)<\infty$,
to $\Pi_\mu(x,B)=Z_\mu(x)^{-1}\int_B p(x\mid\theta)\,\dd\mu(\theta)$
jest \emph{poprawnym rozkładem a posteriori z prioru niewłaściwego}.
W tym przypadku $\mu$ nie jest rozkładem prawdopodobieństwa,
a obliczenie jest formalną normalizacją, nie warunkowaniem we wspólnym
rozkładzie probabilistycznym parametru i obserwacji. Nie wolno z tego
wnioskować o prawidłowej zintegrowanej wiarygodności modelu.
\end{definicja}

\begin{przyklad}[Prior płaski: posterior poprawny albo niewłaściwy]
W modelu $X_1,\ldots,X_n\mid\mu\iid\mathcal N(\mu,\sigma^2)$
przy znanym $\sigma^2>0$, $n\ge1$, formalny prior
$\dd\lambda_1(\mu)$ na $\R$ daje posterior
$\mathcal N(\overline X,\sigma^2/n)$, ponieważ w funkcji $\mu$
wiarygodność jest proporcjonalna do gęstości normalnej.
Natomiast przy priorze $\dd\lambda_1(\theta)$ na $\R$ i danych
$X\mid\theta\sim\operatorname{Bernoulli}(g(\theta))$, gdzie
$g(\theta)=(1+e^{-\theta})^{-1}$, dla pojedynczej obserwacji
$X=1$ mamy $\int_\R g(\theta)\,\dd\lambda_1(\theta)=\infty$:
posterior nie istnieje. Utrzymujące się dodatnie prawdopodobieństwo
na ogonie $\theta\to+\infty$ uniemożliwia normalizację.
Pierwszy z przykładów dotyczy posterioru dla parametrów przy
ustalonym $\sigma$; nie można rozszerzyć wniosku bez sprawdzenia
całkowalności na model ze wspólnie nieznanymi $\mu,\sigma$.
\end{przyklad}

\begin{uwaga}[Dlaczego niewłaściwy prior nie daje zwykle ilorazu Bayesa?]
Jeśli $\mu$ jest niewłaściwa, zastąpienie jej przez $c\mu$, $c>0$,
nie zmienia posterioru, ale mnoży całkę $Z_\mu(x)$ przez $c$.
Gdy porównujemy dwa modele z \emph{oddzielnymi} nieustalonymi
stałymi, ich iloraz zintegrowanych wiarygodności staje się
arbitralny. Istnieją wyspecjalizowane konstrukcje eliminujące część
tych trudności, ale prosty iloraz nie jest wtedy poprawną procedurą.
\end{uwaga}

\section{Estymator Bayesa i decyzja optymalna względem posterioru}
Po obserwacji $x$ możemy podejmować decyzję $a$ z przestrzeni
$\mathsf A$. Nie ma uniwersalnej „najlepszej” punktowej reprezentacji
posterioru: zależy ona od funkcji straty.

\begin{twierdzenie}[Średnia, mediana i dominanta posterioru]
\label{thm:bayes-actions}
Niech $\Pi_x$ będzie rozkładem probabilistycznym na $\R$.
\begin{enumerate}
\item Jeżeli $\int_\R\theta^2\,\dd\Pi_x(\theta)<\infty$, to
jedynym minimizatorem $\int_\R(\theta-a)^2\,\dd\Pi_x(\theta)$
po $a\in\R$ jest $a=\int_\R\theta\,\dd\Pi_x(\theta)$.
\item Jeżeli $\int_\R|\theta|\,\dd\Pi_x(\theta)<\infty$,
to zbiór minimizatorów $\int_\R|\theta-a|\,\dd\Pi_x(\theta)$
po $a\in\R$ jest zbiorem median $\Pi_x$, tj. tych $a$, dla których
$\Pi_x(({-}\infty,a])\ge1/2$ i $\Pi_x([a,\infty))\ge1/2$.
\end{enumerate}
Dominanta (MAP, \emph{maximum a posteriori}) nie jest ogólnie
rozwiązaniem żadnego z tych dwóch problemów. Jeśli rozkład ma gęstość
względem ustalonej miary, MAP oznacza punkt maksymalizujący tę gęstość,
a wynik zależy od wybranej parametryzacji i miary odniesienia.
\end{twierdzenie}
\begin{proof}
Dla pierwszej części oznaczmy $m=\int_\R\theta\,\dd\Pi_x(\theta)$.
Z tożsamości $(\theta-a)^2=(\theta-m)^2+(m-a)^2
+2(\theta-m)(m-a)$ i centrowania otrzymujemy
$\int_\R(\theta-a)^2\,\dd\Pi_x(\theta)
=\int_\R(\theta-m)^2\,\dd\Pi_x(\theta)+(m-a)^2$.
W drugiej części niech $F$ będzie dystrybuantą $\Pi_x$.
Dla $a<b$ rachunek różnicy funkcji wypukłej daje
$\int_\R (|\theta-b|-|\theta-a|)\,\dd\Pi_x(\theta)
=\int_a^b(2F(t)-1)\,\dd\lambda_1(t)$;
wartość dystrybuanty w punktach skoku nie ma wpływu na całkę.
Zatem funkcja jest niemalejąca na prawo od każdej mediany
i nierosnąca na lewo; odwrotnie, gdyby punkt nie był medianą,
jedna z jednostronnych pochodnych wskazywałaby kierunek spadku.
\end{proof}

\begin{przyklad}[Próg decyzji zależy od kosztów]
Po obserwacji danych interesuje nas zdarzenie $H=\{\theta>\theta_0\}$
i jego prawdopodobieństwo a posteriori $q=\Pi_x(H)$.
Są dwie decyzje: $a=1$ (ogłaszamy zagrożenie) i $a=0$ (brak alarmu).
Niech fałszywy alarm kosztuje $C_{\rm FP}>0$, a przeoczenie
$C_{\rm FN}>0$; poprawna decyzja ma stratę zerową.
Ryzyko posteriorowe alarmu to $C_{\rm FP}(1-q)$,
a milczenia $C_{\rm FN}q$. Wybieramy alarm wtedy i tylko wtedy, gdy
$q>C_{\rm FP}/(C_{\rm FP}+C_{\rm FN})$;
przy równości obie decyzje są równoważne.
W zastosowaniu nie zastępujemy jednak przemyślanej funkcji kosztów
uniwersalnym progiem $0{,}95$.
\end{przyklad}

\begin{przyklad}[MAP nie jest niezmienniczy względem zmiany parametru]
Niech posterior parametru $p\in(0,1)$ ma rozkład
$\operatorname{Beta}(2,2)$ o gęstości $f(p)=6p(1-p)$.
Dominanta w parametryzacji $p$ wynosi $p_{\rm MAP}=1/2$.
Wprowadzamy jednak iloraz szans $u=p/(1-p)\in(0,\infty)$.
Z wzoru na zmianę zmiennych gęstość $u$ jest równa
$g(u)=6u/(1+u)^4$. Jej pochodna ma znak $1-3u$,
więc maksimum przypada na $u_{\rm MAP}=1/3$,
co odpowiada $p=1/4$, a nie $1/2$.
Średnia posterioru i prawdopodobieństwo zdarzenia są
jednoznacznymi funkcjami rozkładu, ale dominanta
\emph{gęstości względem konkretnej miary} zależy od współrzędnych.
\end{przyklad}

\begin{stwierdzenie}[Aktualizacja sekwencyjna bez ponownego używania danych]
\label{prop:sequential-bayes}
Niech prior $\pi$ będzie właściwy, a $X$ i $Y$ przy znanym
$\theta$ będą warunkowo niezależne i mają gęstości
$p(x\mid\theta)$ oraz $r(y\mid\theta)$ względem ustalonych
miar $\sigma$-skończonych. Załóżmy, że dla danych $(x,y)$
wszystkie mianowniki aktualizacji są dodatnie i skończone.
Wtedy posterior uzyskany najpierw z $x$, a później z $y$,
jest identyczny z posteriorem uzyskanym jednorazowo z $(x,y)$.
\end{stwierdzenie}
\begin{proof}
Pierwsza aktualizacja daje miarę
$\Pi_x(\dd\theta)=p(x\mid\theta)\pi(\dd\theta)/m(x)$.
Druga, z wykorzystaniem wyłącznie nowej wiarygodności,
daje jądro proporcjonalne do
$r(y\mid\theta)p(x\mid\theta)\pi(\dd\theta)$.
Ponieważ warunkowa niezależność zapewnia, że wiarygodność
połączonych danych jest iloczynem $p(x\mid\theta)r(y\mid\theta)$,
po normalizacji otrzymujemy ten sam posterior.
Jeżeli jednak $Y$ jest funkcją wcześniej użytych danych albo
zachodzi zależność warunkowa, takie mnożenie byłoby niepoprawne.
\end{proof}

\section{Przedziały wiarygodności i ich interpretacja}
\begin{definicja}[Zbiór wiarygodności i HPD]
Niech $\Pi_x$ będzie rozkładem posterioru na $\Theta_0$.
Zbiór mierzalny $C_x\subset\Theta_0$ taki, że
$\Pi_x(C_x)\ge1-\alpha$, $0<\alpha<1$, nazywamy
\emph{zbiorem wiarygodności na poziomie $1-\alpha$}.
W przypadku rozkładu na $\R$, przy ciągłej i ściśle rosnącej
dystrybuancie $F_x$, \emph{centralny przedział wiarygodności}
ma końce $F_x^{-1}(\alpha/2)$ i $F_x^{-1}(1-\alpha/2)$.
Jeżeli posterior ma gęstość $q_x$ względem ustalonej miary $\mu$,
zbiór \emph{HPD} (highest posterior density) przy odpowiednim
progu $c\ge0$ ma postać $\{\theta:q_x(\theta)\ge c\}$
i posteriorową masę co najmniej $1-\alpha$. Nie zawsze istnieje
próg dający równość, a zbiór HPD może być niespójny.
Przy płaskich fragmentach gęstości próg może dawać wiele równoważnych
wyborów zbioru: nie należy wówczas ogłaszać jedynego HPD bez
ustalenia reguły rozstrzygania remisów. Własność HPD zależy ponadto
od miary odniesienia i nie jest niezmiennicza względem dowolnej
nieliniowej reparametryzacji.
\end{definicja}

\begin{uwaga}[Dwie odmienne gwarancje]
Klasyczny zbiór ufności $C(X)$ ma gwarancję
$P_\theta(\theta\in C(X))\ge1-\alpha$ dla każdego $\theta$
w określonym modelu. Bayesowski zbiór wiarygodności $C_x$ ma
$\Pi_x(C_x)\ge1-\alpha$ \emph{po zaobserwowaniu $x$}.
Gwarancja bayesowska nie implikuje automatycznie gwarancji pokrycia
częstościowego dla każdego ustalonego $\theta$. Przy poprawnym
modelu wspólnym ma jednak interpretację średniego pokrycia
względem prioru, co zaraz udowodnimy.
\end{uwaga}

\begin{twierdzenie}[Średnie priorowe pokrycie zbiorów wiarygodności]
\label{thm:prior-coverage}
Niech $(\Theta,X)$ ma wspólny rozkład $Q$ z definicji
\ref{def:bayes-model}, a $C:\mathsf X\to\mathcal T$ będzie
mierzalną wielowartościową regułą w sensie, że
$\{(x,\theta):\theta\in C(x)\}\in\mathcal X\otimes\mathcal T$.
Jeśli $\Pi(X,C(X))\ge1-\alpha$ $M$-prawie na pewno,
to $Q(\Theta\in C(X))\ge1-\alpha$. Jeśli masa zbioru
jest równa $1-\alpha$ prawie na pewno, zachodzi równość.
\end{twierdzenie}
\begin{proof}
Z definicji rozkładu warunkowego i aproksymacji funkcji
mierzalnych funkcjami prostymi
$\E_Q[\ind{\Theta\in C(X)}\mid X]=\Pi(X,C(X))$
prawie na pewno. Całkując obie strony względem wspólnego prawa,
otrzymujemy $Q(\Theta\in C(X))=\E_M\Pi(X,C(X))\ge1-\alpha$.
Zapis jest tożsamością uśrednioną po priorze, a nie gwarancją
jednostajną w $\theta$.
\end{proof}

\section{Posterior predictive: prognozowanie, nie tylko estymacja}
\begin{definicja}[Rozkład predykcyjny a posteriori]
Niech warunkowo względem $\Theta=\theta$ obserwacja przyszła
$Y$ ma jądro $L(\theta,\cdot)$, a $X$ i $Y$ są warunkowo
niezależne przy znanym $\Theta$. Wówczas rozkład
\emph{predykcyjny a posteriori} wynosi
$P(Y\in A\mid X=x)=\int_{\Theta_0}L(\theta,A)\,\dd\Pi_x(\theta)$.
Dla liczbowego $Y$ o skończonym drugim momencie zachodzi
\[
 \Var(Y\mid X)=\E[\Var(Y\mid\Theta,X)\mid X]
 +\Var(\E[Y\mid\Theta,X]\mid X).
\]
Pierwszy składnik to niepewność przyszłej obserwacji przy ustalonym
parametrze; drugi uwzględnia niepewność samego parametru.
\end{definicja}

\begin{przyklad}[Dlaczego prognoza ma szerszą niepewność?]
Dla $X_i\mid\mu\iid\mathcal N(\mu,\sigma^2)$ przy znanym
$\sigma^2$ i normalnym posteriorze $\mu\mid X=x\sim
\mathcal N(m_n,v_n)$ przyszła obserwacja $Y\mid\mu\sim
\mathcal N(\mu,\sigma^2)$, warunkowo niezależna od danych,
ma rozkład predykcyjny $Y\mid X=x\sim\mathcal N(m_n,\sigma^2+v_n)$.
Przedział dla przyszłego pomiaru musi więc uwzględniać zarówno
$\sigma^2$, jak i $v_n$; użycie samego posterioru parametru
sztucznie zaniżyłoby niepewność prognozy.
\end{przyklad}

\begin{center}
\begin{minipage}{\textwidth}
\centering
\begin{tikzpicture}[
 node distance=1.3cm and 1.5cm,
 every node/.style={font=\small,align=center},
 block/.style={draw=blue!45!black,rounded corners=3pt,
 fill=blue!5,minimum height=9mm,text width=4.1cm},
 source/.style={draw=gray!60,rounded corners=3pt,
 fill=gray!6,minimum height=9mm,text width=4.1cm},
 >={Stealth[length=2.4mm]}]
\node[source] (prior) at (-3,1.8) {Prior: $\pi(\dd\theta)$};
\node[source] (data) at (3,1.8) {Dane: $p(x\mid\theta)$};
\node[block] (post) at (0,0.25) {Posterior: $\Pi(\dd\theta\mid x)$};
\node[source] (predict) at (-3,-1.35) {Predykcja nowych danych};
\node[source] (decision) at (3,-1.35) {Decyzja i funkcja straty};
\draw[->,thick] (prior.south) -- (post.north west);
\draw[->,thick] (data.south) -- (post.north east);
\draw[->,thick] (post.south west) -- (predict.north);
\draw[->,thick] (post.south east) -- (decision.north);
\end{tikzpicture}

{\small Schemat wnioskowania: obliczenie posterioru nie kończy analizy;
prognoza i decyzja rozwiązują odrębne zadania.}
\end{minipage}
\end{center}

\section{Jak korzystać z metod bayesowskich w praktyce?}
\begin{center}\small
\begin{tabular}{>{\raggedright\arraybackslash}p{.25\textwidth}>{\raggedright\arraybackslash}p{.32\textwidth}>{\raggedright\arraybackslash}p{.30\textwidth}}
\toprule
\textbf{Problem} & \textbf{Wybór} & \textbf{Sprawdzenie / ograniczenie}\\
\midrule
Znane wcześniejsze informacje & Prior uzasadniony danymi zewnętrznymi & Czy informacje nie dublują bieżącej próby?\\
Mało danych, stabilne parametry & Model sprzężony i analiza wrażliwości & Jak silnie wynik zależy od prioru?\\
Parametry rzadkich zdarzeń & Posterior i rozkład predykcyjny & Czy model Bernoulliego lub Poissona jest właściwy?\\
Porównanie modeli & Właściwe priory, zintegrowane wiarygodności & Niewłaściwy prior uniemożliwia prosty Bayes factor.\\
Prognoza nowych obserwacji & Posterior predictive & Rozróżnić zmienność obserwacji i niepewność parametru.\\
Decyzja przy asymetrycznych kosztach & Reguła minimalizująca stratę posteriorową & Jak wyznaczono koszty i konsekwencje błędów?\\
\bottomrule
\end{tabular}
\end{center}

\begin{zadanie}[Ryzyko alarmu i próg decyzji]
Załóżmy, że posteriorowe prawdopodobieństwo awarii wynosi $q=0{,}08$.
Fałszywy alarm kosztuje $3$ jednostki, a przeoczenie awarii $50$.
Jaką decyzję rekomenduje reguła Bayesa?
\par\emph{Rozwiązanie.} Alarm kosztuje średnio
$3(1-0{,}08)=2{,}76$ jednostki, brak alarmu
$50\cdot0{,}08=4$. Zatem alarm minimalizuje ryzyko
posteriorowe, choć prawdopodobieństwo awarii jest niskie.
\end{zadanie}

% ================================================================
% ROZDZIAL 18 -- RODZINY SPRZEZONE
% ================================================================
\chapter{Rodziny sprzężone: rachunek posterioru i prognoza}
\label{ch:conjugate-families}
\paragraph{Powiązania i odsyłacze.} Wzór Bayesa i warunki normalizacji wprowadziliśmy w rozdziale~\ref{ch:bayes-foundations}; predykcję rozwijamy w rozdziale~\ref{ch:bayes-estimation-prediction}.


\section{Dlaczego sprzężoność jest użyteczna?}
Sprzężoność pozwala przekształcić wzór Bayesa w jawny rachunek:
po obserwacji danych rozkład a posteriori pozostaje w tej samej
rodzinie co prior. W praktyce otrzymujemy dzięki temu proste
prognozy, przedziały wiarygodności i wygodne aktualizowanie wiedzy
po napływie nowych obserwacji. Nie jest to jednak argument za
wyborem rodziny jako \emph{prawdziwego} modelu danych:
wygoda obliczeniowa nie zastępuje diagnostyki ani analizy wrażliwości.

\begin{definicja}[Rodzina sprzężona]
Dla jądra wiarygodności $p(x\mid\theta)$ na
$\Theta_0$ rodzinę \emph{właściwych} priorów
$\{\pi_\eta:\eta\in H\}$ nazywamy sprzężoną, jeśli dla
każdego dopuszczalnego $\eta$ i każdego $x$ o dodatniej,
skończonej gęstości predykcyjnej rozkład a posteriori jest
równy $\pi_{T(\eta,x)}$ dla pewnego odwzorowania aktualizacji
$T$. Sprzężoność określa parę (model wiarygodności, prior),
a nie własność samej rodziny priorów.
\end{definicja}

\section{Beta--Bernoulli i beta--dwumianowy: szacowanie odsetka}
\begin{definicja}[Rozkład beta]
Dla $a,b>0$ rozkład $\operatorname{Beta}(a,b)$ na $(0,1)$
ma gęstość
\[
 f_{a,b}(p)=\frac{p^{a-1}(1-p)^{b-1}}{B(a,b)},
 \qquad B(a,b)=\int_0^1u^{a-1}(1-u)^{b-1}\,\dd\lambda_1(u)
 =\frac{\Gamma(a)\Gamma(b)}{\Gamma(a+b)}.
\]
Funkcja $\Gamma(t)=\int_0^\infty u^{t-1}e^{-u}\,\dd\lambda_1(u)$
dla $t>0$ jest funkcją gamma. Tożsamość beta--gamma wynika
ze zmiany zmiennych w podwójnej całce gamma; wykorzystamy
ją do obliczania momentów.
\end{definicja}

\begin{lemat}[Tożsamość beta--gamma]
Dla wszystkich $a,b>0$ zachodzi
$B(a,b)=\Gamma(a)\Gamma(b)/\Gamma(a+b)$.
\end{lemat}
\begin{proof}
Korzystając z twierdzenia Tonellego, zapisujemy iloczyn
$\Gamma(a)\Gamma(b)$ jako całkę funkcji
$x^{a-1}y^{b-1}e^{-x-y}$ po $(0,\infty)^2$
względem $\lambda_2$. Podstawiamy $x=ru$, $y=r(1-u)$
dla $r>0$, $u\in(0,1)$. Jakobian ma wartość $r$,
więc nowa całka jest iloczynem
$\int_0^\infty r^{a+b-1}e^{-r}\,\dd\lambda_1(r)$
i $\int_0^1u^{a-1}(1-u)^{b-1}\,\dd\lambda_1(u)$.
Pierwsza całka równa się $\Gamma(a+b)$, druga $B(a,b)$.
\end{proof}

\begin{twierdzenie}[Sprzężoność i rozkład predykcyjny beta--dwumianowy]
\label{thm:beta-binomial}
Niech $p\sim\operatorname{Beta}(a,b)$, $a,b>0$, a warunkowo
względem $p$ zmienne $X_1,\ldots,X_n$ będą niezależne
i $X_i\mid p\sim\operatorname{Bernoulli}(p)$, gdzie $n\ge1$.
Niech $S=\sum_{i=1}^nX_i$ i $s\in\{0,\ldots,n\}$.
Wtedy
\[
 p\mid S=s\sim\operatorname{Beta}(a+s,b+n-s),\qquad
 \E[p\mid S=s]=\frac{a+s}{a+b+n},
\]
\[
 \Var(p\mid S=s)=\frac{(a+s)(b+n-s)}{(a+b+n)^2(a+b+n+1)}.
\]
Jeżeli po próbie otrzymamy $m\ge1$ nowych, warunkowo niezależnych
obserwacji Bernoulliego o tym samym parametrze i
$T$ oznacza liczbę sukcesów, to dla $j=0,\ldots,m$
\[
 \Prob(T=j\mid S=s)
 =\binom mj\frac{B(a+s+j,b+n-s+m-j)}{B(a+s,b+n-s)}.
\]
W szczególności przyszły pojedynczy sukces ma prawdopodobieństwo
$(a+s)/(a+b+n)$.
\end{twierdzenie}
\begin{proof}
Prawdopodobieństwo $S=s$ przy znanym $p$ wynosi
$\binom ns p^s(1-p)^{n-s}$. Iloczyn gęstości prioru i
wiarygodności jest proporcjonalny do
$p^{a+s-1}(1-p)^{b+n-s-1}$, więc normalizacja przez
$B(a+s,b+n-s)$ daje rozkład wskazany w tezie.
Dla $r>0$, $t>0$, $k\in\{1,2\}$, z definicji beta
$\E[U^k]=B(r+k,t)/B(r,t)$, gdzie
$U\sim\operatorname{Beta}(r,t)$. Tożsamość
$\Gamma(z+1)=z\Gamma(z)$ daje
$\E U=r/(r+t)$ oraz $\Var(U)=rt/[(r+t)^2(r+t+1)]$.
Po wstawieniu $r=a+s$, $t=b+n-s$ otrzymujemy wzory
momentowe. Dla nowych obserwacji
$\Prob(T=j\mid p)=\binom mjp^j(1-p)^{m-j}$.
Całkując to prawdopodobieństwo względem wyprowadzonego
posterioru, otrzymujemy ostatni wzór.
\end{proof}

\begin{stwierdzenie}[Średnia posterioru jako kompromis MLE i prioru]
\label{prop:beta-shrinkage}
W założeniach twierdzenia~\ref{thm:beta-binomial} oznaczmy
$\widehat p=S/n$, $m_0=a/(a+b)$ i $w=n/(n+a+b)$.
Wówczas $\E[p\mid S]=w\widehat p+(1-w)m_0$.
\end{stwierdzenie}
\begin{proof}
Po podstawieniu dostajemy
$\frac{n}{n+a+b}\frac Sn+
\frac{a+b}{n+a+b}\frac{a}{a+b}=(S+a)/(n+a+b)$.
\end{proof}

\begin{przyklad}[Kontrola jakości: 7 wad na 100 sztuk]
Dla $S=7$, $n=100$ estymator MLE wynosi $0{,}07$.
Z priorem $\operatorname{Beta}(1,1)$ posterior ma rozkład
$\operatorname{Beta}(8,94)$, średnią $8/102\approx0{,}0784$,
a prawdopodobieństwo wady następnej sztuki wynosi także $8/102$.
Dla prioru $\operatorname{Beta}(20,80)$ posterior ma rozkład
$\operatorname{Beta}(27,173)$ i średnią $0{,}135$.
Obydwa modele używają identycznych danych, lecz różnią się
założeniami; prior o łącznej masie parametrów $a+b=100$
ma wyraźny wpływ na wynik. „Siła prioru” $a+b$ jest tu
wygodnym skrótem algebraicznym, nie dokładną liczbą
rzeczywistych dodatkowych niezależnych obserwacji.
Przedział wiarygodności wyznaczymy z kwantyli rozkładu beta,
a jego interpretację odróżnimy od przedziału Wilsona z
rozdziału~\ref{ch:confidence-intervals}.
\end{przyklad}

\begin{uwaga}[MAP i przypadki na brzegu]
Dla posterioru $\operatorname{Beta}(r,t)$, $r,t>1$,
moda leży w $(r-1)/(r+t-2)$. Jeśli $r\le1$ lub $t\le1$,
maksimum gęstości może wypadać na brzegu lub nie być osiągane
na otwartym $(0,1)$. Gdy $r=t=1$, gęstość jest stała:
każdy punkt $(0,1)$ jest jej dominantą. Przy $r<1$ lub
$t<1$ gęstość może być nieograniczona w pobliżu brzegu,
a wówczas \emph{maksymalizator w otwartym przedziale nie istnieje}.
Nie wolno bezwarunkowo używać wzoru na modę dla dowolnych
$r,t>0$. Sam MAP nie jest niezmienniczy względem nieliniowej
reparametryzacji; podobny problem dotyczy definicji HPD.
\end{uwaga}

\section{Gamma--Poisson: częstość awarii i ekspozycja}
\begin{definicja}[Rozkład gamma w parametryzacji shape--rate]
Dla $a,b>0$ przez $\operatorname{Gamma}(a,b)$ rozumiemy
rozkład na $(0,\infty)$ o gęstości
$g_{a,b}(\lambda)=b^a\lambda^{a-1}e^{-b\lambda}/\Gamma(a)$.
Pierwszy parametr to \emph{kształt} (shape), drugi to
\emph{intensywność zaniku} (rate), a nie skala.
Średnia wynosi $a/b$, a wariancja $a/b^2$.
Jeżeli $\lambda$ jest liczbą awarii na godzinę, wtedy
$b$ ma jednostkę godzin, tak aby iloczyn $b\lambda$
był bezwymiarowy.
\end{definicja}

\begin{twierdzenie}[Aktualizacja i predykcja gamma--Poisson]
\label{thm:gamma-poisson}
Niech $\lambda\sim\operatorname{Gamma}(a,b)$,
$a,b>0$. Warunkowo przy $\lambda$ obserwujemy niezależne
zliczenia $Y_i\sim\operatorname{Pois}(e_i\lambda)$,
gdzie ekspozycje $e_i>0$ są znanymi liczbami,
$i=1,\ldots,n$, $n\ge1$. Niech $S=\sum_{i=1}^nY_i$
i $E=\sum_{i=1}^ne_i$. Wtedy
$\lambda\mid Y_1,\ldots,Y_n\sim\operatorname{Gamma}(a+S,b+E)$.
Dla nowej ekspozycji $e_*>0$ i nowego zliczenia
$Y_*\mid\lambda\sim\operatorname{Pois}(e_*\lambda)$,
warunkowo niezależnego od danych, zachodzi dla $k\in\N\cup\{0\}$
\[
 \Prob(Y_*=k\mid Y_1,\ldots,Y_n)
 =\frac{\Gamma(a+S+k)}{\Gamma(a+S)k!}
 \left(\frac{b+E}{b+E+e_*}\right)^{a+S}
 \left(\frac{e_*}{b+E+e_*}\right)^k.
\]
Ponadto $\E[Y_*\mid Y]=e_*(a+S)/(b+E)$ i
$\Var(Y_*\mid Y)=e_*(a+S)/(b+E)
+e_*^2(a+S)/(b+E)^2$.
\end{twierdzenie}
\begin{proof}
Warunkowa wiarygodność zliczeń, z dokładnością do
czynnika niezależnego od $\lambda$, wynosi
$\lambda^Se^{-E\lambda}$. Mnożenie przez prior
$\lambda^{a-1}e^{-b\lambda}$ daje jądro gamma
$\lambda^{a+S-1}e^{-(b+E)\lambda}$, które jest
normalizowalne, gdy $a+S,b+E>0$.
Piszmy $r=a+S$, $v=b+E$. Warunkowe prawdopodobieństwo
$Y_*=k$ ma postać
\[
 \frac{e_*^kv^r}{k!\Gamma(r)}
 \int_0^\infty\lambda^{k+r-1}
 e^{-(v+e_*)\lambda}\,\dd\lambda_1(\lambda).
\]
Z definicji funkcji gamma całka wynosi
$\Gamma(r+k)/(v+e_*)^{r+k}$, co daje wzór.
Wreszcie korzystamy z tożsamości
$\E[Y_*\mid\lambda,Y]=e_*\lambda$ i
$\Var(Y_*\mid\lambda,Y)=e_*\lambda$;
wzór na wariancję całkowitą z rozdziału~\ref{ch:bayes-foundations}
daje podane dwa składniki.
\end{proof}

\begin{przyklad}[Awaryjność z różną długością obserwacji]
Obserwujemy 4 awarie podczas 200 godzin działania urządzenia.
Model zliczeń ma ekspozycję $E=200$ godzin i nieznaną
intensywność $\lambda$ na godzinę. Niech prior to
$\operatorname{Gamma}(2,100)$ (shape--rate).
Posterior to $\operatorname{Gamma}(6,300)$ i
$\E[\lambda\mid\mathrm{dane}]=6/300=0{,}02$ na godzinę.
Prognozowana liczba awarii przez następne 50 godzin ma średnią
$50\cdot0{,}02=1$. Jej wariancja wynosi
$1+50^2\cdot 6/300^2=1+1/6$; to więcej niż
wariancja rozkładu Poissona z ustalonym parametrem $1$.
Powodem jest niepewność o $\lambda$.
W praktyce niezależność przyrostów i stała intensywność
w czasie wymagają sprawdzenia; przy sezonowości lub awariach
w klastrach rozważamy model bardziej złożony.
\end{przyklad}

\section{Gamma--wykładniczy: czas oczekiwania}
\begin{twierdzenie}[Posterior intensywności rozkładu wykładniczego]
Niech $\lambda\sim\operatorname{Gamma}(a,b)$, $a,b>0$,
a przy znanym $\lambda$ obserwacje $T_1,\ldots,T_n$
będą niezależne o gęstości
$f(t\mid\lambda)=\lambda e^{-\lambda t}\ind{t\ge0}$
względem miary Lebesgue'a, $n\ge1$.
Dla $s=\sum_{i=1}^n t_i$ posterior ma rozkład
$\operatorname{Gamma}(a+n,b+s)$.
Jeżeli przyszły czas $T_*$ jest warunkowo niezależny od danych,
ma warunkową gęstość wykładniczą o tej samej intensywności,
to dla $u\ge0$
\[
 \Prob(T_*>u\mid T_1=t_1,\ldots,T_n=t_n)
 =\left(\frac{b+s}{b+s+u}\right)^{a+n}.
\]
\end{twierdzenie}
\begin{proof}
Iloczyn gęstości obserwacji daje $\lambda^n e^{-s\lambda}$,
więc posterior jest gamma przez mnożenie przez jego prior.
Po warunkowaniu na $\lambda$ prawdopodobieństwo przetrwania
to $e^{-u\lambda}$. Całkując je względem posterioru
$\operatorname{Gamma}(a+n,b+s)$ i wykorzystując wzór na
funkcję gamma, otrzymujemy wskazaną potęgę.
\end{proof}

\begin{uwaga}[Cenzorowanie a kompletne czasy]
Jeżeli niektórych czasów nie zaobserwowano do końca, zwykły
iloczyn $\prod_i\lambda e^{-\lambda t_i}$ jest błędny.
Przy prawostronnym, niezależnym cenzorowaniu z liczbą
rzeczywiście zaobserwowanych zdarzeń $d$ i łącznym czasem
ekspozycji $E$, część wiarygodności zależna od $\lambda$
ma postać $\lambda^d e^{-\lambda E}$. Posterior
$\operatorname{Gamma}(a+d,b+E)$ pozostaje sprzężony,
ale tylko przy prawidłowym opisaniu mechanizmu cenzorowania.
\end{uwaga}

\section{Normalny model średniej ze znaną wariancją}
\begin{twierdzenie}[Normalny prior dla normalnej średniej]
\label{thm:normal-normal}
Niech $\mu\sim\mathcal N(m_0,v_0)$, $v_0>0$, a warunkowo
względem $\mu$ zmienne $X_1,\ldots,X_n$ będą niezależne
$\mathcal N(\mu,\sigma^2)$ dla znanego $\sigma^2>0$, $n\ge1$.
Wtedy
\[
 \mu\mid X_1=x_1,\ldots,X_n=x_n\sim\mathcal N(m_n,v_n),
 \quad v_n=\left(\frac1{v_0}+\frac n{\sigma^2}\right)^{-1},
 \quad m_n=v_n\left(\frac{m_0}{v_0}+\frac{n\overline x}{\sigma^2}\right).
\]
Średnia posterioru jest średnią ważoną priorowej średniej
$m_0$ i średniej z próby $\overline x$, z wagami
proporcjonalnymi do ich precyzji (odwrotności wariancji).
\end{twierdzenie}
\begin{proof}
Zależność od $\mu$ w priorze i wiarygodności ma postać
$\exp\{-\tfrac12[(\mu-m_0)^2/v_0+
\sum_{i=1}^n(x_i-\mu)^2/\sigma^2]\}$.
Po rozpisaniu kwadratów współczynnik przy $\mu^2$
wynosi $v_0^{-1}+n\sigma^{-2}=v_n^{-1}$,
a współczynnik przy $-2\mu$ jest
$m_0/v_0+n\overline x/\sigma^2=v_n^{-1}m_n$.
Uzupełnienie kwadratu daje jądro rozkładu
$\mathcal N(m_n,v_n)$; normalizacja jest skończona.
\end{proof}

\begin{przyklad}[Pomiar z kalibracją a nowe obserwacje]
Z wcześniejszej kalibracji wynika $\mu\sim\mathcal N(10,1)$.
Pięć nowych niezależnych pomiarów ma znaną wariancję
błędu $\sigma^2=4$ i średnią $12$. Posteriorowa precyzja
wynosi $1+5/4=2{,}25$, zatem
$v_n=4/9\approx0{,}4444$ i
$m_n=(10+5\cdot12/4)/(1+5/4)=100/9\approx11{,}111$.
Wartość średniej została przesunięta od priorowego $10$ ku
obserwowanemu $12$. Jeżeli wariancja pomiaru jest nieznana,
ten rozkład przestaje być automatycznie właściwym modelem.
\end{przyklad}

\section{Nieznane średnia i wariancja: normalny--odwrotny gamma}
\begin{definicja}[Rozkład odwrotny gamma]
Dla $\alpha,\beta>0$ rozkład
$\operatorname{InvGamma}(\alpha,\beta)$ na $(0,\infty)$
ma gęstość
$h(v)=\beta^\alpha v^{-\alpha-1}e^{-\beta/v}/\Gamma(\alpha)$.
Parametr $\beta$ jest parametrem \emph{skali w wykładniku},
a nie intensywnością gamma na $v$. Jeżeli
$V\sim\operatorname{InvGamma}(\alpha,\beta)$,
to $V^{-1}\sim\operatorname{Gamma}(\alpha,\beta)$
w parametryzacji shape--rate.
\end{definicja}

\begin{twierdzenie}[Sprzężony model normalny z nieznanymi parametrami]
\label{thm:normal-invgamma}
Niech $X_1,\ldots,X_n$ będą warunkowo niezależne
$\mathcal N(\mu,v)$, gdzie $\mu\in\R$, $v>0$, $n\ge1$.
Niech prior ma postać
$v\sim\operatorname{InvGamma}(\alpha_0,\beta_0)$,
$\mu\mid v\sim\mathcal N(m_0,v/\kappa_0)$,
przy $\alpha_0,\beta_0,\kappa_0>0$, $m_0\in\R$.
Piszmy $\overline x=n^{-1}\sum_ix_i$,
$S=\sum_i(x_i-\overline x)^2$. Wtedy posterior ma tę samą
postać z parametrami
\[
 \kappa_n=\kappa_0+n,\quad
 m_n=\frac{\kappa_0m_0+n\overline x}{\kappa_0+n},\quad
 \alpha_n=\alpha_0+\frac n2,
\]
\[
 \beta_n=\beta_0+\frac S2+
 \frac{\kappa_0n}{2(\kappa_0+n)}(\overline x-m_0)^2.
\]
Dokładniej $v\mid x\sim\operatorname{InvGamma}(\alpha_n,\beta_n)$
i $\mu\mid v,x\sim\mathcal N(m_n,v/\kappa_n)$.
\end{twierdzenie}
\begin{proof}
Wspólny prior ma gęstość względem miary Lebesgue'a $\lambda_2$
na $\R\times(0,\infty)$
proporcjonalną do
$v^{-\alpha_0-3/2}\exp\{-\beta_0/v-
\kappa_0(\mu-m_0)^2/(2v)\}$.
Wiarygodność jest proporcjonalna do
$v^{-n/2}\exp\{-\sum_i(x_i-\mu)^2/(2v)\}$.
Korzystamy z dwóch tożsamości:
$\sum_i(x_i-\mu)^2=S+n(\mu-\overline x)^2$ oraz
\[
 \kappa_0(\mu-m_0)^2+n(\mu-\overline x)^2
 =\kappa_n(\mu-m_n)^2+
 \frac{\kappa_0n}{\kappa_n}(\overline x-m_0)^2.
\]
Po podstawieniu wspólne jądro posterioru jest równe
$v^{-\alpha_n-3/2}\exp\{-\beta_n/v-
\kappa_n(\mu-m_n)^2/(2v)\}$.
Jest to jądro gęstości normalnej warunkowej pomnożonej przez
gęstość odwrotną gamma z parametrami $\alpha_n,\beta_n$;
obydwa czynniki są normalizowalne.
\end{proof}

\begin{wniosek}[Predykcja t-Studenta]
\label{cor:student-predictive}
W założeniach twierdzenia~\ref{thm:normal-invgamma}
niech $X_*$ będzie nową obserwacją warunkowo niezależną
od próby przy znanych $\mu,v$ i niech $X_*\mid\mu,v\sim
\mathcal N(\mu,v)$. Wtedy $X_*\mid x$ ma rozkład
$t$-Studenta z $2\alpha_n$ stopniami swobody, środkiem $m_n$
i parametrem skali
$s_n=\sqrt{\beta_n(\kappa_n+1)/(\alpha_n\kappa_n)}$,
w tej konwencji, że $(X_*-m_n)/s_n$ ma standardowy rozkład
$t_{2\alpha_n}$.
\end{wniosek}
\begin{proof}
Przy ustalonym $v$ całkujemy najpierw po $\mu$;
w wyniku splotu normalnego dostajemy
$X_*\mid v,x\sim\mathcal N(m_n,v(1+\kappa_n^{-1}))$.
Zatem przy niezależnych $Z\sim\mathcal N(0,1)$
i $G\sim\operatorname{Gamma}(\alpha_n,\beta_n)$
można przedstawić $X_*-m_n=\sqrt{(1+\kappa_n^{-1})/G}\,Z$.
Ponieważ $2\beta_n G\sim\chi^2_{2\alpha_n}$
(z własności transformacji rozkładu gamma), stosunek
$Z/\sqrt{2\beta_nG/(2\alpha_n)}$ ma standardowy rozkład
$t_{2\alpha_n}$. Po przeliczeniu współczynnika otrzymujemy
$s_n^2=\beta_n(\kappa_n+1)/(\alpha_n\kappa_n)$.
\end{proof}

\begin{uwaga}[Uwaga na istnienie momentów]
Dla posterioru $\operatorname{InvGamma}(\alpha_n,\beta_n)$
wartość oczekiwana wariancji $v$ istnieje dopiero przy
$\alpha_n>1$ i wynosi $\beta_n/(\alpha_n-1)$.
Wariancja rozkładu predykcyjnego t-Studenta jest skończona
dopiero przy $2\alpha_n>2$. Nie należy podawać nieistniejących
momentów w modelach z bardzo słabą informacją.
\end{uwaga}

\section{Dirichlet--wielomianowy: więcej niż dwie kategorie}
\begin{definicja}[Rozkład Dirichleta]
Dla $K\ge2$ i $a_1,\ldots,a_K>0$ niech
$\Delta_{K-1}=\{(p_1,\ldots,p_K)\in(0,1)^K:\sum_{j=1}^Kp_j=1\}$.
Rozkład $\operatorname{Dirichlet}(a_1,\ldots,a_K)$ ma
gęstość względem $\dd\lambda_{K-1}(p_1,\ldots,p_{K-1})$
na obszarze $p_j>0$, $\sum_{j<K}p_j<1$, po podstawieniu
$p_K=1-\sum_{j<K}p_j$, proporcjonalną do
$\prod_{j=1}^Kp_j^{a_j-1}$. Stała normalizacji jest
$\Gamma(a_\bullet)/\prod_{j=1}^K\Gamma(a_j)$,
gdzie $a_\bullet=\sum_{j=1}^Ka_j$.
\end{definicja}

\begin{lemat}[Stała normalizująca rozkładu Dirichleta]
Dla $K\ge2$ oraz $a_1,\ldots,a_K>0$ zachodzi
\[
 \int_{D_{K-1}}\prod_{j=1}^{K}p_j^{a_j-1}
 \,\dd\lambda_{K-1}(p_1,\ldots,p_{K-1})
 =\frac{\prod_{j=1}^K\Gamma(a_j)}{\Gamma(a_\bullet)},
\]
gdzie $D_{K-1}=\{(p_1,\ldots,p_{K-1}):p_j>0,
\ \sum_{j<K}p_j<1\}$, $p_K=1-\sum_{j<K}p_j$ oraz
$a_\bullet=\sum_{j=1}^Ka_j$.
\end{lemat}
\begin{proof}
Iloczyn $\prod_{j=1}^K\Gamma(a_j)$ jest całką funkcji
$\prod_{j=1}^Kt_j^{a_j-1}e^{-t_j}$ po $(0,\infty)^K$
względem $\lambda_K$. Wprowadzamy
$r=\sum_{j=1}^Kt_j$ i $p_j=t_j/r$ dla $j<K$;
wtedy $t_K=r(1-\sum_{j<K}p_j)$.
Jakobian tej zamiany zmiennych ma wartość $r^{K-1}$.
Po zastosowaniu Tonellego całka rozdziela się na
$\int_0^\infty r^{a_\bullet-1}e^{-r}\,\dd\lambda_1(r)
=\Gamma(a_\bullet)$ i całkę z tezy. Dzielenie
przez dodatnią wartość $\Gamma(a_\bullet)$ kończy dowód.
\end{proof}

\begin{twierdzenie}[Sprzężoność dla zliczeń wielokategorialnych]
Niech $p=(p_1,\ldots,p_K)$ ma rozkład
$\operatorname{Dirichlet}(a_1,\ldots,a_K)$, a przy znanym $p$
obserwujemy $n\ge1$ niezależnych wyników z $K$ kategorii,
przy czym prawdopodobieństwo kategorii $j$ wynosi $p_j$.
Jeśli w danych kategoria $j$ wystąpiła $N_j$ razy,
$\sum_{j=1}^KN_j=n$, to posterior ma rozkład
$\operatorname{Dirichlet}(a_1+N_1,\ldots,a_K+N_K)$.
Prawdopodobieństwo, że następny wynik trafi do kategorii $j$,
wynosi $(a_j+N_j)/(a_\bullet+n)$.
\end{twierdzenie}
\begin{proof}
Wiarygodność jako funkcja $p$ jest proporcjonalna do
$\prod_{j=1}^Kp_j^{N_j}$. Po pomnożeniu przez gęstość
Dirichleta otrzymujemy jądro
$\prod_{j=1}^Kp_j^{a_j+N_j-1}$, którego normalizacja
jest gęstością posterioru. Średnią współrzędnej
$p_j$ obliczamy przez iloraz stałych normalizujących
$\Gamma(a_j+N_j+1)/\Gamma(a_j+N_j)$ oraz
$\Gamma(a_\bullet+n)/\Gamma(a_\bullet+n+1)$.
\end{proof}

\section{Obliczenia numeryczne i porównanie z metodami klasycznymi}
Posterior warunkowy opisuje parametry \emph{po zobaczeniu danych},
ale jeszcze przed ich obserwacją także prior nakłada konkretne
ograniczenia na możliwą próbę. W rozdziale~\ref{ch:bayes-estimation-prediction}
wyprowadzimy predykcyjny rozkład liczby sukcesów z modelu
beta--dwumianowego i pokażemy, jak może on służyć do oceny
realizmu założeń a priori. Nie wolno zastępować tej kontroli
stwierdzeniem, że algebraiczna sprzężoność jest wygodna.

W obliczeniach posteriorów korzystamy ze stabilnych funkcji
bibliotecznych kwantyli beta, gamma i t-Studenta; bezpośrednie
mnożenie dużej liczby gęstości prowadzi do niedomiaru
zmiennoprzecinkowego. Poniższe programy odtwarzają
parametry posterioru i predykcję, ale nie zastępują
diagnostyki poprawności modelu.

\begin{lstlisting}[language=Python,caption={Kompletny program: beta--dwumianowy, przedziały i test wrażliwości na prior (Python).}]
from scipy.stats import beta

n, defects = 100, 7
for a, b in [(1, 1), (2, 18), (20, 80)]:
    post_a, post_b = a + defects, b + n - defects
    mean = post_a / (post_a + post_b)
    lo, hi = beta.ppf([0.025, 0.975], post_a, post_b)
    prob_over_10pct = beta.sf(0.10, post_a, post_b)
    print((a, b), "mean:", round(mean, 5),
          "95% credible:", (round(lo, 5), round(hi, 5)),
          "P(p>0.10|data):", round(prob_over_10pct, 5))

# Probability of zero defects in the next ten devices.
a1, b1 = 1 + defects, 1 + n - defects
from scipy.special import betaln
import numpy as np
prob_zero = np.exp(betaln(a1, b1 + 10) - betaln(a1, b1))
print("predictive P(0 defects among next 10):", prob_zero)
\end{lstlisting}

\begin{lstlisting}[language=Julia,caption={Kompletny program: aktualizacja sprzężona i predykcja gamma--Poisson (Julia, biblioteka standardowa).}]
using Printf

function gamma_poisson_update(a::Real, b::Real,
                              counts::AbstractVector{<:Integer},
                              exposures::AbstractVector{<:Real})
    length(counts) == length(exposures) ||
        throw(DimensionMismatch("different lengths"))
    a > 0 && b > 0 || throw(ArgumentError("prior must be proper"))
    all(x -> x >= 0, counts) || throw(ArgumentError("negative count"))
    all(x -> x > 0, exposures) || throw(ArgumentError("exposure <= 0"))
    a_new = a + sum(counts)
    b_new = b + sum(exposures)
    return a_new, b_new
end

a, b = gamma_poisson_update(2.0, 100.0, [1, 0, 3],
                            [50.0, 50.0, 100.0])
e_next = 50.0
mean_next = e_next * a / b
var_next = mean_next + e_next^2 * a / b^2
@printf("posterior Gamma(shape=%.1f, rate=%.1f)\n", a, b)
@printf("predictive mean=%.6f, variance=%.6f\n",
        mean_next, var_next)
\end{lstlisting}

\begin{center}\small
\begin{tabular}{>{\raggedright\arraybackslash}p{.26\textwidth}>{\raggedright\arraybackslash}p{.29\textwidth}>{\raggedright\arraybackslash}p{.32\textwidth}}
\toprule
\textbf{Problem / dane} & \textbf{Klasyczne narzędzie} & \textbf{Sprzężony model bayesowski i uwaga}\\
\midrule
Odsetek wad, mała próba & MLE, Wilson, Clopper--Pearson & Beta--Bernoulli; posterior zależy od prioru.\\
Zliczenia z ekspozycją & MLE Poissona, dokładny przedział & Gamma--Poisson; łatwa predykcja z nadmierną dyspersją.\\
Średnia przy znanym $\sigma$ & Średnia, przedział normalny & Normalny--normalny; precyzyjna aktualizacja.\\
Średnia i wariancja przy normalności & $t$-Studenta, $\chi^2$ & Normalny--odwrotny gamma; predykcja t.\\
Wiele kategorii & Udziały empiryczne, wielomianowy & Dirichlet--wielomianowy; prosty posterior.\\
Nieustalony mechanizm zależności & Model czasowy lub hierarchiczny & Sama sprzężoność iid nie usuwa zależności.\\
\bottomrule
\end{tabular}
\end{center}

\begin{zadanie}[Brak awarii nie oznacza zerowej intensywności]
Podczas 50 godzin działania nie zanotowano awarii. Przy priorze
$\lambda\sim\operatorname{Gamma}(1,50)$ i modelu Poissona
z intensywnością $\lambda$ na godzinę wyznacz posterior oraz
prawdopodobieństwo braku awarii w ciągu kolejnych 50 godzin.
\par\emph{Rozwiązanie.} W twierdzeniu~\ref{thm:gamma-poisson}
$S=0$, $E=50$, więc posterior to $\operatorname{Gamma}(1,100)$.
Prawdopodobieństwo $Y_*=0$ dla $e_*=50$ wynosi
$(100/(100+50))^1=2/3$.
Estymator MLE intensywności wynosi $0$, ale nie wynika z tego
zerowa predykcyjna szansa przyszłej awarii.
\end{zadanie}

\begin{zadanie}[Dane i wrażliwość na założenia]
Dla danych $S=7$, $n=100$ rozważ priory
$\operatorname{Beta}(1,1)$ i $\operatorname{Beta}(20,80)$.
Porównaj średnie posteriorów i zinterpretuj różnicę.
\par\emph{Rozwiązanie.} Otrzymujemy odpowiednio
$8/102\approx0{,}0784$ i $27/200=0{,}135$.
Duża różnica przy stałych danych wynika z innej ilości i treści
informacji zakodowanej w priorach. W praktyce należy opisać
pochodzenie prioru i pokazać analizę wrażliwości; nie wystarczy
wybrać prioru, który daje oczekiwany wynik.
\end{zadanie}


% ================================================================
% ROZDZIAL 19 -- ESTYMACJA I PREDYKCJA BAYESOWSKA
% ================================================================
\chapter{Estymacja i predykcja bayesowska: od posterioru do decyzji}
\label{ch:bayes-estimation-prediction}
\paragraph{Powiązania i odsyłacze.} Konstrukcje posteriorów pochodzą z rozdziałów~\ref{ch:bayes-foundations} i~\ref{ch:conjugate-families}; funkcje straty z rozdziału~\ref{ch:decision}.


\section{Co właściwie chcemy oszacować?}
W rozdziałach~\ref{ch:bayes-foundations}--\ref{ch:conjugate-families}
wyprowadziliśmy rozkłady a posteriori. Sam posterior nie rozstrzyga
jednak, \emph{jaką liczbę powinniśmy podać} ani \emph{jaką decyzję
podjąć}. Te dwa pytania mają różne odpowiedzi: pierwsze wymaga
wyboru funkcji straty, drugie dodatkowo określenia rzeczywistych
konsekwencji działań. Jeszcze inne pytanie dotyczy prognozy
nowej obserwacji. Dla przykładu w kontroli jakości rozróżniamy
trzy cele: oszacowanie nieznanego odsetka wad $p$, przewidzenie
liczby wad w następnej partii i podjęcie decyzji o zatrzymaniu produkcji.

\begin{definicja}[Reguła bayesowska i ryzyko posteriorowe]
\label{def:bayes-rule-ch19}
Niech $(\Theta_0,\mathcal T)$, $(\mathsf X,\mathcal X)$ i
$(\mathsf A,\mathcal G)$ będą przestrzeniami mierzalnymi,
$\pi$ właściwym priorem, $K(\theta,\cdot)$ jądrem modelu danych,
a $\Pi(x,\cdot)$ jego posteriorem. Niech
$L:\Theta_0\times\mathsf A\to[0,\infty]$ będzie mierzalną
funkcją straty. Dla obserwacji $x$ i działania $a$ definiujemy
\emph{ryzyko posteriorowe}
$\rho_x(a)=\int_{\Theta_0}L(\theta,a)\,\dd\Pi(x,\theta)$.
Mierzalną regułę $\delta:\mathsf X\to\mathsf A$ nazywamy
\emph{regułą Bayesa}, jeśli $\rho_x(\delta(x))=\inf_{a\in\mathsf A}\rho_x(a)$
dla $M$-prawie każdego $x$, gdzie $M$ jest rozkładem
predykcyjnym danych przed ich obserwacją. Istnienie mierzalnego
minimizatora wymaga osobnego uzasadnienia; sama definicja go
nie gwarantuje.
\end{definicja}

\begin{twierdzenie}[Optymalizacja posteriorowa minimalizuje ryzyko Bayesa]
\label{thm:bayes-risk-tower}
W założeniach definicji~\ref{def:bayes-rule-ch19} niech
$\delta_*$ będzie mierzalną regułą Bayesa. Wtedy dla każdej
mierzalnej reguły $\delta:\mathsf X\to\mathsf A$ zachodzi
\[
\int_{\Theta_0}\int_{\mathsf X}
 L(\theta,\delta_*(x))\,\dd K(\theta,x)\,\dd\pi(\theta)
\le
\int_{\Theta_0}\int_{\mathsf X}
 L(\theta,\delta(x))\,\dd K(\theta,x)\,\dd\pi(\theta).
\]
Dla obu stron dopuszczamy wartość $+\infty$.
\end{twierdzenie}
\begin{proof}
Nieujemność i mierzalność straty pozwalają zastosować Tonellego.
Z własności warunkowania, najpierw dla funkcji prostych, a potem
przez monotoniczną zbieżność,
\[
\int_{\Theta_0}\int_{\mathsf X}L(\theta,\delta(x))
\,\dd K(\theta,x)\,\dd\pi(\theta)
=\int_{\mathsf X}\rho_x(\delta(x))\,\dd M(x).
\]
Przy każdym $x$ poza zbiorem $M$-miary zero reguła $\delta_*$
minimalizuje wyrażenie pod całką. Całkowanie tej nierówności
daje tezę. To minimalizacja ryzyka \emph{uśrednionego względem
prioru}; nie wynika z niej dominacja ryzyka częstotliwościowego
przy każdym ustalonym $\theta$. Wynik dotyczy wyłącznie
porównywanych reguł, ustalonego modelu i przyjętej funkcji straty.
\end{proof}

\section{Funkcja straty wyznacza estymator}
W twierdzeniu~\ref{thm:bayes-actions} wykazaliśmy, że strata
kwadratowa prowadzi do średniej posterioru, a bezwzględna do
mediany. Do decyzji o rezerwie lub zapasie części zamiennych
bardziej pasuje niekiedy \emph{strata asymetryczna}: zaniżenie
przewidywanej liczby awarii kosztuje więcej niż jej zawyżenie.

\begin{definicja}[Asymetryczna strata liniowa]
Dla $\tau\in(0,1)$ i rzeczywistych $\theta,a$ definiujemy
$L_\tau(\theta,a)=\tau(\theta-a)_++(1-\tau)(a-\theta)_+$,
gdzie $t_+=\max(t,0)$. Parametr $\tau$ określa względny
koszt niedoszacowania, a $1-\tau$ przeszacowania.
\end{definicja}

\begin{twierdzenie}[Kwantyl posterioru jako optymalna decyzja]
\label{thm:bayes-quantile-loss}
Niech $\Pi_x$ będzie rozkładem na $\R$ o skończonym
pierwszym momencie, a $F_x(t)=\Pi_x(({-}\infty,t])$ jego
dystrybuantą. Dla $\tau\in(0,1)$ liczba $a$ minimalizuje
$\int_\R L_\tau(\theta,a)\,\dd\Pi_x(\theta)$ po $a\in\R$
wtedy i tylko wtedy, gdy
$F_x(a-)\le\tau\le F_x(a)$.
W szczególności przy ciągłej, ściśle rosnącej dystrybuancie
jedyną optymalną decyzją jest kwantyl $F_x^{-1}(\tau)$.
\end{twierdzenie}
\begin{proof}
Strata jest wypukła względem $a$ i ograniczona przez
$|\theta|+|a|$, więc jej wartość oczekiwana jest skończona.
Dla $a$ niebędącego atomem rozkładu pochodna po $a$ jest
równa $F_x(a)-\tau$: przy $\theta<a$ pochodna wynosi
$1-\tau$, a przy $\theta>a$ jest równa $-\tau$.
Ogólnie jednostronne pochodne funkcji ryzyka to
$F_x(a-)-\tau$ oraz $F_x(a)-\tau$, co wynika
z ilorazów różnicowych i zbieżności zdominowanej.
Wypukła funkcja osiąga minimum w punkcie $a$ dokładnie
wtedy, gdy lewa pochodna jest nie większa od zera,
a prawa nie mniejsza od zera. To daje tezę.
\end{proof}

\begin{przyklad}[Wielkość zapasu części na kolejny tydzień]
Niech predykcyjna liczba awarii $Y$ ma rozkład, dla którego
kwantyl rzędu $0{,}9$ wynosi $q_{0{,}9}$. Jeśli koszt
brakującej części jest dziewięciokrotnie większy niż koszt
zbędnej części, strata liniowa ma $\tau=0{,}9$.
Optymalna liczba części jest \emph{całkowitoliczbowym}
kwantylem predykcyjnym rzędu $0{,}9$, a nie średnią
posterioru intensywności. W praktyce należy również sprawdzić
koszty magazynowania, ograniczenia pojemności i możliwość
dokupienia części w trakcie tygodnia.
\end{przyklad}

\section{Rozkład parametru a rozkład przyszłych obserwacji}
Niech $\lambda$ oznacza intensywność awarii na godzinę,
a nowe zliczenie $Y_*$ dotyczy $e_*>0$ godzin. Gdy znamy
$\lambda$, losowość $Y_*$ jest poissonowska. Gdy jej nie
znamy, predykcja musi scałkować tę losowość po posteriorze.

\begin{twierdzenie}[Predykcja gamma--Poisson ma większą wariancję]
\label{thm:gamma-poisson-dispersion}
Jeżeli $\lambda\mid X=x\sim\operatorname{Gamma}(r,v)$,
$r,v>0$, a $Y_*\mid\lambda,X=x\sim\operatorname{Pois}(e\lambda)$
dla $e>0$, to
\[
\E[Y_*\mid X=x]=\frac{er}{v},\qquad
\Var(Y_*\mid X=x)=\frac{er}{v}+\frac{e^2r}{v^2}.
\]
W szczególności $\Var(Y_*\mid X=x)>\E[Y_*\mid X=x]$.
\end{twierdzenie}
\begin{proof}
Korzystamy kolejno z
$\E[Y_*\mid\lambda,x]=e\lambda$,
$\Var(Y_*\mid\lambda,x)=e\lambda$ i z prawa całkowitej
wariancji względem posterioru. Pierwszy składnik daje
$e\E[\lambda\mid x]=er/v$, drugi
$e^2\Var(\lambda\mid x)=e^2r/v^2$.
Dodatniość ostatniego składnika wynika z $e,r,v>0$.
\end{proof}

\begin{przyklad}[Ile awarii przewidzimy?]
Z rozdziału~\ref{ch:conjugate-families} weźmy posterior
$\operatorname{Gamma}(6,300)$ (intensywność na godzinę).
Przy ekspozycji $e=50$ predykcja ma średnią $1$ i wariancję
$7/6$. Dokładny rozkład liczby awarii jest ujemny
dwumianowy; z twierdzenia~\ref{thm:gamma-poisson}
$\Prob(Y_*=0\mid x)=(6/7)^6\approx0{,}39657$,
a $\Prob(Y_*\ge3\mid x)\approx0{,}09356$.
Jeśli konieczne jest zabezpieczenie przed co najmniej trzema
awariami, stosujemy ostatnie prawdopodobieństwo, a nie tylko
średnią prognozy równą jeden.
\end{przyklad}

\section{Predykcja optymalna względem właściwych funkcji oceny}
Przy ocenianiu prognoz probabilistycznych chcemy nagradzać
podawanie całego rozkładu, a nie tylko liczby. Oczekiwane
ryzyko obliczamy zawsze względem jawnie wskazanego prawa
przyszłej obserwacji; nie jest to automatycznie to samo, co
częstościowa ocena modelu na niezależnej próbie testowej. \emph{Właściwa
funkcja oceny} (ang. \emph{proper scoring rule}) jest taką funkcją
straty, dla której prawdziwy rozkład predykcyjny minimalizuje
oczekiwaną stratę. Rozróżnijmy dwa elementarne przypadki.

\begin{twierdzenie}[Poprawność miary Briera dla zdarzenia binarnego]
\label{thm:brier}
Niech $Y\in\{0,1\}$ przy zadanej informacji ma rozkład
$\operatorname{Bernoulli}(p)$, $p\in[0,1]$. Prognozujemy
$p$ liczbą $q\in[0,1]$ i stosujemy stratę Briera
$S(q,Y)=(q-Y)^2$. Wtedy
\[
\E[S(q,Y)]=(q-p)^2+p(1-p).
\]
Jedynym optymalnym wyborem jest $q=p$.
\end{twierdzenie}
\begin{proof}
Ponieważ $Y^2=Y$, mamy
$\E(q-Y)^2=q^2-2qp+p=(q-p)^2+p-p^2$.
Zatem nadwyżka ryzyka względem prawidłowej prognozy
jest dokładnie kwadratem błędu prognozowanego
prawdopodobieństwa.
\end{proof}

\begin{twierdzenie}[Reguła logarytmiczna i dywergencja KL]
\label{thm:log-score-bayes}
Niech $\mathsf Y=\{1,\ldots,k\}$, $k\ge2$, a
$P=(p_1,\ldots,p_k)$ będzie rzeczywistym rozkładem przyszłej
obserwacji. Dla prognozy $Q=(q_1,\ldots,q_k)$ z $q_j>0$
gdy $p_j>0$, strata logarytmiczna ma postać
$S(Q,Y)=-\log q_Y$. Wówczas
\[
\E_P S(Q,Y)-\E_P S(P,Y)
=\sum_{j:p_j>0}p_j\log\frac{p_j}{q_j}
= D_{\mathrm{KL}}(P\Vert Q)\ge0,
\]
a równość zachodzi dokładnie, gdy $Q=P$.
\end{twierdzenie}
\begin{proof}
Pierwsza równość to bezpośrednie odjęcie skończonych sum.
Do dowodu nierówności stosujemy $\log t\le t-1$ do
$t=q_j/p_j$ dla $p_j>0$:
$\sum_{p_j>0}p_j\log(q_j/p_j)
\le\sum_{p_j>0}(q_j-p_j)\le0$.
Jeśli zachodzi równość, to każdy składnik wymusza
$q_j=p_j$ dla $p_j>0$; normalizacja wymusza też
$q_j=0$ poza nośnikiem $P$.
\end{proof}

\begin{uwaga}[Walidacja predykcji]
Porównując modele na wspólnym zbiorze testowym, możemy
obliczyć średnią stratę Briera dla prognoz prawdopodobieństw
lub średnią stratę logarytmiczną. Nie należy dobierać priorów
i parametrów na tym samym zbiorze, który ma następnie służyć
za niezależną ocenę uogólnienia. Przy szeregach czasowych
konieczne jest zachowanie chronologii walidacji.
\end{uwaga}

\section{Przedział wiarygodności, przedział ufności i przedział predykcyjny}
\begin{przyklad}[Trzy różne pytania dla siedmiu wad na sto]
Załóżmy $S=7$, $n=100$ w modelu Bernoulliego. Estymator MLE
odsetka wad to $\widehat p=0{,}07$. Klasyczny 95\% przedział
Wilsona z rozdziału~\ref{ch:confidence-intervals} jest w przybliżeniu
$[0{,}03432,0{,}13750]$. Przy priorze
$\operatorname{Beta}(1,1)$ posterior to
$\operatorname{Beta}(8,94)$, którego centralny 95\%
przedział wiarygodności wynosi w przybliżeniu
$[0{,}03482,0{,}13760]$.
Oba przedziały są liczbami bliskie, ale oznaczają co innego:
pierwszy ma (przy przyjętym modelu) przybliżoną gwarancję
powtarzalnego pokrycia, drugi ma posteriorową masę $0{,}95$
po obserwacji danych. Predykcja liczby wad w kolejnych
dziesięciu sztukach nie jest żadnym z tych przedziałów:
ma rozkład beta--dwumianowy, a
$\Prob(T=0\mid S=7)=B(8,104)/B(8,94)\approx0{,}45772$.
Do planowania liczby kontroli lub zapasu części należy używać
rozkładu predykcyjnego, a nie wyłącznie niepewności parametru.
\end{przyklad}

\section{Sprawdzanie modelu przez replikację danych}
Analiza posterioru zakłada, że określony model generowania
danych jest wystarczająco wiarygodny. \emph{Posterior predictive
check} (PPC) polega na wylosowaniu $\theta\sim\Pi(\cdot\mid x)$,
a potem nowych danych $X^{\rm rep}\sim P_\theta$ oraz
porównaniu statystyki diagnostycznej $T(X^{\rm rep})$ z $T(x)$.
Statystyka $T$ powinna celować w charakterystyczne odchylenie:
liczbę zer, rozrzut, ogony, autokorelację lub zależność między
grupami. Nie jest to automatycznie klasyczny test o jednostajnym
rozkładzie wartości p.

\begin{przyklad}[Dlaczego posterior predictive p-value nie jest jednostajna?]
Niech $p\sim\operatorname{Beta}(1,1)$ i
$X\mid p\sim\operatorname{Bernoulli}(p)$, zaś
$X^{\rm rep}\mid p$ niech będzie warunkowo niezależną kopią.
Jako statystykę przyjmijmy $T(X)=X$ i zdefiniujmy
$q(X)=\Prob(X^{\rm rep}\ge X\mid X)$.
Dla $X=0$ zachodzi $q(0)=1$. Dla $X=1$
posterior jest $\operatorname{Beta}(2,1)$, więc
$q(1)=\E[p\mid X=1]=2/3$.
W rozkładzie predykcyjnym przed obserwacją
$\Prob(X=0)=\Prob(X=1)=1/2$, stąd $q(X)$ przyjmuje
wyłącznie wartości $1$ i $2/3$, każdą z prawdopodobieństwem
$1/2$. Nie jest jednostajna na $[0,1]$.
To konkretny powód, aby nie interpretować automatycznie PPC
jako standardowej wartości p z gwarancją poziomu testu.
\end{przyklad}

\begin{stwierdzenie}[Rozkład predykcyjny beta--dwumianowy przed obserwacją]
\label{prop:prior-predictive-betabinomial}
Niech $p\sim\operatorname{Beta}(a,b)$, $a,b>0$,
a warunkowo przy znanym $p$ zmienna $S$ ma rozkład
$\Bin(n,p)$, $n\ge1$. Wtedy dla $s=0,\ldots,n$
\[
\Prob(S=s)=\binom ns\frac{B(a+s,b+n-s)}{B(a,b)}.
\]
Przy priorze jednostajnym $\operatorname{Beta}(1,1)$
rozkład $S$ jest jednostajny na $\{0,\ldots,n\}$.
\end{stwierdzenie}
\begin{proof}
Pierwszy wzór wynika z całkowania
$\binom ns p^s(1-p)^{n-s}$ względem gęstości prioru,
co daje iloraz całek beta. Dla $a=b=1$ korzystamy
z $B(s+1,n-s+1)=s!(n-s)!/(n+1)!$ i otrzymujemy
$\binom ns B(s+1,n-s+1)=1/(n+1)$.
\end{proof}
\begin{uwaga}
Jednostajność liczby sukcesów przed obserwacją \emph{nie oznacza},
że wszystkie wartości $p$ są równie wiarygodne po zobaczeniu danych.
Rozkład predykcyjny przed obserwacją może ujawnić, czy prior
przewiduje realistyczne dane: w modelu beta--dwumianowym
$\operatorname{Beta}(1,1)$ dopuszcza bardzo duży rozrzut liczby
wad, mimo że pojedyncza obserwacja ma prawdopodobieństwo
wady $1/2$. Przy aktualnej wiedzy o rzadkich awariach
taki prior może być nieodpowiedni.
\end{uwaga}

\section{Praktyczny schemat: od celu do metody}
\begin{center}\small
\begin{tabular}{>{\raggedright\arraybackslash}p{.27\textwidth}>{\raggedright\arraybackslash}p{.31\textwidth}>{\raggedright\arraybackslash}p{.27\textwidth}}
\toprule
\textbf{Pytanie} & \textbf{Właściwe narzędzie} & \textbf{Co koniecznie sprawdzić}\tabularnewline
\midrule
Jaka jest nieznana wartość parametru? & Średnia lub kwantyl posterioru, zależnie od straty & Rozkład a priori i identyfikowalność.\\
Ile awarii będzie w kolejnej partii? & Rozkład posterior predictive, jego ogon i kwantyle & Ekspozycja i zmienność między okresami.\\
Jak dobrać zapas? & Kwantyl predykcyjny przy asymetrycznej stracie & Rzeczywiste koszty braków i nadmiaru.\\
Który model lepiej prognozuje? & Brier lub log-loss na danych testowych & Przeciek danych i zależność czasowa.\\
Czy model jest wiarygodny? & PPC, diagnostyka reszt, walidacja predykcyjna & PPC nie jest automatycznie testem poziomu $\alpha$.\\
Czy mamy gwarancję pokrycia dla każdego parametru? & Klasyczny przedział ufności, np. Wilsona & Rodzaj gwarancji i założenia losowania.\\
\bottomrule
\end{tabular}
\end{center}

\Needspace{26\baselineskip}
\begin{lstlisting}[language=Python,caption={Kompletny program: estymacja, predykcja, ryzyko alarmu i przedziały (Python, SciPy).}]
from scipy.stats import beta, betabinom, nbinom

n, defects = 100, 7
a, b = 1, 1
post_a, post_b = a + defects, b + n - defects
print("Posterior mean:", post_a / (post_a + post_b))
print("95% credible interval:", beta.ppf([0.025, 0.975], post_a, post_b))

# Defects in the next ten independent devices.
future_n = 10
print("P(no defects):", betabinom.pmf(0, future_n, post_a, post_b))
print("P(at least two defects):", betabinom.sf(1, future_n, post_a, post_b))
print("Predictive 90% quantile:", betabinom.ppf(0.9, future_n, post_a, post_b))

# Failures in 50 hours, given posterior Gamma(shape=6, rate=300).
shape, rate, exposure = 6, 300, 50
prob_3plus = nbinom.sf(2, shape, rate / (rate + exposure))
print("P(three or more failures):", prob_3plus)

# A warning costs 3 if unjustified; a missed event costs 50.
false_alarm, missed_event = 3.0, 50.0
threshold = false_alarm / (false_alarm + missed_event)
print("Alarm threshold:", threshold, "Warn:", prob_3plus > threshold)
\end{lstlisting}

\begin{zadanie}[Rezerwa dla nierównej straty]
W posteriorze $p\mid x\sim\operatorname{Beta}(8,94)$ chcemy
podać jedną liczbę reprezentującą $p$. Za jednostkową wartość
niedoszacowania płacimy dziewięć razy więcej niż za przeszacowanie.
Jaki estymator Bayesa stosujemy?
\par\emph{Rozwiązanie.} Korzystamy z asymetrycznej straty
liniowej o $\tau=9/(9+1)=0{,}9$. Twierdzenie
\ref{thm:bayes-quantile-loss} daje $F_{\operatorname{Beta}(8,94)}^{-1}(0{,}9)$.
Średnia $8/102$ minimalizowałaby inną, kwadratową stratę,
więc nie jest tu rozwiązaniem zadania decyzyjnego.
\end{zadanie}

% ================================================================
% ROZDZIAL 20 -- BAYESOWSKIE TESTOWANIE HIPOTEZ
% ================================================================
\chapter{Bayesowskie testowanie hipotez, iloraz Bayesa i paradoksy}
\label{ch:bayes-model-testing}
\paragraph{Powiązania i odsyłacze.} Porównaj definicję klasycznej wartości p w rozdziale~\ref{ch:testing} i decyzje bayesowskie w rozdziale~\ref{ch:bayes-estimation-prediction}.


\section{Prawdopodobieństwo hipotezy a zgodność danych z hipotezą}
W testowaniu klasycznym z rozdziału~\ref{ch:testing}
obliczamy prawdopodobieństwo odpowiednio skrajnych danych
\emph{pod warunkiem} hipotezy zerowej. Bayesowskie podejście
porównuje prawdopodobieństwa samych modeli \emph{po danych},
co wymaga zadania rozkładów a priori również na \emph{modele}.
Są to inne wielkości. W szczególności wartość p nie jest
prawdopodobieństwem prawdziwości $H_0$, a iloraz Bayesa nie
jest bez dodatkowych założeń prawdopodobieństwem błędu.

\begin{definicja}[Porównanie dwóch modeli i iloraz Bayesa]
\label{def:bayes-factor}
Niech $\mathcal M_0$ i $\mathcal M_1$ będą dwoma modelami
obserwacji $X$ w tej samej przestrzeni mierzalnej, a
$w_0,w_1>0$, $w_0+w_1=1$, ich wcześniejszymi
prawdopodobieństwami. W modelu $j\in\{0,1\}$
niech $\Theta_j$ będzie przestrzenią parametrów,
$\pi_j$ \emph{właściwym} priorem i
$p_j(x\mid\theta_j)$ gęstością lub funkcją masy
względem tej samej, ustalonej miary $\nu$ na przestrzeni
danych. Definiujemy \emph{zintegrowaną wiarygodność}
$m_j(x)=\int_{\Theta_j}p_j(x\mid\theta_j)\,\dd\pi_j(\theta_j)$.
Jeżeli $0<m_j(x)<\infty$ dla obu $j$, to
$BF_{10}(x)=m_1(x)/m_0(x)$ nazywamy \emph{ilorazem
Bayesa} na korzyść modelu $\mathcal M_1$. Analogicznie
$BF_{01}=1/BF_{10}$.
\end{definicja}

\begin{twierdzenie}[Posterior odds = prior odds razy Bayes factor]
\label{thm:bayes-odds}
W założeniach definicji~\ref{def:bayes-factor}, dla danych
$x$ z $0<w_0m_0(x)+w_1m_1(x)<\infty$ mamy
\[
\frac{\Prob(\mathcal M_1\mid X=x)}
     {\Prob(\mathcal M_0\mid X=x)}
=BF_{10}(x)\,\frac{w_1}{w_0}.
\]
W szczególności $\Prob(\mathcal M_1\mid x)
=w_1m_1(x)/(w_0m_0(x)+w_1m_1(x))$.
\end{twierdzenie}
\begin{proof}
Zmienna indeksująca model ma rozkład dyskretny
$\Prob(J=j)=w_j$, a warunkowa gęstość danych przy $J=j$
to $m_j(x)$ po scałkowaniu parametru $\theta_j$.
Wzór Bayesa dla tej dwupunktowej przestrzeni
parametrów daje $\Prob(J=j\mid x)=w_jm_j(x)/
(w_0m_0(x)+w_1m_1(x))$. Podzielenie odpowiednich
tożsamości kończy dowód.
\end{proof}

\begin{uwaga}[Znaczenie doboru priorów]
Iloraz Bayesa zależy od \emph{rozkładów parametrów wewnątrz
modeli}, nie tylko od ich najlepszych dopasowań. Zastąpienie
$\pi_1$ rozkładem znacznie szerszym może zmienić
$BF_{10}$ nawet wtedy, gdy maksimum wiarygodności i
obserwowane dane pozostają takie same. Priory niewłaściwe
nie pozwalają na zwykłe porównanie modeli z niezależnymi
nieustalonymi stałymi normalizacji
(por. rozdział~\ref{ch:bayes-foundations}).
\end{uwaga}

\section{Dokładny test punktowej hipotezy o odsetku wad}
\begin{twierdzenie}[Iloraz Bayesa: punktowy Bernoulli kontra beta]
\label{thm:bernoulli-point-bf}
Niech $n\ge1$ i $X_1,\ldots,X_n\in\{0,1\}$. Porównujemy
$\mathcal M_0: X_i\iid\operatorname{Bernoulli}(p_0)$,
gdzie $0<p_0<1$, z
$\mathcal M_1: p\sim\operatorname{Beta}(a,b)$,
$a,b>0$, i $X_i\mid p\iid\operatorname{Bernoulli}(p)$.
Jeżeli $S=\sum_iX_i=s$, to
\[
BF_{10}(s)=\frac{B(a+s,b+n-s)}{B(a,b)\,p_0^s(1-p_0)^{n-s}}.
\]
Ten sam wzór otrzymujemy z obserwacji pełnego wektora
lub wyłącznie sumy $S$.
\end{twierdzenie}
\begin{proof}
Dla pełnej próbki $x=(x_1,\ldots,x_n)$ jej funkcja masy
w modelu zerowym wynosi
$p_0^s(1-p_0)^{n-s}$, a w alternatywnym całka
$\int_0^1p^s(1-p)^{n-s}\,\dd\operatorname{Beta}(a,b)(p)$.
Całka jest równa $B(a+s,b+n-s)/B(a,b)$.
Gdy obserwujemy tylko sumę $S$, obie funkcje masy mają
dodatkowy wspólny czynnik $\binom ns$, który skraca się
w ilorazie. W tym modelu $S$ jest statystyką dostateczną,
a zatem agregacja nie gubi informacji o $p$.
\end{proof}

\begin{przyklad}[Kontrola jakości: siedem wad na sto]
Testujemy $\mathcal M_0:p=0{,}10$ wobec alternatywy
$\mathcal M_1:p\sim\operatorname{Beta}(1,1)$, używając
$S=7$, $n=100$. Z powyższego wzoru
$BF_{10}\approx0{,}11138$, a $BF_{01}\approx8{,}9784$.
Przy równych priorowych prawdopodobieństwach modeli
$\Prob(\mathcal M_0\mid S=7)\approx0{,}8998$.
Nie wolno jednak opisać tego jako prawdopodobieństwa, że
$H_0$ jest prawdziwa \emph{bez podania priorycznego
prawdopodobieństwa samej hipotezy punktowej}.
Dla alternatywy $\operatorname{Beta}(2,18)$ otrzymujemy
$BF_{10}\approx0{,}70745$, czyli znacznie inny wynik,
choć obserwacje i $p_0$ się nie zmieniły.
Przy alternatywie $\operatorname{Beta}(20,80)$
$BF_{10}\approx0{,}03371$. Widać stąd konieczność
analizy wrażliwości na prior alternatywy.
\end{przyklad}

\begin{uwaga}[Hipoteza punktowa potrzebuje dodatniej masy modelowej]
Jeżeli jedynym priorem na ciągły parametr $p$ jest rozkład
beta, to $\Prob(p=p_0)=0$ przed obserwacją i po niej,
niezależnie od danych o dodatniej wiarygodności.
Aby mówić o dodatnim posteriorowym prawdopodobieństwie
punktowej hipotezy $p=p_0$, trzeba wprowadzić model
przypisujący jej dodatnią masę, np. prior mieszany
$w_0\delta_{p_0}+w_1\operatorname{Beta}(a,b)$.
Takie porównanie modeli nazywa się niekiedy priorem
\emph{spike-and-slab}. Nie należy mylić z nim posteriorowego
prawdopodobieństwa nierówności $p>p_0$, które możemy obliczać
bez atomu w $p_0$.
\end{uwaga}

\section{Decyzja oparta na posteriorowych prawdopodobieństwach modeli}
\begin{twierdzenie}[Próg decyzji przy nierównych kosztach błędów]
\label{thm:bayes-hypothesis-decision}
Załóżmy model z twierdzenia~\ref{thm:bayes-odds}.
Wybieramy pomiędzy działaniami $a_0$ (uznajemy model
$\mathcal M_0$) i $a_1$ (uznajemy $\mathcal M_1$).
Strata wynosi zero przy poprawnym wyborze,
$C_{10}>0$ przy wyborze $a_1$, gdy prawdziwy jest
$\mathcal M_0$, i $C_{01}>0$ przy wyborze $a_0$,
gdy prawdziwy jest $\mathcal M_1$. Wtedy $a_1$
ściśle minimalizuje stratę posteriorową wtedy i tylko wtedy, gdy
\[
 BF_{10}(x)>\frac{C_{10}}{C_{01}}\frac{w_0}{w_1}.
\]
Przy równości obie decyzje są optymalne.
\end{twierdzenie}
\begin{proof}
Posteriorowa strata $a_1$ równa się
$C_{10}\Prob(\mathcal M_0\mid x)$, a strata $a_0$ wynosi
$C_{01}\Prob(\mathcal M_1\mid x)$.
Porównanie tych liczb i podstawienie ilorazu z
 twierdzenia~\ref{thm:bayes-odds} daje tezę.
\end{proof}

\begin{przyklad}[Jak koszt błędnej decyzji zmienia próg?]
Niech $w_0=w_1=1/2$. Jeśli koszt fałszywego odrzucenia
$\mathcal M_0$ wynosi $10$ jednostek, a koszt błędnego
zaakceptowania $\mathcal M_0$ wynosi $2$, to wybieramy
$\mathcal M_1$ dopiero przy $BF_{10}>5$.
Żaden uniwersalny próg, np. $BF_{10}>3$ lub $BF_{10}>10$,
nie zastąpi rzeczywistej analizy kosztów.
\end{przyklad}

\section{Paradoks Jeffreysa--Lindleya: to samo z, różny wniosek}
\begin{twierdzenie}[Iloraz Bayesa w normalnym teście średniej]
\label{thm:jeffreys-lindley}
Niech $n\ge1$ oraz $X_1,\ldots,X_n\mid\mu$ będą
niezależne o rozkładzie $\mathcal N(\mu,\sigma^2)$
ze znanym $\sigma>0$. Porównujemy
$\mathcal M_0:\mu=0$ z
$\mathcal M_1:\mu\sim\mathcal N(0,\tau^2)$,
$\tau>0$, przy tym samym rozkładzie warunkowym
obserwacji. Niech $z=\sqrt n\,\overline x/\sigma$
i $r=n\tau^2/\sigma^2$. Wtedy
\[
 BF_{01}(x)=\sqrt{1+r}\,
 \exp\!\left(-\frac{r}{1+r}\frac{z^2}{2}\right).
\]
Dla ustalonego $z\in\R$, $\sigma>0$, $\tau>0$ i $n\to\infty$
wartość $BF_{01}(x)$ dąży do nieskończoności, chociaż
klasyczna dwustronna wartość p równa
$2(1-\Phi(|z|))$ pozostaje stała.
\end{twierdzenie}
\begin{proof}
Przy znanej wariancji $\sigma^2$ resztowa suma kwadratów
$\sum_i(x_i-\overline x)^2$ ma ten sam czynnik gęstości
w obu porównywanych modelach i skraca się w ilorazie.
Wystarczy zatem porównać gęstości $\overline X$.
W modelu zerowym $\overline X\sim\mathcal N(0,\sigma^2/n)$.
W alternatywnym $\overline X=\mu+(\overline X-\mu)$,
więc po scałkowaniu normalnego prioru mamy
$\overline X\sim\mathcal N(0,\sigma^2/n+\tau^2)$.
Iloraz gęstości dwóch rozkładów normalnych w punkcie
$\overline x$ daje podany wzór.
Gdy $n\to\infty$ przy ustalonym $z$, to $r\to\infty$,
czynnik wykładniczy zbiega do $e^{-z^2/2}>0$,
a pierwiastek $\sqrt{1+r}\to\infty$. Natomiast rozkład
statystyki $Z=\sqrt n\,\overline X/\sigma$ pod $\mathcal M_0$
jest standardowo normalny dla każdego $n$, więc
wartość p zależy wyłącznie od $|z|$.
\end{proof}

\begin{przyklad}[Istotność 5\% i coraz większa przewaga modelu zerowego]
Przyjmijmy $\sigma=\tau=1$ i zaobserwowane $z=2$.
Dwustronna klasyczna wartość p to
$2(1-\Phi(2))\approx0{,}04550$ dla każdego $n$.
Jednak przy $n=100$, $1000$ i $10000$ dostajemy
odpowiednio $BF_{01}\approx1{,}3873$, $4{,}2904$ oraz
$13{,}5369$. Przy priorowych szansach $1:1$
posteriorowe prawdopodobieństwa modelu zerowego wynoszą
odpowiednio około $0{,}581$, $0{,}811$ i $0{,}931$.
To nie jest sprzeczność rachunkowa. Test klasyczny i iloraz
Bayesa odpowiadają na różne pytania, a przy wzroście $n$
zachowujemy tę samą \emph{wartość statystyki z}, czyli
obserwowana średnia maleje jak $n^{-1/2}$.
Nie należy mylić tej konstrukcji z zachowaniem
stałej, niezerowej wartości prawdziwego parametru $\mu$.
\end{przyklad}

\begin{center}
\begin{minipage}{.82\textwidth}
\centering
\begin{tikzpicture}[>=Latex,scale=.88,every node/.style={font=\small}]
\draw[->,thick] (0,0)--(6.7,0) node[right] {$n$};
\draw[->,thick] (0,0)--(0,5.15) node[above] {$BF_{01}$};
\foreach \x/\lab in {1/100,3/1000,5/10000} {
  \draw (\x,.08)--(\x,-.08) node[below=3pt] {\lab};
}
\foreach \y/\lab in {0/0,1.5/5,3/10,4.5/15} {
  \draw (.08,\y)--(-.08,\y) node[left=3pt] {\lab};
}
\draw[very thick,blue!60!black] (1,.4162)--(3,1.2871)--(5,4.0611);
\foreach \x/\y in {1/.4162,3/1.2871,5/4.0611} {
  \fill[blue!60!black] (\x,\y) circle (2.4pt);
}
\node[anchor=west,font=\footnotesize,fill=white,inner sep=2pt] at (0.45,4.65)
 {$z=2$, $\sigma=\tau=1$};
\end{tikzpicture}

{\small\emph{Rysunek 20.1.} Przy stałej statystyce $z=2$
wartość p pozostaje równa około $0{,}0455$, podczas gdy
$BF_{01}$ rośnie. Punkty oznaczają $n=100,1000,10000$;
połączono je tylko dla czytelności.}
\end{minipage}
\end{center}

\begin{uwaga}[Zbyt szeroka alternatywa — paradoks Bartletta]
W twierdzeniu~\ref{thm:jeffreys-lindley} także przy
ustalonym $n$ i zaobserwowanym $z$ zachodzi
$BF_{01}\to\infty$ dla $\tau\to\infty$.
Zwiększenie rozproszenia prioru alternatywy obniża jej
zintegrowaną wiarygodność w konkretnym punkcie danych,
choć nie zmienia wartości maksymalnej wiarygodności.
To ważny argument przeciw bezrefleksyjnemu określaniu
bardzo szerokiego prioru jako „neutralnego” w porównywaniu
modeli. Granica $\tau\to\infty$ nie uprawnia do zastąpienia
prioru właściwego priorem niewłaściwym w ilorazie Bayesa.
\end{uwaga}

\section{Sekwencyjne testowanie: kiedy wolno patrzeć na dane wielokrotnie?}
W klasycznym testowaniu bez korekty wielokrotne sprawdzanie,
czy wartość p przekroczyła próg, może znacząco zwiększać
częstość błędu pierwszego rodzaju. W specjalnej sytuacji
\emph{prostej} hipotezy zerowej iloraz Bayesa ma przydatną
własność martyngałową, lecz należy precyzyjnie sprawdzić
warunki, zanim wykorzystamy ją do zatrzymywania eksperymentu.

\begin{twierdzenie}[Iloraz Bayesa jako nieujemny martyngał]
\label{thm:bayes-factor-martingale}
Niech $(\mathsf S,\mathcal S,\nu)$ będzie przestrzenią
z miarą $\sigma$-skończoną. Przy $\mathcal M_0$ zmienne
$X_1,X_2,\ldots$ są niezależne z tą samą dodatnią
$\nu$-prawie wszędzie gęstością $f_0$.
Przy alternatywie $\mathcal M_1$ ustalamy właściwy prior
$\pi$ na przestrzeni parametrów $\Theta_1$ i
wspólnie mierzalne gęstości $f_\theta$ względem $\nu$;
przy znanym $\theta$ obserwacje są niezależne o gęstości
$f_\theta$. Niech $\mathcal F_n=\sigma(X_1,\ldots,X_n)$,
$\mathcal F_0=\{\varnothing,\Omega\}$, oraz
\[
 B_n=\int_{\Theta_1}\prod_{i=1}^n
    \frac{f_\theta(X_i)}{f_0(X_i)}\,\dd\pi(\theta),
 \qquad B_0=1.
\]
Wtedy przy $\mathcal M_0$ proces $(B_n)_{n\ge0}$ jest
nieujemnym martyngałem, $\E_0 B_n=1$, oraz dla
każdego $\alpha\in(0,1)$ zachodzi
\[
 \Prob_0\!\left(\sup_{n\ge0}B_n\ge\frac1\alpha\right)
 \le\alpha.
\]
\end{twierdzenie}
\begin{proof}
Dla ustalonego $\theta$ wyrażenie
$L_n(\theta)=\prod_{i=1}^n f_\theta(X_i)/f_0(X_i)$
ma przy $\mathcal M_0$ wartość oczekiwaną jeden, ponieważ
$\int_{\mathsf S} f_\theta\,\dd\nu=1$.
Z niezależności nowej obserwacji wynika
$\E_0[L_{n+1}(\theta)\mid\mathcal F_n]=L_n(\theta)$.
Tonelli i warunkowa wersja tego twierdzenia pozwalają
zamienić całkę po $\pi$ z warunkową wartością oczekiwaną.
Otrzymujemy $\E_0[B_{n+1}\mid\mathcal F_n]=B_n$,
co wraz z $\E_0B_n=1$ daje pierwszą tezę.
Dla $a>0$ i $N\in\N$ zatrzymajmy proces w
$\tau=\inf\{n\ge0:B_n\ge a\}\wedge N$.
Twierdzenie o zatrzymywaniu dla ograniczonego czasu daje
$1=\E_0 B_\tau\ge a\Prob_0(\max_{0\le k\le N}B_k\ge a)$.
Po podstawieniu $a=1/\alpha$ i przejściu $N\to\infty$
przez ciągłość prawdopodobieństwa dla rosnących zdarzeń
otrzymujemy tezę.
\end{proof}

\begin{uwaga}[Dokładna granica zastosowania]
Próg $BF_{10}\ge1/\alpha$ zapewnia w powyższym
\emph{prostym} modelu zerowym kontrolę prawdopodobieństwa
fałszywego przekroczenia poziomu $\alpha$, nawet przy
kolejnym sprawdzaniu danych. Nie jest to twierdzenie
o dowolnej złożonej hipotezie zerowej, priorycznych
prawdopodobieństwach modeli ani o poziomie klasycznego
testu opartego na zwykłej wartości p. Także zmiana prioru
alternatywy w reakcji na bieżące dane narusza przyjętą
konstrukcję martyngału. Do zastosowań konieczne są
wstępnie ustalone modele i reguły raportowania.
\end{uwaga}

\begin{przyklad}[Adaptacyjny prior niszczy gwarancję sekwencyjną]
\label{ex:adapt-bf-fails}
Niech pod hipotezą zerową $X\sim\operatorname{Bernoulli}(1/2)$.
Rozpatrzmy dwie dopuszczalne alternatywy punktowe:
$P_0(X=0)=1$ oraz $P_1(X=1)=1$. Jeśli \emph{przed}
obserwacją ustalimy jedną z nich, iloraz wiarygodności
$L(X)=P_j(\{X\})/(1/2)$ ma średnią jeden pod hipotezą
zerową. Jeśli jednak \emph{po} zobaczeniu $X$ wybierzemy
alternatywę $P_X$, to otrzymamy $L(X)=2$ przy każdym wyniku,
a zatem $\E_0 L(X)=2$, nie zaś jeden.
Przy progu $1/\alpha=2$ (czyli $\alpha=1/2$) fałszywe
przekroczenie występuje z prawdopodobieństwem jeden.
Jest to konkretny kontrprzykład pokazujący, dlaczego prior
alternatywy w twierdzeniu~\ref{thm:bayes-factor-martingale}
musi być zadany przed obserwacją, a nie dobrany do danych.
\end{przyklad}

\section{Jak w praktyce wybierać między testem klasycznym a Bayes factor?}
\begin{center}\small
\begin{tabular}{>{\raggedright\arraybackslash}p{.27\textwidth}>{\raggedright\arraybackslash}p{.31\textwidth}>{\raggedright\arraybackslash}p{.27\textwidth}}
\toprule
\textbf{Cel analizy} & \textbf{Narzędzie} & \textbf{Interpretacja i ograniczenie}\\
\midrule
Kontrola błędu I rodzaju przy ustalonym $H_0$ & Test klasyczny i wartość p & Poziom obowiązuje przy warunkach testu.\\
Porównanie dwóch opisanych modeli & $BF_{10}$ i posterior modelowy & Trzeba podać oba priory i szanse modeli.\\
Ocena $\Prob(p>p_0\mid x)$ & Posterior ciągły i ogon rozkładu & Nie jest to $BF$ hipotezy punktowej.\\
Decyzja o koszcie różnych błędów & Posterior odds i funkcja straty & Próg zależy od rzeczywistych kosztów.\\
Wielokrotne sprawdzanie w czasie & Z góry ustalony proces ilorazu wiarygodności & Kontrola przekroczenia dla prostej $H_0$; nie dowolny BF.\\
Słaby prior alternatywy & Analiza wrażliwości $BF$ & Zbyt szeroki prior może sprzyjać $H_0$.\\
\bottomrule
\end{tabular}
\end{center}

\Needspace{23\baselineskip}
\begin{lstlisting}[language=Python,caption={Kompletny program: iloraz Bayesa dla wadliwych produktów i paradoks Jeffreysa--Lindleya (Python, SciPy).}]
from math import exp, log, log1p, sqrt
from scipy.special import betaln
from scipy.stats import norm

n, s, p0 = 100, 7, 0.10
for a, b in [(1, 1), (2, 18), (20, 80)]:
    log_bf10 = (betaln(a+s, b+n-s) - betaln(a, b)
                - s*log(p0) - (n-s)*log1p(-p0))
    bf10 = exp(log_bf10)
    post_prob_h0 = 1.0 / (1.0 + bf10)  # equal prior model odds
    print("Alternative Beta", (a, b),
          "BF10=", round(bf10, 5),
          "P(H0|data)=", round(post_prob_h0, 5))

sigma, tau, z = 1.0, 1.0, 2.0
print("Two-sided classical p-value:", 2*norm.sf(abs(z)))
for size in [100, 1000, 10000]:
    r = size * tau**2 / sigma**2
    bf01 = sqrt(1+r) * exp(-0.5*z*z*r/(1+r))
    print("n=", size, "BF01=", round(bf01, 5),
          "P(H0|data)=", round(bf01/(1+bf01), 5))
\end{lstlisting}

\begin{zadanie}[Kiedy alarm jest optymalny?]
Dwa modele mają priorowe prawdopodobieństwa $w_0=0{,}8$
i $w_1=0{,}2$. Otrzymano $BF_{10}=4$.
Fałszywe odrzucenie modelu $\mathcal M_0$ kosztuje $6$,
a przeoczenie modelu $\mathcal M_1$ kosztuje $10$.
Które działanie minimalizuje posteriorową stratę?
\par\emph{Rozwiązanie.} Posteriorowe szanse modelu
$\mathcal M_1$ wynoszą $4\cdot(0{,}2/0{,}8)=1$,
zatem oba modele mają posteriorowe prawdopodobieństwa po $1/2$.
Ryzyko działania $a_1$ wynosi $6\cdot1/2=3$,
a działania $a_0$ wynosi $10\cdot1/2=5$.
Optymalny jest wybór $a_1$. Równoważnie sprawdzamy
$BF_{10}>(6/10)(0{,}8/0{,}2)=2{,}4$.
\end{zadanie}

% ================================================================
% ROZDZIAL 21 -- MODELE HIERARCHICZNE
% ================================================================
\chapter{Modele hierarchiczne: częściowe łączenie informacji między grupami}
\label{ch:hierarchical-bayes}
\paragraph{Powiązania i odsyłacze.} Wzory dla modeli sprzężonych wyprowadzono w rozdziale~\ref{ch:conjugate-families}; teorię ryzyka i decyzji w rozdziałach~\ref{ch:decision} i~\ref{ch:bayes-estimation-prediction}.


\section{Kiedy oddzielne analizy grup zawodzą?}
Załóżmy, że chcemy porównać niezawodność dziesięciu maszyn.
Dla jednej zarejestrowano dwa zdarzenia, a dla innej dwieście.
Estymowanie każdej intensywności zupełnie osobno pomija możliwość,
że maszyny są technologicznie podobne. Zastąpienie wszystkich
intensywności jedną wspólną wartością ukrywa natomiast rzeczywistą
zmienność pomiędzy maszynami. \emph{Model hierarchiczny} opisuje te
źródła niepewności osobno.

Niech $J\ge2$ oznacza liczbę grup, $j\in\{1,\ldots,J\}$ numer grupy,
$\theta_j$ jej nieznany parametr, a $Y_j$ jej obserwację zbiorczą.
Poniżej rozróżniamy trzy strategie: \emph{brak łączenia}
(ang. \emph{no pooling}), \emph{pełne łączenie}
(\emph{complete pooling}) oraz \emph{częściowe łączenie}
(\emph{partial pooling}). Trzecia strategia nie jest prostą średnią
wyników: jej wagi zależą od precyzji obserwacji i od zmienności
międzygrupowej.

\begin{definicja}[Hierarchiczny model normalny]
\label{def:normal-hierarchical}
Ustalmy $J\ge2$, znane liczby $v_j>0$, $\mu\in\R$ i $\tau^2>0$.
Niech $\theta_1,\ldots,\theta_J$ będą warunkowo niezależne przy
$(\mu,\tau^2)$ i $\theta_j\mid\mu,\tau^2\sim\Normal(\mu,\tau^2)$.
Warunkowo względem wektora $\theta=(\theta_1,\ldots,\theta_J)$
niech $Y_1,\ldots,Y_J$ będą niezależne oraz
$Y_j\mid\theta\sim\Normal(\theta_j,v_j)$.
Na przykład, gdy grupa $j$ zawiera $n_j\ge1$ niezależnych
pomiarów $X_{ji}\mid\theta_j\sim\Normal(\theta_j,\sigma^2)$
przy znanym $\sigma^2>0$, możemy przyjąć
$Y_j=n_j^{-1}\sum_{i=1}^{n_j}X_{ji}$ i $v_j=\sigma^2/n_j$.
Liczbę $\tau^2$ nazywamy \emph{wariancją międzygrupową}.
\end{definicja}

Zadanie praktyczne wymaga rozróżnienia poziomów losowości:
$\theta_j$ opisuje grupę, $X_{ji}$ pojedynczy pomiar, a $\mu$
wspólną lokalizację rozkładu parametrów grup. Jeśli $\mu,\tau^2$
są znane, poszczególne grupy pozostają niezależne marginalnie;
nie jest to jednak prawdą po dodaniu wspólnego losowego hiperparametru
$\mu$ i jego scałkowaniu.

\section{Dokładny posterior i pochodzenie częściowego łączenia}
\begin{twierdzenie}[Posterior normalny i współczynnik ściągania]
\label{thm:normal-partial-pooling}
W modelu z definicji~\ref{def:normal-hierarchical}, przy ustalonych
$\mu,\tau^2$, dla każdego $j$ i $y\in\R^J$ zachodzi
\[
\begin{aligned}
 \theta_j\mid (Y=y,\mu,\tau^2)&\sim\Normal(m_j,s_j^2),
 & s_j^2&=\left(\frac1{\tau^2}+\frac1{v_j}\right)^{-1},\\
 m_j&=w_jy_j+(1-w_j)\mu,
 & w_j&=\frac{\tau^2}{\tau^2+v_j}.
\end{aligned}
\]
Współrzędne $\theta_j$ są także niezależne warunkowo względem
$Y=y,\mu,\tau^2$. Przy $v_j=\sigma^2/n_j$ waga $w_j$ rośnie
wraz z liczebnością $n_j$ i wraz ze zmiennością $\tau^2$.
\end{twierdzenie}
\begin{proof}
Warunkowa gęstość $\theta_j=t$ po zaobserwowaniu $Y_j=y_j$
jest proporcjonalna do
$\exp[-(t-\mu)^2/(2\tau^2)-(y_j-t)^2/(2v_j)]$.
Współczynnik przy $t^2$ po rozwinięciu wynosi
$\tau^{-2}+v_j^{-1}$, a liniowy
$\mu\tau^{-2}+y_jv_j^{-1}$. Uzupełnienie kwadratu
daje wariancję $s_j^2$ i średnią
$s_j^2(\mu/\tau^2+y_j/v_j)=m_j$.
Iloczynowy charakter prioru i wiarygodności dowodzi
warunkowej niezależności po aktualizacji. Zapis
$w_j=n_j\tau^2/(n_j\tau^2+\sigma^2)$ wyjaśnia monotoniczność.
\end{proof}

\begin{twierdzenie}[Zysk w średnim błędzie i jego ograniczenie]
\label{thm:shrinkage-bayes-risk}
W założeniach twierdzenia~\ref{thm:normal-partial-pooling}
porównajmy estymatory $\delta_{\rm odd}(Y_j)=Y_j$,
$\delta_{\rm wsp}(Y_j)=\mu$ i
$\delta_{\rm hier}(Y_j)=m_j$ przy stracie kwadratowej.
Ich ryzyka uśrednione względem prioru $\theta_j\sim\Normal(\mu,\tau^2)$
wynoszą odpowiednio $v_j$, $\tau^2$ oraz
$\tau^2v_j/(\tau^2+v_j)$.
Ostatnia liczba jest ściśle mniejsza od obu poprzednich.
Przy \emph{ustalonym} $\theta_j=t$ ryzyko częstościowe estymatora
hierarchicznego wynosi jednak
$(1-w_j)^2(t-\mu)^2+w_j^2v_j$ i może być większe od $v_j$.
\end{twierdzenie}
\begin{proof}
Estymator $Y_j$ ma błąd $Y_j-\theta_j$ o wariancji $v_j$.
Estymator stały $\mu$ ma średni błąd kwadratowy
$\E[(\theta_j-\mu)^2]=\tau^2$.
Z twierdzenia~\ref{thm:normal-partial-pooling} posteriorowy
błąd średniokwadratowy $m_j$ równa się $s_j^2$ niezależnie od
obserwacji, zatem prawo całkowitej wartości oczekiwanej
daje ryzyko Bayesa $s_j^2$.
Alternatywnie, dla ustalonego $t$ podstawiamy
$m_j-t=w_j(Y_j-t)+(1-w_j)(\mu-t)$ i wykorzystujemy
zerową średnią szumu oraz $\Var(Y_j\mid\theta_j=t)=v_j$.
Uśrednienie ostatniego wzoru po $t$ daje
$w_j^2v_j+(1-w_j)^2\tau^2=s_j^2$.
Z uwagi na $v_j,\tau^2>0$ mamy
$s_j^2<v_j$ i $s_j^2<\tau^2$. Ryzyko warunkowe rośnie
natomiast kwadratowo z $|t-\mu|$, co wyklucza dominację
przy wszystkich $t$ w tym jednowymiarowym porównaniu.
\end{proof}

\begin{przyklad}[Dwie grupy o różnej liczebności]
Niech $\mu=10$, $\tau^2=4$, $\sigma^2=16$.
Grupa A ma $n_A=4$ i średnią $Y_A=14$, grupa B ma
$n_B=64$ i średnią $Y_B=14$. W pierwszej grupie
$v_A=4$, $w_A=1/2$, $m_A=12$, $s_A^2=2$.
W drugiej $v_B=1/4$, $w_B=16/17$, $m_B=234/17\approx13{,}765$,
$s_B^2=4/17\approx0{,}235$.
Obie grupy mają tę samą średnią z próby, ale wynikiem
analizy hierarchicznej jest silniejsze ściągnięcie grupy
słabo reprezentowanej i znacznie mniejsza niepewność
dla grupy z wieloma pomiarami. Nie oznacza to automatycznie,
że dane z grupy A są błędne.
\end{przyklad}


\Needspace{92mm}
\begin{center}
\begin{tikzpicture}[>=Latex, every node/.style={font=\small},node distance=12mm]
\node[draw,rounded corners,fill=blue!5,minimum width=30mm,minimum height=8mm] (hyper) {$\mu,\tau^2$};
\node[draw,rounded corners,below left=13mm and 18mm of hyper,minimum width=25mm] (t1) {$\theta_1$};
\node[draw,rounded corners,below right=13mm and 18mm of hyper,minimum width=25mm] (t2) {$\theta_J$};
\node[draw,rounded corners,below=11mm of t1] (y1) {$Y_1$};
\node[draw,rounded corners,below=11mm of t2] (y2) {$Y_J$};
\draw[->] (hyper)--(t1); \draw[->] (hyper)--(t2);
\draw[->] (t1)--(y1); \draw[->] (t2)--(y2);
\node at ($(t1)!0.5!(t2)$) {$\cdots$};
\node[below=2mm of y1,font=\footnotesize] {$v_1$ -- znane};
\node[below=2mm of y2,font=\footnotesize] {$v_J$ -- znane};
\end{tikzpicture}

{\small\emph{Rysunek 21.1.} Wspólne hiperparametry wyznaczają rozkłady parametrów grup, a obserwacje informują o poszczególnych grupach.}
\end{center}

\section{Prognoza dla znanej i nowej grupy}
\begin{twierdzenie}[Dwie różne predykcje]
\label{thm:hierarchical-predictive}
W modelu~\ref{def:normal-hierarchical} niech przyszły pomiar
$Z_j\mid\theta_j\sim\Normal(\theta_j,\sigma^2)$ będzie warunkowo
niezależny od dotychczasowych danych dla znanej grupy $j$,
a $Z_{\rm nowy}\mid\theta_{\rm nowy}\sim\Normal(\theta_{\rm nowy},\sigma^2)$
dla nowej grupy, gdzie $\theta_{\rm nowy}\mid\mu,\tau^2\sim\Normal(\mu,\tau^2)$
i jest niezależny od istniejących grup. Wtedy
\[
 Z_j\mid Y=y,\mu,\tau^2\sim\Normal(m_j,\sigma^2+s_j^2),
 \qquad Z_{\rm nowy}\mid Y=y,\mu,\tau^2\sim\Normal(\mu,\sigma^2+\tau^2).
\]
\end{twierdzenie}
\begin{proof}
Nowy pomiar jest sumą $\theta_j$ i niezależnego
szumu $\Normal(0,\sigma^2)$. Po aktualizacji $\theta_j$
ma rozkład normalny ze średnią $m_j$ i wariancją $s_j^2$,
co daje pierwszy rozkład. Dla nowej grupy obserwacje
$Y$ nie aktualizują jej indywidualnego parametru przy
znanych hiperparametrach: pozostaje rozkład
$\Normal(\mu,\tau^2)$. Splot z szumem daje drugą tezę.
\end{proof}

\begin{uwaga}[Jednostka walidacji]
Jeżeli prognozujemy kolejny pomiar \emph{tej samej} maszyny,
możemy trenować na jej wcześniejszych wynikach. Jeśli chcemy
przewidywać wyniki \emph{zupełnie nowej} maszyny, walidacja
powinna odkładać całe maszyny (\emph{leave-one-group-out}).
Losowy podział pojedynczych pomiarów między trening i test
zwykle przecenia jakość przewidywania dla nowych grup.
\end{uwaga}

\section{Nieznane hiperparametry: pełny Bayes i empirical Bayes}
W praktyce $\mu$ i $\tau^2$ nie są znane. Dotychczasowe
wzory należy wtedy \emph{uśrednić po ich posteriorze},
a nie podstawiać oszacowania i udawać, że jest pewne.

\begin{twierdzenie}[Posterior hiperparametrów w modelu normalnym]
\label{thm:normal-hyperposterior}
Przyjmijmy model~\ref{def:normal-hierarchical} z nieznanymi
hiperparametrami $\mu\in\R$, $u=\tau^2>0$, znanymi
$v_j>0$ oraz niezależnymi, właściwymi priorami
$\mu\sim\Normal(m_0,V_0)$ i
$u\sim\operatorname{InvGamma}(a,b)$, gdzie $V_0,a,b>0$.
Przez $\operatorname{InvGamma}(a,b)$ rozumiemy rozkład
z gęstością $b^a u^{-a-1}e^{-b/u}/\Gamma(a)$ dla $u>0$
względem $\lambda_1$. Po scałkowaniu parametrów grup
$Y_j\mid\mu,u\sim\Normal(\mu,u+v_j)$ są niezależne,
a posterior hiperparametrów ma gęstość proporcjonalną do
\[
 \exp\!\left[-\frac{(\mu-m_0)^2}{2V_0}\right]
 u^{-a-1}e^{-b/u}\prod_{j=1}^J(u+v_j)^{-1/2}
 \exp\!\left[-\frac{(y_j-\mu)^2}{2(u+v_j)}\right].
\]
Dla ustalonego $u$ posterior $\mu\mid u,Y=y$ jest normalny
o wariancji $V(u)=(V_0^{-1}+\sum_j(u+v_j)^{-1})^{-1}$
i średniej
$M(u)=V(u)(m_0/V_0+\sum_jy_j/(u+v_j))$.
Posterior hiperparametrów jest właściwy dla każdego $y\in\R^J$.
\end{twierdzenie}
\begin{proof}
Po scałkowaniu niezależnych $\theta_j$ każda obserwacja
$Y_j$ jest sumą dwóch niezależnych zmiennych normalnych,
stąd wynika jej gęstość i niezależność warunkowa.
Mnożąc iloczyn tych gęstości przez priory otrzymujemy
pierwszy wzór. Jako funkcja $\mu$, wykładnik jest
ujemną połową dodatnio określonej funkcji kwadratowej;
uzupełnienie kwadratu daje $M(u)$ i $V(u)$.
Ponieważ $u+v_j\ge v_j>0$, marginalna wiarygodność
$\prod_j\varphi(y_j;\mu,u+v_j)$ jest ograniczona
z góry przez $\prod_j(2\pi v_j)^{-1/2}$.
Jej całka względem właściwego prioru jest skończona
i dodatnia, co dowodzi normalizowalności posterioru.
\end{proof}

\begin{stwierdzenie}[Dodatkowa niepewność hiperparametrów]
\label{prop:hier-variance-hyper}
Niech $H=(\mu,u)$ ma właściwy posterior z poprzedniego
twierdzenia i niech drugi moment $\theta_j\mid Y$ będzie skończony.
Wtedy
\[
 \E[\theta_j\mid Y]=\E[m_j(H)\mid Y],\qquad
 \Var(\theta_j\mid Y)=\E[s_j^2(H)\mid Y]
                       +\Var(m_j(H)\mid Y).
\]
Tu $m_j(H),s_j^2(H)$ otrzymujemy ze wzorów
z twierdzenia~\ref{thm:normal-partial-pooling} po podstawieniu
$\tau^2=u$. Ostatni składnik jest nieujemny i znika tylko,
gdy posteriorowa średnia grupy nie zależy losowo od $H$.
\end{stwierdzenie}
\begin{proof}
Pierwszy wzór to własność iterowanej warunkowej wartości
oczekiwanej, drugi to prawo całkowitej wariancji
z warunkowaniem najpierw po $H$, a następnie po $Y$.
\end{proof}

\begin{uwaga}[Kiedy wersja empirical Bayes może mylić?]
Metoda \emph{empirical Bayes} szacuje $H$ z danych, a potem
podstawia $\widehat H$ do wzorów dla $m_j$ i $s_j^2$.
Bywa szybka i użyteczna przy wielu grupach, ale naiwny
przedział $m_j(\widehat H)\pm1{,}96s_j(\widehat H)$
nie obejmuje niepewności oszacowania $H$. Pełny model
bayesowski ją uwzględnia, choć rezultat nadal zależy od
modelu i hiperprioru. Przy kilku grupach $u$ jest często
słabo zidentyfikowany; trzeba obejrzeć jego posterior,
a nie tylko wybraną estymatę punktową.
\end{uwaga}

\section{Zliczenia i udziały: hierarchiczny model beta--dwumianowy}
\begin{twierdzenie}[Łączenie informacji o rzadkich wadach]
\label{thm:hierarchical-beta-binomial}
Niech $J\ge2$, $\kappa>0$, $m\in(0,1)$.
Warunkowo przy znanych $m,\kappa$ zmienne $p_j$ są niezależne
i $p_j\sim\operatorname{Beta}(\kappa m,\kappa(1-m))$.
Dla każdego $j$ niech $S_j\mid p_j\sim\Bin(n_j,p_j)$,
$n_j\ge1$, i niech zliczenia są warunkowo niezależne.
Wtedy dla $s_j\in\{0,\ldots,n_j\}$
\[
 p_j\mid S_j=s_j,m,\kappa
 \sim\operatorname{Beta}(\kappa m+s_j,\kappa(1-m)+n_j-s_j),
\]
\[
 \E[p_j\mid S_j=s_j,m,\kappa]
 =\frac{n_j}{n_j+\kappa}\frac{s_j}{n_j}
 +\frac{\kappa}{n_j+\kappa}m.
\]
Jeżeli w jednej grupie obserwujemy wskaźniki
$X_{j1},\ldots,X_{jn_j}$ niezależne warunkowo przy $p_j$
i o rozkładzie $\operatorname{Bernoulli}(p_j)$, to dla $i\ne k$
\[
 \Cov(X_{ji},X_{jk})=\frac{m(1-m)}{\kappa+1},\quad
 \Var(S_j)=n_jm(1-m)\left(1+\frac{n_j-1}{\kappa+1}\right).
\]
\end{twierdzenie}
\begin{proof}
Iloczyn wiarygodności dwumianowej i gęstości beta
jest proporcjonalny do
$p_j^{\kappa m+s_j-1}(1-p_j)^{\kappa(1-m)+n_j-s_j-1}$,
co daje pierwszy wzór i średnią. Dla różnych obserwacji
tej samej grupy, przy warunkowej niezależności,
$\E[X_{ji}X_{jk}]=\E[p_j^2]$ oraz
$\E X_{ji}=\E p_j=m$.
Wobec tego kowariancja wynosi
$\Var(p_j)=m(1-m)/(\kappa+1)$.
Ponieważ marginalnie $X_{ji}$ jest Bernoulliego
z parametrem $m$, suma $n_j$ wariancji i
$n_j(n_j-1)$ kowariancji daje ostatni wzór.
\end{proof}

\begin{przyklad}[Dwie linie produkcyjne]
Przyjmijmy $m=0{,}10$, $\kappa=20$. Jedna linia nie wykazała
wad w dwóch próbach, druga wykazała siedem wad w stu.
Oddzielne estymatory MLE to odpowiednio $0$ i $0{,}07$.
Posteriorowe średnie są równe
$(0+2)/(2+20)=0{,}0909\ldots$ oraz
$(7+2)/(100+20)=0{,}075$.
Pierwsza linia nie jest uznawana za całkowicie bezawaryjną
tylko na podstawie dwóch obserwacji. Przy niewłaściwie
ustalonym $m$ ten efekt może jednak prowadzić do błędnego
ściągania, więc parametry populacyjne należy poddawać
diagnostyce i analizie wrażliwości.
\end{przyklad}

\section{Zaskakujący rezultat: ściąganie może pomagać bez prioru}
Poniższe twierdzenie nie zakłada, że nieznany wektor średnich
został wylosowany z rozkładu. Stanowi częstościowy argument,
dlaczego estymowanie grup w pełnej izolacji może być nieefektywne.

\begin{twierdzenie}[Fenomen Jamesa--Steina]
\label{thm:james-stein}
Niech $d\ge3$, $\sigma^2>0$ i
$X\sim\Normal_d(\theta,\sigma^2 I_d)$, gdzie
$\theta\in\R^d$ jest ustalonym, nieznanym parametrem,
a $\Normal_d$ oznacza $d$-wymiarowy rozkład normalny.
Przy stracie $L(\theta,a)=\norm{a-\theta}^2$ estymator
$\delta_{\rm JS}(x)=(1-(d-2)\sigma^2/\norm{x}^2)x$
dla $x\ne0$ oraz $\delta_{\rm JS}(0)=0$ ma ryzyko
\[
 R(\theta,\delta_{\rm JS})
 =d\sigma^2-(d-2)^2\sigma^4
   \E_\theta\!\left[\frac1{\norm{X}^2}\right]
 <d\sigma^2=R(\theta,X)
\]
dla \emph{każdego} $\theta\in\R^d$.
Wartość oczekiwana we wzorze jest skończona i dodatnia.
\end{twierdzenie}
\begin{proof}
Niech $a=(d-2)\sigma^2$ i
$g(x)=-ax/\norm{x}^2$ dla $x\ne0$.
W przestrzeni $\R^d\setminus\{0\}$ zachodzi
$\norm{g(x)}^2=a^2/\norm{x}^2$ oraz
$\operatorname{div}g(x)=-a(d-2)/\norm{x}^2$.
Dla gęstości normalnej $f_\theta$ mamy
$\partial_i f_\theta(x)=-(x_i-\theta_i)f_\theta(x)/\sigma^2$.
Całkowanie przez części po obszarze
$\{\varepsilon<\norm{x}<R\}$, a następnie przejście
$\varepsilon\downarrow0$, $R\uparrow\infty$, daje
\[
 \E_\theta[(X-\theta)\cdot g(X)]
   =\sigma^2\E_\theta[\operatorname{div}g(X)].
\]
Istotnie, wyraz brzegowy na sferze o promieniu
$\varepsilon$ jest rzędu $O(\varepsilon^{d-2})$,
ponieważ $\norm{g(x)}=a/\varepsilon$, pole sfery ma
rząd $\varepsilon^{d-1}$, a $f_\theta$ jest lokalnie
ograniczona. Zewnętrzny wyraz brzegowy zanika dzięki
gaussowskiemu zanikowi gęstości.
Całkowalność $\norm{X}^{-2}$ w otoczeniu zera
wynika z $\int_0^1r^{d-3}\,\dd\lambda_1(r)<\infty$,
co wymaga $d\ge3$; poza kulą jednostkową funkcja
$\norm{x}^{-2}$ jest ograniczona.
Rozwijając błąd kwadratowy $X+g(X)-\theta$ dostajemy
\[
 R(\theta,X+g)=d\sigma^2
   +2\sigma^2\E_\theta[\operatorname{div}g(X)]
   +\E_\theta\norm{g(X)}^2.
\]
Podstawienie wyliczonych pochodnych daje
$d\sigma^2+[a^2-2\sigma^2a(d-2)]\E_\theta[\norm{X}^{-2}]$.
Dla $a=(d-2)\sigma^2$ współczynnik wynosi
$-(d-2)^2\sigma^4<0$. Ponieważ rozkład normalny nie
przypisuje dodatniej masy punktowi zero, definicja w zerze
nie wpływa na ryzyko.
\end{proof}

\begin{uwaga}[Nie należy nadinterpretować twierdzenia]
Wynik nie mówi, że wszystkie współrzędne osobno są
estymowane lepiej, ani że procedura Jamesa--Steina zawsze
zwraca wartość między zerem a średnią z próby (czynnik
ściągania może być ujemny). Stwierdza poprawę łącznego
częstościowego ryzyka kwadratowego dla $d\ge3$.
Łączenie informacji między grupami może więc przynieść
korzyść nawet bez bayesowskiej interpretacji $\theta$.
\end{uwaga}

\section{Kiedy jaką metodę stosować?}
\begin{center}\small
\begin{tabular}{>{\raggedright\arraybackslash}p{.27\textwidth}>{\raggedright\arraybackslash}p{.31\textwidth}>{\raggedright\arraybackslash}p{.27\textwidth}}
\toprule
\textbf{Sytuacja} & \textbf{Metoda do rozważenia} & \textbf{Główne ryzyko}\\
\midrule
Dużo obserwacji w każdej odrębnej grupie & Estymacja oddzielna lub hierarchiczna & Nadmierne ściąganie rzeczywiście różnych grup.\\
Małe próby w wielu podobnych grupach & Model hierarchiczny z partial pooling & Nietrafny rozkład grup i hiperprior.\\
Jedna wspólna intensywność uzasadniona procesem & Model ze wspólnym parametrem & Ukrycie heterogeniczności.\\
Dane binarne w licznych partiach & Beta--dwumianowy lub logit hierarchiczny & Zależności i zmiany w czasie.\\
Szybka analiza wielu grup & Empirical Bayes & Niedoszacowanie niepewności hiperparametrów.\\
Predykcja nowej grupy & Pełny posterior predykcyjny grupy & Niepoprawny podział trening/test.\\
\bottomrule
\end{tabular}
\end{center}

\Needspace{23\baselineskip}
\begin{lstlisting}[language=Python,caption={Kompletny program: partial pooling dla średnich grup i ryzyko Jamesa--Steina (Python, NumPy).}]
import numpy as np

mu, tau2, sigma2 = 10.0, 4.0, 16.0
n = np.array([4, 64], dtype=float)
y = np.array([14.0, 14.0])
v = sigma2 / n
w = tau2 / (tau2 + v)
post_mean = w * y + (1.0 - w) * mu
post_var = tau2 * v / (tau2 + v)
print("group posterior means:", post_mean)
print("group posterior SD:", np.sqrt(post_var))

rng = np.random.default_rng(2026)
d, sigma, repetitions = 5, 2.0, 100000
theta = np.zeros(d)
X = rng.normal(theta, sigma, size=(repetitions, d))
r2 = np.sum(X**2, axis=1)
js = (1.0 - (d - 2) * sigma**2 / r2)[:, None] * X
risk_raw = np.mean(np.sum((X - theta)**2, axis=1))
risk_js = np.mean(np.sum((js - theta)**2, axis=1))
print("empirical risks raw/JS:", risk_raw, risk_js)
print("theoretical risk raw:", d * sigma**2)
\end{lstlisting}

\begin{zadanie}[Rola precyzji i zmienności międzygrupowej]
Przy znanym $\mu=0$ i $v_j=1$ dwie analizy używają
odpowiednio $\tau^2=1/9$ oraz $\tau^2=9$.
Obliczyć $w_j$ i wyjaśnić różnicę.
\par\emph{Rozwiązanie.} Otrzymujemy $w_j=1/10$ i $9/10$.
W pierwszym przypadku podobieństwo grup jest silnym
założeniem priorowym i szacowanie ściąga prawie całkowicie
do zera. W drugim obserwowana grupa może wyraźnie
różnić się od populacji, więc wynik jest bliski średniej
z jej własnych pomiarów. Przed wyborem $\tau^2$ trzeba
sprawdzić zakres możliwej heterogeniczności i liczbę grup.
\end{zadanie}

% ================================================================
% ROZDZIAL 22 -- ASYMPTOTYKA BAYESOWSKA
% ================================================================
\chapter{Asymptotyka bayesowska: koncentracja posterioru i twierdzenie Bernsteina--von Misesa}
\label{ch:bayesian-asymptotics}
\paragraph{Powiązania i odsyłacze.} Warunki klasycznej asymptotyki zebrano w rozdziale~\ref{ch:asymptotic-estimation}; metody posteriorowe i hierarchiczne w rozdziałach~\ref{ch:bayes-foundations} i~\ref{ch:hierarchical-bayes}.


\section{Dwa rodzaje gwarancji i dwa sposoby granicznego rozumowania}
Rozkład a posteriori jest prawdopodobieństwem warunkowym przy
zadanym priorze i modelu. Nie oznacza to automatycznie, że dla
każdej ustalonej wartości parametru jego przedziały wiarygodności
mają deklarowane pokrycie częstościowe. Zbieżność tych dwóch
opisów pojawia się w wielu \emph{regularnych} modelach przy dużej
próbie, ale nie jest uniwersalnym prawem.

W tym rozdziale rozróżnimy: (i) \emph{zgodność} -- koncentrację
posterioru wokół prawdziwego parametru; (ii) \emph{tempo koncentracji}
-- wielkość otoczeń, w których skupia się posterior;
(iii) \emph{kształt} posterioru po odpowiednim przeskalowaniu.
Dopiero trzeci krok prowadzi do twierdzenia Bernsteina--von Misesa.
Dla przejrzystości ogólne twierdzenie sformułujemy przy jawnych
założeniach lokalnej asymptotycznej normalności i koncentracji
w skali $n^{-1/2}$; wykażemy te warunki w modelach normalnym
i Bernoulliego. Nie będziemy przedstawiać lokalnej normalności
jako wystarczającej bez kontroli ogonów posterioru.

\begin{definicja}[Zgodność posterioru]
\label{def:posterior-consistency}
Niech $(\Theta,d)$ będzie przestrzenią metryczną,
$\theta_0\in\Theta$ prawdziwym parametrem, a
$X_1,X_2,\ldots$ niezależnymi obserwacjami o wspólnym
prawie $P_{\theta_0}$. Niech $\Pi_n(\cdot\mid X_1,\ldots,X_n)$
stanowi posterior na $\Theta$.
Posterior jest \emph{zgodny w punkcie $\theta_0$ według
prawdopodobieństwa}, jeśli dla każdych $\varepsilon,\eta>0$
\[
 P_{\theta_0}\!\left(
 \Pi_n(\{\theta:d(\theta,\theta_0)>\varepsilon\}\mid X_1,\ldots,X_n)
 >\eta\right)\longrightarrow0.
\]
O \emph{zgodności prawie na pewno} mówimy, gdy dla każdego
$\varepsilon>0$ posteriorowa masa powyższego zbioru
zbiega do zera $P_{\theta_0}$-prawie na pewno.
\end{definicja}

\begin{uwaga}
Prawdziwy parametr nie musi być losowany z prioru.
Symbol $P_{\theta_0}$ oznacza rzeczywiste prawo generowania
danych przy ustalonym $\theta_0$; posterior jest losową
miarą, ponieważ zależy od próby. W modelu źle określonym
może nie istnieć żaden $\theta_0$ dający rzeczywisty rozkład
obserwacji; taki przypadek wymaga odrębnej analizy.
\end{uwaga}

\section{Pierwsze dowody zgodności: Bernoulli i skończona przestrzeń parametrów}
\begin{twierdzenie}[Zgodność posterioru beta--Bernoulli]
\label{thm:beta-posterior-consistency}
Niech $a,b>0$, $p_0\in[0,1]$, $X_1,X_2,\ldots$
będą niezależne o rozkładzie $\operatorname{Bernoulli}(p_0)$
i niech prior $p$ będzie równy $\operatorname{Beta}(a,b)$.
Wówczas posterior $\operatorname{Beta}(a+S_n,b+n-S_n)$,
$S_n=\sum_{i=1}^nX_i$, jest zgodny prawie na pewno
w punkcie $p_0$ dla metryki $|p-q|$ na $[0,1]$.
Twierdzenie obejmuje także $p_0=0$ i $p_0=1$,
chociaż sam prior nie przypisuje tym punktom dodatniej masy.
\end{twierdzenie}
\begin{proof}
Oznaczmy przez $m_n$ i $v_n$ średnią oraz wariancję
posterioru. Z twierdzenia~\ref{thm:beta-binomial}
\[
 m_n=\frac{a+S_n}{n+a+b},\qquad
 v_n=\frac{(a+S_n)(b+n-S_n)}{(n+a+b)^2(n+a+b+1)}
 \le\frac1{4(n+a+b+1)}.
\]
Z mocnego prawa wielkich liczb $S_n/n\to p_0$
$P_{p_0}$-prawie na pewno, więc także $m_n\to p_0$.
Dla dowolnego $\varepsilon>0$, gdy
$|m_n-p_0|<\varepsilon/2$, nierówność Czebyszewa
zastosowana \emph{do rozkładu posterioru} daje
\[
 \Pi_n(|p-p_0|>\varepsilon\mid X_1,\ldots,X_n)
 \le\Pi_n(|p-m_n|>\varepsilon/2\mid X_1,\ldots,X_n)
 \le\frac{4v_n}{\varepsilon^2}\longrightarrow0.
\]
Warunek $|m_n-p_0|<\varepsilon/2$ zachodzi ostatecznie
na zdarzeniu prawdopodobieństwa jeden, a oszacowanie
wariancji jest deterministyczne. Dowód nie wykorzystuje
warunku $p_0\in(0,1)$.
\end{proof}

\begin{twierdzenie}[Zgodność przy skończonej liczbie hipotez]
\label{thm:finite-bayes-consistency}
Niech $\Theta=\{\theta_1,\ldots,\theta_K\}$,
$K\ge2$, z priorem $\pi(\theta)>0$ dla każdego $\theta$.
Niech $\theta_0\in\Theta$ będzie prawdziwym parametrem,
$X_i\iid P_{\theta_0}$, a każde $P_\theta$ posiada
gęstość $f_\theta$ względem wspólnej miary $\sigma$-skończonej.
Załóżmy, że dla $\theta\ne\theta_0$ iloraz
$f_\theta(X_1)/f_{\theta_0}(X_1)$ jest dodatni
$P_{\theta_0}$-prawie na pewno, jego logarytm jest
całkowalny oraz
\[
 D(\theta_0\Vert\theta)
 :=\E_{\theta_0}\!\left[
  \log\frac{f_{\theta_0}(X_1)}{f_\theta(X_1)}\right]>0.
\]
Wtedy $\Pi_n(\{\theta_0\}\mid X_1,\ldots,X_n)\to1$
$P_{\theta_0}$-prawie na pewno.
\end{twierdzenie}
\begin{proof}
Dla dowolnego $\theta\ne\theta_0$ posteriorowy iloraz
prawdopodobieństw ma postać
\[
 \frac{\Pi_n(\{\theta\}\mid X)}{\Pi_n(\{\theta_0\}\mid X)}
 =\frac{\pi(\theta)}{\pi(\theta_0)}
 \prod_{i=1}^n\frac{f_\theta(X_i)}{f_{\theta_0}(X_i)}.
\]
Logarytm ilorazu podzielony przez $n$ zbiega
prawie na pewno do $-D(\theta_0\Vert\theta)<0$
na mocy mocnego prawa wielkich liczb. Każdy z $K-1$
ilorazów dąży więc do zera. Sumując te ilorazy,
otrzymujemy
$\Pi_n(\{\theta_0\}\mid X)
 =(1+\sum_{\theta\ne\theta_0}\Pi_n(\{\theta\}\mid X)/
 \Pi_n(\{\theta_0\}\mid X))^{-1}\to1$.
\end{proof}

\begin{uwaga}[Czego nie dowodzi ostatnie twierdzenie]
Skończoność przestrzeni parametrów pozwala przejść od
zbieżności każdego ilorazu do zbieżności ich sumy.
Przy nieskończonej przestrzeni wymagana jest dodatkowa
kontrola złożoności modelu i zachowania prioru. Zgodność
w modelu skończonym nie stanowi dowodu ogólnego twierdzenia
o zgodności dowolnych nieparametrycznych posteriorów.
\end{uwaga}

\section{Dokładny model normalny: co dzieje się z priorem?}
\begin{twierdzenie}[Normalny posterior i asymptotyczna normalność]
\label{thm:normal-conjugate-bvm}
Niech $\sigma^2,\tau^2>0$, $m_0\in\R$ i
$X_1,\ldots,X_n\iid\Normal(\theta_0,\sigma^2)$
przy ustalonym $\theta_0\in\R$.
W analizie bayesowskiej przyjmijmy właściwy prior
$\theta\sim\Normal(m_0,\tau^2)$. Posterior jest normalny
z parametrami
\[
 v_n=(\tau^{-2}+n\sigma^{-2})^{-1},\qquad
 m_n=v_n\left(m_0/\tau^2+n\overline X_n/\sigma^2\right).
\]
Przy $n\to\infty$ zachodzi
$\sqrt n(m_n-\overline X_n)\to0$ w prawdopodobieństwie,
$n v_n\to\sigma^2$ i
\[
 \Law\big(\sqrt n(\theta-\overline X_n)\mid X_1,\ldots,X_n\big)
 \longrightarrow\Normal(0,\sigma^2)
\]
w odległości całkowitej wariacji, w prawdopodobieństwie
$P_{\theta_0}$. Ponadto centralny posteriorowy przedział
wiarygodności poziomu $1-\alpha$ ma pokrycie częstościowe
dążące do $1-\alpha$ dla każdego $\alpha\in(0,1)$.
\end{twierdzenie}
\begin{proof}
Pierwsze wzory pochodzą z uzupełnienia kwadratu
w iloczynie normalnej wiarygodności i prioru.
Zachodzi dokładnie
$m_n-\overline X_n=\sigma^2(m_0-\overline X_n)/
(\sigma^2+n\tau^2)$.
Ponieważ $\overline X_n\to\theta_0$ w prawdopodobieństwie,
mnożenie przez $\sqrt n$ daje granicę zero.
Z kolei $nv_n=\sigma^2n\tau^2/(\sigma^2+n\tau^2)\to\sigma^2$.
Rozkład przeskalowanego parametru jest dokładnie
$\Normal(\sqrt n(m_n-\overline X_n),nv_n)$.
Ciągłość gęstości normalnej względem parametrów
w $L^1(\lambda_1)$, którą otrzymujemy z lematu Scheffégo,
daje zbieżność w całkowitej wariacji.
Centralny przedział ma końce
$m_n\pm z_{1-\alpha/2}\sqrt{v_n}$,
gdzie $z_{1-\alpha/2}$ jest kwantylem standardowego
rozkładu normalnego. Z centralnego twierdzenia granicznego,
$\sqrt n(m_n-\theta_0)/\sigma\Rightarrow\Normal(0,1)$
oraz $\sqrt{nv_n}/\sigma\to1$; zatem prawdopodobieństwo
zdarzenia $|m_n-\theta_0|\le z_{1-\alpha/2}\sqrt{v_n}$
dąży do $1-\alpha$.
\end{proof}

\section{Lokalna asymptotyczna normalność: precyzyjne warunki}
Aby oddzielić własność \emph{wiarygodności} od własności
\emph{posterioru}, wprowadzamy formalną lokalną aproksymację
log-wiarygodności. Jest to warunek, a nie uniwersalna
własność każdego modelu statystycznego.

\begin{definicja}[Lokalna asymptotyczna normalność (LAN)]
\label{def:lan}
Niech $\Theta\subseteq\R^d$ będzie zbiorem otwartym,
$\theta_0\in\Theta$, $d\ge1$, a $\ell_n(\theta)$
logarytmem wiarygodności z próby $n$ obserwacji
w eksperymencie o prawie $P_{\theta_0}^{(n)}$.
Mówimy, że model spełnia warunek LAN w $\theta_0$
z macierzą informacji $I_0\in\R^{d\times d}$ dodatnio
określoną, jeśli istnieją wektory losowe $\Delta_n\in\R^d$
takie, że $\Delta_n\Rightarrow\Normal_d(0,I_0)$ pod
$P_{\theta_0}^{(n)}$ oraz dla każdego $M<\infty$
\[
 \sup_{\norm{h}\le M}\left|
 \ell_n(\theta_0+h/\sqrt n)-\ell_n(\theta_0)
 -h^\top\Delta_n+\frac12h^\top I_0h\right|
 \xrightarrow{P_{\theta_0}}0,
\]
przy czym dla dostatecznie dużego $n$ wszystkie
$\theta_0+h/\sqrt n$ z $\norm h\le M$ należą do $\Theta$.
\end{definicja}

\begin{definicja}[Koncentracja posterioru w skali $n^{-1/2}$]
\label{def:root-n-tightness}
Przy $\Theta\subseteq\R^d$ i prawdziwym $\theta_0$
niech $H_n=\sqrt n(\theta-\theta_0)$ będzie zmienną
posteriorową. Posterior jest \emph{ciasny w skali $n^{-1/2}$
w prawdopodobieństwie}, jeśli dla każdego $\eta>0$
\[
 \lim_{M\to\infty}\limsup_{n\to\infty}
 P_{\theta_0}^{(n)}\!\left(
 \Pi_n(\norm{H_n}>M\mid X^{(n)})>\eta\right)=0.
\]
To warunek mocniejszy niż zgodność w każdym ustalonym
otoczeniu $\theta_0$.
\end{definicja}

\begin{twierdzenie}[Bernstein--von Mises przy jawnych warunkach]
\label{thm:bvm-under-lan}
Niech $\Theta\subseteq\R^d$ będzie otwarty,
$\theta_0\in\Theta$, a dane $X^{(n)}$ pochodziły z
modelu zdominowanego z log-wiarygodnością $\ell_n$.
Zakładamy, że:
\begin{enumerate}
\item model spełnia LAN z macierzą $I_0$ w sensie
definicji~\ref{def:lan};
\item prior ma gęstość $\pi$ względem $\lambda_d$,
ciągłą i ściśle dodatnią w punkcie $\theta_0$;
\item posterior jest właściwy oraz spełnia warunek
z definicji~\ref{def:root-n-tightness}.
\end{enumerate}
Niech $\mathcal Q_n$ oznacza posteriorowy rozkład
$H_n=\sqrt n(\theta-\theta_0)$, a $\mathcal G_n$
rozkład $\Normal_d(I_0^{-1}\Delta_n,I_0^{-1})$.
Wtedy
\[
 \norm{\mathcal Q_n-\mathcal G_n}_{\rm TV}
 :=\sup_{A\in\calB(\R^d)}|\mathcal Q_n(A)-\mathcal G_n(A)|
 \xrightarrow{P_{\theta_0}}0.
\]
Jeżeli ponadto estymator $\widehat\theta_n$ spełnia
$\sqrt n(\widehat\theta_n-\theta_0)=I_0^{-1}\Delta_n+o_{P_{\theta_0}}(1)$,
to posteriorowy rozkład
$\sqrt n(\theta-\widehat\theta_n)$ zbiega w całkowitej
wariacji, w prawdopodobieństwie, do $\Normal_d(0,I_0^{-1})$.
\end{twierdzenie}
\begin{proof}
Dla $h$ takiego, że $\theta_0+h/\sqrt n\in\Theta$,
posterior ma względem $\lambda_d$ gęstość proporcjonalną do
\[
 g_n(h)=\exp\big(\ell_n(\theta_0+h/\sqrt n)-\ell_n(\theta_0)\big)
 \pi(\theta_0+h/\sqrt n).
\]
Czynnik $n^{-d/2}$ po zmianie zmiennych jest wspólny
i skraca się przy normalizacji.
Z warunku LAN i ciągłości prioru, na każdej kuli
$K_M=\{h:\norm h\le M\}$ funkcja $g_n(h)$
po podzieleniu przez stałą $\pi(\theta_0)$ jest
jednostajnie, w prawdopodobieństwie, bliska
$g_n^0(h)=\exp(h^\top\Delta_n-h^\top I_0h/2)$.
Aby sformułować to bez niejawnego ograniczenia $\Delta_n$,
ustalamy $C>0$. Na zdarzeniu $\norm{\Delta_n}\le C$
obie funkcje są na $K_M$ ograniczone z góry i z dołu
przez dodatnie stałe zależne od $C,M,I_0$.
Ponieważ $\Delta_n=O_{P_{\theta_0}}(1)$,
można dobrać $C$ tak, aby dopełnienie tego zdarzenia
miało dowolnie małe prawdopodobieństwo.
Stąd unormowane ograniczenia $\mathcal Q_n$ i $\mathcal G_n$
do $K_M$ (po podzieleniu przez ich masy na $K_M$)
zbliżają się do siebie w całkowitej wariacji.
Istotnie, zbieżność jednostajna ich nieunormowanych
gęstości na $K_M$ pociąga zbieżność całek normalizujących
i następnie gęstości unormowanych w $L^1(\lambda_d)$.
Dla dowolnych miar prawdopodobieństwa $Q,G$ oraz zbioru
$K$ o dodatnich masach mamy oszacowanie
\[
 \norm{Q-G}_{\rm TV}
 \le \norm{Q(\cdot\mid K)-G(\cdot\mid K)}_{\rm TV}
       +Q(K^c)+G(K^c).
\]
Wybieramy $M$ tak duże, aby warunek koncentracji posterioru
zapewniał małą masę $\mathcal Q_n(K_M^c)$ z dużym
prawdopodobieństwem. Jednocześnie
$\mathcal G_n(K_M^c)$ jest mała z dużym prawdopodobieństwem,
ponieważ $\Delta_n=O_P(1)$ i macierz kowariancji
$I_0^{-1}$ jest stała. Następnie stosujemy zbieżność
na kuli, a na końcu pozwalamy zwiększać $M$.
Daje to pierwszą tezę.
Druga wynika przez przesunięcie obu miar o
$-\sqrt n(\widehat\theta_n-\theta_0)$ i z ciągłości
w całkowitej wariacji rodziny normalnej względem wektora
średnich, gdy różnica średnich jest $o_P(1)$.
\end{proof}

\begin{uwaga}[Dlaczego założenie o ogonach jest istotne?]
LAN opisuje wiarygodność tylko dla lokalnych zmian parametru
rzędu $n^{-1/2}$. Nie wyklucza odległych obszarów przestrzeni
parametrów, które zachowują znaczną masę posteriorową.
W zastosowaniach do ogólnych modeli warunek ciasności
w skali $n^{-1/2}$ wyprowadza się zwykle z identyfikowalności,
regularności i istnienia odpowiednich testów lub oszacowań
wiarygodności. Tego głębokiego kroku nie zastępujemy
samym rozwinięciem Taylora.
\end{uwaga}

\section{Model Bernoulliego: pełna weryfikacja warunków LAN i koncentracji}
\begin{twierdzenie}[Bernstein--von Mises dla prioru beta we wnętrzu]
\label{thm:bernoulli-bvm}
Niech $p_0\in(0,1)$, $a,b>0$,
$X_1,\ldots,X_n\iid\operatorname{Bernoulli}(p_0)$,
a prior $p$ ma rozkład $\operatorname{Beta}(a,b)$.
Niech $\widehat p_n=S_n/n$. Wówczas posteriorowy rozkład
$\sqrt n(p-\widehat p_n)$ zbiega w całkowitej wariacji,
w prawdopodobieństwie $P_{p_0}$, do
$\Normal(0,p_0(1-p_0))$.
W szczególności dla $\alpha\in(0,1)$ centralny
posteriorowy przedział wiarygodności poziomu $1-\alpha$
ma asymptotyczne pokrycie częstościowe $1-\alpha$.
\end{twierdzenie}
\begin{proof}
Dla $\ell_n(p)=S_n\log p+(n-S_n)\log(1-p)$
na $(0,1)$ wynik różniczkowania w $p_0$ daje
\[
 \Delta_n=\frac{S_n-np_0}{\sqrt n\,p_0(1-p_0)},
 \qquad I_0=\frac1{p_0(1-p_0)}.
\]
Z centralnego twierdzenia granicznego
$\Delta_n\Rightarrow\Normal(0,I_0)$.
Rozwinięcie Taylora $\ell_n(p_0+h/\sqrt n)$ wokół $p_0$
ma drugi współczynnik
$\ell_n''(p)/n=-S_n/(np^2)-(1-S_n/n)/(1-p)^2$.
Na każdym zwartym otoczeniu $p_0$ wyrażenie to zbiega
jednostajnie w prawdopodobieństwie do
$-p_0/p^2-(1-p_0)/(1-p)^2$, a w $p_0$ przyjmuje
wartość $-I_0$. Dla ograniczonych $h$ otrzymujemy LAN.
Gęstość beta jest ciągła i dodatnia w $p_0$.

Pozostaje sprawdzić \emph{tempo} koncentracji.
Niech $m_n=(a+S_n)/(n+a+b)$ i
$v_n=\Var(p\mid X)\le1/[4(n+a+b+1)]$.
Ponieważ $m_n-\widehat p_n=O(1/n)$ jednostajnie
w $S_n\in\{0,\ldots,n\}$, a
$\sqrt n(\widehat p_n-p_0)=O_{P_{p_0}}(1)$,
mamy $\sqrt n(m_n-p_0)=O_{P_{p_0}}(1)$.
Nierówność Czebyszewa pod posteriorem daje dla $M>0$
\[
 \Pi_n(\sqrt n|p-p_0|>M\mid X)
 \le\frac{4nv_n}{M^2}
 \quad\text{na zdarzeniu }\sqrt n|m_n-p_0|\le M/2.
\]
Ponieważ $nv_n\le1/4$, wybierając najpierw $M$ duże,
a potem $n$ duże otrzymujemy ciasność posterioru
w skali $n^{-1/2}$. Możemy zatem zastosować
twierdzenie~\ref{thm:bvm-under-lan}, a z równości
$I_0^{-1}\Delta_n=\sqrt n(\widehat p_n-p_0)$
wynika pierwsza teza.
Zbieżność całkowitej wariacji do normalnego rozkładu
o dodatniej gęstości implikuje zbieżność posteriorowych
kwantyli po standaryzacji. Z CLT dla $\widehat p_n$
wynika wtedy asymptotyczne pokrycie centralnego
przedziału.
\end{proof}

\section{Co się dzieje na brzegu? Ścisły kontrprzykład normalności}
\begin{twierdzenie}[Posterior przy prawdziwym $p_0=0$]
\label{thm:boundary-beta-gamma}
Niech $a,b>0$ i $X_1,\ldots,X_n\iid\operatorname{Bernoulli}(0)$.
W analizie użyjmy prioru $\operatorname{Beta}(a,b)$,
a więc posterioru $p\mid X\sim\operatorname{Beta}(a,b+n)$.
Wtedy
\[
 \Law(np\mid X_1,\ldots,X_n)
 \Rightarrow\operatorname{Gamma}(a,1),
\]
gdzie $\operatorname{Gamma}(a,1)$ oznacza rozkład gamma
z parametrem kształtu $a$ i intensywnością $1$.
Zatem właściwą skalą w tym modelu brzegowym jest $n^{-1}$,
a granica nie jest normalna.
\end{twierdzenie}
\begin{proof}
Przy $p_0=0$ wszystkie obserwacje są zerowe z
prawdopodobieństwem jeden. Gęstość $U_n=np$ względem
$\lambda_1$ jest dla $0<u<n$ równa
\[
 f_n(u)=\frac{u^{a-1}(1-u/n)^{n+b-1}}{n^a B(a,b+n)}.
\]
Z asymptotyki funkcji gamma, wynikającej np. ze wzoru
Stirlinga, $n^a B(a,b+n)\to\Gamma(a)$.
Dla ustalonego $u>0$ mamy
$(1-u/n)^{n+b-1}\to e^{-u}$, stąd
$f_n(u)\to u^{a-1}e^{-u}/\Gamma(a)$.
Granica jest gęstością probabilistyczną na $(0,\infty)$.
Lemat Scheffégo (zbieżność prawie wszędzie nieujemnych
gęstości o całce jeden do gęstości o całce jeden)
daje nawet zbieżność w $L^1(\lambda_1)$, a zatem
zbieżność w całkowitej wariacji.
\end{proof}

\begin{przyklad}[Próba bez wad nie uzasadnia normalnego przedziału]
Przy $a=b=1$, jeśli nie stwierdzono żadnej wady
w $n$ niezależnych próbach, posterior to
$\operatorname{Beta}(1,n+1)$.
Jego średnia wynosi $1/(n+2)$, natomiast asymptotycznie
$n p$ ma rozkład wykładniczy o średniej jeden.
Symetryczny przedział otrzymany przez naiwną aproksymację
normalną może wychodzić poza $[0,1]$ i pomijać
asymetrię posterioru. W tej sytuacji należy stosować
dokładne kwantyle beta albo odpowiedni jednostronny
przedział, a nie automatycznie wzór oparty na CLT we wnętrzu.
\end{przyklad}

\section{Jak interpretować przedziały i jak rozpoznać zawodność asymptotyki?}
\begin{stwierdzenie}[Zbieżność pokrycia w regularnym modelu jednowymiarowym]
\label{prop:credible-frequentist-coverage}
W założeniach twierdzenia~\ref{thm:bvm-under-lan} dla $d=1$
załóżmy dodatkowo, że istnieje estymator
$\widehat\theta_n$ spełniający jego ostatni warunek.
Niech $[L_n,U_n]$ będzie centralnym posteriorowym
przedziałem wiarygodności o masie $1-\alpha$,
$\alpha\in(0,1)$, gdzie kwantyle są brane względem
posterioru $\theta$.
Wtedy $P_{\theta_0}(\theta_0\in[L_n,U_n])\to1-\alpha$.
\end{stwierdzenie}
\begin{proof}
Po standaryzacji $\sqrt n(\theta-\widehat\theta_n)$
posterior zbiega w całkowitej wariacji do
$\Normal(0,I_0^{-1})$ w prawdopodobieństwie.
Ponieważ dystrybuanta tego rozkładu jest ciągła i ściśle
rosnąca, dolny i górny posteriorowy kwantyl mają postać
$\sqrt n(L_n-\widehat\theta_n)
 =-z_{1-\alpha/2}/\sqrt{I_0}+o_P(1)$ i
$\sqrt n(U_n-\widehat\theta_n)
 = z_{1-\alpha/2}/\sqrt{I_0}+o_P(1)$.
Z założenia dotyczącego estymatora i z LAN
$\sqrt n(\widehat\theta_n-\theta_0)
 \Rightarrow\Normal(0,I_0^{-1})$.
Wobec tego zdarzenie pokrycia jest asymptotycznie
równoważne nierówności
$|\sqrt{nI_0}(\widehat\theta_n-\theta_0)|
 \le z_{1-\alpha/2}+o_P(1)$, której prawdopodobieństwo
zbiega do $1-\alpha$.
\end{proof}

\begin{center}\small
\begin{tabular}{>{\raggedright\arraybackslash}p{.26\textwidth}>{\raggedright\arraybackslash}p{.32\textwidth}>{\raggedright\arraybackslash}p{.27\textwidth}}
\toprule
\textbf{Sytuacja} & \textbf{Czy normalność jest uzasadniona?} & \textbf{Lepszy krok praktyczny}\\
\midrule
Duża próba, parametr we wnętrzu regularnego modelu & Często tak, przy spełnieniu LAN i koncentracji & Porównać kwantyle posterioru z aproksymacją normalną.\\
Prawdopodobieństwo bliskie $0$ lub $1$ & Aproksymacja może być zła, na brzegu bywa fałszywa & Dokładny posterior beta, przedział jednostronny.\\
Model mieszaninowy z zamiennymi etykietami & Lokalna LAN nie wyklucza wielomodalnego posterioru & Parametry identyfikowalne lub funkcjonały niezmiennicze.\\
Mało grup, nieznana wariancja międzygrupowa & Hyperposterior może być bardzo szeroki & Pełny Bayes i analiza wrażliwości hiperprioru.\\
Liczba parametrów rośnie z $n$ & Klasyczne twierdzenie o ustalonym $d$ nie wystarcza & Osobne warunki wysokowymiarowe.\\
Źle określony model & Koncentracja może zajść wokół pseudo-prawdy & Diagnostyka i odporna niepewność.\\
\bottomrule
\end{tabular}
\end{center}

\section{Laboratorium: porównanie skończonej próby z granicą}
W programie porównujemy dokładny posterior beta z rozkładem
normalnym wynikającym z teorii asymptotycznej. Dla małych
$n$ i rzadkich zdarzeń rozbieżności bywają widoczne.
Nie wolno na podstawie samej symulacji wnioskować o
spełnieniu warunków twierdzenia w dowolnym modelu.

\Needspace{23\baselineskip}
\begin{lstlisting}[language=Python,caption={Kompletny program: zbieżność posterioru beta do aproksymacji normalnej i przykład brzegowy (Python, SciPy).}]
import numpy as np
from scipy.stats import beta, gamma, norm

rng = np.random.default_rng(2026)
p0, a, b = 0.07, 2.0, 18.0
z = norm.ppf(0.975)
for n in [25, 100, 1000, 10000]:
    s = rng.binomial(n, p0)
    phat = s / n
    post_a, post_b = a + s, b + n - s
    exact_ci = beta.ppf([0.025, 0.975], post_a, post_b)
    # Operational Wald approximation: plug in the observed proportion.
    gaussian_ci = phat + np.array([-1, 1]) * z * np.sqrt(phat*(1-phat)/n)
    print("n", n, "s", s, "exact", exact_ci, "Wald", gaussian_ci)

# At the boundary p0=0, n*p converges to Gamma(a, rate=1).
n = 10000
post_a, post_b = a, b + n
print("n * posterior median:", n * beta.ppf(0.5, post_a, post_b))
print("limiting Gamma median:", gamma.ppf(0.5, a, scale=1.0))
\end{lstlisting}

\Needspace{23\baselineskip}
\begin{lstlisting}[language=Julia,caption={Kompletny program: posterior normalny i asymptotyczna niepewność średniej (Julia, Distributions.jl).}]
using Random, Statistics, Distributions

Random.seed!(2026)
sigma, tau, m0, theta0 = 2.0, 3.0, 0.0, 1.2
for n in (20, 200, 2000)
    x = rand(Normal(theta0, sigma), n)
    xbar = mean(x)
    vpost = inv(inv(tau^2) + n / sigma^2)
    mpost = vpost * (m0 / tau^2 + n * xbar / sigma^2)
    exact = quantile(Normal(mpost, sqrt(vpost)), [0.025, 0.975])
    asymptotic = xbar .+ [-1.0, 1.0] .* (1.9599639845 * sigma / sqrt(n))
    println((n=n, posterior=exact, asymptotic=asymptotic))
end
\end{lstlisting}

\begin{uwaga}[Aproksymacja operacyjna a znany tylko w symulacji parametr]
W listingu powyżej czynnik wariancji w przedziale Walda obliczamy
z $\widehat p_n$, a nie z $p_0$. W rzeczywistej analizie
$p_0$ jest nieznany. Przy $S_n=0$ operacyjny przedział
Walda degeneruje się do $[0,0]$, podczas gdy dokładny
posterior beta pozostaje nietrywialny. Znajomość $p_0$
w kodzie symulacyjnym służy jedynie generowaniu danych
i ocenie błędu, nie konstruowaniu przedziału.
\end{uwaga}

\begin{zadanie}[Asymptotyka przy nieznanej jakości]
W modelu Bernoulliego z priorem $\operatorname{Beta}(1,1)$
zaobserwowano $S_n$ wad na $n$ niezależnych prób.
Zakładając $p_0\in(0,1)$ i $S_n/n\to p_0$,
obliczyć granicę $n\Var(p\mid S_n)$.
\par\emph{Rozwiązanie.} Z wzoru na wariancję posterioru
\[
 n\Var(p\mid S_n)=
 \frac{n(S_n+1)(n-S_n+1)}{(n+2)^2(n+3)}.
\]
Dzieląc licznik i mianownik przez $n^3$ otrzymujemy
granicę $p_0(1-p_0)$.
To zgadza się z odwrotnością jednostkowej informacji
Fishera $I_0^{-1}=p_0(1-p_0)$; samo takie dopasowanie
wariancji nie dowodzi jeszcze zbieżności całego rozkładu
posterioru do normalnego.
\end{zadanie}

\chapter{Metody Monte Carlo: symulacja, błąd i redukcja wariancji}
\label{ch:monte-carlo}
\paragraph{Powiązania i odsyłacze.} Prawo wielkich liczb i CTG pochodzą z tomu I (rozdziały~\ref{I-ch:lln} i~\ref{I-ch:clt}); symulacje procesów dyfuzyjnych omawia tom III, rozdział~\ref{III-ch:sde-properties}.


\section{Dlaczego losować, skoro mamy całkę?}
W statystyce często chcemy obliczyć całkę względem rozkładu a posteriori,
wartość predykcyjną albo ryzyko decyzji, a gęstość tego rozkładu jest
znana tylko z dokładnością do stałej normalizującej. Nawet gdy wartość
całki jest jednoznacznie zdefiniowana, obliczenie jej metodami
analitycznymi może być bardzo trudne. Metody Monte Carlo zastępują
całkowanie uśrednianiem odpowiednio dobranych zmiennych losowych.
W odróżnieniu od liczby danych $n$ z eksperymentu statystycznego,
liczbę \emph{symulacji} będziemy oznaczać przez $N$. Wzrost $N$ zmniejsza
błąd \emph{numeryczny}, lecz sam w sobie nie zwiększa informacji
statystycznej zawartej w pierwotnych danych.

\begin{definicja}[Podstawowy estymator Monte Carlo]
\label{def:mc-basic}
Niech $(\mathsf X,\mathcal A,P)$ będzie przestrzenią probabilistyczną,
$f\colon\mathsf X\to\R$ funkcją mierzalną i $f\in L^1(P)$.
Dla niezależnych $X_1,\ldots,X_N\sim P$ definiujemy
$\mu=\E_P[f(X_1)]=\int_{\mathsf X}f(x)\,\dd P(x)$ oraz
$\widehat\mu_N=N^{-1}\sum_{i=1}^N f(X_i)$.
Jeżeli $f\in L^2(P)$, oznaczamy
$\sigma_f^2=\Var_P(f(X_1))$.
\end{definicja}

\begin{twierdzenie}[Nieobciążoność, zgodność i centralne twierdzenie graniczne]
\label{thm:mc-lln-clt}
W warunkach definicji~\ref{def:mc-basic}, dla każdego $N\ge1$
$\E_P\widehat\mu_N=\mu$ oraz $\widehat\mu_N\to\mu$
prawie na pewno. Jeżeli ponadto $0<\sigma_f^2<\infty$, to
\[
 \Var_P(\widehat\mu_N)=\frac{\sigma_f^2}{N},\qquad
 \frac{\sqrt N(\widehat\mu_N-\mu)}{\sigma_f}
 \Rightarrow\Normal(0,1).
\]
Jeżeli $\sigma_f^2=0$, wówczas $\widehat\mu_N=\mu$ prawie na pewno.
\end{twierdzenie}
\begin{proof}
Liniowość wartości oczekiwanej daje nieobciążoność, a mocne prawo
wielkich liczb dla niezależnych zmiennych o jednakowym rozkładzie
-- zgodność. W przypadku $f\in L^2(P)$ niezależność zeruje
kowariancje między składnikami średniej, więc wariancja wynosi
$N^{-2}N\sigma_f^2=\sigma_f^2/N$.
Ostatnia zbieżność jest klasycznym centralnym twierdzeniem granicznym
z tomu I, zastosowanym do $f(X_i)$. Gdy wariancja jest zerowa,
$f(X_i)=\mu$ prawie na pewno dla każdego $i$; przeliczalność
rodziny pozwala uzyskać wspólne zdarzenie prawdopodobieństwa jeden.
\end{proof}

\begin{wniosek}[Przybliżony błąd Monte Carlo i liczba losowań]
\label{cor:mc-se}
W założeniach twierdzenia~\ref{thm:mc-lln-clt}, przy $N\ge2$, niech
$S_N^2=(N-1)^{-1}\sum_{i=1}^N(f(X_i)-\widehat\mu_N)^2$.
Wtedy $S_N^2\to\sigma_f^2$ prawie na pewno i
$\widehat{\mathrm{SE}}_{\rm MC}=S_N/\sqrt N$ jest zgodnym
oszacowaniem błędu standardowego Monte Carlo. Dla dużego $N$
przedział $\widehat\mu_N\pm z_{1-\alpha/2}S_N/\sqrt N$
ma asymptotyczne pokrycie $1-\alpha$ dla ustalonej całki $\mu$.
Aby zmniejszyć błąd standardowy dwukrotnie, przy niezmienionej
wariancji całkowanej funkcji trzeba około \emph{czterokrotnie}
zwiększyć $N$.
\end{wniosek}
\begin{proof}
Z mocnego prawa wielkich liczb dla $f(X_i)^2$ oraz dla $f(X_i)$
wynika $S_N^2\to\sigma_f^2$ prawie na pewno.
Twierdzenie~\ref{thm:mc-lln-clt} i twierdzenie Slutsky'ego
dają normalną granicę statystyki po podstawieniu $S_N$.
Zależność od $N^{-1/2}$ wyjaśnia ostatnie stwierdzenie.
\end{proof}

\begin{uwaga}[Błąd Monte Carlo nie jest błędem estymacji parametru]
Jeśli $\pi(\theta\mid x)$ jest posteriorem po danych $x$, to
całka $\mu(x)=\int_\Theta f(\theta)\pi(\theta\mid x)\,
\dd\lambda_d(\theta)$ jest \emph{ustaloną funkcją danych}.
Błąd przybliżenia $\widehat\mu_N(x)$ jest warunkowym błędem
numerycznym. Szerokość posterioru opisuje natomiast niepewność
$\theta$ przy danym $x$. Długi przebieg symulacji może bardzo dobrze
obliczyć szeroki posterior, lecz go nie zawęzi.
\end{uwaga}

\begin{przyklad}[Rzadkie zdarzenie -- ograniczenia naiwnego Monte Carlo]
\label{ex:mc-rare}
Jeśli $A$ ma prawdopodobieństwo $p$, bierzemy
$f=\ind{A}$; wtedy $\widehat p_N$ jest średnią z prób Bernoulliego,
a $\Var(\widehat p_N)=p(1-p)/N$. Względny błąd standardowy wynosi
$\sqrt{(1-p)/(Np)}$. Jeżeli $p=10^{-6}$ i $N=10^4$,
to oczekiwana liczba trafień to $0{,}01$, a prawdopodobieństwo
\emph{ani jednego trafienia} wynosi $(1-p)^N\approx0{,}99005$.
Otrzymane zero nie dowodzi, że zdarzenie jest niemożliwe.
\end{przyklad}

\section{Dwa sposoby zmniejszania kosztu: warunkowanie i zmienna kontrolna}
Idea warunkowania nawiązuje do twierdzenia Rao--Blackwella z rozdziału~\ref{ch:sufficiency}; tym razem służy redukcji błędu numerycznego, a nie ocenie ryzyka estymatora parametru.
\begin{twierdzenie}[Rao--Blackwell dla symulacji]
\label{thm:mc-rao-blackwell}
Niech $Z\in L^2(P)$ i $\mathcal G\subseteq\mathcal A$ będzie
pod-$\sigma$-ciałem. Wtedy
$\E[\E(Z\mid\mathcal G)]=\E Z$ oraz
$\Var(\E(Z\mid\mathcal G))\le\Var(Z)$,
z równością wtedy i tylko wtedy, gdy
$Z=\E(Z\mid\mathcal G)$ prawie na pewno.
Dlatego, jeżeli można obliczyć warunkową wartość oczekiwaną
bez kosztownej symulacji pozostałej losowości, warto symulować
$\E(Z\mid\mathcal G)$ zamiast $Z$.
\end{twierdzenie}
\begin{proof}
Pierwsza teza wynika z własności wieży, druga z rozkładu
$\Var(Z)=\E[\Var(Z\mid\mathcal G)]
 +\Var(\E[Z\mid\mathcal G])$.
Pierwszy składnik jest nieujemny i znika dokładnie wtedy,
gdy $Z=\E(Z\mid\mathcal G)$ prawie na pewno.
\end{proof}

\begin{twierdzenie}[Zmienna kontrolna]
\label{thm:control-variate}
Niech $Y,C\in L^2(P)$, $\E C=c$ jest dokładnie znane,
a $\Var(C)>0$. Dla ustalonej liczby $\beta\in\R$
zdefiniujmy $Z_\beta=Y-\beta(C-c)$.
Wtedy $\E Z_\beta=\E Y$ i
\[
 \Var(Z_\beta)=\Var(Y)-2\beta\Cov(Y,C)
                      +\beta^2\Var(C).
\]
Wariancję minimalizuje
$\beta_*=\Cov(Y,C)/\Var(C)$, a jej minimum wynosi
$\Var(Y)(1-\rho_{Y,C}^2)$, gdy $\Var(Y)>0$ i
$\rho_{Y,C}$ oznacza współczynnik korelacji.
\end{twierdzenie}
\begin{proof}
Odjęcie zmiennej o średniej zero nie zmienia wartości oczekiwanej.
Rozwinięcie wariancji kombinacji liniowej daje wielomian kwadratowy.
Po uzupełnieniu kwadratu
$\Var(Z_\beta)=\Var(C)(\beta-\beta_*)^2+
\Var(Y)-\Cov(Y,C)^2/\Var(C)$.
Wzór na korelację daje postać końcową.
\end{proof}

\begin{uwaga}[Pułapka doboru współczynnika na tych samych symulacjach]
Twierdzenie~\ref{thm:control-variate} dotyczy \emph{ustalonego}
$\beta$. Jeżeli wybieramy $\widehat\beta$ na podstawie tych
samych symulacji, które uśredniamy, nie wolno bez dowodu
przypisywać otrzymanej średniej dokładnej nieobciążoności.
Prosta metoda zachowania nieobciążoności polega na oszacowaniu
$\widehat\beta$ z niezależnej próby pilotażowej, a następnie
użyciu go jako stałej w nowej próbie.
\end{uwaga}

\begin{przyklad}[Średnia zmiennej lognormalnej]
Niech $Z\sim\Normal(0,1)$ i
$Y=e^Z$, $C=Z$. Wiemy, że $\E C=0$,
$\E Y=e^{1/2}$, $\Cov(Y,Z)=\E[Ze^Z]=e^{1/2}$
oraz $\Var(e^Z)=e^2-e$.
Najlepszy współczynnik kontrolny to $\beta_*=e^{1/2}$,
a wariancja maleje z $e^2-e$ do $e^2-2e$.
W przykładzie traktujemy $\E Y$ jako szukaną całkę,
mimo że znamy jej dokładną wartość i możemy sprawdzić błąd symulacji.
\end{przyklad}

\begin{stwierdzenie}[Zmienne antytetyczne]
\label{prop:antithetic}
Niech $U\sim\operatorname{Unif}(0,1)$, a
$f\in L^2((0,1),\lambda_1)$ będzie monotoniczna.
Wtedy $\Cov(f(U),f(1-U))\le0$.
Zatem estymator $[f(U)+f(1-U)]/2$ ma wariancję
nie większą niż $\Var(f(U))/2$.
Dla $N$ niezależnych par dzielimy tę wariancję przez $N$.
\end{stwierdzenie}
\begin{proof}
Niech $V$ będzie niezależną kopią $U$. Dla monotonicznej
$f$ funkcja $u\mapsto f(u)$ jest monotoniczna, a
$u\mapsto f(1-u)$ ma monotoniczność przeciwną.
Wobec tego
$[f(U)-f(V)][f(1-U)-f(1-V)]\le0$ prawie na pewno.
Po rozwinięciu i wykorzystaniu niezależności $U,V$
wartość oczekiwana lewej strony jest dwukrotnością
$\Cov(f(U),f(1-U))$, co daje tezę.
Wariancja średniej dwóch zmiennych o tej samej wariancji
wynosi $(\Var(f(U))+\Cov(f(U),f(1-U)))/2$.
\end{proof}

\section{Importance sampling: zmiana rozkładu losowania}
\begin{definicja}[Próbkowanie ważone]
\label{def:importance-sampling}
Niech $p,q$ będą gęstościami probabilistycznymi względem
tej samej miary $\sigma$-skończonej $\nu$ na $(\mathsf X,\mathcal A)$,
a $f\colon\mathsf X\to\R$ będzie mierzalna z
$\int_{\mathsf X}|f(x)|p(x)\,\dd\nu(x)<\infty$.
Zakładamy, że $q(x)>0$ dla $\nu$-prawie wszystkich $x$,
dla których $|f(x)|p(x)>0$. Definiujemy wagę
$w(x)=p(x)/q(x)$ dla $q(x)>0$ oraz $w(x)=0$ dla $q(x)=0$.
Dla niezależnych $X_i\sim q\,\nu$ estymatorem \emph{importance sampling}
jest $\widehat\mu_{N,\mathrm{IS}}=N^{-1}\sum_{i=1}^N f(X_i)w(X_i)$.
\end{definicja}

\begin{twierdzenie}[Nieobciążoność i wariancja importance sampling]
\label{thm:is-variance}
Przy założeniach definicji~\ref{def:importance-sampling} zachodzi
$\E_q\widehat\mu_{N,\mathrm{IS}}=\mu=
\int_{\mathsf X}f(x)p(x)\,\dd\nu(x)$.
Jeżeli ponadto
$\int_{\mathsf X}f(x)^2p(x)^2/q(x)\,\dd\nu(x)<\infty$
(z ilorazem rozumianym jako zero w punktach, w których $f(x)p(x)=0$), to
\[
 \Var_q(\widehat\mu_{N,\mathrm{IS}})=\frac1N
 \left(\int_{\mathsf X}\frac{f(x)^2p(x)^2}{q(x)}\,
       \dd\nu(x)-\mu^2\right).
\]
Spośród wszystkich gęstości $q$ dodatnich na $\{|f|p>0\}$
minimum drugiego momentu osiąga $q_*(x)=|f(x)|p(x)/A$,
gdzie $A=\int_{\mathsf X}|f(x)|p(x)\,\dd\nu(x)>0$.
Dla nieujemnego $f$ wariancja przy $q_*$ wynosi zero,
\emph{ale} konstrukcja $q_*$ wymaga znajomości całki $A=\mu$
i możliwości losowania z tego rozkładu.
\end{twierdzenie}
\begin{proof}
Z założenia nośnikowego $f(x)p(x)=0$ na $\{q=0\}$
prawie wszędzie. Zatem
$\E_q[fw]=\int_{\{q>0\}} f(x)p(x)\,\dd\nu(x)=\mu$.
Rachunek drugiego momentu i niezależność $X_i$ dają wariancję.
Dla dowolnej dopuszczalnej gęstości $q$, z nierówności
Cauchy'ego--Schwarza,
\[
 A^2=\left(\int_{\mathsf X}
       \frac{|f|p}{\sqrt q}\sqrt q\,\dd\nu(x)\right)^2
 \le\left(\int_{\mathsf X}\frac{f^2p^2}{q}\,\dd\nu(x)\right)
       \int_{\mathsf X}q(x)\,\dd\nu(x)
 =\int_{\mathsf X}\frac{f^2p^2}{q}\,\dd\nu(x).
\]
Równość zachodzi, gdy $|f|p/q$ jest stałe na nośniku $q$;
stąd $q_*=|f|p/A$. Gdy $f\ge0$, ważona wartość
$f(X)w(X)$ przy tym $q_*$ jest prawie na pewno równa $A$.
\end{proof}

\begin{uwaga}[Kiedy optymalny rozkład propozycji jest praktycznie nieosiągalny]
Gęstość $q_*$ z twierdzenia~\ref{thm:is-variance} jest optymalna dla
\emph{konkretnego} funkcjonału $f$, a nie jednocześnie dla wszystkich
całek względem $p$. Dla $f$ zmieniającej znak minimalna wariancja
wynosi $A^2-\mu^2$, która jest dodatnia, chyba że $f$ ma stały znak
$p$-prawie wszędzie. W praktycznych algorytmach przybliżamy $q_*$
rozwiązaniem parametrycznym, a potem oceniamy drugi moment wag:
duże ESS nie zastępuje weryfikacji pokrycia nośnika.
\end{uwaga}

\begin{uwaga}[Zero w mianowniku i nieistniejąca wariancja]
Może zachodzić $\E_q|fw|<\infty$, ale
$\E_q[(fw)^2]=\infty$. Wtedy estymator jest zgodny
na mocy mocnego prawa, lecz zwykły błąd standardowy
oparty na CLT nie ma uzasadnienia. W praktyce należy
kontrolować ogony wag i to, czy rozkład propozycji
$q$ jest dostatecznie szeroki na obszarach ważnych dla całki.
\end{uwaga}

\begin{twierdzenie}[Samonormalizowane importance sampling]
\label{thm:sn-is}
Niech $p$ będzie gęstością probabilistyczną,
$q>0$ na $\{p>0\}$ prawie wszędzie,
$f\in L^1(p\,\nu)$ oraz $X_i\iid q\,\nu$.
Zdefiniujmy $w_i=p(X_i)/q(X_i)$ oraz
$\widetilde\mu_N=\sum_{i=1}^Nw_i f(X_i)/\sum_{i=1}^Nw_i$
na zdarzeniu dodatniego mianownika (poza nim przyjmujemy
umownie wartość zero). Wtedy
$\widetilde\mu_N\to\mu$ prawie na pewno.
Jeżeli $\E_q[w^2(f-\mu)^2]<\infty$ i
$0<\Var_q(w(f-\mu))$, to
\[
 \sqrt N(\widetilde\mu_N-\mu)
 \Rightarrow\Normal\left(0,\E_q[w^2(f-\mu)^2]\right).
\]
Ten estymator nie jest na ogół dokładnie nieobciążony.
\end{twierdzenie}
\begin{proof}
Z warunku nośnikowego $\E_q w=1$ i
$\E_q(wf)=\mu$. MLLN dla obu średnich daje zgodność i
w szczególności dodatniość mianownika dla dużych $N$
prawie na pewno. Ponadto
\[
 \widetilde\mu_N-\mu=
 \frac{N^{-1}\sum_{i=1}^N w_i(f(X_i)-\mu)}
      {N^{-1}\sum_{i=1}^N w_i}.
\]
Licznik jest średnią niezależnych, wycentrowanych
zmiennych o skończonej wariancji wskazanej w tezie.
Z CLT i twierdzenia Slutsky'ego uzyskujemy normalną granicę.
Iloraz losowych średnich nie musi mieć wartości oczekiwanej
równej ilorazowi ich wartości oczekiwanych; już dla $N=1$
i różnych $p,q$ samonormalizowana średnia równa się
$f(X_1)$ przy $w_1>0$, więc jej wartość oczekiwana
może być inna niż $\mu$.
\end{proof}

\begin{uwaga}[Efektywna liczebność próby ważonej]
W praktyce oblicza się
$\mathrm{ESS}_{\rm w}=(\sum_iw_i)^2/\sum_iw_i^2\in[1,N]$
na zdarzeniu, że suma wag jest dodatnia.
Jest to \emph{wskaźnik koncentracji wag}, a nie dokładna
gwarancja błędu $\widetilde\mu_N$; przy wagach zerowych
liczymy tylko dodatnie składniki. Mała wartość oznacza,
że kilka losowań dominuje wynik. Sama duża wartość ESS
nie dowodzi, że rozkład propozycji pokrywa wszystkie
istotne obszary ani że istniała wariancja wag.
\end{uwaga}

\begin{przyklad}[Ekspozycyjna zmiana miary dla rzadkiego ogona]
Niech $Z\sim\Normal(0,1)$, $a>0$, a
$p=\Prob(Z>a)=\int_{(a,\infty)}\varphi(z)\,\dd\lambda_1(z)$,
gdzie $\varphi$ oznacza standardową gęstość normalną.
W naiwnej symulacji liczymy $\ind{\{Z>a\}}$.
W importance sampling losujemy $Y\sim\Normal(a,1)$ i liczymy
\[
 W(Y)=\ind{\{Y>a\}}\exp(-aY+a^2/2).
\]
Ponieważ $\varphi(y)/\varphi(y-a)=\exp(-ay+a^2/2)$,
$\E W(Y)=p$. Wybór przesunięcia $a$ sprawia, że
przekroczenie progu występuje pod propozycją
z prawdopodobieństwem $1/2$, zamiast bardzo rzadko.
Dobre ustawienie propozycji trzeba jednak ocenić przez
wariancję wag, a nie samą częstość przekroczeń.
\end{przyklad}

\section{Co kiedy stosować? Porównanie w jednym miejscu}
\begin{center}\small
\begin{tabular}{>{\raggedright\arraybackslash}p{.25\textwidth}>{\raggedright\arraybackslash}p{.30\textwidth}>{\raggedright\arraybackslash}p{.29\textwidth}}
\toprule
\textbf{Zadanie} & \textbf{Metoda do rozważenia} & \textbf{Kontrola jakości}\\
\midrule
Znamy rozkład, umiemy losować niezależnie & Zwykłe Monte Carlo & Błąd $S_N/\sqrt N$ przy skończonej wariancji\\
Rzadkie zdarzenie lub dominujące ogony & Importance sampling & Nośnik propozycji, ogony wag i ESS\\
Znamy średnią skorelowanej zmiennej & Zmienna kontrolna & Zysk wariancyjny; osobna próba pilotażowa\\
Umiemy scałkować część losowości & Warunkowanie Rao--Blackwella & Koszt wyliczenia warunkowej średniej\\
Monotoniczny funkcjonał jednego kwantyla & Losowania antytetyczne & Faktyczna kowariancja w parach\\
Posterior znany tylko do normalizacji & MCMC z rozdziału~\ref{ch:mcmc} & Diagnostyka wielu łańcuchów, autokorelacja\\
\bottomrule
\end{tabular}
\end{center}

\section{Laboratorium: błąd i rzadkie zdarzenia}
\begin{lstlisting}[language=Python,caption={Kompletny program: Monte Carlo, zmienna kontrolna i importance sampling dla ogona normalnego (Python, NumPy i SciPy).}]
import numpy as np
from scipy.stats import norm

rng = np.random.default_rng(2026)
n = 100_000
z = rng.standard_normal(n)
y = np.exp(z)
naive = y.mean()
control = (y - np.exp(0.5) * z).mean()
print("E[exp(Z)] exact:", np.exp(0.5))
print("naive estimate:", naive, "SE:", y.std(ddof=1) / np.sqrt(n))
print("control estimate:", control)

a = 4.0
x = rng.standard_normal(n)
rare_plain = (x > a).astype(float)
y_is = rng.normal(loc=a, scale=1.0, size=n)
w_is = (y_is > a) * np.exp(-a * y_is + a*a/2)
print("tail exact:", norm.sf(a))
print("plain:", rare_plain.mean(), "hits:", int(rare_plain.sum()))
print("importance sampling:", w_is.mean())
print("IS estimated SE:", w_is.std(ddof=1) / np.sqrt(n))

# Self-normalized importance sampling has generally nonzero bias.
weights = np.exp(-a * y_is + a*a/2)
snis = np.sum(weights * (y_is > a)) / weights.sum()
ess = weights.sum()**2 / np.sum(weights**2)
print("SNIS estimate:", snis, "weight ESS:", ess)
\end{lstlisting}

\begin{lstlisting}[language=Julia,caption={Kompletny program: redukcja wariancji przez zmienną kontrolną i estymacja rzadkiego ogona (Julia, Random i Statistics).}]
using Random, Statistics
Random.seed!(2026)
N = 100_000
z = randn(N)
y = exp.(z)
beta = exp(0.5)
println("exact E[exp(Z)]: ", beta)
println("plain: ", mean(y))
println("control: ", mean(y .- beta .* z))

a = 4.0
x = randn(N)
plain = mean(x .> a)
yis = randn(N) .+ a
wis = (yis .> a) .* exp.(-a .* yis .+ a^2 / 2)
println("plain tail estimate: ", plain)
println("importance sampling tail estimate: ", mean(wis))
println("IS estimated standard error: ", std(wis) / sqrt(N))
\end{lstlisting}

\begin{zadanie}[Dwie metody dla tej samej całki]
Niech $U\sim\operatorname{Unif}(0,1)$ i $f(u)=u^2$.
Porównać wariancje pojedynczej wartości $f(U)$ i
średniej antytetycznej $[f(U)+f(1-U)]/2$.
\par\emph{Rozwiązanie.} Mamy $\E U^2=1/3$,
$\E U^4=1/5$, więc $\Var(U^2)=4/45$.
Ponieważ $\E[U^2(1-U)^2]=1/30$,
$\Cov(U^2,(1-U)^2)=1/30-1/9=-7/90$.
Wariancja średniej pary jest
$\tfrac12(4/45-7/90)=1/180$.
Przy dwóch niezależnych próbach wariancja ich średniej
wynosiłaby $2/45$. Para antytetyczna jest więc przy
tej samej liczbie dwóch ewaluacji znacznie dokładniejsza.
\end{zadanie}


\chapter{Metody MCMC: łańcuchy Markowa, estymacja i diagnostyka}
\label{ch:mcmc}
\paragraph{Powiązania i odsyłacze.} Podstawy łańcuchów Markowa omawiają tom I, rozdział~\ref{I-ch:markov}, i tom II, rozdział~\ref{II-ch:ctmc}; metodę Monte Carlo — rozdział~\ref{ch:monte-carlo}.


\section{Gdy nie umiemy losować z rozkładu a posteriori}
Niech dane $x$ wyznaczają posterior o gęstości
$\pi(\theta\mid x)=\widetilde\pi(\theta\mid x)/Z(x)$,
gdzie $\widetilde\pi\ge0$ można obliczać punktowo, lecz
$Z(x)=\int_\Theta\widetilde\pi(\theta\mid x)\,
\dd\lambda_d(\theta)$ nie ma wygodnego wzoru.
Zakładamy $0<Z(x)<\infty$, ponieważ niewłaściwy posterior
nie jest rozkładem, z którego można poprawnie losować.
Metody \emph{Markov chain Monte Carlo} (MCMC) konstruują
łańcuch o rozkładzie niezmienniczym $\pi$ i estymują
$\mu=\int_\Theta f(\theta)\pi(\theta\mid x)\,
\dd\lambda_d(\theta)$ przez średnią wzdłuż trajektorii.
Nie należy zakładać, że kolejne stany łańcucha są niezależne.
Związek długookresowej wariancji ze zjawiskiem autokorelacji omówiliśmy wcześniej w rozdziale~\ref{ch:time-series}.
Wszystkie te pojęcia odnoszą się do ustalonych danych $x$;
parametr $N$ poniżej oznacza długość łańcucha, nie liczbę
danych wykorzystanych do obliczenia posterioru.

\begin{definicja}[Jądro Markowa i miara niezmiennicza]
\label{def:mcmc-kernel}
Na przestrzeni mierzalnej $(\Theta,\mathcal A)$ \emph{jądrem Markowa}
nazywamy odwzorowanie $P\colon\Theta\times\mathcal A\to[0,1]$
takie, że $P(\theta,\cdot)$ jest prawdopodobieństwem dla
każdego $\theta$ oraz $P(\cdot,A)$ jest mierzalna dla każdego $A$.
Miara probabilistyczna $\pi$ jest \emph{niezmiennicza}, jeśli
$\int_\Theta P(\theta,A)\,\dd\pi(\theta)=\pi(A)$ dla każdego $A$.
Jądro spełnia \emph{warunek szczegółowej równowagi} względem $\pi$,
gdy miary na $\Theta\times\Theta$ zadane przez
$\pi(\dd\theta)P(\theta,\dd\vartheta)$ i
$\pi(\dd\vartheta)P(\vartheta,\dd\theta)$ są równe.
Taki łańcuch nazywamy \emph{odwracalnym} względem $\pi$.
\end{definicja}

\begin{stwierdzenie}[Szczegółowa równowaga implikuje niezmienniczość]
\label{prop:balance-invariant}
Jeżeli jądro $P$ spełnia szczegółową równowagę względem
prawdopodobieństwa $\pi$, to $\pi P=\pi$.
\end{stwierdzenie}
\begin{proof}
Dla każdego $A\in\mathcal A$ całkujemy równość miar
szczegółowej równowagi po $\Theta\times A$.
Lewa strona jest $\int_\Theta P(\theta,A)\,\dd\pi(\theta)$,
a prawa $\int_A P(\vartheta,\Theta)\,\dd\pi(\vartheta)=\pi(A)$.
\end{proof}

\begin{uwaga}[Niezmienniczość nie jest jeszcze zbieżnością]
Łańcuch $X_{k+1}=X_k$ ma \emph{każdy} rozkład jako
niezmienniczy, ale nie bada rozkładu docelowego.
Aby średnie MCMC były użyteczne, potrzebujemy dodatkowo
odpowiednich własności ergodycznych: przede wszystkim
osiągalności regionów o dodatniej masie docelowej,
a w typowych twierdzeniach również nieokresowości.
\end{uwaga}

\section{Metropolis--Hastings: dowód poprawności algorytmu}
\begin{twierdzenie}[Jądro Metropolisa--Hastingsa]
\label{thm:mh-detailed}
Niech $\Theta\subseteq\R^d$ będzie mierzalna,
$\pi$ gęstością prawdopodobieństwa względem $\lambda_d$,
a $q(y\mid x)$ mierzalną gęstością propozycji względem
$\lambda_d$ dla każdego $x\in\Theta$.
Dla $\pi(x)q(y\mid x)>0$ zdefiniujmy
\[
 \alpha(x,y)=\min\left\{1,
 \frac{\pi(y)q(x\mid y)}{\pi(x)q(y\mid x)}\right\};
\]
na pozostałych parach z $q(y\mid x)>0$ można ustalić
$\alpha(x,y)$ w dowolny mierzalny sposób zgodny z
zachowaniem na nośniku $\pi$; w szczególności dla
$\pi(x)>0$ i $\pi(y)=0$ przyjmujemy $\alpha(x,y)=0$.
Algorytm: z punktu $x$ losujemy $Y\sim q(\cdot\mid x)$,
a następnie przechodzimy do $Y$ z prawdopodobieństwem
$\alpha(x,Y)$, w przeciwnym razie pozostajemy w $x$.
Tak otrzymane jądro $P$ ma $\pi$ jako rozkład niezmienniczy.
Do obliczania $\alpha$ wystarczy nieunormowana
gęstość $\widetilde\pi=Z\pi$, gdy $0<Z<\infty$.
\end{twierdzenie}
\begin{proof}
Dla $x\ne y$ gęstość przepływu zaakceptowanych przejść
względem $\lambda_{2d}$ jest równa
\[
 \pi(x)q(y\mid x)\alpha(x,y)
 =\min\{\pi(x)q(y\mid x),\pi(y)q(x\mid y)\},
\]
co jest symetryczne w $x,y$. Pozostała masa przejść
znajduje się na przekątnej $\{(x,x):x\in\Theta\}$,
a ta miara także pozostaje niezmieniona po zamianie
współrzędnych. Dla stanów spoza nośnika $\pi$
wybór jądra nie wpływa na powyższe równości miar.
Otrzymaliśmy szczegółową równowagę i możemy zastosować
stwierdzenie~\ref{prop:balance-invariant}.
Stała $Z$ skraca się w ilorazie $\pi(y)/\pi(x)$.
\end{proof}

\begin{przyklad}[Dwie propozycje i ich konsekwencje]
Jeżeli $q(y\mid x)=q(x\mid y)$, to
$\alpha(x,y)=\min\{1,\widetilde\pi(y)/\widetilde\pi(x)\}$.
Tak jest np. dla propozycji $Y=X+\varepsilon$,
$\varepsilon\sim\Normal_d(0,s^2I_d)$; metoda nazywa się
\emph{random-walk Metropolis}. Dla małego $s$ akceptacja
może być wysoka, lecz łańcuch porusza się powoli. Dla
bardzo dużego $s$ częste odrzucenia mogą powodować
powtarzanie tych samych stanów.
Jeżeli natomiast $q(y\mid x)=g(y)$ nie zależy od $x$,
to $\alpha(x,y)=\min\{1,\pi(y)g(x)/(\pi(x)g(y))\}$.
Nie wolno pominąć ilorazu propozycji, chyba że
propozycja jest rzeczywiście symetryczna.
\end{przyklad}

\section{Dlaczego średnie MCMC zbieżne są w szczególnych warunkach?}
\begin{twierdzenie}[Jednoznaczność i zbieżność z warunku minoracji]
\label{thm:mcmc-minorization}
Niech $P$ będzie jądrem Markowa na mierzalnej przestrzeni
$(\Theta,\mathcal A)$. Załóżmy, że istnieją
$m\ge1$, $\varepsilon\in(0,1]$ i prawdopodobieństwo $\nu$
takie, że dla każdego $x\in\Theta$ i $A\in\mathcal A$
$P^m(x,A)\ge\varepsilon\nu(A)$.
Wtedy $P$ ma dokładnie jedną miarę niezmienniczą $\pi$
i dla każdej początkowej miary $\mu$,
\[
 \norm{\mu P^{km}-\pi}_{\rm TV}
 \le(1-\varepsilon)^k\norm{\mu-\pi}_{\rm TV},
 \qquad k\ge0.
\]
W szczególności całkowita wariacja rozkładu łańcucha
dąży do zera również dla ogólnych $N\to\infty$.
\end{twierdzenie}
\begin{proof}
Jeżeli $0<\varepsilon<1$, zapisujemy
$P^m(x,\cdot)=\varepsilon\nu(\cdot)
 +(1-\varepsilon)R(x,\cdot)$, gdzie $R$ jest jądrem
Markowa. Dla dwóch miar probabilistycznych $\mu_1,\mu_2$
mamy
$\mu_1P^m-\mu_2P^m=(1-\varepsilon)(\mu_1R-\mu_2R)$.
Ponieważ jądro Markowa nie zwiększa odległości
w całkowitej wariacji, otrzymujemy kontrakcję o czynnik
$1-\varepsilon$. Przestrzeń miar probabilistycznych
z odległością całkowitej wariacji jest zupełna,
więc twierdzenie Banacha o punkcie stałym daje
jedyny stały punkt $\pi$ operatora $\mu\mapsto\mu P^m$.
Ponieważ $\pi P$ także jest jego punktem stałym,
$\pi P=\pi$. Każda miara niezmiennicza dla $P$
jest też niezmiennicza dla $P^m$, więc musi być $\pi$.
Iterowanie kontrakcji daje oszacowanie.
Dla $N=km+r$, $0\le r<m$, korzystamy z
$\pi P^r=\pi$ i nieekspansywności jądra $P^r$.
Przy $\varepsilon=1$ mamy $P^m(x,\cdot)=\nu$
dla wszystkich $x$ i tezy są natychmiastowe.
\end{proof}

\begin{uwaga}[Warunek jest mocny, ale przejrzysty]
Dla wielu ciągłych przestrzeni stanów, zwłaszcza
nieograniczonych, globalna minoracja nie zachodzi.
Używa się wtedy warunków Harrisa, funkcji Lyapunova
i małych zbiorów. Nie zakładamy, że każda propozycja
Metropolisa--Hastingsa daje zbieżność geometryczną.
Co więcej, samo trafienie w pobliże jednego maksimum
posterioru nie gwarantuje odwiedzenia innych modów.
\end{uwaga}

\begin{twierdzenie}[Prawo wielkich liczb dla skończonego łańcucha]
\label{thm:mcmc-finite-slln}
Niech $\Theta=\{1,\ldots,K\}$ i $P$ będzie macierzą
przejścia nieprzywiedlnego łańcucha Markowa na $\Theta$.
Niech $\pi$ oznacza jego jedyny rozkład stacjonarny.
Dla dowolnej funkcji $f\colon\Theta\to\R$ oraz
\emph{dowolnego} rozkładu początkowego,
$N^{-1}\sum_{k=0}^{N-1}f(X_k)\to\sum_{x=1}^Kf(x)\pi(x)$
prawie na pewno. Nieokresowość nie jest potrzebna
do tego prawa wielkich liczb.
\end{twierdzenie}
\begin{proof}
Niech $g=f-\pi(f)$, gdzie
$\pi(f)=\sum_xf(x)\pi(x)$. Dla skończonego,
nieprzywiedlnego łańcucha równanie Poissona
$(I-P)h=g$ ma rozwiązanie $h$, ponieważ przestrzeń
wektorów o średniej $\pi$ równej zero jest obrazem
$I-P$. Istotnie, wektor w jądrze $I-P$ jest stały:
zasada maksimum i nieprzywiedlność wymuszają
równość współrzędnych, a stąd rząd $I-P$ wynosi $K-1$.
Zatem $g(X_k)=h(X_k)-Ph(X_k)$ i
\[
 \sum_{k=0}^{N-1}g(X_k)
 =h(X_0)-h(X_N)+\sum_{k=1}^{N}
   \big(h(X_k)-Ph(X_{k-1})\big).
\]
Suma po prawej stronie jest martyngałem z przyrostami
o zerowej warunkowej średniej, a wszystkie przyrosty
są ograniczone stałą $2\norm h_\infty$.
Z nierówności Azumy--Hoeffdinga dla takich martyngałów
wynika, że dla każdego $\delta>0$ prawdopodobieństwo
przekroczenia przez wartość bezwzględną martyngału
progu $\delta N$ jest ograniczone przez
$2\exp[-c_\delta N]$ dla pewnego $c_\delta>0$.
Pierwszy lemat Borela--Cantellego daje zatem
zbieżność martyngału podzielonego przez $N$ do zera
prawie na pewno. Wyraz brzegowy jest ograniczony
przez $2\norm h_\infty/N\to0$. Sumując po
przeliczalnym zbiorze dodatnich wymiernych $\delta$
otrzymujemy tezę.
\end{proof}

\begin{lemat}[CLT dla ograniczonych różnic martyngałowych]
\label{lem:bounded-martingale-clt}
Niech $(D_k,\mathcal F_k)_{k\ge1}$ będzie ciągiem różnic
martyngałowych na jednej przestrzeni probabilistycznej,
$\E[D_k\mid\mathcal F_{k-1}]=0$, oraz niech
$|D_k|\le C$ prawie na pewno dla pewnego $C<\infty$
i wszystkich $k$. Oznaczmy
$V_k=\E[D_k^2\mid\mathcal F_{k-1}]$.
Jeżeli $N^{-1}\sum_{k=1}^NV_k\to v$ według
prawdopodobieństwa dla pewnej stałej $v\ge0$, to
$N^{-1/2}\sum_{k=1}^ND_k\Rightarrow\Normal(0,v)$.
Dla $v=0$ przez $\Normal(0,0)$ rozumiemy masę Diraca w zerze.
\end{lemat}
\begin{proof}
Ustalmy $t\in\R$ i napiszmy
$a_{N,k}=\E[\exp(itD_k/\sqrt N)\mid\mathcal F_{k-1}]$.
Z rozwinięcia wykładnika do drugiego rzędu, ograniczoności
$D_k$ i zerowej warunkowej średniej wynika jednostajnie
w $k\le N$ i niemal wszystkich realizacjach
\[
 a_{N,k}=1-\frac{t^2V_k}{2N}+O(N^{-3/2}),
 \quad
 \log a_{N,k}=-\frac{t^2V_k}{2N}+O(N^{-3/2}).
\]
Dla dużego $N$ wszystkie $a_{N,k}$ są niezerowe i ich
moduły są co najmniej $1-C_t/N$ dla pewnej stałej $C_t$.
Dla $S_k=\sum_{j=1}^kD_j$ proces
\[
 M_{N,k}=\frac{\exp(itS_k/\sqrt N)}
 {\prod_{j=1}^k a_{N,j}},\qquad 0\le k\le N,
\]
jest martyngałem zespolonym względem $\mathcal F_k$,
gdyż mianownik po dodaniu czynnika $a_{N,k}$ jest
$\mathcal F_{k-1}$-mierzalny. Mamy $\E M_{N,N}=1$,
a ponadto $|M_{N,N}|\le(1-C_t/N)^{-N}\le e^{2C_t}$
dla dostatecznie dużych $N$.
Zsumowanie rozwinięć logarytmu daje
$\prod_{k=1}^N a_{N,k}=
\exp[-t^2(2N)^{-1}\sum_{k=1}^NV_k+O(N^{-1/2})]$.
Ponieważ $N^{-1}\sum V_k\to v$ według prawdopodobieństwa,
ten iloczyn dąży do $\exp(-t^2v/2)$ według
prawdopodobieństwa i jest jednostajnie ograniczony.
Korzystając z ograniczoności $M_{N,N}$ otrzymujemy
\[
 \E\exp(itS_N/\sqrt N)
 =\E\left[M_{N,N}\prod_{k=1}^Na_{N,k}\right]
 \longrightarrow e^{-t^2v/2}\E M_{N,N}
 =e^{-t^2v/2}.
\]
Twierdzenie Lévy'ego o ciągłości funkcji charakterystycznych
kończy dowód.
\end{proof}

\section{Autokorelacja, długookresowa wariancja i ESS}
\begin{stwierdzenie}[Wariancja średniej ze stacjonarnego łańcucha]
\label{prop:mcmc-autocorrelation}
Niech $(X_k)_{k\ge0}$ będzie ściśle stacjonarnym
łańcuchem Markowa o rozkładzie $\pi$ i niech
$Y_k=f(X_k)$ ma skończoną, dodatnią wariancję
$\gamma_0=\Var_\pi(Y_0)$. Oznaczmy
$\gamma_h=\Cov_\pi(Y_0,Y_h)$ oraz
$\rho_h=\gamma_h/\gamma_0$. Wtedy dla każdego $N\ge1$
\[
 \Var\left(\frac1N\sum_{k=0}^{N-1}Y_k\right)
 =\frac{\gamma_0}{N}\left(1+2\sum_{h=1}^{N-1}
 (1-h/N)\rho_h\right).
\]
Jeżeli $\sum_{h\ge1}|\gamma_h|<\infty$, to
$N\Var(N^{-1}\sum_{k=0}^{N-1}Y_k)\to
\sigma_\infty^2=\gamma_0+2\sum_{h\ge1}\gamma_h\ge0$.
Gdy $\sigma_\infty^2>0$, definiujemy
$\tau_{\rm int}=\sigma_\infty^2/\gamma_0$ i
$N_{\rm eff}=N/\tau_{\rm int}$ jako asymptotyczną
\emph{efektywną liczebność próby dla funkcji $f$}.
\end{stwierdzenie}
\begin{proof}
Wariancja sumy wynosi
$\sum_{i,j=0}^{N-1}\Cov(Y_i,Y_j)$.
Ze stacjonarności, dla każdego opóźnienia $h\ge1$
występuje $2(N-h)$ par o kowariancji $\gamma_h$
oraz $N$ par z kowariancją $\gamma_0$.
Podzielenie przez $N^2$ daje pierwszy wzór.
Przy bezwzględnej sumowalności można zastosować
twierdzenie o zbieżności zdominowanej dla szeregu,
ponieważ $|1-h/N|\le1$ dla $h<N$.
Nieujemność granicy wynika z nieujemności wariancji.
\end{proof}

\begin{twierdzenie}[Skończony odwracalny łańcuch: widmo i CLT]
\label{thm:mcmc-reversible-clt}
Niech $P$ będzie nieprzywiedlną i nieokresową
macierzą przejścia na skończonej przestrzeni $\Theta$,
odwracalną względem stacjonarnego $\pi$.
Dla $f\colon\Theta\to\R$ o $\pi$-średniej $\mu$ niech
$g=f-\mu$. Niech $(\lambda_j,e_j)_{j=2}^K$
stanowi układ ortonormalnych wektorów własnych
operatora $P$ w $L^2(\pi)$ prostopadłych do stałych,
tak że $g=\sum_{j=2}^Kc_je_j$.
Wtedy $-1<\lambda_j<1$ i
\[
 \sigma_\infty^2
 =\sum_{j=2}^K c_j^2\frac{1+\lambda_j}{1-\lambda_j}.
\]
Jeżeli $g\not\equiv0$, to $\sigma_\infty^2>0$.
Przy \emph{dowolnym} rozkładzie początkowym zachodzi
\[
 \sqrt N\left(\frac1N\sum_{k=0}^{N-1}f(X_k)-\mu\right)
 \Rightarrow\Normal(0,\sigma_\infty^2).
\]
\end{twierdzenie}
\begin{proof}
Szczegółowa równowaga oznacza, że $P$ jest samosprzężony
w $L^2(\pi)$. Jest operatorem kontrakcji, więc jego
wartości własne leżą w $[-1,1]$. Nieprzywiedlność i
nieokresowość w skończonym stanie dają prostą wartość
własną $1$ i wykluczają $-1$.
Dla startu stacjonarnego
$\gamma_h=\langle g,P^hg\rangle_\pi=
\sum_{j=2}^Kc_j^2\lambda_j^h$.
Sumując zbieżny szereg geometryczny otrzymujemy
$\gamma_0+2\sum_{h\ge1}\gamma_h
=\sum_jc_j^2(1+\lambda_j)/(1-\lambda_j)$.
Aby uzasadnić CLT, rozwiązujemy $(I-P)h=g$ tak jak
w dowodzie twierdzenia~\ref{thm:mcmc-finite-slln}.
Suma $g(X_k)$ różni się o ograniczony wyraz brzegowy
od sumy różnic martyngałowych
$D_k=h(X_k)-Ph(X_{k-1})$. Ich warunkowe wariancje
wynoszą $V(X_{k-1})=Ph^2(X_{k-1})-(Ph(X_{k-1}))^2$.
Z udowodnionego prawa wielkich liczb
$N^{-1}\sum_{k=1}^NV(X_{k-1})\to\pi(V)$
prawie na pewno. Przyrosty $D_k$ są ograniczone,
więc spełniają warunek Lindeberga. Stosujemy lemat~\ref{lem:bounded-martingale-clt},
którego założenia zostały tu sprawdzone.
Otrzymujemy normalną granicę z wariancją $\pi(V)$. Ta wariancja
równa się obliczonej wcześniej $\sigma_\infty^2$,
gdyż oba wzory są granicą $N^{-1}\Var_\pi(
\sum_{k=0}^{N-1}g(X_k))$; dla ograniczonych przyrostów
różnice od wyrazu brzegowego są $o(N)$ w $L^2$.
Wyjściowy rozkład nie wpływa na CLT: po skończonej
liczbie kroków jego wpływ na znormalizowaną sumę
zanika, co można formalnie uzasadnić sprzężeniem
ze stacjonarną kopią łańcucha (zbieżność geometryczna
w skończonej nieprzywiedlnej nieokresowej przestrzeni).
\end{proof}

\begin{uwaga}[Czego ESS nie gwarantuje]
Dla silnie dodatniej autokorelacji $\tau_{\rm int}\gg1$,
więc $N_{\rm eff}\ll N$. Przy ujemnych korelacjach
może być też $\tau_{\rm int}<1$, a $N_{\rm eff}>N$;
nie jest to liczba faktycznie niezależnych obserwacji.
Efektywna liczebność zależy od funkcji $f$ i od
stacjonarnego rozkładu. Oszacowanie ESS z krótkiego,
źle mieszającego się łańcucha może być bardzo mylące.
\end{uwaga}

\section{Gibbs sampling i modele hierarchiczne}
\begin{twierdzenie}[Kroki Gibbsa zachowują posterior]
\label{thm:gibbs-invariant}
Niech $\pi(x,y)$ będzie wspólną, właściwą gęstością
na $\R^r\times\R^s$ względem $\lambda_{r+s}$.
Załóżmy, że warunkowe gęstości $\pi(x\mid y)$ oraz
$\pi(y\mid x)$ są określone dla prawie wszystkich
wartości odpowiednich zmiennych. Jądro aktualizujące
$x$ przez losowanie z $\pi(\cdot\mid y)$ przy
ustalonym $y$ oraz jądro aktualizujące $y$ przez
losowanie z $\pi(\cdot\mid x)$ przy ustalonym $x$
mają $\pi$ jako miarę niezmienniczą. Ich złożenie,
\emph{systematyczny skan Gibbsa}, również zachowuje $\pi$.
Nie musi jednak samo być odwracalne względem $\pi$.
\end{twierdzenie}
\begin{proof}
Dla zbiorów mierzalnych $A\subset\R^r$, $B\subset\R^s$
po aktualizacji pierwszej współrzędnej otrzymujemy
\[
 \int_{\R^s}\pi_Y(y)
 \left(\int_A\pi(x\mid y)\,\dd\lambda_r(x)\right)
 \ind{B}(y)\,\dd\lambda_s(y)
 =\int_{A\times B}\pi(x,y)\,\dd\lambda_{r+s}(x,y).
\]
Jest to dokładnie $\pi(A\times B)$. Zbiory prostokątne
generują borelowskie $\sigma$-ciało, więc pierwsze
jądro zachowuje $\pi$; identycznie rozumujemy dla drugiego.
Z równości $\pi P_x=\pi$ i $\pi P_y=\pi$
wynika $\pi(P_xP_y)=\pi$.
\end{proof}

\begin{przyklad}[Gibbs dla skorelowanej pary normalnej]
Niech $(X,Y)$ ma normalny rozkład dwuwymiarowy o średniej
$(0,0)$, wariancjach $1$ i korelacji $\rho\in(-1,1)$.
Wtedy
$X\mid Y=y\sim\Normal(\rho y,1-\rho^2)$ oraz
$Y\mid X=x\sim\Normal(\rho x,1-\rho^2)$.
Algorytm naprzemiennie losuje z tych dwóch rozkładów.
Gdy $|\rho|$ jest bliskie jedności, kolejne iteracje
są mocno powiązane, mimo że \emph{każdy krok warunkowy}
ma dokładny rozkład. Może być wtedy lepiej wykonać
aktualizację blokową lub zmienić parametryzację.
\end{przyklad}

\section{Hamiltonian Monte Carlo: dlaczego gradient pomaga?}
\label{sec:hmc-intro}
Niech posterior ma dodatnią, różniczkowalną gęstość
$\pi(q)\propto\exp[-U(q)]$ na $\R^d$, a $M$
będzie dodatnio określoną macierzą masy. Wprowadzamy
pomocniczy pęd $p\sim\Normal_d(0,M)$ oraz
\emph{hamiltonian} $H(q,p)=U(q)+\frac12p^\top M^{-1}p$.
Wspólny rozkład docelowy na przestrzeni fazowej
jest proporcjonalny do $\exp[-H(q,p)]$.
Ideą jest wykonywanie długich, skorelowanych ruchów
wzdłuż przybliżonych trajektorii hamiltonowskich,
zamiast małych, losowych kroków bez informacji o gradiencie.

\begin{stwierdzenie}[Zachowanie rozkładu przez krok HMC]
\label{prop:hmc-step}
Załóżmy, że $U\in C^1(\R^d)$, $M$ jest dodatnio
określona i $\int_{\R^d}e^{-U(q)}\,\dd\lambda_d(q)<\infty$.
Niech $\Phi$ będzie złożeniem skończonej liczby
kroków algorytmu leapfrog o stałym kroku $\epsilon>0$:
\[
 \begin{aligned}
 p_{1/2}&=p-\tfrac\epsilon2\nabla U(q),\\
 q'&=q+\epsilon M^{-1}p_{1/2},\\
 p'&=p_{1/2}-\tfrac\epsilon2\nabla U(q').
 \end{aligned}
\]
Po złożeniu z odwróceniem pędu $R(q,p)=(q,-p)$
propozycja $T=R\circ\Phi$ jest odwracalną inwolucją
zachowującą miarę Lebesgue'a na $\R^{2d}$.
Akceptowanie jej z prawdopodobieństwem
$\min\{1,\exp[H(q,p)-H(T(q,p))]\}$ zachowuje
wspólny rozkład o gęstości proporcjonalnej do $e^{-H}$.
Odświeżenie pędu z $\Normal_d(0,M)$ również zachowuje
ten rozkład. Łańcuch położeń $q$ ma więc $\pi$
jako rozkład niezmienniczy.
\end{stwierdzenie}
\begin{proof}
Każda podoperacja leapfrog jest ścinaniem:
$(q,p)\mapsto(q,p-\frac\epsilon2\nabla U(q))$
lub $(q,p)\mapsto(q+\epsilon M^{-1}p,p)$.
Dla ustalonej drugiej współrzędnej pierwsze ścinanie
przesuwa $p$, a drugie dla ustalonego $p$ przesuwa $q$;
obie transformacje zachowują miarę Lebesgue'a
(bez potrzeby zakładania istnienia drugich pochodnych).
Każde ma odwrotność przez zmianę $\epsilon$ na $-\epsilon$.
Symetria kompozycji daje $R\Phi R=\Phi^{-1}$,
stąd $(R\Phi)^2=I$. Również $R$ zachowuje objętość.
Dla odwracalnej, zachowującej objętość inwolucji $T$
przepływ masy między $z$ i $T(z)$ po odrzuceniach
jest symetryczny:
$e^{-H(z)}\min(1,e^{H(z)-H(Tz)})
=\min(e^{-H(z)},e^{-H(Tz)})$.
Jest to szczegółowa równowaga na przestrzeni fazowej.
Odświeżenie pędu zachowuje produkt rozkładu $q$ i
normalnego rozkładu $p$, zatem złożenie zachowuje
wspólny rozkład. Po marginalizacji otrzymujemy $\pi$.
\end{proof}

\begin{uwaga}[Co wymaga kontroli przy HMC]
Samo zachowanie rozkładu nie dowodzi ergodyczności.
Błędny gradient, niestabilne kroki, zbyt duże $\epsilon$,
wielomodalność lub parametryzacja z wąskimi obszarami
mogą prowadzić do problemów. Należy kontrolować
akceptację, rozbieżne trajektorie numeryczne, efektywne
liczebności próby i zgodność wielu łańcuchów.
Właściwa implementacja NUTS wymaga osobnego uzasadnienia
reguły zatrzymywania trajektorii; nie utożsamiamy jej
z opisanym tu algorytmem o stałej liczbie kroków.
\end{uwaga}

\section{Diagnostyka: co można sprawdzić, a czego nie można udowodnić wykresem?}
\begin{definicja}[Dwie proste diagnostyki wielu łańcuchów]
Załóżmy, że mamy $m\ge2$ łańcuchów, każdy długości
$N\ge2$, wartości $z_{ji}$ badanego funkcjonału,
$\overline z_j=N^{-1}\sum_{i=1}^Nz_{ji}$ oraz
$\overline z=m^{-1}\sum_{j=1}^m\overline z_j$.
Definiujemy
\[
 W=\frac1m\sum_{j=1}^m
 \frac1{N-1}\sum_{i=1}^N(z_{ji}-\overline z_j)^2,
 \qquad
 B=\frac N{m-1}\sum_{j=1}^m(\overline z_j-\overline z)^2.
\]
Jeżeli $W>0$, to klasyczna diagnostyka Gelmana--Rubina
ma postać
\[
 \widehat R=\sqrt{\frac{(N-1)W+B}{NW}}.
\]
W praktyce warto stosować podział łańcuchów na połowy
(\emph{split-$\widehat R$}) oraz wersje rangowe
i diagnostykę ogonów rozkładu.
\end{definicja}

\begin{uwaga}[Wygrzewanie i przerzedzanie łańcucha]
Odrzucenie początkowych iteracji (\emph{warm-up}, także \emph{burn-in})
może zmniejszać wpływ startu, ale arbitralna liczba odrzuconych kroków
nie gwarantuje osiągnięcia stacjonarności. Przerzedzanie (\emph{thinning}),
czyli zapisywanie co $k$-tej iteracji, na ogół nie poprawia
precyzji estymacji przy ustalonym koszcie wygenerowania pełnego łańcucha;
może natomiast ograniczać zużycie pamięci. W raporcie należy podać
liczbę iteracji przed i po warm-up, liczbę łańcuchów i sposób diagnostyki.
\end{uwaga}

\begin{przyklad}[Pozorna zgodność łańcuchów i nieodwiedzony mod]
Niech docelowy rozkład na $\R$ będzie mieszaniną
$\tfrac12\Normal(-a,1)+\tfrac12\Normal(a,1)$ dla dużego $a$.
Jeżeli wszystkie łańcuchy rozpoczynamy w pobliżu $-a$ i używamy
małych kroków random-walk Metropolisa, mogą przez cały dostępny
czas odwiedzać wyłącznie lewy mod. Ich średnie i wariancje
międzyłańcuchowe będą podobne, zatem $\widehat R$ może być
bliska jedności, mimo że estymata $\E_\pi X=0$ będzie bliska $-a$.
To nie jest kontrprzykład do poprawności algorytmu Metropolisa,
lecz do \emph{wnioskowania o zbieżności jedynie z diagnostyki}.
Należy stosować różne punkty startowe, w tym odległe mody,
i sprawdzać stabilność oszacowań przy zmianach parametryzacji
oraz propozycji.
\end{przyklad}

\begin{uwaga}[Nie istnieje magiczny próg zapewniający zbieżność]
Wartość $\widehat R$ bliska jeden nie dowodzi,
że łańcuchy odwiedziły wszystkie istotne mody. Może
się zdarzyć, że wszystkie łańcuchy utknęły w tym samym
obszarze. Same wykresy śladu i histogramy również nie
stanowią dowodu osiągnięcia rozkładu stacjonarnego.
Bezpieczna procedura obejmuje różne punkty startowe,
kontrolę $\widehat R$, ESS dla średnich i ogonów,
błąd Monte Carlo, analizę posterior predictive oraz
badanie wrażliwości na parametry algorytmu.
\end{uwaga}

\begin{center}\small
\begin{tabular}{>{\raggedright\arraybackslash}p{.25\textwidth}>{\raggedright\arraybackslash}p{.30\textwidth}>{\raggedright\arraybackslash}p{.29\textwidth}}
\toprule
\textbf{Sytuacja} & \textbf{Algorytm} & \textbf{Na co uważać}\\
\midrule
Dokładne warunkowe rozkłady łatwe do losowania & Gibbs & Zależność współrzędnych, wolne mieszanie\\
Posterior jednowymiarowy, gęstość do stałej & Random-walk MH & Skala propozycji, autokorelacja\\
Gładki posterior w większym wymiarze & HMC/NUTS & Gradient, rozbieżności, parametryzacja\\
Trudny multimodalny posterior & Metody tempering / specjalne propozycje & Pozorne zbieganie wszystkich łańcuchów\\
Proste losowanie niezależne dostępne & Zwykłe MC, nie MCMC & Po co wprowadzać zależność?\\
\bottomrule
\end{tabular}
\end{center}

\begin{center}
\begin{tikzpicture}[>=Latex, every node/.style={font=\small},node distance=16mm]
\node[draw,rounded corners,fill=blue!5,minimum width=22mm] (a1) {Losowanie 1};
\node[draw,rounded corners,fill=blue!5,minimum width=22mm,right=of a1] (a2) {Losowanie 2};
\node[draw,rounded corners,fill=blue!5,minimum width=22mm,right=of a2] (a3) {Losowanie 3};
\node[below=12mm of a1,draw,rounded corners,fill=green!6] (b1) {$X_0$};
\node[below=12mm of a2,draw,rounded corners,fill=green!6] (b2) {$X_1$};
\node[below=12mm of a3,draw,rounded corners,fill=green!6] (b3) {$X_2$};
\draw[->,thick] (b1)--node[above] {$P$} (b2);
\draw[->,thick] (b2)--node[above] {$P$} (b3);
\node[below=7mm of b2,font=\footnotesize,text width=100mm,align=center]
 {Górny wiersz: niezależne losowania z rozkładu docelowego.\\
 Dolny wiersz: łańcuch MCMC, zależność przez jądro $P$.};
\end{tikzpicture}

{\small\emph{Rysunek 24.1.} Zwykłe Monte Carlo i MCMC wymagają odmiennych wzorów na błąd standardowy.}
\end{center}

\section{Laboratorium: posterior beta, cztery łańcuchy i diagnostyka}
\label{sec:mcmc-laboratorium}
Na danych o $7$ wadach ze $100$ prób przy priorze
$\operatorname{Beta}(1,1)$ posterior ma rozkład
$\operatorname{Beta}(8,94)$. Ponieważ umiemy z niego
losować \emph{bezpośrednio}, w rzeczywistym zastosowaniu
wybralibyśmy samodzielne losowania Monte Carlo.
Celowo użyjemy jednak tego przykładu do sprawdzenia
implementacji Metropolisa--Hastingsa: wiemy, jaka
jest dokładna średnia posterioru, i możemy porównać
wynik łańcuchów z wartością $8/102$.
Współrzędna $z=\log[p/(1-p)]$ obejmuje całą $\R$.
Ważne: po transformacji gęstość względem $\dd\lambda_1(z)$
zawiera \emph{jakobian} $p(1-p)$.
Stąd dla $p=(1+e^{-z})^{-1}$ log-gęstość docelowa
współrzędnej $z$ jest (z dokładnością do stałej)
$a\log p+b\log(1-p)$, a nie
$(a-1)\log p+(b-1)\log(1-p)$.

\begin{lstlisting}[language=Python,caption={Kompletny program: cztery łańcuchy Metropolisa--Hastingsa dla posterioru beta, split-R-hat i przybliżony ESS (Python, NumPy i SciPy).}]
import numpy as np
from scipy.special import expit

A, B = 8.0, 94.0
N, BURN, CHAINS = 24_000, 4_000, 4
STEP = 0.65
rng = np.random.default_rng(2026)

def log_target(z):
    # Posterior in logit coordinate: beta density times Jacobian.
    log_p = -np.logaddexp(0.0, -z)
    log_1mp = -np.logaddexp(0.0, z)
    return A * log_p + B * log_1mp

def run_chain(z0):
    z, accepted = float(z0), 0
    samples = np.empty(N)
    current = log_target(z)
    for k in range(N):
        proposed = z + rng.normal(0.0, STEP)
        newlog = log_target(proposed)
        if np.log(rng.random()) < min(0.0, newlog - current):
            z, current = proposed, newlog
            accepted += 1
        samples[k] = expit(z)
    return samples[BURN:], accepted / N

starts = np.array([-5.0, -3.0, -1.0, 1.0])
results = [run_chain(z0) for z0 in starts]
chains = np.stack([r[0] for r in results])
print("acceptances:", [round(r[1], 3) for r in results])
print("posterior mean estimate:", chains.mean())
print("exact posterior mean:", A / (A + B))

# Split R-hat (original moment-based version).
L = chains.shape[1]
half = L // 2
split = np.concatenate([chains[:, :half], chains[:, -half:]], axis=0)
m = split.shape[0]
n = split.shape[1]
means = split.mean(axis=1)
W = split.var(axis=1, ddof=1).mean()
B_between = n * means.var(ddof=1)
rhat = np.sqrt(((n - 1) * W / n + B_between / n) / W)
print("split R-hat:", rhat)

# Descriptive autocorrelation-based ESS (not a convergence proof).
def acf_fft(x):
    x = np.asarray(x) - np.mean(x)
    n = len(x)
    cov = np.fft.irfft(np.abs(np.fft.rfft(x, n=2*n))**2,
                       n=2*n)[:n]
    return cov / cov[0]

acf = np.array([acf_fft(c) for c in chains])
rho = acf.mean(axis=0)
tau = 1.0
# Truncate at the first nonpositive adjacent-pair sum.
for k in range(1, min(1000, len(rho) - 1), 2):
    pair = rho[k] + rho[k + 1]
    if pair <= 0:
        break
    tau += 2 * pair
print("approximate ESS:", chains.size / tau)
print("approximate MC SE:", np.sqrt(chains.var(ddof=1)
                                     * tau / chains.size))
\end{lstlisting}

\begin{lstlisting}[language=Julia,caption={Kompletny program: random-walk Metropolis dla posterioru beta we współrzędnej logit (Julia, Random i Statistics).}]
using Random, Statistics
Random.seed!(2026)
a, b = 8.0, 94.0
N, burn, step = 30_000, 5_000, 0.65

sigmoid(z) = z >= 0 ? inv(1 + exp(-z)) : exp(z) / (1 + exp(z))
softplus(z) = max(z, 0.0) + log1p(exp(-abs(z)))
logtarget(z) = -a * softplus(-z) - b * softplus(z)

z = log(a / b)
oldlog = logtarget(z)
samples = zeros(Float64, N)
accept = 0
for k in 1:N
    prop = z + step * randn()
    newlog = logtarget(prop)
    if log(rand()) < min(0.0, newlog - oldlog)
        z, oldlog = prop, newlog
        accept += 1
    end
    samples[k] = sigmoid(z)
end
println("acceptance: ", accept / N)
println("MCMC posterior mean: ", mean(samples[(burn+1):end]))
println("exact posterior mean: ", a / (a+b))
\end{lstlisting}

\begin{zadanie}[Zależność zmienia wariancję średniej]
Załóżmy, że w stacjonarnym łańcuchu $\Var(Y_0)=4$ oraz
$\rho_h=\rho^h$ dla pewnego $|\rho|<1$.
Wyznaczyć $\sigma_\infty^2$ i asymptotyczny ESS.
\par\emph{Rozwiązanie.} Z sumy szeregu geometrycznego
$\tau_{\rm int}=1+2\rho/(1-\rho)=(1+\rho)/(1-\rho)$,
$\sigma_\infty^2=4(1+\rho)/(1-\rho)$, a
$N_{\rm eff}=N(1-\rho)/(1+\rho)$.
Przy $\rho=0{,}9$ efektywna liczebność jest w przybliżeniu
$N/19$, chociaż samych iteracji jest $N$.
Dla ujemnej korelacji ESS może przewyższać $N$.
\end{zadanie}

\chapter{Bootstrap, jackknife i testy permutacyjne: wnioskowanie przez ponowne próbkowanie}
\label{ch:resampling}
\paragraph{Powiązania i odsyłacze.} Porównaj estymatory z rozdziału~\ref{ch:estimators}, metody asymptotyczne z rozdziału~\ref{ch:asymptotic-estimation} i analizę procesów zależnych w rozdziale~\ref{ch:time-series}.


\section{Dlaczego ponownie losujemy dane?}
W wielu problemach znamy statystykę $T_n=T(X_1,\ldots,X_n)$,
ale nie znamy jej rozkładu lub wyprowadzenie go wymaga trudnych
rachunków. \emph{Resampling} polega na odtwarzaniu rozkładu
statystyki przez wielokrotne konstruowanie sztucznych prób.
Najpierw trzeba jednak ustalić, \emph{co jest jednostką losowania}.
Gdy jednostkami są pacjenci, a każdy pacjent ma dziesięć pomiarów,
nie wolno bez uzasadnienia losować niezależnie wszystkich pomiarów:
niszczy to zależność wewnątrz pacjenta.

W tym rozdziale \emph{liczebność rzeczywistej próby} oznaczamy przez $n$,
a \emph{liczbę replik resamplingowych} przez $B$. Zwiększenie $B$
zmniejsza błąd symulacyjny wyznaczania rozkładu resamplingowego,
ale nie zwiększa informacji statystycznej zawartej w $n$ obserwacjach.

\section{Bootstrap nieparametryczny: definicja i dokładny rachunek}
\begin{definicja}[Dystrybuanta empiryczna i próba bootstrapowa]
\label{def:bootstrap-sample}
Niech $(X_i)_{i\ge1}$ będą niezależnymi zmiennymi losowymi
o jednakowym rozkładzie $P$ na $(\R,\mathcal B(\R))$.
Po zaobserwowaniu $x_1,\ldots,x_n$ definiujemy
$P_n=n^{-1}\sum_{i=1}^n\delta_{x_i}$.
\emph{Próbka bootstrapowa} $(X_1^*,\ldots,X_n^*)$ to,
warunkowo względem danych, $n$ niezależnych losowań z $P_n$.
Piszemy $\Prob^*(\cdot)=\Prob(\cdot\mid X_1,\ldots,X_n)$,
$\E^*$ i $\Var^*$ odpowiednio dla warunkowej wartości
oczekiwanej i wariancji, a $\mathcal L^*$ dla warunkowego
rozkładu (względem ustalonych danych). Gwiazdką oznaczamy statystykę
wyznaczoną z próby bootstrapowej.
\end{definicja}

\begin{twierdzenie}[Dokładna wariancja średniej bootstrapowej]
\label{thm:bootstrap-mean-variance}
Niech $n\ge2$, $\E X_1^2<\infty$,
$\overline X_n=n^{-1}\sum_{i=1}^nX_i$, a
$S_n^2=(n-1)^{-1}\sum_{i=1}^n(X_i-\overline X_n)^2$.
Wówczas, dla niemal wszystkich realizacji danych,
\[
 \E^*\overline X_n^*=\overline X_n,\qquad
 \Var^*(\overline X_n^*)=\frac{n-1}{n^2}S_n^2.
\]
W szczególności bootstrapowy błąd standardowy $\sqrt{\Var^*(\overline X_n^*)}$
nie jest przy skończonym $n$ identyczny z klasycznym $S_n/\sqrt n$,
ale ich iloraz, przy $S_n>0$, wynosi $\sqrt{(n-1)/n}$.
\end{twierdzenie}
\begin{proof}
Warunkowo względem danych każda $X_i^*$ ma średnią
$\overline X_n$ i wariancję
$n^{-1}\sum_{j=1}^n(X_j-\overline X_n)^2=(n-1)S_n^2/n$.
Warunkowa niezależność $n$ losowań daje wskazaną wariancję średniej.
\end{proof}

\begin{uwaga}[Parametryczny i nieparametryczny bootstrap]
W wariancie \emph{nieparametrycznym} losujemy z rozkładu empirycznego $P_n$.
W wariancie \emph{parametrycznym} najpierw dopasowujemy rodzinę
$\{P_\theta:\theta\in\Theta\}$, a następnie losujemy z
$P_{\widehat\theta_n}$. Pierwszy wariant nie wymaga wyboru rodziny,
lecz może być mało wiarygodny dla ekstremalnych kwantyli,
gdy w danych nie ma obserwacji dalekiego ogona. Drugi wymaga
wiarygodnego modelu; symulowanie z rozkładu błędnie dobranego
może prowadzić do błędnych przedziałów mimo ogromnego $B$.
\end{uwaga}

\begin{twierdzenie}[Warunkowe centralne twierdzenie graniczne dla bootstrapu średniej]
\label{thm:bootstrap-clt}
Niech $(X_i)_{i\ge1}$ będą niezależne o tym samym rozkładzie,
$\E|X_1|^3<\infty$ i $\sigma^2=\Var(X_1)>0$.
Dla prawie każdej nieskończonej realizacji $(X_i)_{i\ge1}$ zachodzi
\[
 \mathcal L^*\bigl(\sqrt n(\overline X_n^*-\overline X_n)\bigr)
 \Rightarrow\Normal(0,\sigma^2),\qquad n\to\infty.
\]
Jest to ten sam graniczny rozkład co $\sqrt n(\overline X_n-\E X_1)$.
\end{twierdzenie}
\begin{proof}
Oznaczmy $Y_{n,j}=X_j-\overline X_n$ i
$m_{k,n}=n^{-1}\sum_{j=1}^n|Y_{n,j}|^k$ dla $k=2,3$.
Z mocnego prawa wielkich liczb $m_{2,n}\to\sigma^2$
oraz $m_{3,n}=O(1)$ prawie na pewno.
Warunkowa funkcja charakterystyczna jednej wycentrowanej
obserwacji bootstrapowej, podzielonej przez $\sqrt n$, wynosi
\[
 \psi_n(t)=\frac1n\sum_{j=1}^n
 e^{itY_{n,j}/\sqrt n}
 =1-\frac{t^2m_{2,n}}{2n}+R_n(t),\qquad
 |R_n(t)|\le \frac{|t|^3m_{3,n}}{6n^{3/2}},
\]
gdyż $\sum_jY_{n,j}=0$. Dla ustalonego $t$ i niemal każdej
realizacji $n\log\psi_n(t)\to-t^2\sigma^2/2$.
Warunkowa niezależność daje funkcję charakterystyczną sumy
równą $\psi_n(t)^n$, więc dąży ona do
$\exp(-t^2\sigma^2/2)$. Z twierdzenia Lévy'ego
o ciągłości funkcji charakterystycznych wynika teza.
Zwykłe CLT dla średniej z próby stosujemy oddzielnie.
\end{proof}

\begin{wniosek}[Metoda delty dla gładkich funkcji średniej]
\label{cor:bootstrap-delta}
W założeniach twierdzenia~\ref{thm:bootstrap-clt} niech
$g\colon\R\to\R$ będzie różniczkowalna w
$\mu=\E X_1$, przy czym $g'$ jest ciągła w otoczeniu $\mu$.
Wówczas dla prawie każdej realizacji danych
\[
 \mathcal L^*\bigl(\sqrt n[g(\overline X_n^*)-g(\overline X_n)]\bigr)
 \Rightarrow \Normal(0,\sigma^2g'(\mu)^2).
\]
Jeżeli $g'(\mu)=0$, granica jest zdegenerowana i nie należy
bez dalszej analizy stosować zwykłego przybliżenia normalnego.
\end{wniosek}
\begin{proof}
Mocne prawo wielkich liczb daje $\overline X_n\to\mu$,
a twierdzenie~\ref{thm:bootstrap-clt} zapewnia warunkową
ciasność $\sqrt n(\overline X_n^*-\overline X_n)$.
Z twierdzenia o wartości średniej i ciągłości $g'$ otrzymujemy
\[
 g(\overline X_n^*)-g(\overline X_n)
 =g'(\mu)(\overline X_n^*-\overline X_n)
   +o_{\Prob^*}(n^{-1/2})
\]
dla niemal każdej ustalonej realizacji danych.
Wystarczy zastosować warunkowe twierdzenie Slutsky'ego.
\end{proof}

\section{Przedziały bootstrapowe: trzy różne konstrukcje}
Niech $\widehat\theta_n=T_n(X_1,\ldots,X_n)$,
a $\widehat\theta_n^{*(1)},\ldots,\widehat\theta_n^{*(B)}$
będą statystykami z $B$ niezależnych prób bootstrapowych,
warunkowo względem danych. Przez $q_\gamma^*$ oznaczamy
warunkowy kwantyl rzędu $\gamma$ rozkładu
$\widehat\theta_n^*$, w praktyce szacowany kwantylem
z $B$ replik.

\begin{definicja}[Przedział percentylowy i podstawowy]
\label{def:bootstrap-intervals}
Dla $0<\alpha<1$ definiujemy \emph{przedział percentylowy}
$I_{\rm perc}=[q_{\alpha/2}^*,q_{1-\alpha/2}^*]$
oraz \emph{przedział podstawowy} (basic)
\[
 I_{\rm basic}=
 [2\widehat\theta_n-q_{1-\alpha/2}^*,
  2\widehat\theta_n-q_{\alpha/2}^*].
\]
Jeżeli istnieje dodatnie oszacowanie $\widehat s_n$ błędu
standardowego i analogiczne $\widehat s_n^*$ w próbie bootstrapowej,
można też użyć \emph{przedziału studentyzowanego},
opartego na kwantylach
$(\widehat\theta_n^*-\widehat\theta_n)/\widehat s_n^*$.
Wszystkie te konstrukcje wymagają osobnej kontroli pokrycia;
sama ich definicja nie gwarantuje dokładnego poziomu $1-\alpha$.
\end{definicja}

\begin{przyklad}[Dane o wadach: bootstrap może zawieść na brzegu]
\label{ex:bootstrap-zero}
Załóżmy, że w $n=40$ niezależnych próbach Bernoulliego
nie stwierdzono żadnej wady. Wtedy $\widehat p=0$
i rozkład empiryczny skupia całą masę w zerze.
\emph{Każda} próba nieparametrycznego bootstrapu złożona
jest wyłącznie z zer: $\widehat p^*=0$ prawie na pewno,
a oba przedziały z definicji~\ref{def:bootstrap-intervals}
redukują się do $[0,0]$. Taki wynik nie oznacza,
że prawdziwe prawdopodobieństwo wady wynosi zero.
Dla niezależnej próby dwumianowej można zastosować
przedział dokładny Cloppera--Pearsona lub Wilsona
z rozdziału~\ref{ch:confidence-intervals}.
Jest to przykład, w którym zwiększanie $B$ nie naprawi
wad modelu resamplingowego.
\end{przyklad}

\begin{uwaga}[Przybliżenia i diagnostyka]
Bootstrap percentylowy potrafi uwzględniać asymetrię rozkładu
statystyki, ale przy małych próbach, silnym obciążeniu lub na
brzegu przestrzeni parametrów może mieć bardzo złe pokrycie.
Nawet dla regularnych estymatorów nie zawsze ma pokrycie
drugiego rzędu ani nie musi zachowywać granic przestrzeni parametrów.
W praktyce porównujemy co najmniej rozkład replik, liczbę $B$,
stabilność kwantyli i oszacowanie studentyzowane, jeżeli jest dostępne.
Metoda BCa koryguje obciążenie i przyspieszenie; jej pełna teoria
wymaga dodatkowych warunków regularności i nie wynika wyłącznie
z podstawowego CLT bootstrapowego.
\end{uwaga}

\section{Jackknife: usuwanie po jednej obserwacji}
\begin{definicja}[Estymator wariancji jackknife]
\label{def:jackknife}
Niech $n\ge2$, a $T_{(-i)}$ oznacza wartość statystyki
obliczonej z danych bez obserwacji $i$.
Definiujemy $\overline T_{(-)}=n^{-1}\sum_{i=1}^nT_{(-i)}$
oraz
\[
 \widehat V_{\rm jack}=
 \frac{n-1}{n}\sum_{i=1}^n(T_{(-i)}-\overline T_{(-)})^2.
\]
Pierwiastek z $\widehat V_{\rm jack}$ jest błędem standardowym
jackknife; jego poprawność zależy od rodzaju statystyki.
\end{definicja}
\begin{stwierdzenie}[Jackknife dla średniej]
\label{prop:jackknife-mean}
Jeżeli $T_n=\overline X_n$ i $n\ge2$, to
$\widehat V_{\rm jack}=S_n^2/n$.
\end{stwierdzenie}
\begin{proof}
Mamy $T_{(-i)}=(n\overline X_n-X_i)/(n-1)$,
a średnia wartości usuniętych równa się $\overline X_n$.
Wobec tego
\[
 \widehat V_{\rm jack}=
 \frac{n-1}{n}\sum_{i=1}^n
 \frac{(X_i-\overline X_n)^2}{(n-1)^2}
 =\frac{S_n^2}{n}.
\]
\end{proof}

\begin{przyklad}[Dlaczego jackknife zawodzi dla maksimum]
Jeśli $x_{(n)}>x_{(n-1)}$, to spośród $n$ statystyk
maksimum po usunięciu jednej obserwacji $n-1$ wynosi
$x_{(n)}$, a jedna wynosi $x_{(n-1)}$.
Bezpośredni rachunek daje
$\widehat V_{\rm jack}=((n-1)/n)^2
(x_{(n)}-x_{(n-1)})^2$.
Dla niezależnych losowań z $\operatorname{Unif}(0,\theta)$
wariancja maksimum wynosi dokładnie
$\theta^2 n/[(n+1)^2(n+2)]$, ponieważ
$X_{(n)}/\theta$ ma rozkład $\operatorname{Beta}(n,1)$.
Odstęp $D_n=X_{(n)}-X_{(n-1)}$ spełnia
$\Prob(D_n>t)=(1-t/\theta)^n$ dla $0\le t\le\theta$.
Aby to sprawdzić, całkujemy wspólną gęstość
$n(n-1)u^{n-2}/\theta^n$ pary
$(X_{(n-1)},X_{(n)})$ po obszarze
$0<u<v-t<\theta-t$.
Stąd $nD_n/\theta\Rightarrow\operatorname{Exp}(1)$
i wobec tego
$\widehat V_{\rm jack}/\Var(X_{(n)})
\Rightarrow E^2$ dla $E\sim\operatorname{Exp}(1)$.
Granica nie jest stałą jeden, więc jackknife nie daje
zgodnego estymatora wariancji maksimum. Problemem jest
nieregularność tej statystyki, nie niedostateczna liczba
replik.
\end{przyklad}

\section{Testy permutacyjne: dokładność wynika z wymienności}
\begin{twierdzenie}[Dokładny test permutacyjny dwóch grup]
\label{thm:permutation-exact}
Niech $X_1,\ldots,X_m,Y_1,\ldots,Y_k$ będą niezależne
i niech wszystkie mają \emph{ten sam} rozkład $P$,
przy czym $m,k\ge1$, $n=m+k$.
Niech $t$ będzie dowolną mierzalną statystyką przypisania
$m$ z $n$ zaobserwowanych wartości do pierwszej grupy,
a pozostałych $k$ do drugiej. Oznaczmy przez $\mathcal G$
zbiór wszystkich $\binom nm$ przypisań (traktujemy indeksy
jako różne, nawet jeśli wartości się powtarzają).
Dla przypisania obserwowanego $g_0$ definiujemy
\[
 p_{\rm perm}=
 \frac{\#\{g\in\mathcal G:t(g)\ge t(g_0)\}}{\binom nm}.
\]
Wtedy dla każdego $\alpha\in[0,1]$
$\Prob_{H_0}(p_{\rm perm}\le\alpha)\le\alpha$.
Przy braku remisów wartości $t(g)$ i na odpowiedniej
siatce poziomów zachodzi równość.
\end{twierdzenie}
\begin{proof}
Warunkowo względem nieuporządkowanego multizbioru
zaobserwowanych wartości, pod hipotezą wspólnego rozkładu
wszystkie przypisania \emph{indeksowanych obserwacji}
są jednakowo prawdopodobne. Uporządkujmy wartości statystyki
$t(g)$ malejąco, dopuszczając remisy. Dla dowolnego
$r\in\{0,\ldots,\binom nm\}$ liczba przypisań,
dla których liczba wartości $t$ nie mniejszych od obserwowanej
nie przekracza $r$, wynosi co najwyżej $r$.
Dzieląc przez $\binom nm$ uzyskujemy warunkową nierówność
$\Prob(p_{\rm perm}\le\alpha\mid\text{dane nieuporządkowane})
\le\alpha$, a całkując po danych, tezę.
\end{proof}

\begin{uwaga}[Nie jest to automatycznie test równości średnich]
\label{rem:perm-exchangeability}
\emph{Dokładność} powyższego testu wynika z równości
\emph{całych rozkładów} grup, nie z samej równości średnich.
Jeżeli grupy mają różne wariancje lub inny kształt rozkładu,
to permutacja niewystudentyzowanej różnicy średnich może
mieć niewłaściwy poziom przy $H_0:\E X=\E Y$.
W takiej sytuacji można rozważyć statystykę Welcha,
permutację studentyzowaną z odpowiednimi wynikami
asymptotycznymi albo odporny model z właściwymi
błędami standardowymi. Najpierw trzeba ustalić,
\emph{co dokładnie ma być hipotezą zerową}.
\end{uwaga}

\begin{stwierdzenie}[Test znakowych permutacji dla par]
\label{prop:signflip}
Niech $D_1,\ldots,D_n$ będą niezależnymi zmiennymi,
a ich rozkłady będą symetryczne względem zera,
$D_i\overset d=-D_i$, $i=1,\ldots,n$.
Jeśli oznaczymy obserwowane różnice przez $d_i$,
a przez $t(\varepsilon_1d_1,\ldots,\varepsilon_nd_n)$
statystykę dla $\varepsilon_i\in\{-1,1\}$,
to udział spośród $2^n$ znaków dających wartość
co najmniej obserwowaną definiuje poprawną
(w sensie poziomu nie większego od $\alpha$)
wartość p pod wskazaną hipotezą symetrii.
\end{stwierdzenie}
\begin{proof}
Rozkład wektora $D=(D_1,\ldots,D_n)$ jest niezmienniczy
względem grupy wszystkich zmian znaków, ponieważ
każda współrzędna jest symetryczna i współrzędne są niezależne.
Warunkowo względem orbity wektora pod tą skończoną grupą,
zaobserwowany reprezentant jest jednostajny (licząc
wielokrotności wynikające z zer). Argument rangowy
z dowodu twierdzenia~\ref{thm:permutation-exact}
daje nierówność dla wartości p.
\end{proof}

\begin{uwaga}[Monte Carlo zamiast wszystkich permutacji]
Jeżeli $\binom nm$ lub $2^n$ jest zbyt duże,
losujemy $B$ permutacji niezależnie i jednostajnie.
Przy częstości $b$ permutacji o statystyce co najmniej
obserwowanej używa się wartości
$\widehat p=(b+1)/(B+1)$, nie samego $b/B$.
Korekta ta uwzględnia obserwowaną konfigurację
w porównaniu rangowym; w szczególności nigdy nie
raportujemy wartości p równej zeru tylko dlatego,
że przy skończonym $B$ nie trafiliśmy skrajniejszego wyniku.
Dokładność zależy od poprawnego mechanizmu randomizacji.
\end{uwaga}

\section{Zależność w czasie: dlaczego potrzebny jest bootstrap blokowy?}
\label{sec:block-bootstrap}
Załóżmy, że mamy stacjonarny szereg $(X_t)_{t\in\mathbb Z}$
z sumowalnymi autokowariancjami $\gamma_h=\Cov(X_0,X_h)$.
Wówczas, zgodnie z rozdziałem~\ref{ch:time-series},
\[
 n\Var(\overline X_n)\longrightarrow
 \sigma_\infty^2=\gamma_0+2\sum_{h=1}^\infty\gamma_h.
\]
Naiwny bootstrap losujący pojedyncze obserwacje odtwarza
jedynie marginalny rozkład empiryczny. Jeżeli ponadto
szereg jest ergodyczny i $\E X_0^2<\infty$, to
$n\Var^*(\overline X_n^*)=(n-1)S_n^2/n\to\gamma_0$
prawie na pewno, a prawdziwa granica
$n\Var(\overline X_n)\to\sigma_\infty^2$
może być inna. Same założenia stacjonarności i sumowalności
autokowariancji nie zapewniają ergodyczności ani
zbieżności empirycznej wariancji do $\gamma_0$.
Dla modelu AR(1) z korelacją $\rho^h$, $0<\rho<1$,
iloraz poprawnej wariancji asymptotycznej do wariancji
naiwnej wynosi $(1+\rho)/(1-\rho)$.

\begin{definicja}[Bootstrap ruchomych bloków]
Dla próby $x_1,\ldots,x_n$ oraz długości bloku
$1\le\ell\le n$ definiujemy $n-\ell+1$ bloków
$(x_j,\ldots,x_{j+\ell-1})$, $j=1,\ldots,n-\ell+1$.
Losujemy niezależnie i jednostajnie bloki ze zwracaniem,
łączymy je i obcinamy wynik do długości $n$.
Tak otrzymana próba ma długość $n$ i zachowuje
część lokalnych zależności wewnątrz bloków.
\end{definicja}

\begin{uwaga}[Jak dobrać długość bloku?]
W typowych twierdzeniach dla słabo zależnych,
stacjonarnych szeregów wymaga się
$\ell=\ell_n\to\infty$ oraz $\ell_n/n\to0$,
ale nie są to warunki \emph{wystarczające bez dodatkowych
hipotez} mieszania, momentów i właściwego funkcjonału.
Zbyt krótkie bloki zrywają zależność; bardzo długie
zostawiają niewiele efektywnie różnych bloków.
Pomiary powtarzane od jednego pacjenta wymagają z kolei
bootstrapu \emph{klastrowego}, nie blokowego w czasie,
o ile zależność wynika z przypisania do pacjenta.
\end{uwaga}

\section{Dobór metody do problemu}
\begin{center}\small
\begin{tabular}{>{\raggedright\arraybackslash}p{.25\textwidth}>{\raggedright\arraybackslash}p{.29\textwidth}>{\raggedright\arraybackslash}p{.30\textwidth}}
\toprule
\textbf{Sytuacja} & \textbf{Metoda do rozważenia} & \textbf{Warunek krytyczny}\\
\midrule
Niezależne obserwacje, gładka statystyka & Bootstrap nieparametryczny & Reprezentatywność próby; regularność statystyki\\
Silny model parametryczny, mała próba & Bootstrap parametryczny & Model dopasowany poprawnie\\
Średnia i gładki funkcjonał wielu danych & Jackknife lub bootstrap & Rozkład statystyki nie jest skrajnie nieregularny\\
Dwie grupy o jednakowym rozkładzie pod $H_0$ & Test permutacyjny & Wymienność pod hipotezą zerową\\
Różnice w parach, symetria wokół zera & Sign-flip & Niezależność i symetria różnic\\
Obserwacje zależne w czasie & Bootstrap blokowy & Stacjonarność, słaba zależność i długość bloku\\
Dane klastrowe & Bootstrap klastrów & Jednostką resamplingu jest cały klaster\\
\bottomrule
\end{tabular}
\end{center}

\section{Laboratorium: bootstrap, permutacje i zależne obserwacje}
\begin{lstlisting}[language=Python,caption={Kompletny program: bootstrap średniej, dokładne permutacje i bootstrap blokowy (Python, NumPy i SciPy).}]
import itertools
import numpy as np
from scipy.stats import norm

rng = np.random.default_rng(2026)
x = np.array([9., 11., 10., 12., 10., 9., 13., 8.])
B = 20_000
idx = rng.integers(0, len(x), size=(B, len(x)))
boot = x[idx].mean(axis=1)
print("mean:", x.mean(), "bootstrap SE:", boot.std(ddof=1))
print("percentile CI:", np.quantile(boot, [0.025, 0.975]))
print("classical SE:", x.std(ddof=1) / np.sqrt(len(x)))

# Exact permutation test of equal distributions, small samples.
a = np.array([1.0, 2.0, 4.0])
b = np.array([5.0, 6.0, 8.0])
pooled = np.r_[a, b]
observed = abs(a.mean() - b.mean())
total = extreme = 0
for comb in itertools.combinations(range(len(pooled)), len(a)):
    mask = np.zeros(len(pooled), dtype=bool)
    mask[list(comb)] = True
    t = abs(pooled[mask].mean() - pooled[~mask].mean())
    total += 1
    extreme += (t >= observed - 1e-12)
print("exact two-sided permutation p:", extreme / total)

# AR(1): naive iid bootstrap versus moving-block bootstrap.
rho, n, ell = 0.75, 320, 16
noise = rng.standard_normal(n)
y = np.zeros(n)
for i in range(1, n):
    y[i] = rho * y[i-1] + noise[i]
starts = rng.integers(0, n - ell + 1, size=(B, (n+ell-1)//ell))
block_indices = (starts[:, :, None] + np.arange(ell)).reshape(B, -1)[:, :n]
block_means = y[block_indices].mean(axis=1)
iid_means = y[rng.integers(0, n, size=(B, n))].mean(axis=1)
print("iid bootstrap SE:", iid_means.std(ddof=1))
print("block bootstrap SE:", block_means.std(ddof=1))
print("theoretical AR(1) inflation, variance:", (1+rho)/(1-rho))
\end{lstlisting}

\begin{lstlisting}[language=Julia,caption={Kompletny program: bootstrap, dokładny test permutacyjny i resampling bloków (Julia, Random i Statistics).}]
using Random, Statistics
Random.seed!(2026)
x = Float64[9, 11, 10, 12, 10, 9, 13, 8]
B = 10_000
boot = [mean(x[rand(1:length(x), length(x))]) for _ in 1:B]
println("bootstrap mean SE: ", std(boot))
println("percentile CI: ", quantile(boot, [0.025, 0.975]))

# All 3-element subsets of six indices, implemented without packages.
a = Float64[1, 2, 4]
b = Float64[5, 6, 8]
values = vcat(a, b)
obs = abs(mean(a) - mean(b))
extreme = 0
total = 0
for i in 1:4, j in (i+1):5, k in (j+1):6
    chosen = [i, j, k]
    other = [z for z in 1:6 if !(z in chosen)]
    stat = abs(mean(values[chosen]) - mean(values[other]))
    extreme += stat >= obs - 1e-12
    total += 1
end
println("exact permutation p: ", extreme / total)

rho, n, ell = 0.75, 320, 16
y = zeros(n)
noise = randn(n)
for i in 2:n
    y[i] = rho*y[i-1] + noise[i]
end
nb = cld(n, ell)
block = zeros(B)
for b in 1:B
    starts = rand(1:(n-ell+1), nb)
    sample = reduce(vcat, [y[s:(s+ell-1)] for s in starts])[1:n]
    block[b] = mean(sample)
end
println("block bootstrap SE: ", std(block))
\end{lstlisting}

\begin{zadanie}[Próba z zerową liczbą sukcesów]
W $n=20$ niezależnych próbach Bernoulliego nie obserwowano sukcesów.
Opisać przedział percentylowy nieparametrycznego bootstrapu
i wskazać alternatywę o poprawnym skończonopróbkowym pokryciu.
\par\emph{Rozwiązanie.} Rozkład empiryczny jest skupiony w zerze,
zatem każdy rozkład bootstrapowy $\widehat p^*$ jest skupiony
w zerze i przedział percentylowy wynosi $[0,0]$ dla dowolnego $B$.
Przedział dokładny Cloppera--Pearsona ma dolną granicę zero
i górną $1-(\alpha/2)^{1/n}$ w wersji dwustronnej,
ponieważ $\Prob_p(S_n=0)=(1-p)^n$.
Dla jednostronnego górnego przedziału ufności na poziomie
$1-\alpha$ stosujemy granicę $1-\alpha^{1/n}$.
\end{zadanie}


\chapter{Metody nieparametryczne i odporne: rozkład, rangi i obserwacje odstające}
\label{ch:nonparametric}
\paragraph{Powiązania i odsyłacze.} Dystrybuantę empiryczną omówiliśmy w rozdziale~\ref{ch:statystyki}, a bootstrap i testy permutacyjne w rozdziale~\ref{ch:resampling}.


\section{Co oznacza brak założenia parametrycznego?}
W modelu parametrycznym ustalamy rodzinę rozkładów
$\{P_\theta:\theta\in\Theta\}$, a w modelu nieparametrycznym
przestrzeń dopuszczalnych rozkładów jest znacznie większa,
np. obejmuje wszystkie rozkłady o skończonej wariancji.
Nie znaczy to, że metoda nieparametryczna nie wymaga
\emph{żadnych} założeń: test rangowy może wymagać wymienności,
estymacja gęstości -- gładkości, a test dla par -- odpowiednio
symetrycznego rozkładu różnic. Często wybieramy pomiędzy
\emph{różnymi hipotezami} i różnymi miarami efektu,
a nie tylko pomiędzy dwoma wzorami na wartość p.

\section{Dystrybuanta empiryczna i dokładność jej aproksymacji}
\begin{definicja}[Dystrybuanta empiryczna]
Niech $X_1,\ldots,X_n$ będą niezależne o jednakowej
dystrybuancie $F$. Definiujemy
$F_n(t)=n^{-1}\sum_{i=1}^n\ind{\{X_i\le t\}}$,
$t\in\R$. Dla ustalonego $t$ zmienna $nF_n(t)$
ma rozkład $\Bin(n,F(t))$.
\end{definicja}

\begin{twierdzenie}[Jednopunktowa kontrola dystrybuanty]
\label{thm:empirical-hoeffding}
Dla dowolnego $n\ge1$, $t\in\R$ i $\varepsilon>0$,
\[
 \Prob\bigl(|F_n(t)-F(t)|\ge\varepsilon\bigr)
 \le2\exp(-2n\varepsilon^2).
\]
\end{twierdzenie}
\begin{proof}
Zmienne $\ind{\{X_i\le t\}}$ są niezależne, przyjmują
wartości w $[0,1]$ i mają wspólną średnią $F(t)$.
Nierówność Hoeffdinga dla średniej z ograniczonych
zmiennych daje tezę.
\end{proof}

\begin{uwaga}[Co znaczy jednoczesny pas ufności?]
Nierówność powyższa dotyczy \emph{ustalonego} $t$.
Nie wolno jej bez uzasadnienia stosować równocześnie
do wszystkich $t\in\R$. Silniejsza nierówność
Dvoretzky'ego--Kiefera--Wolfowitza--Massarta głosi,
że dla próby niezależnej o jednakowym rozkładzie
\[
 \Prob\left(\sup_{t\in\R}|F_n(t)-F(t)|>\varepsilon\right)
 \le2\exp(-2n\varepsilon^2).
\]
Jest to głębsze twierdzenie, którego pełnego dowodu
tu nie przeprowadzamy. Daje jednoczesny pas
$F_n(t)\pm\sqrt{\log(2/\alpha)/(2n)}$, obcięty do
$[0,1]$, o pokryciu co najmniej $1-\alpha$.
Nie należy mylić go z punktowym przedziałem dla $F(t)$.
\end{uwaga}

\section{Mediana, kwantyle i odporność na skażenia}
\begin{definicja}[Funkcjonał kwantyla i punkt skażenia]
Dla dystrybuanty $F$ i $p\in(0,1)$ definiujemy
$Q_p(F)=\inf\{x\in\R:F(x)\ge p\}$.
Dla $0\le\varepsilon\le1$ i $z\in\R$ rozważamy
\emph{model skażenia}
$F_{\varepsilon,z}=(1-\varepsilon)F+\varepsilon\delta_z$,
gdzie $\delta_z$ oznacza rozkład punktowy w $z$.
\emph{Funkcją wpływu} funkcjonału $T$ w punkcie $z$ nazywamy,
jeśli granica istnieje,
\[
 \mathrm{IF}(z;T,F)=
 \lim_{\varepsilon\downarrow0}
 \frac{T(F_{\varepsilon,z})-T(F)}{\varepsilon}.
\]
Jest to miara lokalnej wrażliwości na niewielką domieszkę
obserwacji w punkcie $z$, nie zaś dokładny opis
arbitralnie dużego skażenia.
\end{definicja}

\begin{twierdzenie}[Funkcje wpływu średniej i kwantyla]
\label{thm:influence-mean-quantile}
Niech $F$ ma skończoną średnią $\mu$.
Dla funkcjonału $T(F)=\int_\R x\,\dd F(x)$
zachodzi $\mathrm{IF}(z;T,F)=z-\mu$.
Niech ponadto $q=Q_p(F)$ będzie jedynym kwantylem
rzędu $p$, a $F$ ma w otoczeniu $q$ ciągłą,
dodatnią gęstość $f$, tak że $F(q)=p$.
Wówczas dla każdego $z\ne q$
\[
 \mathrm{IF}(z;Q_p,F)=
 \frac{p-\ind{\{z\le q\}}}{f(q)}.
\]
Funkcja wpływu kwantyla jest zatem ograniczona
dla ustalonego $F$ i $p$, podczas gdy funkcja wpływu
średniej jest nieograniczona.
\end{twierdzenie}
\begin{proof}
Średnia przy skażeniu wynosi
$(1-\varepsilon)\mu+\varepsilon z$, co daje pierwszy wzór.
Dla $z\ne q$ przy małym $\varepsilon$ kwantyl
$q_\varepsilon=Q_p(F_{\varepsilon,z})$ pozostaje
w otoczeniu $q$ nie zawierającym $z$.
Rozwiązujemy zatem równanie
$(1-\varepsilon)F(q_\varepsilon)+
\varepsilon\ind{\{z\le q\}}=p$.
Różniczkując jednostronnie w $\varepsilon=0$,
przy $f(q)>0$, otrzymujemy
$-p+f(q)q'_0+\ind{\{z\le q\}}=0$.
To daje drugi wzór i wskazane ograniczenia.
\end{proof}

\begin{przyklad}[Jedna obserwacja odstająca i dwie różne odpowiedzi]
Rozważmy pomiary $9,10,10,11,12$. Średnia wynosi $10{,}4$,
a mediana $10$. Jeśli pomiar $12$ zostanie zastąpiony
przez $120$, nowa średnia wyniesie $32$, ale mediana
pozostanie równa $10$. Mediana ma wysoką odporność
na kilka skrajnych wartości, lecz przy regularnym modelu
normalnym może być mniej efektywna od średniej.
Wybór zależy od celu: jeśli mierzymy oczekiwany
\emph{łączny koszt}, średnia może być właściwym funkcjonałem
nawet dla rozkładu z ogonem; jeśli szukamy typowego
poziomu, mediana może być łatwiejsza do interpretacji.
\end{przyklad}

\begin{center}
\begin{tikzpicture}[x=.072cm,y=1.0cm,>=Latex]
\draw[->] (0,0)--(130,0) node[right] {pomiar};
\foreach \x in {0,20,40,60,80,100,120} {
 \draw (\x,-.06)--(\x,.06) node[below=4pt,font=\scriptsize] {\x};
}
\foreach \x in {9,10,11} {
 \fill[blue!70!black] (\x,0.40) circle[radius=2.7pt];
}
\fill[blue!70!black] (10,.62) circle[radius=2.7pt];
\fill[red!75!black] (120,.40) circle[radius=3.2pt];
\draw[red!70!black,very thick] (32,.18)--(32,.86)
 node[above,font=\scriptsize] {średnia $32$};
\draw[green!50!black,very thick] (10,.86)--(10,1.35)
 node[above,font=\scriptsize] {mediana $10$};
\node[anchor=west,font=\scriptsize] at (52,1.20)
 {Punkt czerwony zastępuje $12$ przez $120$};
\end{tikzpicture}

{\small\emph{Rysunek 26.1.} Jeden skrajny pomiar może silnie przesunąć średnią,\newline
nie zmieniając mediany. Punkty w pobliżu zera nachodzą na siebie w tej skali.}
\end{center}

\section{Testy rangowe dla dwóch niezależnych grup}
\begin{definicja}[Statystyka Manna--Whitneya]
\label{def:mann-whitney}
Dla niezależnych prób $X_1,\ldots,X_m$ i
$Y_1,\ldots,Y_k$ definiujemy
\[
 U=\sum_{i=1}^m\sum_{j=1}^k
 \left(\ind{\{X_i<Y_j\}}+
 \frac12\ind{\{X_i=Y_j\}}\right).
\]
Jeśli grupy są niezależne i każda jest próbą
z odpowiedniego rozkładu $F$ lub $G$, to
$\E U/(mk)=\Prob(X<Y)+\tfrac12\Prob(X=Y)$,
gdzie $X\sim F$, $Y\sim G$ są niezależne.
\end{definicja}

\begin{twierdzenie}[Rozkład rang i momenty $U$ pod hipotezą zerową]
\label{thm:mann-whitney-moments}
Niech wszystkie $m+k$ obserwacji będą niezależne
i mają wspólny \emph{ciągły} rozkład, przy $m,k\ge1$.
Wtedy
\[
 \E U=\frac{mk}{2},\qquad
 \Var(U)=\frac{mk(m+k+1)}{12}.
\]
Ponadto rozkład $U$ zależy tylko od $(m,k)$,
a nie od wspólnego ciągłego rozkładu obserwacji.
\end{twierdzenie}
\begin{proof}
Przez ciągłość wszystkie obserwacje są różne prawie na pewno.
W wyniku wymienności wspólnych obserwacji każda z
$\binom{m+k}{m}$ konfiguracji rang pierwszej grupy
jest jednakowo prawdopodobna. Jeśli $R_1,\ldots,R_m$
oznaczają jej rangi (od najmniejszej), to
$U=mk+m(m+1)/2-\sum_{i=1}^mR_i$ przy naszej
konwencji zliczania $X_i<Y_j$.
Z losowania bez zwracania $m$ rang z $1,\ldots,m+k$
wynika $\E\sum_iR_i=m(m+k+1)/2$.
Wariancja sumy rang jest równa
\[
 \frac{mk}{m+k-1}\cdot
 \frac{(m+k)^2-1}{12}
 =\frac{mk(m+k+1)}{12}.
\]
Odjęcie sumy rang od stałej nie zmienia wariancji.
Rozkład rang jest jednostajny, więc nie zależy od
wspólnej dystrybuanty.
\end{proof}

\begin{uwaga}[Mann--Whitney nie jest automatycznie testem median]
Hipoteza o identycznych rozkładach $F=G$ daje
\emph{dokładną} rangową kalibrację testu.
Przy rodzinie przesunięć o tym samym kształcie można
go interpretować jako test położenia.
Bez tego założenia różnica median nie jest równoważna
zmianie $\Prob(X<Y)$; grupy mogą mieć tę samą medianę,
ale inne wariancje i istotnie różne wartości statystyki $U$.
Gdy celem jest test \emph{równości średnich}, zamiast
Manna--Whitneya warto rozważyć test Welcha,
studentyzowaną permutację lub właściwy model odpornościowy.
\end{uwaga}

\section{Dane sparowane: test znaków i Wilcoxona}
\begin{twierdzenie}[Test Wilcoxona rangowanych znaków]
\label{thm:wilcoxon-signed-rank}
Niech $D_1,\ldots,D_n$ będą niezależne i mają ten sam
ciągły rozkład symetryczny względem zera, a
$\Prob(D_i=0)=0$. Niech $R_i$ będzie rangą
$|D_i|$ wśród $|D_1|,\ldots,|D_n|$, a
$W_+=\sum_{i=1}^nR_i\ind{\{D_i>0\}}$.
Wówczas
\[
 \E W_+=\frac{n(n+1)}4,\qquad
 \Var(W_+)=\frac{n(n+1)(2n+1)}{24}.
\]
Warunkowo względem wartości bezwzględnych $|D_i|$
rozkład $W_+$ jest taki sam jak rozkład
$\sum_{r=1}^nrB_r$, gdzie $B_r$ są niezależnymi
zmiennymi Bernoulliego o parametrze $1/2$.
\end{twierdzenie}
\begin{proof}
Z ciągłości i symetrii znak każdej różnicy jest
niezależny od jej wartości bezwzględnej, przyjmuje
$+1$ i $-1$ z prawdopodobieństwem $1/2$,
a znaki wszystkich różnic są niezależne.
Po warunkowaniu wartości bezwzględnych rangi tworzą pewną permutację
$1,\ldots,n$, zatem $W_+$ ma wskazany rozkład sumy
niezależnych ważonych zmiennych Bernoulliego.
Jego średnia wynosi $\tfrac12\sum_{r=1}^nr=n(n+1)/4$,
a wariancja $\tfrac14\sum_{r=1}^nr^2
=n(n+1)(2n+1)/24$.
\end{proof}

\begin{uwaga}[Warunki testu i postępowanie z remisami]
Test znaków bada same znaki różnic, podczas gdy
Wilcoxon wykorzystuje dodatkowo rangi wartości
bezwzględnych. Dla różnic o symetrycznym rozkładzie
test Wilcoxona bywa mocniejszy, lecz jego
\emph{dokładny} rozkład opisany wyżej wymaga
braku zerowych różnic i remisów rang.
Przy zerach i remisach trzeba jawnie wybrać
konwencję i kalibrację. Test Wilcoxona nie jest
automatycznie dokładnym testem hipotezy
$\E D_i=0$ dla dowolnego niesymetrycznego rozkładu.
\end{uwaga}

\section{Porównywanie więcej niż dwóch grup i danych sparowanych}
\begin{stwierdzenie}[Statystyka Kruskala--Wallisa i jej kalibracja]
\label{prop:kruskal-wallis}
Niech $k\ge2$ grup składa się z niezależnych próbek,
grupa $j$ ma $n_j\ge1$ obserwacji i
$N=\sum_{j=1}^kn_j$. Przy wspólnym \emph{ciągłym}
rozkładzie pod $H_0$ oznaczmy przez $\overline R_j$
średnią rang w grupie $j$.
Statystyka Kruskala--Wallisa, przy braku remisów, wynosi
\[
 H=\frac{12}{N(N+1)}
 \sum_{j=1}^kn_j
 \left(\overline R_j-\frac{N+1}{2}\right)^2.
\]
Jeżeli $k$ jest stałe, $\min_jn_j\to\infty$ i
$\min_j n_j/N$ pozostaje oddalone od zera,
to $H\Rightarrow\chi^2_{k-1}$ pod $H_0$.
Dla małych prób można zastosować dokładny lub
Monte Carlo test permutacyjny na rangach.
\end{stwierdzenie}
\begin{proof}
Wzór jest definicją statystyki jako ważonej sumy
kwadratów odchyleń średnich rang od rangi globalnej.
Ponieważ pod $H_0$ wszystkie przypisania rang do grup
są jednakowo prawdopodobne, permutacyjna kalibracja
wynika z twierdzenia~\ref{thm:permutation-exact}
(po zastąpieniu dwóch grup dowolnym ustalonym
podziałem na $k$ grup). Asymptotyczny rozkład
$\chi^2_{k-1}$ wynika z wielowymiarowego CLT dla
rang przy losowym podziale i jest tu \emph{przywoływany},
a nie dowodzony w pełnej ogólności.
\end{proof}

\begin{uwaga}[Przykład: trzy grupy i pomiary wielokrotne]
Kruskal--Wallis wymaga niezależnych grup.
Przy pomiarach tego samego pacjenta w kilku warunkach
obserwacje nie są niezależne; stosujemy np.
test Friedmana, który rangowuje warunki \emph{wewnątrz}
każdego bloku lub pacjenta, ewentualnie model mieszany.
W sytuacji z dwoma pomiarami sparowanymi
rozważamy Wilcoxona lub test znaków,
a nie Manna--Whitneya.
Przed wyborem testu zawsze ustalamy, która jednostka
jest niezależna i jaka hipoteza ma zostać sprawdzona.
\end{uwaga}

\section{Test Kołmogorowa--Smirnowa i problem dopasowywania parametrów}
\begin{stwierdzenie}[Bezrozkładowość testu KS dla znanego $F_0$]
\label{prop:ks-distributionfree}
Niech $X_1,\ldots,X_n$ będą niezależne o wspólnej,
\emph{znanej pod hipotezą zerową}, ciągłej i ściśle
rosnącej dystrybuancie $F_0$ na jej nośniku.
Wtedy rozkład statystyki
$D_n=\sup_{t\in\R}|F_n(t)-F_0(t)|$
nie zależy od $F_0$, tylko od $n$.
\end{stwierdzenie}
\begin{proof}
Z transformacji całkowej prawdopodobieństwa
$U_i=F_0(X_i)$ mają niezależny rozkład
$\operatorname{Unif}(0,1)$. Po przepisaniu
$F_n(t)=n^{-1}\sum_{i=1}^n\ind{\{U_i\le F_0(t)\}}$
i przejściu do supremum po $u\in(0,1)$
dostajemy statystykę zbudowaną wyłącznie z
niezależnych zmiennych jednostajnych.
\end{proof}

\begin{uwaga}[Błąd często spotykany przy testach normalności]
Jeśli z tych samych danych oszacujemy średnią i
odchylenie standardowe, a potem obliczymy
$\sup_t|F_n(t)-\Phi((t-\overline X_n)/S_n)|$,
to \emph{nie} wolno używać klasycznych wartości
krytycznych testu KS dla ustalonej, znanej $F_0$.
Rozkład statystyki zmienia się przez estymację parametrów.
W takim problemie rozważamy odpowiednią kalibrację
Lillieforsa, bootstrap parametryczny odtwarzający
\emph{całą} procedurę estymacji albo test Shapiro--Wilka,
którego pełne omówienie umieścimy w osobnym
rozdziale diagnostycznym.
\end{uwaga}

\section{Estymacja gęstości metodą jądrową}
\begin{definicja}[Estymator jądrowy KDE]
\label{def:kde}
Niech $X_1,\ldots,X_n$ będą niezależne o wspólnej
gęstości $f$ względem $\lambda_1$. Niech
$K\colon\R\to[0,\infty)$ będzie mierzalna,
$\int_\R K(u)\,\dd\lambda_1(u)=1$,
a $h>0$ oznacza \emph{szerokość pasma}.
Definiujemy
\[
 \widehat f_{n,h}(x)=\frac1{nh}
 \sum_{i=1}^nK\left(\frac{x-X_i}{h}\right).
\]
\end{definicja}

\begin{twierdzenie}[Obciążenie, wariancja i dobór pasma KDE]
\label{thm:kde-bias-var}
Załóżmy, że $f\in C^2(\R)$, $f$ i $f''$ są ograniczone,
a $f$ jest ciągła w ustalonym $x$.
Niech $K\ge0$ spełnia
$\int_\R K(u)\,\dd\lambda_1(u)=1$, $\int_\R uK(u)\,\dd\lambda_1(u)=0$,
$\mu_2(K)=\int_\R u^2K(u)\,\dd\lambda_1(u)<\infty$,
oraz $R(K)=\int_\R K(u)^2\,\dd\lambda_1(u)<\infty$.
Przy $h=h_n\downarrow0$, $nh_n\to\infty$ mamy
\[
 \begin{aligned}
 \E\widehat f_{n,h}(x)-f(x)
   &=\tfrac12h^2\mu_2(K)f''(x)+o(h^2),\\
 \Var(\widehat f_{n,h}(x))
   &=\frac{f(x)R(K)}{nh}+o((nh)^{-1}).
 \end{aligned}
\]
Jeżeli $f(x)>0$, $f''(x)\ne0$ i $\mu_2(K)>0$,
to minimalizacja głównych członów punktowego MSE daje
$h\asymp n^{-1/5}$, a MSE rzędu $n^{-4/5}$.
\end{twierdzenie}
\begin{proof}
Zmiana zmiennej w całce daje
$\E\widehat f_{n,h}(x)=
\int_\R K(u)f(x-hu)\,\dd\lambda_1(u)$.
Ze wzoru Taylora
$f(x-hu)=f(x)-hu f'(x)+
\tfrac12h^2u^2f''(x)+h^2u^2r_h(u)$,
gdzie $r_h(u)\to0$ dla każdego $u$, a dzięki
globalnej ograniczoności $f''$ mamy jednostajne
ograniczenie $|r_h(u)|\le C$.
Symetria pierwszego momentu jądra i zbieżność
zdominowana dla $u^2K(u)$ dają wzór na obciążenie.
Dla wariancji niezależność implikuje
\[
 \Var(\widehat f_{n,h}(x))=
 \frac1n\left[\frac1h
 \int_\R K(u)^2f(x-hu)\,\dd\lambda_1(u)
 -(\E\widehat f_{n,h}(x))^2\right].
\]
Dominacja przez $\norm f_\infty K(u)^2$ i ciągłość $f$
dają zbieżność pierwszej całki do $f(x)R(K)$.
Drugi składnik jest $O(n^{-1})=o((nh)^{-1})$.
Kwadrat dominującego obciążenia ma rząd $h^4$,
a wariancja rząd $(nh)^{-1}$; zrównanie tych rzędów
daje $h\asymp n^{-1/5}$.
\end{proof}

\begin{uwaga}[KDE przy granicy nośnika]
Ponieważ zwykły estymator KDE nie uwzględnia z góry granic
nośnika, jego poprawność punktowa we wnętrzu przestrzeni nie
musi przenosić się na brzeg.
Jeżeli prawdziwy rozkład ma nośnik $[0,\infty)$,
a stosujemy symetryczne jądro bez korekty, to
przy $x=0$ część masy jądra wychodzi poza nośnik,
a zwykły wynik o obciążeniu rzędu $h^2$ nie musi
obowiązywać. Rozważamy transformację (np. logarytm),
jądro jednostronne lub odpowiednią korektę brzegową.
Przed porównywaniem KDE różnych prób należy też
ustalić sposób doboru pasma: zbyt małe daje poszarpaną
gęstość, zbyt duże ukrywa rzeczywiste wielomodalności.
\end{uwaga}

\section{Estymatory M odporne na obserwacje odstające}
\begin{definicja}[Estymator Hubera w modelu położenia]
\label{def:huber}
Niech $c>0$ i $s>0$ będą ustalonymi parametrami.
Definiujemy funkcję
$\psi_c(u)=\min\{c,\max\{-c,u\}\}$.
\emph{Estymatorem położenia Hubera} nazywamy rozwiązanie
względem $t$ równania
\[
 \sum_{i=1}^n\psi_c\left(\frac{X_i-t}{s}\right)=0.
\]
Jest to warunek pierwszego rzędu minimalizacji
sumy strat $\rho_c((X_i-t)/s)$, gdzie
$\rho_c(u)=u^2/2$ dla $|u|\le c$ oraz
$\rho_c(u)=c|u|-c^2/2$ dla $|u|>c$.
\end{definicja}

\begin{stwierdzenie}[Asymptotyczna wariancja estymatora Hubera]
\label{prop:huber-asympt}
Niech $X_i=\theta+\varepsilon_i$, gdzie $\varepsilon_i$
są niezależne o jednakowym rozkładzie symetrycznym
względem zera, $s>0$ jest znane,
$\E\psi_c(\varepsilon_1/s)^2<\infty$ oraz
$A=\E\psi_c'(\varepsilon_1/s)>0$.
Załóżmy, że rozkład $\varepsilon_1$ nie ma atomów
w punktach $\pm cs$, a rozwiązania $\widehat\theta_n$
są zgodne i spełniają lokalne rozwinięcie
$\sum_i\psi_c((X_i-\widehat\theta_n)/s)
=\sum_i\psi_c(\varepsilon_i/s)
-(\widehat\theta_n-\theta) n A/s
+o_\Prob(\sqrt n+ n|\widehat\theta_n-\theta|)$.
Wówczas
\[
 \sqrt n(\widehat\theta_n-\theta)
 \Rightarrow\Normal\left(0,
 \frac{s^2\E[\psi_c(\varepsilon_1/s)^2]}{A^2}\right).
\]
\end{stwierdzenie}
\begin{proof}
Symetria i nieparzystość $\psi_c$ zapewniają
$\E\psi_c(\varepsilon_1/s)=0$.
CLT daje normalną granicę dla
$n^{-1/2}\sum_i\psi_c(\varepsilon_i/s)$.
Wykorzystujemy rozwinięcie z tezy w równaniu estymatora.
Dodatniość $A$ i zgodność umożliwiają przeniesienie
składnika liniowego na drugą stronę oraz otrzymanie
reprezentacji
\[
 \sqrt n(\widehat\theta_n-\theta)
 =\frac{s}{A\sqrt n}\sum_{i=1}^n
 \psi_c(\varepsilon_i/s)+o_\Prob(1).
\]
Twierdzenie Slutsky'ego prowadzi do wskazanej granicy.
Warunek lokalnego rozwinięcia jest standardowym krokiem
teorii M-estymatorów, a jego weryfikacja dla konkretnego
modelu wymaga odrębnej kontroli regularności.
\end{proof}

\begin{uwaga}[Interpretacja parametru $c$]
Małe $c$ silniej ogranicza wpływ skrajnych reszt,
ale może zmniejszać efektywność w dokładnie normalnym
modelu. Duże $c$ zbliża algorytm do średniej i MNK,
lecz zwiększa wrażliwość na skrajne obserwacje.
Przed zastosowaniem trzeba zdefiniować skalę $s$
(np. odpornie przez MAD), ponieważ próg $cs$ działa
w jednostkach oryginalnych danych. Wyboru $c$
nie należy optymalizować jedynie na podstawie
pojedynczej realizacji z obserwacją odstającą.
\end{uwaga}

\section{Praktyczna mapa decyzji}
\begin{center}\small
\begin{tabular}{>{\raggedright\arraybackslash}p{.26\textwidth}>{\raggedright\arraybackslash}p{.29\textwidth}>{\raggedright\arraybackslash}p{.29\textwidth}}
\toprule
\textbf{Cel} & \textbf{Metoda} & \textbf{Sprawdź najpierw}\\
\midrule
Porównać dwa niezależne rozkłady & Mann--Whitney, KS, permutacje & Czy interesuje Cię położenie, średnia czy cała dystrybuanta?\\
Porównać pomiary przed i po & Wilcoxon, test znaków, model par & Czy parowanie jest rzeczywiste i czy rozkład różnic jest symetryczny?\\
Porównać co najmniej trzy grupy & Kruskal--Wallis & Czy grupy są niezależne, czy to powtarzane pomiary?\\
Porównać powtarzane warunki & Friedman, model mieszany & Pacjent/blok jest jednostką niezależną\\
Zbadać zgodność z zadanym rozkładem & KS, Anderson--Darling, Shapiro--Wilk & Czy parametry hipotezy są znane czy estymowane?\\
Oszacować gęstość i mody & KDE & Pasmo, granice nośnika i liczba obserwacji\\
Oszacować typowy poziom przy outlierach & Mediana, estymator Hubera & Czy celem jest mediana/położenie, czy prawdziwa średnia?\\
\bottomrule
\end{tabular}
\end{center}

\section{Laboratorium: rangi, gęstość i odporność}
\begin{lstlisting}[language=Python,caption={Kompletny program: testy rangowe, KS, KDE oraz estymator Hubera (Python, NumPy i SciPy).}]
import numpy as np
from scipy import stats

x = np.array([9., 10., 10., 11., 12.])
y = np.array([10., 11., 12., 13., 14.])
print("Mann-Whitney:", stats.mannwhitneyu(x, y, alternative="two-sided"))

before = np.array([21., 25., 23., 28., 22., 24.])
after = np.array([20., 23., 20., 24., 17., 18.])
print("Wilcoxon paired:", stats.wilcoxon(after-before, method="exact"))

samples = [np.array([6., 7., 6., 8., 7.]),
           np.array([9., 8., 10., 9., 11.]),
           np.array([11., 12., 10., 13., 12.])]
print("Kruskal-Wallis:", stats.kruskal(*samples))

rng = np.random.default_rng(2026)
z = rng.normal(size=80)
# Valid KS calibration: test against a fully specified N(0,1).
print("KS fixed normal:", stats.kstest(z, "norm", args=(0, 1)))
# For fitted mean and variance, this p-value would not be valid.

contaminated = np.r_[rng.normal(size=80), [15.0, 20.0]]
print("mean, median:", np.mean(contaminated), np.median(contaminated))

scale = stats.median_abs_deviation(contaminated, scale="normal")
c = 1.345
center = np.median(contaminated)
for _ in range(100):
    residual = (contaminated - center) / scale
    weights = np.ones_like(residual)
    mask = np.abs(residual) > c
    weights[mask] = c / np.abs(residual[mask])
    next_center = np.sum(weights * contaminated) / np.sum(weights)
    if abs(next_center - center) < 1e-12:
        break
    center = next_center
print("Huber location:", center)

kde = stats.gaussian_kde(contaminated)
for point in [-1.0, 0.0, 1.0]:
    print("KDE at", point, float(kde(point)[0]))
\end{lstlisting}

\begin{lstlisting}[language=Julia,caption={Kompletny program: odporna estymacja i jądrowa estymacja gęstości (Julia, Random i Statistics).}]
using Random, Statistics
Random.seed!(2026)
x = vcat(randn(80), [15.0, 20.0])
println("mean: ", mean(x), " median: ", median(x))

# Huber location with robust initial MAD scale.
scale = 1.4826 * median(abs.(x .- median(x)))
c = 1.345
center = median(x)
for _ in 1:100
    residual = (x .- center) ./ scale
    weight = min.(1.0, c ./ max.(abs.(residual), 1e-12))
    nextcenter = sum(weight .* x) / sum(weight)
    if abs(nextcenter - center) < 1e-12
        center = nextcenter
        break
    end
    center = nextcenter
end
println("Huber location: ", center)

# Gaussian KDE implemented with elementary functions.
n = length(x)
sigma = std(x)
h = 1.06 * sigma * n^(-1/5)  # rule of thumb; sensitive to outliers
kde(t) = sum(exp.(-0.5 .* ((t .- x)./h).^2)) / (n*h*sqrt(2*pi))
for point in (-1.0, 0.0, 1.0)
    println("KDE at ", point, ": ", kde(point))
end
\end{lstlisting}

\begin{zadanie}[Dlaczego test Manna--Whitneya nie zawsze bada średnią?]
Niech $X$ będzie stała równa zero, a $Y$ przyjmuje
wartość $-1$ z prawdopodobieństwem $0{,}9$ i wartość
$9$ z prawdopodobieństwem $0{,}1$. Wyznaczyć
$\E X$, $\E Y$ oraz $\Prob(X<Y)$.
\par\emph{Rozwiązanie.} Mamy $\E X=0$,
$\E Y=-0{,}9+0{,}9=0$, ale
$\Prob(X<Y)=\Prob(Y=9)=0{,}1$.
Wobec tego równość średnich nie wymusza
$\Prob(X<Y)=1/2$. Test oparty na porządku obserwacji
może wykrywać różnicę mimo równości średnich;
nie jest to sprzeczność, lecz różnica między hipotezami.
\end{zadanie}


\chapter{Diagnostyka rozkładu i założeń: normalność, wariancje i zależność}
\label{ch:diagnostics}
\paragraph{Powiązania i odsyłacze.} Testy parametryczne przedstawiono w rozdziale~\ref{ch:parametric-tests}, modele regresji w rozdziale~\ref{ch:linear-regression}, a ANOVA/ANCOVA rozwijamy w rozdziale~\ref{ch:anova-ancova}.


\section{Co sprawdzamy przed wyborem testu?}
Nie istnieje uniwersalny test, który odpowiada na pytanie, czy dane nadają się do analizy statystycznej. Najpierw formułujemy model, cel i jednostkę niezależnego losowania; następnie sprawdzamy wymagania konkretnego wniosku. W modelu regresji $Y_i=x_i^\top\beta+\varepsilon_i$ normalność -- gdy jest potrzebna do dokładnych testów $t$ i $F$ -- dotyczy \emph{warunkowego rozkładu błędów} $\varepsilon_i$ przy danym projekcie $x_i$, a nie marginalnego rozkładu wszystkich wartości $Y_i$. Jeśli grupy mają różne średnie, połączone dane mogą być wielomodalne, mimo że każdy błąd wewnątrz grupy jest normalny.

\begin{przyklad}[Trzy grupy, normalne błędy, nienormalna mieszanina]
Niech $G$ ma rozkład jednostajny na $\{1,2,3\}$, $\varepsilon\sim\Normal(0,1)$ i $G$ będzie niezależne od $\varepsilon$. Dla $(\mu_1,\mu_2,\mu_3)=(-5,0,5)$ definiujemy $Y=\mu_G+\varepsilon$. Każdy rozkład $\Law(Y\mid G=j)$ jest normalny, ale rozkład marginalny $Y$ jest mieszaniną trzech normalnych rozkładów. Zastosowanie Shapiro--Wilka do \emph{połączonego} $Y$ może więc wykryć rzeczywiste odstępstwo od normalności, które nie narusza założenia normalności błędów ANOVA.
\end{przyklad}

\begin{uwaga}[Trzy niezależne decyzje]
\emph{Normalność} dotyczy kształtu rozkładu, \emph{homoskedastyczność} -- równości wariancji warunkowych, a \emph{niezależność} -- wspólnego rozkładu wielu obserwacji. Możliwy jest model normalny i heteroskedastyczny, jak również model nienormalny z równymi wariancjami. Test Shapiro--Wilka nie wykrywa automatycznie zależności w czasie; test Levene'a nie sprawdza niezależności. Dla reszt dopasowanego modelu rozkład statystyki diagnostycznej może różnić się od tego dla \emph{bezpośrednio obserwowanej} próby iid, ponieważ reszty są związane równaniami estymacyjnymi.
\end{uwaga}

\section{Test Shapiro--Wilka: konstrukcja i własności}
\begin{definicja}[Statystyka Shapiro--Wilka]\label{def:shapiro-wilk}
Niech $n\ge3$ i $x_1,\ldots,x_n\in\R$ nie będą wszystkie równe; $x_{(1)}\le\cdots\le x_{(n)}$ oznaczają statystyki porządkowe, $\bar x=n^{-1}\sum_i x_i$. Niech $Z_1,\ldots,Z_n$ będą niezależne o rozkładzie $\Normal(0,1)$, $m_i=\E Z_{(i)}$, $V_{ij}=\Cov(Z_{(i)},Z_{(j)})$, $m=(m_i)_{i=1}^n$, a $V=(V_{ij})_{i,j=1}^n$. Zakładamy odwracalność $V$ i definiujemy wektor współczynników $a=V^{-1}m/\|V^{-1}m\|_2$. Statystyką Shapiro--Wilka nazywamy
\[
 W_n=\frac{\bigl(\sum_{i=1}^n a_i x_{(i)}\bigr)^2}{\sum_{i=1}^n(x_i-\bar x)^2}.
\]
Hipoteza zerowa głosi, że $X_1,\ldots,X_n$ są niezależne i mają ten sam rozkład $\Normal(\mu,\sigma^2)$ z pewnymi $\mu\in\R$ i $\sigma>0$. W praktycznych implementacjach współczynniki i dystrybuantę zerową oblicza się numerycznie lub aproksymacyjnie; nie są one prostymi funkcjami elementarnymi $n$.
\end{definicja}

\begin{twierdzenie}[Niezmienniczość i zakres statystyki]\label{thm:shapiro-properties}
Niech $n\ge3$, $a$ będzie wektorem z definicji~\ref{def:shapiro-wilk}, a $x\in\R^n$ nie będzie stały. Wówczas $\sum_i a_i=0$, $a_{n+1-i}=-a_i$, $0\le W_n\le1$. Dla dowolnych $c\in\R$ i $b\in\R\setminus\{0\}$ zachodzi $W_n(bx_1+c,\ldots,bx_n+c)=W_n(x_1,\ldots,x_n)$.
\end{twierdzenie}
\begin{proof}
Niech $J$ oznacza macierz odwracającą kolejność współrzędnych. Ze względu na symetrię rozkładu normalnego względem zera wektor statystyk porządkowych $(Z_{(1)},\ldots,Z_{(n)})$ ma taki sam rozkład jak $(-Z_{(n)},\ldots,-Z_{(1)})$. Stąd $Jm=-m$ i $JVJ=V$. W konsekwencji $JV^{-1}m=-V^{-1}m$, więc $Ja=-a$ i $\sum_i a_i=0$. Dla wektora $x^\circ=(x_{(i)}-\bar x)_{i=1}^n$, na mocy nierówności Cauchy'ego--Schwarza, $|a^\top x^\circ|^2\le\|a\|_2^2\|x^\circ\|_2^2=\|x^\circ\|_2^2$. Ponieważ $a^\top x^\circ=a^\top x_{()}$, otrzymujemy ograniczenia. Przy $b>0$ uporządkowanie obserwacji pozostaje takie samo; licznik i mianownik mnożą się przez $b^2$, a przesunięcie znika dzięki $\sum_i a_i=0$. Gdy $b<0$, kolejność ulega odwróceniu, a równość $Ja=-a$ daje ten sam wynik.
\end{proof}

\begin{uwaga}[Dlaczego test odrzuca przy małym $W$?]
W statystyce $W$ licznik porównuje uporządkowaną próbę z kierunkiem określonym przez oczekiwane normalne statystyki porządkowe, a mianownik mierzy całkowitą zmienność próby. Mała wartość oznacza, że znacznej części zmienności nie opisuje normalny wzorzec uporządkowania. Kalibracja wartości p wymaga rozkładu $W_n$ pod hipotezą normalności. Wynik $p>\alpha$ oznacza jedynie brak podstaw do odrzucenia; nie dowodzi normalności. Dla bardzo dużych prób test może reagować na nieistotne praktycznie odchylenia, dla bardzo małych mieć niewielką moc. Dla wszystkich jednakowych obserwacji $W$ jest niezdefiniowana (mianownik równy zeru), a nie równa $1$.
\end{uwaga}

\begin{przyklad}[Normalność danych kontra normalność błędów]
Przy porównaniu trzech maszyn testujemy normalność \emph{reszt w obrębie modelu}, a nie połączonych surowych wyników produkcji. Najpierw wykreślamy wykres kwantyl--kwantyl (Q--Q), sprawdzamy wartości odstające, jednostki obserwacji i kształt ogonów; Shapiro--Wilka traktujemy jako dodatkową informację. Mała wartość p nie oznacza automatycznie konieczności zastosowania Kruskala--Wallisa: jeśli badamy \emph{średnie}, możemy rozważyć odporny test Welcha, transformację uzasadnioną modelem lub dobrze skalibrowaną permutację.
\end{przyklad}

\section{Anderson--Darling, Jarque--Bera, Shapiro--Francia i KS}
\begin{definicja}[Statystyka Andersona--Darlinga dla ustalonego rozkładu]\label{def:ad}
Niech $F_0$ będzie w pełni określoną ciągłą dystrybuantą, a $X_1,\ldots,X_n$ niezależną próbą pod hipotezą zerową z $F_0$. Dla $U_{(i)}=F_0(X_{(i)})$ definiujemy, gdy $0<U_{(1)}\le\cdots\le U_{(n)}<1$,
\[
 A_n^2=-n-\frac1n\sum_{i=1}^n(2i-1)\bigl[\log U_{(i)}+\log(1-U_{(n+1-i)})\bigr].
\]
Duże wartości świadczą o niezgodności; logarytmy zwiększają wagę ogonów. Jeśli parametry $F_0$ wyznaczono z tej samej próby, rozkład zerowy zmienia się i trzeba zastosować odpowiednią, osobno skalibrowaną wersję testu.
\end{definicja}

\begin{definicja}[Statystyka Jarque'a--Bery]\label{def:jarque-bera}
Dla obserwacji $x_1,\ldots,x_n$ i $m_r=n^{-1}\sum_i(x_i-\bar x)^r$ ($m_2>0$) niech $g_1=m_3/m_2^{3/2}$ i $g_2=m_4/m_2^2$. Definiujemy $\mathrm{JB}=n[g_1^2/6+(g_2-3)^2/24]$. Przy niezależnej próbie z niedegenerowanego rozkładu normalnego mamy $\mathrm{JB}\Rightarrow\chi^2_2$; jest to kalibracja asymptotyczna, nie dokładna dla małego $n$.
\end{definicja}

\begin{twierdzenie}[Asymptotyczny rozkład JB w modelu normalnym]\label{thm:jb}
Jeżeli $(X_i)_{i\ge1}$ są niezależne o wspólnym rozkładzie $\Normal(\mu,\sigma^2)$, $\sigma>0$, to dla $n\to\infty$ zachodzi $\mathrm{JB}\Rightarrow\chi^2_2$.
\end{twierdzenie}
\begin{proof}
Ze względu na niezmienniczość względem przesunięcia i skali możemy założyć $\mu=0$, $\sigma=1$. Z wielowymiarowego CLT dla empirycznych momentów do rzędu czwartego i metody delty funkcjonały skośności oraz kurtozy mają funkcje wpływu odpowiednio $H_3(z)=z^3-3z$ i $H_4(z)=z^4-6z^2+3$. Można to sprawdzić różniczkując $m_3/m_2^{3/2}$ i $m_4/m_2^2$ po empirycznych momentach surowych w punkcie $(\E Z,\E Z^2,\E Z^3,\E Z^4)=(0,1,0,3)$, z uwzględnieniem centrowania przez $\bar X_n$. W konsekwencji
\[
 \sqrt n\binom{g_1}{g_2-3}=\frac1{\sqrt n}\sum_{i=1}^n\binom{H_3(X_i)}{H_4(X_i)}+o_\Prob(1).
\]
Z momentów standardowego rozkładu normalnego otrzymujemy $\E H_3(Z)^2=6$, $\E H_4(Z)^2=24$ i $\E[H_3(Z)H_4(Z)]=0$. Wielowymiarowe CLT i twierdzenie Slutsky'ego dają granicę $\Normal(0,\diag(6,24))$. Suma kwadratów obu wystandaryzowanych współrzędnych ma więc rozkład $\chi^2_2$.
\end{proof}

\begin{uwaga}[Co rzeczywiście testuje JB?]
Wielkie $\mathrm{JB}$ wskazuje na odchylenie skośności lub kurtozy od wartości normalnych. Nie wykrywa \emph{wszystkich} odstępstw od normalności: istnieją rozkłady nienormalne z taką samą skośnością i kurtozą jak rozkład normalny. Test D'Agostina--Pearsona stosuje inne transformacje miar skośności i kurtozy. Shapiro--Francia opiera się na korelacji uporządkowanych obserwacji z oczekiwanymi normalnymi statystykami porządkowymi, z wagami proporcjonalnymi do wektora ich średnich, zamiast pełnego ważenia macierzą $V^{-1}$. Nie należy mylić tych procedur z testem Kołmogorowa--Smirnowa o \emph{całkowicie znanej} dystrybuancie zerowej; przypadek dopasowywania parametrów omówiliśmy w rozdziale~\ref{ch:nonparametric}.
\end{uwaga}

\section{Równość wariancji: Bartlett, Levene i Brown--Forsythe}
\begin{definicja}[Statystyka Bartletta]\label{def:bartlett}
Niech grupy $j=1,\ldots,k$ zawierają niezależne próby o licznościach $n_j\ge2$. Oznaczmy $N=\sum_j n_j$, przez $S_j^2$ nieobciążoną wariancję w grupie oraz przez $S_p^2=\sum_j(n_j-1)S_j^2/(N-k)$ wariancję łączoną. Jeżeli wszystkie wariancje są dodatnie i $N>k$, definiujemy
\[
 B=\frac{(N-k)\log S_p^2-\sum_{j=1}^k(n_j-1)\log S_j^2}{1+\frac1{3(k-1)}\bigl(\sum_{j=1}^k\frac1{n_j-1}-\frac1{N-k}\bigr)}.
\]
Dla niezależnych prób normalnych, przy wspólnej wariancji pod $H_0$ i ustalonym $k$ oraz $\min_j n_j\to\infty$, $B\Rightarrow\chi^2_{k-1}$. Test może silnie reagować na nienormalność, nawet jeśli wariancje grup są równe.
\end{definicja}

\begin{definicja}[Levene i Brown--Forsythe]\label{def:levene}
Niech $x_{ji}$ oznacza obserwację $i$ w grupie $j$, $n_j\ge2$, $k\ge2$. Dla zadanych centrów $c_j$ definiujemy $Z_{ji}=|x_{ji}-c_j|$, $\bar Z_j=n_j^{-1}\sum_i Z_{ji}$, $\bar Z=N^{-1}\sum_{j,i} Z_{ji}$. Statystyka Levene'a to klasyczna statystyka jednoczynnikowej ANOVA zastosowana do $Z_{ji}$:
\[
 F_L=\frac{\sum_j n_j(\bar Z_j-\bar Z)^2/(k-1)}{\sum_{j,i}(Z_{ji}-\bar Z_j)^2/(N-k)}.
\]
Wersja z $c_j=\bar x_j$ to klasyczny Levene, a z $c_j=\mathrm{Med}(x_{j1},\ldots,x_{jn_j})$ -- wariant Browna--Forsythe'a. W praktyce p-value uzyskuje się przez przybliżenie rozkładem $F_{k-1,N-k}$; przy estymowanych centrach nie jest to ogólnie dokładny skończonopróbkowy rozkład zerowy.
\end{definicja}

\begin{stwierdzenie}[Wspólna skala w rodzinie położenie--skala]\label{prop:levene-scale}
Niech $X_{ji}=\mu_j+\sigma_j\varepsilon_{ji}$, $\sigma_j>0$, przy czym wszystkie $\varepsilon_{ji}$ są niezależne i mają wspólny rozkład symetryczny względem zera, dla którego $\E|\varepsilon_{11}|<\infty$. Wówczas $\E|X_{ji}-\mu_j|=\sigma_j\E|\varepsilon_{11}|$. Jeśli $\varepsilon_{11}$ nie jest zerowa prawie na pewno, równość oczekiwanych odchyleń bezwzględnych od centrów $\mu_j$ jest równoważna równości skal $\sigma_j$.
\end{stwierdzenie}
\begin{proof}
Ze wzoru $X_{ji}-\mu_j=\sigma_j\varepsilon_{ji}$ i dodatniości $\sigma_j$ dostajemy $|X_{ji}-\mu_j|=\sigma_j|\varepsilon_{ji}|$. Biorąc wartości oczekiwane, otrzymujemy pierwszą równość. Ponieważ $\E|\varepsilon_{11}|>0$, druga równoważność wynika z podzielenia przez tę liczbę.
\end{proof}

\begin{uwaga}[Równość wariancji a równość odchyleń bezwzględnych]
Powyższy wniosek jest uzasadniony w \emph{rodzinie wspólnego kształtu i zmiennej skali}. Bez wspólnego kształtu rozkładów równość $\E|X-\mu|$ i równość wariancji nie są równoważne. Levene i Brown--Forsythe są użytecznymi \emph{testami diagnostycznymi}, ale ich interpretacja wymaga sprawdzenia różnic kształtu, skośności i obserwacji odstających. Bartlett jest szczególnie efektywny przy normalności, ale może być zawodny przy grubych ogonach.
\end{uwaga}

\section{Wykresy diagnostyczne, wybór metody i konsekwencje testowania założeń}
\begin{center}
\begin{tikzpicture}[scale=.78,>=Latex]
\draw[->] (0,0)--(5,0) node[right,font=\small]{wartość dopasowana};
\draw[->] (0,-1.8)--(0,1.8) node[above,font=\small]{reszta};
\draw[dashed,gray] (0,0)--(5,0);
\foreach \x/\y in {.5/.10,.9/-.18,1.2/.21,1.6/-.40,2/.53,2.4/-.64,2.7/.32,3.0/-.90,3.3/1.10,3.6/-.50,4/1.28,4.2/-1.31,4.6/1.52}
{\fill[blue!65!black] (\x,\y) circle (1.7pt);}
\draw[red!65!black,densely dashed] (.25,.05)--(4.8,1.55);
\draw[red!65!black,densely dashed] (.25,-.05)--(4.8,-1.55);
\node[align=center,font=\footnotesize] at (2.5,-2.32) {Przykład: rozszerzający się rozrzut reszt\\ sugeruje heteroskedastyczność};
\end{tikzpicture}
\end{center}

\begin{center}\small
\begin{tabular}{>{\raggedright\arraybackslash}p{.29\textwidth}>{\raggedright\arraybackslash}p{.29\textwidth}>{\raggedright\arraybackslash}p{.27\textwidth}}
\toprule
\textbf{Cel} & \textbf{Narzędzie} & \textbf{Co naprawdę wynika?}\\
\midrule
Ocenić normalność modelowych błędów & Q--Q, Shapiro--Wilk, Anderson--Darling & Odstępstwa od rozkładu normalnego; nie test niezależności\\
Sprawdzić skośność i ogony & JB, wykres gęstości i Q--Q & Czułość na określone momenty; JB nie wykrywa wszystkiego\\
Porównać rozrzut w grupach & Brown--Forsythe, Levene; Bartlett przy normalności & Wskazówki o skali; nie zastępują modelu grup\\
Badać autokorelację & Wykres ACF, Ljung--Box, analiza projektu & Zależność w czasie; wymaga właściwej kalibracji\\
Porównać średnie przy nierównych wariancjach & Welch, heteroskedastycznie odporna regresja & Możliwa alternatywa dla klasycznej ANOVA\\
\bottomrule
\end{tabular}
\end{center}

\begin{uwaga}[Dlaczego nie wolno mechanicznie testować, a następnie wybierać?]\label{rem:pretest}
Niech $A$ będzie zdarzeniem ``test normalności nie odrzucił'', a $T$ wybranym po obejrzeniu danych testem średniej. Nominalny poziom osobno skalibrowanego testu $T$ nie daje automatycznej kontroli $\Prob_{H_0}(\text{odrzucenie przez procedurę }T(X)\mid A)$ ani całkowitego poziomu procedury przełączającej testy. Selekcja zależy od tych samych danych, z których liczy się końcową wartość p. Lepiej uzasadnić główną procedurę na etapie planowania, wykorzystać diagnostykę do interpretacji i analizy wrażliwości, a w razie konieczności modelować \emph{całą} procedurę selekcji w symulacji.
\end{uwaga}

\begin{uwaga}[Próba zamiast samej wartości p: przykład praktyczny]\label{rem:diagnostic-workflow}
Dla trzech niezależnych grup przed wykonaniem ANOVA zapisz \emph{jednostkę losowania}, szacowany parametr (np. różnicę średnich) i planowane kontrasty. Następnie zestaw średnie, odchylenia standardowe i liczności grup oraz wykresy reszt względem wartości dopasowanych i Q--Q. Testy Shapiro--Wilka oraz Browna--Forsythe'a można podać jako dodatkową diagnostykę, lecz ich wartości p nie są automatycznym przełącznikiem między ANOVA i Kruskalem--Wallisem. Jeżeli najważniejszy jest efekt w średniej przy nierównych wariancjach, rozważ test Welcha; jeśli badamy różnicę rozkładów lub dominację stochastyczną, pytanie odpowiadające testowi rangowemu jest inne. Przy małej liczbie niezależnych jednostek rozważ czułość wniosku na model i obserwacje skrajne.\end{uwaga}

\section{Laboratorium: normalność i jednorodność wariancji}
\begin{lstlisting}[language=Python,caption={Kompletny program: normalność w grupach, normalność połączonych danych i testy jednorodności wariancji (Python, NumPy, SciPy).}]
import numpy as np
from scipy import stats

rng = np.random.default_rng(2026)
# Same normal noise law, different group means.
groups = [rng.normal(loc=mu, scale=1.0, size=45)
          for mu in (-5.0, 0.0, 5.0)]
pooled = np.concatenate(groups)
residuals = np.concatenate([g - g.mean() for g in groups])
print("Shapiro pooled:", stats.shapiro(pooled))
print("Shapiro residuals (descriptive):", stats.shapiro(residuals))
print("Shapiro within groups:", [stats.shapiro(g).pvalue for g in groups])
print("Brown-Forsythe:", stats.levene(*groups, center="median"))
print("Bartlett:", stats.bartlett(*groups))
print("Jarque-Bera of within-group residuals:", stats.jarque_bera(residuals))

# Heteroscedastic alternative, same expected value.
hetero = [rng.normal(0.0, scale=s, size=45) for s in (1.0, 1.8, 3.0)]
print("Brown-Forsythe heteroscedastic:",
      stats.levene(*hetero, center="median"))
print("Welch ANOVA (SciPy when supported) is different from testing variances")
\end{lstlisting}

\begin{lstlisting}[language=Julia,caption={Kompletny program: diagnostyczne statystyki odchyleń absolutnych od median i wariancje grup (Julia, Random i Statistics).}]
using Random, Statistics
Random.seed!(2026)
groups = [mu .+ randn(45) for mu in (-5.0, 0.0, 5.0)]
all_values = vcat(groups...)
resid = vcat([g .- mean(g) for g in groups]...)
println("pooled variance: ", var(all_values))
println("within-group residual variance: ", var(resid))
for (j, g) in enumerate(groups)
    z = abs.(g .- median(g))
    println("group ", j, ": mean absolute deviation from median = ", mean(z))
end
# For a calibrated Brown-Forsythe p-value use a distribution package
# or the corresponding scipy.stats.levene(..., center="median") routine.
\end{lstlisting}

\begin{zadanie}[Czy Shapiro--Wilk rozstrzyga o użyciu testu Welcha?]
Niech w dwóch niezależnych grupach obserwujemy wyniki o skośnych rozkładach, ale skończonych wariancjach. Obie próby mają po $400$ obserwacji. Wartość p Shapiro--Wilka w obu grupach wynosi mniej niż $0{,}001$. Czy oznacza to, że nie wolno porównywać średnich testem Welcha?\par\emph{Rozwiązanie.} Nie. Odrzucenie dokładnej normalności nie dowodzi zawodności przybliżenia normalnego dla różnicy średnich; o jego adekwatności decydują m.in. ogony, skończoność wariancji, nierównowaga grup i wpływ obserwacji skrajnych. Przy dostatecznych warunkach CLT test Welcha może być sensowny mimo małych wartości p Shapiro--Wilka. Dla ciężkich ogonów należy sprawdzić stabilność, rozważyć odporne oszacowanie średniej albo właściwy bootstrap. Test rangowy odpowiada na inne pytanie, chyba że założono dodatkowy wspólny kształt rozkładów.
\end{zadanie}

\chapter{ANOVA i ANCOVA: modele, interakcje, porównania i zastosowania}
\label{ch:anova-ancova}
\paragraph{Powiązania i odsyłacze.} Podstawą są test F z rozdziału~\ref{ch:parametric-tests} i model liniowy z rozdziału~\ref{ch:linear-regression}; o diagnostyce założeń zobacz rozdział~\ref{ch:diagnostics}.


\section{Jaki problem ma rozwiązać analiza wariancji?}
Jeśli porównujemy średnie trzech technologii, wykonanie wszystkich trzech testów $t$ bez korekty wielokrotności nie kontroluje w zwykły sposób rodzinnego błędu I rodzaju. Jednoczynnikowa ANOVA formułuje jedną hipotezę o wszystkich średnich. Dla dwóch czynników (np. technologia i temperatura) badamy dodatkowo, czy wpływ technologii \emph{zależy} od temperatury. Gdy grupy różnią się także istotną zmienną ilościową (np. wiekiem pacjenta, pomiarem początkowym, obciążeniem urządzenia), rozważamy ANCOVA. Wszystkie te konstrukcje są szczególnymi przypadkami modelu liniowego z rozdziału~\ref{ch:linear-regression}.

\begin{uwaga}[Jednostka niezależna i projekt badania]
Przed wyborem ANOVA należy zapisać, co jest obserwacją, a co jednostką niezależną. Dziesięć pomiarów jednej maszyny nie jest tym samym co dziesięć niezależnych maszyn. Dla zależnych obserwacji stosuje się modele mieszane, analizę powtarzanych pomiarów lub odpowiednio zaprojektowane testy permutacyjne. ANOVA nie usuwa problemu pseudoreplikacji. Jeśli wartości czynników nie były randomizowane, istotny statystycznie współczynnik nie dowodzi skutku przyczynowego.
\end{uwaga}

\section{Jednoczynnikowa ANOVA w języku rzutów}
Rozważmy $k\ge2$ grup z licznościami $n_j\ge1$, $n=\sum_jn_j>k$. Model zakłada $Y_{ji}=\mu_j+\varepsilon_{ji}$, gdzie $\varepsilon_{ji}$ są niezależne o wspólnym rozkładzie $\Normal(0,\sigma^2)$, $\sigma^2>0$. Hipoteza zerowa ma postać $H_0:\mu_1=\cdots=\mu_k$. Niech $X\in\R^{n\times k}$ będzie macierzą wskaźników grup, $P_X=X(X^\top X)^{-1}X^\top$, a $P_1=n^{-1}\mathbf1\mathbf1^\top$. Wtedy
\[
 SS_B=\|P_XY-P_1Y\|_2^2,\qquad
 SS_W=\|(I-P_X)Y\|_2^2,\qquad
 F=\frac{SS_B/(k-1)}{SS_W/(n-k)}.
\]
Rozdział~\ref{ch:parametric-tests} zawiera elementarny dowód odpowiadający sumom kwadratów. Następujące twierdzenie uogólnia go do dowolnych dwóch zagnieżdżonych podprzestrzeni modelowych i będzie podstawą ANCOVA.

\begin{twierdzenie}[Dokładny test $F$ dla modeli zagnieżdżonych]\label{thm:nested-F-anova}
Niech $Y\sim\Normal_n(\mu,\sigma^2 I_n)$, $\sigma^2>0$. Niech $M_R\subset M_F\subset\R^n$ będą ustalonymi podprzestrzeniami liniowymi o wymiarach $p_R<p_F<n$, a $P_R,P_F$ rzutami ortogonalnymi na te przestrzenie. Załóżmy $H_0:\mu\in M_R$. Wtedy przy $q=p_F-p_R$ zmienna
\[
 F=\frac{(\|(I-P_R)Y\|_2^2-\|(I-P_F)Y\|_2^2)/q}{\|(I-P_F)Y\|_2^2/(n-p_F)}
\]
ma dokładny rozkład Fishera $F_{q,n-p_F}$.
\end{twierdzenie}
\begin{proof}
Ponieważ $M_R\subset M_F$, zachodzi $P_FP_R=P_RP_F=P_R$. Stąd $Q=P_F-P_R$ jest symetryczną projekcją ortogonalną o rzędzie $q$, a $H=I-P_F$ jest projekcją o rzędzie $n-p_F$; ponadto $QH=HQ=0$. Z twierdzenia Pitagorasa $\|(I-P_R)Y\|^2-\|(I-P_F)Y\|^2=\|QY\|^2$. Pod $H_0$ mamy $Q\mu=0$ i $H\mu=0$, a więc obie formy są kwadratami norm rzutów $\varepsilon/\sigma$ na wzajemnie ortogonalne podprzestrzenie o wymiarach $q$ i $n-p_F$. Ortogonalna niezmienniczość standardowego wektora gaussowskiego daje niezależność i rozkłady odpowiednio $\chi_q^2$ oraz $\chi_{n-p_F}^2$. Ich stosunek po podzieleniu przez stopnie swobody ma rozkład $F_{q,n-p_F}$.
\end{proof}

\begin{uwaga}[Kiedy test $F$ nie jest dokładny?]
Powyższy wynik \emph{nie} obejmuje błędów o różnych wariancjach, obserwacji zależnych ani modelu dobranego adaptacyjnie na podstawie tej samej próby bez odpowiedniej korekty selekcji. Przy heteroskedastyczności można rozważyć regresję z odporną macierzą kowariancji i test Walda lub ANOVA Welcha. W dużych próbach wyniki asymptotyczne często pozostają użyteczne bez dokładnej normalności, ale należy osobno uzasadnić ich warunki.
\end{uwaga}

\section{Dwuczynnikowa ANOVA i interakcje}
\begin{definicja}[Model dwuczynnikowy z interakcją]\label{def:two-way-anova}
Dla poziomów czynnika $A\in\{1,\ldots,a\}$ i $B\in\{1,\ldots,b\}$ oraz obserwacji $Y_{ijr}$, $1\le r\le n_{ij}$, rozważamy model
\[
 Y_{ijr}=\mu+\alpha_i+\beta_j+\gamma_{ij}+\varepsilon_{ijr}.
\]
Parametry identyfikujemy np. przez warunki $\sum_i\alpha_i=\sum_j\beta_j=0$ oraz $\sum_i\gamma_{ij}=0$ dla każdego $j$ i $\sum_j\gamma_{ij}=0$ dla każdego $i$. Zakładamy, że $n_{ij}\ge1$ dla wszystkich $1\le i\le a$, $1\le j\le b$, oraz że błędy są niezależne, mają zerowe wartości oczekiwane i wspólną dodatnią wariancję; normalność błędów jest dodatkowym warunkiem dokładnej kalibracji testów $F$. \emph{Brak interakcji} oznacza $\gamma_{ij}=0$ dla wszystkich $i,j$ w tej parametryzacji. Hipotezę testujemy przez porównanie modelu pełnego z modelem addytywnym.
\end{definicja}

\begin{stwierdzenie}[Interakcja to różnica różnic]\label{prop:interaction-did}
Niech dwa czynniki mają po dwa poziomy i oznaczmy średnie komórkowe przez $m_{ij}=\E(Y\mid A=i,B=j)$, $i,j\in\{0,1\}$. W modelu bez interakcji zachodzi $m_{11}-m_{10}-m_{01}+m_{00}=0$. Odwrotnie, jeśli wszystkie cztery średnie istnieją i zachodzi ta równość, można je zapisać w postaci addytywnej $m_{ij}=u+a_i+b_j$.
\end{stwierdzenie}
\begin{proof}
Dla addytywnego modelu $m_{11}-m_{10}-m_{01}+m_{00}$ jest zerem po skróceniu składników. Odwrotnie, wybierzmy $u=m_{00}$, $a_0=b_0=0$, $a_1=m_{10}-m_{00}$, $b_1=m_{01}-m_{00}$. Wtedy $u+a_1+b_1=m_{10}+m_{01}-m_{00}=m_{11}$ na mocy założenia. Pozostałe trzy równości są automatyczne.
\end{proof}

\begin{przyklad}[Dlaczego interakcja ma znaczenie]\label{ex:interaction-practical}
Dla dwóch technologii $A=0,1$ i dwóch warunków pracy $B=0,1$ średnie wynoszą odpowiednio $m_{00}=10$, $m_{10}=14$, $m_{01}=12$, $m_{11}=12$. W warunku $B=0$ zmiana technologii zwiększa średnią o $4$, a w warunku $B=1$ nie zmienia jej wcale. Interakcja ma wartość $12-14-12+10=-4$. Raportowanie wyłącznie \emph{jednego} ``efektu technologii'' ukrywa tę zależność. Należy podać oszacowane średnie i kontrasty \emph{w obrębie warunków}, z przedziałami ufności i korektą na wielokrotność, jeśli testujemy kilka kontrastów.
\end{przyklad}

\begin{center}
\begin{tikzpicture}[x=1.65cm,y=.29cm,>=Latex]
\draw[->] (-.2,8.5)--(2.4,8.5) node[right,font=\small] {poziom $B$};
\draw[->] (0,8)--(0,16) node[above,font=\small] {średnia};
\foreach \v in {10,12,14} {\draw (-.06,\v)--(.06,\v) node[left=4pt,font=\scriptsize] {\v};}
\draw[blue!75!black,thick] (.5,10)--(1.7,12);
\draw[red!70!black,thick] (.5,14)--(1.7,12);
\foreach \x/\y in {.5/10,1.7/12} {\fill[blue!75!black] (\x,\y) circle (2.5pt);}
\foreach \x/\y in {.5/14,1.7/12} {\draw[red!70!black,thick,fill=white] (\x,\y) circle (2.7pt);}
\node[below,font=\small] at (.5,8.5) {$B=0$};
\node[below,font=\small] at (1.7,8.5) {$B=1$};
\node[font=\scriptsize,blue!75!black,anchor=west] at (1.8,14.5) {$A=0$};
\node[font=\scriptsize,red!70!black,anchor=west] at (1.8,13.1) {$A=1$};
\end{tikzpicture}

{\small\emph{Rysunek 28.1.} Nierównoległe profile średnich ilustrują interakcję czynników.}
\end{center}

\section{Niezrównoważone plany i sumy kwadratów typów I, II i III}
Przy równych licznościach każdej komórki i pełnym planie efekty główne i interakcje tworzą często ortogonalne podprzestrzenie modelu, a podział sum kwadratów jest jednoznaczny. Przy \emph{niezrównoważonych} danych składniki na ogół nie są ortogonalne i różne hipotezy odpowiadają różnym sumom kwadratów.

\begin{definicja}[Przyrostowa suma kwadratów dla dwóch modeli]\label{def:partial-ss}
Dla ustalonych, zagnieżdżonych macierzy projektu $X_R,X_F$ o pełnych rangach $p_R<p_F$ definiujemy $SS(X_F\mid X_R)=RSS(X_R)-RSS(X_F)$, gdzie $RSS(X)=\min_\beta\|Y-X\beta\|_2^2$. \emph{Typ I} testuje kolejne przyrosty zgodnie z kolejnością wprowadzania składników do modelu. \emph{Typ II} testuje efekt główny po uwzględnieniu pozostałych efektów głównych, zwykle nie dostosowując do wyższych interakcji zawierających ten efekt. \emph{Typ III} próbuje testować efekt po uwzględnieniu wszystkich pozostałych składników danego modelu; przy interakcjach interpretacja zależy m.in. od kodowania, ograniczeń i estymowalności kontrastów.
\end{definicja}

\begin{uwaga}[Który typ ANOVA wybrać?]
Najpierw sformułuj \emph{konkretny kontrast lub porównanie modeli}, a dopiero potem wybieraj typ sum kwadratów. Przy silnej interakcji często więcej sensu mają proste efekty warunkowe niż pojedynczy test ``efektu głównego''. Przy pustych komórkach niektóre kontrasty nie są estymowalne; procedura programowa nie tworzy informacji, której nie ma w danych. W kodzie z testami typu III trzeba świadomie dobrać kontrasty (np. sum-to-zero), a także sprawdzić, co faktycznie testowana hipoteza oznacza w tej parametryzacji.
\end{uwaga}

\section{ANCOVA: wspólne nachylenie, interakcja z kowariantą i wartości skorygowane}
\begin{definicja}[Model ANCOVA z dwiema grupami]\label{def:ancova}
Dla $i=1,\ldots,n$ obserwujemy zmienną wynikową $Y_i\in\R$, znacznik grupy $G_i\in\{0,1\}$ oraz kowariantę $X_i\in\R$. Dla ustalonej wartości odniesienia $x_0$ przyjmujemy model
\[
 Y_i=\beta_0+\tau G_i+\beta_1(X_i-x_0)+\delta G_i(X_i-x_0)+\varepsilon_i.
\]
Przy $\delta=0$ \emph{wspólne nachylenie} względem $X$ wynosi $\beta_1$, a $\tau$ jest różnicą średnich warunkowych między grupami przy \emph{tej samej} wartości kowarianty $X=x_0$. Przy $\delta\ne0$ różnica wynosi $\tau+\delta(x-x_0)$ i zależy od $x$. Zakładamy pełną rangę macierzy projektu do estymacji wszystkich zadanych współczynników.
\end{definicja}

\begin{stwierdzenie}[Różnica średnich warunkowych w ANCOVA]\label{prop:ancova-contrast}
W modelu~\ref{def:ancova}, przy $\E(\varepsilon_i\mid G_i,X_i)=0$ dla wszystkich $i$, zachodzi
$\E(Y_i\mid G_i=1,X_i=x)-\E(Y_i\mid G_i=0,X_i=x)=\tau+\delta(x-x_0)$
wszędzie tam, gdzie odpowiednie wartości warunkowe są określone. W szczególności dla $\delta=0$ różnica jest stała względem $x$.
\end{stwierdzenie}
\begin{proof}
Wystarczy podstawić $G_i=1$ i $G_i=0$ do modelu, wziąć warunkowe wartości oczekiwane i odjąć, korzystając z założenia o średniej błędu.
\end{proof}

\begin{uwaga}[Brak nakładania zakresów kowarianty]
Jeżeli grupa kontrolna ma $X\in[20,30]$, a grupa badana $X\in[60,70]$, to porównanie \emph{przy wspólnym} $x_0=45$ jest ekstrapolacją i może być niemiarodajne. Potrzebujemy wspólnego zakresu lub dodatkowego, uzasadnionego modelu ekstrapolacji. Korygowanie względem kowarianty zmierzonej \emph{po} zastosowaniu interwencji może ponadto zmienić interpretowany efekt, a nawet wprowadzić błąd przez kontrolowanie mediatora lub zmiennej kolizyjnej. W badaniach obserwacyjnych ANCOVA nie usuwa automatycznie wszystkich czynników zakłócających.
\end{uwaga}

\begin{przyklad}[ANCOVA kontra porównanie surowych średnich]
Badamy dwa sposoby nauczania, a $X$ oznacza wynik testu wstępnego i $Y$ wynik końcowy. Jeśli jeden sposób częściej trafia do uczniów z wysokimi wynikami początkowymi, różnica surowych średnich $Y$ może odzwierciedlać różnice w $X$, a nie efekt sposobu nauczania. ANCOVA porównuje warunkowe średnie $Y$ dla tej samej wartości $X$. Trzeba jednak sprawdzić zgodność zakresów $X$ w grupach, możliwą interakcję $G\times X$, liniowość zależności i projekt doboru uczniów.
\end{przyklad}


\begin{twierdzenie}[Kiedy korekta o kowariantę zmniejsza wariancję]\label{thm:ancova-variance}
Niech $n_0,n_1\ge1$. W grupie $g\in\{0,1\}$ obserwujemy niezależne pary $(X_{gi},Y_{gi})$, $1\le i\le n_g$, przy czym
$Y_{gi}=\mu+\tau g+\beta X_{gi}+\varepsilon_{gi}$.
Załóżmy, że wszystkie $X_{gi}$ są niezależne o jednakowym rozkładzie, $\E X_{gi}=0$, $\Var(X_{gi})=v<\infty$, wszystkie $\varepsilon_{gi}$ są niezależne o średniej zero i wariancji $\sigma^2<\infty$, a rodziny $X$ i $\varepsilon$ są niezależne. Niech $D=\bar Y_1-\bar Y_0$ oznacza surową różnicę średnich, a $D_{\rm adj}=D-\beta(\bar X_1-\bar X_0)$ korektę przy \emph{znanym} współczynniku $\beta$. Wówczas
\[
 \Var(D)=(\beta^2v+\sigma^2)\left(\frac1{n_1}+\frac1{n_0}\right),\qquad
 \Var(D_{\rm adj})=\sigma^2\left(\frac1{n_1}+\frac1{n_0}\right).
\]
W szczególności korekta zmniejsza wariancję, gdy $\beta^2v>0$. Wniosek zakłada znane nachylenie i nie daje automatycznej gwarancji takiej poprawy dla każdego estymatora nachylenia w małej próbie.
\end{twierdzenie}
\begin{proof}
Po uśrednieniu równania modelu dla obu grup otrzymujemy
$D-\tau=\beta(\bar X_1-\bar X_0)+(\bar\varepsilon_1-\bar\varepsilon_0)$.
Oba składniki po prawej stronie są niezależne i mają wariancje odpowiednio
$\beta^2v(1/n_1+1/n_0)$ oraz $\sigma^2(1/n_1+1/n_0)$.
Odjęcie pierwszego z nich od $D$ pozostawia tylko składnik błędów, co daje oba wzory.
\end{proof}

\begin{uwaga}[Gdy $\beta$ trzeba oszacować]
W praktyce $\beta$ nie jest znane. Jeśli $\widehat\beta-\beta=o_\Prob(1)$, a $n_0/n$ i $n_1/n$ są oddalone od zera przy $n=n_0+n_1\to\infty$, to $\bar X_1-\bar X_0=O_\Prob(n^{-1/2})$ i dodatkowy błąd korekty
$(\widehat\beta-\beta)(\bar X_1-\bar X_0)=o_\Prob(n^{-1/2})$.
Dlatego przy odpowiedniej regularności estymacja nachylenia nie zmienia głównego składnika wariancji asymptotycznej. Ten argument opiera się na wspólnym modelu dla obu grup; nie usprawiedliwia dowolnego doboru kowariant po obejrzeniu wyników.
\end{uwaga}

\begin{przyklad}[Paradoks Simpsona: odwrócenie porównania po agregacji]
W dwóch grupach porównujemy częstość sukcesu, przy czym pacjenci należą do grupy niskiego albo wysokiego ryzyka. W grupie niskiego ryzyka sukces obserwowano u $80$ z $100$ leczonych i u $9$ z $10$ kontrolnych; w grupie wysokiego ryzyka u $1$ z $10$ leczonych i u $20$ ze $100$ kontrolnych. Leczenie ma niższy odsetek sukcesów \emph{w obu warstwach}: $80\%<90\%$ oraz $10\%<20\%$. Po złączeniu warstw leczonych sukces ma $81$ na $110$ osób, a kontrolnych tylko $29$ na $110$: porządek odsetków \emph{odwraca się}. Przyczyną jest bardzo różny udział warstw ryzyka, a nie sprzeczność rachunkowa. Jest to przestroga przed bezrefleksyjnym raportowaniem surowych średnich, ale samo warstwowanie nie zapewnia wnioskowania przyczynowego, jeśli dobór do grup jest obciążony.
\end{przyklad}

\section{Porównania post hoc i wielkość efektu}
\begin{definicja}[Wielkości efektu w ANOVA]\label{def:anova-effectsize}
Dla jednoczynnikowej ANOVA z $k$ grupami i $n>k$, niech $SS_B$ i $SS_W$ oznaczają sumy kwadratów między grupami i wewnątrz nich. Definiujemy $SS_T=SS_B+SS_W$, $MS_W=SS_W/(n-k)$, $\eta^2=SS_B/SS_T$ przy $SS_T>0$ oraz
\[
 \omega^2=\frac{SS_B-(k-1)MS_W}{SS_T+MS_W}
\]
przy dodatnim mianowniku. Estymator $\omega^2$ może być ujemny dla danych o bardzo małym efekcie; nie należy bez komentarza ukrywać tego przez obcięcie do zera. Dla porównania \emph{dwóch zagnieżdżonych modeli} definiujemy $\eta_{\mathrm{partial}}^2=(RSS_R-RSS_F)/RSS_R$, gdy $RSS_R>0$.
\end{definicja}

\begin{uwaga}[Wielkość efektu nie zastępuje przedziału ufności]
Duża statystyka $F$ mówi o stosunku sygnału do szumu, nie o praktycznej wielkości zmiany w jednostkach wyniku. Dla pojedynczych różnic średnich podawaj przedziały w oryginalnej skali. Po wykryciu różnic między $k$ grupami rozważ planowane kontrasty i kontrolę wielokrotności: \emph{Tukey HSD} dla wszystkich par w klasycznym normalnym modelu o wspólnej wariancji (prosta wersja dla jednakowych liczności), \emph{Games--Howell} przy nierównych wariancjach dla porównań par, \emph{Dunnett} przy porównywaniu każdej terapii wyłącznie z kontrolą; lub Holm dla osobno poprawnie skalibrowanych wartości p. Nie twierdź, że globalna istotna ANOVA dowodzi różnicy \emph{każdej} pary grup.
\end{uwaga}

\begin{center}\small
\begin{tabular}{>{\raggedright\arraybackslash}p{.28\textwidth}>{\raggedright\arraybackslash}p{.27\textwidth}>{\raggedright\arraybackslash}p{.31\textwidth}}
\toprule
\textbf{Pytanie badawcze} & \textbf{Metoda do rozważenia} & \textbf{Sprawdź i zaraportuj}\\
\midrule
Więcej niż dwie niezależne grupy, wspólna wariancja & Jednoczynnikowa ANOVA & Średnie, reszty, wielkości efektu, kontrasty\\
Nierówne wariancje, różne liczności & ANOVA Welcha & Zróżnicowanie wariancji, przedziały, Games--Howell\\
Dwa czynniki & Dwuczynnikowa ANOVA & Interakcja, średnie komórkowe, niezrównoważenie\\
Grupa i kowarianta ilościowa & ANCOVA & Nakładanie $X$, nachylenie, moment pomiaru kowarianty\\
Pomiar wielokrotny w tej samej jednostce & Model mieszany / ANOVA powtarzanych pomiarów & Kowariancja w obrębie jednostki\\
Kilka współzależnych wyników & MANOVA / modele wielowymiarowe & Wspólna macierz kowariancji i sens łącznej hipotezy\\
\bottomrule
\end{tabular}
\end{center}

\section{Laboratorium: modele zagnieżdżone, ANCOVA i interakcje}
\begin{lstlisting}[language=Python,caption={Kompletny program: dwuczynnikowa ANOVA, ANCOVA, interakcja i testy zagnieżdżone (Python, NumPy, pandas, statsmodels).}]
import numpy as np
import pandas as pd
import statsmodels.api as sm
from statsmodels.formula.api import ols

rng = np.random.default_rng(2026)
rows = []
for a in ("A0", "A1"):
    for b in ("B0", "B1"):
        for _ in range(32):
            group = int(a == "A1")
            condition = int(b == "B1")
            baseline = rng.normal(50, 9)
            # Difference A1-A0 falls from +4 to 0 across B levels.
            y = (10 + 2*condition + 4*group
                 - 4*group*condition + 0.35*(baseline-50)
                 + rng.normal(0, 2))
            rows.append((a, b, baseline, y))
df = pd.DataFrame(rows, columns=["A", "B", "baseline", "outcome"])

# Two-way ANOVA with interaction, also adjusting for baseline.
m = ols("outcome ~ C(A, Sum)*C(B, Sum) + baseline", data=df).fit()
print(sm.stats.anova_lm(m, typ=2))

# ANCOVA: do slopes vary between groups?
m_common = ols("outcome ~ C(A) + baseline", data=df).fit()
m_slopes = ols("outcome ~ C(A)*baseline", data=df).fit()
print("ANCOVA slope-interaction nested F test:")
print(sm.stats.anova_lm(m_common, m_slopes))

# Factor A effect at a chosen condition, not global marginal meaning.
for b in ("B0", "B1"):
    predictions = pd.DataFrame({"A": ["A0", "A1"],
                                "B": [b, b], "baseline": [50, 50]})
    print("adjusted A1-A0 at", b,
          np.diff(m.predict(predictions))[0])

# Examine robust covariance inference if heteroscedasticity is suspected.
print("HC3 standard errors:", m.get_robustcov_results(cov_type="HC3").bse)
\end{lstlisting}

\begin{lstlisting}[language=Julia,caption={Kompletny program: dokładny rachunek zagnieżdżonej statystyki F metodą najmniejszych kwadratów (Julia, Random, Statistics i LinearAlgebra).}]
using Random, Statistics, LinearAlgebra
Random.seed!(2026)
n = 80
x = randn(n)
g = vcat(zeros(Int, div(n, 2)), ones(Int, div(n, 2)))
y = 3 .+ 2 .* g .+ 0.7 .* x .+ randn(n)
Xr = hcat(ones(n), x)
Xf = hcat(ones(n), x, g)
betar = Xr \ y
betaf = Xf \ y
rssr = sum(abs2, y - Xr * betar)
rssf = sum(abs2, y - Xf * betaf)
q = size(Xf, 2) - size(Xr, 2)
F = ((rssr - rssf) / q) / (rssf / (n - size(Xf, 2)))
println("nested-model F statistic: ", F)
println("adjusted group difference: ", betaf[end])
# For an exact p-value add the F-distribution CDF from Distributions.jl.
\end{lstlisting}

\begin{zadanie}[Czy nieistotna interakcja dowodzi wspólnego nachylenia?]
W modelu ANCOVA test hipotezy $H_0:\delta=0$ ma wartość p równą $0{,}28$. Czy wolno twierdzić, że nachylenia na pewno są jednakowe?\par\emph{Rozwiązanie.} Nie. Brak odrzucenia $H_0$ nie dowodzi jej prawdziwości. Przy małych próbach test interakcji może mieć niską moc. Należy oszacować $\delta$ i przedział ufności, rozważyć sens praktyczny możliwych różnic nachyleń oraz preferowany model uzasadnić teorią i projektem badania. Mechaniczne testowanie interakcji, po którym raportuje się wyłącznie wynik modelu wybranego na tych samych danych, może zniekształcić nominalną niepewność.
\end{zadanie}


\chapter{MANOVA i MANCOVA: wnioskowanie o wielu wynikach jednocześnie}
\label{ch:manova-mancova}
\paragraph{Powiązania i odsyłacze.} Teorię modeli liniowych i ANOVA znajdziesz w rozdziałach~\ref{ch:linear-regression} i~\ref{ch:anova-ancova}; dobór testu omawia rozdział~\ref{ch:stat-laboratory}.


\section{Motywacja: kilka wyników nie oznacza kilku niezależnych eksperymentów}
Załóżmy, że porównujemy trzy sposoby rehabilitacji. U każdego pacjenta mierzymy czas chodu, siłę mięśni i wynik skali funkcjonalnej. Oddzielne testowanie każdej zmiennej może nie wykryć zmiany ich \emph{wspólnego} profilu; może także prowadzić do problemu wielokrotnego testowania. \emph{Wielowymiarowa analiza wariancji} (MANOVA) testuje hipotezę dotyczącą wektora średnich, uwzględniając kowariancje między wynikami. Gdy dodatkowo uwzględniamy kowarianty, otrzymujemy \emph{wielowymiarową analizę kowariancji} (MANCOVA).

Nie należy mylić wielu \emph{zmiennych wynikowych u tej samej osoby} z wieloma \emph{pomiarami tej samej cechy w kolejnych chwilach}. W szczególności odrzucenie globalnej hipotezy MANOVA nie oznacza, że różnią się wszystkie poszczególne wyniki; do wskazania konkretnych zmiennych lub kontrastów potrzebna jest osobna, odpowiednio zaplanowana analiza. W obu sytuacjach występuje korelacja, lecz w tej drugiej zwykle ważne są porządek czasu, trajektoria i brakujące wizyty; zajmiemy się tym w rozdziale~\ref{ch:repeated-measures}.

\begin{center}
\begin{tikzpicture}[x=1.12cm,y=.82cm,>=Latex]
\draw[->] (0,0)--(7.5,0) node[right,font=\small] {wynik 1};
\draw[->] (0,0)--(0,4.6) node[above,font=\small] {wynik 2};
\draw[blue!70!black,thick,rotate around={27:(2.7,2.15)}] (2.7,2.15) ellipse (1.5 and .48);
\draw[red!65!black,thick,rotate around={27:(4.5,2.95)}] (4.5,2.95) ellipse (1.5 and .48);
\fill[blue!75!black] (2.7,2.15) circle (2.8pt);
\fill[red!65!black] (4.5,2.95) circle (2.8pt);
\draw[->,thick] (2.7,2.15)--(4.5,2.95);
\node[below,font=\footnotesize,blue!70!black] at (2.7,1.98) {grupa A};
\node[above,font=\footnotesize,red!65!black] at (4.65,3.15) {grupa B};
\node[align=center,font=\footnotesize] at (5.9,1.1) {kierunek różnicy\\średnich};
\end{tikzpicture}

{\small\emph{Rysunek 29.1.} Schemat: MANOVA porównuje wektory średnich z uwzględnieniem kierunków współzmienności wyników; elipsy są poglądowe, nie oznaczają konkretnych przedziałów ufności.}
\end{center}

\begin{definicja}[Wielowymiarowy model liniowy]\label{def:multivariate-linear}
Niech $n,p,r\ge1$, $X\in\R^{n\times r}$ będzie ustaloną macierzą pełnej rangi $r<n$, $B\in\R^{r\times p}$ macierzą nieznanych współczynników, a $\Sigma\in\R^{p\times p}$ nieznaną dodatnio określoną macierzą. Obserwujemy $Y\in\R^{n\times p}$ w modelu
\[
 Y=XB+U,\qquad U_{i\cdot}\ \text{są niezależne o rozkładzie }\Normal_p(0,\Sigma),\quad 1\le i\le n.
\]
Symbol $U_{i\cdot}$ oznacza $i$-ty wiersz macierzy $U$ traktowany jako wektor w $\R^p$; $\Normal_p(0,\Sigma)$ jest $p$-wymiarowym rozkładem normalnym o zerowym wektorze średniej i macierzy kowariancji $\Sigma$. Niezależność wierszy dotyczy \emph{jednostek}, natomiast współrzędne jednego wiersza mogą być skorelowane.
\end{definicja}

Przy kodowaniu grup kolumnami wskaźnikowymi jest to MANOVA; po dodaniu kolumn odpowiadających kowariantom -- MANCOVA. Jeżeli kowarianty są losowe, wszystkie poniższe dokładne rozkłady należy rozumieć \emph{warunkowo względem macierzy $X$}, przy założeniu o wskazanym rozkładzie błędów warunkowych.

\begin{stwierdzenie}[Estymator i macierz reszt]\label{prop:multiv-ols}
W modelu~\ref{def:multivariate-linear} estymator metodą najmniejszych kwadratów i macierz sum iloczynów reszt mają postać
\[
 \widehat B=(X^\top X)^{-1}X^\top Y,\qquad E=Y^\top(I-P_X)Y,\quad P_X=X(X^\top X)^{-1}X^\top.
\]
Ponadto $\E\widehat B=B$ i $\Cov(\widehat B_{j\cdot},\widehat B_{l\cdot})=((X^\top X)^{-1})_{jl}\Sigma$.
\end{stwierdzenie}
\begin{proof}
Minimalizujemy $\operatorname{tr}((Y-XC)^\top(Y-XC))$ po $C\in\R^{r\times p}$. Równania normalne to $X^\top XC=X^\top Y$; macierz $X^\top X$ jest odwracalna. Wstawienie $Y=XB+U$ daje $\widehat B-B=(X^\top X)^{-1}X^\top U$. Kowariancje wierszy $U$ są zerowe dla różnych jednostek i równe $\Sigma$ w obrębie wiersza, więc po rozwinięciu sum otrzymujemy tezę o kowariancji. Reszty wynoszą $(I-P_X)Y$.
\end{proof}

\section{Macierze hipotezy i błędu: ścisła konstrukcja}
\begin{definicja}[Hipoteza wielowymiarowa i macierze $H,E$]\label{def:manova-he}
Niech $M_R\subsetneq M_F\subset\R^n$ będą \emph{ustalonymi} podprzestrzeniami wymiarów odpowiednio $r_R<r_F<n$. Niech $P_R,P_F$ będą rzutami ortogonalnymi na te podprzestrzenie i oznaczmy $q=r_F-r_R$, $\nu=n-r_F$. Rozważamy model $Y=\mu+U$, w którym $\mu\in\R^{n\times p}$ jest nieznaną macierzą wartości oczekiwanych, każdy słupek $\mu$ należy do $M_F$, a wiersze $U$ są niezależne $\Normal_p(0,\Sigma)$. Testujemy hipotezę $H_0:\ \text{każdy słupek }\mu\text{ należy do }M_R$. Definiujemy
\[
 H=Y^\top(P_F-P_R)Y,\qquad E=Y^\top(I-P_F)Y.
\]
Są to odpowiednio \emph{macierz sum iloczynów hipotezy} i \emph{macierz sum iloczynów błędu}. Obie są symetryczne i nieujemnie określone. Dalej zakładamy $\nu\ge p$, aby $E$ była dodatnio określona niemal na pewno w modelu normalnym.
\end{definicja}

\begin{definicja}[Rozkład Wisharta]\label{def:wishart-manova}
Dla $p,m\ge1$ i dodatnio określonej $\Sigma\in\R^{p\times p}$ rozkładem $W_p(\Sigma,m)$ nazywamy rozkład macierzy $\sum_{i=1}^m Z_iZ_i^\top$, gdzie $Z_i$ są niezależne o rozkładzie $\Normal_p(0,\Sigma)$. Macierz taka ma wartość oczekiwaną $m\Sigma$ i jest dodatnio określona niemal na pewno, jeżeli $m\ge p$.
\end{definicja}

\begin{twierdzenie}[Dokładny rozkład macierzy $H,E$ pod hipotezą zerową]\label{thm:wishart-he}
W warunkach definicji~\ref{def:manova-he}, pod $H_0$ macierze $H$ i $E$ są niezależne oraz
\[
 H\sim W_p(\Sigma,q),\qquad E\sim W_p(\Sigma,\nu).
\]
\end{twierdzenie}
\begin{proof}
Macierze $Q=P_F-P_R$ i $R=I-P_F$ są rzutami ortogonalnymi o rzędach $q$ i $\nu$ oraz $QR=0$. Wybierzmy ortonormalne bazy ich obrazów i dopełnijmy je do bazy $\R^n$. Lewostronne pomnożenie $U$ przez otrzymaną macierz ortogonalną nie zmienia jego rozkładu: po transformacji wiersze wciąż są niezależne o rozkładzie $\Normal_p(0,\Sigma)$. Pod $H_0$ mamy $Q\mu=R\mu=0$. Zatem $Y^\top QY$ i $Y^\top RY$ są sumami iloczynów odpowiednio $q$ i $\nu$ \emph{rozłącznych} podrodzin niezależnych normalnych wierszy. Definicja~\ref{def:wishart-manova} daje oba rozkłady i niezależność.
\end{proof}

\begin{uwaga}[Co znaczy dokładność?]
Twierdzenie dotyczy stałego, z góry określonego projektu i wspólnej macierzy kowariancji we wszystkich jednostkach. Nie obejmuje niezależności pozornej (np. dwóch wizyt jednego pacjenta), różnych $\Sigma$ w grupach ani selekcji wyników i kowariant po obejrzeniu danych. Przy liczbie wyników bliskiej liczbie niezależnych jednostek macierz $E$ może być źle uwarunkowana nawet wtedy, gdy formalnie jest odwracalna.
\end{uwaga}

\section{Cztery klasyczne statystyki wielowymiarowe}
\begin{definicja}[Pierwiastki charakterystyczne MANOVA]\label{def:manova-roots}
Niech $H,E\in\R^{p\times p}$ będą symetryczne nieujemnie określone, a $E$ dodatnio określona. Niech $\lambda_1\ge\cdots\ge\lambda_p\ge0$ będą wartościami własnymi macierzy symetrycznej $E^{-1/2}HE^{-1/2}$. Definiujemy
\[
 \Lambda_{\mathrm W}=\frac{\det E}{\det(E+H)},\quad
 V_{\mathrm P}=\operatorname{tr}((E+H)^{-1}H),\quad
 T_{\mathrm{HL}}=\operatorname{tr}(E^{-1}H),\quad
 \Theta_{\mathrm R}=\lambda_1.
\]
Nazywamy je odpowiednio \emph{lambdą Wilksa}, \emph{śladem Pillaia}, \emph{śladem Hotellinga--Lawleya} i \emph{największym pierwiastkiem Roya}. Zdarza się także konwencja Roya $\lambda_1/(1+\lambda_1)$, dlatego zawsze trzeba podać zastosowaną definicję.
\end{definicja}

\begin{twierdzenie}[Postać spektralna i niezmienniczość]\label{thm:manova-spectra}
Dla macierzy z definicji~\ref{def:manova-roots} zachodzi
\[
 \Lambda_{\mathrm W}=\prod_{j=1}^p(1+\lambda_j)^{-1},\quad
 V_{\mathrm P}=\sum_{j=1}^p\frac{\lambda_j}{1+\lambda_j},\quad
 T_{\mathrm{HL}}=\sum_{j=1}^p\lambda_j,\quad \Theta_{\mathrm R}=\max_j\lambda_j.
\]
Jeżeli jednocześnie przekształcimy wszystkie wyniki za pomocą odwracalnej macierzy $A\in\R^{p\times p}$, zastępując $H,E$ przez $A^\top HA,A^\top EA$, wszystkie cztery statystyki pozostają niezmienione.
\end{twierdzenie}
\begin{proof}
Zapiszmy $S=E^{-1/2}HE^{-1/2}$. Wówczas $\det(E+H)=\det E\det(I+S)$, a cykliczność śladu daje $\operatorname{tr}(E^{-1}H)=\operatorname{tr}S$ oraz $\operatorname{tr}((E+H)^{-1}H)=\operatorname{tr}((I+S)^{-1}S)$. Diagonalizacja symetrycznej macierzy $S$ prowadzi do wzorów. Przy zmianie $E\mapsto A^\top EA$, $H\mapsto A^\top HA$ wielomian $\det(H-tE)$ mnoży się przez $(\det A)^2$, więc jego pierwiastki, równe $\lambda_j$, nie zmieniają się.
\end{proof}

\begin{twierdzenie}[Wilks jako iloraz największych wiarygodności]\label{thm:wilks-lrt-manova}
Niech model wielowymiarowy będzie normalny jak w definicji~\ref{def:manova-he}, niech $\nu\ge p$ i niech porównywane podprzestrzenie są ustalone. Maksymalizując wiarygodność względem macierzy średnich w obu modelach i dodatnio określonej $\Sigma$, otrzymujemy iloraz największych wiarygodności modelu ograniczonego i pełnego równy $\Lambda_{\mathrm W}^{n/2}$.
\end{twierdzenie}
\begin{proof}
Dla zadanej macierzy średnich $M$ logarytm wiarygodności, z pominięciem stałej, to $-\frac n2\log\det\Sigma-\frac12\operatorname{tr}(\Sigma^{-1}(Y-M)^\top(Y-M))$. Dla ustalonej $\Sigma$ najlepsza macierz $M$ jest rzutem $P_RY$ albo $P_FY$. Odpowiednie macierze reszt to $E+H$ i $E$. Dla dowolnej dodatnio określonej macierzy reszt $S$ maksimum po $\Sigma$ osiąga się przy $\widehat\Sigma=S/n$, ponieważ dla $A=\Sigma^{-1/2}(S/n)\Sigma^{-1/2}$ nierówność $\operatorname{tr}A-\log\det A\ge p$ zachodzi po diagonalizacji. Zostaje czynnik proporcjonalny do $(\det S)^{-n/2}$. Zatem stosunek wynosi $(\det E/\det(E+H))^{n/2}$.
\end{proof}

\begin{uwaga}[Kalibracja wartości p]
Dla wielu wymiarów i wielu niezależnych kontrastów $\Lambda_{\mathrm W}$, $V_{\mathrm P}$, $T_{\mathrm{HL}}$ i $\Theta_{\mathrm R}$ mają na ogół \emph{różne} rozkłady zerowe; nie wolno automatycznie podstawiać ich do rozkładu $F$ z jedną uniwersalną liczbą stopni swobody. Programy wykorzystują dokładne wzory w wybranych przypadkach, przybliżenia $F$ albo metody numeryczne. Warto podać metodę kalibracji i stopnie swobody, a nie tylko samą wartość p. \emph{Nie} należy mylić rozkładu lambda Wilksa z twierdzeniem Wilksa o asymptotycznym rozkładzie ilorazu wiarygodności, omówionym w rozdziale~\ref{ch:asymptotic-estimation}.
\end{uwaga}

\section{Dwie grupy i dokładny test Hotellinga}
\begin{twierdzenie}[Dwuprobkowy test $T^2$ Hotellinga]\label{thm:hotelling-two-sample}
Niech $p\ge1$, $n_0,n_1\ge2$, $N=n_0+n_1>p+1$. Niech $Y_{gi}$, $g\in\{0,1\}$, $1\le i\le n_g$, będą niezależne, $Y_{gi}\sim\Normal_p(\mu_g,\Sigma)$, z tą samą dodatnio określoną $\Sigma$. Niech $S_p$ będzie wspólną, nieobciążoną macierzą kowariancji
\[
 S_p=\frac{\sum_{g=0}^1\sum_{i=1}^{n_g}(Y_{gi}-\bar Y_g)(Y_{gi}-\bar Y_g)^\top}{N-2}.
\]
Przy $H_0:\mu_0=\mu_1$ statystyka $T^2=\frac{n_0n_1}{N}(\bar Y_1-\bar Y_0)^\top S_p^{-1}(\bar Y_1-\bar Y_0)$ spełnia
\[
 \frac{N-p-1}{p(N-2)}T^2\sim F_{p,N-p-1}.
\]
\end{twierdzenie}
\begin{proof}
Niech $\nu=N-2$ i $z=\sqrt{n_0n_1/N}\,\Sigma^{-1/2}(\bar Y_1-\bar Y_0)$. Pod hipotezą zerową $z\sim\Normal_p(0,I)$, natomiast $W=\Sigma^{-1/2}(\nu S_p)\Sigma^{-1/2}\sim W_p(I,\nu)$. Z twierdzenia~\ref{thm:wishart-he} oraz niezależności normalnej średniej od reszt mamy $z\perp W$. Aby ustalić rozkład $z^\top W^{-1}z$, dla $z\ne0$ obróćmy układ współrzędnych tak, by pierwszy wektor osiowy był równy $z/\|z\|$. Dzięki ortogonalnej niezmienniczości $W$ rozkład elementu $(W^{-1})_{11}$ nie zależy od obrotu. Zapisując $W=A^\top A$ dla macierzy $A\in\R^{\nu\times p}$ o niezależnych standardowych elementach normalnych i stosując wzór na odwrotność macierzy blokowej, otrzymujemy
\[
 (W^{-1})_{11}^{-1}=a_1^\top(I-P_{A_2})a_1\sim\chi^2_{\nu-p+1},
\]
gdzie $a_1$ jest pierwszą kolumną $A$, $A_2$ zawiera pozostałe kolumny, a $P_{A_2}$ rzutuje na ich rozpiętość (ma ona niemal na pewno wymiar $p-1$). Ponieważ $\|z\|^2\sim\chi^2_p$ jest niezależne od $W$, mamy $z^\top W^{-1}z\stackrel d=\chi_p^2/\chi_{\nu-p+1}^2$ z niezależnymi licznikami. Ponieważ $T^2=\nu z^\top W^{-1}z$, otrzymujemy zadany rozkład $F$.
\end{proof}

\begin{wniosek}[Jeden kontrast grupowy]\label{cor:manova-rankone}
W warunkach definicji~\ref{def:manova-he}, jeśli $q=1$ i $\nu\ge p$, wszystkie niezerowe pierwiastki charakterystyczne sprowadzają się do jednego $\lambda$. Wówczas $\Lambda_{\mathrm W}=(1+\lambda)^{-1}$, $V_{\mathrm P}=\lambda/(1+\lambda)$, $T_{\mathrm{HL}}=\Theta_{\mathrm R}=\lambda$ i statystyka $\frac{\nu-p+1}{p}\lambda$ ma pod $H_0$ rozkład $F_{p,\nu-p+1}$.
\end{wniosek}
\begin{proof}
Dla $q=1$ macierz $H$ ma rząd co najwyżej jeden, stąd tożsamości spektralne. Rozkład wynika z dowodu twierdzenia~\ref{thm:hotelling-two-sample}, którego argument wykorzystuje jedynie niezależny standardowy wektor normalny i macierz Wisharta o $\nu$ stopniach swobody, a nie szczególną postać macierzy wskaźników grup.
\end{proof}

\begin{zadanie}[Hotelling a oddzielne testy dwóch wyników]
Dla dwóch niezależnych grup po $10$ osób zmierzono dwa wyniki. Załóżmy poprawny model normalny o wspólnej kowariancji i przyjmijmy zaobserwowaną różnicę średnich $\bar Y_1-\bar Y_0=(1,1)^\top$ oraz połączoną macierz kowariancji
$S_p=\begin{pmatrix}1&0{,}4\\0{,}4&1\end{pmatrix}$.
Wyznacz statystykę testu globalnego i porównaj decyzję z oddzielnymi testami t przy poziomie $0{,}05$.\par
\emph{Rozwiązanie.} Mamy $N=20$, $p=2$, $S_p^{-1}(1,1)^\top=(1/1{,}4)(1,1)^\top$, zatem $T^2=5\cdot 2/1{,}4=50/7\approx7{,}143$. Dokładna statystyka $F=\frac{17}{36}T^2\approx3{,}373$ ma pod $H_0$ rozkład $F_{2,17}$, co daje wartość p około $0{,}0584$. Test globalny nie odrzuca na poziomie $0{,}05$. Każdy osobny test t ma natomiast statystykę $1/\sqrt{1/10+1/10}\approx2{,}236$ i dwustronną wartość p około $0{,}0382$. To nie jest sprzeczność: hipotezy i statystyki są różne, a dwie nieskorygowane wartości p nie zapewniają kontroli błędu rodzinnego na poziomie $0{,}05$. Wniosek o konkretnych wynikach wymaga zaplanowanych kontrastów i uwzględnienia wielokrotności.
\end{zadanie}

\section{MANCOVA: korekta o kowarianty, kontrasty i interpretacja}
W MANCOVA dodajemy do macierzy $X$ kolumny zmiennych objaśniających mierzonych przed interwencją. Na przykład dla trzech grup i wieku $x_i$ porównujemy modele $M_R$ i $M_F$, gdzie do obu należy wyraz wolny oraz wiek, a tylko do pełnego należą dwa wskaźniki grup. Wówczas $q=2$ i test dotyczy \emph{łącznie wszystkich $p$ wyników}. Gdy nachylenie względem wieku różni się między grupami, należy modelować interakcję grupa$\times$wiek; wspólny współczynnik nie wynika z samej nazwy ANCOVA.

\begin{stwierdzenie}[Korekta o kowarianty jako test częściowy]\label{prop:mancova-adjusted}
Niech $Z\in\R^{n\times s}$ ma pełną rangę i zawiera wyraz wolny oraz kowarianty, a $G\in\R^{n\times q}$ koduje porównywane grupy. Jeśli $X_F=[Z\ G]$ ma pełną rangę, to w wielowymiarowym modelu liniowym $Y=X_FB+U$ test hipotezy, że wszystkie współczynniki odpowiadające kolumnom $G$ są zerowe, można przeprowadzić przez parę
$E=Y^\top(I-P_{X_F})Y$, $H=Y^\top(P_{X_F}-P_Z)Y$.
\end{stwierdzenie}
\begin{proof}
Hipoteza ustala, że kolumny średniej $\E Y$ leżą w $\operatorname{im}Z$. Model pełny dopuszcza całą przestrzeń $\operatorname{im}X_F$. Podprzestrzenie są zagnieżdżone, więc bezpośrednio stosujemy definicję~\ref{def:manova-he}.
\end{proof}

\begin{uwaga}[Kiedy MANOVA/MANCOVA nie są dobrym wyborem?]
Model wymaga sensownej wspólnej analizy kilku wyników i odpowiedniego stosunku liczby niezależnych jednostek do $p$. Silna kolinearność wyników, duże $p$, bardzo nierówne macierze kowariancji między grupami i obserwacje odstające mogą zniszczyć stabilność wyników. Przy pojedynczym głównym wyniku lepiej zwykle z góry zdefiniować ten wynik niż sztucznie mnożyć wskaźniki. Przy wielu interpretowanych później testach jednowymiarowych należy kontrolować wielokrotność i wyraźnie odróżniać \emph{test globalny} od konkretnych kontrastów. Jeśli wyniki mają odmienne skale lub typy (np. ciągły wynik i liczba zdarzeń), wspólny model normalny może być nieadekwatny.
\end{uwaga}

\begin{center}\small
\begin{tabular}{>{\raggedright\arraybackslash}p{.29\textwidth}>{\raggedright\arraybackslash}p{.27\textwidth}>{\raggedright\arraybackslash}p{.29\textwidth}}
\toprule
\textbf{Pytanie} & \textbf{Metoda do rozważenia} & \textbf{Co sprawdzić}\\
\midrule
Jedna zmienna wynikowa, wiele grup & ANOVA lub Welch & Wariancje, jednostki, kontrasty\\
Wiele skorelowanych wyników, grupy & MANOVA & Kowariancje, $p$ względem $n$, globalna hipoteza\\
Wiele wyników i wspólna kowarianta & MANCOVA & Nakładanie zakresów, interakcje, czas pomiaru\\
Wielokrotne pomiary u tych samych osób & RM--ANOVA albo model mieszany & Kowariancja powtórzeń, braki danych\\
Wiele wyników o różnych typach & Odpowiednie GLM / modele wspólne & Rozkłady i łączny cel analizy\\
\bottomrule
\end{tabular}
\end{center}

\section{Laboratorium: globalny test i interpretacja wyników}
\begin{lstlisting}[language=Python,caption={Kompletny program: MANOVA, MANCOVA i pierwiastki charakterystyczne (Python, NumPy, pandas, statsmodels).}]
import numpy as np
import pandas as pd
from statsmodels.multivariate.manova import MANOVA

rng = np.random.default_rng(2026)
n_per_group = 40
covariance = np.array([[1.0, 0.65, 0.35],
                       [0.65, 1.0, 0.40],
                       [0.35, 0.40, 1.0]])
rows = []
for group, effect in zip(("A", "B", "C"), (0.0, 0.45, 0.8)):
    for _ in range(n_per_group):
        age = rng.normal(50, 9)
        noise = rng.multivariate_normal(np.zeros(3), covariance)
        outcomes = np.array([2.0, 5.0, 8.0])
        outcomes += effect * np.array([1.0, 0.7, -0.45])
        outcomes += (age - 50) * np.array([0.06, -0.04, 0.03])
        outcomes += noise
        rows.append((group, age, *outcomes))
df = pd.DataFrame(rows, columns=["group", "age", "y1", "y2", "y3"])

print("MANOVA without age:")
print(MANOVA.from_formula("y1 + y2 + y3 ~ C(group)", data=df).mv_test())
print("MANCOVA adjusted for age:")
print(MANOVA.from_formula("y1 + y2 + y3 ~ C(group) + age", data=df).mv_test())

# E and H for testing the two added group columns after adjustment.
Y = df[["y1", "y2", "y3"]].to_numpy()
G = pd.get_dummies(df["group"], drop_first=True).to_numpy(dtype=float)
Z = np.column_stack([np.ones(len(df)), df["age"].to_numpy()])
Xfull = np.column_stack([Z, G])
def residual_sscp(X):
    coef = np.linalg.lstsq(X, Y, rcond=None)[0]
    residuals = Y - X @ coef
    return residuals.T @ residuals
E = residual_sscp(Xfull)
H = residual_sscp(Z) - E
L = np.linalg.cholesky(E)
W = np.linalg.solve(L, H)
S = np.linalg.solve(L, W.T).T
roots = np.clip(np.linalg.eigvalsh((S + S.T) / 2), 0, None)
print("Wilks lambda:", np.prod(1 / (1 + roots)))
print("Pillai trace:", np.sum(roots / (1 + roots)))
print("Hotelling-Lawley:", roots.sum(), "Roy:", roots.max())
\end{lstlisting}

\begin{lstlisting}[language=Julia,caption={Kompletny program: macierze $H,E$ i cztery statystyki MANOVA z rozkładu próby (Julia, Random, Statistics, LinearAlgebra).}]
using Random, Statistics, LinearAlgebra
Random.seed!(2026)
p, n = 3, 90
g = repeat(1:3, inner=30)
Xr = ones(n, 1)
Xf = hcat(ones(n), Float64.(g .== 2), Float64.(g .== 3))
A = [1.0 0.0 0.0; 0.6 0.8 0.0; 0.2 0.4 0.9]
noise = randn(n, p) * transpose(A)
mu = hcat(0.6 .* Float64.(g .== 2),
          0.8 .* Float64.(g .== 3),
          0.5 .* Float64.(g .== 3))
Y = mu + noise
function sscp(X, Y)
    R = Y - X * (X \ Y)
    return transpose(R) * R
end
E = sscp(Xf, Y)
H = sscp(Xr, Y) - E
L = cholesky(Symmetric(E)).L
S = L \ H / transpose(L)
roots = max.(eigvals(Symmetric((S + transpose(S))/2)), 0.0)
println("Wilks: ", prod(1 ./ (1 .+ roots)))
println("Pillai: ", sum(roots ./ (1 .+ roots)))
println("Hotelling-Lawley: ", sum(roots))
println("Roy: ", maximum(roots))
# For calibrated p-values use a MANOVA implementation and its stated method.
\end{lstlisting}

\begin{zadanie}[Dlaczego osobne testy nie rozstrzygają hipotezy MANOVA?]
W badaniu dwóch metod rehabilitacji zmierzono siłę, prędkość chodu i zakres ruchu. Każdy pojedynczy test t daje wartość p większą niż $0{,}05$. Czy oznacza to, że MANOVA nie może odrzucić hipotezy zerowej?\par
\emph{Rozwiązanie.} Nie. MANOVA bierze pod uwagę wektor średnich i jego macierz kowariancji: odpowiednia \emph{kombinacja liniowa} wyników może mieć wysoki stosunek sygnału do szumu, mimo że każda zmienna osobno daje słabą informację. Odwrotna sytuacja także jest możliwa. Test globalny nie mówi automatycznie, które współrzędne odpowiadają za różnicę; potrzebne są wcześniej zaplanowane kontrasty i analiza wielokrotności.
\end{zadanie}

\paragraph{Źródła do dalszej weryfikacji implementacji.}
Opisy rozkładów i sposobów aproksymacji statystyk wielowymiarowych można zestawić z dokumentacją: \href{https://online.stat.psu.edu/stat505/Lesson08}{Penn State STAT 505, MANOVA} oraz \href{https://support.sas.com/documentation/cdl/en/statug/68162/HTML/default/statug_introreg_sect038.htm}{SAS, Multivariate Tests}. Różne programy mogą prezentować inny wariant pierwiastka Roya lub inne przybliżenia wartości p.

\chapter{Pomiary powtarzane, sferyczność i modele mieszane}
\label{ch:repeated-measures}
\paragraph{Powiązania i odsyłacze.} Zależność pomiarów i procesy stochastyczne omawia tom II, rozdział~\ref{II-ch:procesy}; ANOVA oraz MANOVA znajdują się w rozdziałach~\ref{ch:anova-ancova} i~\ref{ch:manova-mancova}.


\section{Jednostka badania, powtórzenie i pseudoreplikacja}
Mierzymy ciśnienie u $n$ pacjentów przed leczeniem i w trzech kolejnych terminach. Otrzymujemy $4n$ liczb, ale nadal tylko $n$ \emph{niezależnych osób}. Pomiary u jednej osoby mogą być silnie skorelowane, a ich zależność jest częścią modelu, nie przeszkodą możliwą do pominięcia. Jednoczynnikowa ANOVA niezależnych obserwacji zastosowana do wszystkich $4n$ liczb ma na ogół błędny mianownik statystyki testowej.

\begin{center}
\begin{tikzpicture}[x=1.25cm,y=.65cm,>=Latex]
\draw[->] (0,0)--(5,0) node[right,font=\small] {wizyta};
\draw[->] (0,0)--(0,4.9) node[above,font=\small] {wynik};
\foreach \x/\lab in {1/1,2/2,3/3,4/4} {\draw (\x,.1)--(\x,-.1) node[below,font=\scriptsize] {\lab};}
\draw[blue!65!black,thick] (1,1.2)--(2,1.9)--(3,2.4)--(4,2.2);
\draw[red!70!black,thick] (1,2.2)--(2,2.0)--(3,3.1)--(4,3.4);
\draw[green!55!black,thick] (1,3.1)--(2,3.6)--(3,3.3)--(4,4.1);
\foreach \x/\y in {1/1.2,2/1.9,3/2.4,4/2.2} {\fill[blue!65!black] (\x,\y) circle (2pt);}
\foreach \x/\y in {1/2.2,2/2,3/3.1,4/3.4} {\fill[red!70!black] (\x,\y) circle (2pt);}
\foreach \x/\y in {1/3.1,2/3.6,3/3.3,4/4.1} {\fill[green!55!black] (\x,\y) circle (2pt);}
\node[font=\footnotesize,blue!65!black,right] at (4.1,2.2) {osoba 1};
\node[font=\footnotesize,red!70!black,right] at (4.1,3.4) {osoba 2};
\node[font=\footnotesize,green!55!black,right] at (4.1,4.1) {osoba 3};
\end{tikzpicture}

{\small\emph{Rysunek 30.1.} Schemat trajektorii: pomiary tej samej osoby tworzą parę lub całą serię, a nie kilka nowych niezależnych jednostek.}
\end{center}

\begin{definicja}[Model powtarzanych pomiarów]\label{def:rm-model}
Niech $n\ge2$, $k\ge2$. Dla osoby $i=1,\ldots,n$ obserwujemy wektor $Y_i=(Y_{i1},\ldots,Y_{ik})^\top\in\R^k$ zawierający $k$ pomiarów w ustalonych terminach. W podstawowym modelu zakładamy, że $Y_1,\ldots,Y_n$ są niezależne, $Y_i\sim\Normal_k(\mu,\Sigma)$, gdzie $\mu\in\R^k$ jest nieznane, a $\Sigma$ jest dodatnio określona. Testujemy hipotezę
\[
 H_0:\mu_1=\cdots=\mu_k.
\]
Nie zakładamy niezależności współrzędnych $Y_i$. Projekt określający kolejność wizyt i brakujące pomiary będzie wymagał osobnego komentarza.
\end{definicja}

\begin{przyklad}[Ile wynosi wariancja zmiany?]
Dla dwóch pomiarów $Y_1,Y_2$ jednej osoby o wariancjach $\sigma_1^2,\sigma_2^2$ i korelacji $\rho$ mamy $\Var(Y_2-Y_1)=\sigma_1^2+\sigma_2^2-2\rho\sigma_1\sigma_2$. Przy dodatniej korelacji porównanie \emph{sparowanych różnic} może być znacznie dokładniejsze niż test zakładający dwie niezależne próby. Nie należy jednak zastępować wielu niezależnych osób wieloma replikatami tej samej osoby.
\end{przyklad}

\section{Ortogonalne kontrasty i istota sferyczności}
\begin{definicja}[Kontrasty i sferyczność]\label{def:sphericity}
Niech $k\ge2$, $d=k-1$ i niech $C\in\R^{d\times k}$ ma ortonormalne wiersze, $CC^\top=I_d$, $C\mathbf1_k=0$. Macierz $C$ nazywamy \emph{macierzą ortonormalnych kontrastów czasowych}. Dla dodatnio określonej macierzy $\Sigma\in\R^{k\times k}$ niech $\Omega=C\Sigma C^\top$. Mówimy, że zachodzi \emph{sferyczność}, jeżeli istnieje $\sigma_c^2>0$ takie, że $\Omega=\sigma_c^2 I_d$. Definicja nie zależy od wyboru macierzy $C$.
\end{definicja}

\begin{twierdzenie}[Sferyczność a wariancje różnic]\label{thm:sphericity-differences}
Niech $k\ge2$, $\Sigma$ będzie dodatnio określona, a $Y\in\R^k$ wektorem losowym o macierzy kowariancji $\Sigma$. Sferyczność jest równoważna istnieniu liczby $v>0$ takiej, że
\[
 \Var(Y_j-Y_l)=v\quad\text{dla wszystkich }j\ne l.
\]
Dla $k=2$ warunek zachodzi automatycznie.
\end{twierdzenie}
\begin{proof}
Niech $P=I_k-k^{-1}\mathbf1_k\mathbf1_k^\top=C^\top C$ będzie rzutem na przestrzeń kontrastów. Jeśli $C\Sigma C^\top=\sigma_c^2 I_d$, to $P\Sigma P=\sigma_c^2P$. Ponieważ $e_j-e_l$ należy do obrazu $P$ i ma normę $\sqrt2$, otrzymujemy $\Var(Y_j-Y_l)=2\sigma_c^2$.

Odwrotnie, niech $D\in\R^{k\times k}$ będzie macierzą o wyrazach $D_{jl}=\Var(Y_j-Y_l)$, czyli $D_{jl}=\Sigma_{jj}+\Sigma_{ll}-2\Sigma_{jl}$. Po obustronnym przemnożeniu przez $P$ składniki postaci $a\mathbf1_k^\top$ i $\mathbf1_k a^\top$ znikają, stąd $PDP=-2P\Sigma P$. Jeżeli $D_{jj}=0$ i $D_{jl}=v$ dla $j\ne l$, to $D=v(\mathbf1_k\mathbf1_k^\top-I_k)$ i $PDP=-vP$. Wynika stąd $P\Sigma P=(v/2)P$, a po pomnożeniu przez $C$ i $C^\top$ otrzymujemy $\Omega=(v/2)I_d$. Gdy $k=2$, przestrzeń kontrastów jest jednowymiarowa, więc każda dodatnia forma kwadratowa jest w niej wielokrotnością identyczności.
\end{proof}

\begin{uwaga}[Sferyczność to nie to samo co niezależność ani compound symmetry]
Silniejszym warunkiem jest \emph{symetria złożona} (ang. \emph{compound symmetry}), czyli $\Sigma=\sigma^2I_k+\tau^2\mathbf1_k\mathbf1_k^\top$ w dopuszczalnym zakresie parametrów. Oznacza ona wspólne wariancje i wspólne kowariancje wszystkich par; implikuje sferyczność, ale nie jest z nią równoważna. Na przykład dla $k=3$ wektor $v=(0{,}1,0,-0{,}1)^\top$ i $\Sigma=I_3+\mathbf1_3v^\top+v\mathbf1_3^\top$ dają dodatnio określoną macierz (norma operatorowa perturbacji jest mniejsza od $2\sqrt3\|v\|_2<1$), dla której $P\Sigma P=P$, choć przekątne i pozaprzekątne wyrazy $\Sigma$ nie są stałe.
\end{uwaga}

\section{Dokładny test \texorpdfstring{$F$}{F} w klasycznej RM--ANOVA}
\begin{twierdzenie}[Rozkład statystyki powtarzanych pomiarów przy sferyczności]\label{thm:rm-F}
Niech $n\ge2$, $k\ge2$, $d=k-1$, a $Y_i\sim\Normal_k(\mu,\Sigma)$ będą niezależne jak w definicji~\ref{def:rm-model}. Załóżmy sferyczność $C\Sigma C^\top=\sigma_c^2 I_d$. Oznaczmy $Z_i=CY_i$ i $\bar Z=n^{-1}\sum_i Z_i$. Jeśli $H_0:C\mu=0$, to
\[
 F_{\mathrm{RM}}=
 \frac{n\|\bar Z\|_2^2/d}{\bigl(\sum_{i=1}^n\|Z_i-\bar Z\|_2^2\bigr)/((n-1)d)}
 \sim F_{d,(n-1)d}.
\]
\end{twierdzenie}
\begin{proof}
W modelu kontrastów $Z_i$ są niezależne o rozkładzie $\Normal_d(C\mu,\sigma_c^2 I_d)$. Pod $H_0$ wektory te są centrowane. Ortogonalny obrót współrzędnych w przestrzeni $\R^n$ osób rozdziela ich macierz na wiersz $\sqrt n\,\bar Z$ oraz $n-1$ wierszy będących kontrastami między osobami. Ze względu na wspólną izotropową macierz kowariancji te $n$ wierszy są niezależne, a ich współrzędne mają rozkład $\Normal(0,\sigma_c^2)$. Zatem $n\|\bar Z\|_2^2/\sigma_c^2\sim\chi_d^2$ oraz $\sum_i\|Z_i-\bar Z\|_2^2/\sigma_c^2\sim\chi_{(n-1)d}^2$, niezależnie. Ich odpowiednio przeskalowany iloraz ma deklarowany rozkład.
\end{proof}

\begin{uwaga}[Co się zmienia bez sferyczności?]
Jeżeli $\Omega=C\Sigma C^\top$ ma różne wartości własne $\lambda_1,\ldots,\lambda_d$, to pod $H_0$ kwadrat normy średniej z kontrastów ma rozkład sumy $\sum_{j=1}^d\lambda_j Z_j^2$ z niezależnymi $Z_j\sim\Normal(0,1)$. Odpowiedni iloraz w ogólności nie jest dokładnie rozkładem $F$. Dla dwóch terminów ($d=1$) problem sferyczności nie występuje i test odpowiada sparowanemu testowi $t$ po podniesieniu statystyki do kwadratu.
\end{uwaga}

\section{Test Mauchly'ego i poprawki Greenhouse'a--Geissera oraz Huynha--Feldta}
\begin{definicja}[Współczynnik $\epsilon$ i estymator GG]\label{def:gg-epsilon}
Niech $d=k-1\ge2$, $\Omega=C\Sigma C^\top$ będzie dodatnio określona. Definiujemy \emph{populacyjny współczynnik niesferyczności}
\[
 \epsilon=\frac{(\operatorname{tr}\Omega)^2}{d\,\operatorname{tr}(\Omega^2)}.
\]
Dla $n>d$ i niezależnych pełnych obserwacji $Y_1,\ldots,Y_n$ definiujemy empiryczną macierz kowariancji w przestrzeni kontrastów
$S_C=(n-1)^{-1}\sum_i (CY_i-\overline{CY})(CY_i-\overline{CY})^\top$.
Jeśli $S_C$ jest dodatnio określona, estymator \emph{Greenhouse'a--Geissera} wynosi $\widehat\epsilon_{\mathrm{GG}}=(\operatorname{tr}S_C)^2/[d\operatorname{tr}(S_C^2)]$.
\end{definicja}

\begin{twierdzenie}[Zakres współczynnika $\epsilon$]\label{thm:gg-bound}
Dla każdej dodatnio określonej macierzy $\Omega\in\R^{d\times d}$, $d\ge1$, zachodzi $1/d\le\epsilon\le1$. Równość $\epsilon=1$ występuje wtedy i tylko wtedy, gdy $\Omega$ jest wielokrotnością $I_d$.
\end{twierdzenie}
\begin{proof}
Dla dodatnich wartości własnych $\lambda_1,\ldots,\lambda_d$ otrzymujemy $\epsilon=(\sum_j\lambda_j)^2/(d\sum_j\lambda_j^2)$. Górne ograniczenie jest nierównością Cauchy'ego--Schwarza i zamienia się w równość dokładnie przy wszystkich $\lambda_j$ jednakowych. Dolne ograniczenie wynika z $\sum_j\lambda_j^2\le(\sum_j\lambda_j)^2$; dla dodatnio określonej macierzy o $d>1$ jest ono w rzeczywistości ostre jedynie w granicy degeneracji. Dla $d=1$ mamy $\epsilon=1$.
\end{proof}

\begin{definicja}[Statystyka Mauchly'ego]\label{def:mauchly}
Niech $n>d$, $d=k-1\ge2$, $S_C$ będzie dodatnio określona jak w definicji~\ref{def:gg-epsilon}. Definiujemy $W_M=\det(S_C)/[\operatorname{tr}(S_C)/d]^d$. Z nierówności między średnią geometryczną i arytmetyczną wartości własnych wynika $0<W_M\le1$. Małe wartości sygnalizują odstępstwo od sferyczności. Przy niezależnych normalnych wektorach kontrastów, pod hipotezą \emph{sferyczności macierzy kowariancji}, dla ustalonego $d$ i $n\to\infty$ stosuje się przybliżenie
\[
 T_M=-\left[(n-1)-\frac{2d^2+d+2}{6d}\right]\log W_M.
\]
Rozkład zerowy statystyki $T_M$ przybliża rozkład $\chi^2_{d(d+1)/2-1}$. Jest to \emph{przybliżenie}, nie skończonopróbkowa równość rozkładów.
\end{definicja}

\begin{uwaga}[Dwa problemy z automatycznym testowaniem sferyczności]
Test Mauchly'ego wymaga modelowych warunków kalibracji i ma ograniczoną moc przy niewielkiej liczbie osób. Brak odrzucenia nie dowodzi sferyczności. Wybór korekty wyłącznie na podstawie testu Mauchly'ego może dodatkowo komplikować kontrolę poziomu błędu całej procedury. Warto obok testu podać oszacowanie $\widehat\epsilon$, liczbę terminów, wykres średnich i oszacowania różnic wraz z niepewnością.
\end{uwaga}

\begin{definicja}[Korekty stopni swobody]\label{def:gg-hf}
W klasycznej jednoczynnikowej RM--ANOVA statystykę liczbową $F_{\mathrm{RM}}$ pozostawia się bez zmian, ale do oceny wartości p używa się rozkładu $F$ ze skorygowanymi stopniami swobody $d_1=\widehat\epsilon\,d$ i $d_2=\widehat\epsilon\,(n-1)d$. Korekta GG stosuje $\widehat\epsilon=\widehat\epsilon_{\mathrm{GG}}$. Dla $d\ge2$, $n>d+1$ i dodatniego mianownika, prosty wariant estymatora Huynha--Feldta to
\[
 \widehat\epsilon_{\mathrm{HF}}=\min\left\{1,\;\frac{nd\,\widehat\epsilon_{\mathrm{GG}}-2}{d\bigl(n-1-d\widehat\epsilon_{\mathrm{GG}}\bigr)}\right\}.
\]
W praktyce, gdy wzór HF prowadzi poza przedział, używa się ograniczenia do $[1/d,1]$ i trzeba sprawdzić przyjętą w programie odmianę poprawki (niektóre implementacje stosują również modyfikację Lecoutre'a). Te procedury mają naturę przybliżeń; nie odzyskują dokładnego rozkładu $F$ przy dowolnej macierzy kowariancji.
\end{definicja}

\begin{przyklad}[Dlaczego stopnie swobody bywają niecałkowite?]
Przy $k=4$ terminach i $n=30$ osobach mamy $d=3$, klasyczne stopnie swobody $3$ i $87$. Jeżeli $\widehat\epsilon_{\mathrm{GG}}=0{,}65$, przy ocenie p pozostawiamy obserwowaną statystykę $F$, ale używamy $F_{1{,}95;\,56{,}55}$. Nie oznacza to, że liczba obserwacji została zmieniona: zmieniamy przybliżenie rozkładu referencyjnego w celu uwzględnienia nierównej zmienności kontrastów.
\end{przyklad}

\section{Modele mieszane i niejednakowa liczba wizyt}
\begin{definicja}[Liniowy model mieszany]\label{def:lmm}
Niech $i=1,\ldots,n$ indeksuje niezależne jednostki, a $m_i\ge1$ ich liczbę pomiarów. Dla każdej jednostki zakładamy model wektorowy
\[
 Y_i=X_i\beta+Z_i b_i+\varepsilon_i,
\]
gdzie $Y_i\in\R^{m_i}$, $X_i\in\R^{m_i\times r}$, $Z_i\in\R^{m_i\times q}$ są zadane, $\beta\in\R^r$ jest wektorem efektów stałych, $b_i\sim\Normal_q(0,D)$ losowym efektem jednostki, a $\varepsilon_i\sim\Normal_{m_i}(0,R_i)$ błędem wewnątrz jednostki. Zakładamy, że rodziny $(b_i,\varepsilon_i)$ są niezależne między jednostkami i $b_i\perp\varepsilon_i$, przy dodatnio półokreślonej $D$ i dodatnio określonych $R_i$.
\end{definicja}

\begin{twierdzenie}[Marginalny rozkład w modelu mieszanym]\label{thm:lmm-marginal}
W modelu~\ref{def:lmm}, dla każdej jednostki $i$ zachodzi
\[
 Y_i\sim\Normal_{m_i}(X_i\beta,V_i),\qquad V_i=Z_iDZ_i^\top+R_i.
\]
Jeżeli $V_i$ są znane i dodatnio określone, a $X$ powstała przez pionowe zestawienie $X_i$ ma pełną rangę, uogólniony estymator najmniejszych kwadratów
\[
 \widehat\beta_{\mathrm{GLS}}=\left(\sum_iX_i^\top V_i^{-1}X_i\right)^{-1}\sum_iX_i^\top V_i^{-1}Y_i
\]
jest nieobciążony i ma macierz kowariancji $\left(\sum_iX_i^\top V_i^{-1}X_i\right)^{-1}$.
\end{twierdzenie}
\begin{proof}
Wektor $Z_ib_i+\varepsilon_i$ jest sumą niezależnych normalnych wektorów, więc jest normalny, a jego kowariancja jest sumą macierzy $Z_iDZ_i^\top$ i $R_i$. Wstawienie $Y_i=X_i\beta+\xi_i$ do wzoru GLS daje różnicę $\widehat\beta_{\mathrm{GLS}}-\beta=A^{-1}\sum_iX_i^\top V_i^{-1}\xi_i$, gdzie $A=\sum_iX_i^\top V_i^{-1}X_i$. Wartość oczekiwana jest zerowa. Ponieważ $\xi_i$ są niezależne między jednostkami, wariancja sumy jest $\sum_iX_i^\top V_i^{-1}V_iV_i^{-1}X_i=A$. Ostatecznie $\Cov(\widehat\beta_{\mathrm{GLS}})=A^{-1}$.
\end{proof}

\begin{przyklad}[Losowy punkt początkowy i losowe nachylenie]
Dla czasu $t_{ij}$ można przyjąć $Y_{ij}=\beta_0+\beta_1t_{ij}+b_{0i}+b_{1i}t_{ij}+\varepsilon_{ij}$, gdzie $(b_{0i},b_{1i})$ jest dwuwymiarowym normalnym wektorem losowym. Wariancja $Y_{ij}$ i kowariancja między wizytami zależą teraz od czasów $t_{ij}$, co nie odpowiada na ogół sferyczności. Model mieszany dopuszcza różną liczbę pomiarów na osobę i rzeczywiste, nierównomiernie rozłożone chwile pomiaru. Dla danych bez kompletu wizyt wnioskowanie na podstawie obserwowanych danych wymaga jednak uzasadnionego mechanizmu brakowania, np. MAR, oraz poprawnego modelu warunkowego; sam program nie usuwa problemu MNAR.
\end{przyklad}

\begin{uwaga}[ML, REML i estymowana kowariancja]
W praktyce macierze $D,R_i$ są nieznane i estymuje się je zwykle przez ML lub \emph{restricted maximum likelihood} (REML), co zmienia rozkłady statystyk opartych na estymowanym $V_i$. Testy efektów stałych w małych próbach mogą wymagać przybliżeń Satterthwaite'a lub Kenwarda--Rogera, bootstrapu parametrycznego albo innych jawnie uzasadnionych procedur. Porównując modele o \emph{różnych efektach stałych} za pomocą ilorazu wiarygodności, należy używać dopasowania ML, nie bezpośredniego porównania wartości ograniczonej wiarygodności REML. Gdy liczba grup jest mała albo estymator wariancji efektu losowego leży na brzegu przestrzeni parametrów, standardowe asymptotyczne reguły $\chi^2$ mogą nie obowiązywać.
\end{uwaga}

\begin{center}\small
\begin{tabular}{>{\raggedright\arraybackslash}p{.28\textwidth}>{\raggedright\arraybackslash}p{.27\textwidth}>{\raggedright\arraybackslash}p{.30\textwidth}}
\toprule
\textbf{Projekt i cel} & \textbf{Metoda do rozważenia} & \textbf{Krytyczne pytanie}\\
\midrule
Dwa pomiary na osobę & Test $t$ dla par; model różnic & Czy pary są niezależne między osobami?\\
$k\ge3$ terminów, kompletne dane & RM--ANOVA, GG/HF & Czy sferyczność jest wiarygodna?\\
Czasy nierówne, różne liczby wizyt & Liniowy model mieszany & Jaki model kowariancji i braków danych?\\
Dane binarne lub zliczenia w czasie & GLMM albo GEE & Rozkład warunkowy, cel marginalny/indywidualny\\
Kilka odrębnych wyników na osobę & MANOVA/MANCOVA & Czy pytamy o wspólny wektor średnich?\\
Zależność szeregów w czasie & Modele ARMA, GLS/HAC & Czy obserwacje tworzą jeden szereg czy wiele osób?\\
\bottomrule
\end{tabular}
\end{center}

\section{Laboratorium: sferyczność, korekty i model mieszany}
\begin{lstlisting}[language=Python,caption={Kompletny program: RM--ANOVA, współczynniki GG/HF i przybliżenie testu Mauchly'ego (Python, NumPy, SciPy).}]
import numpy as np
from scipy import stats

rng = np.random.default_rng(2026)
n, k = 36, 4
mu = np.array([0.0, 0.4, 0.8, 1.0])
# Non-spherical covariance: first-order autoregressive correlations.
cov = np.fromfunction(lambda i, j: 0.65 ** np.abs(i-j), (k, k))
Y = rng.multivariate_normal(mu, cov, size=n)
# Orthonormal rows spanning the subspace perpendicular to (1,...,1).
Q, _ = np.linalg.qr(np.column_stack((np.ones(k),
                                     np.eye(k)[:, :k-1])))
C = Q[:, 1:].T
Z = Y @ C.T
zbar = Z.mean(axis=0)
d = k - 1
SS_t = n * np.sum(zbar ** 2)
SS_e = np.sum((Z - zbar) ** 2)
F = (SS_t / d) / (SS_e / ((n-1)*d))
S = np.cov(Z, rowvar=False)
epsGG = np.trace(S)**2 / (d * np.trace(S @ S))
epsHF = min(1.0, (n*d*epsGG - 2)/(d*(n-1-d*epsGG)))
epsHF = max(1/d, epsHF)
W = np.linalg.det(S) / (np.trace(S)/d)**d
chi_m = -((n-1) - (2*d*d+d+2)/(6*d))*np.log(W)
p_m = stats.chi2.sf(chi_m, d*(d+1)//2 - 1)
print("RM F:", F, "uncorrected p:", stats.f.sf(F, d, (n-1)*d))
print("GG epsilon:", epsGG, "GG p:",
      stats.f.sf(F, epsGG*d, epsGG*(n-1)*d))
print("HF epsilon:", epsHF, "HF p:",
      stats.f.sf(F, epsHF*d, epsHF*(n-1)*d))
print("Mauchly W:", W, "approximate p:", p_m)
\end{lstlisting}

\begin{lstlisting}[language=Python,caption={Kompletny program: liniowy model mieszany z losowym punktem początkowym i różną liczbą wizyt (Python, pandas, statsmodels).}]
import numpy as np
import pandas as pd
import statsmodels.formula.api as smf

rng = np.random.default_rng(2027)
rows = []
for subject in range(80):
    intercept = rng.normal(0, 1.5)
    # Different visit counts; missingness here is generated independently.
    times = np.sort(rng.choice(np.arange(5),
                               size=rng.integers(2, 6), replace=False))
    for time in times:
        y = 10 + 0.6*time + intercept + rng.normal(0, 0.8)
        rows.append((subject, float(time), y))
df = pd.DataFrame(rows, columns=["subject", "time", "response"])
model = smf.mixedlm("response ~ time", data=df,
                    groups=df["subject"], re_formula="1")
fit = model.fit(reml=True, method="lbfgs")
print(fit.summary())
print("estimated time effect:", fit.params["time"])
# Random intercept alone induces compound symmetry, not arbitrary AR(1).
# For AR-structured within-subject residuals use a suitable GLS/GEE model.
\end{lstlisting}

\begin{lstlisting}[language=Julia,caption={Kompletny program: empiryczny współczynnik Greenhouse'a--Geissera z ortonormalnych kontrastów (Julia, Random, Statistics, LinearAlgebra).}]
using Random, Statistics, LinearAlgebra
Random.seed!(2026)
n, k = 36, 4
C = transpose(Matrix(qr(hcat(ones(k), Matrix{Float64}(I, k, k-1))).Q)[:, 2:k])
# Data matrix: each row is one independent subject, each column a visit.
Y = randn(n, k) .+ reshape([0.0, 0.4, 0.8, 1.0], 1, :)
Z = Y * transpose(C)
S = cov(Z; dims=1)
d = k - 1
epsilon_gg = tr(S)^2 / (d * tr(S*S))
println("GG epsilon: ", epsilon_gg)
println("GG degrees of freedom: ", epsilon_gg*d, ", ", epsilon_gg*(n-1)*d)
# With k=2 sphericity holds automatically, and epsilon is exactly 1.
\end{lstlisting}

\begin{przyklad}[Wartości p bez i z korektą sferyczności]
W laboratorium powyżej zasymulowano $n=36$ niezależnych osób z $k=4$ wizytami i macierzą kowariancji o postaci autoregresyjnej. Wartości p uzyskane po korektach GG/HF są aproksymacjami; ich ważność wymaga założeń dotyczących wspólnej struktury wewnątrzosobowej, nie jest gwarantowana przy dowolnym procesie brakowania wizyt. Dla ustalonego ziarna otrzymujemy $F_{\mathrm{RM}}\approx9{,}259$ i $\widehat\epsilon_{\mathrm{GG}}\approx0{,}782$. Nieskorygowane przybliżenie daje $p\approx1{,}73\cdot10^{-5}$; korekta GG daje $p\approx1{,}04\cdot10^{-4}$, a wersja HF $p\approx6{,}35\cdot10^{-5}$. Wszystkie trzy metody wskazują tutaj odrzucenie, ale niepewność wyraźnie się zmienia. Są to wartości obliczone na \emph{jednym konkretnym zbiorze danych}, a nie uniwersalne gwarancje tych korekt.
\end{przyklad}

\begin{zadanie}[Trzy wizyty, sześć pomiarów -- ile niezależnych jednostek?]
U $40$ pacjentów wykonano po trzy pomiary w sześciu technicznie powtarzanych próbach. Badacz raportuje $720$ ``niezależnych obserwacji'' i klasyczną ANOVA trzech grup. Oceń ten wybór.\par
\emph{Rozwiązanie.} Niezależnych pacjentów jest $40$, a nie $720$. Pomiary są zagnieżdżone: repliki techniczne w wizytach, wizyty w pacjentach. Model musi uwzględniać oba poziomy zależności albo najpierw uzasadnić agregację replik technicznych. Test ANOVA niezależnych obserwacji może znacznie zaniżać błąd standardowy. Wybór modelu zależy od tego, czy celem jest zmiana średniej w czasie, porównanie osób czy analiza dokładności technicznego pomiaru.
\end{zadanie}

\paragraph{Dalsze źródła dotyczące kalibracji.}
Dokładną konwencję używaną do korekt $\epsilon$ i testów wielowymiarowych należy sprawdzić w dokumentacji programu. Zarówno test Mauchly'ego, jak i korekty GG/HF odnoszą się do struktury kowariancji \emph{wewnątrz osoby}; nie kontrolują zależności pomiędzy osobami, selektywnych braków danych ani wyboru terminów pomiaru po obejrzeniu wyników. Zobacz także np. \href{https://support.sas.com/documentation/cdl/en/statug/68162/HTML/default/statug_glm_details46.htm}{SAS, Repeated Measures Analysis of Variance} albo \href{https://www.mathworks.com/help/stats/repeatedmeasuresmodel.epsilon.html}{MATLAB, Epsilon Adjustment}. Podstawowe twierdzenie~\ref{thm:rm-F} jest dokładne pod jawnie zapisanymi hipotezami; procedury GG/HF i Mauchly'ego wykorzystują odpowiednie aproksymacje.

% --- Rozdzial 31: eksperymenty i wnioskowanie przyczynowe ---
\chapter{Eksperymenty, A/B testing i wnioskowanie przyczynowe}
\label{ch:causal-ab}
\paragraph{Powiązania i odsyłacze.} Podstawy testowania i wielokrotności znajdziesz w rozdziałach~\ref{ch:testing} i~\ref{ch:multiple-testing}; regresję do korekty zmiennych w rozdziale~\ref{ch:linear-regression}.


\section{Od związku statystycznego do efektu działania}
W badaniu obserwacyjnym pacjenci przyjmujący nowy lek mogą różnić się od pozostałych wiekiem, stanem zdrowia i dostępem do opieki. Samo porównanie ich wyników nie musi zatem mierzyć działania leku. Podobny problem występuje w A/B testingu: użytkownicy, którzy sami włączyli nową funkcję, mogą być bardziej aktywni już \emph{przed} jej włączeniem. Wnioskowanie przyczynowe pyta o wynik, jaki wystąpiłby u tej samej jednostki po zastosowaniu różnych, jasno określonych działań.

\begin{definicja}[Wyniki potencjalne, spójność i jednostka eksperymentu]\label{def:potential-outcomes}
Niech $i=1,\ldots,N$ numeruje ustalone jednostki. Dla każdej jednostki definiujemy dwie \emph{ustalone} wartości potencjalne $Y_i(0),Y_i(1)\in\R$: wyniki odpowiednio bez interwencji i z interwencją. Zmienna $A_i\in\{0,1\}$ oznacza przydział, a obserwowany wynik spełnia warunek \emph{spójności}
\[
Y_i^{\mathrm{obs}}=(1-A_i)Y_i(0)+A_iY_i(1).
\]
W tym modelu wynik jednostki nie zależy od przydziałów innych jednostek, a sposób przeprowadzenia interwencji jest jednoznaczny. Są to elementy założenia braku interferencji i dobrze określonych interwencji; ich niespełnienie wymaga innego modelu wyników potencjalnych.
Definiujemy indywidualny efekt $\tau_i=Y_i(1)-Y_i(0)$ i średni efekt w badanej populacji skończonej $\tau_N=N^{-1}\sum_{i=1}^N\tau_i$.
\end{definicja}

Nie obserwujemy jednocześnie $Y_i(0)$ i $Y_i(1)$ dla tej samej jednostki. Jest to \emph{fundamentalny problem wnioskowania przyczynowego}; nie można go usunąć przez większą liczbę pomiarów tej samej zrealizowanej interwencji. Randomizacja pozwala natomiast estymować \emph{średni} efekt bez konieczności odtwarzania wszystkich efektów indywidualnych.

\begin{uwaga}[ATE, ATT i ryzyko wyboru niewłaściwego celu]
Przy losowej jednostce z populacji docelowej definiujemy $\mathrm{ATE}=\E[Y(1)-Y(0)]$ oraz $\mathrm{ATT}=\E[Y(1)-Y(0)\mid A=1]$, o ile odpowiednie wartości oczekiwane istnieją. Nie muszą być równe: ATT dotyczy osób rzeczywiście leczonych, a ATE całej populacji docelowej. W eksperymencie na ustalonych $N$ jednostkach podstawowym parametrem będzie $\tau_N$. Przenoszenie wniosku z uczestników na inną populację wymaga osobnego założenia o reprezentatywności lub modelu transportowalności.
\end{uwaga}

\section{Pełna randomizacja i jej gwarancje skończonopróbkowe}
\begin{definicja}[Kompletny plan losowania]\label{def:complete-randomization}
Niech $N\ge4$, $2\le n_1\le N-2$, $n_0=N-n_1$. Losujemy wektor $A=(A_1,\ldots,A_N)$ jednostajnie ze zbioru
$\mathcal W=\{a\in\{0,1\}^N:\sum_i a_i=n_1\}$, niezależnie od ustalonych wyników potencjalnych. Estymator różnicy średnich ma postać
\[
\widehat\tau=\frac1{n_1}\sum_{i:A_i=1}Y_i^{\mathrm{obs}}-\frac1{n_0}\sum_{i:A_i=0}Y_i^{\mathrm{obs}}.
\]
Dla $a\in\{0,1\}$ oznaczamy $\bar Y(a)=N^{-1}\sum_iY_i(a)$, a $S_a^2=(N-1)^{-1}\sum_i(Y_i(a)-\bar Y(a))^2$. Niech także $S_\tau^2=(N-1)^{-1}\sum_i(\tau_i-\tau_N)^2$.
\end{definicja}

\begin{twierdzenie}[Nieobciążoność i wariancja Neymana]\label{thm:neyman-variance}
W planie~\ref{def:complete-randomization}, przy ustalonych wynikach potencjalnych, zachodzi
\[
\E_A\widehat\tau=\tau_N,\qquad
\Var_A(\widehat\tau)=\frac{S_1^2}{n_1}+\frac{S_0^2}{n_0}-\frac{S_\tau^2}{N},
\]
gdzie $\E_A,\Var_A$ oznaczają wartość oczekiwaną i wariancję wyłącznie względem randomizacji.
\end{twierdzenie}
\begin{proof}
Dla każdego $i$ prawdopodobieństwo $A_i=1$ wynosi $n_1/N$, więc liniowość wartości oczekiwanej daje $\E_A[n_1^{-1}\sum_i A_iY_i(1)]=\bar Y(1)$ i analogicznie $\E_A[n_0^{-1}\sum_i(1-A_i)Y_i(0)]=\bar Y(0)$.

Dla ustalonego wektora liczb $u=(u_1,\ldots,u_N)$ z wariancją skończonopopulacyjną $S_u^2$, średnia w losowej próbie bez zwracania o liczności $m$ ma wariancję $(m^{-1}-N^{-1})S_u^2$. Ponadto, dla dwóch wektorów $u,v$, kowariancja średniej $u$ w grupie losowo wybranej i średniej $v$ w grupie dopełniającej wynosi $-S_{uv}/N$, gdzie $S_{uv}=(N-1)^{-1}\sum_i(u_i-\bar u)(v_i-\bar v)$. Ostatni wzór wynika z $\Cov_A(A_i,1-A_j)=-(n_1n_0/N^2)$ dla $i=j$ i $n_1n_0/[N^2(N-1)]$ dla $i\ne j$, po rozwinięciu podwójnej sumy. Otrzymujemy
$\Var_A(\widehat\tau)=(n_1^{-1}-N^{-1})S_1^2+(n_0^{-1}-N^{-1})S_0^2+2S_{10}/N$.
Ponieważ $S_\tau^2=S_1^2+S_0^2-2S_{10}$, wynika teza.
\end{proof}

\begin{wniosek}[Konserwatywna ocena niepewności]\label{cor:neyman-se}
W planie~\ref{def:complete-randomization} niech $s_a^2$ będzie wariancją z mianownikiem $n_a-1$, obliczoną z wyników obserwowanych w grupie $a$, $a=0,1$. Estymator
$\widehat V_N=s_1^2/n_1+s_0^2/n_0$ spełnia $\E_A\widehat V_N=\Var_A(\widehat\tau)+S_\tau^2/N\ge\Var_A(\widehat\tau)$.
\end{wniosek}
\begin{proof}
Każda grupa jest próbą prostą bez zwracania z odpowiedniego skończonego zbioru wyników potencjalnych, więc $\E_As_a^2=S_a^2$. Podstawiamy do twierdzenia~\ref{thm:neyman-variance}.
\end{proof}

\begin{uwaga}[Czego nie obiecuje ten wniosek?]
Konserwatywność zachodzi dla \emph{wartości oczekiwanej estymatora wariancji}, nie stanowi automatycznej dokładnej gwarancji pokrycia dowolnego przedziału $\widehat\tau\pm z_{1-\alpha/2}\sqrt{\widehat V_N}$ przy małym $N$. Przedziały normalne wymagają odpowiednich warunków asymptotycznych. Ponadto, gdy skutki indywidualne nie są jednakowe, $S_\tau^2$ nie jest identyfikowalne z samej pojedynczej realizacji randomizacji.
\end{uwaga}

\section{Dokładny test randomizacyjny i dwa różne rodzaje hipotezy zerowej}
\begin{definicja}[Hipoteza ostra Fishera]\label{def:fisher-sharp}
W skończonopopulacyjnym modelu~\ref{def:potential-outcomes} \emph{ostrą hipotezą zerową Fishera} nazywamy $H_0^{\mathrm F}:Y_i(1)=Y_i(0)$ dla wszystkich $i=1,\ldots,N$. W odróżnieniu od hipotezy słabej Neymana $H_0^{\mathrm N}:\tau_N=0$, hipoteza ostra pozwala odtworzyć wszystkie brakujące wyniki potencjalne z zaobserwowanych wyników.
\end{definicja}

\begin{twierdzenie}[Dokładna kontrola poziomu testu randomizacyjnego]\label{thm:fisher-exact}
Niech losowanie ma plan~\ref{def:complete-randomization}, niech pod $H_0^{\mathrm F}$ wyniki są stałym wektorem $y=(y_1,\ldots,y_N)$, a $T(a,y)\in\R$ dowolną ustaloną przed obejrzeniem danych statystyką, w której większa wartość oznacza silniejsze odstępstwo od hipotezy. Dla obserwowanego $A\in\mathcal W$ zdefiniujmy
\[
 p_{\mathrm{rand}}(A,y)=\frac1{|\mathcal W|}\sum_{a\in\mathcal W}\mathbf 1_{\{T(a,y)\ge T(A,y)\}}.
\]
Wtedy dla każdego $\alpha\in[0,1]$ zachodzi $\Prob_A(p_{\mathrm{rand}}(A,y)\le\alpha)\le\alpha$.
\end{twierdzenie}
\begin{proof}
Niech $M=|\mathcal W|$ i uszeregujmy możliwe wartości statystyki nierosnąco, uwzględniając powtórzenia. Dla każdego wyniku na pozycji $j$ liczba wartości nie mniejszych od niego jest co najmniej $j$. Zatem warunek $p_{\mathrm{rand}}\le\alpha$ spełnia co najwyżej $\lfloor\alpha M\rfloor$ z $M$ jednakowo prawdopodobnych przydziałów. Losowość dotyczy jedynie $A$, więc prawdopodobieństwo odrzucenia nie przekracza $\alpha$.
\end{proof}

\begin{przyklad}[Mały eksperyment: wszystkie przydziały]
Dla $N=6$ i $n_1=n_0=3$ jest $\binom63=20$ dopuszczalnych przydziałów. Po zapisaniu wyników możemy dla każdego przydziału obliczyć różnicę średnich. Jeżeli obserwowana wartość bezwzględna jest co najmniej tak duża jak w dwóch spośród dwudziestu przydziałów (wliczając obserwowany), dokładna dwustronna wartość p wynosi $2/20=0{,}1$. Nie należy tutaj wstawiać dowolnego asymptotycznego rozkładu t: rozkład randomizacyjny jest dostępny z samego planu eksperymentu.
\end{przyklad}

\begin{uwaga}[Ostra hipoteza a brak średniego efektu]
Dokładność testu~\ref{thm:fisher-exact} dotyczy $H_0^{\mathrm F}$, a nie automatycznie $H_0^{\mathrm N}$. Przy zróżnicowanych efektach indywidualnych średni efekt może być zerowy, chociaż hipoteza ostra jest fałszywa. Trzeba rozdzielać pytanie „czy żadnej jednostce interwencja nie pomogła ani nie zaszkodziła?” od pytania „czy przeciętny efekt wynosi zero?”.
\end{uwaga}

\section{A/B testing: projekt, moc i pułapki sekwencyjnego oglądania danych}
\begin{przyklad}[Projekt eksperymentu na odsetku konwersji]
Niech konwersja będzie zmienną Bernoulliego, $p_0,p_1\in(0,1)$ oznaczają prawdopodobieństwa odpowiednio w grupie kontrolnej i testowej, a $\Delta=p_1-p_0$ jest efektem na skali bezwzględnej. W modelu dwóch niezależnych prób po $m$ osób, z rozkładów Bernoulliego, $\widehat\Delta=\widehat p_1-\widehat p_0$ ma wariancję $p_1(1-p_1)/m+p_0(1-p_0)/m$. Przybliżony przedział poziomu $1-\alpha$ to $\widehat\Delta\pm z_{1-\alpha/2}\sqrt{\widehat p_1(1-\widehat p_1)/m+\widehat p_0(1-\widehat p_0)/m}$. Dla niewielu konwersji należy rozważyć przedział oparty na wynikach dwóch przedziałów Wilsona lub metodę dokładniejszą, zamiast bezkrytycznie używać Walda.
\end{przyklad}

\begin{stwierdzenie}[Przybliżona liczebność próby]\label{prop:ab-samplesize}
Jeżeli planujemy dwie niezależne próby Bernoulliego po $m$ obserwacji, chcemy wykryć ustaloną różnicę $\Delta\ne0$, zakładamy przybliżenie normalne i oznaczamy $v=p_0(1-p_0)+p_1(1-p_1)$, to orientacyjna liczebność na grupę przy dwustronnym poziomie $\alpha$ i mocy $1-\beta$ wynosi
\[
 m\approx\frac{(z_{1-\alpha/2}+z_{1-\beta})^2\,v}{\Delta^2}.
\]
Wzór jest heurystycznym przybliżeniem planistycznym, a nie dokładną gwarancją przy skończonym $m$.
\end{stwierdzenie}
\begin{proof}
W przybliżeniu $\widehat\Delta\sim\Normal(\Delta,v/m)$. Dla testu dwustronnego odrzucenie następuje w przybliżeniu wtedy, gdy $|\widehat\Delta|$ przekroczy próg równy $z_{1-\alpha/2}$ razy zakładany błąd standardowy. Zignorowanie mniejszego ogona w rozkładzie przesuniętym daje równanie $|\Delta|\sqrt{m/v}\approx z_{1-\alpha/2}+z_{1-\beta}$. Po podniesieniu do kwadratu otrzymujemy wzór. Dokładne planowanie powinno uwzględnić wariancję pod hipotezą zerową, rozkład testowy, dyskretność danych i zamierzony schemat analiz.
\end{proof}

\begin{center}\small
\begin{tabular}{>{\raggedright\arraybackslash}p{.28\textwidth}>{\raggedright\arraybackslash}p{.27\textwidth}>{\raggedright\arraybackslash}p{.29\textwidth}}
\toprule
\textbf{Pytanie} & \textbf{Decyzja projektowa} & \textbf{Ryzyko, które kontrolujemy}\\
\midrule
Kto jest niezależną jednostką? & Losuj osobę albo klaster użytkowników & Interferencja i pseudoreplikacja\\
Jaki jest główny wynik? & Ustal metrykę i okno pomiaru z góry & Wybór „najlepszej” metryki po teście\\
Jaka różnica ma znaczenie? & Ustal minimalny istotny efekt praktyczny & Duże $n$ i błahy efekt\\
Kiedy kończymy zbieranie danych? & Z góry ustal liczność i regułę stopu & Inflacja błędu przy podglądaniu\\
Czy mamy kilka grup/testów? & Zaplanuj kontrolę FWER lub FDR & Fałszywe odkrycia\\
Czy obserwacje są zależne? & Rozważ randomizację klastrową, modele mieszane & Zaniżone błędy standardowe\\
\bottomrule
\end{tabular}
\end{center}

\begin{uwaga}[Podglądanie danych i niedotrzymanie planu]
Klasyczny test o poziomie $\alpha$ zwykle kontroluje błąd przy ustalonym wcześniej terminie analizy. Powtarzane kończenie eksperymentu przy pierwszym $p<\alpha$ podnosi częstość fałszywych alarmów. Jeśli potrzebujemy analiz pośrednich, należy z góry przyjąć plan sekwencyjny (np. granice alfa-spending, e-wartości lub inne procedury z udowodnioną kontrolą błędu). Warto też odróżnić \emph{intention to treat} (porównanie według przydziału) od analizy tylko użytkowników, którzy faktycznie skorzystali z funkcji: takie przejście może zniszczyć randomizację.
\end{uwaga}

\section{Badania obserwacyjne: identyfikacja, propensity score i podwójna odporność}
\begin{definicja}[Warunki identyfikacji efektu średniego]\label{def:causal-id}
Niech $(X,A,Y(0),Y(1))$ będzie wektorem losowym, gdzie $X$ jest mierzalną przed interwencją kowariantą o wartościach w przestrzeni $\mathsf X$, $A\in\{0,1\}$, $\E|Y(a)|<\infty$ dla $a=0,1$ oraz $Y=AY(1)+(1-A)Y(0)$. Zakładamy:
\begin{enumerate}
\item \emph{Wymienność warunkową} $(Y(0),Y(1))\perp A\mid X$.
\item \emph{Dodatniość} $0<e(X)<1$ niemal na pewno, gdzie $e(x)=\Prob(A=1\mid X=x)$ jest wersją warunkowego prawdopodobieństwa (\emph{propensity score}).
\item \emph{Spójność} obserwowany wynik odpowiada rzeczywiście otrzymanej, jednoznacznie określonej interwencji, a brak interferencji zapewnia sens wynikom potencjalnym.
\end{enumerate}
Zakładamy istnienie regularnych rozkładów warunkowych i odpowiednich wersji warunkowych wartości oczekiwanych.
\end{definicja}

\begin{twierdzenie}[Standaryzacja i ważenie odwrotnością prawdopodobieństwa]\label{thm:causal-ipw}
Przy założeniach definicji~\ref{def:causal-id}, jeżeli dodatkowo $\E[|Y|/e(X)\mid A=1]$ i $\E[|Y|/(1-e(X))\mid A=0]$ są skończone (co zachodzi np. przy ograniczonych wynikach i $e(X)$ oddzielonym od $0,1$), to dla $m_a(x)=\E[Y\mid A=a,X=x]$ zachodzi
\[
 \mathrm{ATE}=\E[m_1(X)-m_0(X)]
 =\E\!\left[\frac{AY}{e(X)}-\frac{(1-A)Y}{1-e(X)}\right].
\]
\end{twierdzenie}
\begin{proof}
Wymienność i spójność dają $m_a(X)=\E[Y(a)\mid X]$ niemal na pewno. Wobec tego $\E[m_1(X)-m_0(X)]=\E[Y(1)-Y(0)]$. Dalej, warunkując względem $X$ i korzystając z $\E[A Y\mid X]=e(X)m_1(X)$ oraz $\E[(1-A)Y\mid X]=(1-e(X))m_0(X)$, otrzymujemy równanie z wagami. Dodatkowe warunki całkowalności uprawniają do użycia własności wieży.
\end{proof}

\begin{twierdzenie}[Tożsamość podwójnej odporności]\label{thm:dr-id}
Przy założeniach definicji~\ref{def:causal-id} niech $\widetilde e:\mathsf X\to(0,1)$ i $\widetilde m_0,\widetilde m_1:\mathsf X\to\R$ będą mierzalnymi funkcjami, dla których poniższe składniki są całkowalne. Zdefiniujmy
\[
 \psi=\widetilde m_1(X)-\widetilde m_0(X)
 +\frac{A}{\widetilde e(X)}(Y-\widetilde m_1(X))
 -\frac{1-A}{1-\widetilde e(X)}(Y-\widetilde m_0(X)).
\]
Jeżeli $\widetilde e=e$ niemal na pewno \emph{albo} $\widetilde m_0=m_0$ i $\widetilde m_1=m_1$ niemal na pewno, to $\E\psi=\mathrm{ATE}$.
\end{twierdzenie}
\begin{proof}
Oznaczmy rzeczywiste funkcje przez $e,m_0,m_1$. Po warunkowaniu względem $X$ mamy
\[
 \E[\psi\mid X]=\widetilde m_1-\widetilde m_0
 +\frac e{\widetilde e}(m_1-\widetilde m_1)
 -\frac{1-e}{1-\widetilde e}(m_0-\widetilde m_0),
\]
gdzie wszystkie funkcje po prawej obliczono w $X$. Jeśli $\widetilde e=e$, prawa strona upraszcza się do $m_1-m_0$; jeśli obie $\widetilde m_a=m_a$, składniki resztowe znikają. W obu sytuacjach twierdzenie~\ref{thm:causal-ipw} daje wynik po uśrednieniu.
\end{proof}

\begin{uwaga}[Czego podwójna odporność nie gwarantuje?]
Twierdzenie dotyczy \emph{funkcji populacyjnych}, a nie automatycznej nieobciążoności skończonopróbkowej po podstawieniu funkcji dopasowanych na tych samych danych. Do uzasadnienia asymptotycznego potrzebne są dodatkowe warunki dotyczące estymacji funkcji pomocniczych, np. cross-fitting, ograniczeń wag i zbieżności iloczynów błędów. Żadna metoda nie naprawia niezmierzonego czynnika zakłócającego ani braku dodatniości.
\end{uwaga}

\section{Różnica w różnicach: gdy mamy dane przed i po}
\begin{twierdzenie}[Identyfikacja ATT w modelu dwóch terminów]\label{thm:did-att}
Niech $A\in\{0,1\}$ oznacza przynależność do grupy poddanej interwencji po chwili $0$, a $Y_0(0),Y_1(0),Y_1(1)$ są całkowalnymi wynikami potencjalnymi w chwilach $0$ i $1$; przed interwencją wszystkie jednostki mają wynik $Y_0=Y_0(0)$, po interwencji zaś $Y_1=A Y_1(1)+(1-A)Y_1(0)$. Załóżmy $0<\Prob(A=1)<1$ i \emph{równoległe trendy} bez leczenia:
$\E[Y_1(0)-Y_0(0)\mid A=1]=\E[Y_1(0)-Y_0(0)\mid A=0]$.
Wtedy efekt średni w grupie leczonej $\mathrm{ATT}=\E[Y_1(1)-Y_1(0)\mid A=1]$ spełnia
\[
\mathrm{ATT}=\big(\E[Y_1\mid A=1]-\E[Y_0\mid A=1]\big)
 -\big(\E[Y_1\mid A=0]-\E[Y_0\mid A=0]\big).
\]
\end{twierdzenie}
\begin{proof}
Dla $A=1$ obserwowana zmiana jest sumą $Y_1(1)-Y_1(0)$ i nieobserwowanej zmiany $Y_1(0)-Y_0(0)$. Dla $A=0$ obserwujemy dokładnie tę drugą zmianę. Po odjęciu średnich warunek równoległych trendów usuwa nieobserwowany składnik, pozostawiając ATT.
\end{proof}

\begin{uwaga}[Trendy równoległe nie są testowalne z dwóch chwil]
Zaledwie jeden pomiar przed interwencją nie pozwala zweryfikować kontrfaktycznego trendu po interwencji. Wiele pomiarów przed interwencją umożliwia ocenę zgodności trendów historycznych, ale nadal nie dowodzi założenia dotyczącego nieobserwowanej przyszłości. W analizach panelowych należy dodatkowo uwzględniać autokorelację i klastry; proste błędy standardowe dla niezależnych wierszy są zwykle niewłaściwe.
\end{uwaga}

\section{Laboratorium: od randomizacji do oszacowania efektu}
\begin{lstlisting}[language=Python,caption={Kompletny program: dokładny test randomizacyjny dla małej próby (Python, itertools, NumPy).}]
from itertools import combinations
import numpy as np

outcomes = np.array([4., 5., 7., 8., 10., 11.])
assignment = np.array([0, 0, 0, 1, 1, 1], dtype=int)

def contrast(a, y):
    return y[a == 1].mean() - y[a == 0].mean()

observed = abs(contrast(assignment, outcomes))
values = []
for idx in combinations(range(len(outcomes)), 3):
    a = np.zeros(len(outcomes), dtype=int)
    a[list(idx)] = 1
    values.append(abs(contrast(a, outcomes)))
p_value = np.mean(np.array(values) >= observed - 1e-12)
print("number of assignments:", len(values))
print("observed difference:", contrast(assignment, outcomes))
print("two-sided Fisher randomization p:", p_value)
# Exact for the sharp no-effect null under complete randomization.
\end{lstlisting}

\begin{lstlisting}[language=Python,caption={Kompletny program: dwie grupy A/B, estymacja różnicy konwersji, przedział i kontrola zrównoważenia (Python, NumPy, SciPy).}]
import numpy as np
from scipy import stats

rng = np.random.default_rng(2031)
n = 800
x = rng.normal(0, 1, n)  # pre-treatment engagement score
assignment = np.array([0] * (n//2) + [1] * (n//2))
rng.shuffle(assignment)
logit = np.log(0.08 / 0.92) + 0.65*x + 0.38*assignment
prob = 1 / (1 + np.exp(-logit))
y = rng.binomial(1, prob)
n0, n1 = (assignment == 0).sum(), (assignment == 1).sum()
p0, p1 = y[assignment == 0].mean(), y[assignment == 1].mean()
rd = p1 - p0
se = np.sqrt(p0*(1-p0)/n0 + p1*(1-p1)/n1)
z = stats.norm.ppf(0.975)
print("conversions control, treatment:", p0, p1)
print("risk difference:", rd)
print("approximate 95% CI:", (rd - z*se, rd + z*se))
print("baseline mean imbalance:", x[assignment == 1].mean() -
      x[assignment == 0].mean())
print("unit of randomization: user; analysis: intention to treat")
# For rare conversions compare this approximation with a score-based CI.
\end{lstlisting}

\begin{lstlisting}[language=Julia,caption={Kompletny program: estymator różnicy średnich w doświadczeniu randomizowanym (Julia, Random, Statistics).}]
using Random, Statistics
Random.seed!(2031)
n = 500
Y0 = randn(n)
Y1 = Y0 .+ 0.3 .+ 0.2 .* randn(n)
A = shuffle(vcat(zeros(Int, div(n, 2)), ones(Int, div(n, 2))))
observed = ifelse.(A .== 1, Y1, Y0)
treatment = observed[A .== 1]
control = observed[A .== 0]
estimate = mean(treatment) - mean(control)
se = sqrt(var(treatment) / length(treatment) +
          var(control) / length(control))
println("ATE in this finite population: ", mean(Y1 .- Y0))
println("randomized estimate: ", estimate, "; SE: ", se)
\end{lstlisting}

\begin{zadanie}[Która analiza zachowuje sens randomizacji?]
W eksperymencie A/B wylosowano wariant funkcji dla $1000$ użytkowników, ale tylko część użytkowników rzeczywiście ją otworzyła. Badacz porównuje tych, którzy otworzyli nową wersję, z wszystkimi użytkownikami kontrolnymi, ponieważ uzyskuje niższą wartość p. Czy estymuje efekt przydziału?\par
\emph{Rozwiązanie.} Nie. Otwarcie funkcji jest zdarzeniem \emph{po przydziale} i może zależeć od zachowań, które wpływają na wynik. Selekcja według otwarcia może wprowadzić obciążenie. Pierwotnym, identyfikowalnym dzięki randomizacji celem jest efekt \emph{przydziału} (intention to treat), obliczany po wszystkich jednostkach zgodnie z wylosowanym wariantem. Efekt faktycznego użycia wymaga odrębnych założeń i metod, np. analizy z instrumentem, nie zaś arbitralnego usuwania obserwacji.
\end{zadanie}

\paragraph{Dalsza lektura.}
O rozróżnieniu identyfikacji, estymacji i decyzji w badaniach interwencyjnych oraz obserwacyjnych zob. \href{https://www.hsph.harvard.edu/miguel-hernan/causal-inference-book/}{Hernán i Robins, \emph{Causal Inference: What If}}. Konstrukcje rozdziału opierają się na jawnych założeniach spójności, wymienności i dodatniości, a nie na utożsamieniu korelacji z przyczynowością.


% --- Rozdzial 32: praktyczne laboratorium porownawcze ---
\chapter{Laboratorium statystyczne: wybór metody, analiza i raportowanie}
\label{ch:stat-laboratory}
\paragraph{Powiązania i odsyłacze.} Laboratorium łączy teorię testowania z rozdziału~\ref{ch:testing}, regresję z rozdziałów~\ref{ch:linear-regression}--\ref{ch:glm}, podejście bayesowskie z rozdziału~\ref{ch:bayes-foundations} oraz projektowanie eksperymentów z rozdziału~\ref{ch:causal-ab}.


\section{Przed uruchomieniem testu: co trzeba wiedzieć o danych?}
Ten rozdział nie wprowadza jednego nowego testu. Łączy teorię z poprzednich rozdziałów w kompletny tok postępowania. W praktyce najgroźniejsze błędy powstają \emph{przed} obliczeniem wartości p: niewłaściwy cel, pozorna niezależność jednostek, zmiana hipotezy po analizie, selekcja danych albo zastosowanie poprawnego testu do nieodpowiedniego projektu.

\begin{definicja}[Plan analizy statystycznej]\label{def:stat-analysis-plan}
\emph{Planem analizy} nazywamy ustalony przed obejrzeniem wyników docelowych zestaw reguł określających: populację i jednostkę analizy, parametr będący celem wnioskowania (\emph{estimand}), główny wynik, schemat losowania lub przydziału, główny estymator/test, sposób liczenia niepewności, sposób obchodzenia się z brakami danych, wielokrotnością porównań i analizami pośrednimi. Plan może wskazywać również wcześniej ustalone analizy wrażliwości i kryterium decyzji praktycznej.
\end{definicja}

\begin{center}
\begin{tikzpicture}[node distance=5mm and 6mm,>=Latex,
box/.style={draw=blue!55!black,rounded corners=2pt,align=center,text width=3.15cm,font=\footnotesize,inner sep=5pt},
arrow/.style={->,line width=.65pt}]
\node[box] (goal) {Zdefiniuj cel: efekt, średnia,\\prognoza czy związek?};
\node[box,right=of goal] (design) {Rozpoznaj projekt:\\jednostka i zależność};
\node[box,right=of design] (method) {Wybierz model oraz\\procedurę niepewności};
\node[box,below=of method] (check) {Zbadaj diagnostykę\\i wrażliwość};
\node[box,left=of check] (report) {Raportuj efekt, przedział,\\założenia i ograniczenia};
\draw[arrow] (goal)--(design);
\draw[arrow] (design)--(method);
\draw[arrow] (method)--(check);
\draw[arrow] (check)--(report);
\end{tikzpicture}

{\small\emph{Rysunek 32.1.} Schemat pracy z danymi. Diagnostyka może uzasadnić analizę wrażliwości, lecz nie powinna automatycznie zastępować celu ustalonego w planie.}
\end{center}

\begin{center}\small
\begin{tabular}{>{\raggedright\arraybackslash}p{.30\textwidth}>{\raggedright\arraybackslash}p{.26\textwidth}>{\raggedright\arraybackslash}p{.29\textwidth}}
\toprule
\textbf{Struktura danych i pytanie} & \textbf{Procedura do rozważenia} & \textbf{Krytyczny warunek}\\
\midrule
Dwie niezależne grupy, średnie & Welch; odporny test/studentyzowany bootstrap & Niezależność jednostek, ciężkie ogony\\
Dwa pomiary tej samej osoby & Test t dla różnic; test zmian znaków & Pary są poprawnie skojarzone\\
Kilka grup, jedna zmienna ciągła & ANOVA, Welch, kontrasty & Projekt i wariancje błędów\\
Jedna zmienna i kowarianty & ANCOVA lub regresja & Kowarianty sprzed interwencji; interakcje\\
Wiele powiązanych wyników & MANOVA/MANCOVA & Wspólny cel i macierz kowariancji\\
Wiele wizyt tej samej osoby & Model mieszany; RM--ANOVA w prostym planie & Struktura korelacji i braki danych\\
Sukces/porażka & Model dwumianowy, logit lub test proporcji & Prawidłowy mianownik i poziom jednostki\\
Liczba zdarzeń w różnym czasie & Poisson z offsetem; model ujemny dwumianowy & Ekspozycja i nadmierna dyspersja\\
Szereg czasowy & ARMA, GLS/HAC, walidacja chronologiczna & Stacjonarność i zależność\\
Obserwacje nierandomizowane, efekt działania & Standaryzacja, IPW/AIPW, DiD & Wymienność lub trendy równoległe\\
\bottomrule
\end{tabular}
\end{center}

\section{Studium I: eksperyment A/B od losowania do wniosku}
\begin{przyklad}[Pytanie biznesowe i parametr]
Platforma wylosowała jeden z dwóch wariantów interfejsu dla każdego z $5000$ użytkowników. Interesuje nas, czy nowy wariant zwiększa prawdopodobieństwo dokonania zakupu w ciągu ustalonego okna obserwacji. Jednostką randomizacji i analizy jest \emph{użytkownik}, nie odsłona strony. Głównym parametrem jest różnica prawdopodobieństw konwersji $\Delta=p_1-p_0$, a nie iloraz szans z regresji logistycznej. Przed rozpoczęciem eksperymentu ustalamy też minimalną praktycznie istotną poprawę, np. $\Delta_\star=0{,}01$; sam fakt otrzymania $p<0{,}05$ nie dowodzi przekroczenia tego progu.
\end{przyklad}

Dla dwóch niezależnych prób Bernoulliego o licznościach $n_0,n_1$, przy $0<p_0,p_1<1$, niech $K_a$ oznacza liczbę zakupów i $\widehat p_a=K_a/n_a$. Przy niezależności jednostek oraz rosnących $n_a$, dwuwymiarowe centralne twierdzenie graniczne i twierdzenie Slutsky'ego prowadzą do przybliżenia
\[
\frac{(\widehat p_1-\widehat p_0)-(p_1-p_0)}{\sqrt{\widehat p_1(1-\widehat p_1)/n_1+\widehat p_0(1-\widehat p_0)/n_0}}\ \overset{d}{\longrightarrow}\ \Normal(0,1).
\]
Gdy konwersji jest bardzo mało, gdy $\widehat p_a$ leży na brzegu $\{0,1\}$ albo gdy jednostki są zależne, powyższa normalizacja nie uzasadnia stosowania naiwnego przedziału Walda. Warto wówczas rozważyć dokładniejsze metody dla danych binarnych albo procedury uwzględniające zależność.

\begin{stwierdzenie}[Efekt standardyzowany z regresji logistycznej]\label{prop:logit-std}
Niech dla każdego $x\in\R^q$ i $a\in\{0,1\}$ model warunkowej średniej ma postać $m_a(x;\beta)=g(v_a(x)^\top\beta)$, gdzie $g(u)=(1+e^{-u})^{-1}$, $v_a(x)\in\R^r$ jest zadanym wektorem cech zawierającym m.in. wartość $a$, a $\beta\in\R^r$. Dla ustalonych kowariant $x_1,\ldots,x_n$ definiujemy \emph{standaryzowaną różnicę ryzyka}
$D(\beta)=n^{-1}\sum_i[m_1(x_i;\beta)-m_0(x_i;\beta)]$.
Wtedy gradient ma postać
\[
\nabla D(\beta)=\frac1n\sum_{i=1}^n\left[g'(v_1(x_i)^\top\beta)v_1(x_i)-g'(v_0(x_i)^\top\beta)v_0(x_i)\right],
\]
gdzie $g'(u)=g(u)(1-g(u))$. Jeśli $\widehat\beta$ jest asymptotycznie normalny, a $\widehat\Sigma_\beta$ estymuje jego macierz kowariancji w skali estymatora, to warunkowy, pierwszorzędowy przybliżony błąd standardowy $D(\widehat\beta)$ można obliczyć jako $\sqrt{\nabla D(\widehat\beta)^\top\widehat\Sigma_\beta\nabla D(\widehat\beta)}$.
\end{stwierdzenie}
\begin{proof}
Pochodna funkcji logistycznej wynika bezpośrednio z reguły łańcuchowej. Skończoną sumę można różniczkować wyraz po wyrazie, otrzymując wzór na gradient. Drugie stwierdzenie jest zastosowaniem wielowymiarowej metody delty do różniczkowalnego odwzorowania $D$. Dla losowych kowariant z próby populacyjnej dochodzi również niepewność wynikająca z estymacji rozkładu $X$; nie wolno bez dodatkowych warunków utożsamiać powyższego rachunku z pełnym bezwarunkowym błędem standardowym populacyjnego ATE.
\end{proof}

\begin{uwaga}[Współczynnik logitowy to nie różnica ryzyka]
Regresja logistyczna podaje zwykle $e^{\widehat\beta_A}$, czyli \emph{warunkowy iloraz szans}. Nawet w poprawnym modelu nie jest on tym samym co iloraz ryzyka $p_1/p_0$ ani różnica ryzyka $p_1-p_0$. Dla komunikacji biznesowej różnica ryzyka w punktach procentowych bywa bardziej zrozumiała, ale trzeba zawsze podać także wyjściowy poziom ryzyka i horyzont obserwacji.
\end{uwaga}

\begin{lstlisting}[language=Python,caption={Kompletny program: A/B testing z analizą podstawową, standaryzacją logistyczną i posteriorowym porównaniem odsetków (Python, NumPy, SciPy, statsmodels).}]
import numpy as np
from scipy import stats
from scipy.special import expit
import statsmodels.api as sm

rng = np.random.default_rng(2032)
n = 5000
A = np.repeat([0, 1], n // 2)
rng.shuffle(A)
x = rng.normal(size=n)  # measured before randomization
risk = expit(np.log(.08 / .92) + .55*x + .45*A)
y = rng.binomial(1, risk)

n0, n1 = np.sum(A == 0), np.sum(A == 1)
k0, k1 = y[A == 0].sum(), y[A == 1].sum()
p0, p1 = k0 / n0, k1 / n1
rd = p1 - p0
se = np.sqrt(p0*(1-p0)/n0 + p1*(1-p1)/n1)
ci = (rd - 1.96*se, rd + 1.96*se)
pooled = (k0 + k1) / (n0 + n1)
z = rd / np.sqrt(pooled*(1-pooled)*(1/n0 + 1/n1))
pvalue = 2*stats.norm.sf(abs(z))
print("control, treatment:", p0, p1)
print("risk difference, 95% Wald CI:", rd, ci)
print("two-sided pooled z-test p:", pvalue)
print("pre-treatment baseline imbalance:", x[A == 1].mean() -
      x[A == 0].mean())

# Pre-specified adjusted analysis, with A and baseline x.
X = sm.add_constant(np.column_stack((A, x)))
fit = sm.GLM(y, X, family=sm.families.Binomial()).fit(
    cov_type="HC3")
X0, X1 = X.copy(), X.copy()
X0[:, 1], X1[:, 1] = 0, 1
eta0, eta1 = X0 @ fit.params, X1 @ fit.params
m0, m1 = expit(eta0), expit(eta1)
adjusted = (m1 - m0).mean()
# This Delta-method SE conditions on the observed covariate values.
grad = ((m1*(1-m1))[:, None] * X1 -
        (m0*(1-m0))[:, None] * X0).mean(axis=0)
adjusted_se = np.sqrt(grad @ fit.cov_params() @ grad)
print("conditional-on-X adjusted difference:", adjusted)
print("conditional-on-X delta-method SE:", adjusted_se)
print("adjusted conditional odds ratio:", np.exp(fit.params[1]))

# Independent Beta(1, 1) priors for the two conversion rates.
draw0 = rng.beta(k0 + 1, n0 - k0 + 1, size=100000)
draw1 = rng.beta(k1 + 1, n1 - k1 + 1, size=100000)
difference = draw1 - draw0
print("posterior 95% credible interval:",
      np.quantile(difference, [0.025, 0.975]))
print("posterior P(p1 > p0 | data):", np.mean(difference > 0))
# Randomization targets effect of assignment (intention to treat).
\end{lstlisting}

\begin{przyklad}[Interpretacja wyników z wygenerowanego eksperymentu A/B]
Przy ustalonym ziarnie powyższy program daje $200/2500=0{,}080$ konwersji w grupie kontrolnej i $305/2500=0{,}122$ w grupie testowej. Różnica wynosi $0{,}042$, czyli $4{,}2$ punktu procentowego, a normalny przedział ufności $95\%$ wynosi w przybliżeniu $[0{,}0253,0{,}0587]$. Test dwóch proporcji daje wartość p około $8{,}31\cdot10^{-7}$. Po uwzględnieniu mierzonej przed interwencją cechy $x$ model logistyczny i standaryzacja dają różnicę około $0{,}0415$; iloraz szans około $1{,}61$ jest wielkością na innej skali, nie zastępuje różnicy ryzyka. W tym sztucznym przykładzie cały przedział dla różnicy leży powyżej wcześniej ustalonego progu praktycznego $0{,}01$. Nie oznacza to jednak, że identyczna dokładność wystąpi przy mniejszych próbach, rzadkich zdarzeniach lub zależnych użytkownikach.
\end{przyklad}

\begin{uwaga}[Porównywanie podejść klasycznego i bayesowskiego]
Przedział ufności w pierwszej części programu ma interpretację częstościową w hipotetycznie powtarzanych próbach, natomiast przedział z losowań posterioru opisuje niepewność parametru warunkową względem danych i konkretnego prioru. Nie należy bezpośrednio porównywać wartości p z $\Prob(p_1>p_0\mid\text{dane})$ jako dwóch estymatorów tej samej wielkości. Powyższy rozkład beta przyjmuje niezależne prawdopodobieństwa konwersji; nie jest to model warunkowy z kowariantą $x$.
\end{uwaga}

\section{Studium II: obserwacyjne porównanie i paradoks Simpsona}
\begin{zadanie}[Ta sama tabela, sprzeczny wniosek po agregacji]\label{ex:simpson-decision}
W badaniu obserwacyjnym porównano sukces dwóch metod w dwóch warstwach o różnym bazowym prawdopodobieństwie sukcesu. Zliczenia przedstawia tabela:
\begin{center}\small
\begin{tabular}{lcc}
\toprule
\textbf{Warstwa} & \textbf{Metoda 0: sukcesy / łącznie} & \textbf{Metoda 1: sukcesy / łącznie}\\
\midrule
Niska skuteczność & $1/10$ & $18/90$\\
Wysoka skuteczność & $72/90$ & $9/10$\\
\bottomrule
\end{tabular}
\end{center}
Oblicz różnicę sukcesu bez uwzględnienia warstw oraz różnice wewnątrz warstw. Zastanów się, czy wynik jest przyczynowy.\par
\emph{Rozwiązanie.} Dla metody 0 otrzymujemy łącznie $73/100=0{,}73$, a dla metody 1 $27/100=0{,}27$. Naiwna różnica wynosi $-0{,}46$. Jednak w warstwie o niskiej bazowej skuteczności metoda 1 daje $0{,}20-0{,}10=0{,}10$, a w warstwie o wysokiej bazowej skuteczności $0{,}90-0{,}80=0{,}10$. Po standaryzacji do populacji, w której każda warstwa ma wagę $1/2$, różnica wynosi $0{,}10$. Agregacja zmieniła kierunek porównania, ponieważ metody stosowano u pacjentów z bardzo różnym bazowym prawdopodobieństwem sukcesu. Samo skorygowanie o znane warstwy nie gwarantuje przyczynowości: potrzeba dodatkowo wiarygodnej wymienności warunkowej, dodatniości i spójności interwencji.
\end{zadanie}

\begin{center}
\begin{tikzpicture}[x=1cm,y=1cm,>=Latex]
% Rysunek odtwarza dokładnie odsetki z tabeli (1/10,18/90,72/90,9/10,73/100,27/100).
% Oddzielna legenda nad wykresem, aby próbki kolorów nie zakrywały opisów.
\begin{scope}[every node/.style={font=\scriptsize}]
  \fill[blue!65!black] (1.50,4.56) rectangle +(0.27,0.16);
  \node[anchor=west] at (1.87,4.64) {metoda 0};
  \fill[red!70!black] (3.60,4.56) rectangle +(0.27,0.16);
  \node[anchor=west] at (3.97,4.64) {metoda 1};
\end{scope}
\draw[->] (0.55,0.45)--(0.55,4.28);
\draw[->] (0.55,0.45)--(7.40,0.45);
\node[anchor=west,font=\footnotesize] at (0.55,4.93) {Odsetek sukcesów};
\foreach \v/\lab in {0/0,0.5/50,1/100} {
   \draw (0.48,{0.45+3.5*\v})--(0.62,{0.45+3.5*\v});
   \node[anchor=east,font=\scriptsize] at (0.42,{0.45+3.5*\v}) {\lab\%};
}
\draw[densely dashed,gray!70] (5.15,0.48)--(5.15,4.06);
\foreach \x/\label in {1.70/{niska},3.85/{wysoka},6.25/{łącznie}} {
  \node[font=\footnotesize] at (\x,0.13) {\label};
}
% Dwie serie o znanej wysokości; jawnie podpisane procenty.
\foreach \x/\v/\lab in {1.48/0.10/10,3.63/0.80/80,6.03/0.73/73} {
   \fill[blue!65!black] (\x,0.45) rectangle +(.34,{3.5*\v});
   \node[font=\scriptsize,blue!55!black,anchor=south] at ({\x+.17},{0.49+3.5*\v}) {\lab\%};
}
\foreach \x/\v/\lab in {1.87/0.20/20,4.02/0.90/90,6.42/0.27/27} {
   \fill[red!70!black] (\x,0.45) rectangle +(.34,{3.5*\v});
   \node[font=\scriptsize,red!65!black,anchor=south] at ({\x+.17},{0.49+3.5*\v}) {\lab\%};
}
\end{tikzpicture}

{\small\emph{Rysunek 32.2.} Paradoks Simpsona. Metoda 1 ma wyższy odsetek sukcesów osobno w każdej warstwie, lecz niższy w próbie połączonej z powodu różnych proporcji warstw.}
\end{center}

\section{Studium III: pary, klastry i wielokrotność -- nie zamieniaj jednostek}
\begin{przyklad}[Pomiary tej samej osoby]
Jeżeli $40$ osób otrzymuje po dwa pomiary, testujemy średnią różnicę na \emph{czterdziestu} wartościach $D_i=Y_{i,2}-Y_{i,1}$, a nie na osiemdziesięciu niezależnych pomiarach. Wnioskowanie za pomocą sparowanego testu t wymaga niezależnych osób i odpowiednich warunków normalności różnic przy małej próbie. Test Wilcoxona wymaga innych założeń i dotyczy funkcjonału rozkładu różnic, nie po prostu tej samej średniej. Przy więcej niż dwóch wizytach lub nierównych odstępach lepszy może być model mieszany.
\end{przyklad}

\begin{przyklad}[Wybór post hoc a wielokrotność]
Załóżmy, że porównujemy pięć wariantów interfejsu, z których żaden nie różni się od kontrolnego. Wykonując niezależnie pięć testów o poziomie $0{,}05$, otrzymalibyśmy prawdopodobieństwo co najmniej jednego fałszywego odrzucenia równe $1-0{,}95^5\approx0{,}226$ (przy niezależności testów). Jeśli zależności nie znamy, nierówność Bonferroniego daje uniwersalne ograniczenie $\mathrm{FWER}\le\sum_{j=1}^5\alpha_j$. Dobór korekty zależy od tego, czy kontrolujemy ryzyko choć jednego fałszywego odkrycia, czy oczekiwany odsetek błędnych odkryć.
\end{przyklad}

\begin{lstlisting}[language=Python,caption={Kompletny program: dane sparowane, poprawny test różnic i ilustracja wielokrotności (Python, NumPy, SciPy, statsmodels).}]
import numpy as np
from scipy import stats
from statsmodels.stats.multitest import multipletests

rng = np.random.default_rng(2033)
subjects = 40
baseline = rng.normal(10, 2.5, subjects)
followup = baseline + 0.6 + rng.normal(0, 1.0, subjects)
diff = followup - baseline
paired = stats.ttest_rel(followup, baseline)
wrong = stats.ttest_ind(followup, baseline, equal_var=False)
print("mean change, paired t p:", diff.mean(), paired.pvalue)
print("wrong independent Welch p:", wrong.pvalue)
print("difference variance:", np.var(diff, ddof=1))

pvalues = np.array([.008, .02, .045, .09, .17])
reject, adjusted, _, _ = multipletests(pvalues, .05,
                                      method="holm")
print("raw p:", pvalues)
print("Holm adjusted p:", adjusted)
print("Holm rejection:", reject)
\end{lstlisting}

\section{Raportowanie niepewności i audyt decyzji}
\begin{definicja}[Minimalny raport wyniku]
Dla analizy porównawczej \emph{minimalnym raportem} nazywamy zapis zawierający: definicję populacji i jednostki, liczbę jednostek i braków, definicję zmiennej wynikowej, oszacowanie parametru i jednostkę jego miary, przedział niepewności z określoną interpretacją, nazwę testu lub modelu z założeniami, liczbę wykonanych porównań oraz informację o analizach zaplanowanych i eksploracyjnych. Przedział, wartość p lub posterior nie zastępują jawnego opisu jakości projektu.
\end{definicja}

\begin{center}\small
\begin{tabular}{>{\raggedright\arraybackslash}p{.26\textwidth}>{\raggedright\arraybackslash}p{.55\textwidth}}
\toprule
\textbf{Element decyzji} & \textbf{Pytanie kontrolne}\\
\midrule
Cel & Czy estymujemy efekt interwencji, związek czy predykcję?\\
Projekt & Czy jednostki były losowane / przydzielane niezależnie?\\
Braki danych & Czy mechanizm brakowania może zależeć od nieobserwowanego wyniku?\\
Estymacja & Czy estymator jest zgodny z celem i typem danych?\\
Niepewność & Czy błąd standardowy uwzględnia klastry, dobór i selekcję?\\
Diagnostyka & Czy analizujemy model reszt i wpływ obserwacji, nie tylko p-value Shapiro--Wilka?\\
Wielokrotność & Czy wynik jest pierwotny, czy wybrany spośród wielu prób?\\
Praktyka & Czy oszacowana różnica przekracza próg istotności praktycznej?\\
Replikacja & Czy kod, ziarno losowania, wersje bibliotek i dane są możliwe do odtworzenia?\\
\bottomrule
\end{tabular}
\end{center}

\begin{zadanie}[Sformułuj obronny wniosek]
W eksperymencie A/B uzyskano oszacowaną różnicę konwersji $0{,}006$, asymptotyczny przedział ufności $95\%$ równy $[-0{,}002,0{,}014]$, a biznesowy próg opłacalności to $0{,}01$. Jak sformułować wniosek?\par
\emph{Rozwiązanie.} Dane są zgodne zarówno z niewielką stratą, jak i poprawą przekraczającą próg biznesowy. Nie ma podstaw do stwierdzenia, że udowodniono wzrost konwersji; nie można też uznać braku odrzucenia za dowód nieistnienia efektu. Jeśli celem jest wykazanie poprawy co najmniej $0{,}01$, należy zaplanować jednostronny test lub przedział dla takiego progu z odpowiednią mocą i regułą decyzji. Decyzja może dodatkowo wymagać analizy kosztów, ryzyka i kolejnych danych.
\end{zadanie}

\paragraph{Podsumowanie praktyczne.}
\emph{Nie istnieje jeden ``najlepszy test statystyczny'' niezależny od pytania i projektu.} Nawet matematycznie dokładny test Shapiro--Wilka, Wilksa lub t-Studenta odpowiada na określoną hipotezę i wymaga określonych warunków. Zadaniem analityka jest najpierw zdefiniować problem, a potem dopasować do niego model, estymator i kontrolę niepewności. Statystyka nie zastępuje projektu badania; dobrze wykonany projekt czyni późniejszą statystykę wiarygodną.

\end{document}
